% Generated mechanically from frozen input inputs/p05/ja/more-morphisms.tex.
% Do not edit this generated file; edit the bound locale source instead.

% OK, start here.
%

\setcounter{chapter}{36}
\chapter{射についてさらに}



\phantomsection
\label{more-morphisms-section-phantom}


\section{序論}
\label{more-morphisms-section-introduction}

\noindent
本章では、スキームの射の性質に関する研究を続ける。
基本的な参考文献は \cite{EGA} である。






\section{肥厚化}
\label{more-morphisms-section-thickenings}

\noindent
以下の用語法は必ずしも完全に標準的ではないが、便利である。

\begin{definition}
\label{more-morphisms-definition-thickening}
肥厚化について次のように定める。
\begin{enumerate}
\item スキーム $X'$ がスキーム $X$ の {\it 肥厚化} であるとは、
$X$ が $X'$ の閉部分スキームであり、両者の基礎位相空間が
等しいことをいう。
\item スキーム $X'$ がスキーム $X$ の {\it 一次肥厚化} であるとは、
$X$ が $X'$ の閉部分スキームであり、$X$ を定める準連接イデアル層
$\mathcal{I} \subset \mathcal{O}_{X'}$ の平方が零であることをいう。
\item スキーム $X'$ がスキーム $X$ の {\it 有限次肥厚化} であるとは、
$X$ が $X'$ の閉部分スキームであり、$X$ を定める準連接イデアル層
$\mathcal{I} \subset \mathcal{O}_{X'}$ が冪零、すなわち
$\mathcal{I}^{n + 1} = 0$ を満たす整数 $n \geq 0$ が存在することをいう。
% P05-EN-CH37-0001
\item スキーム $X'$ がスキーム $X$ の {\it $n$ 次肥厚化} であるとは、
$X$ が $X$ の閉部分スキームであり、$X$ を定める準連接イデアル層
$\mathcal{I} \subset \mathcal{O}_{X'}$ に対して
$\mathcal{I}^{n + 1} = 0$ が成り立つことをいう。
\item 二つの肥厚化 $X \subset X'$ および $Y \subset Y'$ が与えられたとき、
{\it 肥厚化の射} とは、$f'(X) \subset Y$ を満たす射
$f' : X' \to Y'$、すなわち $f'|_X$ が閉部分スキーム $Y$ を経由する
ような射をいう。この状況では $f = f'|_X : X \to Y$ とおき、
$(f, f') : (X \subset X') \to (Y \subset Y')$ を肥厚化の射と呼ぶ。
\item $S$ をスキームとする。同様に {\it $S$ 上の肥厚化} および
{\it $S$ 上の肥厚化の射} を定める。これは、上のスキーム
$X, X', Y, Y'$ が $S$ 上のスキームであり、射
$X \to X'$、$Y \to Y'$ および $f' : X' \to Y'$ が
$S$ 上の射であることを意味する。
\end{enumerate}
\end{definition}

\noindent
$i_X : X \to X'$ を肥厚化とする。核
$\mathcal{I} = \Ker(i_X^\sharp)$ の任意の局所切断は局所的に冪零である。
これにより $\mathcal{O}_{X'}$ の構造をある程度制御できるが、技術的には
有限次肥厚化 $X \subset X'$ の類の方がはるかに扱いやすい。
% P05-EN-CH37-0002
実際、$X'$ が $n$ 次肥厚化なら、フィルトレーション
$$
0 = \mathcal{I}^{n + 1} \subset
\mathcal{I}^n \subset
\mathcal{I}^{n - 1} \subset \ldots \subset
\mathcal{I} \subset \mathcal{O}_{X'}
$$
があり、$X'$ は閉部分空間の列
$$
X = X_1 \subset X_2 \subset \ldots \subset X_n \subset X_{n + 1} = X'
$$
でフィルター付けされ、各対 $X_i \subset X_{i + 1}$ は $S$ 上の
一次肥厚化となる。簡単な帰納法を用いれば、一次肥厚化について証明された
多くの結果を有限次肥厚化についての結果として言い換えることができる。

\medskip\noindent
% P05-EN-CH37-0003
一次肥厚化は次のように記述される
（Modules, Lemma \ref{modules-lemma-double-structure-gives-derivation} を参照）。

\begin{lemma}
\label{more-morphisms-lemma-first-order-thickening}
% P05-EN-CH37-0004
$X$ を基礎 $S$ 上のスキームとする。$X$ 上の層の短完全列
$$
0 \to \mathcal{I} \to \mathcal{A} \to \mathcal{O}_X \to 0
$$
を考える。ここで $\mathcal{A}$ は $f^{-1}\mathcal{O}_S$-代数の層、
$\mathcal{A} \to \mathcal{O}_X$ は
$f^{-1}\mathcal{O}_S$-代数の層の全射であり、$\mathcal{I}$ はその核である。
次を仮定する。
\begin{enumerate}
\item $\mathcal{I}$ は $\mathcal{A}$ の平方零イデアルである。
\item $\mathcal{I}$ は $\mathcal{O}_X$-加群として準連接である。
\end{enumerate}
このとき $X' = (X, \mathcal{A})$ はスキームであり、$X \to X'$ は
$S$ 上の一次肥厚化である。さらに、$S$ 上の任意の一次肥厚化は
この形をしている。
\end{lemma}

\begin{proof}
$X'$ が局所環付き空間であることは明らかである。$U = \Spec(B)$ を
$X$ のアフィン開集合とし、$A = \Gamma(U, \mathcal{A})$ とおく。
$H^1(U, \mathcal{I}) = 0$ であるから
（Cohomology of Schemes, Lemma
\ref{coherent-lemma-quasi-coherent-affine-cohomology-zero})
$A \to B$ は全射である。仮定により、その核
$I = \mathcal{I}(U)$ は環 $A$ の平方零イデアルである。
Schemes, Lemma \ref{schemes-lemma-morphism-into-affine}
により、局所環付き空間の標準的な射
$$
(U, \mathcal{A}|_U) \longrightarrow \Spec(A)
$$
% P05-EN-CH37-0005
が写像 $B \to \Gamma(U, \mathcal{A})$ から得られる。この射は可換図式
$$
\xymatrix{
(U, \mathcal{O}_X|_U) \ar[d] \ar[r] & \Spec(B) \ar[d] \\
(U, \mathcal{A}|_U) \ar[r] & \Spec(A)
}
$$
に入るので、基礎位相空間上の同相写像であることが分かる。
したがって、これが同型であることを示すには、局所環上に同型を
誘導することを確かめれば十分である。
素イデアル $\mathfrak p \subset A$ に対応する $u \in U$ に対し、
短完全列の可換図式
$$
\xymatrix{
0 \ar[r] &
I_{\mathfrak p} \ar[r] \ar[d] &
A_{\mathfrak p} \ar[r] \ar[d] &
B_{\mathfrak p} \ar[r] \ar[d] & 0 \\
0 \ar[r] &
\mathcal{I}_u \ar[r] &
\mathcal{A}_u \ar[r] &
\mathcal{O}_{X, u} \ar[r] & 0.
}
$$
を得る。$\mathcal{I}$ と $\mathcal{O}_X$ は準連接層なので、
左右の縦矢印は同型である。したがって中央の写像も同型である。
ゆえに $X' = (X, \mathcal{A})$ の各点はアフィン近傍をもち、
$X'$ は所望のスキームである。
\end{proof}

\begin{lemma}
\label{more-morphisms-lemma-thickening-affine-scheme}
\begin{slogan}
アフィン性は肥厚化によって変わらない
\end{slogan}
\begin{reference}
有限次肥厚化の場合は \cite[Proposition 5.1.9]{EGA1} である。
\end{reference}
アフィンスキームの任意の肥厚化はアフィンである。
\end{lemma}

\begin{proof}
これは Limits, Proposition \ref{limits-proposition-affine}
の特別な場合である。
\end{proof}

\begin{proof}[有限次肥厚化に対する証明]
$X \subset X'$ を、$X$ がアフィンである有限次肥厚化とする。
このとき Serre の判定法を用いて $X'$ がアフィンであることを示せる。
より正確には、Cohomology of Schemes, Lemma
\ref{coherent-lemma-quasi-compact-h1-zero-covering} を用いる。
$\mathcal{F}$ を準連接 $\mathcal{O}_{X'}$-加群とする。
$H^1(X', \mathcal{F}) = 0$ を示せば十分である。与えられた閉埋め込みを
$i : X \to X'$ と書き、
$\mathcal{I} = \Ker(i^\sharp : \mathcal{O}_{X'} \to i_*\mathcal{O}_X)$
とおく。有限次肥厚化についての先の議論
（Definition \ref{more-morphisms-definition-thickening} の直後）により、
ある $n \geq 0$ と、準連接部分加群によるフィルトレーション
$$
0 = \mathcal{F}_{n + 1} \subset \mathcal{F}_n \subset
\mathcal{F}_{n - 1} \subset \ldots \subset
\mathcal{F}_0 = \mathcal{F}
$$
が存在し、$\mathcal{F}_a/\mathcal{F}_{a + 1}$ は
$\mathcal{I}$ によって零化される。実際、
$\mathcal{F}_a = \mathcal{I}^a\mathcal{F}$ ととれる。このとき、ある準連接
$\mathcal{O}_X$-加群 $\mathcal{G}_a$ に対して
$\mathcal{F}_a/\mathcal{F}_{a + 1} = i_*\mathcal{G}_a$ となる。
Morphisms, Lemma \ref{morphisms-lemma-i-star-equivalence} を参照せよ。
したがって
$$
H^1(X', \mathcal{F}_a/\mathcal{F}_{a + 1}) =
H^1(X', i_*\mathcal{G}_a) = H^1(X, \mathcal{G}_a) = 0
$$
を得る。第二の等号は Cohomology of Schemes, Lemma
\ref{coherent-lemma-relative-affine-cohomology}
から、最後の等号は Cohomology of Schemes, Lemma
\ref{coherent-lemma-quasi-coherent-affine-cohomology-zero} から従う。
ゆえに $\mathcal{F}$ は、各逐次商の第一コホモロジーが消える
有限フィルトレーションをもつので、簡単な帰納法により
$H^1(X', \mathcal{F}) = 0$ が従う。
\end{proof}

\begin{lemma}
\label{more-morphisms-lemma-base-change-thickening}
$S \subset S'$ をスキームの肥厚化とする。射 $X' \to S'$ をとり、
$X = S \times_{S'} X'$ とおく。このとき
$(X \subset X') \to (S \subset S')$ は肥厚化の射である。
$S \subset S'$ が一次（それぞれ有限次）肥厚化ならば、
$X \subset X'$ も一次（それぞれ有限次）肥厚化である。
\end{lemma}

\begin{proof}
省略する。
\end{proof}

\begin{lemma}
\label{more-morphisms-lemma-composition-thickening}
\begin{slogan}
肥厚化の合成は肥厚化である
\end{slogan}
$S \subset S'$ および $S' \subset S''$ が肥厚化ならば、
$S \subset S''$ も肥厚化である。
\end{lemma}

\begin{proof}
省略する。
\end{proof}

\begin{lemma}
\label{more-morphisms-lemma-descending-property-thickening}
肥厚化であるという性質は fpqc 局所的である。
一次肥厚化についても同様である。
\end{lemma}

\begin{proof}
主張の意味は次の通りである。$X \to X'$ をスキームの射とし、
$\{g_i : X'_i \to X'\}$ を fpqc 被覆で、任意の $i$ に対して
基底変換 $X_i \to X'_i$ が肥厚化であるものとする。
このとき $X \to X'$ は肥厚化である。射 $g_i$ は共同全射なので、
$X \to X'$ は全射である。さらに
Descent, Lemma \ref{descent-lemma-descending-property-closed-immersion}
により $X \to X'$ は閉埋め込みである。したがって
$X \to X'$ は肥厚化である。一次肥厚化の場合の証明は省略する。
\end{proof}









\section{肥厚化の射}
\label{more-morphisms-section-morphisms-thickenings}

\noindent
$(f, f') : (X \subset X') \to (Y \subset Y')$ がスキームの肥厚化の射ならば、
射 $f$ の性質が $f'$ に受け継がれることが多い。いくつかの型がある。

\begin{lemma}
\label{more-morphisms-lemma-thicken-property-morphisms}
$(f, f') : (X \subset X') \to (S \subset S')$ を肥厚化の射とする。このとき
\begin{enumerate}
\item $f$ がアフィン射であることと $f'$ がアフィン射であることは同値である。
\item $f$ が全射であることと $f'$ が全射であることは同値である。
% P05-EN-CH37-0006
\item $f$ が準コンパクトであることと $f'$ が準コンパクトであることは同値である。
\item $f$ が普遍閉であることと $f'$ が普遍閉であることは同値である。
\item $f$ が整であることと $f'$ が整であることは同値である。
\item $f$ が（準）分離的であることと $f'$ が（準）分離的であることは同値である。
\item $f$ が普遍単射であることと $f'$ が普遍単射であることは同値である。
\item $f$ が普遍開であることと $f'$ が普遍開であることは同値である。
\item $f$ が準アフィンであることと $f'$ が準アフィンであることは同値である。
% P05-EN-CH37-0007
\item ここにさらに追加する。
\end{enumerate}
\end{lemma}

\begin{proof}
$S \to S'$ と $X \to X'$ は普遍同相写像であることに注意する
（例えば Morphisms, Lemma
\ref{morphisms-lemma-reduction-universal-homeomorphism} を参照）。
これから直ちに (2), (3), (4), (7), (8) が従う。
(1) は Lemma \ref{more-morphisms-lemma-thickening-affine-scheme} から従う。実際この補題は、
$S$ と $S'$ のアフィン開集合の間、および $X$ と $X'$ のアフィン開集合の間に
一対一対応があることを述べている。(9) は
Limits, Lemma \ref{limits-lemma-thickening-quasi-affine}
と、いま述べた $S$ と $S'$ のアフィン開集合に関する注意から従う。
(5) は (1), (4) と Morphisms, Lemma
\ref{morphisms-lemma-integral-universally-closed} から従う。
最後に
% P05-EN-CH37-0008
$$
S \times_X S = S \times_{X'} S \to S \times_{X'} S' \to S' \times_{X'} S'
$$
が肥厚化であることに注意する（二つの矢印は Lemma
\ref{more-morphisms-lemma-base-change-thickening} により肥厚化である）。
したがって、肥厚化の射
$(S \subset S') \to (S \times_X S \to S' \times_{X'} S')$
に (3), (4) を適用すると (6) を得る。
\end{proof}

\begin{lemma}
\label{more-morphisms-lemma-thicken-property-relatively-ample}
$(f, f') : (X \subset X') \to (S \subset S')$ を肥厚化の射とする。
$\mathcal{L}'$ を $X'$ 上の可逆層とし、その $X$ への制限を
$\mathcal{L}$ と書く。このとき、$\mathcal{L}'$ が $f'$-豊富であることと
$\mathcal{L}$ が $f$-豊富であることは同値である。
\end{lemma}

\begin{proof}
相対的豊富性は基礎の各アフィン開集合上の条件であることを思い出そう。
Morphisms, Definition \ref{morphisms-definition-relatively-ample} を参照せよ。
Lemma \ref{more-morphisms-lemma-thickening-affine-scheme} により、$S$ と $S'$ の
アフィン開集合の間には一対一対応がある。したがって $S$ と $S'$ は
アフィンであると仮定してよく、$\mathcal{L}'$ が豊富であることと
$\mathcal{L}$ が豊富であることが同値であると示すことに帰着する。
これは Limits, Lemma \ref{limits-lemma-ample-on-reduction} である。
\end{proof}

\begin{lemma}
\label{more-morphisms-lemma-thicken-property-morphisms-cartesian}
$(f, f') : (X \subset X') \to (S \subset S')$ を
$X = S \times_{S'} X'$ を満たす肥厚化の射とする。
$S \subset S'$ が有限次肥厚化ならば、次が成り立つ。
\begin{enumerate}
\item $f$ が閉埋め込みであることと $f'$ が閉埋め込みであることは同値である。
\item $f$ が局所有限型であることと $f'$ が局所有限型であることは同値である。
\item $f$ が局所準有限であることと $f'$ が局所準有限であることは同値である。
\item $f$ が相対次元 $d$ の局所有限型であることと、$f'$ が
相対次元 $d$ の局所有限型であることは同値である。
\item $\Omega_{X/S} = 0$ であることと $\Omega_{X'/S'} = 0$ であることは
同値である。
\item $f$ が不分岐であることと $f'$ が不分岐であることは同値である。
\item $f$ が固有であることと $f'$ が固有であることは同値である。
\item $f$ が有限であることと $f'$ が有限であることは同値である。
\item $f$ がモノ射であることと $f'$ がモノ射であることは同値である。
\item $f$ が埋め込みであることと $f'$ が埋め込みであることは同値である。
\item ここにさらに追加する。
\end{enumerate}
\end{lemma}

\begin{proof}
補題に挙げた性質 $\mathcal{P}$ はすべて基底変換で安定である。
したがって $f'$ が性質 $\mathcal{P}$ をもてば $f$ もこれをもつ。
Schemes, Lemmas \ref{schemes-lemma-base-change-immersion} および
\ref{schemes-lemma-base-change-monomorphism}
および
Morphisms, Lemmas
\ref{morphisms-lemma-base-change-finite-type},
\ref{morphisms-lemma-base-change-quasi-finite},
\ref{morphisms-lemma-base-change-relative-dimension-d},
\ref{morphisms-lemma-base-change-differentials},
\ref{morphisms-lemma-base-change-unramified},
\ref{morphisms-lemma-base-change-proper}, および
\ref{morphisms-lemma-base-change-finite}
を参照せよ。

\medskip\noindent
したがって各場合に本質的なのは、$f$ がその性質をもつと仮定し、
$f'$ もそれをもつと導く向きである。肥厚化の次数に関する帰納法により、
$S \subset S'$ は一次肥厚化であると仮定してよい。
Definition \ref{more-morphisms-definition-thickening} 直後の議論を参照せよ。

\medskip\noindent
証明の大部分ではアフィンの場合への帰着を用いる。
% P05-EN-CH37-0009
$U' \subset S'$ をアフィン開集合、$V' \subset X'$ を $U'$ の上にある
アフィン開集合とする。$U' = \Spec(A')$ とし、閉部分スキーム
$U' \cap S$ を定めるイデアルを $I \subset A'$ と書く。
$V' = \Spec(B')$ とする。このとき $V' \cap X = \Spec(B'/IB')$ である。
$A = A'/I$ および $B = B'/IB'$ とおくと、完全な行をもつ可換図式
$$
\xymatrix{
0 \ar[r] &
IB' \ar[r] &
B' \ar[r] &
B \ar[r] & 0 \\
0 \ar[r] &
IA' \ar[r] \ar[u] &
A' \ar[r] \ar[u] &
A \ar[r] \ar[u] & 0
}
$$
を得て、$I^2 = 0$ である。

\medskip\noindent
(1) を代数に翻訳すると次のようになる。$A \to B$ が全射ならば
$A' \to B'$ も全射である。これは Nakayama の補題
（Algebra, Lemma \ref{algebra-lemma-NAK}）から従う。

\medskip\noindent
(2) を代数に翻訳すると次のようになる。$A \to B$ が有限型の環準同型ならば、
$A' \to B'$ も有限型の環準同型である。これは
$A[x_1, \ldots, x_n] \to B$ が全射となるような写像
$A'[x_1, \ldots, x_n] \to B'$ に Nakayama の補題
（Algebra, Lemma \ref{algebra-lemma-NAK}）を適用すれば従う。

\medskip\noindent
(3) の証明。(2) と、局所有限型射の準有限性はファイバー上で
判定できることから従う。Morphisms, Lemma
\ref{morphisms-lemma-quasi-finite-at-point-characterize} を参照せよ。

\medskip\noindent
% P05-EN-CH37-0010
(4) の証明。(2) と、「相対次元 $d$ である」という付加的性質が
ファイバー上で判定できることから従う
（定義により、Morphisms, Definition
\ref{morphisms-definition-relative-dimension-d} を参照）。

\medskip\noindent
(5) を代数に翻訳すると次のようになる。$\Omega_{B/A} = 0$ ならば
$\Omega_{B'/A'} = 0$ である。
Algebra, Lemma \ref{algebra-lemma-differentials-base-change}
により $0 = \Omega_{B/A} = \Omega_{B'/A'}/I\Omega_{B'/A'}$ である。
したがって Nakayama の補題（Algebra, Lemma
\ref{algebra-lemma-NAK}）により $\Omega_{B'/A'} = 0$ となる。

\medskip\noindent
% P05-EN-CH37-0011
(6) を代数に翻訳すると次のようになる。$A \to B$ が不分岐ならば
$A' \to B'$ も不分岐である。$A \to B$ は有限型なので、上の (2) により
$A' \to B'$ も有限型である。$A \to B$ は不分岐だから
$\Omega_{B/A} = 0$ であり、(5) により $\Omega_{B'/A'} = 0$ である。
したがって $A' \to B'$ は不分岐である。

\medskip\noindent
(7) の証明。(2)、Lemma \ref{more-morphisms-lemma-thicken-property-morphisms} の結果、
および「固有」が
準コンパクト $+$ 分離的 $+$ 局所有限型 $+$ 普遍閉
と同値であることを組み合わせれば従う。

\medskip\noindent
(8) の証明。(2)、Lemma \ref{more-morphisms-lemma-thicken-property-morphisms} の結果、
および「有限」が 整 $+$ 局所有限型 と同値であることを用いれば従う
（Morphisms, Lemma \ref{morphisms-lemma-finite-integral}）。

\medskip\noindent
(9) の証明。$f$ はモノ射なので $X = X \times_S X$ である。
これまでに示した結果を肥厚化の射
$(X \subset X') \to (X \times_S X \subset X' \times_{S'} X')$
に適用できる。(1) により $X' \to X' \times_{S'} X'$ は閉埋め込みである。
実際、これは一次肥厚化である。というのも、閉埋め込み
$X' \to X' \times_{S'} X'$ を定めるイデアルは、$S'$ 内で $S$ を
切り出すイデアル $\mathcal{I} \subset \mathcal{O}_{S'}$ の引き戻しに
含まれるからである。実際、
$X = X \times_S X = (X' \times_{S'} X') \times_{S'} S$ は
$X'$ に含まれる。したがって
Morphisms, Lemma \ref{morphisms-lemma-differentials-diagonal}
により、$\Omega_{X'/S'} = 0$ を示せば十分である。これは (5) と
$X/S$ に対する対応する主張から従う。

\medskip\noindent
(10) の証明。$f : X \to S$ が埋め込みならば、
$X \to U \to S$ と分解する。ここで $U \to S$ は開埋め込み、
$X \to U$ は閉埋め込みである。$U' \subset S'$ を、その基礎位相空間が
$U$ と同じである開部分スキームとする。このとき $X' \to S'$ は
$U'$ を経由し、(1) により $X' \to U'$ は閉埋め込みである。
これで証明が完了する。
\end{proof}

\noindent
次の補題は直前の補題の変種である。関係する肥厚化が有限次であると
仮定する代わりに（この仮定なら局所有限型という性質を $f$ から
$f'$ へ移せる）、ここでは $f$ と $f'$ がともに局所有限型であることを
仮定する。

\begin{lemma}
\label{more-morphisms-lemma-properties-that-extend-over-thickenings}
% P05-EN-CH37-0012
$(f, f') : (X \subset X') \to (Y \to Y')$ を肥厚化の射とする。
$f$ と $f'$ は局所有限型であり、$X = Y \times_{Y'} X'$ であると仮定する。
このとき次が成り立つ。
\begin{enumerate}
\item $f$ が局所準有限であることと $f'$ が局所準有限であることは同値である。
\item $f$ が有限であることと $f'$ が有限であることは同値である。
\item $f$ が閉埋め込みであることと $f'$ が閉埋め込みであることは同値である。
\item $\Omega_{X/Y} = 0$ であることと $\Omega_{X'/Y'} = 0$ であることは
同値である。
\item $f$ が不分岐であることと $f'$ が不分岐であることは同値である。
\item $f$ がモノ射であることと $f'$ がモノ射であることは同値である。
\item $f$ が埋め込みであることと $f'$ が埋め込みであることは同値である。
\item $f$ が固有であることと $f'$ が固有であることは同値である。
% P05-EN-CH37-0013
\item ここにさらに追加する。
\end{enumerate}
\end{lemma}

\begin{proof}
補題に挙げた性質 $\mathcal{P}$ はすべて基底変換で安定であるから、
$f'$ が性質 $\mathcal{P}$ をもてば $f$ もこれをもつ。
Schemes, Lemmas \ref{schemes-lemma-base-change-immersion} および
\ref{schemes-lemma-base-change-monomorphism}
ならびに
Morphisms, Lemmas
\ref{morphisms-lemma-base-change-quasi-finite},
\ref{morphisms-lemma-base-change-relative-dimension-d},
\ref{morphisms-lemma-base-change-differentials},
\ref{morphisms-lemma-base-change-unramified},
\ref{morphisms-lemma-base-change-proper}, および
\ref{morphisms-lemma-base-change-finite}
を参照せよ。したがって各場合、$f$ が所望の性質をもてば $f'$ も
それをもつことだけを示せばよい。

\medskip\noindent
射が局所準有限であることと、それが局所有限型でありスキーム論的
ファイバーが離散空間であることは同値である。Morphisms, Lemma
\ref{morphisms-lemma-locally-quasi-finite-fibres} を参照せよ。
肥厚化へ移っても基礎位相空間は変わらないので、$f'$ が局所準有限である
ことと $f$ が局所準有限であることは同値である。これで (1) が示された。

\medskip\noindent
(2) は Lemma \ref{more-morphisms-lemma-thicken-property-morphisms} の (5) と、
有限射が局所有限型の整射にほかならないことから従う
（Morphisms, Lemma \ref{morphisms-lemma-finite-integral}）。

\medskip\noindent
(3)。これは Morphisms, Lemma
\ref{morphisms-lemma-check-closed-infinitesimally}
から直ちに従う。

\medskip\noindent
(4) は次の代数的主張から従う。$A$ を環、$I \subset A$ を
局所的に冪零なイデアル、$B$ を有限型 $A$-代数とする。
$\Omega_{(B/IB)/(A/I)} = 0$ ならば $\Omega_{B/A} = 0$ である。
% P05-EN-CH37-0014
実際、仮定は $I\Omega_{B/A} = 0$ を意味する。
Algebra, Lemma \ref{algebra-lemma-differentials-base-change}
を参照せよ。一方、$\Omega_{B/A}$ は有限 $B$-加群である。
Algebra, Lemma \ref{algebra-lemma-differentials-finitely-generated}
を参照せよ。したがって $\Omega_{B/A}$ の消滅は Nakayama の補題
（Algebra, Lemma \ref{algebra-lemma-NAK}）と、$IB$ が $B$ の
Jacobson 根基に含まれることから従う。

\medskip\noindent
(5) は (4) と Morphisms, Lemma
\ref{morphisms-lemma-unramified-omega-zero} から直ちに従う。

\medskip\noindent
(6) の証明。$f$ はモノ射なので $X = X \times_Y X$ である。
これまでに示した結果を肥厚化の射
$(X \subset X') \to (X \times_Y X \subset X' \times_{Y'} X')$
に適用する。(3) により
$\Delta_{X'/Y'} : X' \to X' \times_{Y'} X'$ は閉埋め込みである。
実際 $\Delta_{X'/Y'}$ は基礎集合上の全単射なので、
$\Delta_{X'/Y'}$ は肥厚化である。
一方、$\Delta_{X'/Y'}$ は Morphisms, Lemma
\ref{morphisms-lemma-diagonal-morphism-finite-type}
により局所有限表示である。言い換えると、
$\Delta_{X'/Y'}(X')$ は有限型の準連接イデアル層
$\mathcal{J} \subset \mathcal{O}_{X' \times_{Y'} X'}$ によって
切り出される。(5) により $\Omega_{X'/Y'} = 0$ なので、
$X' \to X' \times_{Y'} X'$ の余法層は Morphisms, Lemma
\ref{morphisms-lemma-differentials-diagonal}
により零である。すなわち $\mathcal{J}/\mathcal{J}^2 = 0$ である。
これから $\Delta_{X'/Y'}$ は同型であることが従う。例えば
Algebra, Lemma \ref{algebra-lemma-ideal-is-squared-union-connected}
を用いればよい。

\medskip\noindent
(7) の証明。$f : X \to Y$ が埋め込みならば、
$X \to V \to Y$ と分解する。ここで $V \to Y$ は開埋め込み、
$X \to V$ は閉埋め込みである。$V' \subset Y'$ を、その基礎位相空間が
$V$ と同じである開部分スキームとする。
% P05-EN-CH37-0015
このとき $X' \to V'$ は $V'$ を経由し、(3) により
$X' \to V'$ は閉埋め込みである。

\medskip\noindent
(8) は Lemma \ref{more-morphisms-lemma-thicken-property-morphisms} と、
固有射とは準コンパクト、普遍閉かつ分離的な局所有限型射であるという
定義から従う。
\end{proof}








\section{肥厚化の Picard 群}
\label{more-morphisms-section-picard-group-thickening}

\noindent
肥厚化の Picard 群に関するいくつかの事実を述べる。

\begin{lemma}
\label{more-morphisms-lemma-picard-group-first-order-thickening}
$X \subset X'$ をイデアル層 $\mathcal{I}$ をもつ一次肥厚化とする。
このときアーベル群の標準的完全列
$$
\xymatrix{
0 \ar[r] &
H^0(X, \mathcal{I}) \ar[r] &
H^0(X', \mathcal{O}_{X'}^*) \ar[r] &
H^0(X, \mathcal{O}^*_X) \ar `r[d] `d[l] `l[llld] `d[dll] [dll] \\
& H^1(X, \mathcal{I}) \ar[r] &
\Pic(X') \ar[r] &
\Pic(X) \ar `r[d] `d[l] `l[llld] `d[dll] [dll] \\
& H^2(X, \mathcal{I}) \ar[r] & \ldots \ar[r] & \ldots
}
$$
がある。
\end{lemma}

\begin{proof}
これはアーベル群の層の短完全列
$$
0 \to \mathcal{I} \to \mathcal{O}_{X'}^* \to \mathcal{O}_X^* \to 0
$$
に付随する長完全コホモロジー列である。ここで最初の写像は、
$\mathcal{I}$ の局所切断 $f$ を $\mathcal{O}_{X'}$ の可逆切断
$1 + f$ に送る。また、環付き空間の Picard 群と可逆関数の層の
第一コホモロジー群との同一視を用いる。Cohomology, Lemma
\ref{cohomology-lemma-h1-invertible} を参照せよ。
\end{proof}

\begin{lemma}
\label{more-morphisms-lemma-torsion-pic-thickening}
$X \subset X'$ を肥厚化とし、$n$ を $\mathcal{O}_X$ において
可逆な整数とする。このとき写像
$\Pic(X')[n] \to \Pic(X)[n]$ は全単射である。
\end{lemma}

\begin{proof}[有限次肥厚化に対する証明]
Definition \ref{more-morphisms-definition-thickening} の後で説明した一般原理により、
一次肥厚化の場合に帰着する。
% P05-EN-CH37-0016
Lemma \ref{more-morphisms-lemma-picard-group-first-order-thickening} を用いると、
$H^1(X, \mathcal{I})[n]$, $H^1(X, \mathcal{I})/n$, および
$H^2(X, \mathcal{I})[n]$ が零であることを示せば十分だと分かる。
これは $\mathcal{I}$ 上の $n$ 倍写像が
$\mathcal{O}_X$-加群としての同型であることから従う。
\end{proof}

\begin{proof}[一般の場合の証明]
$\mathcal{I} \subset \mathcal{O}_{X'}$ を $X$ を切り出す
準連接イデアル層とする。このときアーベル群の短完全列
$$
0 \to (1 + \mathcal{I})^* \to \mathcal{O}_{X'}^* \to \mathcal{O}_X^* \to 0
$$
がある。Lemma \ref{more-morphisms-lemma-picard-group-first-order-thickening} の主張において
$H^i(X, \mathcal{I})$ を $H^i(X, (1 + \mathcal{I})^*)$ で置き換えた
長完全コホモロジー列を得る。したがって $n$ 乗写像が
同型 $(1 + \mathcal{I})^* \to (1 + \mathcal{I})^*$ であると示せばよい。
アフィン開集合上の切断をとれば、これは Algebra, Lemma
\ref{algebra-lemma-lift-nth-roots} から従う。
\end{proof}





\section{無限小近傍}
\label{more-morphisms-section-first-order-infinitesimal-neighbourhood}

\noindent
有限次肥厚化の自然な構成は次の通りである。
% P05-EN-CH37-0017
$i : Z \to X$ をスキームの埋め込みとする。$i$ によって $Z$ が
閉部分スキーム $Z \subset U$ と同一視されるような開部分スキーム
$U \subset X$ を選ぶ。$\mathcal{I} \subset \mathcal{O}_U$ を
$U$ 内で $Z$ を定める準連接イデアル層とする。$n \geq 1$ に対し、
準連接イデアル層 $\mathcal{I}^{n + 1}$ によって定まる閉部分スキーム
$Z_n \subset U$ を考えることができる。

\begin{definition}
\label{more-morphisms-definition-first-order-infinitesimal-neighbourhood}
$i : Z \to X$ をスキームの埋め込みとする。
\begin{enumerate}
\item $X$ における $Z$ の {\it 一次無限小近傍} とは、上で記述した
$X$ 上の一次肥厚化 $Z \subset Z_1$ のことである。
\item $X$ における $Z$ の {\it $n$ 次無限小近傍} とは、上で記述した
$X$ 上の $n$ 次肥厚化 $Z \subset Z_n$ のことである。
\end{enumerate}
\end{definition}

\noindent
これらの肥厚化は次の普遍性をもつ（これにより、上の構成が開集合
$U$ の選択に依存するのではないかという懸念は解消される）。

\begin{lemma}
\label{more-morphisms-lemma-first-order-infinitesimal-neighbourhood}
$i : Z \to X$ をスキームの埋め込みとする。
\begin{enumerate}
\item $X$ における $Z$ の一次無限小近傍 $Z'$ は次の普遍性をもつ。
任意の可換図式
$$
\xymatrix{
Z \ar[d]_i & T \ar[l]^a \ar[d] \\
X & T' \ar[l]_b
}
$$
で $T \subset T'$ が $X$ 上の一次肥厚化であるものが与えられたとき、
% P05-EN-CH37-0018
$X$ 上の肥厚化の射
$(a', a) : (T \subset T') \to (Z \subset Z')$ が一意に存在する。
\item $n \geq 1$ に対し、$X$ における $Z$ の $n$ 次無限小近傍 $Z_n$ は
次の普遍性をもつ。任意の可換図式
$$
\xymatrix{
Z \ar[d]_i & T \ar[l]^a \ar[d] \\
X & T' \ar[l]_b
}
$$
で $T \subset T'$ が $X$ 上の $n$ 次肥厚化であるものが与えられたとき、
$X$ 上の肥厚化の射
$(a', a) : (T \subset T') \to (Z \subset Z_n)$ が一意に存在する。
\end{enumerate}
\end{lemma}

\begin{proof}
(1) のみを証明する。$U \subset X$ を $Z'$ の構成に用いた開集合、
すなわち $Z$ が準連接イデアル層 $\mathcal{I}$ で切り出される
$U$ の閉部分スキームと同一視されるような開集合とする。
$|T| = |T'|$ なので $b(T') \subset U$ である。したがって $b$ を
$U$ への射とみなせる。$\mathcal{J} \subset \mathcal{O}_{T'}$ を
$T$ を切り出すイデアルとする。上の図式により $b(T) \subset Z$ なので
$b^\sharp(b^{-1}\mathcal{I}) \subset \mathcal{J}$ である。
$T'$ は $T$ の一次肥厚化だから $\mathcal{J}^2 = 0$ であり、
したがって $b^\sharp(b^{-1}(\mathcal{I}^2)) = 0$ となる。
Schemes, Lemma \ref{schemes-lemma-characterize-closed-subspace}
により $b$ は $Z'$ を経由する。
% P05-EN-CH37-0019
この分解を $a' : T' \to Z'$ と書けば、すべて明らかである。
\end{proof}

\begin{lemma}
\label{more-morphisms-lemma-infinitesimal-neighbourhood-conormal}
$i : Z \to X$ をスキームの埋め込みとし、$Z \subset Z'$ を
$X$ における $Z$ の一次無限小近傍とする。このとき図式
$$
\xymatrix{
Z \ar[r] \ar[d] & Z' \ar[d] \\
Z \ar[r] & X
}
$$
は Morphisms, Lemma \ref{morphisms-lemma-conormal-functorial} により
余法層の写像 $\mathcal{C}_{Z/X} \to \mathcal{C}_{Z/Z'}$ を誘導する。
この写像は同型である。
\end{lemma}

\begin{proof}
これは上の $Z'$ の構成から明らかである。
\end{proof}
















\section{形式的不分岐射}
\label{more-morphisms-section-formally-unramified}

\noindent
環準同型 $R \to A$ が {\it 形式的不分岐} であるとは、任意の可換な実線図式
$$
\xymatrix{
A \ar[r] \ar@{-->}[rd] & B/I \\
R \ar[r] \ar[u] & B \ar[u]
}
$$
で $I \subset B$ が平方零イデアルであるものに対し、図式を可換にする
点線矢印が高々一つ存在することをいう
（Algebra, Definition \ref{algebra-definition-formally-unramified} を参照）。
これに動機づけられて、スキームの射に対する次の類似概念を定める。

\begin{definition}
\label{more-morphisms-definition-formally-unramified}
$f : X \to S$ をスキームの射とする。任意の可換な実線図式
$$
\xymatrix{
X \ar[d]_f & T \ar[d]^i \ar[l] \\
S & T' \ar[l] \ar@{-->}[lu]
}
$$
で $T \subset T'$ が $S$ 上のアフィンスキームの一次肥厚化であるものに対し、
図式を可換にする点線矢印が高々一つ存在するとき、$f$ は
{\it 形式的不分岐} であるという。
\end{definition}

\noindent
まず、対応する図式を描くだけで証明できるいくつかの形式的な補題を示す。

\begin{lemma}
\label{more-morphisms-lemma-formally-unramified-not-affine}
$f : X \to S$ を形式的不分岐射とする。このとき、任意の可換な実線図式
$$
\xymatrix{
X \ar[d]_f & T \ar[d]^i \ar[l] \\
S & T' \ar[l] \ar@{-->}[lu]
}
$$
で $T \subset T'$ が $S$ 上のスキームの一次肥厚化であるものに対し、
図式を可換にする点線矢印は高々一つ存在する。言い換えると、
Definition \ref{more-morphisms-definition-formally-unramified} において
$T$ がアフィンであるという条件を除いてよい。
\end{lemma}

\begin{proof}
射はアフィン開集合への制限によって決定されるからである。
\end{proof}

\begin{lemma}
\label{more-morphisms-lemma-composition-formally-unramified}
形式的不分岐射の合成は形式的不分岐である。
\end{lemma}

\begin{proof}
これは形式的である。
\end{proof}

\begin{lemma}
\label{more-morphisms-lemma-base-change-formally-unramified}
形式的不分岐射の基底変換は形式的不分岐である。
\end{lemma}

\begin{proof}
これは形式的である。
\end{proof}

\begin{lemma}
\label{more-morphisms-lemma-formally-unramified-on-opens}
$f : X \to S$ をスキームの射とする。$U \subset X$ と $V \subset S$ を
$f(U) \subset V$ を満たす開集合とする。$f$ が形式的不分岐ならば、
$f|_U : U \to V$ も形式的不分岐である。
\end{lemma}

\begin{proof}
Definition \ref{more-morphisms-definition-formally-unramified} に現れる実線図式
$$
\xymatrix{
U \ar[d]_{f|_U} & T \ar[d]^i \ar[l]^a \\
V & T' \ar[l] \ar@{-->}[lu]
}
$$
を考える。
% P05-EN-CH37-0020
$f$ が形式的に分岐しているならば、$a'|_T = a$ を満たす
$S$-射 $a' : T' \to X$ は高々一つ存在する。
したがって、そのような $U$ への射も明らかに高々一つしか存在しない。
\end{proof}

\begin{lemma}
\label{more-morphisms-lemma-affine-formally-unramified}
$f : X \to S$ をスキームの射とし、$X$ と $S$ はアフィンであると仮定する。
このとき $f$ が形式的不分岐であることと、環準同型
$\mathcal{O}_S(S) \to \mathcal{O}_X(X)$ が形式的不分岐であることは
同値である。
\end{lemma}

\begin{proof}
環の圏とアフィンスキームの圏の同値により、これは定義
（Definition \ref{more-morphisms-definition-formally-unramified} および
Algebra, Definition \ref{algebra-definition-formally-unramified}）
から直ちに従う。Schemes, Lemma
\ref{schemes-lemma-category-affine-schemes} を参照せよ。
\end{proof}

\noindent
微分層による特徴づけは次の通りである。

\begin{lemma}
\label{more-morphisms-lemma-formally-unramified-differentials}
$f : X \to S$ をスキームの射とする。このとき $f$ が形式的不分岐で
あることと $\Omega_{X/S} = 0$ であることは同値である。
\end{lemma}

\begin{proof}
Morphisms, Lemma \ref{morphisms-lemma-differentials-affine} の証明の
議論の一部を思い出そう。$W \subset X \times_S X$ を、
$\Delta : X \to X \times_S X$ が $W$ への閉埋め込みを誘導するような
開集合とする。$\mathcal{J} \subset \mathcal{O}_W$ をこの閉埋め込みの
イデアル層とし、$X' \subset W$ を準連接イデアル層
$\mathcal{J}^2$ によって定まる閉部分スキームとする。二つの射影
$X \times_S X \to X$ によって誘導される二つの射
$p_1, p_2 : X' \to X$ を考える。$p_1$ と $p_2$ は
$\Delta : X \to X'$ と合成すると一致すること、また
% P05-EN-CH37-0021
$X \to X'$ は平方が零であるイデアルによって定まる閉埋め込みであることに
注意する。さらに短完全列
$$
0 \to \mathcal{J}/\mathcal{J}^2 \to \mathcal{O}_{X'} \to \mathcal{O}_X \to 0
$$
があり、$\Omega_{X/S} = \mathcal{J}/\mathcal{J}^2$ である。
さらに $\mathcal{J}/\mathcal{J}^2$ は、$\mathcal{O}_X$ の局所切断
$f$ に対する局所切断 $p_1^\sharp(f) - p_2^\sharp(f)$ で生成される。

\medskip\noindent
$f : X \to S$ が形式的不分岐であると仮定する。仮定により、
任意のアフィン開集合 $T' \subset X'$ への制限上で $p_1 = p_2$ である。
したがって $p_1 = p_2$ であり、上の議論から
$\Omega_{X/S} = \mathcal{J}/\mathcal{J}^2 = 0$ と結論できる。

\medskip\noindent
逆に $\Omega_{X/S} = 0$ と仮定する。このとき $X' = X$ である。
Definition \ref{more-morphisms-definition-formally-unramified} の図式に点線矢印として
入る任意の二つの射 $f'_1, f'_2 : T' \to X$ をとる。これらは射
$(f'_1, f'_2) : T' \to X \times_S X$ を与える。
$f'_1|_T = f'_2|_T$ かつ $|T| =|T'|$ なので、
$(f'_1, f'_2)$ による $T'$ の像は上で選んだ開集合 $W$ に含まれる。
$(f'_1, f'_2)(T) \subset \Delta(X)$ であり、$T$ は $T'$ 内で
平方零イデアルにより定まるので、$(f'_1, f'_2)$ は $X'$ を経由する。
$X' = X$ だから、所望の通り $f_1' = f'_2$ を得る。
\end{proof}

\begin{lemma}
\label{more-morphisms-lemma-formally-unramified}
$f : X \to Y$ をスキームの射とする。次は同値である。
\begin{enumerate}
\item $f$ は形式的不分岐である。
\item 任意の $x \in X$ に対し、$f(U) \subset V$ かつ
$f|_U : U \to V$ が形式的不分岐となるような開集合
$x \in U \subset X$ および $f(x) \in V \subset Y$ が存在する。
\item $f(U) \subset V$ を満たす任意のアフィン開集合の対
$U \subset X$, $V \subset Y$ に対し、環準同型
$\mathcal{O}_Y(V) \to \mathcal{O}_X(U)$ は形式的不分岐である。
\item アフィン開被覆 $Y = \bigcup V_j$ と、各 $j$ に対するアフィン開被覆
$f^{-1}(V_j) = \bigcup U_{ji}$ が存在して、すべての $j$, $i$ に対し
% P05-EN-CH37-0022
環準同型 $\mathcal{O}_Y(V) \to \mathcal{O}_X(U)$ が形式的不分岐である。
\end{enumerate}
\end{lemma}

\begin{proof}
Lemmas \ref{more-morphisms-lemma-formally-unramified-on-opens} と
\ref{more-morphisms-lemma-affine-formally-unramified} を組み合わせると
(1) $\Rightarrow$ (3) が分かる。(3) $\Rightarrow$ (4) は明らかであり、
Lemma \ref{more-morphisms-lemma-affine-formally-unramified} により
(4) $\Rightarrow$ (2) である。(2) が成り立つとする。このとき
% P05-EN-CH37-0023
Lemma \ref{more-morphisms-lemma-formally-unramified-differentials} により
$\Omega_{X/S}$ は各点の近傍で消える
（ここでは Morphisms, Lemma
\ref{morphisms-lemma-differentials-restrict-open} も用いる）。
そこで Lemma \ref{more-morphisms-lemma-formally-unramified-differentials} により
$f$ は形式的不分岐である。
\end{proof}

\begin{lemma}
\label{more-morphisms-lemma-unramified-formally-unramified}
\begin{slogan}
% P05-EN-CH37-0024
不分岐射とは、局所有限型な形式的不分岐射にほかならない。
\end{slogan}
$f : X \to S$ をスキームの射とする。次は同値である。
\begin{enumerate}
\item 射 $f$ は不分岐（それぞれ G-不分岐）である。
\item 射 $f$ は局所有限型（それぞれ局所有限表示）かつ形式的不分岐である。
\end{enumerate}
\end{lemma}

\begin{proof}
Lemma \ref{more-morphisms-lemma-formally-unramified-differentials} および
Morphisms, Lemma \ref{morphisms-lemma-unramified-omega-zero} を用いる。
\end{proof}









\section{普遍一次肥厚化}
\label{more-morphisms-section-universal-thickening}

\noindent
$h : Z \to X$ をスキームの射とする。$X$ 上の $Z$ の
{\it 普遍一次肥厚化} とは $X$ 上の一次肥厚化 $Z \subset Z'$ であって、
$X$ 上の任意の一次肥厚化 $T \subset T'$ と可換な実線図式
$$
\xymatrix{
& Z \ar[ld] & & T \ar[rd] \ar[ll]^a \\
Z' \ar[rrd] & & & & T' \ar@{..>}[llll]_{a'} \ar[lld]^b \\
 & & X
}
$$
が与えられたとき、図式を可換にする点線矢印が一意に存在するものをいう。
この状況では $(a, a') : (T \subset T') \to (Z \subset Z')$ は
$X$ 上の肥厚化の射であることに注意する。したがって普遍一次肥厚化が
存在すれば、一意的同型を除いて一意である。一般には普遍一次肥厚化は
存在しないが、$h$ が形式的不分岐ならば存在する。

\begin{lemma}
\label{more-morphisms-lemma-universal-thickening}
$h : Z \to X$ をスキームの形式的不分岐射とする。このとき
$X$ 上の $Z$ の普遍一次肥厚化 $Z \subset Z'$ が存在する。
\end{lemma}

\begin{proof}
この証明では、$Z \subset Z'$ が補題の条件を満たすとき、
$X$ 上の $Z$ の普遍一次肥厚化であるという。まず
Algebra, Lemma \ref{algebra-lemma-universal-thickening} である
アフィンの場合から始め、貼り合わせによって $X$ 上の普遍一次肥厚化
$Z \subset Z'$ を構成する。いくつかの一般的注意から始めよう。

\medskip\noindent
$X$ 上の $Z$ の普遍一次肥厚化が存在すれば、一意的同型を除いて一意である。
さらに $V \subset Z$ と $U \subset X$ を $h(V) \subset U$ を満たす
開部分スキームとし、$Z \subset Z'$ を $X$ 上の $Z$ の普遍一次肥厚化とする。
$V' \subset Z'$ を $V = Z \cap V'$ を満たす開部分スキームとする。
このとき $V \subset V'$ は $U$ 上の $V$ の普遍一次肥厚化であると主張する。
実際、任意の図式
$$
\xymatrix{
V \ar[d]_h & T \ar[l]^a \ar[d] \\
U & T' \ar[l]_b
}
$$
で $T \subset T'$ が $U$ 上の一次肥厚化であるものを考える。
$Z'$ の普遍性により
$(a, a') : (T \subset T') \to (Z \subset Z')$ を得る。
基礎位相空間の等式 $|T| = |T'|$ があるので $a'(T') \subset V'$ である。
したがって $(a, a')$ を $U$ 上の肥厚化の射
$(a, a') : (T \subset T') \to (V \subset V')$ とみなせる。
一意性も明らかである。まったく同様に、$h(Z) \subset U$ であり
$Z \subset Z'$ が $U$ 上の普遍一次肥厚化ならば、$Z \subset Z'$ は
$X$ 上の普遍一次肥厚化でもあることが示される。

\medskip\noindent
アフィン部分を貼り合わせる前に、$Z$ と $X$ がアフィンの場合に
補題を示そう。$X = \Spec(R)$, $Z = \Spec(S)$ とする。
Algebra, Lemma \ref{algebra-lemma-universal-thickening}
により、図式
$$
\xymatrix{
Z \ar[d]_h & T \ar[l]^a \ar[d] \\
X & T' \ar[l]_b
}
$$
で $T, T'$ がアフィンであるものに対して補題の普遍性をもつ
$X$ 上の一次肥厚化 $Z \subset Z'$ が存在する。一般の図式に対しては
アフィン開被覆 $T' = \bigcup T'_i$ を選べば、
$a'_i|_{T_i} = a|_{T_i}$ を満たす $X$ 上の射
$a'_i : T'_i \to Z'$ を得る。一意性により $a'_i$ と $a'_j$ は
$T'_i \cap T'_j$ の任意のアフィン開集合上で一致する。
したがって射 $a'_i$ は所望の $X$ 上の大域射
$a' : T' \to Z'$ に貼り合わさる。ゆえに $X$ と $Z$ がアフィンなら
補題は成り立つ。

\medskip\noindent
$Z_i$ がそれぞれ $X$ のアフィン開集合 $U_i$ に写るような
アフィン開被覆 $Z = \bigcup Z_i$ を選ぶ。
Lemma \ref{more-morphisms-lemma-formally-unramified-on-opens}
により射 $Z_i \to U_i$ は形式的不分岐である。したがってアフィンの場合から、
$U_i$ 上の普遍一次肥厚化 $Z_i \subset Z_i'$ を得る。
上の一般的注意により、$Z_i \subset Z_i'$ は $X$ 上の $Z_i$ の
普遍一次肥厚化でもある。$Z'_{i, j} \subset Z'_i$ を
$Z_i \cap Z_j = Z'_{i, j} \cap Z_i$ を満たす開部分スキームとする。
一般的注意により、$Z'_{i, j}$ と $Z'_{j, i}$ はともに $X$ 上の
$Z_i \cap Z_j$ の普遍一次肥厚化である。したがって最初の一般的注意から、
$Z_i \cap Z_j$ 上に恒等写像を誘導する標準的同型
$\varphi_{ij} : Z'_{i, j} \to Z'_{j, i}$ がある。
これらの射は Schemes, Section \ref{schemes-section-glueing-schemes} の
コサイクル条件を満たすと主張する
（確認は省略する。ヒント：$Z'_{i, j} \cap Z'_{i, k}$ は
$Z_i \cap Z_j \cap Z_k$ の普遍一次肥厚化であり、上の議論により
一意的同型を除いて一意に定まることを用いる）。
したがって Schemes, Section \ref{schemes-section-glueing-schemes} の結果から、
% P05-EN-CH37-0025
$Z_i = Z'_i \cap Z$ を満たす開部分スキーム $Z'_i \subset Z'$ が
$X$ 上の $Z_i$ の普遍一次肥厚化となるという性質をもつ
$X$ 上の一次肥厚化 $Z \subset Z'$ を得る。

\medskip\noindent
これから形式的に $Z'$ が $X$ 上の $Z$ の普遍一次肥厚化であることが従う。
実際、図式
$$
\xymatrix{
Z \ar[d]_h & T \ar[l]^a \ar[d] \\
X & T' \ar[l]_b
}
$$
で $a(T)$ がある $Z_i$ に含まれるものに対しては普遍性が成り立つ。
一般の図式に対し、$a(T_i) \subset Z_i$ となるような開被覆
$T' = \bigcup T'_i$ を選べる。$a'_i|_{T_i} = a|_{T_i}$ を満たす
$X$ 上の射 $a'_i : T'_i \to Z'$ を得る。$a'_i|_{T'_i \cap T'_j}$ と
$a'_j|_{T'_i \cap T'_j}$ はともに、$X$ 上の $Z_i \cap Z_j$ の
普遍一次肥厚化 $Z'_i \cap Z'_j$ への射を求める問題の解なので、
$a'_i$ と $a'_j$ は $T'_i \cap T'_j$ 上で必ず一致する。
したがって射 $a'_i$ は所望の $X$ 上の大域射
$a' : T' \to Z'$ に貼り合わさる。これで証明が完了する。
\end{proof}

\begin{definition}
\label{more-morphisms-definition-universal-thickening}
$h : Z \to X$ をスキームの形式的不分岐射とする。
\begin{enumerate}
\item $X$ 上の $Z$ の {\it 普遍一次肥厚化} とは、
Lemma \ref{more-morphisms-lemma-universal-thickening} で構成した肥厚化
$Z \subset Z'$ のことである。
\item {\it $X$ 上の $Z$ の余法層} とは、$X$ 上の普遍一次肥厚化
$Z'$ における $Z$ の余法層のことである。
\end{enumerate}
この状況では余法層をしばしば $\mathcal{C}_{Z/X}$ と書く。
\end{definition}

\noindent
したがって $Z$ 上の層の短完全列
$$
0 \to \mathcal{C}_{Z/X} \to \mathcal{O}_{Z'} \to \mathcal{O}_Z \to 0
$$
がある。次の補題は、この定義と $Z \to X$ が埋め込みの場合の定義との間に
矛盾がないことを示す。

\begin{lemma}
\label{more-morphisms-lemma-immersion-universal-thickening}
$i : Z \to X$ をスキームの埋め込みとする。このとき
\begin{enumerate}
\item $i$ は形式的不分岐である。
\item $X$ 上の $Z$ の普遍一次肥厚化は、Definition
\ref{more-morphisms-definition-first-order-infinitesimal-neighbourhood} における
$X$ 内の $Z$ の一次無限小近傍である。
\item Morphisms, Definition \ref{morphisms-definition-conormal-sheaf}
の意味での $i$ の余法層は、Definition
\ref{more-morphisms-definition-universal-thickening} の意味での $i$ の余法層と一致する。
\end{enumerate}
\end{lemma}

\begin{proof}
Morphisms, Lemmas \ref{morphisms-lemma-open-immersion-unramified} および
\ref{morphisms-lemma-closed-immersion-unramified}
により埋め込みは不分岐であり、したがって
Lemma \ref{more-morphisms-lemma-unramified-formally-unramified} により形式的不分岐である。
他の主張は Lemmas \ref{more-morphisms-lemma-first-order-infinitesimal-neighbourhood} および
\ref{more-morphisms-lemma-infinitesimal-neighbourhood-conormal}
を定義と組み合わせれば従う。
\end{proof}

\begin{lemma}
\label{more-morphisms-lemma-universal-thickening-unramified}
$Z \to X$ をスキームの形式的不分岐射とする。このとき普遍一次肥厚化
$Z'$ は $X$ 上形式的不分岐である。
\end{lemma}

\begin{proof}
二つの証明がある。第一の証明では、アフィン局所的に議論して
Algebra, Lemma \ref{algebra-lemma-differentials-universal-thickening}
を適用し、$\Omega_{Z'/X} = 0$ を示す。その後
Lemma \ref{more-morphisms-lemma-formally-unramified-differentials} により所望の結果を得る。
第二の証明は次の直接的議論である。

\medskip\noindent
$T \subset T'$ を一次肥厚化とし、可換図式
$$
\xymatrix{
Z' \ar[d] & T \ar[l]^c \ar[d] \\
X & T' \ar[l] \ar[lu]^{a, b}
}
$$
をとる。図式に入る二つの射 $a, b : T' \to Z'$ を考える。
$T_0 = c^{-1}(Z) \subset T$ および $T'_a = a^{-1}(Z)$
（スキーム論的に）とおく。$Z'$ は $Z$ の一次肥厚化なので、
$T'$ は $T'_a$ の一次肥厚化である。さらに $c = a|_T$ だから
$T_0 = T \cap T'_a$（スキーム論的に）である。$T'$ は $T$ の
一次肥厚化なので、$T'_a$ は $T_0$ の一次肥厚化である。
$a|_{T'_a}$ と $b|_{T'_a}$ は $T'_a$ から $Z'$ への $X$ 上の射であり、
$T_0$ 上では $Z$ への射として一致する。したがって $Z'$ の普遍性により
$a|_{T'_a} = b|_{T'_a}$ である。
% P05-EN-CH37-0026
ゆえに $a$ と $b$ は、$T'_a$ の一次肥厚化 $T'$ からの射であって、
$T'_a$ への制限が $Z$ への射として一致する。そこで $Z'$ の普遍性を
もう一度用いると $a = b$ を得る。言い換えると、形式的不分岐射の
定義条件が所望の通り $Z' \to X$ に対して成り立つ。
\end{proof}

\begin{lemma}
\label{more-morphisms-lemma-universal-thickening-functorial}
スキームの可換図式
$$
\xymatrix{
Z \ar[r]_h \ar[d]_f & X \ar[d]^g \\
W \ar[r]^{h'} & Y
}
$$
を考え、$h$ と $h'$ は形式的不分岐であるとする。$Z \subset Z'$ を
$X$ 上の $Z$ の普遍一次肥厚化、$W \subset W'$ を $Y$ 上の $W$ の
普遍一次肥厚化とする。
% P05-EN-CH37-0027
次の可換図式に入る $Y$ 上の肥厚化の標準的射
$(f, f') : (Z, Z') \to (W, W')$ が存在する。
$$
\xymatrix{
& & & Z' \ar[ld] \ar[d]^{f'} \\
Z \ar[rr] \ar[d]_f \ar[rrru] & & X \ar[d] & W' \ar[ld] \\
W \ar[rrru]|!{[rr];[rruu]}\hole \ar[rr] & & Y
}
$$
特に、肥厚化の射 $(f, f')$ は余法層の射
$f^*\mathcal{C}_{W/Y} \to \mathcal{C}_{Z/X}$ を誘導する。
\end{lemma}

\begin{proof}
第一の主張は $W'$ の普遍性から明らかである。余法層上に誘導される写像は、
Morphisms, Lemma \ref{morphisms-lemma-conormal-functorial}
を $(Z \subset Z') \to (W \subset W')$ に適用して得られる写像である。
\end{proof}

\begin{lemma}
\label{more-morphisms-lemma-universal-thickening-fibre-product}
図式
$$
\xymatrix{
Z \ar[r]_h \ar[d]_f & X \ar[d]^g \\
W \ar[r]^{h'} & Y
}
$$
をスキームの圏におけるファイバー積図式とし、$h'$ は形式的不分岐であるとする。
このとき $h$ は形式的不分岐である。また $W \subset W'$ が
$Y$ 上の $W$ の普遍一次肥厚化ならば、
$Z = X \times_Y W \subset X \times_Y W'$ は $X$ 上の $Z$ の
普遍一次肥厚化である。特に Lemma
\ref{more-morphisms-lemma-universal-thickening-functorial} の標準的写像
$f^*\mathcal{C}_{W/Y} \to \mathcal{C}_{Z/X}$
は全射である。
\end{lemma}

\begin{proof}
$h$ は Lemma \ref{more-morphisms-lemma-base-change-formally-unramified} により
形式的不分岐である。$X \times_Y W'$ が一次肥厚化であることは明らかである。
$W'$ が普遍性をもつこととファイバー積の写像性質から、
これが普遍性をもつことも直ちに確かめられる。余法層の写像が
全射になることについては
Morphisms, Lemma \ref{morphisms-lemma-conormal-functorial-flat}
を参照せよ。
\end{proof}

\begin{lemma}
\label{more-morphisms-lemma-universal-thickening-fibre-product-flat}
図式
$$
\xymatrix{
Z \ar[r]_h \ar[d]_f & X \ar[d]^g \\
W \ar[r]^{h'} & Y
}
$$
をスキームの圏におけるファイバー積図式とし、$h'$ は形式的不分岐、
$g$ は平坦であるとする。このとき、対応する普遍一次肥厚化の写像
$Z' \to W'$ は平坦であり、
$f^*\mathcal{C}_{W/Y} \to \mathcal{C}_{Z/X}$ は同型である。
\end{lemma}

\begin{proof}
平坦性は基底変換で保たれる。Morphisms, Lemma
\ref{morphisms-lemma-base-change-flat} を参照せよ。したがって第一の主張は、
% P05-EN-CH37-0028
Lemma \ref{more-morphisms-lemma-universal-thickening-fibre-product} における
$W'$ の記述から従う。$X \times_Y W'$ が一次肥厚化であることは明らかである。
$W'$ が普遍性をもつこととファイバー積の写像性質から、
これが普遍性をもつことも直ちに確かめられる。
これにより余法層の写像が同型となる理由については
Morphisms, Lemma \ref{morphisms-lemma-conormal-functorial-flat}
を参照せよ。
\end{proof}

\begin{lemma}
\label{more-morphisms-lemma-universal-thickening-localize}
普遍一次肥厚化をとる操作は開部分をとる操作と可換である。
より正確には、$h : Z \to X$ をスキームの形式的不分岐射とし、
$V \subset Z$, $U \subset X$ を $h(V) \subset U$ を満たす開集合とする。
$Z'$ を $X$ 上の $Z$ の普遍一次肥厚化とする。このとき
$h|_V : V \to U$ は形式的不分岐であり、$U$ 上の $V$ の普遍一次肥厚化は
$V = Z \cap V'$ を満たす開部分スキーム $V' \subset Z'$ である。
特に $\mathcal{C}_{Z/X}|_V = \mathcal{C}_{V/U}$ である。
\end{lemma}

\begin{proof}
第一の主張は Lemma \ref{more-morphisms-lemma-formally-unramified-on-opens} である。
普遍肥厚化の整合性は Lemma \ref{more-morphisms-lemma-universal-thickening} の証明、
Algebra, Lemma \ref{algebra-lemma-universal-thickening-localize}、
または Lemma \ref{more-morphisms-lemma-universal-thickening-fibre-product-flat}
から導ける。
\end{proof}

\begin{lemma}
\label{more-morphisms-lemma-differentials-universally-unramified}
$h : Z \to X$ を $S$ 上のスキームの形式的不分岐射とする。
$Z \subset Z'$ を $X$ 上の $Z$ の普遍一次肥厚化とし、その構造射を
$h' : Z' \to X$ とする。標準的写像
$$
c_{h'} : (h')^*\Omega_{X/S} \longrightarrow \Omega_{Z'/S}
$$
は同型 $h^*\Omega_{X/S} \to \Omega_{Z'/S} \otimes \mathcal{O}_Z$
を誘導する。
\end{lemma}

\begin{proof}
写像 $c_{h'}$ は Morphisms, Lemma
\ref{morphisms-lemma-functoriality-differentials} で定義された写像である。
$i : Z \to Z'$ を与えられた閉埋め込みとすると、$i^*c_{h'}$ は写像
$h^*\Omega_{X/S} \to \Omega_{Z'/S} \otimes \mathcal{O}_Z$ である。
これが同型であることの確認は、局所化によりアフィンの場合に帰着する。
Lemma \ref{more-morphisms-lemma-universal-thickening-localize} および Morphisms, Lemma
\ref{morphisms-lemma-differentials-restrict-open} を参照せよ。この場合の結果は
Algebra, Lemma \ref{algebra-lemma-differentials-universal-thickening} である。
\end{proof}

\begin{lemma}
\label{more-morphisms-lemma-universally-unramified-differentials-sequence}
$h : Z \to X$ を $S$ 上のスキームの形式的不分岐射とする。
標準的完全列
$$
\mathcal{C}_{Z/X} \to h^*\Omega_{X/S} \to \Omega_{Z/S} \to 0.
$$
がある。ここで第一の矢印は $\text{d}_{Z'/S}$ によって誘導され、
$Z'$ は $X$ 上の $Z$ の普遍一次近傍である。
\end{lemma}

\begin{proof}
標準的完全列
$$
\mathcal{C}_{Z/Z'} \to
\Omega_{Z'/S} \otimes \mathcal{O}_Z \to
\Omega_{Z/S} \to 0.
$$
があることが分かっている。Morphisms, Lemma
\ref{morphisms-lemma-differentials-relative-immersion} を参照せよ。
したがって Lemma \ref{more-morphisms-lemma-differentials-universally-unramified}
を適用すれば結果が従う。
\end{proof}

\begin{lemma}
\label{more-morphisms-lemma-two-unramified-morphisms}
図式
$$
\xymatrix{
Z \ar[r]_i \ar[rd]_j & X \ar[d] \\
& Y
}
$$
をスキームの可換図式とし、$i$ と $j$ は形式的不分岐であるとする。
このとき標準的完全列
$$
\mathcal{C}_{Z/Y} \to
\mathcal{C}_{Z/X} \to
i^*\Omega_{X/Y} \to 0
$$
がある。第一の矢印は Lemma \ref{more-morphisms-lemma-universal-thickening-functorial}
から、第二の矢印は
Lemma \ref{more-morphisms-lemma-universally-unramified-differentials-sequence}
から得られる。
\end{lemma}

\begin{proof}
% P05-EN-CH37-0029
$Z \to Z'$ を $X$ 上の $Z$ の普遍一次肥厚化、
$Z \to Z''$ を $Y$ 上の $Z$ の普遍一次肥厚化と書く。
% P05-EN-CH37-0030
Lemma \ref{more-morphisms-lemma-universally-unramified-differentials-sequence}
により標準的射 $Z' \to Z''$ があり、可換図式
$$
\xymatrix{
Z \ar[r]_{i'} \ar[rd]_{j'} & Z' \ar[r] \ar[d] & X \ar[d] \\
& Z'' \ar[r] & Y
}
$$
を得る。左の三角形に Morphisms, Lemma
\ref{morphisms-lemma-two-immersions} を適用すると完全列
$$
\mathcal{C}_{Z/Z''} \to
\mathcal{C}_{Z/Z'} \to
(i')^*\Omega_{Z'/Z''} \to 0
$$
を得る。$Z''$ は $Y$ 上形式的不分岐なので
（Lemma \ref{more-morphisms-lemma-universal-thickening-unramified} を参照）、
% P05-EN-CH37-0031
$\Omega_{Z'/Z''} = \Omega_{Z/Y}$
と分かる（Lemma \ref{more-morphisms-lemma-formally-unramified-differentials} と
Morphisms, Lemma \ref{morphisms-lemma-triangle-differentials} を
組み合わせる）。さらに Lemma
\ref{more-morphisms-lemma-differentials-universally-unramified} により
$(i')^*\Omega_{Z'/Y} = i^*\Omega_{X/Y}$ である。
\end{proof}

\begin{lemma}
\label{more-morphisms-lemma-transitivity-conormal}
$Z \to Y \to X$ をスキームの形式的不分岐射とする。
\begin{enumerate}
\item $Z \subset Z'$ が $X$ 上の $Z$ の普遍一次肥厚化であり、
$Y \subset Y'$ が $X$ 上の $Y$ の普遍一次肥厚化ならば、
射 $Z' \to Y'$ が存在し、$Y \times_{Y'} Z'$ は
$Y$ 上の $Z$ の普遍一次肥厚化である。
\item 標準的完全列
$$
i^*\mathcal{C}_{Y/X} \to
\mathcal{C}_{Z/X} \to
\mathcal{C}_{Z/Y} \to 0
$$
がある。ここで写像は Lemma
\ref{more-morphisms-lemma-universal-thickening-functorial} から得られ、
$i : Z \to Y$ は最初の射である。
\end{enumerate}
\end{lemma}

\begin{proof}
(1) の写像 $h : Z' \to Y'$ は
Lemma \ref{more-morphisms-lemma-universal-thickening-functorial} から得られる。
$Y \times_{Y'} Z'$ が $Y$ 上の $Z$ の普遍一次肥厚化であることは、
$Z'$ と $Y'$ の普遍性から明らかである。
Morphisms, Lemma \ref{morphisms-lemma-transitivity-conormal}
により完全列
$$
(i')^*\mathcal{C}_{Y \times_{Y'} Z'/Z'} \to
\mathcal{C}_{Z/Z'} \to
\mathcal{C}_{Z/Y \times_{Y'} Z'} \to 0
$$
がある。ここで $i' : Z \to Y \times_{Y'} Z'$ は与えられた射である。
Morphisms, Lemma \ref{morphisms-lemma-conormal-functorial-flat}
により全射
$h^*\mathcal{C}_{Y/Y'} \to \mathcal{C}_{Y \times_{Y'} Z'/Z'}$
が存在する。これを等式
$\mathcal{C}_{Y/Y'} = \mathcal{C}_{Y/X}$,
$\mathcal{C}_{Z/Z'} = \mathcal{C}_{Z/X}$, および
$\mathcal{C}_{Z/Y \times_{Y'} Z'} = \mathcal{C}_{Z/Y}$
と組み合わせれば補題が従う。
\end{proof}











\section{形式的エタール射}
\label{more-morphisms-section-formally-etale}

\noindent
環準同型 $R \to A$ が {\it 形式的エタール} であるとは、任意の可換な実線図式
$$
\xymatrix{
A \ar[r] \ar@{-->}[rd] & B/I \\
R \ar[r] \ar[u] & B \ar[u]
}
$$
で $I \subset B$ が平方零イデアルであるものに対し、図式を可換にする
点線矢印がちょうど一つ存在することをいう
（Algebra, Definition \ref{algebra-definition-formally-etale} を参照）。
これに動機づけられて、スキームの射に対する次の類似概念を定める。

\begin{definition}
\label{more-morphisms-definition-formally-etale}
$f : X \to S$ をスキームの射とする。任意の可換な実線図式
$$
\xymatrix{
X \ar[d]_f & T \ar[d]^i \ar[l] \\
S & T' \ar[l] \ar@{-->}[lu]
}
$$
で $T \subset T'$ が $S$ 上のアフィンスキームの一次肥厚化であるものに対し、
図式を可換にする点線矢印がちょうど一つ存在するとき、$f$ は
{\it 形式的エタール} であるという。
\end{definition}

\noindent
形式的エタール射が形式的不分岐であることは明らかである。
したがって $f : X \to S$ が形式的エタールならば
$\Omega_{X/S}$ は零である。Lemma
\ref{more-morphisms-lemma-formally-unramified-differentials} を参照せよ。

\begin{lemma}
\label{more-morphisms-lemma-formally-etale-not-affine}
$f : X \to S$ を形式的エタール射とする。このとき任意の可換な実線図式
$$
\xymatrix{
X \ar[d]_f & T \ar[d]^i \ar[l] \\
S & T' \ar[l] \ar@{-->}[lu]
}
$$
で $T \subset T'$ が $S$ 上のスキームの一次肥厚化であるものに対し、
図式を可換にする点線矢印がちょうど一つ存在する。言い換えると、
Definition \ref{more-morphisms-definition-formally-etale} において
$T$ がアフィンであるという条件を除いてよい。
\end{lemma}

\begin{proof}
$T' = \bigcup T'_i$ をアフィン開被覆とし、$T_i = T \cap T'_i$ とおく。
図式に入る射 $a'_i : T'_i \to X$ を得る。一意性により
$a'_i$ と $a'_j$ は $T'_i \cap T'_j$ の任意のアフィン開部分スキーム上で
一致する。したがって $a'_i$ と $a'_j$ は $T'_i \cap T'_j$ 上で一致するので、
射 $a'_i$ は
大域射 $a' : T' \to X$ に貼り合わさる。$a'$ の一意性は
Lemma \ref{more-morphisms-lemma-formally-unramified-not-affine} で示した。
\end{proof}

\begin{lemma}
\label{more-morphisms-lemma-composition-formally-etale}
形式的エタール射の合成は形式的エタールである。
\end{lemma}

\begin{proof}
これは形式的である。
\end{proof}

\begin{lemma}
\label{more-morphisms-lemma-base-change-formally-etale}
形式的エタール射の基底変換は形式的エタールである。
\end{lemma}

\begin{proof}
これは形式的である。
\end{proof}

\begin{lemma}
\label{more-morphisms-lemma-formally-etale-on-opens}
$f : X \to S$ をスキームの射とする。$U \subset X$ と $V \subset S$ を
$f(U) \subset V$ を満たす開部分スキームとする。$f$ が形式的エタールならば、
$f|_U : U \to V$ も形式的エタールである。
\end{lemma}

\begin{proof}
Definition \ref{more-morphisms-definition-formally-etale} に現れる実線図式
$$
\xymatrix{
U \ar[d]_{f|_U} & T \ar[d]^i \ar[l]^a \\
V & T' \ar[l] \ar@{-->}[lu]
}
$$
を考える。
% P05-EN-CH37-0032
$f$ が形式的に分岐しているならば、$a'|_T = a$ を満たす
$S$-射 $a' : T' \to X$ がちょうど一つ存在する。
$|T'| = |T|$ なので $a'(T') \subset U$ となり、$T'$ から $U$ への
一意な射が得られる。
\end{proof}

\begin{lemma}
\label{more-morphisms-lemma-characterize-formally-etale}
$f : X \to S$ をスキームの射とする。次は同値である。
\begin{enumerate}
\item $f$ は形式的エタールである。
\item $f$ は形式的不分岐であり、$S$ 上の $X$ の普遍一次肥厚化は
$X$ に等しい。
\item $f$ は形式的不分岐であり、$\mathcal{C}_{X/S} = 0$ である。
\item $\Omega_{X/S} = 0$ かつ $\mathcal{C}_{X/S} = 0$ である。
\end{enumerate}
\end{lemma}

\begin{proof}
% P05-EN-CH37-0033
実際、最後の主張が意味をなすのは、$\Omega_{X/S} = 0$ ならば
Lemma \ref{more-morphisms-lemma-formally-unramified-differentials}
および Definition \ref{more-morphisms-definition-universal-thickening} によって
$\mathcal{C}_{X/S}$ が定義されるからである。これから (3) と (4) が
同値であることも明らかになる。

\medskip\noindent
% P05-EN-CH37-0034
(1), (2), (3) のいずれの仮定からも $f$ が形式的不分岐であることが従う。
したがって $f$ は形式的不分岐であると仮定してよい。
% P05-EN-CH37-0035
(1), (2), (3) の同値性は、$S$ 上の $X$ の普遍一次肥厚化 $X'$ の普遍性と
$X = X' \Leftrightarrow \mathcal{C}_{X/S} = 0$ という事実から従う。
実際、定義により $\mathcal{C}_{X/S} = \mathcal{C}_{X/X'}$ は
$X'$ における $X$ のイデアル層である。
\end{proof}

\begin{lemma}
\label{more-morphisms-lemma-unramified-flat-formally-etale}
不分岐平坦射は形式的エタールである。
\end{lemma}

\begin{proof}
$X \to S$ は不分岐かつ平坦であるとする。このとき
$\Delta : X \to X \times_S X$ は開埋め込みである。Morphisms, Lemma
\ref{morphisms-lemma-diagonal-unramified-morphism} を参照せよ。
$\mathcal{C}_{X/S}$ が零であることを示さなければならない。
二つの射影 $p, q : X \times_S X \to X$ を考える。
$f$ は形式的不分岐なので
（Lemma \ref{more-morphisms-lemma-unramified-formally-unramified} を参照）、
$q$ も形式的不分岐である
（Lemma \ref{more-morphisms-lemma-base-change-formally-unramified} を参照）。
$f$ は平坦なので $p$ も平坦である。Morphisms, Lemma
\ref{morphisms-lemma-base-change-flat} を参照せよ。したがって
Lemma \ref{more-morphisms-lemma-universal-thickening-fibre-product-flat} により
$p^*\mathcal{C}_{X/S} = \mathcal{C}_q$ である。ここで
$\mathcal{C}_q$ は形式的不分岐射 $q : X \times_S X \to X$ の余法層を表す。
しかし $\Delta(X) \subset X \times_S X$ は $q$ によって $X$ に
同型に写る開部分スキームである。したがって
Lemma \ref{more-morphisms-lemma-universal-thickening-localize}
により $\mathcal{C}_q|_{\Delta(X)} = \mathcal{C}_{X/X} = 0$ である。
言い換えると、恒等射による $\mathcal{C}_{X/S}$ の $X$ への引き戻しは
零である。すなわち $\mathcal{C}_{X/S} = 0$ である。
\end{proof}

\begin{lemma}
\label{more-morphisms-lemma-affine-formally-etale}
$f : X \to S$ をスキームの射とする。
$X$ と $S$ はアフィンであると仮定する。
このとき $f$ が形式的エタールであることと、環準同型
$\mathcal{O}_S(S) \to \mathcal{O}_X(X)$ が形式的エタールであることは
同値である。
\end{lemma}

\begin{proof}
環の圏とアフィンスキームの圏の同値により、これは定義
（Definition \ref{more-morphisms-definition-formally-etale} および
Algebra, Definition \ref{algebra-definition-formally-etale}）から直ちに従う。
Schemes, Lemma \ref{schemes-lemma-category-affine-schemes} を参照せよ。
\end{proof}

\begin{lemma}
\label{more-morphisms-lemma-formally-etale}
$f : X \to Y$ をスキームの射とする。次は同値である。
\begin{enumerate}
\item $f$ は形式的エタールである。
\item 任意の $x \in X$ に対し、$f(U) \subset V$ を満たす開部分スキーム
$x \in U \subset X$ と $f(x) \in V \subset Y$ が存在して、
$f|_U : U \to V$ は形式的エタールである。
\item $f(U) \subset V$ を満たす任意のアフィン開部分スキームの組
$U \subset X$ と $V \subset Y$ に対し、環準同型
$\mathcal{O}_Y(V) \to \mathcal{O}_X(U)$ は形式的エタールである。
\item アフィン開被覆 $Y = \bigcup V_j$ が存在し、各 $j$ に対して
$f^{-1}(V_j) = \bigcup U_{ji}$ はアフィン開被覆であり、すべての $j$ と $i$ に
% P05-EN-CH37-0036
対して環準同型 $\mathcal{O}_Y(V) \to \mathcal{O}_X(U)$ は形式的エタールである。
\end{enumerate}
\end{lemma}

\begin{proof}
Lemmas \ref{more-morphisms-lemma-formally-etale-on-opens} と
\ref{more-morphisms-lemma-affine-formally-etale} を合わせると (1) $\Rightarrow$ (3) を得る。
(3) $\Rightarrow$ (4) は直ちに分かる。
Lemma \ref{more-morphisms-lemma-affine-formally-etale} により (4) $\Rightarrow$ (2) である。
(2) が成り立つならば、Lemma \ref{more-morphisms-lemma-characterize-formally-etale}
（ここでは Morphisms, Lemma
\ref{morphisms-lemma-differentials-restrict-open} および
Lemma \ref{more-morphisms-lemma-universal-thickening-localize} も用いる）により、
% P05-EN-CH37-0037
$\Omega_{X/S}$ と $\mathcal{C}_{X/S}$ は各点の近傍で消える。そこで
% P05-EN-CH37-0038
Lemma \ref{more-morphisms-lemma-characterize-formally-etale} により $f$ は形式的不分岐である。
\end{proof}

\begin{lemma}
\label{more-morphisms-lemma-etale-formally-etale}
$f : X \to S$ をスキームの射とする。次は同値である。
\begin{enumerate}
\item 射 $f$ はエタールである。
\item 射 $f$ は局所有限表示かつ形式的エタールである。
\end{enumerate}
\end{lemma}

\begin{proof}
$f$ がエタールであると仮定する。エタール射は局所有限表示、平坦、かつ
不分岐である。Morphisms, Section \ref{morphisms-section-etale} を参照せよ。
したがって $f$ は局所有限表示かつ形式的エタールである。Lemma
\ref{more-morphisms-lemma-unramified-flat-formally-etale} を参照せよ。

\medskip\noindent
逆に、$f$ は局所有限表示かつ形式的エタールであると仮定する。
エタールであることは $X$ と $S$ 上の Zariski 位相に関して局所的である。
Morphisms, Lemma \ref{morphisms-lemma-etale-characterize} を参照せよ。
Lemma \ref{more-morphisms-lemma-formally-etale-on-opens} により、$X$ をアフィン開部分スキーム
$U$ で被覆でき、各々はアフィン開部分スキーム $V$ に写り、
$U \to V$ は形式的エタールである（また有限表示でもある。Morphisms,
Lemma \ref{morphisms-lemma-locally-finite-presentation-characterize} を参照せよ）。
Lemma \ref{more-morphisms-lemma-affine-formally-etale} により、環準同型
$\mathcal{O}(V) \to \mathcal{O}(U)$ は形式的エタールである
（また有限表示でもある）。したがって Algebra, Lemma
\ref{algebra-lemma-formally-etale-etale} から結論を得る。
（この含意については、形式的滑らかな射を論じる際に別の証明も与える。）
\end{proof}













\section{写像の無限小変形}
\label{more-morphisms-section-action-by-derivations}

\noindent
本節では、導分を用いて写像を無限小に動かす方法を説明する。
本節を通して、$X$ の肥厚化 $X'$ 上の層を $X$ 上の層とみなせることを用いる。

\begin{lemma}
\label{more-morphisms-lemma-difference-derivation}
$S$ をスキームとする。
$X \subset X'$ と $Y \subset Y'$ を $S$ 上の二つの一次肥厚化とする。
$(a, a'), (b, b') : (X \subset X') \to (Y \subset Y')$ を
$S$ 上の肥厚化の二つの射とする。次を仮定する。
\begin{enumerate}
\item $a = b$ である。
\item 二つの写像
$a^*\mathcal{C}_{Y/Y'} \to \mathcal{C}_{X/X'}$
（Morphisms, Lemma \ref{morphisms-lemma-conormal-functorial}）
は等しい。
\end{enumerate}
このとき写像 $(a')^\sharp - (b')^\sharp$ は
$$
\mathcal{O}_{Y'} \to \mathcal{O}_Y \xrightarrow{D}
a_*\mathcal{C}_{X/X'} \to a_*\mathcal{O}_{X'}
$$
と因子分解する。ここで $D$ は $\mathcal{O}_S$-導分である。
\end{lemma}

\begin{proof}
$Y$ 上ではなく $X$ 上で考える。こうすると引き戻し関手
$a^{-1}$ が完全であるという利点がある。(1) と (2) を用いると、
各行が完全な可換図式
$$
\xymatrix{
0 \ar[r] &
\mathcal{C}_{X/X'} \ar[r] &
\mathcal{O}_{X'} \ar[r] &
\mathcal{O}_X \ar[r] & 0 \\
0 \ar[r] &
a^{-1}\mathcal{C}_{Y/Y'} \ar[r] \ar[u] &
a^{-1}\mathcal{O}_{Y'}
\ar[r] \ar@<1ex>[u]^{(a')^\sharp} \ar@<-1ex>[u]_{(b')^\sharp} &
a^{-1}\mathcal{O}_Y \ar[r] \ar[u] & 0
}
$$
を得る。このような状況では、$\mathcal{O}_S$-代数準同型
$(a')^\sharp$ と $(b')^\sharp$ の差が
$a^{-1}\mathcal{O}_Y$ から $\mathcal{C}_{X/X'}$ への
$\mathcal{O}_S$-導分になるという一般的事実がある。
関手 $a^{-1}$ と $a_*$ の随伴性により、これは
$\mathcal{O}_Y$ から $a_*\mathcal{C}_{X/X'}$ への
$\mathcal{O}_S$-導分と同じものである。詳細の一部は省略する。
\end{proof}

\noindent
上の補題の状況では、$D$ を
\begin{equation}
\label{more-morphisms-equation-D}
% P05-EN-CH37-0039
D = \text{d}_{Y/S} \circ \theta
\end{equation}
と書けることに注意する。ここで $\theta$ は
$\mathcal{O}_Y$-線形写像
$\theta : \Omega_{Y/S} \to a_*\mathcal{C}_{X/X'}$ である。
もちろん、このとき再び随伴性により $\theta$ を
$\mathcal{O}_X$-線形写像
$\theta : a^*\Omega_{Y/S} \to \mathcal{C}_{X/X'}$ とみなせる。

\begin{lemma}
\label{more-morphisms-lemma-action-by-derivations}
$S$ をスキームとする。
$(a, a') : (X \subset X') \to (Y \subset Y')$ を
$S$ 上の一次肥厚化の射とする。
$$
\theta : a^*\Omega_{Y/S} \to \mathcal{C}_{X/X'}
$$
を $\mathcal{O}_X$-線形写像とする。このとき、Lemma
\ref{more-morphisms-lemma-difference-derivation} の (1) と (2) を満たし、
導分 $D$ と $\theta$ が Equation (\ref{more-morphisms-equation-D}) で関係づけられるような
対の射 $(b, b') : (X \subset X') \to (Y \subset Y')$ が一意に存在する。
\end{lemma}

\begin{proof}
$b = a$ とおき、$(b')^\sharp$ を写像
$$
(a')^\sharp + D : a^{-1}\mathcal{O}_{Y'} \to \mathcal{O}_{X'}
$$
と定めればよい。ここで $D$ は Equation (\ref{more-morphisms-equation-D}) のものである。
$(b')^\sharp$ が $\mathcal{O}_S$-代数の層の写像であり、Lemma
\ref{more-morphisms-lemma-difference-derivation} の (1) と (2) が成り立つことの確認は省略する。
Equation (\ref{more-morphisms-equation-D}) は構成により成り立つ。
\end{proof}

\begin{remark}
\label{more-morphisms-remark-action-by-derivations}
仮定と記法は Lemma \ref{more-morphisms-lemma-action-by-derivations} と同じとする。
局所切断 $\theta$ の $a'$ への作用を $\theta \cdot a'$ と表すことがある。
これは、$(a')^\sharp$ と $(\theta \cdot a')^\sharp$ の差が、
Equation (\ref{more-morphisms-equation-D}) により $\theta$ と関係づけられる導分 $D$ であることを
意味するにすぎない。
\end{remark}

\begin{lemma}
\label{more-morphisms-lemma-sheaf}
$S$ をスキームとする。
$X \subset X'$ と $Y \subset Y'$ を $S$ 上の一次肥厚化とする。
射 $a : X \to Y$ と $\mathcal{O}_X$-加群の写像
$A : a^*\mathcal{C}_{Y/Y'} \to \mathcal{C}_{X/X'}$ が与えられているとする。
開部分スキーム $U' \subset X'$ に対し、次を満たす射
$a' : U' \to Y'$ を考える。
\begin{enumerate}
\item $a'$ は $S$ 上の射である。
\item $a'|_U = a|_U$ である。
\item 誘導される写像
$a^*\mathcal{C}_{Y/Y'}|_U \to \mathcal{C}_{X/X'}|_U$
は $A$ の $U$ への制限である。
\end{enumerate}
ここで $U = X \cap U'$ とする。このとき規則
\begin{equation}
\label{more-morphisms-equation-sheaf}
U' \mapsto
\{a' : U' \to Y'\text{ で (1), (2), (3) を満たすもの}\}
\end{equation}
は $X'$ 上の集合の層を定める。
\end{lemma}

\begin{proof}
% P05-EN-CH37-0040
補題の規則を $\mathcal{F}$ と表す。
$V' \subset U' \subset X'$ に対する $\mathcal{F}$ の制限写像
$\mathcal{F}(U') \to \mathcal{F}(V')$ は、まさに制限写像
$a' \mapsto a'|_{V'}$ である。この定義の下で、射は局所的に定まるので
$\mathcal{F}$ が層であることは明らかである。
\end{proof}

\noindent
次の補題では、$X$ 上の層と $X$ の任意の肥厚化上の層とを同一視する。

\begin{lemma}
\label{more-morphisms-lemma-action-sheaf}
記法と仮定は Lemma \ref{more-morphisms-lemma-sheaf} と同じとする。
層
$$
\SheafHom_{\mathcal{O}_X}(a^*\Omega_{Y/S}, \mathcal{C}_{X/X'})
$$
は層 (\ref{more-morphisms-equation-sheaf}) に作用する。さらに、層
(\ref{more-morphisms-equation-sheaf}) が切断をもつ任意の開部分スキーム
$U' \subset X'$ 上で、この作用は単純推移的である。
\end{lemma}

\begin{proof}
これは Lemmas \ref{more-morphisms-lemma-difference-derivation},
\ref{more-morphisms-lemma-action-by-derivations},
および \ref{more-morphisms-lemma-sheaf} を合わせたものである。
\end{proof}

\begin{remark}
\label{more-morphisms-remark-special-case}
Lemmas \ref{more-morphisms-lemma-difference-derivation},
\ref{more-morphisms-lemma-action-by-derivations},
\ref{more-morphisms-lemma-sheaf}, および
\ref{more-morphisms-lemma-action-sheaf}
の特別な場合として $Y = Y'$ の場合がある。この場合、写像 $A$ は常に零である。
Lemma \ref{more-morphisms-lemma-sheaf} の層は単に規則
$$
U' \mapsto
\{a' : U' \to Y\text{ は }S\text{ 上で } a'|_U = a|_U\text{ を満たす}\}
$$
で与えられ、層
$\SheafHom_{\mathcal{O}_X}(a^*\Omega_{Y/S}, \mathcal{C}_{X/X'})$
がこれに作用する。
\end{remark}

\begin{remark}
\label{more-morphisms-remark-another-special-case}
Lemmas \ref{more-morphisms-lemma-difference-derivation},
\ref{more-morphisms-lemma-action-by-derivations},
\ref{more-morphisms-lemma-sheaf}, および
\ref{more-morphisms-lemma-action-sheaf}
のもう一つの特別な場合として、$S$ 自身が肥厚化 $Z \subset Z' = S$ であり、
$Y = Z \times_{Z'} Y'$ である場合がある。図式は
$$
\xymatrix{
(X \subset X') \ar@{..>}[rr]_{(a, ?)} \ar[rd]_{(g, g')} & &
(Y \subset Y') \ar[ld]^{(h, h')} \\
& (Z \subset Z')
}
$$
この場合、写像 $A : a^*\mathcal{C}_{Y/Y'} \to \mathcal{C}_{X/X'}$ は
$a$ によって決まる。実際、写像
$h^*\mathcal{C}_{Z/Z'} \to \mathcal{C}_{Y/Y'}$ は全射である
（$Y = Z \times_{Z'} Y'$ と仮定したためである）。したがって引き戻し
$g^*\mathcal{C}_{Z/Z'} = a^*h^*\mathcal{C}_{Z/Z'} \to
a^*\mathcal{C}_{Y/Y'}$ は全射であり、合成
$g^*\mathcal{C}_{Z/Z'} \to a^*\mathcal{C}_{Y/Y'} \to \mathcal{C}_{X/X'}$
は $g'$ によって誘導される標準写像でなければならない。したがって
Lemma \ref{more-morphisms-lemma-sheaf} の層は単に規則
$$
U' \mapsto
\{a' : U' \to Y'\text{ は }Z'\text{ 上で } a'|_U = a|_U\text{ を満たす}\}
$$
で与えられ、層
$\SheafHom_{\mathcal{O}_X}(a^*\Omega_{Y/Z}, \mathcal{C}_{X/X'})$
がこれに作用する。
\end{remark}

\begin{lemma}
\label{more-morphisms-lemma-omega-deformation}
$S$ をスキームとする。$X \subset X'$ を $S$ 上の一次肥厚化とし、
$Y$ を $S$ 上のスキームとする。$a', b' : X' \to Y$ を、
$a = a'|_X = b'|_X$ を満たす $S$ 上の二つの射とする。これにより可換図式
$$
\xymatrix{
X \ar[r] \ar[d]_a & X' \ar[d]^{(b', a')} \\
Y \ar[r]^-{\Delta_{Y/S}} & Y \times_S Y
}
$$
を得る。横の矢印は余法層がそれぞれ
$\mathcal{C}_{X/X'}$ と $\Omega_{Y/S}$ である埋め込みなので、
Morphisms, Lemma \ref{morphisms-lemma-conormal-functorial} により写像
$\theta : a^*\Omega_{Y/S} \to \mathcal{C}_{X/X'}$ を得る。この $\theta$ と
Lemma \ref{more-morphisms-lemma-difference-derivation} の導分 $D$ は
Equation (\ref{more-morphisms-equation-D}) により関係づけられる。
\end{lemma}

\begin{proof}
省略する。ヒント：等式はアフィン開部分スキーム上で確認でき、そこでは
次の計算から従う。$f$ を $\mathcal{O}_Y$ の局所切断とすると、
$1 \otimes f - f \otimes 1$ は $\text{d}_{Y/S}(f)$ に対応する
$\mathcal{C}_{Y/(Y \times_S Y)}$ の局所切断である。これは
$\mathcal{C}_{X/X'}$ の局所切断
$(a')^\sharp(f) - (b')^\sharp(f) = D(f)$ に写る。言い換えると、
$\theta(\text{d}_{Y/S}(f)) = D(f)$ である。
\end{proof}

\noindent
後で用いるため、Lemma \ref{more-morphisms-lemma-action-by-derivations} の構成が
エタール局所化と両立することを大まかに述べる結果が必要である。

\begin{lemma}
\label{more-morphisms-lemma-sheaf-differentials-etale-localization}
可換図式
$$
\xymatrix{
X_1 \ar[d] & X_2 \ar[l]^f \ar[d] \\
S_1 & S_2 \ar[l]
}
$$
を考え、$X_2 \to X_1$ と $S_2 \to S_1$ はエタールであるとする。
このとき Morphisms, Lemma
\ref{morphisms-lemma-functoriality-differentials} の写像
$c_f : f^*\Omega_{X_1/S_1} \to \Omega_{X_2/S_2}$ は同型である。
\end{lemma}

\begin{proof}
エタール射 $U \to V$ は $\Omega_{U/V} = 0$ を満たす滑らかな射であることを
思い出す。これを用いると、Morphisms, Lemma
\ref{morphisms-lemma-triangle-differentials} から
$\Omega_{X_2/S_2} = \Omega_{X_2/S_1}$ が従い、Morphisms, Lemma
\ref{morphisms-lemma-triangle-differentials-smooth} から写像
$f^*\Omega_{X_1/S_1} \to \Omega_{X_2/S_1}$ が同型であることが従う
（ここで射 $f$ は $S_1$ 上の射とみなす）。よって補題が従う。
\end{proof}

\begin{lemma}
\label{more-morphisms-lemma-action-by-derivations-etale-localization}
一次肥厚化の可換図式
$$
\vcenter{
\xymatrix{
(T_2 \subset T_2') \ar[d]_{(h, h')} \ar[rr]_{(a_2, a_2')} & &
(X_2 \subset X_2') \ar[d]^{(f, f')} \\
(T_1 \subset T_1') \ar[rr]^{(a_1, a_1')} & &
(X_1 \subset X_1')
}
}
\quad
\begin{matrix}
\text{およびスキームの} \\
\text{可換図式}
\end{matrix}
\quad
\vcenter{
\xymatrix{
X_2' \ar[r] \ar[d] & S_2 \ar[d] \\
X_1' \ar[r] & S_1
}
}
$$
を考え、$X_2 \to X_1$ と $S_2 \to S_1$ はエタールであるとする。
任意の $\mathcal{O}_{T_1}$-線形写像
$\theta_1 : a_1^*\Omega_{X_1/S_1} \to \mathcal{C}_{T_1/T'_1}$ に対し、
$\theta_2$ を合成
$$
\xymatrix{
a_2^*\Omega_{X_2/S_2} \ar@{=}[r] &
h^*a_1^*\Omega_{X_1/S_1} \ar[r]^-{h^*\theta_1} &
h^*\mathcal{C}_{T_1/T'_1} \ar[r] &
\mathcal{C}_{T_2/T'_2}
}
$$
とする（等号は証明で説明する）。このとき図式
$$
\xymatrix{
T_2' \ar[rr]_{\theta_2 \cdot a_2'} \ar[d] & & X'_2 \ar[d] \\
T_1' \ar[rr]^{\theta_1 \cdot a_1'} & & X'_1
}
$$
は可換である。ここで作用 $\theta_2 \cdot a_2'$ と
$\theta_1 \cdot a_1'$ は Remark \ref{more-morphisms-remark-action-by-derivations}
のものである。
\end{lemma}

\begin{proof}
等号は Lemma \ref{more-morphisms-lemma-sheaf-differentials-etale-localization} の同一視
$f^*\Omega_{X_1/S_1} = \Omega_{X_2/S_2}$ から来る。すなわち、これを用いると
$a_2^*\Omega_{X_2/S_2} = a_2^*f^*\Omega_{X_1/S_1} =
h^*a_1^*\Omega_{X_1/S_1}$ である。なぜなら
$f \circ a_2 = a_1 \circ h$ だからである。以上を踏まえると、図式の可換性は
アフィン開部分スキーム上で確認できる。したがって、冒頭の大きな図式に現れる
スキームはアフィンであると仮定してよい。こうして可換図式
$$
\vcenter{
\xymatrix{
(B'_2, I_2) & & (A'_2, J_2) \ar[ll]^{a_2'} \\
(B'_1, I_1) \ar[u]^{h'} & & (A'_1, J_1) \ar[ll]_{a_1'} \ar[u]_{f'}
}
}
\quad\text{および}\quad
\vcenter{
\xymatrix{
A'_2 & & R_2 \ar[ll] \\
A'_1 \ar[u] & & R_1 \ar[ll] \ar[u]
}
}
$$
を得る。この記法は、$I_1, I_2, J_1, J_2$ が平方零イデアルであり、
対の写像がイデアルをイデアルへ送る環準同型であることを表す。
$A_1 = A'_1/J_1$, $A_2 = A'_2/J_2$, $B_1 = B'_1/I_1$, および
$B_2 = B'_2/I_2$ とおく。写像
$$
A_2 \otimes_{A_1} \Omega_{A_1/R_1} \longrightarrow \Omega_{A_2/R_2}
$$
が同型であることが与えられている。このとき
$\theta_1 : B_1 \otimes_{A_1} \Omega_{A_1/R_1} \to I_1$ は
$B_1$-線形である。これにより $R_1$-導分
$D_1 = \theta_1 \circ \text{d}_{A_1/R_1} : A_1 \to I_1$ を得る。
同様に、
$\theta_2 : B_2 \otimes_{A_2} \Omega_{A_2/R_2} \to I_2$ から
$R_2$-導分
$D_2 = \theta_2 \circ \text{d}_{A_2/R_2} : A_2 \to I_2$ を得る。
$\theta_2$ の構成から $\theta_1$ と $\theta_2$ の間に次の両立性が従う。
任意の $x \in A_1$ に対して
$$
h'(D_1(x)) = D_2(f'(x))
$$
$I_2$ の元として等しい。$D_1$ は、写像
% P05-EN-CH37-0041
$A'_1 \to A_1 \xrightarrow{D_1} I_1 \to B_1$ を用いて
$A'_1 \to B'_1$ なる写像とみなせる。同様に $D_2$ を
$A'_2 \to B'_2$ なる写像とみなせる。このとき表示した等式は
$x \in A'_1$ に対して成り立つ。
Lemma \ref{more-morphisms-lemma-action-by-derivations} および
Remark \ref{more-morphisms-remark-action-by-derivations} における作用の構成により、
$\theta_1 \cdot a_1'$ は環準同型 $a_1' + D_1 : A'_1 \to B'_1$ に対応し、
$\theta_2 \cdot a_2'$ は環準同型 $a_2' + D_2 : A'_2 \to B'_2$ に対応する。
表示した等式により
$h' \circ (a_1' + D_1) = (a_2' + D_2) \circ f'$
を得る。これが求めるものである。
\end{proof}

\begin{remark}
\label{more-morphisms-remark-tiny-improvement}
Lemma \ref{more-morphisms-lemma-action-by-derivations-etale-localization} は
次のように改良できる。Lemma
\ref{more-morphisms-lemma-action-by-derivations-etale-localization} と同じ可換図式があるが、
$X_2 \to X_1$ と $S_2 \to S_1$ がエタールであるとは仮定しないとする。
さらに、
$\theta_1 : a_1^*\Omega_{X_1/S_1} \to \mathcal{C}_{T_1/T'_1}$
および
$\theta_2 : a_2^*\Omega_{X_2/S_2} \to \mathcal{C}_{T_2/T'_2}$
があり、
$$
\xymatrix{
f_*\mathcal{O}_{X_2} \ar[rr]_{f_*D_2} & &
f_*a_{2, *}\mathcal{C}_{T_2/T_2'} \\
\mathcal{O}_{X_1} \ar[rr]^{D_1} \ar[u]^{f^\sharp} & &
a_{1, *}\mathcal{C}_{T_1/T_1'} \ar[u]_{\text{次により誘導：}(h')^\sharp}
}
$$
が可換であるとする。ここで $D_i$ は Equation (\ref{more-morphisms-equation-D}) と同様に
$\theta_i$ に対応する。このとき Lemma
\ref{more-morphisms-lemma-action-by-derivations-etale-localization} の結論が成り立つ。
$X_2 \to X_1$ と $S_2 \to S_1$ の両方がエタールであるという条件の重要性は、
$\theta_1$ から $\theta_2$ を構成できるようにする点にある。
\end{remark}








\section{スキームの無限小変形}
\label{more-morphisms-section-deform}

\noindent
次の簡単な補題は、写像の無限小変形が平坦かどうかを確認する際に
便利な道具となることが多い。

\begin{lemma}
\label{more-morphisms-lemma-deform}
$(f, f') : (X \subset X') \to (S \subset S')$ を一次肥厚化の射とする。
$f$ は平坦であると仮定する。このとき次は同値である。
\begin{enumerate}
\item $f'$ は平坦であり、$X = S \times_{S'} X'$ である。
\item 標準写像 $f^*\mathcal{C}_{S/S'} \to \mathcal{C}_{X/X'}$ は同型である。
\end{enumerate}
\end{lemma}

\begin{proof}
問題は $X'$ 上で局所的なので、$X, X', S, S'$ はアフィンスキームであると
仮定してよい。$S' = \Spec(A')$, $X' = \Spec(B')$,
$S = \Spec(A)$, $X = \Spec(B)$ とし、ある平方零イデアルに対して
$A = A'/I$ および $B = B'/J$ であるとする。このとき可換図式
$$
\xymatrix{
0 \ar[r] &
J \ar[r] &
B' \ar[r] &
B \ar[r] & 0 \\
0 \ar[r] &
I \ar[r] \ar[u] &
A' \ar[r] \ar[u] &
A \ar[r] \ar[u] & 0
}
$$
を得、その各行は完全である。補題の標準写像は写像
$$
I \otimes_A B = I \otimes_{A'} B' \longrightarrow J.
$$
である。$f$ が平坦であるという仮定は、$A \to B$ が平坦であることを意味する。

\medskip\noindent
(1) を仮定する。このとき $A' \to B'$ は平坦であり、$J = IB'$ である。
平坦性から $\text{Tor}_1^{A'}(B', A) = 0$ が従う
（Algebra, Lemma \ref{algebra-lemma-characterize-flat} を参照）。
これは $I \otimes_{A'} B' \to B'$ が単射であることを意味する
（Algebra, Remark \ref{algebra-remark-Tor-ring-mod-ideal} を参照）。
したがって $I \otimes_A B \to J$ は同型である。

\medskip\noindent
(2) を仮定する。このとき $J = IB'$ が従うので
$X = S \times_{S'} X'$ である。さらに、前段落の含意を逆にたどると
$\text{Tor}_1^{A'}(B', A'/I) = 0$ を得る。したがって Algebra, Lemma
\ref{algebra-lemma-what-does-it-mean} により $B'$ は $A'$ 上平坦である。
\end{proof}

\noindent
次の補題は「ファイバーによる平坦性判定法」の「冪零」版である。
Section \ref{more-morphisms-section-criterion-flat-fibres} を参照せよ。

\begin{lemma}
\label{more-morphisms-lemma-flatness-morphism-thickenings}
肥厚化の可換図式
$$
\xymatrix{
(X \subset X') \ar[rr]_{(f, f')} \ar[rd] & & (Y \subset Y') \ar[ld] \\
& (S \subset S')
}
$$
を考える。次を仮定する。
\begin{enumerate}
\item $X'$ は $S'$ 上平坦である。
\item $f$ は平坦である。
\item $S \subset S'$ は有限次肥厚化である。
\item $X = S \times_{S'} X'$ かつ $Y = S \times_{S'} Y'$ である。
\end{enumerate}
このとき $f'$ は平坦であり、$Y'$ は $f'$ の像に属するすべての点で
$S'$ 上平坦である。
\end{lemma}

\begin{proof}
Algebra, Lemma \ref{algebra-lemma-criterion-flatness-fibre-nilpotent}
の直ちに分かる帰結である。
\end{proof}

\noindent
スキームの射の多くの性質は平坦変形で保たれる。

\begin{lemma}
\label{more-morphisms-lemma-deform-property}
肥厚化の可換図式
$$
\xymatrix{
(X \subset X') \ar[rr]_{(f, f')} \ar[rd] & & (Y \subset Y') \ar[ld] \\
& (S \subset S')
}
$$
を考える。$S \subset S'$ は有限次肥厚化、$X'$ は $S'$ 上平坦、
$X = S \times_{S'} X'$, かつ $Y = S \times_{S'} Y'$ であると仮定する。
このとき次が成り立つ。
\begin{enumerate}
\item $f$ が平坦であることと $f'$ が平坦であることは同値である。
\label{more-morphisms-item-flat}
\item $f$ が同型であることと $f'$ が同型であることは同値である。
\label{more-morphisms-item-isomorphism}
\item $f$ が開埋め込みであることと $f'$ が開埋め込みであることは同値である。
\label{more-morphisms-item-open-immersion}
\item $f$ が準コンパクトであることと $f'$ が準コンパクトであることは同値である。
\label{more-morphisms-item-quasi-compact}
\item $f$ が普遍閉であることと $f'$ が普遍閉であることは同値である。
\label{more-morphisms-item-universally-closed}
\item $f$ が（準）分離的であることと $f'$ が（準）分離的であることは同値である。
\label{more-morphisms-item-separated}
\item $f$ がモノ射であることと $f'$ がモノ射であることは同値である。
\label{more-morphisms-item-monomorphism}
\item $f$ が全射であることと $f'$ が全射であることは同値である。
\label{more-morphisms-item-surjective}
\item $f$ が普遍単射であることと $f'$ が普遍単射であることは同値である。
\label{more-morphisms-item-universally-injective}
\item $f$ がアフィンであることと $f'$ がアフィンであることは同値である。
\label{more-morphisms-item-affine}
\item
\label{more-morphisms-item-finite-type}
$f$ が局所有限型であることと $f'$ が局所有限型であることは同値である。
\item $f$ が局所準有限であることと $f'$ が局所準有限であることは同値である。
\label{more-morphisms-item-quasi-finite}
\item
\label{more-morphisms-item-finite-presentation}
$f$ が局所有限表示であることと $f'$ が局所有限表示であることは同値である。
\item
\label{more-morphisms-item-relative-dimension-d}
$f$ が相対次元 $d$ の局所有限型であることと
$f'$ が相対次元 $d$ の局所有限型であることは同値である。
\item $f$ が普遍開であることと $f'$ が普遍開であることは同値である。
\label{more-morphisms-item-universally-open}
\item $f$ がシントミックであることと $f'$ がシントミックであることは同値である。
\label{more-morphisms-item-syntomic}
\item $f$ が滑らかであることと $f'$ が滑らかであることは同値である。
\label{more-morphisms-item-smooth}
\item $f$ が不分岐であることと $f'$ が不分岐であることは同値である。
\label{more-morphisms-item-unramified}
\item $f$ がエタールであることと $f'$ がエタールであることは同値である。
\label{more-morphisms-item-etale}
\item $f$ が固有であることと $f'$ が固有であることは同値である。
\label{more-morphisms-item-proper}
\item $f$ が整であることと $f'$ が整であることは同値である。
\label{more-morphisms-item-integral}
\item $f$ が有限であることと $f'$ が有限であることは同値である。
\label{more-morphisms-item-finite}
\item
\label{more-morphisms-item-finite-locally-free}
$f$ が（階数 $d$ の）有限局所自由であることと
$f'$ が（階数 $d$ の）有限局所自由であることは同値である。
% P05-EN-CH37-0042
\item ここにさらに追加する。
\end{enumerate}
\end{lemma}

\begin{proof}
仮定された $X$ と $Y$ の条件は、$f$ が $f'$ の $X \to X'$ による
基底変換であることを意味する。上の (1) -- (23) に挙げた性質
$\mathcal{P}$ はすべて基底変換で安定である。したがって $f'$ が性質
$\mathcal{P}$ をもてば $f$ もそれをもつ。次を参照せよ。
Schemes, Lemmas \ref{schemes-lemma-base-change-immersion},
\ref{schemes-lemma-quasi-compact-preserved-base-change},
\ref{schemes-lemma-separated-permanence}, および
\ref{schemes-lemma-base-change-monomorphism}
ならびに
Morphisms, Lemmas
\ref{morphisms-lemma-base-change-surjective},
\ref{morphisms-lemma-base-change-universally-injective},
\ref{morphisms-lemma-base-change-affine},
\ref{morphisms-lemma-base-change-finite-type},
\ref{morphisms-lemma-base-change-quasi-finite},
\ref{morphisms-lemma-base-change-finite-presentation},
\ref{morphisms-lemma-base-change-relative-dimension-d},
\ref{morphisms-lemma-base-change-syntomic},
\ref{morphisms-lemma-base-change-smooth},
\ref{morphisms-lemma-base-change-unramified},
\ref{morphisms-lemma-base-change-etale},
\ref{morphisms-lemma-base-change-proper},
\ref{morphisms-lemma-base-change-finite}, および
\ref{morphisms-lemma-base-change-finite-locally-free}.

\medskip\noindent
したがって各場合で興味深い向きは、$f$ がその性質をもつと仮定して
$f'$ もそれをもつことを導く向きである。肥厚化の次数に関する帰納法により、
$S \subset S'$ は一次肥厚化であると仮定してよい。Definition
\ref{more-morphisms-definition-thickening} の直後の議論を参照せよ。以下の議論で断りなく用いる
一般的注意を二つ述べる。
(I) $W' \subset S'$ をアフィン開部分スキームとし、
$U' \subset X'$ と $V' \subset Y'$ を $W'$ 上にあり
$f'(U') \subset V'$ を満たすアフィン開部分スキームとする。
$W' = \Spec(R')$ とおき、
% P05-EN-CH37-0043
$I \subset R'$ を閉部分スキーム $W' \cap S$ を定義するイデアルと表す。
$U' = \Spec(B')$ および $V' = \Spec(A')$ とする。このとき可換図式
$$
\xymatrix{
0 \ar[r] &
IB' \ar[r] &
B' \ar[r] &
B \ar[r] & 0 \\
0 \ar[r] &
IA' \ar[r] \ar[u] &
A' \ar[r] \ar[u] &
A \ar[r] \ar[u] & 0
}
$$
を得、その各行は完全である。さらに $IB' \cong I \otimes_R B$ である。
Lemma \ref{more-morphisms-lemma-deform} の証明を参照せよ。
(II) 射 $X \to X'$ と $Y \to Y'$ は普遍同相である。したがって写像
$f$ と $f'$ の位相は（任意の基底変換の後でも）同一である。
(III) $f$ が平坦ならば $f'$ は平坦であり、$Y' \to S'$ は $f'$ の像に属する
すべての点で平坦である。Lemma
\ref{more-morphisms-lemma-flatness-morphism-thickenings} を参照せよ。

\medskip\noindent
(\ref{more-morphisms-item-flat}) について。これは一般的注意 (III) である。

\medskip\noindent
(\ref{more-morphisms-item-isomorphism}) について。$f$ は同型であると仮定する。
(III) により $Y' \to S'$ は平坦である。アフィン開部分スキーム
$V' \subset Y'$ を選び、$U' = (f')^{-1}(V')$ とおく。このとき
$V = Y \cap V'$ はアフィンであるから
$V \cong f^{-1}(V) = U = Y \times_{Y'} U'$ はアフィンである。Lemma
\ref{more-morphisms-lemma-thickening-affine-scheme} により $U'$ はアフィンである。
したがって一般的注意 (I) のような図式を得る。さらに
$R' \to A'$ が平坦なので $IA \cong I \otimes_R A$ である。よって
$IB' \cong I \otimes_R B \cong I \otimes_R A \cong IA'$
かつ $A \cong B$ である。上の図式の各行の完全性から
$A' \cong B'$、すなわち $U' \cong V'$ を得る。したがって $f'$ は同型である。

\medskip\noindent
(\ref{more-morphisms-item-open-immersion}) について。$f$ は開埋め込みであると仮定する。
このとき $f$ は $X$ と開部分スキーム $V \subset Y$ との同型である。
$V' \subset Y'$ を、その台となる位相空間が $V$ である開部分スキームとする。
すると $f'$ は $X'$ から $V'$ への写像であり、
(\ref{more-morphisms-item-isomorphism}) により同型である。したがって $f'$ は開埋め込みである。

\medskip\noindent
(\ref{more-morphisms-item-quasi-compact}) について。注意 (II) から直ちに従う。
より一般的な主張については Lemma
\ref{more-morphisms-lemma-thicken-property-morphisms} も参照せよ。

\medskip\noindent
(\ref{more-morphisms-item-universally-closed}) について。注意 (II) から直ちに従う。
より一般的な主張については Lemma
\ref{more-morphisms-lemma-thicken-property-morphisms} も参照せよ。

\medskip\noindent
(\ref{more-morphisms-item-separated}) について。
$X \times_Y X = Y \times_{Y'} (X' \times_{Y'} X')$ なので
$X' \times_{Y'} X'$ は $X \times_Y X$ の肥厚化であることに注意する。
したがって写像 $\Delta_{X/Y}$ と $\Delta_{X'/Y'}$ の位相は一致し、
結論を得る。より一般的な主張については Lemma
\ref{more-morphisms-lemma-thicken-property-morphisms} も参照せよ。

\medskip\noindent
(\ref{more-morphisms-item-monomorphism}) について。$f$ はモノ射であると仮定する。
対角射 $\Delta_{X'/Y'} : X' \to X' \times_{Y'} X'$ を考える。
$\Delta_{X'/Y'}$ の $S \to S'$ による基底変換は $\Delta_{X/Y}$ であり、
仮定により同型である。(\ref{more-morphisms-item-isomorphism}) により
$\Delta_{X'/Y'}$ は同型であると結論する。

\medskip\noindent
(\ref{more-morphisms-item-surjective}) について。これは明らかである。より一般的な主張については
Lemma \ref{more-morphisms-lemma-thicken-property-morphisms} も参照せよ。

\medskip\noindent
(\ref{more-morphisms-item-universally-injective}) について。注意 (II) から直ちに従う。
より一般的な主張については Lemma
\ref{more-morphisms-lemma-thicken-property-morphisms} も参照せよ。

\medskip\noindent
(\ref{more-morphisms-item-affine}) について。$f$ はアフィンであると仮定する。
アフィン開部分スキーム $V' \subset Y'$ を選び、
$U' = (f')^{-1}(V')$ とおく。このとき $V = Y \cap V'$ はアフィンなので
$U = Y \times_{Y'} U'$ もアフィンである。Lemma
\ref{more-morphisms-lemma-thickening-affine-scheme} により $U'$ はアフィンである。
したがって $f'$ はアフィンである。より一般的な主張については Lemma
\ref{more-morphisms-lemma-thicken-property-morphisms} も参照せよ。

\medskip\noindent
(\ref{more-morphisms-item-finite-type}) について。注意 (I) により、$A \to B$ が有限型ならば
$A' \to B'$ も有限型であることを示す問題に帰着する。
$x_1, \ldots, x_n \in B'$ を、その $B$ における像が $A$-代数として
$B$ を生成する元とする。
% P05-EN-CH37-0044
このとき $A'[x_1, \ldots, x_n] \to B$ は全射である。なぜなら
$A'[x_1, \ldots, x_n] \to B$ と
$I \otimes_R A[x_1, \ldots, x_n] \to I \otimes_R B$ がともに全射だからである。
より一般的な主張については Lemma
\ref{more-morphisms-lemma-thicken-property-morphisms-cartesian} も参照せよ。

\medskip\noindent
(\ref{more-morphisms-item-quasi-finite}) について。(\ref{more-morphisms-item-finite-type}) と、
有限型射の準有限性はファイバー上で確認できるという事実から従う。
Morphisms, Lemma \ref{morphisms-lemma-quasi-finite-at-point-characterize}
を参照せよ。より一般的な主張については Lemma
\ref{more-morphisms-lemma-thicken-property-morphisms-cartesian} も参照せよ。

\medskip\noindent
(\ref{more-morphisms-item-finite-presentation}) について。注意 (I) により、
$A \to B$ が有限表示ならば $A' \to B'$ も有限表示であることを示す問題に帰着する。
(\ref{more-morphisms-item-finite-type}) により、あるイデアル $K'$ に対して
$B' = A'[x_1, \ldots, x_n]/K'$ であると仮定してよい。短完全列
$$
0 \to K' \to A'[x_1, \ldots, x_n] \to B' \to 0
$$
を得る。$B'$ は $R'$ 上平坦なので、$K' \otimes_{R'} R$ は全射
$A[x_1, \ldots, x_n] \to B$ の核である。$A \to B$ に関する仮定により、
$A[x_1, \ldots, x_n]$ における像がこの核を生成する有限個の元
$f'_1, \ldots, f'_m \in K'$ が存在する。$I$ は冪零なので、Nakayama の補題により
$f'_1, \ldots, f'_m$ は $K'$ を生成する。Algebra, Lemma
\ref{algebra-lemma-NAK} を参照せよ。

\medskip\noindent
(\ref{more-morphisms-item-relative-dimension-d}) について。(\ref{more-morphisms-item-finite-type}) と
一般的注意 (II) から従う。より一般的な主張については Lemma
\ref{more-morphisms-lemma-thicken-property-morphisms-cartesian} も参照せよ。

\medskip\noindent
(\ref{more-morphisms-item-universally-open}) について。一般的注意 (II) から直ちに従う。
より一般的な主張については Lemma
\ref{more-morphisms-lemma-thicken-property-morphisms} も参照せよ。

\medskip\noindent
(\ref{more-morphisms-item-syntomic}) について。$f$ はシントミックであると仮定する。
(\ref{more-morphisms-item-finite-presentation}) により $f'$ は局所有限表示であり、
一般的注意 (III) により $f'$ は平坦であり、$f'$ のファイバーは
$f$ のファイバーである。したがって Morphisms, Lemma
\ref{morphisms-lemma-syntomic-flat-fibres} により $f'$ はシントミックである。

\medskip\noindent
(\ref{more-morphisms-item-smooth}) について。$f$ は滑らかであると仮定する。
(\ref{more-morphisms-item-finite-presentation}) により $f'$ は局所有限表示であり、
一般的注意 (III) により $f'$ は平坦であり、$f'$ のファイバーは
$f$ のファイバーである。したがって Morphisms, Lemma
\ref{morphisms-lemma-smooth-flat-smooth-fibres} により $f'$ は滑らかである。

\medskip\noindent
(\ref{more-morphisms-item-unramified}) について。$f$ は不分岐であると仮定する。
(\ref{more-morphisms-item-finite-type}) により $f'$ は局所有限型であり、
$f'$ のファイバーは $f$ のファイバーである。したがって Morphisms, Lemma
\ref{morphisms-lemma-unramified-etale-fibres} により $f'$ は不分岐である。
より一般的な主張については Lemma
\ref{more-morphisms-lemma-thicken-property-morphisms-cartesian} も参照せよ。

\medskip\noindent
(\ref{more-morphisms-item-etale}) について。$f$ はエタールであると仮定する。
(\ref{more-morphisms-item-finite-presentation}) により $f'$ は局所有限表示であり、
一般的注意 (III) により $f'$ は平坦であり、$f'$ のファイバーは
$f$ のファイバーである。したがって Morphisms, Lemma
\ref{morphisms-lemma-etale-flat-etale-fibres} により $f'$ はエタールである。

\medskip\noindent
(\ref{more-morphisms-item-proper}) について。これは (\ref{more-morphisms-item-separated}),
(\ref{more-morphisms-item-finite-type}), (\ref{more-morphisms-item-quasi-compact}), および
(\ref{more-morphisms-item-universally-closed}) を合わせると従う。より一般的な主張については
Lemma \ref{more-morphisms-lemma-thicken-property-morphisms-cartesian} も参照せよ。

\medskip\noindent
(\ref{more-morphisms-item-integral}) について。(\ref{more-morphisms-item-universally-closed}) と
(\ref{more-morphisms-item-affine}) を Morphisms, Lemma
\ref{morphisms-lemma-integral-universally-closed} と合わせる。
より一般的な主張については Lemma
\ref{more-morphisms-lemma-thicken-property-morphisms} も参照せよ。

\medskip\noindent
(\ref{more-morphisms-item-finite}) について。(\ref{more-morphisms-item-integral}) と
(\ref{more-morphisms-item-finite-type}) を Morphisms, Lemma
\ref{morphisms-lemma-finite-integral} と合わせる。より一般的な主張については
Lemma \ref{more-morphisms-lemma-thicken-property-morphisms-cartesian} も参照せよ。

\medskip\noindent
(\ref{more-morphisms-item-finite-locally-free}) について。$f$ は有限局所自由であると仮定する。
(\ref{more-morphisms-item-finite}) により $f'$ は有限であり、一般的注意 (III) により
$f'$ は平坦であり、(\ref{more-morphisms-item-finite-presentation}) により $f'$ は
局所有限表示である。したがって Morphisms, Lemma
\ref{morphisms-lemma-finite-flat} により $f'$ は有限局所自由である。
\end{proof}

\noindent
次の補題は「ファイバーによる平坦性判定法」の「局所冪零」版である。
Section \ref{more-morphisms-section-criterion-flat-fibres} を参照せよ。

\begin{lemma}
\label{more-morphisms-lemma-flatness-morphism-thickenings-fp-over-ft}
肥厚化の可換図式
$$
\xymatrix{
(X \subset X') \ar[rr]_{(f, f')} \ar[rd] & & (Y \subset Y') \ar[ld] \\
& (S \subset S')
}
$$
を考える。次を仮定する。
\begin{enumerate}
\item $Y' \to S'$ は局所有限型である。
\item $X' \to S'$ は平坦かつ局所有限表示である。
\item $f$ は平坦である。
\item $X = S \times_{S'} X'$ かつ $Y = S \times_{S'} Y'$ である。
\end{enumerate}
このとき $f'$ は平坦であり、$f'$ の像に属するすべての
$y' \in Y'$ に対して、局所環 $\mathcal{O}_{Y', y'}$ は
$\mathcal{O}_{S', s'}$ 上平坦かつ本質的有限表示である。
\end{lemma}

\begin{proof}
Algebra, Lemma
\ref{algebra-lemma-criterion-flatness-fibre-locally-nilpotent}
の直ちに分かる帰結である。
\end{proof}

\noindent
スキームの射の多くの性質は、上の補題のような平坦変形で保たれる。

\begin{lemma}
\label{more-morphisms-lemma-deform-property-fp-over-ft}
肥厚化の可換図式
$$
\xymatrix{
(X \subset X') \ar[rr]_{(f, f')} \ar[rd] & & (Y \subset Y') \ar[ld] \\
& (S \subset S')
}
$$
を考える。$Y' \to S'$ は局所有限型、$X' \to S'$ は平坦かつ局所有限表示、
$X = S \times_{S'} X'$, かつ $Y = S \times_{S'} Y'$ であると仮定する。
このとき次が成り立つ。
\begin{enumerate}
\item $f$ が平坦であることと $f'$ が平坦であることは同値である。
\label{more-morphisms-item-flat-fp-over-ft}
\item $f$ が同型であることと $f'$ が同型であることは同値である。
\label{more-morphisms-item-isomorphism-fp-over-ft}
\item $f$ が開埋め込みであることと $f'$ が開埋め込みであることは同値である。
\label{more-morphisms-item-open-immersion-fp-over-ft}
\item $f$ が準コンパクトであることと $f'$ が準コンパクトであることは同値である。
\label{more-morphisms-item-quasi-compact-fp-over-ft}
\item $f$ が普遍閉であることと $f'$ が普遍閉であることは同値である。
\label{more-morphisms-item-universally-closed-fp-over-ft}
\item $f$ が（準）分離的であることと $f'$ が（準）分離的であることは同値である。
\label{more-morphisms-item-separated-fp-over-ft}
\item $f$ がモノ射であることと $f'$ がモノ射であることは同値である。
\label{more-morphisms-item-monomorphism-fp-over-ft}
\item $f$ が全射であることと $f'$ が全射であることは同値である。
\label{more-morphisms-item-surjective-fp-over-ft}
\item $f$ が普遍単射であることと $f'$ が普遍単射であることは同値である。
\label{more-morphisms-item-universally-injective-fp-over-ft}
\item $f$ がアフィンであることと $f'$ がアフィンであることは同値である。
\label{more-morphisms-item-affine-fp-over-ft}
\item $f$ が局所準有限であることと $f'$ が局所準有限であることは同値である。
\label{more-morphisms-item-quasi-finite-fp-over-ft}
\item
\label{more-morphisms-item-relative-dimension-d-fp-over-ft}
$f$ が相対次元 $d$ の局所有限型であることと
$f'$ が相対次元 $d$ の局所有限型であることは同値である。
\item $f$ が普遍開であることと $f'$ が普遍開であることは同値である。
\label{more-morphisms-item-universally-open-fp-over-ft}
\item $f$ がシントミックであることと $f'$ がシントミックであることは同値である。
\label{more-morphisms-item-syntomic-fp-over-ft}
\item $f$ が滑らかであることと $f'$ が滑らかであることは同値である。
\label{more-morphisms-item-smooth-fp-over-ft}
\item $f$ が不分岐であることと $f'$ が不分岐であることは同値である。
\label{more-morphisms-item-unramified-fp-over-ft}
\item $f$ がエタールであることと $f'$ がエタールであることは同値である。
\label{more-morphisms-item-etale-fp-over-ft}
\item $f$ が固有であることと $f'$ が固有であることは同値である。
\label{more-morphisms-item-proper-fp-over-ft}
\item $f$ が有限であることと $f'$ が有限であることは同値である。
\label{more-morphisms-item-finite-fp-over-ft}
\item
\label{more-morphisms-item-finite-locally-free-fp-over-ft}
$f$ が（階数 $d$ の）有限局所自由であることと
$f'$ が（階数 $d$ の）有限局所自由であることは同値である。
% P05-EN-CH37-0045
\item ここにさらに追加する。
\end{enumerate}
\end{lemma}

\begin{proof}
仮定された $X$ と $Y$ の条件は、$f$ が $f'$ の $X \to X'$ による
基底変換であることを意味する。上の (1) -- (20) に挙げた性質
$\mathcal{P}$ はすべて基底変換で安定である。したがって $f'$ が性質
$\mathcal{P}$ をもてば $f$ もそれをもつ。次を参照せよ。
Schemes, Lemmas \ref{schemes-lemma-base-change-immersion},
\ref{schemes-lemma-quasi-compact-preserved-base-change},
\ref{schemes-lemma-separated-permanence}, および
\ref{schemes-lemma-base-change-monomorphism}
ならびに
Morphisms, Lemmas
\ref{morphisms-lemma-base-change-surjective},
\ref{morphisms-lemma-base-change-universally-injective},
\ref{morphisms-lemma-base-change-affine},
\ref{morphisms-lemma-base-change-quasi-finite},
\ref{morphisms-lemma-base-change-relative-dimension-d},
\ref{morphisms-lemma-base-change-syntomic},
\ref{morphisms-lemma-base-change-smooth},
\ref{morphisms-lemma-base-change-unramified},
\ref{morphisms-lemma-base-change-etale},
\ref{morphisms-lemma-base-change-proper},
\ref{morphisms-lemma-base-change-finite}, および
\ref{morphisms-lemma-base-change-finite-locally-free}.

\medskip\noindent
したがって各場合で興味深い向きは、$f$ がその性質をもつと仮定して
$f'$ もそれをもつことを導く向きである。以下の議論で断りなく用いる
一般的注意を二つ述べる。
(I) $W' \subset S'$ をアフィン開部分スキームとし、
$U' \subset X'$ と $V' \subset Y'$ を $W'$ 上にあり
$f'(U') \subset V'$ を満たすアフィン開部分スキームとする。
$W' = \Spec(R')$ とおき、
% P05-EN-CH37-0046
$I \subset R'$ を閉部分スキーム $W' \cap S$ を定義するイデアルと表す。
$U' = \Spec(B')$ および $V' = \Spec(A')$ とする。このとき可換図式
$$
\xymatrix{
0 \ar[r] &
IB' \ar[r] &
B' \ar[r] &
B \ar[r] & 0 \\
0 \ar[r] &
IA' \ar[r] \ar[u] &
A' \ar[r] \ar[u] &
A \ar[r] \ar[u] & 0
}
$$
を得、その各行は完全である。
(II) 射 $X \to X'$ と $Y \to Y'$ は普遍同相である。したがって写像
$f$ と $f'$ の位相は（任意の基底変換の後でも）同一である。
(III) $f$ が平坦ならば $f'$ は平坦であり、$Y' \to S'$ は $f'$ の像に属する
すべての点で平坦である。次を参照せよ。
% P05-EN-CH37-0047
Lemma \ref{more-morphisms-lemma-flatness-morphism-thickenings}.

\medskip\noindent
(\ref{more-morphisms-item-flat-fp-over-ft}) について。これは一般的注意 (III) である。

\medskip\noindent
(\ref{more-morphisms-item-isomorphism-fp-over-ft}) について。$f$ は同型であると仮定する。
アフィン開部分スキーム $V' \subset Y'$ を選び、
$U' = (f')^{-1}(V')$ とおく。このとき $V = Y \cap V'$ はアフィンであるから
$V \cong f^{-1}(V) = U = Y \times_{Y'} U'$ はアフィンである。Lemma
\ref{more-morphisms-lemma-thickening-affine-scheme} により $U'$ はアフィンである。
したがって一般的注意 (I) のような図式を得る。Algebra, Lemma
\ref{algebra-lemma-isomorphism-modulo-locally-nilpotent} により
$A' \to B'$ は同型、すなわち $U' \cong V'$ である。
よって $f'$ は同型である。

\medskip\noindent
(\ref{more-morphisms-item-open-immersion-fp-over-ft}) について。
$f$ は開埋め込みであると仮定する。このとき $f$ は $X$ と
開部分スキーム $V \subset Y$ との同型である。$V' \subset Y'$ を、
その台となる位相空間が $V$ である開部分スキームとする。
すると $f'$ は $X'$ から $V'$ への写像であり、
(\ref{more-morphisms-item-isomorphism-fp-over-ft}) により同型である。
したがって $f'$ は開埋め込みである。

\medskip\noindent
(\ref{more-morphisms-item-quasi-compact-fp-over-ft}) について。注意 (II) から直ちに従う。
より一般的な主張については Lemma
\ref{more-morphisms-lemma-thicken-property-morphisms} も参照せよ。

\medskip\noindent
(\ref{more-morphisms-item-universally-closed-fp-over-ft}) について。
注意 (II) から直ちに従う。より一般的な主張については Lemma
\ref{more-morphisms-lemma-thicken-property-morphisms} も参照せよ。

\medskip\noindent
(\ref{more-morphisms-item-separated-fp-over-ft}) について。
$X \times_Y X = Y \times_{Y'} (X' \times_{Y'} X')$ なので
$X' \times_{Y'} X'$ は $X \times_Y X$ の肥厚化であることに注意する。
したがって写像 $\Delta_{X/Y}$ と $\Delta_{X'/Y'}$ の位相は一致し、
結論を得る。より一般的な主張については Lemma
\ref{more-morphisms-lemma-thicken-property-morphisms} も参照せよ。

\medskip\noindent
(\ref{more-morphisms-item-monomorphism-fp-over-ft}) について。$f$ はモノ射であると仮定する。
対角射 $\Delta_{X'/Y'} : X' \to X' \times_{Y'} X'$ を考える。
$X' \times_{Y'} X' \to S'$ は局所有限型であることに注意する。
$\Delta_{X'/Y'}$ の $S \to S'$ による基底変換は $\Delta_{X/Y}$ であり、
仮定により同型である。(\ref{more-morphisms-item-isomorphism-fp-over-ft}) により
$\Delta_{X'/Y'}$ は同型であると結論する。

\medskip\noindent
(\ref{more-morphisms-item-surjective-fp-over-ft}) について。これは明らかである。
より一般的な主張については Lemma
\ref{more-morphisms-lemma-thicken-property-morphisms} も参照せよ。

\medskip\noindent
(\ref{more-morphisms-item-universally-injective-fp-over-ft}) について。
注意 (II) から直ちに従う。より一般的な主張については Lemma
\ref{more-morphisms-lemma-thicken-property-morphisms} も参照せよ。

\medskip\noindent
(\ref{more-morphisms-item-affine-fp-over-ft}) について。$f$ はアフィンであると仮定する。
アフィン開部分スキーム $V' \subset Y'$ を選び、
$U' = (f')^{-1}(V')$ とおく。このとき $V = Y \cap V'$ はアフィンなので
$U = Y \times_{Y'} U'$ もアフィンである。Lemma
\ref{more-morphisms-lemma-thickening-affine-scheme} により $U'$ はアフィンである。
したがって $f'$ はアフィンである。より一般的な主張については Lemma
\ref{more-morphisms-lemma-thicken-property-morphisms} も参照せよ。

\medskip\noindent
(\ref{more-morphisms-item-quasi-finite-fp-over-ft}) について。$f'$ が局所有限型であること
（Morphisms, Lemma \ref{morphisms-lemma-permanence-finite-type} による）と、
有限型射の準有限性はファイバー上で確認できることから従う。
Morphisms, Lemma \ref{morphisms-lemma-quasi-finite-at-point-characterize}
を参照せよ。

\medskip\noindent
(\ref{more-morphisms-item-relative-dimension-d-fp-over-ft}) について。
一般的注意 (II) と $f'$ が局所有限型であること
（Morphisms, Lemma \ref{morphisms-lemma-permanence-finite-type}）
から従う。

\medskip\noindent
(\ref{more-morphisms-item-universally-open-fp-over-ft}) について。
一般的注意 (II) から直ちに従う。より一般的な主張については Lemma
\ref{more-morphisms-lemma-thicken-property-morphisms} も参照せよ。

\medskip\noindent
(\ref{more-morphisms-item-syntomic-fp-over-ft}) について。$f$ はシントミックであると仮定する。
Morphisms, Lemma \ref{morphisms-lemma-finite-presentation-permanence}
により $f'$ は局所有限表示である。一般的注意 (III) により $f'$ は平坦である。
$f'$ のファイバーは $f$ のファイバーである。したがって Morphisms, Lemma
\ref{morphisms-lemma-syntomic-flat-fibres} により $f'$ はシントミックである。

\medskip\noindent
(\ref{more-morphisms-item-smooth-fp-over-ft}) について。$f$ は滑らかであると仮定する。
Morphisms, Lemma \ref{morphisms-lemma-finite-presentation-permanence}
により $f'$ は局所有限表示である。一般的注意 (III) により $f'$ は平坦である。
$f'$ のファイバーは $f$ のファイバーである。したがって Morphisms, Lemma
\ref{morphisms-lemma-smooth-flat-smooth-fibres} により $f'$ は滑らかである。

\medskip\noindent
(\ref{more-morphisms-item-unramified-fp-over-ft}) について。$f$ は不分岐であると仮定する。
Morphisms, Lemma \ref{morphisms-lemma-permanence-finite-type} により
$f'$ は局所有限型である。$f'$ のファイバーは $f$ のファイバーである。
したがって Morphisms, Lemma
\ref{morphisms-lemma-unramified-etale-fibres} により $f'$ は不分岐である。

\medskip\noindent
(\ref{more-morphisms-item-etale-fp-over-ft}) について。$f$ はエタールであると仮定する。
Morphisms, Lemma \ref{morphisms-lemma-finite-presentation-permanence}
により $f'$ は局所有限表示である。一般的注意 (III) により $f'$ は平坦である。
$f'$ のファイバーは $f$ のファイバーである。したがって Morphisms, Lemma
\ref{morphisms-lemma-etale-flat-etale-fibres} により $f'$ はエタールである。

\medskip\noindent
(\ref{more-morphisms-item-proper-fp-over-ft}) について。これは
(\ref{more-morphisms-item-separated-fp-over-ft})、$f$ が局所有限型であるという事実
（Morphisms, Lemma \ref{morphisms-lemma-permanence-finite-type}）、
(\ref{more-morphisms-item-quasi-compact-fp-over-ft})、および
(\ref{more-morphisms-item-universally-closed-fp-over-ft}) を合わせると従う。

\medskip\noindent
(\ref{more-morphisms-item-finite-fp-over-ft}) について。
(\ref{more-morphisms-item-universally-closed-fp-over-ft}),
(\ref{more-morphisms-item-affine-fp-over-ft}),
Morphisms, Lemma \ref{morphisms-lemma-integral-universally-closed},
$f$ が局所有限型であるという事実
（Morphisms, Lemma \ref{morphisms-lemma-permanence-finite-type}）、および
Morphisms, Lemma \ref{morphisms-lemma-finite-integral} を合わせる。

\medskip\noindent
(\ref{more-morphisms-item-finite-locally-free-fp-over-ft}) について。
$f$ は有限局所自由であると仮定する。
(\ref{more-morphisms-item-finite-fp-over-ft}) により $f'$ は有限である。
一般的注意 (III) により $f'$ は平坦である。Morphisms, Lemma
\ref{morphisms-lemma-finite-presentation-permanence} により
$f'$ は局所有限表示である。したがって Morphisms, Lemma
\ref{morphisms-lemma-finite-flat} により $f'$ は有限局所自由である。
\end{proof}

\begin{lemma}[射影スキームの変形]
\label{more-morphisms-lemma-deform-projective}
$f : X \to S$ を固有、平坦、かつ有限表示なスキームの射とする。
$\mathcal{L}$ を $f$-豊富とし、$S$ は準コンパクトであると仮定する。
ある $d_0 \geq 0$ が存在して、任意の Cartesian 図式
$$
\vcenter{
\xymatrix{
X \ar[r]_{i'} \ar[d]_f & X' \ar[d]^{f'} \\
S \ar[r]^i & S'
}
}
\quad\text{および}\quad
\begin{matrix}
\text{可逆 }\mathcal{O}_{X'}\text{-加群}\\
\mathcal{L}'\text{ で }\mathcal{L} \cong (i')^*\mathcal{L}'\text{ を満たすもの}
\end{matrix}
$$
で、$S \subset S'$ が肥厚化であり、$f'$ が固有、平坦、かつ有限表示であるものに
対して、次が成り立つ。
\begin{enumerate}
\item すべての $p > 0$ と $d \geq d_0$ に対して
$R^p(f')_*(\mathcal{L}')^{\otimes d} = 0$ である。
\item $\mathcal{A}'_d = (f')_*(\mathcal{L}')^{\otimes d}$
は $d \geq d_0$ に対して有限局所自由である。
\item $\mathcal{A}' =
\mathcal{O}_{S'} \oplus \bigoplus_{d \geq d_0} \mathcal{A}'_d$
は有限表示な準連接 $\mathcal{O}_{S'}$-代数である。
\item 標準同型
$r' : X' \to \underline{\text{Proj}}_{S'}(\mathcal{A}')$ が存在する。
\item 標準同型
$\theta' : (r')^*\mathcal{O}_{\underline{\text{Proj}}_{S'}(\mathcal{A}')}(1)
\to \mathcal{L}'$ が存在する。
\end{enumerate}
$\mathcal{A}'$, $r'$, $\theta'$ の構成は、データ
$(X', S', i, i', f', \mathcal{L}')$ に関して関手的である。
\end{lemma}

\begin{proof}
まず写像 $r'$ と $\theta'$ を記述する。
$\mathcal{L}'$ は $f'$-豊富であることに注意する。Lemma
\ref{more-morphisms-lemma-thicken-property-relatively-ample} を参照せよ。
準連接次数付き $\mathcal{O}_{S'}$-代数の標準写像
$\mathcal{A}' \to \bigoplus_{d \geq 0} (f')_*(\mathcal{L}')^{\otimes d}$
があり、これは次数 $\geq d_0$ で同型である。
したがって、これは構造層の Serre 捻りと整合する相対 Proj 上の同型を誘導する。
Constructions, Lemma
\ref{constructions-lemma-eventual-iso-graded-rings-map-relative-proj}
を参照せよ。ゆえに Morphisms, Lemma
\ref{morphisms-lemma-characterize-relatively-ample}
（これはさらに Constructions, Lemma
\ref{constructions-lemma-invertible-map-into-relative-proj}
で与えられる構成を用いる）により射 $r'$ を得る。そしてこれは
Morphisms, Lemma \ref{morphisms-lemma-proper-ample-is-proj}
により同型である。写像 $\theta'$ は Constructions, Lemma
\ref{constructions-lemma-invertible-map-into-relative-proj} から得られる。
Properties, Lemma \ref{properties-lemma-ample-gcd-is-one} により
$\theta'$ は同型であることが分かる（ここではさらに、Morphisms, Lemma
\ref{morphisms-lemma-proper-ample-is-proj} の証明で説明されているように、
$S'$ のアフィン開集合上の射 $r'$ が Properties, Lemma
\ref{properties-lemma-map-into-proj} の射と同じであることを用いる）。

\medskip\noindent
(1) と (2) で述べた消滅と局所自由性を仮定すれば、構成の関手性は
次のように分かる。$h : T \to S'$ をスキームの射とし、
$f_T : X'_T \to T$ を $f'$ の基底変換、
% P05-EN-CH37-0048
$\mathcal{L}_T$ を $X'_T$ への $\mathcal{L}$ の引き戻しと書く。
コホモロジーと基底変換（例えば Derived Categories of Schemes,
Lemma \ref{perfect-lemma-compare-base-change} の形）により、$T$ 上でも
対応する消滅が成り立ち、さらに
$h^*\mathcal{A}'_d = f_{T, *}\mathcal{L}_T^{\otimes d}$
である（したがって、押し出しの局所自由性と、対応する次数付き
$\mathcal{O}_T$-代数 $\mathcal{A}_T$ の有限生成性も成り立つ）。
ゆえに射
% P05-EN-CH37-0049
$r_T : X_T \to
\underline{\text{Proj}}_T(\bigoplus f_{T, *}\mathcal{L}_T^{\otimes d})$
は単に $r'$ の $T$ への基底変換であり、$\theta'$ の引き戻しは写像
$\theta_T$ である。

\medskip\noindent
以上から、(1)、(2)、(3) を証明すれば十分である。
すべての $d \geq d_0$ と $p > 0$ に対して
$R^pf_*\mathcal{L}^{\otimes d} = 0$ となるように $d_0$ を選ぶ。
Cohomology of Schemes, Lemma
\ref{coherent-lemma-coherent-proper-ample} を参照せよ。
この $d_0$ が条件を満たすと主張する。

\medskip\noindent
コホモロジーと基底変換（Derived Categories of Schemes,
Lemma \ref{perfect-lemma-flat-proper-perfect-direct-image-general}）により、
$E'_d = Rf'_*(\mathcal{L}')^{\otimes d}$ は
$D(\mathcal{O}_{S'})$ の完全対象であり、その形成は任意の基底変換と可換である。
特に $E_d = Li^*E'_d = Rf_*\mathcal{L}^{\otimes d}$ である。
Derived Categories of Schemes, Lemma
\ref{perfect-lemma-vanishing-implies-locally-free} により、$d \geq d_0$ に対して
複体 $E_d$ は、コホモロジー次数 $0$ に置かれた有限局所自由
$\mathcal{O}_S$-加群 $f_*\mathcal{L}^{\otimes d}$ と同型である。
すると Derived Categories of Schemes, Lemma
\ref{perfect-lemma-open-where-cohomology-in-degree-i-rank-r} により、
$E'_d$ はコホモロジー次数 $0$ に置かれた有限局所自由加群と同型である。
もちろん、これは $E'_d = \mathcal{A}'_d[0]$、
$p > 0$ に対して $R^pf'_*(\mathcal{L}')^{\otimes d} = 0$、かつ
$\mathcal{A}'_d$ が有限局所自由であることを意味する。
これで (1) と (2) が証明された。

\medskip\noindent
最後に示すべきことは、$\mathcal{A}'$ が $\mathcal{O}_{S'}$-代数の層として
有限表示であることである（この概念は Properties, Section
\ref{properties-section-extending-quasi-coherent-sheaves} で導入された）。
$U' = \Spec(R') \subset S'$ をアフィン開集合とする。
このとき $A' = \mathcal{A}'(U')$ は次数付き $R'$-代数であり、
各次数部分は有限射影 $R'$-加群である。
$A'$ が有限表示 $R'$-代数であることを示さなければならない。
これを Noetherian の場合への帰着によって証明する。実際、有限型
$\mathbf{Z}$-部分代数 $R'_0 \subset R'$ と、$R'_0$ 上の組
\footnote{$X' \to S'$ と $\mathcal{L}'$ がもつものと同じ性質、すなわち
$X'_0 \to \Spec(R'_0)$ は平坦かつ固有で、$\mathcal{L}'_0$ は豊富である。}
$(X'_0, \mathcal{L}'_0)$ で、その基底変換が
$(X'_{U'}, \mathcal{L}'|_{X'_{U'}})$ となるものを見つけられる。Limits, Lemmas
\ref{limits-lemma-descend-modules-finite-presentation},
\ref{limits-lemma-descend-invertible-modules},
\ref{limits-lemma-eventually-proper},
\ref{limits-lemma-descend-flat-finite-presentation}, および
\ref{limits-lemma-limit-ample} を参照せよ。
Cohomology of Schemes, Lemma \ref{coherent-lemma-coherent-proper-ample} により
$A'_0 = \bigoplus_{d \geq 0} H^0(X'_0, (\mathcal{L}'_0)^{\otimes d})$
は有限生成次数付き $R'_0$-代数であり、さらに、$d \geq d'_0$ に対する
$p > 0$ について
$H^p(X'_0, (\mathcal{L}'_0)^{\otimes d}) = 0$ となる $d'_0$ が存在する。
上の議論を $X'_0 \to \Spec(R'_0)$ と $\mathcal{L}'_0$ に適用すると、
$(A'_0)_d$ は有限射影 $R'_0$-加群であり、
$$
A'_d = \mathcal{A}'_d(U') =
H^0(X'_{U'}, (\mathcal{L}')^{\otimes d}|_{X'_{U'}}) =
H^0(X'_0, (\mathcal{L}'_0)^{\otimes d}) \otimes_{R'_0} R' =
(A'_0)_d \otimes_{R'_0} R'
$$
が $d \geq d'_0$ に対して成り立つことが分かる。ここで議論上の小さな注意は、
$d'_0$ を $d_0$ に等しく選べるかどうかが分からないことである
\footnote{実際には、さらに極限の議論を行えばこの場合に帰着できる。}。
これを回避するため、次の一連の議論で証明を終える。
\begin{enumerate}
\item[(a)] 代数
$B = R'_0 \oplus \bigoplus_{d \geq \max(d_0, d'_0)} (A'_0)_d$
は有限型 $R'_0$-代数である。$B \subset A'_0$ に Artin--Tate の補題を適用する。
Algebra, Lemma \ref{algebra-lemma-Artin-Tate} を参照せよ。
\item[(b)] $R'_0$ は Noetherian なので、$B$ は有限表示 $R'_0$-代数である。
\item[(c)] テンソル積の右完全性により、$B \otimes_{R'_0} R'$ は
有限表示 $R'$-代数である。
\item[(d)] 表示した等式により、これはちょうど
$C = R' \oplus \bigoplus_{d \geq \max(d_0, d'_0)} A'_d$
が有限表示 $R'$-代数であるということである。
% P05-EN-CH37-0050
\item[(e)] 商 $A'/C$ は有限射影 $R'$-加群 $A'_d$、
$d_0 \leq d \leq \max(d_0, d'_0)$ の直和であるから、
$R'$-加群として有限表示である。
\item[(f)] Algebra, Lemma
\ref{algebra-lemma-finitely-presented-over-subring} により、商 $A'/C$ は
$C$-加群として有限表示である。
\item[(g)] したがって Algebra, Lemma \ref{algebra-lemma-extension} により、
$A'$ は $C$-加群として有限表示である。
\item[(h)] Algebra, Lemma \ref{algebra-lemma-finite-finite-type} により、
これは $A'$ が有限表示 $C$-代数であることを意味する。
\item[(i)] 最後に、Algebra, Lemma
\ref{algebra-lemma-compose-finite-type} を $R' \to C \to A'$ に適用すると、
$A'$ は有限表示 $R'$-代数であると分かる。
\end{enumerate}
これで証明が完了する。
\end{proof}









\section{形式的滑らかな射}
\label{more-morphisms-section-formally-smooth}

\noindent
滑らかさの微分判定法（例えば Morphisms, Lemma
\ref{morphisms-lemma-smooth-at-point}）に対する Michael Artin の立場は、
それらは（実際上）基本的に役に立たない、というものである。
この節では、形式的滑らかな射 $X \to S$ の概念を導入する。
このような射は、$T$ がアフィンならば $X$ の $T$-値点が $T$ の
無限小肥厚化へ持ち上がるという性質によって特徴づけられる。
主結果は、形式的滑らかかつ局所有限表示な射が滑らかであるというものである。
Lemma \ref{more-morphisms-lemma-smooth-formally-smooth} を参照せよ。
この判定法は、上で述べた微分判定法よりもしばしば使いやすいことが分かる。

\medskip\noindent
環準同型 $R \to A$ が {\it 形式的滑らか} であるとは、任意の実線からなる可換図式
$$
\xymatrix{
A \ar[r] \ar@{-->}[rd] & B/I \\
R \ar[r] \ar[u] & B \ar[u]
}
$$
で $I \subset B$ が平方零イデアルであるものに対し、図式を可換にする点線矢印が
存在することをいう（Algebra, Definition
\ref{algebra-definition-formally-smooth} を参照）。これはスキームの射に対する
次の類似概念を動機づける。

\begin{definition}
\label{more-morphisms-definition-formally-smooth}
$f : X \to S$ をスキームの射とする。任意の実線からなる可換図式
$$
\xymatrix{
X \ar[d]_f & T \ar[d]^i \ar[l] \\
S & T' \ar[l] \ar@{-->}[lu]
}
$$
で $T \subset T'$ が $S$ 上のアフィンスキームの一次肥厚化であるものに対し、
図式を可換にする点線矢印が存在するとき、$f$ は {\it 形式的滑らか} であるという。
\end{definition}

\noindent
形式的不分岐射および形式的エタール射の場合には、$T'$ がアフィンであるという
条件を外すことができた。Lemmas
\ref{more-morphisms-lemma-formally-unramified-not-affine} および
\ref{more-morphisms-lemma-formally-etale-not-affine} を参照せよ。
形式的滑らかな射の場合、これはもはや正しくない。
実際、点線矢印を $T'$ 上 Zariski 局所的に補えるという条件の方が少し自然であろう。
Lemma \ref{more-morphisms-lemma-formally-smooth} の証明を解析すると、この条件は現在の定義と
同値であることが分かる。特に、形式的滑らかであることは始域上 Zariski 局所的である
（実際には始域上で滑らか局所的である。
% P05-EN-CH37-0051
ここに将来の参照を挿入する）。

\begin{lemma}
\label{more-morphisms-lemma-composition-formally-smooth}
形式的滑らかな射の合成は形式的滑らかである。
\end{lemma}

\begin{proof}
省略する。
\end{proof}

\begin{lemma}
\label{more-morphisms-lemma-base-change-formally-smooth}
形式的滑らかな射の基底変換は形式的滑らかである。
\end{lemma}

\begin{proof}
省略する。ただし代数版については Algebra, Lemma
\ref{algebra-lemma-base-change-fs} を参照せよ。
\end{proof}

\begin{lemma}
\label{more-morphisms-lemma-formally-etale-unramified-smooth}
$f : X \to S$ をスキームの射とする。このとき $f$ が形式的エタールであることと、
$f$ が形式的滑らかかつ形式的不分岐であることは同値である。
\end{lemma}

\begin{proof}
省略する。
\end{proof}

\begin{lemma}
\label{more-morphisms-lemma-formally-smooth-on-opens}
$f : X \to S$ をスキームの射とする。$U \subset X$ と $V \subset S$ を、
$f(U) \subset V$ を満たす開部分スキームとする。$f$ が形式的滑らかならば、
$f|_U : U \to V$ も形式的滑らかである。
\end{lemma}

\begin{proof}
Definition \ref{more-morphisms-definition-formally-smooth} にあるような実線図式
$$
\xymatrix{
U \ar[d]_{f|_U} & T \ar[d]^i \ar[l]^a \\
V & T' \ar[l] \ar@{-->}[lu]
}
$$
を考える。$f$ が形式的滑らかならば、$a'|_T = a$ を満たす $S$-射
$a' : T' \to X$ が存在する。$T$ と $T'$ の基礎集合は同じなので、
$a'$ は $U$ への射である（Schemes, Section
\ref{schemes-section-open-immersion} を参照）。また、これは明らかに
$V$-射でもある。したがって望ましい上の点線矢印を得る。
% P05-EN-CH37-0052
\end{proof}

\begin{lemma}
\label{more-morphisms-lemma-affine-formally-smooth}
$f : X \to S$ をスキームの射とし、$X$ と $S$ はアフィンであると仮定する。
このとき $f$ が形式的滑らかであることと、
$\mathcal{O}_S(S) \to \mathcal{O}_X(X)$ が形式的滑らかな環準同型であることは
同値である。
\end{lemma}

\begin{proof}
環とアフィンスキームの圏の同値により、これは定義（Definition
\ref{more-morphisms-definition-formally-smooth} および Algebra, Definition
\ref{algebra-definition-formally-smooth}）から直ちに従う。
Schemes, Lemma \ref{schemes-lemma-category-affine-schemes} を参照せよ。
\end{proof}

\noindent
次の補題はこの節の主結果である。これは、Limits, Proposition
\ref{limits-proposition-characterize-locally-finite-presentation} と合わせると、
射 $f : X \to S$ が滑らかであるかどうかを関手
$h_X : \Sch/S \to \textit{Sets}$ の「単純な」性質によって認識できることを
意味するという点で、関手的観点の勝利である。

\begin{lemma}[無限小持ち上げ判定法]
\label{more-morphisms-lemma-smooth-formally-smooth}
$f : X \to S$ をスキームの射とする。次は同値である。
\begin{enumerate}
\item 射 $f$ は滑らかである。
\item 射 $f$ は局所有限表示かつ形式的滑らかである。
\end{enumerate}
\end{lemma}

\begin{proof}
$f : X \to S$ は局所有限表示かつ形式的滑らかであると仮定する。
$f(U) \subset V$ を満たすアフィン開集合の組
$\Spec(A) = U \subset X$ と $\Spec(R) = V \subset S$ を考える。
Lemma \ref{more-morphisms-lemma-formally-smooth-on-opens} により $U \to V$ は形式的滑らかである。
Lemma \ref{more-morphisms-lemma-affine-formally-smooth} により $R \to A$ は形式的滑らかである。
Morphisms, Lemma
\ref{morphisms-lemma-locally-finite-presentation-characterize} により
$R \to A$ は有限表示である。Algebra, Proposition
\ref{algebra-proposition-smooth-formally-smooth} により $R \to A$ は滑らかである。
したがって滑らかな射の定義により $X \to S$ は滑らかである。

\medskip\noindent
逆に、$f : X \to S$ は滑らかであると仮定する。Definition
\ref{more-morphisms-definition-formally-smooth} にあるような実線からなる可換図式
$$
\xymatrix{
X \ar[d]_f & T \ar[d]^i \ar[l]^a \\
S & T' \ar[l] \ar@{-->}[lu]
}
$$
を考える。点線矢印が存在することを示し、それにより $f$ が形式的滑らかであることを
証明する。

\medskip\noindent
Remark \ref{more-morphisms-remark-special-case} で論じた特別な場合について、
$\mathcal{F}$ を Lemma \ref{more-morphisms-lemma-sheaf} の $T'$ 上の集合の層とする。
また
$$
\mathcal{H} =
\SheafHom_{\mathcal{O}_T}(a^*\Omega_{X/S}, \mathcal{C}_{T/T'})
$$
を、Lemma \ref{more-morphisms-lemma-action-sheaf} の作用
$\mathcal{H} \times \mathcal{F} \to \mathcal{F}$ をもつ
$\mathcal{O}_T$-加群の層とする。目標は単に
$\mathcal{F}(T) \not = \emptyset$ を示すことである。言い換えれば、
$\mathcal{F}$ が $T$ 上の自明な $\mathcal{H}$-捩子であることを示そうとしている
（Cohomology, Section \ref{cohomology-section-h1-torsors} を参照）。
二つの段階がある。(I) $\mathcal{F}$ が捩子であることを示すには、すべての
$t \in T$ に対して $\mathcal{F}_t \not = \emptyset$ を示さなければならない
（Cohomology, Definition \ref{cohomology-definition-torsor} を参照）。
(II) $\mathcal{F}$ が自明な捩子であることを示すには、
$H^1(T, \mathcal{H}) = 0$ を示せば十分である（Cohomology, Lemma
\ref{cohomology-lemma-torsors-h1} を参照。ここでは $\mathcal{H}$ を Abel 層として
みたコホモロジーと $\mathcal{O}_T$-加群としてみたコホモロジーのどちらを用いてもよい。
Cohomology, Lemma \ref{cohomology-lemma-modules-abelian} を参照せよ）。

\medskip\noindent
まず (I) を証明する。各 $t \in T$ に対し、$a(U)$ が、あるアフィン開集合
$\Spec(R) = V \subset S$ へ写るアフィン開集合
$\Spec(A) = W \subset X$ に含まれるような、$t$ のアフィン開近傍
$U \subset T$ を選べる。Morphisms, Lemma
\ref{morphisms-lemma-smooth-characterize} により環準同型 $R \to A$ は滑らかである。
したがって Algebra, Proposition
\ref{algebra-proposition-smooth-formally-smooth} により環準同型 $R \to A$ は
形式的滑らかである。さらに Lemma \ref{more-morphisms-lemma-affine-formally-smooth} により
$W \to V$ は形式的滑らかである。ゆえに $a|_U : U \to W$ を $V$-射
$a' : U' \to W \subset X$ に持ち上げることができ、
$\mathcal{F}(U) \not = \emptyset$ が示される。

\medskip\noindent
最後に (II) を証明する。Morphisms, Lemma
\ref{morphisms-lemma-finite-presentation-differentials} により
$\Omega_{X/S}$ は有限表示である（実際、Morphisms, Lemma
\ref{morphisms-lemma-smooth-omega-finite-locally-free} により有限局所自由でさえある）。
したがって $a^*\Omega_{X/S}$ は有限表示である（Modules, Lemma
\ref{modules-lemma-pullback-finite-presentation} を参照）。ゆえに層
$\mathcal{H} =
\SheafHom_{\mathcal{O}_T}(a^*\Omega_{X/S}, \mathcal{C}_{T/T'})$
は、Schemes, Section \ref{schemes-section-quasi-coherent} の議論により準連接である。
したがって Cohomology of Schemes, Lemma
\ref{coherent-lemma-quasi-coherent-affine-cohomology-zero} により、望み通り
$H^1(T, \mathcal{H}) = 0$ である。
\end{proof}

\noindent
局所射影準連接加群は Properties, Section
\ref{properties-section-locally-projective} で定義される。

\begin{lemma}
\label{more-morphisms-lemma-formally-smooth-sheaf-differentials}
$f : X \to Y$ を形式的滑らかなスキームの射とする。
このとき $\Omega_{X/Y}$ は $X$ 上局所射影である。
\end{lemma}

\begin{proof}
アフィン開集合 $U \subset X$ と $V \subset Y$ を、
$f(U) \subset V$ となるように選ぶ。Lemma
\ref{more-morphisms-lemma-formally-smooth-on-opens} により
$f|_U : U \to V$ は形式的滑らかである。したがって Lemma
\ref{more-morphisms-lemma-affine-formally-smooth} により
$\Gamma(V, \mathcal{O}_V) \to \Gamma(U, \mathcal{O}_U)$ は形式的滑らかな
環準同型である。ゆえに Algebra, Lemma
\ref{algebra-lemma-characterize-formally-smooth-again} により、
$\Gamma(U, \mathcal{O}_U)$-加群
$\Omega_{\Gamma(U, \mathcal{O}_U)/\Gamma(V, \mathcal{O}_V)}$
は射影的である。したがって $\Omega_{U/V}$ は局所射影である。
Properties, Section \ref{properties-section-locally-projective} を参照せよ。
\end{proof}

\begin{lemma}
\label{more-morphisms-lemma-h1-is-zero}
$T$ をアフィンスキームとする。$\mathcal{F}$ と $\mathcal{G}$ を準連接
$\mathcal{O}_T$-加群とする。
$\mathcal{H} = \SheafHom_{\mathcal{O}_T}(\mathcal{F}, \mathcal{G})$.
を考える。$\mathcal{F}$ が局所射影ならば $H^1(T, \mathcal{H}) = 0$ である。
\end{lemma}

\begin{proof}
スキーム上の局所射影層の定義（Properties, Definition
\ref{properties-definition-locally-projective} を参照）により、
$\mathcal{F}$ は自由 $\mathcal{O}_T$-加群の直和因子である。
したがって $\mathcal{F} = \bigoplus_{i \in I} \mathcal{O}_T$ は自由加群であると
仮定してよい。この場合 $\mathcal{H} = \prod_{i \in I} \mathcal{G}$ は
準連接加群の積である。Cohomology, Lemma
\ref{cohomology-lemma-cohomology-products} により $H^1 = 0$ と結論する。
実際、アフィンスキーム上の準連接層のコホモロジーは零である。
Cohomology of Schemes, Lemma
\ref{coherent-lemma-quasi-coherent-affine-cohomology-zero} を参照せよ。
\end{proof}

\begin{lemma}
\label{more-morphisms-lemma-formally-smooth}
$f : X \to Y$ をスキームの射とする。次は同値である。
\begin{enumerate}
\item $f$ は形式的滑らかである。
\item すべての $x \in X$ に対し、$f(U) \subset V$ かつ
$f|_U : U \to V$ が形式的滑らかとなるような開集合
$x \in U \subset X$ と $f(x) \in V \subset Y$ が存在する。
\item $f(U) \subset V$ を満たす任意のアフィン開集合の組
$U \subset X$ と $V \subset Y$ に対し、環準同型
$\mathcal{O}_Y(V) \to \mathcal{O}_X(U)$ は形式的滑らかである。
\item アフィン開被覆 $Y = \bigcup V_j$ と、各 $j$ に対するアフィン開被覆
$f^{-1}(V_j) = \bigcup U_{ji}$ が存在し、すべての $j$ と $i$ に対して
% P05-EN-CH37-0053
$\mathcal{O}_Y(V) \to \mathcal{O}_X(U)$ が形式的滑らかな環準同型である。
\end{enumerate}
\end{lemma}

\begin{proof}
含意 (1) $\Rightarrow$ (2)、(1) $\Rightarrow$ (3)、および
(2) $\Rightarrow$ (4) は Lemma
% P05-EN-CH37-0054
\ref{more-morphisms-lemma-formally-smooth-on-opens} から従う。
含意 (3) $\Rightarrow$ (4) は直ちに従う。

\medskip\noindent
(4) を仮定する。$f$ が形式的滑らかであることの証明は Lemma
\ref{more-morphisms-lemma-smooth-formally-smooth} の証明の後半と同じである。
Definition \ref{more-morphisms-definition-formally-smooth} にあるような実線からなる可換図式
$$
\xymatrix{
X \ar[d]_f & T \ar[d]^i \ar[l]^a \\
Y & T' \ar[l] \ar@{-->}[lu]
}
$$
を考える。点線矢印が存在することを示し、それにより $f$ が形式的滑らかであることを
証明する。Remark \ref{more-morphisms-remark-special-case} で論じた特別な場合について、
$\mathcal{F}$ を Lemma \ref{more-morphisms-lemma-sheaf} の $T'$ 上の集合の層とする。また
$$
\mathcal{H} =
\SheafHom_{\mathcal{O}_T}(a^*\Omega_{X/Y}, \mathcal{C}_{T/T'})
$$
を、Lemma \ref{more-morphisms-lemma-action-sheaf} の作用
$\mathcal{H} \times \mathcal{F} \to \mathcal{F}$ をもつ $T$ 上の
$\mathcal{O}_T$-加群の層とする。作用
$\mathcal{H} \times \mathcal{F} \to \mathcal{F}$ によって
$\mathcal{F}$ は擬 $\mathcal{H}$-捩子となる。Cohomology, Definition
\ref{cohomology-definition-torsor} を参照せよ。
目標は $\mathcal{F}$ が自明な $\mathcal{H}$-捩子であることを示すことである。
二つの段階がある。(I) $\mathcal{F}$ が捩子であることを示すには、
$\mathcal{F}$ が局所的に切断をもつことを示さなければならない。
(II) $\mathcal{F}$ が自明な捩子であることを示すには、
$H^1(T, \mathcal{H}) = 0$ を示せば十分である。Cohomology, Lemma
\ref{cohomology-lemma-torsors-h1} を参照せよ。

\medskip\noindent
まず (I) を証明する。各 $t \in T$ に対し、ある $i,j$ について
$a(W)$ が $U_{ji}$ に含まれるような $t$ のアフィン開近傍
$W \subset T$ を選べる。$W' \subset T'$ を対応する開部分スキームとする。
仮定 (4) により $a|_W : W \to U_{ji}$ を $V_j$-射
$a' : W' \to U_{ji}$ に持ち上げることができ、
$\mathcal{F}(W')$ が空でないことが示される。

\medskip\noindent
最後に (II) を証明する。Lemma
\ref{more-morphisms-lemma-formally-smooth-sheaf-differentials} により、
$\Omega_{U_{ji}/V_j}$ は局所射影である。
% P05-EN-CH37-0055
したがって Properties, Lemma
\ref{properties-lemma-locally-projective} により
$\Omega_{X/Y}$ は局所射影である。さらに Properties, Lemma
\ref{properties-lemma-locally-projective-pullback} により
$a^*\Omega_{X/Y}$ は局所射影である。ゆえに望み通り、Lemma
\ref{more-morphisms-lemma-h1-is-zero} により
$$
H^1(T, \mathcal{H}) =
H^1(T, \SheafHom_{\mathcal{O}_T}(a^*\Omega_{X/Y}, \mathcal{C}_{T/T'})) = 0
$$
である。
\end{proof}

\begin{lemma}
\label{more-morphisms-lemma-triangle-differentials-formally-smooth}
$f : X \to Y$ と $g : Y \to S$ をスキームの射とする。
$f$ は形式的滑らかであると仮定する。このとき
$$
0 \to f^*\Omega_{Y/S} \to \Omega_{X/S} \to \Omega_{X/Y} \to 0
$$
は短完全列である（Morphisms, Lemma
\ref{morphisms-lemma-triangle-differentials} を参照）。
\end{lemma}

\begin{proof}
この補題の代数版は次の通りである。環準同型 $A \to B \to C$ で
$B \to C$ が形式的滑らかなものが与えられたとき、Algebra, Lemma
\ref{algebra-lemma-exact-sequence-differentials} の列
$$
0 \to C \otimes_B \Omega_{B/A} \to \Omega_{C/A} \to \Omega_{C/B} \to 0
$$
は完全である。これは Algebra, Lemma
\ref{algebra-lemma-ses-formally-smooth} である。
\end{proof}

\begin{lemma}
\label{more-morphisms-lemma-differentials-formally-unramified-formally-smooth}
$h : Z \to X$ を $S$ 上の形式的不分岐なスキームの射とする。
$Z$ は $S$ 上形式的滑らかであると仮定する。このとき Lemma
\ref{more-morphisms-lemma-universally-unramified-differentials-sequence} の標準完全列
$$
0 \to \mathcal{C}_{Z/X} \to h^*\Omega_{X/S} \to \Omega_{Z/S} \to 0
$$
は短完全列である。
\end{lemma}

\begin{proof}
$Z \to Z'$ を $X$ 上の $Z$ の普遍一次肥厚化とする。
Lemma \ref{more-morphisms-lemma-universally-unramified-differentials-sequence} の証明から、
上の列は列
$$
\mathcal{C}_{Z/Z'} \to \Omega_{Z'/S} \otimes \mathcal{O}_Z \to
\Omega_{Z/S} \to 0.
$$
と同一視される。$Z \to S$ は形式的滑らかなので、$Z'$ 上局所的に、
包含写像 $Z \to Z'$ の $S$ 上の左逆 $Z' \to Z$ を見つけられる。
したがってこの列は局所分裂する。Morphisms, Lemma
\ref{morphisms-lemma-differentials-relative-immersion-section} を参照せよ。
\end{proof}

\begin{lemma}
\label{more-morphisms-lemma-two-unramified-morphisms-formally-smooth}
可換図式
$$
\xymatrix{
Z \ar[r]_i \ar[rd]_j & X \ar[d]^f \\
& Y
}
$$
で $i$ と $j$ が形式的不分岐、$f$ が形式的滑らかであるものを考える。
このとき Lemma \ref{more-morphisms-lemma-two-unramified-morphisms} の標準完全列
$$
0 \to
\mathcal{C}_{Z/Y} \to
\mathcal{C}_{Z/X} \to
i^*\Omega_{X/Y} \to 0
$$
は完全かつ局所分裂する。
\end{lemma}

\begin{proof}
$Z \to Z'$ を $X$ 上の $Z$ の普遍一次肥厚化、
$Z \to Z''$ を $Y$ 上の $Z$ の普遍一次肥厚化と書く。
% P05-EN-CH37-0056
Lemma \ref{more-morphisms-lemma-universally-unramified-differentials-sequence} により、
標準射 $Z' \to Z''$ があり、可換図式
% P05-EN-CH37-0057
$$
\xymatrix{
Z \ar[r]_{i'} \ar[rd]_{j'} & Z' \ar[r]_a \ar[d]^k & X \ar[d]^f \\
& Z'' \ar[r]^b & Y
}
$$
を得る。Lemma \ref{more-morphisms-lemma-two-unramified-morphisms} の証明において、
上の列を列
$$
\mathcal{C}_{Z/Z''} \to
\mathcal{C}_{Z/Z'} \to
(i')^*\Omega_{Z'/Z''} \to 0
$$
と同一視した。$U'' \subset Z''$ をアフィン開集合とする。
$U \subset Z$ と $U' \subset Z'$ を対応するアフィン開部分スキームと書く。
$f$ は形式的滑らかなので、$U$ 上で $i$ と一致し、
$f \circ h$ が $b|_{U''}$ に等しい射 $h : U'' \to X$ が存在する。
$Z'$ は普遍一次肥厚化なので、
% P05-EN-CH37-0058
$g = a \circ h$ を満たす一意な射 $g : U'' \to Z'$ を得る。
$Z''$ の普遍性により $k \circ g$ は包含写像 $U'' \to Z''$ である。
したがって $g$ は $k$ の左逆である。図式で表すと
$$
\xymatrix{
U \ar[d] \ar[r] & Z' \ar[d]^k \\
U'' \ar[r] \ar[ru]^g & Z''
}
$$
である。したがって $g$ は写像
$\mathcal{C}_{Z/Z'}|_U \to \mathcal{C}_{Z/Z''}|_U$ を誘導し、これは $U$ 上の写像
$\mathcal{C}_{Z/Z''} \to \mathcal{C}_{Z/Z'}$ の左逆である。
\end{proof}



































\section{Noetherian 基底上の滑らかさ}
\label{more-morphisms-section-smooth-Noetherian}

\noindent
基底が Noetherian ならば、形式的滑らかさの定式化における条件を弱められることが
分かる。ある意味で、次の補題群は変形理論の始まりである。

\begin{lemma}
\label{more-morphisms-lemma-lifting-along-artinian-at-point}
$f : X \to S$ をスキームの射とし、$x \in X$ とする。
$S$ は局所 Noetherian であり、$f$ は局所有限型であると仮定する。
% P05-EN-CH37-0059
次は同値である。
\begin{enumerate}
\item $f$ は $x$ において滑らかである。
\item 任意の実線からなる可換図式
$$
\xymatrix{
X \ar[d]_f & \Spec(B) \ar[d]^i \ar[l]^-\alpha \\
S & \Spec(B') \ar[l]_-{\beta} \ar@{-->}[lu]
}
$$
で、$B' \to B$ が局所環の全射、$\Ker(B' \to B)$ が平方零、かつ
$\alpha$ が $\Spec(B)$ の閉点を $x$ に写すものに対し、図式を可換にする
点線矢印が存在する。
\item (2) と同じだが、$B' \to B$ は小拡大全体を動く
（Algebra, Definition \ref{algebra-definition-small-extension} を参照）。
\item (2) と同じだが、$B' \to B$ は、$\mathfrak m \subset B$ を極大イデアルとして
$\alpha$ が同型 $\kappa(x) \to \kappa(\mathfrak m)$ を誘導するような小拡大全体を
動く。
\end{enumerate}
\end{lemma}

\begin{proof}
$f(x)$ のアフィン近傍 $V \subset S$ と、$f(U) \subset V$ を満たす
$x$ のアフィン近傍 $U \subset X$ を選ぶ。(2) にある任意の「試験」図式について、
射 $\alpha$ は $\Spec(B)$ を $U$ に写し、射 $\beta$ は $\Spec(B')$ を $V$ に写す
（Schemes, Section \ref{schemes-section-points} を参照）。
したがって、この補題はアフィン間の射 $f|_U : U \to V$ に帰着する。
（実際、$V$ は Noetherian で $f|_U$ は有限型である。Properties, Lemma
\ref{properties-lemma-locally-Noetherian} および Morphisms, Lemma
\ref{morphisms-lemma-locally-finite-type-characterize} を参照せよ。）
このアフィンの場合、補題は Algebra, Lemma
\ref{algebra-lemma-smooth-test-artinian} と同一である。
\end{proof}

\noindent
持ち上げ判定法を、有限型点を「中心とする」小拡大についてだけ確認すればよいと
知ることが有用な場合がある（Morphisms, Section
\ref{morphisms-section-points-finite-type} を参照）。

\begin{lemma}
\label{more-morphisms-lemma-lifting-along-artinian}
$f : X \to S$ をスキームの射とする。$S$ は局所 Noetherian であり、
$f$ は局所有限型であると仮定する。次は同値である。
\begin{enumerate}
\item $f$ は滑らかである。
\item 任意の実線からなる可換図式
$$
\xymatrix{
X \ar[d]_f & \Spec(B) \ar[d]^i \ar[l]^-\alpha \\
S & \Spec(B') \ar[l]_-{\beta} \ar@{-->}[lu]
}
$$
で、$B' \to B$ が Artin 局所環の小拡大、かつ $\beta$ が有限型（！）であるものに
対し、図式を可換にする点線矢印が存在する。
% P05-EN-CH37-0060
\end{enumerate}
\end{lemma}

\begin{proof}
$f$ が滑らかならば、無限小持ち上げ判定法（Lemma
\ref{more-morphisms-lemma-smooth-formally-smooth}）により $f$ は形式的滑らかであり、(2) が成り立つ。

\medskip\noindent
(2) を仮定する。$f$ が滑らかでない点 $x \in X$ の集合は $X$ の閉部分集合
$T$ をなす。Morphisms, Section
\ref{morphisms-section-points-finite-type} の議論により、
$T \not = \emptyset$ ならば射
$$
\Spec(\kappa(x)) \to X \to S
$$
が有限型となる点 $x \in T \subset X$ が存在する（すなわち、$X$ のあるアフィン開集合で
閉じている $T$ の任意の点 $x$ を選ぶ）。Morphisms, Lemma
\ref{morphisms-lemma-artinian-finite-type} により、剰余体が $\kappa(x)$ である
任意の Artin 局所環 $B'$ に対し、任意の射
$\beta : \Spec(B') \to S$ は有限型である。したがって Lemma
\ref{more-morphisms-lemma-lifting-along-artinian-at-point} の (4) で用いるすべての図式は、
現在の補題の (2) にある図式に対応することが分かる。
ゆえに $X \to S$ は $x$ において滑らかであり、矛盾する。
% P05-EN-CH37-0061
\end{proof}

\noindent
次は有用な応用である。

\begin{lemma}
\label{more-morphisms-lemma-check-smoothness-on-infinitesimal-nbhds}
$f : X \to S$ を局所 Noetherian スキームの有限型射とする。
$Z \subset S$ を閉部分スキームとし、その $n$ 次無限小近傍を
$Z_n \subset S$ とする。$X_n = Z_n \times_S X$ とおく。
\begin{enumerate}
\item すべての $n$ に対して $X_n \to Z_n$ が滑らかならば、$f$ は
$f^{-1}(Z)$ のすべての点で滑らかである。
\item すべての $n$ に対して $X_n \to Z_n$ がエタールならば、$f$ は
$f^{-1}(Z)$ のすべての点でエタールである。
\end{enumerate}
\end{lemma}

\begin{proof}
すべての $n$ に対して $X_n \to Z_n$ は滑らかであると仮定する。
$x \in X$ を $Z$ のある点の上にある点とする。Lemma
\ref{more-morphisms-lemma-lifting-along-artinian-at-point} の (3) にある小拡大
$B' \to B$ と射 $\alpha$, $\beta$ が与えられたとする。
$B'$ の極大イデアルは冪零である（$B'$ は Artin 環なので）。したがって、適当な $n$ に対し
射 $\beta$ は $Z_n$ を経由し、$\alpha$ は $X_n$ を経由する。
そこで $X_n \to Z_n$ の持ち上げ性を使えば、図式内の望ましい点線矢印を得る。
これで (1) が証明された。(2) は (1) と、射がエタールであることと相対次元 $0$ で
滑らかであることが同値であるという事実から従う。
\end{proof}

\begin{lemma}
\label{more-morphisms-lemma-check-flatness-on-infinitesimal-nbhds}
$f : X \to S$ を局所 Noetherian スキームの射とする。
$Z \subset S$ を閉部分スキームとし、その $n$ 次無限小近傍を
$Z_n \subset S$ とする。$X_n = Z_n \times_S X$ とおく。
すべての $n$ に対して $X_n \to Z_n$ が平坦ならば、$f$ は
$f^{-1}(Z)$ のすべての点で平坦である。
\end{lemma}

\begin{proof}
これは Algebra, Lemma \ref{algebra-lemma-flat-module-powers} を
スキームの言葉に翻訳したものである。
\end{proof}








\section{素朴な余接複体}
\label{more-morphisms-section-netherlander}

\noindent
この節は Modules, Section \ref{modules-section-netherlander} の続きであり、
その節はさらに Algebra, Section \ref{algebra-section-netherlander} の議論を
継続するものである。

\begin{definition}
\label{more-morphisms-definition-netherlander}
$f : X \to Y$ をスキームの射とする。{\it $f$ の素朴な余接複体}とは、
Modules, Definition
\ref{modules-definition-cotangent-complex-morphism-ringed-topoi} で定義された
複体をいう。記法は $\NL_f$ または $\NL_{X/Y}$ とする。
\end{definition}

\begin{lemma}
\label{more-morphisms-lemma-NL-affine}
$f : X \to Y$ をスキームの射とする。
$\Spec(A) = U \subset X$ と
% P05-EN-CH37-0062
$\Spec(R) = V \subset S$ を $f(U) \subset V$ を満たすアフィン開集合とする。
複体の標準写像
$$
\widetilde{\NL_{A/R}} \longrightarrow \NL_{X/Y}|_U
$$
があり、これは $D(\mathcal{O}_U)$ で同型である。
\end{lemma}

\begin{proof}
Modules, Section \ref{modules-section-netherlander} における
$\NL_{X/Y}$ の構成から、$A = \mathcal{O}_X(U)$-加群の複体の標準写像
$\NL_{\mathcal{O}_X(U)/f^{-1}\mathcal{O}_Y(U)} \to \NL_{X/Y}(U)$
があり、さらなる制限と整合することが分かる。標準写像
$R \to f^{-1}\mathcal{O}_Y(U)$ を用いて、$A$-加群の複体の標準写像
$\NL_{A/R} \to \NL_{\mathcal{O}_X(U)/f^{-1}\mathcal{O}_Y(U)}$
を得る。関手 $\widetilde{\ }$ の普遍性（Schemes, Lemma
\ref{schemes-lemma-compare-constructions} を参照）を用いると、補題の主張にある
写像を得る。この写像がコホモロジー層上で同型であることは、茎上で同型を誘導することを
確認すれば判定できる。これは Algebra, Lemmas
\ref{algebra-lemma-NL-localize-bottom} および
\ref{algebra-lemma-localize-NL}、ならびに Modules, Lemma
\ref{modules-lemma-stalk-NL} から従う（また、点
$\mathfrak p \in \Spec(A)$ における $\mathcal{O}_X$ と
$f^{-1}\mathcal{O}_Y$ の茎はそれぞれ $A_\mathfrak p$ と
$R_\mathfrak q$ である。ここで $\mathfrak q = R \cap \mathfrak p$ であり、
用いる参照は Schemes, Lemma \ref{schemes-lemma-spec-sheaves} および
Sheaves, Lemma \ref{sheaves-lemma-stalk-pullback} である）。
\end{proof}

\begin{lemma}
\label{more-morphisms-lemma-netherlander-quasi-coherent}
$f : X \to Y$ をスキームの射とする。複体 $\NL_{X/Y}$ のコホモロジー層は
準連接であり、次数 $-1$, $0$ 以外では零であり、次数 $0$ では
$\Omega_{X/Y}$ に等しい。
\end{lemma}

\begin{proof}
Modules, Section \ref{modules-section-netherlander} における素朴な余接複体の
構成により、$\NL_{X/Y}$ は次数 $-1$, $0$ にある複体であり、その次数 $0$ の
コホモロジーは $\Omega_{X/Y}$ である。微分層は準連接である
（Morphisms, Lemma \ref{morphisms-lemma-differentials-diagonal} による）。
証明を終えるには $H^{-1}(\NL_{X/Y})$ が準連接であることを示せば十分である。
これは Lemma \ref{more-morphisms-lemma-NL-affine} を用いてアフィン上で確認すれば従う。
\end{proof}

\begin{lemma}
\label{more-morphisms-lemma-netherlander-fp}
$f : X \to Y$ をスキームの射とする。$f$ が局所有限表示ならば、
$\NL_{X/Y}$ は $X$ 上局所的に準連接 $\mathcal{O}_X$-加群の複体
$$
\ldots \to 0 \to \mathcal{F}^{-1} \to \mathcal{F}^0 \to 0 \to \ldots
$$
で、$\mathcal{F}^0$ が有限表示、$\mathcal{F}^{-1}$ が有限型であるものと
擬同型である。
\end{lemma}

\begin{proof}
Lemma \ref{more-morphisms-lemma-NL-affine} により、$R \to A$ が有限表示な環準同型ならば
$\NL_{A/R}$ がこの形をもつことを示せば十分である。
$I$ が有限生成となるように $A = R[x_1, \ldots, x_n]/I$ と書く。
このとき $I/I^2$ は有限 $A$-加群であり、Algebra, Section
\ref{algebra-section-netherlander}、特に Algebra, Lemma
\ref{algebra-lemma-NL-homotopy} により $\NL_{A/R}$ は
$$
\ldots \to 0 \to I/I^2 \to
\bigoplus\nolimits_{i = 1, \ldots, n} A\text{d}x_i \to 0 \to \ldots
$$
と擬同型である。
\end{proof}

\begin{lemma}
\label{more-morphisms-lemma-NL-formally-smooth}
$f : X \to Y$ をスキームの射とする。次は同値である。
\begin{enumerate}
\item $f$ は形式的滑らかである。
\item $H^{-1}(\NL_{X/Y}) = 0$ であり、
$H^0(\NL_{X/Y}) = \Omega_{X/Y}$ は局所射影である。
\end{enumerate}
\end{lemma}

\begin{proof}
これは Algebra, Proposition
\ref{algebra-proposition-characterize-formally-smooth} および
Lemma \ref{more-morphisms-lemma-formally-smooth} から従う。
\end{proof}

\begin{lemma}
\label{more-morphisms-lemma-NL-formally-etale}
$f : X \to Y$ をスキームの射とする。次は同値である。
\begin{enumerate}
\item $f$ は形式的エタールである。
\item $H^{-1}(\NL_{X/Y}) = H^0(\NL_{X/Y}) = 0$.
\end{enumerate}
\end{lemma}

\begin{proof}
形式的エタール射は形式的滑らかなので、Lemma
\ref{more-morphisms-lemma-NL-formally-smooth} により
$H^{-1}(\NL_{X/Y}) = 0$ である。一方、Lemma
\ref{more-morphisms-lemma-characterize-formally-etale} により $\Omega_{X/Y} = 0$ である。
逆に (2) が成り立つならば、$f$ は Lemma
\ref{more-morphisms-lemma-NL-formally-smooth} により形式的滑らかであり、Lemma
\ref{more-morphisms-lemma-formally-unramified-differentials} により形式的不分岐である。
したがって Lemmas
% P05-EN-CH37-0063
\ref{more-morphisms-lemma-formally-etale-unramified-smooth} により形式的エタールである。
\end{proof}

\begin{lemma}
\label{more-morphisms-lemma-NL-smooth}
$f : X \to Y$ をスキームの射とする。次は同値である。
\begin{enumerate}
\item $f$ は滑らかである。
\item $f$ は局所有限表示であり、$H^{-1}(\NL_{X/Y}) = 0$ であり、
$H^0(\NL_{X/Y}) = \Omega_{X/Y}$ は有限局所自由である。
\end{enumerate}
\end{lemma}

\begin{proof}
これは滑らかな環準同型の定義（Algebra, Definition
\ref{algebra-definition-smooth}）、Lemma \ref{more-morphisms-lemma-NL-affine}、および
滑らかなスキームの射の定義（Morphisms, Definition
\ref{morphisms-definition-smooth}）から従う。また、環上の加群について
有限局所自由と有限射影が同じであることも用いる（Algebra, Lemma
\ref{algebra-lemma-finite-projective}）。
\end{proof}

\begin{lemma}
\label{more-morphisms-lemma-NL-etale}
$f : X \to Y$ をスキームの射とする。次は同値である。
\begin{enumerate}
\item $f$ はエタールである。
\item $f$ は局所有限表示であり、
$H^{-1}(\NL_{X/Y}) = H^0(\NL_{X/Y}) = 0$.
\end{enumerate}
\end{lemma}

\begin{proof}
これはエタール環準同型の定義（Algebra, Definition
\ref{algebra-definition-etale}）、Lemma \ref{more-morphisms-lemma-NL-affine}、および
エタールなスキームの射の定義（Morphisms, Definition
\ref{morphisms-definition-etale}）から従う。
\end{proof}


\begin{lemma}
\label{more-morphisms-lemma-NL-immersion}
$i : Z \to X$ をスキームの埋め込みとする。このとき $\NL_{Z/X}$ は
$D(\mathcal{O}_Z)$ において $\mathcal{C}_{Z/X}[1]$ と同型である。
ここで $\mathcal{C}_{Z/X}$ は $X$ における $Z$ の余法層である。
\end{lemma}

\begin{proof}
これは Algebra, Lemma \ref{algebra-lemma-NL-surjection}、
Morphisms, Lemma \ref{morphisms-lemma-affine-conormal}、および
Lemma \ref{more-morphisms-lemma-NL-affine} から従う。
\end{proof}

\begin{lemma}
\label{more-morphisms-lemma-exact-sequence-NL}
$f : X \to Y$ と $g : Y \to Z$ をスキームの射とする。
コホモロジー層の標準六項完全列
$$
H^{-1}(f^*\NL_{Y/Z}) \to
H^{-1}(\NL_{X/Z}) \to
H^{-1}(\NL_{X/Y}) \to
f^*\Omega_{Y/Z} \to \Omega_{X/Z} \to \Omega_{X/Y} \to 0
$$
が存在する。
\end{lemma}

\begin{proof}
Modules, Lemma \ref{modules-lemma-exact-sequence-NL-ringed-topoi} の特別な場合である。
\end{proof}

\begin{lemma}
\label{more-morphisms-lemma-get-triangle-NL}
$f : X \to Y$ と $Y \to Z$ をスキームの射とする。
$X \to Y$ は完全交叉射であると仮定する。このとき $D(\mathcal{O}_X)$ に標準識別三角形
$$
f^*\NL_{Y/Z} \to \NL_{X/Z} \to \NL_{X/Y} \to f^*\NL_{Y/Z}[1]
$$
があり、これは Lemma \ref{more-morphisms-lemma-exact-sequence-NL} の $6$-項完全列を復元する。
\end{lemma}

\begin{proof}
Modules, Lemma \ref{modules-lemma-exact-sequence-NL-ringed-topoi} の標準写像
$$
f^*\NL_{Y/Z} \to \text{Cone}(\NL_{X/Z} \to \NL_{X/Y})[-1]
$$
が $D(\mathcal{O}_X)$ で同型であることを示せば十分である。
これを示すには、$6$-項列の左端が零であること、すなわち
$H^{-1}(f^*\NL_{Y/Z}) \to H^{-1}(\NL_{X/Z})$ が単射であることを示せば十分である。
アフィン局所的には、これは More on Algebra, Lemma
\ref{more-algebra-lemma-transitive-lci-at-end} の対応する代数的結果から従う。
代数へ翻訳するには Lemma \ref{more-morphisms-lemma-NL-affine} を用いる。
\end{proof}

\begin{lemma}
\label{more-morphisms-lemma-smooth-etale-permanence}
$X \to Y \to Z$ をスキームの射とする。$X \to Z$ は滑らかで、
$Y \to Z$ はエタールであると仮定する。このとき $X \to Y$ は滑らかである。
% P05-EN-CH37-0064
\end{lemma}

\begin{proof}
Morphisms, Lemma \ref{morphisms-lemma-finite-presentation-permanence} により
射 $X \to Y$ は局所有限表示である。Lemma \ref{more-morphisms-lemma-NL-smooth} により
$H^{-1}(\NL_{X/Z}) = 0$ であり、加群 $\Omega_{X/Z}$ は有限局所自由である。
Lemma \ref{more-morphisms-lemma-NL-etale} により
$H^{-1}(\NL_{Y/Z}) = H^0(\NL_{Y/Z}) = 0$ である。
Lemma \ref{more-morphisms-lemma-exact-sequence-NL} により
$H^{-1}(\NL_{X/Y}) = 0$ を得て、
$\Omega_{X/Y} \cong \Omega_{X/Z}$ は有限局所自由である。
Lemma \ref{more-morphisms-lemma-NL-smooth} により射 $X \to Y$ は滑らかである。
\end{proof}

\begin{lemma}
\label{more-morphisms-lemma-get-NL}
$f : X \to Y$ をスキームの射とし、$i$ は埋め込み、$g : P \to Y$ は
形式的滑らか（例えば滑らか）であって、$f = g \circ i$ と分解するとする。
このとき $D(\mathcal{O}_X)$ に標準同型
$$
\NL_{X/Y} \cong \left(\mathcal{C}_{X/P} \to i^*\Omega_{P/Y}\right)
$$
がある。ここで余法層 $\mathcal{C}_{X/P}$ は次数 $-1$ に置かれている。
\end{lemma}

\begin{proof}
（括弧内の主張については Lemma \ref{more-morphisms-lemma-smooth-formally-smooth} を参照せよ。）
Lemmas \ref{more-morphisms-lemma-NL-immersion} および
\ref{more-morphisms-lemma-NL-formally-smooth} により
$\NL_{X/P} = \mathcal{C}_{X/P}[1]$、$\NL_{P/Y} = \Omega_{P/Y}$ であり、
$\Omega_{P/Y}$ は局所射影である。これは
$i^*\NL_{P/Y} \to i^*\Omega_{P/Y}$ も擬同型であることを意味する
（細部は省略する。その理由は $i^*\NL_{P/Y}$ が
$\tau_{\geq -1}Li^*\NL_{P/Y}$ と同じものであることにある。
More on Algebra, Lemma \ref{more-algebra-lemma-tensor-NL} を参照せよ）。
したがって Modules, Lemma
\ref{modules-lemma-exact-sequence-NL-ringed-topoi} の標準写像
$$
i^*\NL_{P/Y} \to \text{Cone}(\NL_{X/Y} \to \NL_{X/P})[-1]
$$
は $D(\mathcal{O}_X)$ で同型である。実際、上で述べたことによりコホモロジー群
$H^{-1}(i^*\NL_{P/Y})$ は零である。言い換えれば、識別三角形
$$
i^*\NL_{P/Y} \to \NL_{X/Y} \to \NL_{X/P} \to i^*\NL_{P/Y}[1]
$$
がある。明らかに、これは $\NL_{X/Y}$ が写像
$\NL_{X/P}[-1] \to i^*\NL_{P/Y}$ の写像錐であることを意味する。
上で与えた導来圏内のこれらの対象のコホモロジー層の計算により、これは補題の主張と
同値である。
\end{proof}

\begin{lemma}
\label{more-morphisms-lemma-base-change-NL}
スキームの Cartesian 図式
$$
\xymatrix{
X' \ar[r]_{g'} \ar[d] & X \ar[d] \\
Y' \ar[r] & Y
}
$$
を考える。標準写像 $(g')^*\NL_{X/Y} \to \NL_{X'/Y'}$ は
$H^0$ 上で同型を、$H^{-1}$ 上で全射を誘導する。
\end{lemma}

\begin{proof}
代数に翻訳すると、これは More on Algebra, Lemma
\ref{more-algebra-lemma-base-change-NL} である。
翻訳には Lemma \ref{more-morphisms-lemma-NL-affine} を用いる。
\end{proof}

\begin{lemma}
\label{more-morphisms-lemma-flat-base-change-NL}
スキームの Cartesian 図式
$$
\xymatrix{
X' \ar[d] \ar[r]_{g'} & X \ar[d] \\
Y' \ar[r] & Y
}
$$
を考える。$Y' \to Y$ が平坦ならば、標準写像
$(g')^*\NL_{X/Y} \to \NL_{X'/Y'}$ は擬同型である。
\end{lemma}

\begin{proof}
Lemma \ref{more-morphisms-lemma-NL-affine} により、これは Algebra, Lemma
\ref{algebra-lemma-change-base-NL} から従う。
\end{proof}

\begin{lemma}
\label{more-morphisms-lemma-base-change-NL-flat}
スキームの Cartesian 図式
$$
\xymatrix{
X' \ar[r]_{g'} \ar[d] & X \ar[d] \\
Y' \ar[r] & Y
}
$$
を考える。$X \to Y$ が平坦ならば、標準写像
$(g')^*\NL_{X/Y} \to \NL_{X'/Y'}$ は擬同型である。
さらに $\NL_{X/Y}$ が $[-1, 0]$ に Tor 振幅をもつならば、
$L(g')^*\NL_{X/Y} \to \NL_{X'/Y'}$ も擬同型である。
\end{lemma}

\begin{proof}
代数に翻訳すると、これは More on Algebra, Lemma
\ref{more-algebra-lemma-base-change-NL-flat} である。
翻訳には Lemma \ref{more-morphisms-lemma-NL-affine} および
Derived Categories of Schemes, Lemmas
\ref{perfect-lemma-affine-compare-bounded} と
\ref{perfect-lemma-tor-dimension-affine} を用いる。
\end{proof}











\section{スキームの圏における押し出し I}
\label{more-morphisms-section-pushouts}

\noindent
この節では、$X \to Y$ がアフィンで $X \to X'$ が肥厚化である場合に、
$Y \leftarrow X \rightarrow X'$ の押し出しを構成する。
これは実際に重要な場合となるため、詳しく論じる価値がある。
Section \ref{more-morphisms-section-pushouts-II} では、より興味深く、より難しい場合を論じる。
任意の圏における押し出しの一般論については Categories, Section
\ref{categories-section-pushouts} を参照せよ。

\begin{lemma}
\label{more-morphisms-lemma-basic-example-pushout}
$A' \to A$ を環の全射、$B \to A$ を環準同型とする。
$B' = B \times_A A'$ を環のファイバー積とし、
$S = \Spec(A)$、$S' = \Spec(A')$、$T = \Spec(B)$、
$T' = \Spec(B')$ とおく。このとき
$$
\vcenter{
\xymatrix{
S \ar[r]_i \ar[d]_f & S' \ar[d]^{f'} \\
T \ar[r]^{i'} & T'
}
}
\quad\text{に対応する}\quad
\vcenter{
\xymatrix{
A & A' \ar[l] \\
B \ar[u] & B' \ar[l] \ar[u]
}
}
$$
はスキームの押し出しである。
\end{lemma}

\begin{proof}
More on Algebra, Lemma \ref{more-algebra-lemma-points-of-fibre-product} により、
位相空間として $T' = T \amalg_S S'$ である。すなわち、この図式は位相空間の圏に
おける押し出しである。次に写像
$$
((i')^\sharp, (f')^\sharp) :
\mathcal{O}_{T'}
\longrightarrow
i'_*\mathcal{O}_T \times_{g_*\mathcal{O}_S} f'_*\mathcal{O}_{S'}
$$
を考える。ここで $g = i' \circ f = f' \circ i$ である。この写像は環の層の同型であると
主張する。実際、両辺を準連接 $\mathcal{O}_{T'}$-加群とみなすことができ
（右辺については Schemes,
% P05-EN-CH37-0065
Lemmas \ref{schemes-lemma-push-forward-quasi-coherent} を用いる）、写像は
$\mathcal{O}_{T'}$-線型である。したがって大域切断の水準で同型であることを示せば
十分である（Schemes, Lemma
\ref{schemes-lemma-equivalence-quasi-coherent}）。大域切断上では、補題の主張にある
同一視 $B' \to B \times_A A'$ を復元する（これは $B'$ の選び方そのものである）。

\medskip\noindent
$X$ をスキームとする。$m' \circ i = n \circ f$（これを $m$ と呼ぶ）を満たす
スキームの射 $m' : S' \to X$ と $n : T \to X$ が与えられたとする。
$T' = T \amalg_S S'$ なので（上を参照）、$m'$ と $n$ と整合する一意な位相空間の写像
$n' : T' \to X$ を得る。前段落の $\mathcal{O}_{T'}$ の記述により、
$(m')^\sharp$ と $n^\sharp$ から与えられる一意な環の層の準同型
% P05-EN-CH37-0066
$$
(n')^\sharp :
\mathcal{O}_X
\longrightarrow
(n')_*\mathcal{O}_{T'} =
m'_*\mathcal{O}_T \times_{m_*\mathcal{O}_T} n_*\mathcal{O}_S
$$
を得る。したがって $(n', (n')^\sharp)$ は $m'$ と $n$ と整合する一意な環付き空間の射
$T' \to X$ である。証明を終えるには、$n'$ がスキームの射、すなわち局所環付き空間の
射であることを示せば十分である。

\medskip\noindent
$t' \in T'$ の像を $x \in X$ とする。
$\mathcal{O}_{X, x} \to \mathcal{O}_{T', t'}$ が局所的であることを示さなければならない。
$t' \not \in T$ ならば、$t'$ は一意な点 $s' \in S'$ の像であり、
$\mathcal{O}_{T', t'} = \mathcal{O}_{S', s'}$ である。実際、$B' \to A'$ は同型
$\Ker(B' \to B) = \Ker(A' \to A)$ を誘導するので、
$S' \setminus S \to T' \setminus T$ はスキームの同型である。
$t'$ が $t \in T$ の像ならば、合成
$\mathcal{O}_{X, x} \to \mathcal{O}_{T', t'} \to \mathcal{O}_{T, t}$
は局所的であることが分かっており、同様に結論する。
% P05-EN-CH37-0067
\end{proof}

\begin{lemma}
\label{more-morphisms-lemma-pushout-fpqc-local}
$\mathcal{I} \to (\Sch/S)_{fppf}$、$i \mapsto X_i$ をスキームの図式とする。
$(W, X_i \to W)$ をスキームの圏におけるこの図式の余錐とする
（Categories, Remark \ref{categories-remark-cones-and-cocones}）。
次を満たすスキームの fpqc 被覆 $\{W_a \to W\}_{a \in A}$ が存在するとする。
\begin{enumerate}
\item すべての $a \in A$ に対して
$W_a = \colim X_i \times_W W_a$
がスキームの圏で成り立つ。
\item すべての $a, b \in A$ に対して
$W_a \times_W W_b = \colim X_i \times_W W_a \times_W W_b$
がスキームの圏で成り立つ。
\end{enumerate}
このとき $W = \colim X_i$ がスキームの圏で成り立つ。
\end{lemma}

\begin{proof}
実際、スキーム $T$ に対し、射 $W \to T$ を与えることは、重なり
$W_a \times_W W_b$ 上で一致する射 $W_a \to T$、$a \in A$ の族を与えることと同じである。
% P05-EN-CH37-0068
Descent, Lemma
\ref{descent-lemma-fpqc-universal-effective-epimorphisms} を参照せよ。
\end{proof}

\begin{lemma}
\label{more-morphisms-lemma-pushout-along-thickening}
$X \to X'$ をスキームの肥厚化、$X \to Y$ をアフィンなスキームの射とする。
このときスキームの圏に押し出し
$$
\xymatrix{
X \ar[r] \ar[d]_f
&
X' \ar[d]^{f'}
\\
Y \ar[r]
&
Y'
}
$$
が存在する。さらに $Y \subset Y'$ は肥厚化であり、
$X = Y \times_{Y'} X'$ であり、
$$
\mathcal{O}_{Y'} = \mathcal{O}_Y \times_{f_*\mathcal{O}_X} f'_*\mathcal{O}_{X'}
$$
が $|Y| = |Y'|$ 上の層として成り立つ。
\end{lemma}

\begin{proof}
まず $Y'$ を環付き空間として構成する。位相空間としては $Y' = Y$ とする。
% P05-EN-CH37-0069
$f' : X' \to Y'$ を $f$ と等しい位相空間の写像と書く。
構造層 $\mathcal{O}_{Y'}$ として補題の等式の右辺をとる。
$Y'$ がスキームであることを示すには、各点がアフィン近傍をもつことを示さなければ
ならない。層のファイバー積の形成は開集合への制限と可換なので、$Y$ はアフィンであると
仮定してよい。このとき $f$ はアフィンなので $X$ はアフィンであり、$X'$ もアフィンである
（Lemma \ref{more-morphisms-lemma-thickening-affine-scheme} を参照）。
$Y \leftarrow X \rightarrow X'$ が $B \rightarrow A \leftarrow A'$ に対応するとする。
$B' = B \times_A A'$ とおく。これは $\mathcal{O}_{Y'}$ の大域切断である。
$A' \to A$ は局所冪零な核をもつ全射なので、$B' \to B$ も局所冪零な核をもつ全射である。
したがって位相空間として $\Spec(B') = \Spec(B)$ である。
$Y' = \Spec(B')$ と主張する。これを示すため、$g' \in B'$ の像を
$g \in B$ として、$\mathcal{O}_{Y'}(D(g)) = B'_{g'}$ を示す。
More on Algebra, Lemma \ref{more-algebra-lemma-diagram-localize} により
$$
(B')_{g'} = B_g \times_{A_h} A'_{h'}
$$
である。ここで $h \in A$ と $h' \in A'$ は $g'$ の像である。
$B_g$、$A_h$、$A'_{h'}$ はそれぞれ $\mathcal{O}_Y(D(g))$、
$f_*\mathcal{O}_X(D(g))$、$f'_*\mathcal{O}_{X'}(D(g))$ に等しいので、主張が従う。

\medskip\noindent
残るのは $Y'$ が押し出しであることを示すことである。上の議論により、スキーム $Y'$ は
アフィン開被覆 $Y' = \bigcup W'_i$ をもち、対応する開集合
$U'_i \subset X'$、$W_i \subset Y$、$U_i \subset X$ はアフィン開集合である。
さらに $A'_i$, $B_i$, $A_i$ をそれぞれ $U'_i$, $W_i$, $U_i$ に対応する環とすると、
$W'_i$ は $B_i \times_{A_i} A'_i$ に対応する。
したがって Lemmas \ref{more-morphisms-lemma-basic-example-pushout} および
\ref{more-morphisms-lemma-pushout-fpqc-local} を適用して、我々の構成がスキームの圏における押し出しで
あると結論する。
\end{proof}

\noindent
次の補題では Categories, Example
\ref{categories-example-2-fibre-product-categories} で定義された圏のファイバー積を用いる。

\begin{lemma}
\label{more-morphisms-lemma-equivalence-categories-schemes-over-pushout}
$X \to X'$ をスキームの肥厚化、$X \to Y$ をアフィンなスキームの射とする。
$Y' = Y \amalg_X X'$ を押し出しとする（Lemma
\ref{more-morphisms-lemma-pushout-along-thickening} を参照）。基底変換は関手
$$
F : (\Sch/Y') \longrightarrow (\Sch/Y) \times_{(\Sch/X)} (\Sch/X')
$$
を与える。これは $V' \longmapsto (V' \times_{Y'} Y, V' \times_{Y'} X', 1)$
によって与えられ、左随伴
$$
G : (\Sch/Y) \times_{(\Sch/X)} (\Sch/X') \longrightarrow (\Sch/Y')
$$
をもつ。これは三つ組 $(V, U', \varphi)$ を押し出し
$V \amalg_{(V \times_Y X)} U'$ に送る。最後に、$F \circ G$ は恒等関手と同型である。
\end{lemma}

\begin{proof}
$(V, U', \varphi)$ をファイバー積圏の対象とする。
$U = U' \times_{X'} X$ とおく。$U \to U'$ は肥厚化であることに注意する。
$\varphi : V \times_Y X \to U' \times_{X'} X = U$ は同型なので、
$X \to Y$ 上の射 $U \to V$ があり、これは $U$ をファイバー積
$X \times_Y V$ と同一視する。特に $U \to V$ はアフィンである。
Morphisms, Lemma \ref{morphisms-lemma-base-change-affine} を参照せよ。
したがって Lemma \ref{more-morphisms-lemma-pushout-along-thickening} を適用して押し出し
$V' = V \amalg_U U'$ を得る。$V'$ が押し出しであり、$Y'$ への射として $U$ 上で一致する
射 $V \to Y$ と $U' \to X'$ が与えられていることから得られる射
% P05-EN-CH37-0070
$V' \to Y'$ をそのまま書く。$G(V, U', \varphi) = V'$ とおくと関手 $G$ を得る。

\medskip\noindent
$G$ が $F$ の左随伴であることを証明しよう。$Z$ を $Y'$ 上のスキームとする。
$$
\Mor(V', Z) = \Mor((V, U', \varphi), F(Z))
$$
を示さなければならない。ここで射集合はそれぞれの圏でとる。
% P05-EN-CH37-0071
射 $g' : V' \to Z$ をとる。$g'$ と射 $V \to V'$、および $U' \to V'$ との合成を
% P05-EN-CH37-0072
それぞれ $\tilde g$、$\tilde f'$ と書く。
$\tilde g$、$\tilde f'$ をそれぞれ $Y \to Y'$、$X' \to Y'$ で基底変換すると、射
$g : V \to Z \times_{Y'} Y$、$f' : U' \to Z \times_{Y'} X'$ を得る。
すると $(g, f')$ は上の等式の右辺の元である（細部は省略する）。
逆に、$(g, f') : (V, U', \varphi) \to F(Z)$ が右辺の元であるとする。
% P05-EN-CH37-0073
$g$、$f$ と、それぞれ $Z \times_{Y'} X' \to Z$、
$Z \times_{Y'} Y \to Z$ との合成
$\tilde g : V \to Z$、$\tilde f' : U' \to Z$ を考えられる。
このとき $\tilde g$ と $\tilde f'$ は $U$ から $Z$ への射として一致する。
% P05-EN-CH37-0074
押し出しの普遍性により射 $g' : V' \to Z$、すなわち左辺の元を得る。
これらの構成が互いに逆であることの確認は省略する。

\medskip\noindent
$F \circ G$ が恒等関手と同型であることを証明するには、随伴写像
$(V, U', \varphi) \to F(G(V, U', \varphi))$ が同型であることを示さなければならない。
これはアフィン局所的に調べてよい。$X = \Spec(A)$、$X' = \Spec(A')$、
$Y = \Spec(B)$ とする。このとき $A' \to A$ と $B \to A$ は More on Algebra,
Lemma \ref{more-algebra-lemma-module-over-fibre-product} にある環準同型であり、
$B' = B \times_A A'$ として $Y' = \Spec(B')$ である。次に
$V = \Spec(D)$、$U' = \Spec(C')$ とし、$\varphi$ は $A$-代数同型
$D \otimes_B A \to C' \otimes_{A'} A = C'/IC'$ から与えられるとする。
$D' = D \times_{C'/IC'} C'$ とおく。この場合、示すべき主張は
$D' \otimes_{B'} B \cong D$ および $D' \otimes_{B'} A' \cong C'$ である。
これは More on Algebra, Lemma
\ref{more-algebra-lemma-module-over-fibre-product} の特別な場合である。
\end{proof}

\begin{lemma}
\label{more-morphisms-lemma-scheme-over-pushout-flat-modules}
$X \to X'$ をスキームの肥厚化、$X \to Y$ をアフィンなスキームの射とする。
$Y' = Y \amalg_X X'$ を押し出しとする（Lemma
\ref{more-morphisms-lemma-pushout-along-thickening} を参照）。$V' \to Y'$ をスキームの射とする。
$V = Y \times_{Y'} V'$、$U' = X' \times_{Y'} V'$、
$U = X \times_{Y'} V'$ とおく。次の圏の間に同値がある。
\begin{enumerate}
\item $Y'$ 上平坦な準連接 $\mathcal{O}_{V'}$-加群。
\item 次を満たす三つ組 $(\mathcal{G}, \mathcal{F}', \varphi)$ の圏。
\begin{enumerate}
\item $\mathcal{G}$ は $Y$ 上平坦な準連接 $\mathcal{O}_V$-加群である。
\item $\mathcal{F}'$ は $X'$ 上平坦な準連接 $\mathcal{O}_{U'}$-加群である。
\item $\varphi : (U \to V)^*\mathcal{G} \to (U \to U')^*\mathcal{F}'$
は $\mathcal{O}_U$-加群の同型である。
\end{enumerate}
\end{enumerate}
この同値は $\mathcal{G}'$ を
$((V \to V')^*\mathcal{G}', (U' \to V')^*\mathcal{G}', can)$ に写す。
$\mathcal{G}'$ が三つ組 $(\mathcal{G}, \mathcal{F}', \varphi)$ に対応するとする。
このとき次が成り立つ。
\begin{enumerate}
\item[(a)] $\mathcal{G}'$ が有限型 $\mathcal{O}_{V'}$-加群であることと、
% P05-EN-CH37-0075
$\mathcal{G}$ と $\mathcal{F}'$ が有限型 $\mathcal{O}_Y$-加群および
$\mathcal{O}_{U'}$-加群であることは同値である。
\item[(b)] $V' \to Y'$ が局所有限表示ならば、$\mathcal{G}'$ が有限表示
$\mathcal{O}_{V'}$-加群であることと、$\mathcal{G}$ と $\mathcal{F}'$ が有限表示
$\mathcal{O}_Y$-加群および $\mathcal{O}_{U'}$-加群であることは同値である。
\end{enumerate}
\end{lemma}

\begin{proof}
擬逆関手は三つ組 $(\mathcal{G}, \mathcal{F}', \varphi)$ にファイバー積
$$
(V \to V')_*\mathcal{G}
\times_{(U \to V')_*\mathcal{F}}
(U' \to V')_*\mathcal{F}'
$$
を対応させる。ここで $\mathcal{F} = (U \to U')^*\mathcal{F}'$ である。
これがうまくいくのは、アフィン上で More on Algebra, Lemma
\ref{more-algebra-lemma-relative-flat-module-over-fibre-product} の同値を復元するからである。
いくつかの細部は省略する。

\medskip\noindent
(a) と (b) は More on Algebra, Lemmas
\ref{more-algebra-lemma-relative-finite-module-over-fibre-product} and
\ref{more-algebra-lemma-relative-finitely-presented-module-over-fibre-product} から従う。
\end{proof}

\begin{lemma}
\label{more-morphisms-lemma-equivalence-categories-schemes-over-pushout-flat}
Lemma \ref{more-morphisms-lemma-equivalence-categories-schemes-over-pushout} の状況を考える。
ある三つ組 $(V, U', \varphi)$ に対して $V' = G(V, U', \varphi)$ ならば、
\begin{enumerate}
\item $V' \to Y'$ が局所有限型であることと、$V \to Y$ と $U' \to X'$ が
局所有限型であることは同値である。
\item $V' \to Y'$ が平坦であることと、$V \to Y$ と $U' \to X'$ が平坦であることは
同値である。
\item $V' \to Y'$ が平坦かつ局所有限表示であることと、$V \to Y$ と $U' \to X'$ が
平坦かつ局所有限表示であることは同値である。
\item $V' \to Y'$ が滑らかであることと、$V \to Y$ と $U' \to X'$ が滑らかであることは
同値である。
\item $V' \to Y'$ がエタールであることと、$V \to Y$ と $U' \to X'$ がエタールであることは
同値である。
% P05-EN-CH37-0076
\item 必要に応じてここにさらに追加する。
\end{enumerate}
$W'$ が $Y'$ 上平坦ならば、随伴写像 $G(F(W')) \to W'$ は同型である。
したがって $F$ と $G$ は、$Y'$ 上平坦なスキームの圏と、$V \to Y$ と
$U' \to X'$ が平坦である三つ組 $(V, U', \varphi)$ の圏との間の互いに擬逆な
関手を定める。
\end{lemma}

\begin{proof}
アフィン片上で調べれば、この補題の主張は More on Algebra, Lemma
\ref{more-algebra-lemma-properties-algebras-over-fibre-product} の対応する主張と
同値である。
\end{proof}







\section{平坦軌跡の開性}
\label{more-morphisms-section-open-flat}

\noindent
この結果を証明するには多少の作業を要し、（おそらく）
独立した節を設けるに値する。以下がその結果である。

\begin{theorem}
\label{more-morphisms-theorem-openness-flatness}
\begin{reference}
\cite[IV Theorem 11.3.1]{EGA}
\end{reference}
$S$ をスキームとする。
$f : X \to S$ を局所有限表示な射とする。
$\mathcal{F}$ を局所有限表示な準連接 $\mathcal{O}_X$-加群とする。このとき
$$
U = \{x \in X \mid \mathcal{F}\text{ は点 }x\text{ において }S\text{ 上平坦}\}
$$
は $X$ の開集合である。
\end{theorem}

\begin{proof}
開性は $X$ 上局所的に判定できるので、$f$ はアフィンスキームの射であると仮定してよい。
この場合、定理はちょうど
Algebra, Theorem \ref{algebra-theorem-openness-flatness}
である。
\end{proof}

\begin{lemma}
\label{more-morphisms-lemma-flat-locus-base-change}
$S$ をスキームとする。次の図式
$$
\xymatrix{
X' \ar[r]_{g'} \ar[d]_{f'} & X \ar[d]^f \\
S' \ar[r]^g & S
}
$$
がスキームのデカルト図式であるとする。
$\mathcal{F}$ を準連接 $\mathcal{O}_X$-加群とする。
$x' \in X'$ とし、その像を
$x = g'(x')$ および $s' = f'(x')$ とする。
\begin{enumerate}
\item $\mathcal{F}$ が点 $x$ において $S$ 上平坦ならば、
$(g')^*\mathcal{F}$ は点 $x'$ において $S'$ 上平坦である。
\item $g$ が点 $s'$ において平坦であり、$(g')^*\mathcal{F}$ が点
$x'$ において $S'$ 上平坦ならば、$\mathcal{F}$ は点 $x$ において $S$ 上平坦である。
\end{enumerate}
特に、$g$ が平坦で、$f$ が局所有限表示であり、
$\mathcal{F}$ が局所有限表示ならば、
Theorem \ref{more-morphisms-theorem-openness-flatness}
の開集合の形成は基底変換と可換である。
\end{lemma}

\begin{proof}
局所環の可換図式
$$
\xymatrix{
\mathcal{O}_{X', x'} & \mathcal{O}_{X, x} \ar[l] \\
\mathcal{O}_{S', s'} \ar[u] & \mathcal{O}_{S, s} \ar[l] \ar[u]
}
$$
を考える。$\mathcal{O}_{X', x'}$ は
$\mathcal{O}_{X, x} \otimes_{\mathcal{O}_{S, s}} \mathcal{O}_{S', s'}$
の局所化であり、$((g')^*\mathcal{F})_{x'}$ は
$\mathcal{F}_x \otimes_{\mathcal{O}_{X, x}} \mathcal{O}_{X', x'}$
に等しいことに注意する。したがって補題は
Algebra, Lemma \ref{algebra-lemma-base-change-flat-up-down}
から従う。
\end{proof}






\section{ファイバーによる平坦性判定法}
\label{more-morphisms-section-criterion-flat-fibres}

\noindent
スキームの可換図式（左の図式）
$$
\xymatrix{
X \ar[rr]_f \ar[dr] & & Y \ar[dl] \\
& S
}
\quad
\xymatrix{
X_s \ar[rr]_{f_s} \ar[rd] & & Y_s \ar[dl] \\
& \Spec(\kappa(s))
}
$$
と準連接 $\mathcal{O}_X$-加群 $\mathcal{F}$ を考える。
$s \in S$ 上にある点 $x \in X$ を取り、その像を $y = f(x)$ とする。
ここで、$\mathcal{F}$ は点 $x$ において $Y$ 上平坦か、という問題を考える。
$\mathcal{F}$ が点 $x$ において $S$ 上平坦ならば、この問題は、
$\mathcal{F}$ のファイバーへの制限
$$
\mathcal{F}_s = (X_s \to X)^*\mathcal{F}
$$
が点 $x$ において $Y_s$ 上平坦かという問題と密接に関係する、と定理は述べる。
以下では「Noetherian」版と「有限表示」版を与える。
「冪零」版は先に扱ったので、
Lemma \ref{more-morphisms-lemma-flatness-morphism-thickenings}
を参照されたい。

\begin{theorem}
\label{more-morphisms-theorem-criterion-flatness-fibre-Noetherian}
$S$ をスキームとする。
$f : X \to Y$ を $S$ 上のスキームの射とする。
$\mathcal{F}$ を準連接 $\mathcal{O}_X$-加群とする。
$x \in X$ とし、$y = f(x)$ とおき、$s \in S$ を $S$ における $x$ の像とする。
% P05-EN-CH37-0077: frozen source omits "are" in the locally Noetherian hypothesis; mathematical content is unchanged.
$S$, $X$, $Y$ は局所 Noetherian、$\mathcal{F}$ は連接、かつ
$\mathcal{F}_x \not = 0$ と仮定する。このとき次は同値である。
\begin{enumerate}
\item $\mathcal{F}$ は点 $x$ において $S$ 上平坦であり、かつ
$\mathcal{F}_s$ は点 $x$ において $Y_s$ 上平坦である。
\item $Y$ は点 $y$ において $S$ 上平坦であり、かつ
$\mathcal{F}$ は点 $x$ において $Y$ 上平坦である。
\end{enumerate}
\end{theorem}

\begin{proof}
環準同型
$$
\mathcal{O}_{S, s} \longrightarrow
\mathcal{O}_{Y, y} \longrightarrow \mathcal{O}_{X, x}
$$
と加群 $\mathcal{F}_x$ を考える。点 $x$ における $\mathcal{F}_s$ の茎は
$\mathcal{F}_x/\mathfrak m_s \mathcal{F}_x$ であり、点 $y$ における
$Y_s$ の局所環は
$\mathcal{O}_{Y, y}/\mathfrak m_s \mathcal{O}_{Y, y}$
である。したがって (1) $\Rightarrow$ (2) は
Algebra, Lemma \ref{algebra-lemma-criterion-flatness-fibre-Noetherian}
である。(2) が成り立つならば、最初の環準同型は忠実平坦であり、
$\mathcal{F}_x$ は $\mathcal{O}_{Y, y}$ 上平坦であるから、
Algebra, Lemma \ref{algebra-lemma-composition-flat}
により $\mathcal{F}_x$ は $\mathcal{O}_{S, s}$ 上平坦である。
さらに、$\mathcal{F}_x/\mathfrak m_s \mathcal{F}_x$ は、平坦加群
$\mathcal{F}_x$ を
$\mathcal{O}_{Y, y} \to \mathcal{O}_{Y, y}/\mathfrak m_s \mathcal{O}_{Y, y}$
によって基底変換したものなので、
Algebra, Lemma \ref{algebra-lemma-flat-base-change}
により平坦である。
\end{proof}

\noindent
次が非 Noetherian 版である。

\begin{theorem}
\label{more-morphisms-theorem-criterion-flatness-fibre}
$S$ をスキームとする。
$f : X \to Y$ を $S$ 上のスキームの射とする。
$\mathcal{F}$ を準連接 $\mathcal{O}_X$-加群とする。次を仮定する。
\begin{enumerate}
\item $X$ は $S$ 上局所有限表示である。
% P05-EN-CH37-0078: frozen source omits "is" in the finite-presentation hypothesis; mathematical content is unchanged.
\item $\mathcal{F}$ は有限表示な $\mathcal{O}_X$-加群である。
\item $Y$ は $S$ 上局所有限型である。
\end{enumerate}
$x \in X$ とする。$y = f(x)$ とおき、$s \in S$ を $S$ における $x$ の像とする。
$\mathcal{F}_x \not = 0$ ならば、次は同値である。
\begin{enumerate}
\item $\mathcal{F}$ は点 $x$ において $S$ 上平坦であり、かつ
$\mathcal{F}_s$ は点 $x$ において $Y_s$ 上平坦である。
\item $Y$ は点 $y$ において $S$ 上平坦であり、かつ
$\mathcal{F}$ は点 $x$ において $Y$ 上平坦である。
\end{enumerate}
さらに、(1) と (2) が成り立つ点 $x$ の集合は
$\text{Supp}(\mathcal{F})$ において開である。
\end{theorem}

\begin{proof}
環準同型
$$
\mathcal{O}_{S, s} \longrightarrow
\mathcal{O}_{Y, y} \longrightarrow \mathcal{O}_{X, x}
$$
と加群 $\mathcal{F}_x$ を考える。点 $x$ における $\mathcal{F}_s$ の茎は
$\mathcal{F}_x/\mathfrak m_s \mathcal{F}_x$ であり、点 $y$ における
$Y_s$ の局所環は
$\mathcal{O}_{Y, y}/\mathfrak m_s \mathcal{O}_{Y, y}$
である。したがって (1) $\Rightarrow$ (2) は
Algebra, Lemma \ref{algebra-lemma-criterion-flatness-fibre-fp-over-ft}
である。(2) が成り立つならば、最初の環準同型は忠実平坦であり、
$\mathcal{F}_x$ は $\mathcal{O}_{Y, y}$ 上平坦であるから、
Algebra, Lemma \ref{algebra-lemma-composition-flat}
により $\mathcal{F}_x$ は $\mathcal{O}_{S, s}$ 上平坦である。
さらに、$\mathcal{F}_x/\mathfrak m_s \mathcal{F}_x$ は、平坦加群
$\mathcal{F}_x$ を
$\mathcal{O}_{Y, y} \to \mathcal{O}_{Y, y}/\mathfrak m_s \mathcal{O}_{Y, y}$
によって基底変換したものなので、
Algebra, Lemma \ref{algebra-lemma-flat-base-change}
により平坦である。

\medskip\noindent
Morphisms, Lemma \ref{morphisms-lemma-finite-presentation-permanence}
により、射 $f$ は局所有限表示である。集合
\begin{equation}
\label{more-morphisms-equation-open}
U = \{x \in X \mid \mathcal{F} \text{ は点 }x
\text{ において }Y\text{ 上かつ }S\text{ 上平坦}\}.
\end{equation}
を考える。この集合は
Theorem \ref{more-morphisms-theorem-openness-flatness}
により $X$ の開集合である。$x \in U$ ならば、$\mathcal{F}_s$ は
射 $Y_s \to Y$ による平坦加群の基底変換として、点 $x$ において
$Y_s$ 上平坦であることに注意する。これについては
Morphisms, Lemma \ref{morphisms-lemma-base-change-module-flat}
を参照されたい。したがって $U \cap \text{Supp}(\mathcal{F})$ の各点で
条件 (1) が成り立つ。一方、$x \in \text{Supp}(\mathcal{F})$ が
(1) と (2) を満たすならば $x \in U$ であることは明らかである。
ゆえに求める開集合は $U \cap \text{Supp}(\mathcal{F})$ である。
\end{proof}

\noindent
これらの定理はしばしば次の簡略化した形で用いられる。
ここでは大域的な主張だけを与えるが、もちろん各点での版もある。

\begin{lemma}
\label{more-morphisms-lemma-morphism-between-flat-Noetherian}
$S$ をスキームとする。
$f : X \to Y$ を $S$ 上のスキームの射とする。次を仮定する。
\begin{enumerate}
\item $S$, $X$, $Y$ は局所 Noetherian である。
\item $X$ は $S$ 上平坦である。
\item 任意の $s \in S$ に対して、射
$f_s : X_s \to Y_s$ は平坦である。
\end{enumerate}
このとき $f$ は平坦である。さらに $f$ が全射ならば、$Y$ は $S$ 上平坦である。
\end{lemma}

\begin{proof}
これは Theorem \ref{more-morphisms-theorem-criterion-flatness-fibre-Noetherian}
の特別な場合である。
\end{proof}

\begin{lemma}
\label{more-morphisms-lemma-morphism-between-flat}
$S$ をスキームとする。
$f : X \to Y$ を $S$ 上のスキームの射とする。次を仮定する。
\begin{enumerate}
\item $X$ は $S$ 上局所有限表示である。
\item $X$ は $S$ 上平坦である。
\item 任意の $s \in S$ に対して、射
$f_s : X_s \to Y_s$ は平坦である。
\item $Y$ は $S$ 上局所有限型である。
\end{enumerate}
このとき $f$ は平坦である。さらに $f$ が全射ならば、$Y$ は $S$ 上平坦である。
\end{lemma}

\begin{proof}
これは Theorem \ref{more-morphisms-theorem-criterion-flatness-fibre}
の特別な場合である。
\end{proof}

\begin{lemma}
\label{more-morphisms-lemma-base-change-criterion-flatness-fibre}
$S$ をスキームとする。$f : X \to Y$ を $S$ 上のスキームの射とする。
$\mathcal{F}$ を準連接 $\mathcal{O}_X$-加群とする。次を仮定する。
\begin{enumerate}
\item $X$ は $S$ 上局所有限表示である。
% P05-EN-CH37-0079: frozen source omits "is" in the finite-presentation hypothesis; mathematical content is unchanged.
\item $\mathcal{F}$ は有限表示な $\mathcal{O}_X$-加群である。
\item $\mathcal{F}$ は $S$ 上平坦である。
\item $Y$ は $S$ 上局所有限型である。
\end{enumerate}
このとき集合
$$
% P05-EN-CH37-0081: frozen source inserts a full stop before the predicate "is open"; punctuation is preserved.
U = \{x \in X \mid \mathcal{F} \text{ は点 }x \text{ において }Y\text{ 上平坦}\}.
$$
は $X$ の開集合であり、その形成は任意の基底変換と可換である。
$S' \to S$ をスキームの射とし、$U'$ を、
% P05-EN-CH37-0080: frozen source writes a fibre product of a sheaf with S'; the intended pullback notation is not silently repaired.
$X' = X \times_S S'$ の点のうち $\mathcal{F}' = \mathcal{F} \times_S S'$
が $Y' = Y \times_S S'$ 上平坦である点全体とすれば、$U' = U \times_S S'$ である。
\end{lemma}

\begin{proof}
Morphisms, Lemma \ref{morphisms-lemma-finite-presentation-permanence}
により、射 $f$ は局所有限表示である。したがって
Theorem \ref{more-morphisms-theorem-openness-flatness}
により $U$ は開である。$\mathcal{F}$ は $S$ 上平坦であると仮定したので、
Theorem \ref{more-morphisms-theorem-criterion-flatness-fibre}
から
$$
% P05-EN-CH37-0082: frozen source ends the display with a full stop but continues with lowercase "where"; punctuation is preserved.
U = \{x \in X \mid \mathcal{F}_s \text{ は点 }x \text{ において }Y_s\text{ 上平坦}\}.
$$
を得る。ここで $s$ は常に $S$ における $x$ の像を表す。
（この記述は $\mathcal{F}_x = 0$ の場合にも自明に成り立つ。）
さらに、射 $f' : X' \to Y'$ と層 $\mathcal{F}'$ に対しても
補題の仮定は保たれる。したがって $U'$ も同様に記述される。
言い換えれば、$s \in S$ に写る $s' \in S'$ を与えたとき、
$$
\{x' \in X'_{s'} \mid
\mathcal{F}'_{s'} \text{ は点 }x' \text{ において }Y'_{s'}\text{ 上平坦}\}
$$
が $X_s$ の対応する軌跡の逆像であることを示せば十分である。
これは Lemma \ref{more-morphisms-lemma-flat-locus-base-change}
から従う。実際、デカルト図式
$$
\xymatrix{
X'_{s'} \ar[d] \ar[r] & X_s \ar[d] \\
Y'_{s'} \ar[r] & Y_s
}
$$
において、水平射は平坦射
$\Spec(\kappa(s')) \to \Spec(\kappa(s))$
による基底変換なので平坦である。
\end{proof}

\begin{lemma}
\label{more-morphisms-lemma-base-change-flatness-fibres}
$S$ をスキームとする。$f : X \to Y$ を $S$ 上のスキームの射とする。
次を仮定する。
\begin{enumerate}
\item $X$ は $S$ 上局所有限表示である。
\item $X$ は $S$ 上平坦である。
\item $Y$ は $S$ 上局所有限型である。
\end{enumerate}
このとき集合
$$
% P05-EN-CH37-0083: frozen source inserts a full stop before the predicate "is open"; punctuation is preserved.
U = \{x \in X \mid X\text{ は点 }x \text{ において }Y\text{ 上平坦}\}.
$$
は $X$ の開集合であり、その形成は任意の基底変換と可換である。
\end{lemma}

\begin{proof}
これは Lemma \ref{more-morphisms-lemma-base-change-criterion-flatness-fibre}
の特別な場合である。
\end{proof}

\noindent
次の補題は
Algebra, Lemma \ref{algebra-lemma-free-fibre-flat-free}
の変形である。$(\mathcal{F}_s)_x$ が平坦な
$\mathcal{O}_{X_s, x}$-加群であるという仮定は、
$(\mathcal{F}_s)_x$ が自由 $\mathcal{O}_{X_s, x}$-加群であることを意味する。
$x \in X_s$ が $X_s$ の既約成分の一般点で、かつ $X_s$ が被約ならば、
これは常に成り立つ（実際、この場合 $\mathcal{O}_{X_s, x}$ は体である。
Algebra, Lemma \ref{algebra-lemma-minimal-prime-reduced-ring}
を参照せよ）。

\begin{lemma}
\label{more-morphisms-lemma-flat-and-free-at-point-fibre}
$f : X \to S$ を有限表示なスキームの射とする。
$\mathcal{F}$ を有限表示な $\mathcal{O}_X$-加群とする。
$x \in X$ とし、その像を $s \in S$ とする。
$\mathcal{F}$ が点 $x$ において $S$ 上平坦であり、$(\mathcal{F}_s)_x$ が平坦な
$\mathcal{O}_{X_s, x}$-加群ならば、$\mathcal{F}$ は $x$ のある近傍上で有限自由である。
\end{lemma}

\begin{proof}
$\mathcal{F}_x \otimes \kappa(x)$ が零ならば、Nakayama の補題
（Algebra, Lemma \ref{algebra-lemma-NAK}）により $\mathcal{F}_x = 0$ である。
したがって $\mathcal{F}$ は $x$ のある近傍で零であり
（Modules, Lemma \ref{modules-lemma-finite-type-stalk-zero}）、
補題が成り立つ。ゆえに $\mathcal{F}_x \otimes \kappa(x)$ は零でないと仮定してよい。
このとき Theorem \ref{more-morphisms-theorem-criterion-flatness-fibre}
を $f = \text{id} : X \to X$ として適用できる。したがって
$\mathcal{F}_x$ は $\mathcal{O}_{X, x}$ 上平坦である。ゆえに
$\mathcal{F}_x$ は自由である。例えば
Algebra, Lemma \ref{algebra-lemma-finite-flat-local}
を参照されたい。開近傍 $x \in U \subset X$ と、$\mathcal{F}_x$ の基底に写る切断
$s_1, \ldots, s_r \in \mathcal{F}(U)$ を選ぶ。対応する写像
$\psi : \mathcal{O}_U^{\oplus r} \to \mathcal{F}|_U$ は、$U$ を縮小すれば全射になる
（Modules, Lemma \ref{modules-lemma-finite-type-stalk-zero}）。
このとき $\Ker(\psi)$ は有限型であり（Modules, Lemma
\ref{modules-lemma-kernel-surjection-finite-free-onto-finite-presentation} 参照）、
$\Ker(\psi)_x = 0$ である。したがって $U$ をもう一度縮小すれば、
$\psi$ は同型になる。
\end{proof}

\begin{lemma}
\label{more-morphisms-lemma-finite-free-open}
$f : X \to S$ を局所有限表示なスキームの射とする。
$\mathcal{F}$ を $S$ 上平坦な有限表示 $\mathcal{O}_X$-加群とする。このとき集合
$$
\{x \in X : \mathcal{F}\text{ は }x\text{ のある近傍で自由}\}
$$
は $X$ の開集合であり、その形成は任意の基底変換
$S' \to S$ と可換である。
\end{lemma}

\begin{proof}
開性は明らかである。$x \in X$ とし、その像を $s \in S$ とする。
Lemma \ref{more-morphisms-lemma-flat-and-free-at-point-fibre}
により、$x$ がこの集合に属することと、$\mathcal{F}|_{X_s}$ が点 $x$ において
$X_s$ 上平坦であることは同値である。これは明らかに、$\mathcal{F}$ が点 $x$ において
$X$ 上平坦であることとも同値である（この条件は $\mathcal{F}_x$ の自由性から従い、
また $\mathcal{F}|_{X_s}$ が点 $x$ において $X_s$ 上平坦であることを含意する）。
したがって基底変換に関する主張は、$S$ 上の $\text{id} : X \to X$ に
Lemma \ref{more-morphisms-lemma-base-change-criterion-flatness-fibre}
を適用して従う。
\end{proof}






\section{滑らかなスキーム間の閉埋め込み}
\label{more-morphisms-section-closed-immersions-smooth}

\noindent
他の箇所にはうまく収まらないいくつかの結果を述べる。

\begin{lemma}
\label{more-morphisms-lemma-etale-local-structure}
$S$ をスキームとする。$Y \to X$ を、$S$ 上滑らかなスキームの
閉埋め込みとする。任意の $y \in Y$ に対して、整数
$0 \leq m, n$ と可換図式
% P05-EN-CH37-0084: frozen source omits the comma in the final zero-coordinate list; the formula is preserved.
$$
\xymatrix{
Y \ar[d] &
V \ar[l] \ar[d] \ar[r] &
\mathbf{A}^m_S
\ar[d]^{(a_1, \ldots, a_m) \mapsto (a_1, \ldots, a_m, 0 \ldots, 0)} \\
X &
U \ar[l] \ar[r]^-\pi &
\mathbf{A}^{m + n}_S
}
$$
が存在し、ここで $U \subset X$ は開、$V = Y \cap U$、
$\pi$ はエタール、$V = \pi^{-1}(\mathbf{A}^m_S)$、かつ $y \in V$ である。
\end{lemma}

\begin{proof}
問題は $X$ 上局所的なので、$X$ を $y$ の開近傍で置き換えてよい。
Divisors, Lemma
\ref{divisors-lemma-immersion-smooth-into-smooth-regular-immersion}
により $Y \to X$ は正則埋め込みであるから、$X = \Spec(A)$ はアフィンであり、
$Y = V(f_1, \ldots, f_n)$ となる正則列
$f_1, \ldots, f_n \in A$ が存在すると仮定してよい。
$X$（したがって $Y$）をさらに縮小すれば、エタール射
$Y \to \mathbf{A}^m_S$ が存在すると仮定できる。これについては
Morphisms, Lemma \ref{morphisms-lemma-smooth-etale-over-affine-space}
を参照されたい。$\overline{g}_1, \ldots, \overline{g}_m$ を
$\mathcal{O}_Y(Y)$ に属するこのエタール射の座標関数とする。これらの関数の持ち上げ
$g_1, \ldots, g_m \in A$ を選び、射
$$
(g_1, \ldots, g_m, f_1, \ldots, f_n) :
X
\longrightarrow
\mathbf{A}^{m + n}_S
$$
を $S$ 上で考える。これは $S$ 上局所有限表示なスキーム間の射なので、
局所有限表示である
（Morphisms, Lemma \ref{morphisms-lemma-finite-presentation-permanence}）。
構成により、この射の
$\mathbf{A}^m_S \subset \mathbf{A}^{m + n}_S$
への制限はエタールである。したがって
$X \to \mathbf{A}^{m + n}_S$ が点 $y$ においてエタールであることを示すには、
$X \to \mathbf{A}^{m + n}_S$ が点 $y$ において平坦であることを示せば十分である。これについては
Morphisms, Lemma \ref{morphisms-lemma-etale-at-point}
を参照されたい。$s \in S$ を $y$ の像とする。
Theorem \ref{more-morphisms-theorem-criterion-flatness-fibre}
により、$X_s \to \mathbf{A}^{m + n}_s$ が点 $y$ において平坦であることを
確認すれば十分である。$z \in \mathbf{A}^{m + n}_s$ を $y$ の像とする。
局所環準同型
$$
\mathcal{O}_{\mathbf{A}^{m + n}_s, z}
\longrightarrow
\mathcal{O}_{X_s, y}
$$
は Algebra, Lemma \ref{algebra-lemma-CM-over-regular-flat}
により平坦である。実際、体上滑らかなスキームは正則であり、正則環は
Cohen--Macaulay である。Varieties, Lemma \ref{varieties-lemma-smooth-regular}
および Algebra, Lemma \ref{algebra-lemma-regular-ring-CM}
を参照されたい。したがって始域と終域はいずれも正則局所環である
（ゆえに CM である）。始域と終域の次元は等しい。実際、
More on Algebra, Lemma \ref{more-algebra-lemma-dimension-etale-extension}
により
$\dim(\mathcal{O}_{Y_s, y}) = \dim(\mathcal{O}_{\mathbf{A}^m_s, z})$、
また $\dim(\mathcal{O}_{\mathbf{A}^{m + n}_s, z}) = n +
\dim(\mathcal{O}_{\mathbf{A}^m_s, z})$ であり、さらに
$\mathcal{O}_{Y_s, y}$ は $\mathcal{O}_{X_s, y}$ を長さ $n$ の正則列
$f_1, \ldots, f_n$ で割った商なので、
$\dim(\mathcal{O}_{X_s, y}) = n + \dim(\mathcal{O}_{Y_s, y})$
である（Divisors, Remark
\ref{divisors-remark-relative-regular-immersion-elements} 参照）。
最後に、$Y_s \to \mathbf{A}^m_s$ は点 $y$ においてエタールなので、
表示した矢印のファイバー環は $\kappa(z)$ 上有限である。
これで証明が完了する。
\end{proof}

\begin{remark}
\label{more-morphisms-remark-relative-blowup}
環 $R$ を固定し、$S = \Spec(R)$ とおく。整数 $0 \leq m$ および
$1 \leq n$ を固定する。閉埋め込み
% P05-EN-CH37-0084: the same frozen zero-coordinate punctuation occurs here and remains unchanged.
$$
Z = \mathbf{A}^m_S \longrightarrow \mathbf{A}^{m + n}_S = X,\quad
(a_1, \ldots, a_m) \mapsto (a_1, \ldots, a_m, 0, \ldots 0).
$$
を考える。閉部分スキーム $Z$ に沿う $X$ の吹き上げ $X'$ を考える。次のように書く。
$$
X = \Spec(A)
\quad\text{with}\quad
A = R[x_1, \ldots, x_m, y_1, \ldots, y_n]
$$
このとき $X'$ はイデアル $(y_1, \ldots, y_n)$ に関する $A$ の
Rees 代数の Proj である。この Rees 代数は
$B = A[T_1, \ldots, T_n]/(y_iT_j - y_jT_i)$ に等しい。詳細は省略する。
したがって $X' = \text{Proj}(B)$ は次のアフィン開集合で被覆されるので、
$S$ 上滑らかである。
% P05-EN-CH37-0085: frozen source omits the comma in "\ldots t_n" in the first chart presentation; the formula is preserved.
\begin{align*}
D_+(T_i)
& =
\Spec(B_{(T_i)}) \\
& =
\Spec(A[t_1, \ldots, \hat t_i, \ldots t_n]/(y_j - y_i t_j)) \\
& =
\Spec(R[x_1, \ldots, x_m, y_i, t_1, \ldots, \hat t_i, \ldots, t_n])
\end{align*}
これらは $\mathbf{A}^{n + m}_S$ と同型である。
このチャートでは例外因子は $y_i = 0$ によって切り出される。
したがって例外因子も $S$ 上滑らかである。
\end{remark}

\begin{lemma}
\label{more-morphisms-lemma-blowup}
$S$ をスキームとする。$Z \to X$ を、$S$ 上滑らかなスキームの
閉埋め込みとする。$b : X' \to X$ を $Z$ の吹き上げとし、その例外因子を
$E \subset X'$ とする。このとき $X'$ と $E$ は $S$ 上滑らかである。
射 $p : E \to Z$ は射影空間束
$$
\mathbf{P}(\mathcal{I}/\mathcal{I}^2) \longrightarrow Z
$$
と標準的に同型である。ここで $\mathcal{I} \subset \mathcal{O}_X$ は
$Z$ のイデアル層である。射影空間束の構造から得られる相対的な
$\mathcal{O}_E(1)$ は、$\mathcal{O}_{X'}(-E)$ の $E$ への制限と同型である。
\end{lemma}

\begin{proof}
Divisors, Lemma
\ref{divisors-lemma-immersion-smooth-into-smooth-regular-immersion}
により、$Z \to X$ は正則埋め込みである。
% P05-EN-CH37-0086: frozen source spells "immersion" with three m's; Japanese avoids the typo.
したがってイデアル層 $\mathcal{I}$ は有限型であり、ゆえに $b$ は射影的な射で、
相対的に豊富な可逆層
$\mathcal{O}_{X'}(1) = \mathcal{O}_{X'}(-E)$ をもつ。これについては
Divisors, Lemmas
\ref{divisors-lemma-blowing-up-gives-effective-Cartier-divisor} および
\ref{divisors-lemma-blowing-up-projective}
を参照されたい。標準写像 $\mathcal{I} \to b_*\mathcal{O}_{X'}(1)$ は、
吹き上げの構成そのものにより閉埋め込み
$$
X' \longrightarrow
\mathbf{P}\left(\bigoplus\nolimits_{n \geq 0}
\text{Sym}^n_{\mathcal{O}_X}(\mathcal{I})\right)
$$
を与える。この射を $E$ に制限すると、$Z$ 上の標準写像
$$
E \longrightarrow
\mathbf{P}\left(\bigoplus\nolimits_{n \geq 0}
\text{Sym}^n_{\mathcal{O}_Z}(\mathcal{I}/\mathcal{I}^2)\right)
$$
を得る。$\mathcal{I}/\mathcal{I}^2$ は有限局所自由であるから、
% P05-EN-CH37-0087: frozen source omits the clause-separating comma before "if"; Japanese supplies normal syntax.
この標準写像が同型ならば補題の最後の主張が成り立つ。
以上を踏まえると、問題は $X$ 上エタール局所的である。
実際、Divisors, Lemma \ref{divisors-lemma-flat-base-change-blowing-up}
により吹き上げは平坦基底変換と可換であり、全射なエタール射との前合成の後で
滑らかさを判定できる。したがって Lemma
\ref{more-morphisms-lemma-etale-local-structure} が与える $S$ 上のスキームの閉埋め込みの
エタール局所構造により、Remark \ref{more-morphisms-remark-relative-blowup}
で論じた場合に帰着する。
\end{proof}








\section{平坦加群と相対随伴点集合}
\label{more-morphisms-section-flat-relative-assassin}

\noindent
この節では、幾何学的な状況において、平坦加群の支持が基底上
（ある意味で）等次元であることを証明する。Noetherian の場合については
\cite[IV Proposition 12.1.1.5]{EGA}
を参照されたい。まず二つの補助補題を証明する。

\begin{lemma}
\label{more-morphisms-lemma-mod-injective-valuation-ring}
$A$ を付値環とする。
% P05-EN-CH37-0088: frozen source says "Let A -> B is"; Japanese supplies the grammatical predicate.
$A \to B$ を、本質的有限型な局所環の局所準同型とする。
$u : N \to M$ を有限 $B$-加群の写像とする。
$M$ は $A$ 上平坦であり、
$\overline{u} : N/\mathfrak m_A N \to M/\mathfrak m_A M$ は単射であると仮定する。
このとき $u$ は単射であり、$M/u(N)$ は $A$ 上平坦である。
\end{lemma}

\begin{proof}
この補題を Algebra, Lemma \ref{algebra-lemma-mod-injective-general}
から導く（$M$ と $N$ の役割を入れ替えていることに注意されたい）。
この帰着には More on Algebra, Lemma
\ref{more-algebra-lemma-valuation-ring-flat-essentially-finite-type}
を用いて有限表示の場合に帰着する。

\medskip\noindent
仮定により、$B$ は多項式環 $P = A[x_1, \ldots, x_n]$ を
素イデアル $\mathfrak q$ で局所化した環の商として書ける。
したがって $u : N \to M$ を有限 $P_\mathfrak q$-加群の写像とみなせる。
ゆえに $B$ は $A$ 上本質的有限表示であると仮定する。

\medskip\noindent
次に、$B$-加群 $N$ は有限だが有限表示とは限らない。
$N = \colim N_\lambda$ を、全射な遷移写像をもつ有限表示
$B$-加群のフィルター付き余極限として書く。例えば表示
$0 \to K \to B^{\oplus r} \to N \to 0$ を選び、$K$ を有限部分加群
$K_\lambda$ の合併として書き、$N_\lambda = \Coker(K_\lambda \to B^{\oplus r})$
とおけばよい。加群 $N/\mathfrak m_A N$ は Noetherian 環
$B/\mathfrak m_A B$ 上有限表示である。したがって十分大きい $\lambda$ に対して
$N_\lambda/\mathfrak m_A N_\lambda = N/\mathfrak m_A N$ となる。
ここで合成 $u_\lambda : N_\lambda \to M$ に対して補題を示せれば、
$N_\lambda = N$ が従い、$u$ に対して結果が成り立つ。
ゆえに $N$ は $B$ 上有限表示であると仮定してよい。

\medskip\noindent
More on Algebra, Lemma
\ref{more-algebra-lemma-valuation-ring-flat-essentially-finite-type}
により、加群 $M$ は $B$ 上有限表示である。したがって
Algebra, Lemma \ref{algebra-lemma-mod-injective-general}
のすべての仮定が成り立ち、結論を得る。
\end{proof}

\begin{lemma}
\label{more-morphisms-lemma-make-base-change}
\begin{reference}
これは次の命題の証明に見られる。
\cite[IV Proposition 12.1.1.5]{EGA}
\end{reference}
$f : X \to S$ をスキームの射とする。$y \in X$ を、その像が
$t \in S$ である点とする。
% P05-EN-CH37-0089: frozen source omits "by" in the naming construction for Y; Japanese supplies normal syntax.
$Y \subset X$ を、$\{y\}$ の閉包を $X$ の整な閉部分スキームとみなしたものとする。
$s \in S$ とし、$x \in Y_s$ を $Y_s$ の既約成分の一般点とする。
次の性質をもつデカルト図式
$$
\xymatrix{
X' \ar[r]_{g'} \ar[d]_{f'} & X \ar[d]^f \\
S' \ar[r]^g & S
}
$$
が存在する。
\begin{enumerate}
\item $S'$ は、一般点 $t'$ と閉点 $s'$ をもつ付値環のスペクトルである。
\item $g(t') = t$ かつ $g(s') = s$ である。
\item 点 $y' \in X'_{t'}$ が存在し、これは
$(S' \times_S Y)_{t'} = Y_t \times_t t'$ の既約成分の一般点であり、
$g'(y') = y$ を満たす。
\item $Y' \subset X'$ を、$\{y'\}$ の閉包を $X'$ の整な閉部分スキームと
みなしたものとすると、点 $x' \in Y'_{s'}$ が存在し、これは
$Y'_{s'}$ の既約成分の一般点であり、$g'(x') = x$ を満たす。
\end{enumerate}
\end{lemma}

\begin{proof}
付値環 $R$ を選び、$S' = \Spec(R)$ とおく。その一般点を $t'$、閉点を $s'$ とし、
$h(t') = y$ および $h(s') = x$ を満たす射 $h : S' \to X$ を選ぶ。
Schemes, Lemma \ref{schemes-lemma-points-specialize}
を参照されたい。$g = f \circ h$ とおけば、$g(t') = t$ かつ $g(s') = s$ である。
基底変換
$$
\xymatrix{
X' \ar[r]_{g'} \ar[d] & X \ar[d] \\
S' \ar@/^1em/[u]^\sigma \ar[r]^-g & S
}
$$
を考える。$h = g' \circ \sigma$ を満たすこの基底変換の切断 $\sigma$ を得る。

\medskip\noindent
もちろん $h$ は $Y$ を経由するので、$\sigma$ は $Y$ の基底変換
$S' \times_S Y$ を経由する。$y' \in X'_{t'} \subset X'$ を、点
$\sigma(t')$ を含むファイバー
$$
(S' \times_S Y)_{t'} = Y_t \times_t t'
$$
の既約成分の一般点、すなわち $y' \leadsto \sigma(t')$ を満たす点とする。
% P05-EN-CH37-0090: frozen source applies g:S'->S to points of X'; the ill-typed formulas are preserved.
$g'(y') \in Y_t$ かつ $g(y') \leadsto g(\sigma(t')) = y$ なので、
$y$ がファイバー $Y_t$ の一般点であることから $g'(y') = y$ を得る。
% P05-EN-CH37-0091: frozen source omits "by" in the naming construction for Y'; Japanese supplies normal syntax.
$Y' \subset X'$ を、$X'$ における $\{y'\}$ の閉包を整な閉部分スキームと
みなしたものとする。$\sigma(t') \in Y'$ なので、$\sigma$ は $Y'$ を経由する。
$Y'_{s'}$ のうち $\sigma(s')$ を含む既約成分の一般点
$x' \in Y'_{s'}$ を選ぶ。すなわち $x' \leadsto \sigma(s')$ であり、したがって
$g'(x') \leadsto g'(\sigma(s')) = x$ となる。仮定により $x$ は $Y_s$ の
既約成分の一般点であり、また $g'(Y') \subset Y$ なので、$g'(x') = x$ を得る。
\end{proof}

\begin{lemma}
\label{more-morphisms-lemma-associated-point-specializes}
$f : X \to S$ を局所有限型なスキームの射とする。
$\mathcal{F}$ を準連接な有限型 $\mathcal{O}_X$-加群とする。
$y \in \text{Ass}_{X/S}(\mathcal{F})$ とし、その像を $t \in S$ とする。
% P05-EN-CH37-0092: frozen source omits "by" in the naming construction for Y; Japanese supplies normal syntax.
$Y \subset X$ を、$X$ における $\{y\}$ の閉包を整な閉部分スキームと
みなしたものとする。$s \in S$ とし、$x \in Y_s$ を $Y_s$ の
既約成分の一般点とする。$\mathcal{F}$ が点 $x$ において $S$ 上平坦ならば、
$x \in \text{Ass}_{X/S}(\mathcal{F})$ であり、
$\dim_x(Y_s) = \dim(Y_t)$ である。
\end{lemma}

\begin{proof}
Lemma \ref{more-morphisms-lemma-make-base-change} のような図式を選ぶ。
$\mathcal{F}' = (g')^*\mathcal{F}$ とおく。
Divisors, Lemma \ref{divisors-lemma-base-change-relative-assassin}
により $y' \in \text{Ass}_{X'/S'}(\mathcal{F}')$ である。
$y'$ の選び方から $\dim(Y'_{t'}) = \dim(Y_t)$ でもある。例えば Algebra, Lemma
\ref{algebra-lemma-inequalities-under-field-extension}.
を参照されたい。Algebra, Lemma
\ref{algebra-lemma-finite-type-domain-over-valuation-ring-dim-fibres}
により、$Y'_{s'}$ は次元 $\dim(Y_t)$ の等次元スキームである。
$\mathcal{F}$ は点 $x$ において $S$ 上平坦なので、$\mathcal{F}'$ は点 $x'$ において
$S'$ 上平坦である。Morphisms, Lemma
\ref{morphisms-lemma-base-change-module-flat} を参照されたい。

\medskip\noindent
% P05-EN-CH37-0093: frozen source switches the relative-assassin base from S' to S; the formula is preserved.
$x' \in \text{Ass}_{X'/S}(\mathcal{F}')$ を示せると仮定する。このとき
Divisors, Lemma \ref{divisors-lemma-base-change-relative-assassin}
により $x \in \text{Ass}_{X/S}(\mathcal{F})$ である。また、$x'$ を含む
$Y'_{s'}$ の既約成分 $C'$ は、$x$ を含む既約成分
$C \subset Y_s$ に対する $C \times_s s'$ の既約成分である。
したがって $\dim(C) = \dim(C') = \dim(Y_t)$ であり（上記参照）、
証明が完了する。ゆえに次の段落で論じる場合に帰着する。

\medskip\noindent
$S = \Spec(A)$ と仮定する。ここで $A$ は付値環であり、$t$ と $s$ は
$S$ の一般点と閉点である。矛盾を導くため、
$x \not \in \text{Ass}_{X/S}(\mathcal{F})$ と仮定する。言い換えれば、
$x \not \in \text{Ass}_{X_s}(\mathcal{F}_s)$ と仮定する。ここで
$\mathcal{F}_s$ は $\mathcal{F}$ の $X_s$ への引き戻しである。
環準同型
$$
A \longrightarrow \mathcal{O}_{X, x} = B
$$
と、$B = \mathcal{O}_{X, x}$ 上の加群 $N = \mathcal{F}_x$ を考える。
このとき $B/\mathfrak m_A B = \mathcal{O}_{X_s, x}$ であり、
$N/\mathfrak m_A N$ は点 $x$ における $\mathcal{F}_s$ の茎である。
% P05-EN-CH37-0094: frozen source omits "by" in the naming construction for q; Japanese supplies normal syntax.
$\mathfrak q \subset B$ を点 $y$ に対応する素イデアルとする。Schemes, Lemma
\ref{schemes-lemma-specialize-points} を参照されたい。
$x$ は $Y_s$ の一般点なので、$\mathfrak q + \mathfrak m_A B$ の根基は
$\mathfrak m_B$ である。ここで
$\text{Ass}_{B/\mathfrak m_A B}(N/\mathfrak m_A N)$ は素イデアルの有限集合であり
（Algebra, Lemma \ref{algebra-lemma-finite-ass}）、
$x \not \in \text{Ass}_{X/S}(\mathcal{F})$ なので、$B/\mathfrak m_A B$ の
極大イデアルを含まない。
% P05-EN-CH37-0095: frozen source repeats "of" in "the image of of q"; Japanese avoids the duplication.
したがって $\mathfrak q$ の $B/\mathfrak m_A B$ における像は、これらの
どの素イデアルにも含まれない。ゆえに素イデアル回避
（Algebra, Lemma \ref{algebra-lemma-silly}）により、元 $g \in \mathfrak q$ で、その
$B/\mathfrak m_A B$ における像が $N/\mathfrak m_A N$ 上の非零因子となるものを
選べる（ここでは Algebra, Lemma \ref{algebra-lemma-ass-zero-divisors}
の零因子の記述を用いた）。
% P05-EN-CH37-0096: frozen source grammatically attributes A-flatness to a lemma that instead uses it; the citation is reattached without changing the claim.
$N = \mathcal{F}_x$ は $A$ 上平坦なので、
Lemma \ref{more-morphisms-lemma-mod-injective-valuation-ring} により
$$
g : N \longrightarrow N
$$
は単射である。特に、$K = \text{Frac}(A)$ を $A$ の分数体とすれば、
$$
g : N \otimes_A K \longrightarrow N \otimes_A K
$$
も単射である。$\mathfrak q$ は $B \otimes_A K$ の素イデアルに対応することに注意する。
% P05-EN-CH37-0097: frozen source omits "by" in the naming construction for F_t; Japanese supplies normal syntax.
$\mathcal{F}_t$ を $\mathcal{F}$ の一般ファイバー $X_t$ への制限とする。このとき
$(B \otimes_A K)_{\mathfrak q} = \mathcal{O}_{X_t, y}$
であり、$(N \otimes_A K)_\mathfrak q$ は点 $y$ における
$\mathcal{F}_t$ の茎である。したがって
$g \in \mathfrak m_y \subset \mathcal{O}_{X_t, y}$ は茎
$(\mathcal{F}_t)_y$ 上の非零因子である。これは
$y \in \text{Ass}_{X/S}(\mathcal{F})$ という仮定に矛盾する。
\end{proof}

\begin{lemma}
\label{more-morphisms-lemma-relative-dimension-support-flat}
$f : X \to S$ を局所有限型なスキームの射とする。
$\mathcal{F}$ を $S$ 上平坦な有限型準連接 $\mathcal{O}_X$-加群とする。
$S$ は一般点 $\eta$ をもつ既約スキームであると仮定する。
$\dim(\text{Supp}(\mathcal{F}_\eta)) \leq r$ ならば、任意の $s \in S$ に対して
$\dim(\text{Supp}(\mathcal{F}_s)) \leq r$ である。
\end{lemma}

\begin{proof}
$x \in \text{Supp}(\mathcal{F}_s)$ を $\text{Supp}(\mathcal{F}_s)$ の
既約成分の一般点とする。
Algebra, Lemma \ref{algebra-lemma-going-down-flat-module}
により、$f(y) = \eta$ を満たす $\text{Supp}(\mathcal{F})$ 内の特殊化
$y \leadsto x$ を選べる。もちろん $y$ は
$\text{Supp}(\mathcal{F}_\eta)$ の既約成分の一般点であると仮定してよい。
Lemma \ref{more-morphisms-lemma-associated-point-specializes}
から、$\overline{\{x\}}$ の次元は高々 $r$ であると結論する。
\end{proof}

\begin{lemma}
\label{more-morphisms-lemma-flat-associated-equidimensional}
$f : X \to S$ を局所有限型なスキームの射とする。
$\mathcal{F}$ を有限型準連接 $\mathcal{O}_X$-加群とする。
$y \in \text{Ass}_{X/S}(\mathcal{F})$ とする。
% P05-EN-CH37-0098: frozen source omits "by" in the two naming constructions for Y and T; Japanese supplies normal syntax.
$Y \subset X$ を、$X$ における $\{y\}$ の閉包を整な閉部分スキームと
みなしたものとする。$T \subset S$ を、$\{f(y)\}$ の閉包を整な閉部分スキームと
みなしたものとする。可換図式
$$
\xymatrix{
Y \ar[r] \ar[d] & X \ar[d] \\
T \ar[r] & S
}
$$
を得る。ここで $Y \to T$ は支配的である。$\mathcal{F}$ は、$Y \to T$ の
ファイバーの既約成分のすべての一般点において $S$ 上平坦であると仮定する
（例えば $\mathcal{F}$ が $S$ 上平坦ならば、この仮定は成り立つ）。このとき
\begin{enumerate}
\item $s \in S$ とし、$x \in Y_s$ が $Y_s$ の既約成分の一般点ならば、
$x \in \text{Ass}_{X/S}(\mathcal{F})$ である。
\item 整数 $d \geq 0$ が存在し、$Y \to T$ は相対次元 $d$ である。次を参照せよ。
Morphisms, Definition \ref{morphisms-definition-relative-dimension-d}。
\end{enumerate}
\end{lemma}

\begin{proof}
% P05-EN-CH37-0099: frozen source omits "of" before the cited lemma in "the pointwise version Lemma"; Japanese supplies normal syntax.
これは Lemma \ref{more-morphisms-lemma-associated-point-specializes}
の各点での版から直ちに従う。局所代数的スキーム $Y_s$ の次元を計算するには、
一般点の近傍だけを調べれば十分であることに注意する。Varieties, Section
\ref{varieties-section-algebraic-schemes} を参照されたい。
\end{proof}

\begin{remark}
\label{more-morphisms-remark-flat-equidimensional}
上の結果、とりわけ Lemma \ref{more-morphisms-lemma-flat-associated-equidimensional}
により、スキームの射 $f : X \to S$ が
Morphisms, Definition \ref{morphisms-definition-relative-dimension-d}
で定義される相対次元 $d$ をもつと結論できるいくつかの場合を挙げる。
例えば、次の条件のもとで成り立つ。
\begin{enumerate}
\item $X$ は一般点 $\xi$ をもつ整スキームである。
\item $\kappa(f(\xi))$ 上の $\kappa(\xi)$ の超越次数は $d$ である。
\item $f$ は局所有限型である。
\item 有限型準連接 $\mathcal{O}_X$-加群 $\mathcal{F}$ が存在し、これは
$S$ 上平坦で、$\text{Supp}(\mathcal{F}) = X$ を満たす。
\end{enumerate}
別の十分な仮定は次のとおりである。
\begin{enumerate}
\item $S$ は一般点 $\eta$ をもつ既約スキームである。
\item $X_\eta$ は $X$ で稠密である。
\item $X_\eta$ のすべての既約成分の次元は $d$ である。
\item $f$ は局所有限型である。
\item 有限型準連接 $\mathcal{O}_X$-加群 $\mathcal{F}$ が存在し、これは
$S$ 上平坦で、$\text{Supp}(\mathcal{F}) = X$ を満たす。
\end{enumerate}
もちろん、$\mathcal{F}$ に対する平坦性の条件を弱め、$\mathcal{F}$ が $S$ 上
余次元 $0$ で平坦、すなわち $\mathcal{F}$ がすべてのファイバーの各一般点において
$S$ 上平坦であることだけを要求してもよい。
% P05-EN-CH37-0100: frozen source exposes a future-work editorial note instead of a present mathematical statement; it remains visible in translation.
これらの結果が必要になったならば、ここで注意深く定式化して証明することにする。
\end{remark}











\section{正規化再考}
\label{more-morphisms-section-normalization}

\noindent
正規化は滑らかな基底変換と可換である。

\begin{lemma}
\label{more-morphisms-lemma-integral-closure-smooth-pullback}
$f : Y \to X$ をスキームの滑らかな射とする。$\mathcal{A}$ を準連接な
$\mathcal{O}_X$-代数の層とする。$f^*\mathcal{A}$ における
$\mathcal{O}_Y$ の整閉包は $f^*\mathcal{A}'$ に等しい。ここで
$\mathcal{A}' \subset \mathcal{A}$ は $\mathcal{A}$ における
$\mathcal{O}_X$ の整閉包である。
\end{lemma}

\begin{proof}
Algebra, Lemma \ref{algebra-lemma-integral-closure-commutes-smooth}
をスキームの言葉に翻訳したものである。詳細は省略する。
\end{proof}

\begin{lemma}[正規化は滑らかな基底変換と可換である]
\label{more-morphisms-lemma-normalization-smooth-localization}
次の図式
$$
\xymatrix{
Y_2 \ar[r] \ar[d]_{f_2} & Y_1 \ar[d]^{f_1} \\
X_2 \ar[r]^\varphi & X_1
}
$$
をスキームの圏におけるファイバー平方とする。$f_1$ は準コンパクトかつ
準分離であり、$\varphi$ は滑らかであると仮定する。
$Y_i \to X_i' \to X_i$ を $Y_i$ における $X_i$ の正規化とする。
このとき $X_2' \cong X_2 \times_{X_1} X_1'$ である。
\end{lemma}

\begin{proof}
$Y_1 \to X_1' \to X_1$ を $X_2$ へ基底変換すると、分解
$Y_2 \to X_2 \times_{X_1} X_1' \to X_2$ が得られ、
$X_2 \times_{X_1} X_1' \to X_2$ は整である
（Morphisms, Lemma \ref{morphisms-lemma-base-change-finite}）。
したがって Morphisms, Lemma \ref{morphisms-lemma-characterize-normalization}
の普遍性により射 $h : X_2' \to X_2 \times_{X_1} X_1'$ を得る。
$X_2'$ は $f_{2, *}\mathcal{O}_{Y_2}$ における $\mathcal{O}_{X_2}$ の
整閉包の相対スペクトルであることに注意する。
$\mathcal{A}' \subset f_{1, *}\mathcal{O}_{Y_1}$ が $\mathcal{O}_{X_1}$ の
整閉包を表すならば、$X_2 \times_{X_1} X_1'$ は
$\varphi^*\mathcal{A}'$ の相対スペクトルである。Constructions, Lemma
\ref{constructions-lemma-spec-properties} を参照されたい。
Cohomology of Schemes, Lemma \ref{coherent-lemma-flat-base-change-cohomology}
により $f_{2, *}\mathcal{O}_{Y_2} = \varphi^*f_{1, *}\mathcal{O}_{Y_1}$ である。
したがって結果は Lemma \ref{more-morphisms-lemma-integral-closure-smooth-pullback}
から従う。
\end{proof}

\begin{lemma}[正規化と滑らかな射]
\label{more-morphisms-lemma-normalization-and-smooth}
$X \to Y$ をスキームの滑らかな射とする。$Y$ の任意の準コンパクト開集合は
有限個の既約成分をもつと仮定する。このとき $X$ に対しても同じことが成り立ち、
$X$ 上の一意な同型 $X^\nu = X \times_Y Y^\nu$ が存在する。ここで
$X^\nu$, $Y^\nu$ はそれぞれ $X$, $Y$ の正規化である。
\end{lemma}

\begin{proof}
Descent, Lemma \ref{descent-lemma-locally-finite-nr-irred-local-fppf}
により、$X$ の任意の準コンパクト開集合は有限個の既約成分をもつ。
被約スキーム上滑らかなスキームは被約なので、
$X_{red} = X \times_Y Y_{red}$ であることに注意する。次を参照せよ。
Descent, Lemma \ref{descent-lemma-reduced-local-smooth}.
したがって $X$ と $Y$ は被約であると仮定してよい（定義により、スキームの
正規化はその被約化の正規化に等しい）。次に Descent, Lemma
\ref{descent-lemma-normal-local-smooth} により、
$X' = X \times_Y Y^\nu$ は正規スキームである。
射 $X' \to Y^\nu$ は滑らか（したがって平坦）なので、$X'$ の既約成分の一般点は
$Y^\nu$ の既約成分の一般点の上にある。$Y^\nu \to Y$ は双有理なので、
$X' \to X$ も双有理であると結論する（$X' \to Y^\nu$ は $Y$ の一般点上の
ファイバーに同型を誘導するためである）。したがって分解
$X^\nu \to X' \to X$ が存在する。次を参照せよ。
Morphisms, Lemma \ref{morphisms-lemma-normalization-normal}
$X'$ は正規であり $X$ 上整なので、この分解は同型である。
\end{proof}

\begin{lemma}[正規化と Hensel 化]
\label{more-morphisms-lemma-normalization-and-henselization}
$X$ を局所 Noetherian スキームとする。$\nu : X^\nu \to X$ を正規化射とする。
このとき任意の点 $x \in X$ に対して、基底変換
$$
X^\nu \times_X \Spec(\mathcal{O}_{X, x}^h) \to \Spec(\mathcal{O}_{X, x}^h),
\quad\text{resp.}\quad
X^\nu \times_X \Spec(\mathcal{O}_{X, x}^{sh}) \to \Spec(\mathcal{O}_{X, x}^{sh})
$$
はそれぞれ $\Spec(\mathcal{O}_{X, x}^h)$ および
$\Spec(\mathcal{O}_{X, x}^{sh})$ の正規化である。
\end{lemma}

\begin{proof}
$\eta_1, \ldots, \eta_r$ を、$x$ を通る $X$ の既約成分の一般点とする。
正規化を $\Spec(\mathcal{O}_{X, x})$ へ基底変換したものは、
$\prod \kappa(\eta_i)$ における $\mathcal{O}_{X, x}$ の整閉包のスペクトルである。
これは $X$ の正規化の構成と
Morphisms, Definition \ref{morphisms-definition-normalization}
および Morphisms, Lemma \ref{morphisms-lemma-integral-closure}
から従う。Morphisms, Lemma \ref{morphisms-lemma-description-normalization}
の正規化の記述を用いてもよい。したがって次の代数の問題に帰着する。
$A$ を Noetherian 局所環とする。このとき Hensel 化 $A^h$ と狭義 Hensel 化
$A^{sh}$ も Noetherian であることを思い出そう（More on Algebra, Lemma
\ref{more-algebra-lemma-henselization-noetherian}）。
$\mathfrak p_1, \ldots, \mathfrak p_r$ をその極小素イデアルとする。
$A'$ を $\prod \kappa(\mathfrak p_i)$ における $A$ の整閉包とする。
問題は、$A' \otimes_A A^h$ および $A' \otimes_A A^{sh}$ がそれぞれ
Noetherian 局所環 $A^h$ および $A^{sh}$ から同じ仕方で構成されることを
示すことである。

\medskip\noindent
$A^h$ および $A^{sh}$ はエタール $A$-代数の余極限なので、
% P05-EN-CH37-0101: frozen source names the minimal primes of A instead of A^h in the first paired clause; the formula is preserved.
$A$ および $A^{sh}$ の極小素イデアルはちょうど、$A$ の極小素イデアルの上にある
$A^h$ および $A^{sh}$ の素イデアルであることが分かる（下降による。Algebra, Lemmas
\ref{algebra-lemma-flat-going-down} および
\ref{algebra-lemma-minimal-prime-image-minimal-prime} を参照せよ）。
したがって More on Algebra, Lemma \ref{more-algebra-lemma-fibres-henselization}
により、
$A^h \otimes_A \prod \kappa(\mathfrak p_i)$,
および
$A^{sh} \otimes_A \prod \kappa(\mathfrak p_i)$
はそれぞれ $A^h$ および $A^{sh}$ の極小素イデアルにおける剰余体の積である。
上環における整閉包を取る操作はエタール基底変換と可換である。Algebra, Lemma
\ref{algebra-lemma-integral-closure-commutes-etale} を参照されたい。
% P05-EN-CH37-0102: frozen source calls the henselizations a limit of etale algebras after correctly calling them colimits; the prose is preserved.
$A^h$ と $A^{sh}$ をエタール $A$-代数の極限として書けば、$A^h$ および
$A^{sh}$ への基底変換についても同じことが成り立つと分かる
（より一般的な Algebra, Lemma
\ref{algebra-lemma-integral-closure-commutes-colim-smooth} を用いてもよい）。
\end{proof}








\section{正規射}
\label{more-morphisms-section-normal}

\noindent
Deligne と Mumford の論文 \cite{DM} では正規射という概念が言及されている。
これは同様に定義できる一連の射の型の一つにすぎない\footnote{
他の型は coprof $\leq k$、Cohen--Macaulay、$(S_k)$、正則、$(R_k)$、
および被約である。\cite[IV Definition 6.8.1.]{EGA} を参照せよ。
Gorenstein 射は Duality for Schemes, Section
\ref{duality-section-gorenstein} で定義する。}。
% P05-EN-CH37-0103: frozen source exposes a rolling future-work note; it remains visible in translation.
必要に応じて、これらを今後それぞれ独立の節として追加することにする。

\begin{definition}
\label{more-morphisms-definition-normal}
$f : X \to Y$ をスキームの射とする。すべてのファイバー $X_y$ は
局所 Noetherian スキームであると仮定する。
\begin{enumerate}
\item $x \in X$ とし、$y = f(x)$ とおく。$f$ が点 $x$ で{\it 正規}であるとは、
$f$ が点 $x$ で平坦であり、スキーム $X_y$ が $\kappa(y)$ 上、点 $x$ で
幾何学的正規であることをいう（Varieties, Definition
\ref{varieties-definition-geometrically-normal} 参照）。
\item $f$ が $X$ のすべての点で正規であるとき、$f$ を{\it 正規射}という。
\end{enumerate}
\end{definition}

\noindent
したがって射 $X \to Y$ が正規であるという条件は、すべてのファイバーが
正規な局所 Noetherian スキームであることだけを要求する条件より強い。

\begin{lemma}
\label{more-morphisms-lemma-normal}
$f : X \to Y$ をスキームの射とする。$f$ のすべてのファイバーは
局所 Noetherian であると仮定する。次は同値である。
\begin{enumerate}
\item $f$ は正規である。
\item $f$ は平坦であり、そのファイバーは幾何学的正規スキームである。
\end{enumerate}
\end{lemma}

\begin{proof}
定義から直ちに従う。
\end{proof}

\begin{lemma}
\label{more-morphisms-lemma-smooth-normal}
滑らかな射は正規である。
\end{lemma}

\begin{proof}
$f : X \to Y$ を滑らかな射とする。$f$ は局所有限表示なので（
Morphisms, Lemma \ref{morphisms-lemma-smooth-locally-finite-presentation}
参照）、ファイバー $X_y$ は体上局所有限型であり、したがって局所 Noetherian である。
さらに、Morphisms, Lemma \ref{morphisms-lemma-smooth-flat}
により $f$ は平坦である。最後に、ファイバー $X_y$ は体上滑らかであり
（Morphisms, Lemma \ref{morphisms-lemma-base-change-smooth}）、したがって
Varieties, Lemma \ref{varieties-lemma-smooth-geometrically-normal}
により幾何学的正規である。ゆえに Lemma \ref{more-morphisms-lemma-normal}
により $f$ は正規である。
\end{proof}

\noindent
この概念が滑らかな位相について始域と終域上局所的であることを示したい。
まずファイバーが局所 Noetherian であるという性質を扱う。

\begin{lemma}
\label{more-morphisms-lemma-locally-Noetherian-fibres-fppf-local-source-and-target}
性質 $\mathcal{P}(f)=$「$f$ のファイバーは局所 Noetherian である」は、
始域と終域上の fppf 位相について局所的である。
\end{lemma}

\begin{proof}
$f : X \to Y$ をスキームの射とする。
$\{\varphi_i : Y_i \to Y\}_{i \in I}$ を $Y$ の fppf 被覆とする。
% P05-EN-CH37-0104: frozen source omits "by" in the naming construction for f_i; Japanese supplies normal syntax.
$f_i : X_i \to Y_i$ を $\varphi_i$ による $f$ の基底変換とする。
$i \in I$ とし、$y_i \in Y_i$ を点とする。$y = \varphi_i(y_i)$ とおく。このとき
$$
X_{i, y_i} = \Spec(\kappa(y_i)) \times_{\Spec(\kappa(y))} X_y.
$$
さらに $\varphi_i$ は有限表示なので、体拡大 $\kappa(y_i)/\kappa(y)$ は有限生成である。
したがってこの状況では、$X_y$ が局所 Noetherian であることと
$X_{i, y_i}$ が局所 Noetherian であることは同値である。Varieties, Lemma
\ref{varieties-lemma-locally-Noetherian-base-change} を参照されたい。
これにより終域上の局所性が従う。

\medskip\noindent
$\{X_i \to X\}$ を $X$ の fppf 被覆とし、$y \in Y$ とする。この場合
$\{X_{i, y} \to X_y\}$ はファイバーの fppf 被覆である。したがって始域上の局所性は
Descent, Lemma \ref{descent-lemma-Noetherian-local-fppf}
から従う。
\end{proof}

\begin{lemma}
\label{more-morphisms-lemma-normal-fppf-local-source-and-target}
性質
$\mathcal{P}(f)=$「$f$ のファイバーは局所 Noetherian であり、$f$ は正規である」は、
終域上の fppf 位相について局所的であり、始域上の滑らかな位相について局所的である。
\end{lemma}

\begin{proof}
次の分解がある。
$\mathcal{P}(f) =
\mathcal{P}_1(f) \wedge \mathcal{P}_2(f) \wedge \mathcal{P}_3(f)$
ここで、$\mathcal{P}_1(f)=$「$f$ のファイバーは局所 Noetherian である」、
$\mathcal{P}_2(f)=$「$f$ は平坦である」、および
$\mathcal{P}_3(f)=$「$f$ のファイバーは幾何学的正規である」とする。
$\mathcal{P}_1$ と $\mathcal{P}_2$ が始域と終域上の fppf 位相について
局所的であることはすでに見た。次を参照せよ。
Lemma \ref{more-morphisms-lemma-locally-Noetherian-fibres-fppf-local-source-and-target},
および Descent, Lemmas \ref{descent-lemma-descending-property-flat} および
\ref{descent-lemma-flat-fpqc-local-source}。したがって $\mathcal{P}_3$ を扱えばよい。

\medskip\noindent
$f : X \to Y$ をスキームの射とする。
$\{\varphi_i : Y_i \to Y\}_{i \in I}$ を $Y$ の fpqc 被覆とする。
% P05-EN-CH37-0105: frozen source omits "by" in the naming construction for f_i; Japanese supplies normal syntax.
$f_i : X_i \to Y_i$ を $\varphi_i$ による $f$ の基底変換とする。
$i \in I$ とし、$y_i \in Y_i$ を点とする。$y = \varphi_i(y_i)$ とおく。このとき
$$
X_{i, y_i} = \Spec(\kappa(y_i)) \times_{\Spec(\kappa(y))} X_y.
$$
したがってこの状況では、$X_y$ が幾何学的正規であることと
$X_{i, y_i}$ が幾何学的正規であることは同値である。Varieties, Lemma
\ref{varieties-lemma-geometrically-normal-upstairs} を参照されたい。
これにより $\mathcal{P}_3$ は終域上 fpqc 局所的である。

\medskip\noindent
$\{X_i \to X\}$ を $X$ の滑らかな被覆とし、$y \in Y$ とする。この場合
$\{X_{i, y} \to X_y\}$ はファイバーの滑らかな被覆である。したがって
始域上の滑らかな位相に関する $\mathcal{P}_3$ の局所性は
Descent, Lemma \ref{descent-lemma-normal-local-smooth}
から従う。以上を合わせて補題を得る。
\end{proof}



\section{正則射}
\label{more-morphisms-section-regular}

\noindent
Section \ref{more-morphisms-section-normal} と比較せよ。この概念の代数的な版は
More on Algebra, Section \ref{more-algebra-section-regular} で論じる。

\begin{definition}
\label{more-morphisms-definition-regular}
$f : X \to Y$ をスキームの射とする。すべてのファイバー $X_y$ は
局所 Noetherian スキームであると仮定する。
\begin{enumerate}
\item $x \in X$ とし、$y = f(x)$ とおく。$f$ が点 $x$ で平坦であり、
スキーム $X_y$ が $\kappa(y)$ 上、点 $x$ で幾何学的正則であるとき、
$f$ は{\it 点 $x$ で正則}であるという（Varieties, Definition
\ref{varieties-definition-geometrically-regular} 参照）。
\item $f$ が $X$ のすべての点で正則であるとき、$f$ を{\it 正則射}という。
\end{enumerate}
\end{definition}

\noindent
射 $X \to Y$ が正則であるという条件は、すべてのファイバーが
正則な局所 Noetherian スキームであることだけを要求する条件より強い。

\begin{lemma}
\label{more-morphisms-lemma-regular}
$f : X \to Y$ をスキームの射とする。$f$ のすべてのファイバーは
局所 Noetherian であると仮定する。次は同値である。
\begin{enumerate}
\item $f$ は正則である。
\item $f$ は平坦であり、そのファイバーは幾何学的正則スキームである。
\item $f(U) \subset V$ を満たす任意のアフィン開 $U \subset X$、$V \subset Y$
に対し、環準同型 $\mathcal{O}_Y(V) \to \mathcal{O}_X(U)$ は正則である。
\item 開被覆 $Y = \bigcup_{j \in J} V_j$ と開被覆
$f^{-1}(V_j) = \bigcup_{i \in I_j} U_i$ で、各射 $U_i \to V_j$ が
正則となるものが存在する。
\item アフィン開被覆 $Y = \bigcup_{j \in J} V_j$ とアフィン開被覆
$f^{-1}(V_j) = \bigcup_{i \in I_j} U_i$ で、環準同型
$\mathcal{O}_Y(V_j) \to \mathcal{O}_X(U_i)$ が正則となるものが存在する。
\end{enumerate}
\end{lemma}

\begin{proof}
(1) と (2) の同値性は定義から直ちに従う。
$x \in X$ とし、$y = f(x)$ とおく。定義により、$f$ が点 $x$ で
平坦であることと $\mathcal{O}_{Y, y} \to \mathcal{O}_{X, x}$ が
平坦な環準同型であることは同値である。また、$X_y$ が $\kappa(y)$ 上、
点 $x$ で幾何学的正則であることと
$\mathcal{O}_{X_y, x} = \mathcal{O}_{X, x}/\mathfrak m_y\mathcal{O}_{X, x}$
が $\kappa(y)$ 上幾何学的正則な代数であることも同値である。
% P05-EN-CH37-0106: frozen source capitalizes "Whether" mid-sentence; Japanese uses normal sentence syntax.
したがって $f$ が点 $x$ で正則であるかどうかは、局所環の局所準同型
$\mathcal{O}_{Y, y} \to \mathcal{O}_{X, x}$ のみに依存する。
ゆえに (1) と (4) の同値性は明らかである。

\medskip\noindent
環準同型 $A \to B$ が正則であることと、それが平坦で、かつすべての素イデアル
$\mathfrak p \subset A$ に対してファイバー環
$B \otimes_A \kappa(\mathfrak p)$ が Noetherian かつ幾何学的正則であることは
同値であることを思い出そう（More on Algebra, Definition
\ref{more-algebra-definition-regular}）。Varieties, Lemma
\ref{varieties-lemma-geometrically-regular} により、これは
$\Spec(B \otimes_A \kappa(\mathfrak p))$ が $\kappa(\mathfrak p)$ 上
幾何学的正則なスキームであることと同値である。したがって (2) は
(3) を含意する。また (3) が (5) を含意することは明らかである。
最後に (5) を仮定する。このとき $f$ は平坦である（Morphisms, Lemma
\ref{morphisms-lemma-flat-characterize} 参照）。さらに $y \in Y$ とすれば、
ある $j$ に対して $y \in V_j$ であり、
$X_y = \bigcup_{i \in I_j} U_{i, y}$ となる。各 $U_{i, y}$ は
Varieties, Lemma \ref{varieties-lemma-geometrically-regular} により
$\kappa(y)$ 上幾何学的正則である。Varieties, Lemma
\ref{varieties-lemma-geometrically-regular} をもう一度適用すれば
$X_y$ は幾何学的正則である。よって (2) が成り立ち、証明が完了する。
\end{proof}

\begin{lemma}
\label{more-morphisms-lemma-smooth-regular}
滑らかな射は正則である。
\end{lemma}

\begin{proof}
$f : X \to Y$ を滑らかな射とする。$f$ は局所有限表示なので（
Morphisms, Lemma \ref{morphisms-lemma-smooth-locally-finite-presentation}
参照）、ファイバー $X_y$ は体上局所有限型であり、したがって局所 Noetherian である。
さらに、Morphisms, Lemma \ref{morphisms-lemma-smooth-flat}
により $f$ は平坦である。最後に、ファイバー $X_y$ は体上滑らかであり
（Morphisms, Lemma \ref{morphisms-lemma-base-change-smooth}）、したがって
Varieties, Lemma \ref{varieties-lemma-smooth-geometrically-normal}
により幾何学的正則である。ゆえに Lemma \ref{more-morphisms-lemma-regular}
により $f$ は正則である。
\end{proof}

\begin{lemma}
\label{more-morphisms-lemma-regular-fppf-local-source-and-target}
性質
$\mathcal{P}(f)=$「$f$ のファイバーは局所 Noetherian であり、$f$ は正則である」は、
終域上の fppf 位相について局所的であり、始域上の滑らかな位相について局所的である。
\end{lemma}

\begin{proof}
次の分解がある。
$\mathcal{P}(f) =
\mathcal{P}_1(f) \wedge \mathcal{P}_2(f) \wedge \mathcal{P}_3(f)$
ここで、$\mathcal{P}_1(f)=$「$f$ のファイバーは局所 Noetherian である」、
$\mathcal{P}_2(f)=$「$f$ は平坦である」、および
$\mathcal{P}_3(f)=$「$f$ のファイバーは幾何学的正則である」とする。
$\mathcal{P}_1$ と $\mathcal{P}_2$ が始域と終域上の fppf 位相について
局所的であることはすでに見た。次を参照せよ。
Lemma \ref{more-morphisms-lemma-locally-Noetherian-fibres-fppf-local-source-and-target},
および Descent, Lemmas \ref{descent-lemma-descending-property-flat} および
\ref{descent-lemma-flat-fpqc-local-source}。したがって $\mathcal{P}_3$ を扱えばよい。

\medskip\noindent
$f : X \to Y$ をスキームの射とする。
$\{\varphi_i : Y_i \to Y\}_{i \in I}$ を $Y$ の fpqc 被覆とする。
% P05-EN-CH37-0107: frozen source omits "by" in the naming construction for f_i; Japanese supplies normal syntax.
$f_i : X_i \to Y_i$ を $\varphi_i$ による $f$ の基底変換とする。
$i \in I$ とし、$y_i \in Y_i$ を点とする。$y = \varphi_i(y_i)$ とおく。このとき
$$
X_{i, y_i} = \Spec(\kappa(y_i)) \times_{\Spec(\kappa(y))} X_y.
$$
したがってこの状況では、$X_y$ が幾何学的正則であることと
$X_{i, y_i}$ が幾何学的正則であることは同値である。Varieties, Lemma
\ref{varieties-lemma-geometrically-regular-upstairs} を参照されたい。
これにより $\mathcal{P}_3$ は終域上 fpqc 局所的である。

\medskip\noindent
$\{X_i \to X\}$ を $X$ の滑らかな被覆とし、$y \in Y$ とする。この場合
$\{X_{i, y} \to X_y\}$ はファイバーの滑らかな被覆である。したがって
始域上の滑らかな位相に関する $\mathcal{P}_3$ の局所性は
Descent, Lemma \ref{descent-lemma-regular-local-smooth}
から従う。以上を合わせて補題を得る。
\end{proof}




\section{Cohen--Macaulay 射}
\label{more-morphisms-section-CM}

\noindent
Section \ref{more-morphisms-section-normal} と比較せよ。Algebra, Section
\ref{algebra-section-geometrically-CM} および Varieties, Section
\ref{varieties-section-CM} で指摘したように、「幾何学的
Cohen--Macaulay」は単に Cohen--Macaulay であることと同じである。

\begin{definition}
\label{more-morphisms-definition-CM}
$f : X \to Y$ をスキームの射とする。すべてのファイバー $X_y$ は
局所 Noetherian スキームであると仮定する。
\begin{enumerate}
\item $x \in X$ とし、$y = f(x)$ とおく。$f$ が点 $x$ で平坦であり、
スキーム $X_y$ の点 $x$ における局所環が Cohen--Macaulay であるとき、
$f$ は{\it 点 $x$ で Cohen--Macaulay} であるという。
\item $f$ が $X$ のすべての点で Cohen--Macaulay であるとき、$f$ を
{\it Cohen--Macaulay 射}という。
\end{enumerate}
\end{definition}

\noindent
この定義を言い換えると次のようになる。

\begin{lemma}
\label{more-morphisms-lemma-CM}
$f : X \to Y$ をスキームの射とする。$f$ のすべてのファイバーは
局所 Noetherian であると仮定する。次は同値である。
\begin{enumerate}
\item $f$ は Cohen--Macaulay である。
\item $f$ は平坦であり、そのファイバーは Cohen--Macaulay スキームである。
\end{enumerate}
\end{lemma}

\begin{proof}
定義から直ちに従う。
\end{proof}

\begin{lemma}
\label{more-morphisms-lemma-CM-dimension}
$f : X \to Y$ を局所 Noetherian スキームの局所有限型かつ
Cohen--Macaulay な射とする。$X$ の任意の点 $x$ と、その $Y$ における像 $y$ に対し、
$$
\dim_x(X) = \dim_y(Y) + \dim_x(X_y),
$$
が成り立つ。ここで $X_y$ は $y$ 上のファイバーを表す。
\end{lemma}

\begin{proof}
$X$ を $x$ の開近傍で置き換えると、$X \to Y$ のすべてのファイバーが
各点で次元 $d$ をもつような自然数 $d$ が存在する。Morphisms, Lemma
\ref{morphisms-lemma-flat-finite-presentation-CM-fibres-relative-dimension}
を参照せよ。このとき $f$ は平坦、局所有限型、かつ相対次元 $d$ である。
したがって Morphisms, Lemma
\ref{morphisms-lemma-rel-dimension-dimension} から結論が従う。
\end{proof}

\begin{lemma}
\label{more-morphisms-lemma-composition-CM}
$f : X \to Y$ および $g : Y \to Z$ をスキームの射とする。$f$、$g$、および
$g \circ f$ のファイバーは局所 Noetherian であると仮定する。
$x \in X$ とし、その像をそれぞれ $y \in Y$、$z \in Z$ とする。
\begin{enumerate}
\item $f$ が点 $x$ で Cohen--Macaulay であり、$g$ が点 $f(x)$ で
Cohen--Macaulay ならば、$g \circ f$ は点 $x$ で Cohen--Macaulay である。
\item $f$ と $g$ が Cohen--Macaulay ならば、$g \circ f$ は Cohen--Macaulay である。
\item $g \circ f$ が点 $x$ で Cohen--Macaulay であり、$f$ が点 $x$ で
平坦ならば、$f$ は点 $x$ で Cohen--Macaulay であり、$g$ は点 $f(x)$ で
Cohen--Macaulay である。
\item $g \circ f$ が Cohen--Macaulay であり、$f$ が平坦ならば、$f$ は
Cohen--Macaulay であり、$g$ は $f$ の像に属するすべての点で
Cohen--Macaulay である。
\end{enumerate}
\end{lemma}

\begin{proof}
Noetherian 局所環の準同型
$$
\mathcal{O}_{Y_z, y} \to \mathcal{O}_{X_z, x}
$$
を考える。そのファイバーは
$$
\mathcal{O}_{X_z, x}/\mathfrak m_{Y_z, y}\mathcal{O}_{X_z, x} =
\mathcal{O}_{X_y, x}
$$
であることに注意せよ。
% P05-EN-CH37-0108: frozen source has the malformed phrase "Thus the lemma this follows from"; Japanese uses normal syntax.
したがって補題は Algebra, Lemma \ref{algebra-lemma-CM-goes-up} から従う。
\end{proof}

\begin{lemma}
\label{more-morphisms-lemma-flat-morphism-from-CM-scheme}
$f : X \to Y$ を局所 Noetherian スキームの平坦射とする。$X$ が
Cohen--Macaulay ならば、$f$ は Cohen--Macaulay であり、すべての $x \in X$
に対して $\mathcal{O}_{Y, f(x)}$ は Cohen--Macaulay である。
\end{lemma}

\begin{proof}
代数の言葉に移せば、Algebra, Lemma \ref{algebra-lemma-CM-goes-up}
から従う。
\end{proof}

\begin{lemma}
\label{more-morphisms-lemma-base-change-CM}
$f : X \to Y$ をスキームの射とする。すべてのファイバー $X_y$ は
局所 Noetherian スキームであると仮定する。$Y' \to Y$ は局所有限型とし、
$f' : X' \to Y'$ を $f$ の基底変換とする。
$x' \in X'$ を点とし、その像を $x \in X$ とする。
\begin{enumerate}
\item $f$ が点 $x$ で Cohen--Macaulay ならば、$f' : X' \to Y'$ は点 $x'$ で
Cohen--Macaulay である。
\item $f$ が点 $x$ で平坦であり、$f'$ が点 $x'$ で Cohen--Macaulay ならば、
$f$ は点 $x$ で Cohen--Macaulay である。
\item $Y' \to Y$ が点 $f'(x')$ で平坦であり、$f'$ が点 $x'$ で
Cohen--Macaulay ならば、$f$ は点 $x$ で Cohen--Macaulay である。
\end{enumerate}
\end{lemma}

\begin{proof}
$Y' \to Y$ に関する仮定から、$y' \in Y'$ が $y \in Y$ に写るとき、
体拡大 $\kappa(y')/\kappa(y)$ は有限生成であることに注意せよ。したがって
すべてのファイバー $X'_{y'} = (X_y)_{\kappa(y')}$ も局所 Noetherian である。
Varieties, Lemma \ref{varieties-lemma-locally-Noetherian-base-change}
を参照せよ。ゆえに補題の主張は意味をもつ。$y' = f'(x')$ および
$y = f(x)$ とおく。このとき局所環の可換図式
$$
\xymatrix{
\mathcal{O}_{X', x'} & \mathcal{O}_{X, x} \ar[l] \\
\mathcal{O}_{Y', y'} \ar[u] & \mathcal{O}_{Y, y} \ar[l] \ar[u]
}
$$
を得る。ここで左上の環は、右上と左下の環を右下の環上でテンソル積したものの
局所化である。

\medskip\noindent
$f$ が点 $x$ で Cohen--Macaulay であると仮定する。
$\mathcal{O}_{Y, y} \to \mathcal{O}_{X, x}$ の平坦性から
$\mathcal{O}_{Y', y'} \to \mathcal{O}_{X', x'}$ の平坦性が従う。Algebra,
Lemma \ref{algebra-lemma-base-change-flat-up-down} を参照せよ。
$\mathcal{O}_{X, x}/\mathfrak m_y\mathcal{O}_{X, x}$ が
Cohen--Macaulay であることから
$\mathcal{O}_{X', x'}/\mathfrak m_{y'}\mathcal{O}_{X', x'}$
も Cohen--Macaulay であることが従う。Varieties, Lemma
\ref{varieties-lemma-CM-base-change} を参照せよ。したがって $f'$ は
点 $x'$ で Cohen--Macaulay である。

\medskip\noindent
$f$ が点 $x$ で平坦であり、$f'$ が点 $x'$ で Cohen--Macaulay であると仮定する。
$\mathcal{O}_{X', x'}/\mathfrak m_{y'}\mathcal{O}_{X', x'}$ が
Cohen--Macaulay であることから
$\mathcal{O}_{X, x}/\mathfrak m_y\mathcal{O}_{X, x}$
も Cohen--Macaulay であることが従う。Varieties, Lemma
\ref{varieties-lemma-CM-base-change} を参照せよ。したがって $f$ は
点 $x$ で Cohen--Macaulay である。

\medskip\noindent
$Y' \to Y$ が点 $y'$ で平坦であり、$f'$ が点 $x'$ で Cohen--Macaulay
であると仮定する。$\mathcal{O}_{Y', y'} \to \mathcal{O}_{X', x'}$ と
$\mathcal{O}_{Y, y} \to \mathcal{O}_{Y', y'}$ の平坦性から
$\mathcal{O}_{Y, y} \to \mathcal{O}_{X, x}$ の平坦性が従う。Algebra,
Lemma \ref{algebra-lemma-base-change-flat-up-down} を参照せよ。
$\mathcal{O}_{X', x'}/\mathfrak m_{y'}\mathcal{O}_{X', x'}$ が
Cohen--Macaulay であることから
$\mathcal{O}_{X, x}/\mathfrak m_y\mathcal{O}_{X, x}$
も Cohen--Macaulay であることが従う。Varieties, Lemma
\ref{varieties-lemma-CM-base-change} を参照せよ。したがって $f$ は
点 $x$ で Cohen--Macaulay である。
\end{proof}

\begin{lemma}
\label{more-morphisms-lemma-flat-finite-presentation-CM-open}
\begin{reference}
\cite[IV Corollary 12.1.7(iii)]{EGA}
\end{reference}
$f : X \to S$ を平坦かつ局所有限表示なスキームの射とする。次をおく。
$$
W = \{x \in X \mid f\text{ は点 }x\text{ で Cohen--Macaulay である}\}
$$
このとき次が成り立つ。
\begin{enumerate}
\item $W = \{x \in X \mid \mathcal{O}_{X_{f(x)}, x}\text{ は Cohen--Macaulay である}\}$。
\item $W$ は $X$ の開部分集合である。
\item $W$ は $X \to S$ の各ファイバーで稠密である。
\item $W$ の形成は $f$ の任意の基底変換と可換である。任意の射
$g : S' \to S$ に対し、$f$ の基底変換 $f' : X' \to S'$ と射影
$g' : X' \to X$ を考える。
% P05-EN-CH37-0109: frozen source redundantly says the corresponding set W' is equal to W' = ...; Japanese states the identity once.
このとき射 $f'$ に対応する集合 $W'$ は $W' = (g')^{-1}(W)$ に等しい。
\end{enumerate}
\end{lemma}

\begin{proof}
$f$ は平坦で、そのファイバーは局所 Noetherian なので、(1) の等式は
定義から成り立つ。(2) と (3) は Algebra, Lemma
\ref{algebra-lemma-generic-CM-flat-finite-presentation} から従う。
(4) は Algebra, Lemma \ref{algebra-lemma-CM-locus-commutes-base-change}
または Varieties, Lemma \ref{varieties-lemma-CM-base-change} から従う。
\end{proof}

\begin{lemma}
\label{more-morphisms-lemma-flat-finite-presentation-characterize-CM}
$f : X \to S$ を平坦かつ局所有限表示なスキームの射とする。
$x \in X$ とし、その像を $s \in S$ とする。$d = \dim_x(X_s)$ とおく。
次は同値である。
\begin{enumerate}
\item $f$ は点 $x$ で Cohen--Macaulay である。
\item $x$ の開近傍 $U \subset X$ と、点 $x$ で平坦な $S$ 上の
局所準有限射 $U \to \mathbf{A}^d_S$ が存在する。
\item $x$ の開近傍 $U \subset X$ と、$S$ 上の局所準有限平坦射
$U \to \mathbf{A}^d_S$ が存在する。
\item $x$ の任意の開近傍 $U \subset X$ と任意の $S$-射
$g : U \to \mathbf{A}^d_S$ に対し、次が成り立つ：
$g$ が点 $x$ で準有限 $\Rightarrow$ $g$ が点 $x$ で平坦。
\end{enumerate}
\end{lemma}

\begin{proof}
平坦性の開性から (2) と (3) は同値である。Theorem
\ref{more-morphisms-theorem-openness-flatness} を参照せよ。

\medskip\noindent
$x \in U$ を満たすアフィン開 $U = \Spec(A) \subset X$ と、
$f(U) \subset V$ を満たすアフィン開 $V = \Spec(R) \subset S$ を選ぶ。
このとき $R \to A$ は有限表示な平坦環準同型である。$x$ に対応する素イデアルを
$\mathfrak p \subset A$ とする。$A$ を主局所化で置き換えることにより、準有限準同型
$R[x_1, \ldots, x_d]  \to A$
が存在すると仮定してよい。Algebra, Lemma
\ref{algebra-lemma-quasi-finite-over-polynomial-algebra} を参照せよ。
したがって、$x$ の開近傍 $U \subset X$ と局所\footnote{$S$ が準分離ならば、
$g$ は準有限となる。}準有限射 $g : U \to \mathbf{A}^d_S$ からなる組
$(U, g)$ が少なくとも一つ存在する。

\medskip\noindent
主張：$R \to A$ が平坦かつ有限表示で、$\mathfrak p \subset A$ が素イデアル、
$\varphi : R[x_1, \ldots, x_d] \to A$ が $\mathfrak p$ で準有限ならば、
$\Spec(\varphi)$ が $\mathfrak p$ で平坦であることと
$\Spec(A) \to \Spec(R)$ が $\mathfrak p$ で Cohen--Macaulay であることは
同値である。実際、Theorem \ref{more-morphisms-theorem-criterion-flatness-fibre}
により平坦性はファイバー上で確認できる。Cohen--Macaulay であることについても
同じことが成り立つ（$A$ はすでに $R$ 上平坦と仮定されている）。したがって主張は
Algebra, Lemma \ref{algebra-lemma-where-CM} から従う。

\medskip\noindent
この主張から (1) と (4) の同値性が従う。また第 2 段落で適切な $(U, g)$ を
構成したことと合わせれば、この主張から (1) と (2) の同値性も従う。
\end{proof}

\begin{lemma}
\label{more-morphisms-lemma-flat-finite-presentation-CM-pieces}
$f : X \to S$ を平坦かつ局所有限表示なスキームの射とする。各 $d \geq 0$
に対し、次の性質をもつ開部分集合 $U_d \subset X$ が存在する。
\begin{enumerate}
\item $W = \bigcup_{d \geq 0} U_d$ は $f$ の各ファイバーで稠密である。
\item $U_d \to S$ は相対次元 $d$ である（Morphisms, Definition
\ref{morphisms-definition-relative-dimension-d} 参照）。
\end{enumerate}
\end{lemma}

\begin{proof}
Lemma \ref{more-morphisms-lemma-flat-finite-presentation-CM-open} と Morphisms, Lemma
\ref{morphisms-lemma-flat-finite-presentation-CM-fibres-relative-dimension}
を組み合わせれば従う。
\end{proof}

\begin{lemma}
\label{more-morphisms-lemma-flat-finite-presentation-specialization-dimension}
$f : X \to S$ を平坦かつ局所有限表示なスキームの射とする。
$X$ の点の特殊化 $x' \leadsto x$ の $S$ における像が $s' \leadsto s$
であると仮定する。$x$ が $X_s$ の既約成分の一般点ならば、
$\dim_{x'}(X_{s'}) = \dim_x(X_s)$ である。
\end{lemma}

\begin{proof}
% P05-EN-CH37-0110: frozen proof states only that x lies in some U_d and omits the final openness/specialization and relative-dimension steps.
ある $d$ に対して点 $x$ は $U_d$ に含まれる。ここで $U_d$ は Lemma
\ref{more-morphisms-lemma-flat-finite-presentation-CM-pieces} におけるものとする。
\end{proof}

\begin{lemma}
\label{more-morphisms-lemma-CM-local-source-and-target}
性質
$\mathcal{P}(f)=$「$f$ のファイバーは局所 Noetherian であり、$f$ は
Cohen--Macaulay である」は、終域上の fppf 位相について局所的であり、
始域上の syntomic 位相について局所的である。
\end{lemma}

\begin{proof}
次の分解がある。
$\mathcal{P}(f) =
\mathcal{P}_1(f) \wedge \mathcal{P}_2(f)$
ここで、$\mathcal{P}_1(f)=$「$f$ は平坦である」、および
$\mathcal{P}_2(f)=$「$f$ のファイバーは局所 Noetherian かつ
Cohen--Macaulay である」とする。$\mathcal{P}_1$ が始域と終域上の
fppf 位相について局所的であることは分かっている。Descent, Lemmas
\ref{descent-lemma-descending-property-flat} および
\ref{descent-lemma-flat-fpqc-local-source} を参照せよ。
したがって $\mathcal{P}_2$ を扱えばよい。

\medskip\noindent
$f : X \to Y$ をスキームの射とする。
$\{\varphi_i : Y_i \to Y\}_{i \in I}$ を $Y$ の fppf 被覆とする。
% P05-EN-CH37-0111: frozen source omits "by" in the naming construction for f_i; Japanese supplies normal syntax.
$f_i : X_i \to Y_i$ を $\varphi_i$ による $f$ の基底変換とする。
$i \in I$ とし、$y_i \in Y_i$ を点とする。$y = \varphi_i(y_i)$ とおく。このとき
$$
X_{i, y_i} = \Spec(\kappa(y_i)) \times_{\Spec(\kappa(y))} X_y.
$$
% P05-EN-CH37-0112: frozen display ends with a period before the continuation "and that"; Japanese uses a new sentence.
また $\kappa(y_i)/\kappa(y)$ は有限生成体拡大である。したがって $X_y$ が
局所 Noetherian ならば、$X_{i, y_i}$ も局所 Noetherian である。Varieties,
Lemma \ref{varieties-lemma-locally-Noetherian-base-change} を参照せよ。
さらに $X_y$ が Cohen--Macaulay ならば、$X_{i, y_i}$ も
Cohen--Macaulay である。Varieties, Lemma
\ref{varieties-lemma-CM-base-change} を参照せよ。
したがって $\mathcal{P}_2$ は終域上 fppf 局所的である。

\medskip\noindent
$\{X_i \to X\}$ を $X$ の syntomic 被覆とし、$y \in Y$ とする。この場合
$\{X_{i, y} \to X_y\}$ はファイバーの syntomic 被覆である。したがって
始域上の syntomic 位相に関する $\mathcal{P}_2$ の局所性は
Descent, Lemma \ref{descent-lemma-CM-local-syntomic}
から従う。以上を合わせて補題を得る。
\end{proof}






\section{Cohen--Macaulay 射の切断}
\label{more-morphisms-section-slicing-CM}

\noindent
この節の結果から最終的に、fppf 位相は「有限表示、平坦、準有限」位相と
同じであることが従う。次の補題は Divisors, Lemma
\ref{divisors-lemma-fibre-Cartier} と非常に密接に関係する。

\begin{lemma}
\label{more-morphisms-lemma-slice-given-element}
$f : X \to S$ をスキームの射とする。$x \in X$ を点とし、その像を
$s \in S$ とする。$h \in \mathfrak m_x \subset \mathcal{O}_{X, x}$ とする。
次を仮定する。
\begin{enumerate}
\item $f$ は局所有限表示である。
\item $f$ は点 $x$ で平坦である。
\item $h$ の
$\mathcal{O}_{X_s, x} = \mathcal{O}_{X, x}/\mathfrak m_s\mathcal{O}_{X, x}$
における像 $\overline{h}$ は非零因子である。
\end{enumerate}
このとき $x$ のアフィン開近傍 $U \subset X$ で、$h$ が
$h \in \Gamma(U, \mathcal{O}_U)$ から来て、$D = V(h)$ が $U$ の
有効 Cartier 因子となり、$x \in D$ かつ $D \to S$ が平坦かつ
局所有限表示となるものが存在する。
\end{lemma}

\begin{proof}
Noetherian の場合に帰着して証明する。平坦性の開性（Theorem
\ref{more-morphisms-theorem-openness-flatness} 参照）により、$X$ を $x$ の開近傍で
置き換えれば $X \to S$ は平坦であると仮定してよい。さらに $X$ と $S$ は
アフィンであると仮定してよい。必要なら $X$ を少し縮小することにより、
与えられた $h$ に写る $h \in \Gamma(X, \mathcal{O}_X)$ が存在すると仮定してよい。

\medskip\noindent
$S = \Spec(A)$ と書き、$A = \colim_i A_i$ を有限型 $\mathbf{Z}$-代数の
有向余極限として書く。Algebra, Lemma
\ref{algebra-lemma-flat-finite-presentation-limit-flat}、または Limits, Lemmas
\ref{limits-lemma-descend-finite-presentation}、
\ref{limits-lemma-descend-affine-finite-presentation}、および
% P05-EN-CH37-0113: frozen source repeats limits-lemma-descend-finite-presentation as both the first and third item in one citation list.
\ref{limits-lemma-descend-finite-presentation} により、Cartesian 図式
$$
\xymatrix{
X \ar[r] \ar[d]_f & X_0 \ar[d]^{f_0} \\
S \ar[r] & S_0
}
$$
で、$f_0$ が平坦かつ有限表示、$X_0$ がアフィン、$S_0$ がアフィンかつ
Noetherian となるものを得る。$x$ の像を $x_0 \in X_0$、$s$ の像を
$s_0 \in S_0$ とする。また、$X$ 上で $h$ に制限される元
$h_0 \in \Gamma(X_0, \mathcal{O}_{X_0})$ が存在すると仮定してよい。
（上の代数の参照結果を用いた場合は明らかである。極限の章の参照結果を
用いた場合は、$h$ を射 $X \to \mathbf{A}^1_S$ と考えれば Limits, Lemma
\ref{limits-lemma-descend-finite-presentation} から従う。）
$\mathcal{O}_{X_s, x}$ は
$\mathcal{O}_{(X_0)_{s_0}, x_0} \otimes_{\kappa(s_0)} \kappa(s)$ の局所化である。
したがって $\mathcal{O}_{(X_0)_{s_0}, x_0} \to \mathcal{O}_{X_s, x}$ は
平坦な局所環準同型であり、特に忠実平坦である。ゆえに像
$\overline{h}_0 \in \mathcal{O}_{(X_0)_{s_0}, x_0}$
は $\mathfrak m_{(X_0)_{s_0}, x_0}$ に属し、非零因子である。
$X_0$ を $x_0$ の主開近傍で置き換えると、元 $h_0$ は
$B_0 = \Gamma(X_0, \mathcal{O}_{X_0})$ の非零因子であり、かつ
$B_0/h_0B_0$ は $A_0 = \Gamma(S_0, \mathcal{O}_{S_0})$ 上平坦になると主張する。
実際、そうであれば
$$
0 \to B_0 \xrightarrow{h_0} B_0 \to B_0/h_0B_0 \to 0
$$
は平坦 $A_0$-加群の短完全列である。したがって $A$ とテンソル積しても完全性を保つ
（Algebra, Lemma \ref{algebra-lemma-flat-tor-zero} による）ので、補題が従う。

\medskip\noindent
上の主張を証明すればよい。対応する代数の主張は次のとおりである
（ここでは下添字 ${}_0$ を省く）。$A \to B$ を Noetherian 環の有限型平坦環準同型とする。
$\mathfrak q \subset B$ を $\mathfrak p \subset A$ の上にある素イデアルとする。
$h \in \mathfrak q$ の $B_{\mathfrak q}/\mathfrak p B_{\mathfrak q}$ における像は
非零因子であると仮定する。目標は、ある $g \in B$、$g \not \in \mathfrak q$
に対して $B$ を $B_g$ で置き換えた後、$h$ が非零因子となり、$B/hB$ が
$A$ 上平坦となることを示すことである。
% P05-EN-CH37-0114: frozen source says "after possible replacing" instead of "after possibly replacing".
Algebra, Lemma \ref{algebra-lemma-grothendieck} により、$h$ は
$B_{\mathfrak q}$ の非零因子であり、$B_{\mathfrak q}/hB_{\mathfrak q}$ は
$A$ 上平坦である。平坦性の開性、すなわち Algebra, Theorem
\ref{algebra-theorem-openness-flatness} または Theorem
\ref{more-morphisms-theorem-openness-flatness} により、ある $g \in B$、$g \not \in \mathfrak q$
に対して $B$ を $B_g$ で置き換えれば、$B/hB$ は $A$ 上平坦になる。
最後に $I = \{b \in B \mid hb = 0\}$ を $h$ の零化イデアルとする。
$h$ は $B_{\mathfrak q}$ の非零因子なので $IB_{\mathfrak q} = 0$ である。
また $B$ は Noetherian なので $I$ は有限生成である。したがって
$IB_g = 0$ となる $g \in B$、$g \not \in \mathfrak q$ が存在する。
$B$ を $B_g$ で置き換えれば、$h$ が非零因子であることが分かる。
\end{proof}

\begin{lemma}
\label{more-morphisms-lemma-slice-given-elements}
$f : X \to S$ をスキームの射とする。$x \in X$ を点とし、その像を
$s \in S$ とする。$h_1, \ldots, h_r \in \mathcal{O}_{X, x}$ とする。
次を仮定する。
\begin{enumerate}
\item $f$ は局所有限表示である。
\item $f$ は点 $x$ で平坦である。
\item $h_1, \ldots, h_r$ の
$\mathcal{O}_{X_s, x} = \mathcal{O}_{X, x}/\mathfrak m_s\mathcal{O}_{X, x}$
における像は正則列をなす。
\end{enumerate}
このとき $x$ のアフィン開近傍 $U \subset X$ で、$h_1, \ldots, h_r$ が
$h_1, \ldots, h_r \in \Gamma(U, \mathcal{O}_U)$ から来て、
$Z = V(h_1, \ldots, h_r) \to U$ が正則埋め込みとなり、$x \in Z$ かつ
$Z \to S$ が平坦かつ局所有限表示となるものが存在する。さらに、$S$ 上の
任意のスキーム $S'$ に対し、基底変換 $Z_{S'} \to U_{S'}$ は正則埋め込みである。
\end{lemma}

\begin{proof}
（正則列に関する本章の規約から、各 $i$ に対して $h_i \in \mathfrak m_x$
である。）$r = 1$ の場合は Lemma \ref{more-morphisms-lemma-slice-given-element} と
Divisors, Lemma \ref{divisors-lemma-relative-Cartier} を組み合わせ、
$V(h_1)$ が基底変換後も有効 Cartier 因子であることを用いれば従う。
$r > 1$ の場合は $r$ に関する直接的な帰納法から従う（$r = 1$ の結果を
ちょうど $r$ 回適用する。詳細は省略する）。

\medskip\noindent
Divisors, Section \ref{divisors-section-relative-regular-immersion} の内容を
用いて補題を証明することもできる。まず平坦性の開性（Theorem
\ref{more-morphisms-theorem-openness-flatness} 参照）により、$X$ を $x$ の開近傍で置き換えれば
$X \to S$ は平坦であると仮定してよい。また $X$ と $S$ はアフィンであると
仮定してよい。必要なら $X$ を少し縮小して
$h_1, \ldots, h_r \in \Gamma(X, \mathcal{O}_X)$ と仮定する。
$Z = V(h_1, \ldots, h_r)$ とおく。$X_s$ は代数的 $\kappa(s)$-スキームなので
Noetherian スキームであり（Varieties, Section
\ref{varieties-section-algebraic-schemes} 参照）、$X_s$ の位相は $X$ の位相から
誘導される（Schemes, Lemma \ref{schemes-lemma-fibre-topological} 参照）。
したがって $X$ をさらに少し縮小すれば、$Z_s \subset X_s$ は $r$ 個の元
$h_i|_{X_s}$ で切り出される正則埋め込みであると仮定してよい。Divisors, Lemma
\ref{divisors-lemma-Noetherian-scheme-regular-ideal} とその証明を参照せよ。
さらに、次の等式から $r = \dim_x(X_s) - \dim_x(Z_s)$ は明らかである。
\begin{align*}
\dim_x(X_s) & = \dim(\mathcal{O}_{X_s, x}) +
\text{trdeg}_{\kappa(s)}(\kappa(x)), \\
\dim_x(Z_s) & = \dim(\mathcal{O}_{Z_s, x}) +
\text{trdeg}_{\kappa(s)}(\kappa(x)), \\
\dim(\mathcal{O}_{X_s, x}) & = \dim(\mathcal{O}_{Z_s, x}) + r
\end{align*}
最初の二つの等式は Algebra, Lemma
\ref{algebra-lemma-dimension-at-a-point-finite-type-field} による。
% P05-EN-CH37-0115: frozen source says "the second" although the r-fold one-equation argument proves the third displayed equality.
第 2 の等式は Algebra, Lemma \ref{algebra-lemma-one-equation} を $r$ 回
適用することによる。したがって Divisors, Lemma
\ref{divisors-lemma-fibre-quasi-regular} の (3) を適用すると、$X$ を Zariski
開集合に縮小した後、射 $Z \to X$ は正則埋め込みであることが分かる。これに
Divisors, Lemma
\ref{divisors-lemma-relative-regular-immersion-flat-in-neighbourhood}
を適用でき、平坦性と基底変換に関する主張が得られる。
\end{proof}

\begin{lemma}
\label{more-morphisms-lemma-slice-once}
$f : X \to S$ をスキームの射とする。$x \in X$ を点とし、その像を
$s \in S$ とする。次を仮定する。
\begin{enumerate}
\item $f$ は局所有限表示である。
\item $f$ は点 $x$ で平坦である。
\item $\mathcal{O}_{X_s, x}$ は $\text{depth} \geq 1$ を満たす。
\end{enumerate}
このとき $x$ のアフィン開近傍 $U \subset X$ と、$x$ を含む有効 Cartier 因子
$D \subset U$ で、$D \to S$ が平坦かつ有限表示となるものが存在する。
\end{lemma}

\begin{proof}
$\mathcal{O}_{X_s, x}$ で非零因子に写る任意の
$h \in \mathfrak m_x \subset \mathcal{O}_{X, x}$ を選び、Lemma
\ref{more-morphisms-lemma-slice-given-element} を適用する。
\end{proof}

\begin{lemma}
\label{more-morphisms-lemma-slice-CM}
\begin{reference}
\cite[IV Proposition 17.16.1]{EGA}
\end{reference}
$f : X \to S$ をスキームの射とする。$x \in X$ を点とし、その像を
$s \in S$ とする。次を仮定する。
\begin{enumerate}
\item $f$ は局所有限表示である。
\item $f$ は点 $x$ で Cohen--Macaulay である。
\item $x$ は $X_s$ の閉点である。
\end{enumerate}
このとき $x$ を含む正則埋め込み $Z \to X$ で、次を満たすものが存在する。
\begin{enumerate}
\item[(a)] $Z \to S$ は平坦かつ局所有限表示である。
\item[(b)] $Z \to S$ は局所準有限である。
\item[(c)] 集合論的に $Z_s = \{x\}$ である。
\end{enumerate}
\end{lemma}

\begin{proof}
$S$ を $s$ のアフィン開近傍で置き換える。アフィン $S$ に対し、
$d = \dim_x(X_s)$ に関する帰納法で補題を証明する。

\medskip\noindent
$d = 0$ の場合。この場合、$Z$ として $x$ の開近傍を取れることを示す
（開埋め込みは正則埋め込みであることに注意せよ）。実際、$d = 0$ ならば
$X \to S$ は点 $x$ で準有限である。Morphisms, Lemma
\ref{morphisms-lemma-locally-quasi-finite-rel-dimension-0} を参照せよ。
したがって $U \to S$ が準有限となるアフィン開近傍 $U \subset X$ が存在する。
Morphisms, Lemma \ref{morphisms-lemma-quasi-finite-points-open} を参照せよ。
$X$ を $U$ で置き換えると、ファイバー $X_s$ は有限離散集合である。
% P05-EN-CH37-0116: frozen source says "a further affine open neighbourhood of X" where the required neighbourhood is of x.
そこで $X$ を $X$ のさらなるアフィン開近傍で置き換えると、
$f^{-1}(\{s\}) = \{x\}$ となる（$X_s$ の位相は $X$ の位相から誘導される。
Schemes, Lemma \ref{schemes-lemma-fibre-topological} を参照せよ）。
これでこの場合の補題が証明された。

\medskip\noindent
次に $d > 0$ と仮定する。$x$ はそのファイバーの閉点なので、Hilbert の零点定理により
拡大 $\kappa(x)/\kappa(s)$ は有限である。Morphisms, Lemma
\ref{morphisms-lemma-closed-point-fibre-locally-finite-type} を参照せよ。
したがって
$$
\text{depth}(\mathcal{O}_{X_s, x}) = \dim(\mathcal{O}_{X_s, x}) = d > 0
$$
を得る。最初の等式は $\mathcal{O}_{X_s, x}$ が Cohen--Macaulay であることによる。
第 2 の等式は Morphisms, Lemma
\ref{morphisms-lemma-dimension-fibre-at-a-point} による。したがって Lemma
\ref{more-morphisms-lemma-slice-once} を適用して図式
$$
\xymatrix{
D \ar[r] \ar[rrd] & U \ar[r] \ar[rd] & X \ar[d] \\
& & S
}
$$
で $x \in D$ となるものを得る。ある非零因子 $\overline{h}$ に対し
$\mathcal{O}_{D_s, x} = \mathcal{O}_{X_s, x}/(\overline{h})$
である。Divisors, Lemma \ref{divisors-lemma-relative-Cartier} を参照せよ。
したがって $\mathcal{O}_{D_s, x}$ は Cohen--Macaulay で、その次元は
$\mathcal{O}_{X_s, x}$ の次元より 1 小さい。例えば Algebra, Lemma
\ref{algebra-lemma-reformulate-CM} を参照せよ。ゆえに射 $D \to S$ は平坦、
局所有限表示であり、点 $x$ で Cohen--Macaulay であって、
$\dim_x(D_s) = \dim_x(X_s) - 1 = d - 1$ を満たす。帰納法の仮定により、
(a)、(b)、(c) を満たす正則埋め込み $Z \to D$ が存在する。
$Z \to D \to U$ はともに正則埋め込みなので、Divisors, Lemma
\ref{divisors-lemma-composition-regular-immersion} により $Z \to U$ も
正則埋め込みである。これで証明が完了する。
\end{proof}

\begin{lemma}
\label{more-morphisms-lemma-qf-fp-flat-neighbourhood-dominates-fppf}
$f : X \to S$ を局所有限表示なスキームの平坦射とする。$s \in S$ を
$f$ の像に属する点とする。このとき可換図式
$$
\xymatrix{
S' \ar[rr] \ar[rd]_g & & X \ar[ld]^f \\
& S
}
$$
で、$g : S' \to S$ が平坦、局所有限表示、局所準有限であり、
$s \in g(S')$ となるものが存在する。
\end{lemma}

\begin{proof}
仮定によりファイバー $X_s$ は空でない。したがって $f$ が Cohen--Macaulay
となる閉点 $x \in X_s$ が存在する。Lemma
\ref{more-morphisms-lemma-flat-finite-presentation-CM-open} を参照せよ。Lemma
\ref{more-morphisms-lemma-slice-CM} を適用し、$S' = Z$ とおく。
\end{proof}

\noindent
次の補題は、fppf 位相に対する層が「準有限、平坦、有限表示」位相に対する層と
同じものであることを示す。

\begin{lemma}
\label{more-morphisms-lemma-qf-fp-flat-dominates-fppf}
$S$ をスキームとする。$\mathcal{U} = \{S_i \to S\}_{i \in I}$ を $S$ の
fppf 被覆とする。Topologies, Definition
\ref{topologies-definition-fppf-covering} を参照せよ。このとき、
$\mathcal{U}$ を細分する fppf 被覆 $\mathcal{V} = \{T_j \to S\}_{j \in J}$
で、各 $T_j \to S$ が局所準有限となるものが存在する（Sites, Definition
\ref{sites-definition-morphism-coverings} 参照）。
\end{lemma}

\begin{proof}
各 $s \in S$ に対し、$s$ が $S_i \to S$ の像に属するような $i \in I$ が存在する。
Lemma \ref{more-morphisms-lemma-qf-fp-flat-neighbourhood-dominates-fppf} により、
$s \in g_s(T_s)$ を満たす射 $g_s : T_s \to S$ が存在する。この $g_s$ は
平坦、局所有限表示、局所準有限であり、$S_i \to S$ を経由する。したがって
$\{T_s \to S\}$ は $\mathcal{U}$ を細分する所望の $S$ の被覆である。
\end{proof}



\section{一般ファイバー}
\label{more-morphisms-section-generic}

\noindent
一般ファイバーと近傍のファイバーとの関係に関するいくつかの結果を示す。

\begin{lemma}
\label{more-morphisms-lemma-empty-generic-fibre}
$f : X \to Y$ をスキームの有限型射とする。$Y$ は既約で、一般点を
$\eta$ と仮定する。$X_\eta = \emptyset$ ならば、ある非空開集合
$V \subset Y$ が存在して $X_V = V \times_Y X = \emptyset$ となる。
\end{lemma}

\begin{proof}
より一般的な Morphisms, Lemma
\ref{morphisms-lemma-quasi-compact-generic-point-not-in-image} から直ちに従う。
\end{proof}

\begin{lemma}
\label{more-morphisms-lemma-nonempty-generic-fibre}
$f : X \to Y$ をスキームの有限型射とする。$Y$ は既約で、一般点を
$\eta$ と仮定する。$X_\eta \not = \emptyset$ ならば、ある非空開集合
$V \subset Y$ が存在して $X_V = V \times_Y X \to V$ は全射となる。
\end{lemma}

\begin{proof}
アフィン開集合を取れば、これは Algebra, Lemma
\ref{algebra-lemma-characterize-image-finite-type} から従う。
（もちろん一般平坦性からも従う。）
\end{proof}

\begin{lemma}
\label{more-morphisms-lemma-nowhere-dense-generic-fibre}
$f : X \to Y$ をスキームの有限型射とする。$Y$ は既約で、一般点を
$\eta$ と仮定する。$Z \subset X$ が閉部分集合で、$Z_\eta$ が
$X_\eta$ において疎であるならば、ある非空開集合 $V \subset Y$ が存在して、
すべての $y \in V$ に対し $Z_y$ は $X_y$ において疎となる。
\end{lemma}

\begin{proof}
$Y' \subset Y$ を $Y$ の被約化とする。
$X' = Y' \times_Y X$ および $Z' = Y' \times_Y Z$ とおく。
Morphisms, Lemma \ref{morphisms-lemma-reduction-universal-homeomorphism}
により $Y' \to Y$ は普遍同相なので、$Z' \subset X' \to Y'$ について
補題を証明すれば十分である。したがって Properties, Lemma
\ref{properties-lemma-characterize-integral} により、$Y$ は整であると仮定してよい。
Morphisms, Proposition \ref{morphisms-proposition-generic-flatness} により、
$X_V \to V$ と $Z_V \to V$ が平坦かつ有限表示となる非空アフィン開集合
$V \subset Y$ が存在する。この $V$ が所望のものであると主張する。
$y \in V$ を取る。$Z_y$ の内部が空でなければ、$Z_y$ は $X_y$ のある
既約成分の一般点 $\xi$ を含む。$\eta \leadsto f(\xi)$ であることに注意する。
$Z_V \to V$ は平坦なので、$f(\xi') = \eta$ を満たす特殊化
$\xi' \leadsto \xi$, $\xi' \in Z$ を選べる。Morphisms, Lemma
\ref{morphisms-lemma-generalizations-lift-flat} を参照せよ。Lemma
\ref{more-morphisms-lemma-flat-finite-presentation-specialization-dimension} により
$$
\dim_{\xi'}(Z_\eta) = \dim_{\xi}(Z_y) = \dim_{\xi}(X_y) = \dim_{\xi'}(X_\eta).
$$
したがって、$\xi'$ を通る $Z_\eta$ のある既約成分は次元
$\dim_{\xi'}(X_\eta)$ をもつ。これは $Z_\eta$ が $X_\eta$ において疎である
という仮定に反するので、証明が完了する。
\end{proof}

\begin{lemma}
\label{more-morphisms-lemma-scheme-theoretically-dense-generic-fibre}
$f : X \to Y$ をスキームの有限型射とする。$Y$ は既約で、一般点を
$\eta$ と仮定する。$U \subset X$ を、$U_\eta$ が $X_\eta$ において
スキーム論的に稠密となる開部分スキームとする。このとき、ある非空開集合
$V \subset Y$ が存在して、すべての $y \in V$ に対し $U_y$ は $X_y$ において
スキーム論的に稠密となる。
\end{lemma}

\begin{proof}
$Y' \subset Y$ を $Y$ の被約化とする。
$X' = Y' \times_Y X$ および $U' = Y' \times_Y U$ とおく。
$Y' \to Y$ は点集合上の全単射を誘導し、$U' \to U$ と $X' \to X$ は
スキーム論的ファイバーの同型を誘導するので、$Y$ を $Y'$ で、$X$ を
$X'$ で置き換えてよい。したがって Properties, Lemma
\ref{properties-lemma-characterize-integral} により、$Y$ は整であると仮定してよい。
さらに $Y$ を非空アフィン開集合で置き換えてよい。すなわち、$A$ が分数体
$K$ をもつ整域であって $Y = \Spec(A)$ であると仮定してよい。

\medskip\noindent
$f$ は有限型なので $X$ は準コンパクトである。あるアフィン開集合 $X_i$ を用いて
$X = X_1 \cup \ldots \cup X_n$ と書く。Morphisms, Definition
\ref{morphisms-definition-scheme-theoretically-dense} により、
$U_i = X_i \cap U$ は $X_i$ の開部分スキームで、$U_{i, \eta}$ は
$X_{i, \eta}$ においてスキーム論的に稠密である。したがって各組
$(X_i, U_i)$ について結果を証明すれば十分であり、言い換えれば $X$ は
アフィンであると仮定してよい。

\medskip\noindent
$X = \Spec(B)$ と書く。$B_K$ は有限型 $K$-代数なので Noetherian であることに
注意する。したがって $U_\eta$ は準コンパクトである。ゆえに、
$D(g_j) \subset U$ かつ
$U_\eta = D(g_1)_\eta \cup \ldots \cup D(g_m)_\eta$ となる有限個の元
$g_1, \ldots, g_m \in B$ を取れる。$U_\eta$ が $X_\eta$ において
スキーム論的に稠密であるという条件は、写像
$B_K \to \bigoplus_j (B_K)_{g_j}$ が単射であることを意味する。
Morphisms, Example \ref{morphisms-example-scheme-theoretic-closure} を参照せよ。
Algebra, Lemma \ref{algebra-lemma-when-injective-covering} により、これは写像
$B_K \to \bigoplus\nolimits_{j = 1, \ldots, m} B_K$,
$b \mapsto (g_1b, \ldots, g_mb)$ が単射であることと同値である。この写像を
$A$ 上で考えた余核を $M$ とする。すなわち、完全列
$$
0 \to I \to B \xrightarrow{(g_1, \ldots, g_m)}
\bigoplus\nolimits_{j = 1, \ldots, m} B \to M \to 0
$$
をもつ。ある非零元 $h$ に対して $A$ を $A_h$ で置き換えることにより、
$B$ は平坦な有限表示 $A$-代数であり、$M$ は $A$ 上平坦であると仮定してよい。
Algebra, Lemma \ref{algebra-lemma-generic-flatness} を参照せよ。
$B$ が $A$ 上平坦であることから、$B$ は $A$-加群として捩れなしである。
More on Algebra, Lemma \ref{more-algebra-lemma-flat-torsion-free} を参照せよ。
したがって $B \subset B_K$ である。仮定より $I_K = 0$ であり、これから
$I = 0$ が従う（$I \subset B \subset B_K$ は $I_K$ の部分集合だからである）。
ゆえに、いま短完全列
$$
0 \to B \xrightarrow{(g_1, \ldots, g_m)}
\bigoplus\nolimits_{j = 1, \ldots, m} B \to M \to 0
$$
があり、$M$ は $A$ 上平坦である。したがって、$\kappa$ が体である任意の準同型
$A \to \kappa$ に対し、短完全列
$$
0 \to B \otimes_A \kappa \xrightarrow{(g_1 \otimes 1, \ldots, g_m \otimes 1)}
\bigoplus\nolimits_{j = 1, \ldots, m} B \otimes_A \kappa \to
M \otimes_A \kappa \to 0
$$
を得る。Algebra, Lemma \ref{algebra-lemma-flat-tor-zero} を参照せよ。
上の議論を逆にたどると、これは $\bigcup D(g_j \otimes 1)$ が
$\Spec(B \otimes_A \kappa)$ においてスキーム論的に稠密であることを意味する。
$\bigcup D(g_j \otimes 1) = \bigcup D(g_j)_\kappa \subset U_\kappa$ なので、
$U_\kappa$ は $X_\kappa$ においてスキーム論的に稠密であり、これが示すべきことである。
\end{proof}

\noindent
スキームの射 $f : X \to Y$ と点 $y \in Y$ が与えられたとする。
ファイバー $X_y$ は、$X$ の部分集合 $f^{-1}(\{y\})$ に誘導位相を入れたものと
同相であることを思い出そう。Schemes, Lemma
\ref{schemes-lemma-fibre-topological} を参照せよ。閉部分集合
$T(y) \subset X_y$ が与えられたとし、$T$ を $X$ における $T(y)$ の閉包とする。
$T$ には誘導される被約スキーム構造を入れる。このとき $T$ は $X$ の閉部分スキームで、
集合論的に $T_y = T(y)$ を満たす。実際、$T$ はこの性質をもつ $X$ の最小の閉部分スキームである。
したがって、望むなら $X_y$ の閉部分集合を $T_y$ と記しても「差し支えない」。
次の補題では、これを $f$ の一般ファイバーに適用する。

\begin{lemma}
\label{more-morphisms-lemma-cover-generic-fibre-neighbourhood}
$f : X \to Y$ をスキームの有限型射とする。$Y$ は既約で、一般点を
$\eta$ と仮定する。$X_\eta$ の閉部分集合による一般ファイバーの被覆を
$X_\eta = Z_{1, \eta} \cup \ldots \cup Z_{n, \eta}$ とする。
$Z_i$ を $X$ における $Z_{i, \eta}$ の閉包とする（上の議論を参照せよ）。
このとき、ある非空開集合 $V \subset Y$ が存在して、すべての $y \in V$ に対し
$X_y = Z_{1, y} \cup \ldots \cup Z_{n, y}$ となる。
\end{lemma}

\begin{proof}
$Y$ が Noetherian ならば $U = X \setminus (Z_1 \cup \ldots \cup Z_n)$ は
$Y$ 上有限型であるから、Lemma \ref{more-morphisms-lemma-empty-generic-fibre} を直接適用して、
ある非空開集合 $V \subset Y$ に対し $U_V = \emptyset$ を得る。一般の場合は
次のように論じる。問題は位相的なので、$Y$ をその被約化で置き換えてよい。
したがって Properties, Lemma \ref{properties-lemma-characterize-integral} により、
$Y$ は整である。$Y$ を縮小した後、$X \to Y$ は平坦であると仮定してよい。
Morphisms, Proposition \ref{morphisms-proposition-generic-flatness} を参照せよ。
この場合、Morphisms, Lemma
\ref{morphisms-lemma-generalizations-lift-flat} により、$X_y$ の各点 $x$ は
ある点 $x' \in X_\eta$ の特殊化である。$Z_i$ は $X$ で閉じており一般ファイバーを
被覆するので、これは所望のとおり各 $y \in Y$ に対し
$X_y = \bigcup Z_{i, y}$ であることを意味する。
\end{proof}

\noindent
次の補題は、始域が被約である射の一般ファイバーは被約であることを述べる。

\begin{lemma}
\label{more-morphisms-lemma-reduction-generic-fibre}
$f : X \to Y$ をスキームの射とする。$\eta \in Y$ を $Y$ のある既約成分の
一般点とする。このとき
$(X_\eta)_{red} = (X_{red})_\eta$.
\end{lemma}

\begin{proof}
$\eta$ のアフィン近傍 $\Spec(A) \subset Y$ を選ぶ。射 $f$ によって
$\Spec(A)$ に写るアフィン開集合 $\Spec(B) \subset X$ を選ぶ。
$\mathfrak p \subset A$ を $\eta$ に対応する極小素イデアルとする。
$B_{red}$ を $B$ の冪零根基 $\sqrt{(0)}$ による商とする。この補題の代数的内容は
$C = B_{red} \otimes_A \kappa(\mathfrak p)$ が被約であるということである。
% P05-EN-CH37-0117: the frozen English has two malformed "Denote ..." constructions; see the sidecar.
$I \subset A$ を冪零根基と記せば $A_{red} = A/I$ である。また
$\mathfrak p_{red} = \mathfrak p/I$ と記す。これは $A_{red}$ の極小素イデアルで、
$\kappa(\mathfrak p) = \kappa(\mathfrak p_{red})$ を満たす。
$A \to B_{red}$ と $A \to \kappa(\mathfrak p)$ はともに $A \to A_{red}$ を
経由するので、$C = B_{red} \otimes_{A_{red}} \kappa(\mathfrak p_{red})$ である。
Algebra, Lemma \ref{algebra-lemma-minimal-prime-reduced-ring} により
$\kappa(\mathfrak p_{red}) = (A_{red})_{\mathfrak p_{red}}$ は局所化である。
したがって $C$ は $B_{red}$ の局所化であり（Algebra, Lemma
\ref{algebra-lemma-tensor-localization}）、ゆえに被約である。
\end{proof}

\begin{lemma}
\label{more-morphisms-lemma-make-generic-fibre-geometrically-reduced}
$f : X \to Y$ をスキームの射とする。$Y$ は既約で、$f$ は有限型であると仮定する。
次の図式が存在する。
$$
\xymatrix{
X' \ar[d]_{f'} \ar[r]_{g'} & X_V \ar[r] \ar[d] & X \ar[d]^f \\
Y' \ar[r]^g & V \ar[r] & Y
}
$$
ここで
\begin{enumerate}
\item $V$ は $Y$ の非空開集合であり、
\item $X_V = V \times_Y X$,
\item $g : Y' \to V$ は有限普遍同相であり、
\item $X' = (Y' \times_Y X)_{red} = (Y' \times_V X_V)_{red}$,
\item $g'$ は有限普遍同相であり、
\item $Y'$ は整アフィンスキームであり、
\item $f'$ は平坦かつ有限表示であり、かつ
\item $f'$ の一般ファイバーは幾何学的に被約である。
\end{enumerate}
\end{lemma}

\begin{proof}
$V = \Spec(A)$ を $Y$ の非空アフィン開集合とする。仮定により、$A$ の冪零根基
$\mathfrak p = \sqrt{(0)}$ は素イデアルである。$K = \kappa(\mathfrak p)$ とおく。
$p$ を、$K$ の標数が正ならばその標数、零ならば $1$ とする。
Varieties, Lemma \ref{varieties-lemma-finite-extension-geometrically-reduced}
により、有限純非分離体拡大 $K'/K$ で、$(X_{K'})_{red}$ が $K'$ 上
幾何学的に被約となるものが存在する。$K'$ を $K$ 上生成し、各 $x_i$ のある
$p$-冪が $A/\mathfrak p$ に属するような元 $x_1, \ldots, x_n \in K'$ を選ぶ。
$A' \subset K'$ を、$x_1, \ldots, x_n$ で生成される $K'$ の有限 $A$-部分代数とする。
$A'$ は分数体 $K'$ をもつ整域であることに注意する。Algebra, Lemma
\ref{algebra-lemma-p-ring-map} により、$A \to A'$ はスペクトル上に普遍同相を誘導する。
$Y' = \Spec(A')$ および $X' = (Y' \times_Y X)_{red}$ とおく。Lemma
\ref{more-morphisms-lemma-reduction-generic-fibre} により、$X' \to Y'$ の一般ファイバーは
$(X_{K'})_{red}$ であり、これは構成から幾何学的に被約である。
$X' \to X_V$ は、被約化射 $X' \to Y' \times_Y X$（Morphisms, Lemma
\ref{morphisms-lemma-reduction-universal-homeomorphism} を参照せよ）と $g$ の基底変換との
合成なので、有限普遍同相であることに注意する。この時点で、(7) を除けば補題の
すべての性質が成り立つ。Morphisms, Proposition
\ref{morphisms-proposition-generic-flatness} により、$Y'$ とそれに応じて $V$ を縮小すれば
(7) も達成できる。
\end{proof}

\begin{lemma}
\label{more-morphisms-lemma-make-components-generic-fibre-geometrically-irreducible}
$f : X \to Y$ をスキームの射とする。$Y$ は既約で、$f$ は有限型であると仮定する。
次の図式が存在する。
$$
\xymatrix{
X' \ar[d]_{f'} \ar[r]_{g'} & X_V \ar[r] \ar[d] & X \ar[d]^f \\
Y' \ar[r]^g & V \ar[r] & Y
}
$$
ここで
\begin{enumerate}
\item $V$ は $Y$ の非空開集合であり、
\item $X_V = V \times_Y X$,
\item $g : Y' \to V$ は全射な有限 \'etale 射であり、
\item $X' = Y' \times_Y X = Y' \times_V X_V$,
\item $g'$ は全射な有限 \'etale 射であり、
\item $Y'$ は既約アフィンスキームであり、かつ
\item $f'$ の一般ファイバーのすべての既約成分は幾何学的に既約である。
\end{enumerate}
\end{lemma}

\begin{proof}
$V = \Spec(A)$ を $Y$ の非空アフィン開集合とする。仮定により、$A$ の冪零根基
$\mathfrak p = \sqrt{(0)}$ は素イデアルである。$K = \kappa(\mathfrak p)$ とおく。
Varieties, Lemma
\ref{varieties-lemma-finite-extension-geometrically-irreducible-components}
により、有限分離体拡大 $K'/K$ で、$X_{K'}$ のすべての既約成分が $K'$ 上
幾何学的に既約となるものが存在する。$K'$ を $K$ 上生成する元
$\alpha \in K'$ を選ぶ。Fields, Lemma \ref{fields-lemma-primitive-element} を参照せよ。
$P(T) \in K[T]$ を $\alpha$ の $K$ 上の最小多項式とする。ある
$f \in A$, $f \not \in \mathfrak p$ に対して $\alpha$ を $f \alpha$ で
置き換えることにより、写像 $A[T] \to K[T]$ の下で $P(T) \in K[T]$ に写る
モニック多項式 $T^d + a_1T^{d - 1} + \ldots + a_d \in A[T]$ が存在すると
仮定してよい。
% P05-EN-CH37-0118: P was defined only in K[T], while the frozen source forms A[T]/(P); see the sidecar.
$A' = A[T]/(P)$ とおく。このとき $A \to A'$ は有限自由な環準同型で、
$\mathfrak p$ の上にある素イデアル $\mathfrak q$ がただ一つ存在し、
$K = \kappa(\mathfrak p) \subset \kappa(\mathfrak q) = K'$ は有限分離であり、
$\mathfrak pA'_{\mathfrak q}$ は $A'_{\mathfrak q}$ の極大イデアルとなる。
したがって $g : Y' = \Spec(A') \to V = \Spec(A)$ は $\mathfrak q$ において
\'etale である。Algebra, Lemma \ref{algebra-lemma-characterize-etale} を参照せよ。
これは、$g|_W : W \to \Spec(A)$ が \'etale となる開集合
$W \subset \Spec(A')$ が存在することを意味する。$g$ は有限であり、
$\mathfrak q$ は $\mathfrak p$ の上にある唯一の点なので、
$Z = g(Y' \setminus W)$ は $\mathfrak p$ を含まない $V$ の閉部分集合である。
したがって $V$ を $Z$ と交わらない $V$ の主アフィン開集合で置き換えれば、
$g$ は有限 \'etale となる。
\end{proof}







\section{相対随伴点}
\label{more-morphisms-section-assassin}

\begin{lemma}
\label{more-morphisms-lemma-relative-assassin-in-neighbourhood}
$f : X \to S$ をスキームの射とし、$\mathcal{F}$ を準連接
$\mathcal{O}_X$-加群とする。$\xi \in \text{Ass}_{X/S}(\mathcal{F})$ とし、
$Z = \overline{\{\xi\}} \subset X$ とおく。$f$ が局所有限型で、
$\mathcal{F}$ が有限型 $\mathcal{O}_X$-加群ならば、ある非空開集合
$V \subset Z$ が存在して、すべての $s \in f(V)$ に対し $V_s$ の一般点は
$\text{Ass}_{X/S}(\mathcal{F})$ の元となる。
\end{lemma}

\begin{proof}
$S$ を $f(\xi)$ のアフィン開近傍で、$X$ を $\xi$ のアフィン開近傍で
置き換えてよい。したがって $S = \Spec(A)$, $X = \Spec(B)$ で、$f$ は
有限型環準同型 $A \to B$ によって与えられると仮定してよい。
Morphisms, Lemma \ref{morphisms-lemma-locally-finite-type-characterize} を参照せよ。
さらに、ある有限 $B$-加群 $M$ によって $\mathcal{F} = \widetilde{M}$ と書ける。
Properties, Lemma \ref{properties-lemma-finite-type-module} を参照せよ。
$\mathfrak q \subset B$ を $\xi$ に対応する素イデアルとし、
$\mathfrak p \subset A$ をそれに対応する $A$ の素イデアルとする。仮定より
$\mathfrak q \in \text{Ass}_B(M \otimes_A \kappa(\mathfrak p))$ である。
Algebra, Remark \ref{algebra-remark-bourbaki} および Divisors, Lemma
\ref{divisors-lemma-associated-affine-open} を参照せよ。この記法の下で
$Z = V(\mathfrak q) \subset \Spec(B)$ であり、特に
$f(Z) \subset V(\mathfrak p)$ である。したがって $A$, $B$, $M$ をそれぞれ
$A/\mathfrak pA$, $B/\mathfrak pB$, $M/\mathfrak pM$ で置き換えた後に
補題を証明すれば明らかに十分である。言い換えれば、$A$ は分数体 $K$ をもつ
整域で、$\mathfrak q \subset B$ は $M \otimes_A K$ の随伴素イデアルであると
仮定してよい。

\medskip\noindent
ここで一般平坦性を用いることができる。すなわち Algebra, Lemma
\ref{algebra-lemma-generic-flatness} により、$M_g$ が $A_g$-加群として
平坦となる非零元 $g \in A$ が存在する。$A$ を $A_g$ で置き換えれば、
$M$ は $A$-加群として平坦であると仮定してよい。

\medskip\noindent
この場合、Algebra, Lemma \ref{algebra-lemma-post-bourbaki} により、
$\mathfrak q$ は $M$ の随伴素イデアルでもある。したがって単射な $B$-加群写像
$B/\mathfrak q \to M$ を得る。その余核を $Q$ とすれば、有限 $B$-加群の短完全列
$$
0 \to B/\mathfrak q \to M \to Q \to 0
$$
を得る。一般平坦性 Algebra, Lemma
\ref{algebra-lemma-generic-flatness} を今度は $B$-加群 $Q$ にもう一度適用すれば、
$Q$ は平坦な $A$-加群であると仮定してよい。特に上の短完全列は普遍的に単射であると
仮定してよい。Algebra, Lemma \ref{algebra-lemma-flat-tor-zero} を参照せよ。
この状況では、$A$ の任意の素イデアル $\mathfrak p'$ に対して
$(B/\mathfrak q) \otimes_A \kappa(\mathfrak p')
\subset M \otimes_A \kappa(\mathfrak p')$
である。このとき
$(B/\mathfrak q) \otimes_A \kappa(\mathfrak p')$
の支持の極小素イデアル $\mathfrak q'$ は、Divisors, Lemma
\ref{divisors-lemma-minimal-support-in-ass} により
$(B/\mathfrak q) \otimes_A \kappa(\mathfrak p')$ の随伴素イデアルであるから、
補題が従う。
\end{proof}

\begin{lemma}
\label{more-morphisms-lemma-bad-case}
$f : X \to Y$ をスキームの射とする。$\mathcal{F}$ を準連接
$\mathcal{O}_X$-加群とし、$U \subset X$ を開部分スキームとする。次を仮定する。
\begin{enumerate}
\item $f$ は有限型であり、
\item $\mathcal{F}$ は有限型であり、
\item $Y$ は既約で一般点を $\eta$ とし、かつ
\item $\text{Ass}_{X_\eta}(\mathcal{F}_\eta)$ は $U_\eta$ に含まれない。
\end{enumerate}
このとき、ある非空開部分スキーム $V \subset Y$ が存在して、すべての
$y \in V$ に対し集合 $\text{Ass}_{X_y}(\mathcal{F}_y)$ は $U_y$ に含まれない。
\end{lemma}

\begin{proof}
$Z \subset X$ を $\mathcal{F}$ のスキーム論的支持とする。Morphisms, Definition
\ref{morphisms-definition-scheme-theoretic-support} を参照せよ。このとき $Z_\eta$ は
$\mathcal{F}_\eta$ のスキーム論的支持である（Morphisms, Lemma
\ref{morphisms-lemma-flat-pullback-support}）。したがって Divisors, Lemma
\ref{divisors-lemma-minimal-support-in-ass} により、$Z_\eta$ の既約成分の一般点は
$\text{Ass}_{X_\eta}(\mathcal{F}_\eta)$ に含まれる。
% P05-EN-CH37-0119: the frozen proof makes two incompatible/non-sequitur deductions here; see the sidecar.
ゆえに $Z_\eta \cap U_\eta = \emptyset$ であることが分かる。したがって
$T = Z \setminus U$ は $T_\eta = \emptyset$ を満たす $Z$ の閉部分集合である。
$T$ に誘導される被約スキーム構造を入れれば、$T \to Y$ は有限型射である。
Lemma \ref{more-morphisms-lemma-empty-generic-fibre} により、$T_V = \emptyset$ を満たす
非空開集合 $V \subset Y$ が存在する。この $V$ が所望のものである。
\end{proof}

\begin{lemma}
\label{more-morphisms-lemma-good-case}
$f : X \to Y$ をスキームの射とする。$\mathcal{F}$ を準連接
$\mathcal{O}_X$-加群とし、$U \subset X$ を開部分スキームとする。次を仮定する。
\begin{enumerate}
\item $f$ は有限型であり、
\item $\mathcal{F}$ は有限型であり、
\item $Y$ は既約で一般点を $\eta$ とし、かつ
\item $\text{Ass}_{X_\eta}(\mathcal{F}_\eta) \subset U_\eta$.
\end{enumerate}
このとき、ある非空開部分スキーム $V \subset Y$ が存在して、すべての
$y \in V$ に対し $\text{Ass}_{X_y}(\mathcal{F}_y) \subset U_y$ となる。
\end{lemma}

\begin{proof}
（この証明は Lemma \ref{more-morphisms-lemma-scheme-theoretically-dense-generic-fibre} の証明と
同じである。まずそちらの証明を読むことを勧める。）主張はファイバーに関するものなので、
$Y$ をその被約化で置き換えてよいことは明らかである。したがって Properties, Lemma
\ref{properties-lemma-characterize-integral} により、$Y$ は整であると仮定してよい。
さらに $Y = \Spec(A)$ はアフィンであると仮定してよい。このとき $A$ は分数体
$K$ をもつ整域である。

\medskip\noindent
$f$ は有限型なので $X$ は準コンパクトである。あるアフィン開集合 $X_i$ によって
$X = X_1 \cup \ldots \cup X_n$ と書き、
$\mathcal{F}_i = \mathcal{F}|_{X_i}$ とおく。仮定により
$U_i = X_i \cap U$ の一般ファイバーは
$\text{Ass}_{X_{i, \eta}}(\mathcal{F}_{i, \eta})$ を含む。したがって三つ組
$(X_i, \mathcal{F}_i, U_i)$ について結果を証明すれば十分であり、言い換えれば
$X$ はアフィンであると仮定してよい。

\medskip\noindent
$X = \Spec(B)$ と書く。$\mathcal{F} = \widetilde{N}$ となる有限
$B$-加群 $N$ を取る。$B_K$ は有限型 $K$-代数なので Noetherian であることに
注意する。したがって $U_\eta$ は準コンパクトである。ゆえに、
$D(g_j) \subset U$ かつ
$U_\eta = D(g_1)_\eta \cup \ldots \cup D(g_m)_\eta$ となる有限個の元
$g_1, \ldots, g_m \in B$ を取れる。
$\text{Ass}_{X_\eta}(\mathcal{F}_\eta) \subset U_\eta$ なので、写像
$N_K \to \bigoplus_j (N_K)_{g_j}$ は単射である。Algebra, Lemma
\ref{algebra-lemma-when-injective-covering} により、これは写像
$N_K \to \bigoplus\nolimits_{j = 1, \ldots, m} N_K$,
$n \mapsto (g_1n, \ldots, g_mn)$ が単射であることと同値である。この写像を
$A$ 上で考えた核と余核をそれぞれ $I$ と $M$ とする。すなわち、完全列
$$
0 \to I \to N \xrightarrow{(g_1, \ldots, g_m)}
\bigoplus\nolimits_{j = 1, \ldots, m} N \to M \to 0
$$
をもつ。ある非零元 $h$ に対して $A$ を $A_h$ で置き換えることにより、
$B$ は平坦な有限表示 $A$-代数で、$M$ と $N$ はともに $A$ 上平坦であると
仮定してよい。Algebra, Lemma \ref{algebra-lemma-generic-flatness} を参照せよ。
$N$ が $A$ 上平坦であることから、$N$ は $A$-加群として捩れなしである。
More on Algebra, Lemma \ref{more-algebra-lemma-flat-torsion-free} を参照せよ。
したがって $N \subset N_K$ である。構成より $I_K = 0$ であり、これから
$I = 0$ が従う（$I \subset N \subset N_K$ は $I_K$ の部分集合だからである）。
ゆえに、いま短完全列
$$
0 \to N \xrightarrow{(g_1, \ldots, g_m)}
\bigoplus\nolimits_{j = 1, \ldots, m} N \to M \to 0
$$
があり、$M$ は $A$ 上平坦である。したがって、$\kappa$ が体である任意の準同型
$A \to \kappa$ に対し、短完全列
$$
0 \to N \otimes_A \kappa \xrightarrow{(g_1 \otimes 1, \ldots, g_m \otimes 1)}
\bigoplus\nolimits_{j = 1, \ldots, m} N \otimes_A \kappa \to
M \otimes_A \kappa \to 0
$$
を得る。Algebra, Lemma \ref{algebra-lemma-flat-tor-zero} を参照せよ。
上の議論を逆にたどると、これは $\bigcup D(g_j \otimes 1)$ が
$\text{Ass}_{B \otimes_A \kappa}(N \otimes_A \kappa)$ を含むことを意味する。
$\bigcup D(g_j \otimes 1) = \bigcup D(g_j)_\kappa \subset U_\kappa$ なので、
$U_\kappa$ は $\text{Ass}_{X \otimes \kappa}(\mathcal{F} \otimes \kappa)$ を含み、
これが示すべきことである。
\end{proof}

\begin{lemma}
\label{more-morphisms-lemma-base-change-assassin-in-U}
$f : X \to S$ を局所有限型な射とする。$\mathcal{F}$ を有限型の準連接
$\mathcal{O}_X$-加群とし、$U \subset X$ を開部分スキームとする。
$g : S' \to S$ をスキームの射とし、$f' : X' = X_{S'} \to S'$ を $f$ の
基底変換、$g' : X' \to X$ を射影とする。また
$\mathcal{F}' = (g')^*\mathcal{F}$ および $U' = (g')^{-1}(U)$ とおく。
最後に、$s' \in S'$ とし、その像を $s = g(s')$ とする。このとき
$$
\text{Ass}_{X_s}(\mathcal{F}_s) \subset U_s
\Leftrightarrow
\text{Ass}_{X'_{s'}}(\mathcal{F}'_{s'}) \subset U'_{s'}.
$$
\end{lemma}

\begin{proof}
Divisors, Lemma \ref{divisors-lemma-base-change-relative-assassin} から直ちに従う。
Divisors, Remark \ref{divisors-remark-base-change-relative-assassin} も参照せよ。
\end{proof}

\begin{lemma}
\label{more-morphisms-lemma-relative-assassin-constructible}
$f : X \to Y$ を有限表示射とする。$\mathcal{F}$ を有限表示の準連接
$\mathcal{O}_X$-加群とする。$U \subset X$ を、$U \to Y$ が準コンパクトとなる
開部分スキームとする。このとき集合
$$
E = \{y \in Y \mid \text{Ass}_{X_y}(\mathcal{F}_y) \subset U_y\}
$$
は $Y$ において局所構成可能である。
\end{lemma}

\begin{proof}
$y \in Y$ とする。$Y$ における $y$ のある開近傍 $V$ で、$E \cap V$ が $V$ において
構成可能となるものが存在することを示さなければならない。したがって $Y$ はアフィンであると
仮定してよい。$Y = \Spec(A)$ と書き、この環を有限型 $\mathbf{Z}$-代数の有向余極限
$A = \colim A_i$ として書く。Limits, Lemma
\ref{limits-lemma-descend-finite-presentation} により、ある $i$ と有限表示射
$f_i : X_i \to \Spec(A_i)$ で、その $Y$ への基底変換が $f$ を復元するものが存在する。
$i$ を必要なら大きくすることにより、有限表示の準連接 $\mathcal{O}_{X_i}$-加群
$\mathcal{F}_i$ で、その $X$ への引き戻しが $\mathcal{F}$ と同型になるものが存在すると
仮定してよい。Limits, Lemma
\ref{limits-lemma-descend-modules-finite-presentation} を参照せよ。さらにもう一度
$i$ を必要なら大きくすることにより、$X$ における逆像が $U$ となる開部分スキーム
$U_i \subset X_i$ が存在すると仮定してよい。Limits, Lemma
\ref{limits-lemma-descend-opens} を参照せよ。Lemma
\ref{more-morphisms-lemma-base-change-assassin-in-U} により、$f_i$ について補題を証明すれば十分である。
したがって $Y$ が Noetherian 環のスペクトルである場合に帰着する。

\medskip\noindent
$Y$ が Noetherian スキームの場合に $E$ が構成可能であることを示すため、
Topology, Lemma \ref{topology-lemma-characterize-constructible-Noetherian} の
判定法を用いる。実際、$Z \subset Y$ を既約閉部分スキームとする。
$E \cap Z$ が非空開部分集合を含むか、または $Z$ において稠密でないことを
示さなければならない。これは、$Z$ 上の基底変換
$(X, \mathcal{F}, U) \times_Y Z$ に Lemmas \ref{more-morphisms-lemma-bad-case} および
\ref{more-morphisms-lemma-good-case} を適用すれば従う。
\end{proof}











\section{被約ファイバー}
\label{more-morphisms-section-reduced}

\begin{lemma}
\label{more-morphisms-lemma-nonreduced-in-neighbourhood}
$f : X \to Y$ をスキームの射とする。$Y$ は既約で一般点を $\eta$ とし、
$f$ は有限型であると仮定する。$X_\eta$ が非被約ならば、ある非空開集合
$V \subset Y$ が存在して、すべての $y \in V$ に対しファイバー $X_y$ は非被約となる。
\end{lemma}

\begin{proof}
$Y' \subset Y$ を $Y$ の被約化とし、$X' \to Y'$ を $f$ の基底変換とする。
$Y' \to Y$ は点集合上の全単射を誘導し、$X' \to X$ はファイバーを同一視することに
注意する。
% P05-EN-CH37-0120: the frozen source says Y' rather than Y in this reduction step; see the sidecar.
したがって $Y'$ は被約、すなわち整であると仮定してよい。Properties, Lemma
\ref{properties-lemma-characterize-integral} を参照せよ。さらに $Y$ をアフィン開集合で
置き換えてよい。したがって $A$ が整域で $Y = \Spec(A)$ であると仮定してよい。
% P05-EN-CH37-0121: the frozen English naming construction omits "by" here and at its repeated occurrence; see the sidecar.
$K$ を $A$ の分数体とする。アフィン開集合 $\Spec(B) = U \subset X$ と、
非零かつ冪零な切断 $h_\eta \in \Gamma(U_\eta, \mathcal{O}_{U_\eta}) = B_K$ を選ぶ。
$Y$ を縮小することにより、$h$ は $h \in \Gamma(U, \mathcal{O}_U) = B$ から
来ると仮定してよい。$Y$ をさらに少し縮小すれば、$h$ は冪零であると仮定してよい。
$I = \{b \in B \mid hb = 0\}$ を $h$ の零化イデアルとする。このとき
$C = B/I$ は有限型 $A$-代数で、その一般ファイバー $(B/I)_K$ は非零である
（$h_\eta \not = 0$ だからである）。$A \to C$ と $A \to B/hB$ に一般平坦性を
適用する。Algebra, Lemma \ref{algebra-lemma-generic-flatness} を参照せよ。すると
$C_g$ が自由 $A_g$-加群で、$(B/hB)_g$ が平坦 $A_g$-加群となる
$g \in A$, $g \not = 0$ を得る。$Y$ を $D(g) \subset Y$ で置き換える。
いま短完全列
$$
0 \to C \to B \to B/hB \to 0.
$$
があり、$B/hB$ は $A$ 上平坦で、$C$ は非零な自由 $A$-加群である。
したがって、体への任意の準同型 $A \to \kappa$ に対し、環
$C \otimes_A \kappa$ は非零で、列
$$
0 \to C \otimes_A \kappa \to B \otimes_A \kappa \to
B/hB \otimes_A \kappa \to 0
$$
は完全である。Algebra, Lemma \ref{algebra-lemma-flat-tor-zero} を参照せよ。
テンソル積の右完全性により
$B/hB \otimes_A \kappa = (B \otimes_A \kappa) / h(B \otimes_A \kappa)$
であることに注意する。したがって $B \otimes_A \kappa$ 上の $h$ 倍写像は零でない。
これは明らかに、任意の点 $y \in Y$ に対して元 $h$ の $U_y$ への制限が非零であり、
したがって $X_y$ が非被約であることを意味する。
\end{proof}

\begin{lemma}
\label{more-morphisms-lemma-base-change-fibres-geometrically-reduced}
$f : X \to Y$ をスキームの射とする。$g : Y' \to Y$ を任意の射とし、
$f' : X' \to Y'$ を $f$ の基底変換と記す。このとき
\begin{align*}
\{y' \in Y' \mid X'_{y'}\text{ が幾何学的に被約}\} \\
= g^{-1}(\{y \in Y \mid X_y\text{ が幾何学的に被約}\}).
\end{align*}
\end{lemma}

\begin{proof}
これは、像を $y \in Y$ とする $y' \in Y'$ に対し、ファイバー
$X'_{y'} = X_y \times_y y'$ が $\kappa(y')$ 上幾何学的に被約であることと、
$X_y$ が $\kappa(y)$ 上幾何学的に被約であることが同値だという主張に帰着する。
これは Varieties, Lemma \ref{varieties-lemma-geometrically-reduced-upstairs} から従う。
\end{proof}

\begin{lemma}
\label{more-morphisms-lemma-not-geometrically-reduced-in-neighbourhood}
$f : X \to Y$ をスキームの射とする。$Y$ は既約で一般点を $\eta$ とし、
$f$ は有限型であると仮定する。$X_\eta$ が幾何学的に被約でないならば、
ある非空開集合 $V \subset Y$ が存在して、すべての $y \in V$ に対し
ファイバー $X_y$ は幾何学的に被約でない。
\end{lemma}

\begin{proof}
Lemma \ref{more-morphisms-lemma-make-generic-fibre-geometrically-reduced} を適用して
$$
\xymatrix{
X' \ar[d]_{f'} \ar[r]_{g'} & X_V \ar[r] \ar[d] & X \ar[d]^f \\
Y' \ar[r]^g & V \ar[r] & Y
}
$$
という図式を得る。これは同補題に挙げたすべての性質をもつ。$\eta'$ を $Y'$ の
一般点とする。射 $X' \to X_{Y'}$（これは被約化射である）と、それから得られる
一般ファイバーの射 $X'_{\eta'} \to X_{\eta'}$ を考える。
$X'_{\eta'}$ は幾何学的に被約である一方、$X_\eta$ はそうでないので、この射は
同型ではありえない。Varieties, Lemma
\ref{varieties-lemma-geometrically-reduced-upstairs} を参照せよ。したがって
$X_{\eta'}$ は非被約である。ゆえに Lemma
\ref{more-morphisms-lemma-nonreduced-in-neighbourhood} により、ある非空開集合 $V' \subset Y'$ の
すべての点 $y' \in V'$ において、$X_{Y'} \to Y'$ のファイバーは非被約である。
$g : Y' \to V$ は同相写像なので、Lemma
\ref{more-morphisms-lemma-base-change-fibres-geometrically-reduced} により、$g(V')$ が所望の
開集合である。
\end{proof}

\begin{lemma}
\label{more-morphisms-lemma-geometrically-reduced-generic-fibre}
$f : X \to Y$ をスキームの射とする。次を仮定する。
\begin{enumerate}
\item $Y$ は既約で一般点を $\eta$ とし、
\item $X_\eta$ は幾何学的に被約であり、かつ
\item $f$ は有限型である。
\end{enumerate}
このとき、ある非空開部分スキーム $V \subset Y$ が存在して、
$X_V \to V$ は幾何学的に被約なファイバーをもつ。
\end{lemma}

\begin{proof}
$Y' \subset Y$ を $Y$ の被約化とし、$X' \to Y'$ を $f$ の基底変換とする。
$Y' \to Y$ は点集合上の全単射を誘導し、$X' \to X$ はファイバーを同一視することに
注意する。
% P05-EN-CH37-0120: repeated frozen Y'/Y mismatch in the reduction step; see the sidecar.
したがって $Y'$ は被約、すなわち整であると仮定してよい。Properties, Lemma
\ref{properties-lemma-characterize-integral} を参照せよ。さらに $Y$ をアフィン開集合で
置き換えてよい。したがって $A$ が整域で $Y = \Spec(A)$ であると仮定してよい。
% P05-EN-CH37-0121: repeated malformed "Denote K" construction; see the sidecar.
$K$ を $A$ の分数体とする。$Y$ を少し縮小することにより、さらに $X \to Y$ は
平坦かつ有限表示であると仮定してよい。Morphisms, Proposition
\ref{morphisms-proposition-generic-flatness} を参照せよ。

\medskip\noindent
$X_\eta$ は幾何学的に被約なので、$V \to \Spec(K)$ が滑らかとなる開稠密部分集合
$V \subset X_\eta$ が存在する。Varieties, Lemma
\ref{varieties-lemma-geometrically-reduced-dense-smooth-open} を参照せよ。
$U \subset X$ を $f$ が滑らかとなる点の集合とする。Morphisms, Lemma
\ref{morphisms-lemma-set-points-where-fibres-smooth} により $V \subset U_\eta$ である。
したがって $U$ の一般ファイバーは $X$ の一般ファイバーにおいて稠密である。
$X_\eta$ は被約なので、$U_\eta$ は $X_\eta$ においてスキーム論的に稠密である。
Morphisms, Lemma \ref{morphisms-lemma-reduced-scheme-theoretically-dense} を参照せよ。
$U \to Y$ は滑らかなので、$U \to Y$ のすべてのファイバーは幾何学的に被約であることに注意する。
したがって $Y$ を縮小した後、すべての $y \in Y$ に対してスキーム $U_y$ が
$X_y$ においてスキーム論的に稠密であることを示せば十分である。Morphisms, Lemma
\ref{morphisms-lemma-reduced-subscheme-closure} を参照せよ。これは Lemma
\ref{more-morphisms-lemma-scheme-theoretically-dense-generic-fibre} から従う。
\end{proof}

\begin{lemma}
\label{more-morphisms-lemma-geometrically-reduced-constructible}
$f : X \to Y$ を準コンパクトかつ局所有限表示な射とする。このとき集合
$$
E = \{y \in Y \mid X_y\text{ が幾何学的に被約}\}
$$
は $Y$ において局所構成可能である。
\end{lemma}

\begin{proof}
$y \in Y$ とする。$Y$ における $y$ のある開近傍 $V$ で、$E \cap V$ が $V$ において
構成可能となるものが存在することを示さなければならない。したがって $Y$ はアフィンであると
仮定してよい。このとき $X$ は準コンパクトである。有限アフィン開被覆
$X = U_1 \cup \ldots \cup U_n$ を選ぶ。このとき $y$ における $U_i \to Y$ の
ファイバーは、$y$ における $X \to Y$ のファイバーのアフィン開被覆をなす。
したがって $X$ もアフィンであると仮定してよい。$Y = \Spec(A)$ と書く。
この環を有限型 $\mathbf{Z}$-代数の有向余極限 $A = \colim A_i$ として書く。
Limits, Lemma \ref{limits-lemma-descend-finite-presentation} により、ある $i$ と
有限表示射 $f_i : X_i \to \Spec(A_i)$ で、その $Y$ への基底変換が $f$ を
復元するものが存在する。Lemma
\ref{more-morphisms-lemma-base-change-fibres-geometrically-reduced} により、$f_i$ について補題を
証明すれば十分である。したがって $Y$ が Noetherian 環のスペクトルである場合に帰着する。

\medskip\noindent
$Y$ が Noetherian スキームの場合に $E$ が構成可能であることを示すため、
Topology, Lemma \ref{topology-lemma-characterize-constructible-Noetherian} の
判定法を用いる。実際、$Z \subset Y$ を一般点 $\xi$ をもつ既約閉部分スキームとする。
$E \cap Z$ が非空開部分集合を含むか、または $Z$ において稠密でないことを
示さなければならない。$X_\xi$ が幾何学的に被約ならば、射 $X_Z \to Z$ に適用した
Lemma \ref{more-morphisms-lemma-geometrically-reduced-generic-fibre} により、ある非空開集合
$V \subset Z$ に対してすべてのファイバー $X_y$ は幾何学的に被約である。
$X_\xi$ が幾何学的に被約でないならば、射 $X_Z \to Z$ に適用した Lemma
\ref{more-morphisms-lemma-not-geometrically-reduced-in-neighbourhood} により、ある非空開集合
$V \subset Z$ に対してすべてのファイバー $X_y$ は幾何学的に被約でない。
これで証明が完了する。
\end{proof}

\begin{lemma}
\label{more-morphisms-lemma-proper-flat-over-dvr-reduced-fibre}
$R$ を離散付値環とし、$f : X \to \Spec(R)$ を固有射とする。$X$ の各既約成分が
$\Spec(R)$ を優越すると仮定し（例えば $f$ が平坦ならばそうである）、特殊ファイバーは
被約であると仮定する。このとき $X$ と一般ファイバー $X_\eta$ はともに被約である。
\end{lemma}

\begin{proof}
（平坦射が既約成分に関する仮定を満たすことは、例えば Divisors, Lemma
\ref{divisors-lemma-bourbaki}、または平坦環準同型に対する下降定理から従う。）
特殊ファイバー $X_s$ は被約であると仮定する。任意の点 $x \in X$ を取り、
$\mathcal{O}_{X, x}$ が被約であることを示そう。これにより $X$ と $X_\eta$ が
被約であることが分かる。$x'$ が特殊ファイバーに属するような特殊化
$x \leadsto x'$ を取る。固有射は閉写像なので、このような特殊化は存在する。
局所環 $A = \mathcal{O}_{X, x'}$ を考える。$\mathcal{O}_{X, x}$ は $A$ の
局所化なので、$A$ が被約であることを示せば十分である。$\pi \in R$ を一意化元とする。
$a \in A$ が非零ならば、$a = \pi^n a'$ かつ $a' \not \in \pi A$ となる
$n \geq 0$ と元 $a' \in A$ が存在する。これは Krull の交叉定理
（Algebra, Lemma \ref{algebra-lemma-intersect-powers-ideal-module-zero}）から従う。
% P05-EN-CH37-0122: frozen prose has "the a" instead of "then a"; see the sidecar.
$a$ が冪零ならば $a$ は $A$ のすべての極小素イデアルに含まれ、したがって
$a'$ もそうである。実際、$X$ の既約成分に関する仮定により $\pi$ は $A$ の
どの極小素イデアルにも含まれないからである。ゆえに $a'$ は冪零である。
しかし $a'$ は被約環 $A/\pi A = \mathcal{O}_{X_s, x'}$ の非零元に写る。
$A$ が被約でなければこれは矛盾であり、所望の結論を得る。
\end{proof}

\begin{lemma}
\label{more-morphisms-lemma-geometrically-reduced-open}
$f : X \to Y$ を平坦かつ固有な有限表示射とする。このとき集合
$\{y \in Y \mid X_y\text{ が幾何学的に被約}\}$ は $Y$ において開である。
\end{lemma}

\begin{proof}
$Y$ はアフィンであると仮定してよい。このとき $Y$ は有限型
$\mathbf{Z}$-アフィンスキームの余フィルター付き極限である。したがって
$Y_0$ が有限型 $\mathbf{Z}$-代数のスペクトルで、$X_0 \to Y_0$ が平坦かつ
固有であり、$X \to Y$ は $X_0 \to Y_0$ の基底変換であると仮定してよい。
Limits, Lemma \ref{limits-lemma-descend-finite-presentation},
\ref{limits-lemma-descend-flat-finite-presentation}, および
\ref{limits-lemma-eventually-proper} を参照せよ。ファイバーが幾何学的に被約となる
点の集合の構成は基底変換と可換なので（Lemma
\ref{more-morphisms-lemma-base-change-fibres-geometrically-reduced}）、基底は Noetherian であると
仮定してよい。

\medskip\noindent
$Y$ は Noetherian であると仮定する。Lemma
\ref{more-morphisms-lemma-geometrically-reduced-constructible} により、この集合は構成可能である。
したがって、この集合が一般化で安定であることを示せば十分である（Topology, Lemma
\ref{topology-lemma-characterize-closed-Noetherian}）。Properties, Lemma
\ref{properties-lemma-locally-Noetherian-specialization-dvr} により、
$Y = \Spec(R)$ で $R$ は離散付値環、閉ファイバー $X_y$ は幾何学的に被約である
場合に帰着する。示すべきことは、一般ファイバー $X_\eta$ が幾何学的に被約であることである。

\medskip\noindent
そうでなければ、$R$ の分数体のある有限拡大 $L$ で $X_L$ が非被約となるものが
存在する。Varieties, Lemma \ref{varieties-lemma-geometrically-reduced} を参照せよ。
$R$ を優越し、分数体が $L$ である離散付値環 $R' \subset L$ が存在する。
Algebra, Lemma \ref{algebra-lemma-integral-closure-Dedekind} を参照せよ。
$R$ を $R'$ で置き換えれば、Lemma
\ref{more-morphisms-lemma-proper-flat-over-dvr-reduced-fibre} に帰着する。
\end{proof}







\section{ファイバーの既約成分}
\label{more-morphisms-section-irreducible}

\begin{lemma}
\label{more-morphisms-lemma-irreducible-components-in-neighbourhood}
$f : X \to Y$ をスキームの射とする。$Y$ は既約で一般点を $\eta$ とし、
$f$ は有限型であると仮定する。$X_\eta$ が $n$ 個の既約成分をもつならば、
ある非空開集合 $V \subset Y$ が存在して、すべての $y \in V$ に対しファイバー
$X_y$ は少なくとも $n$ 個の既約成分をもつ。
\end{lemma}

\begin{proof}
問題は純粋に位相的なので、$X$ と $Y$ をそれぞれの被約化で置き換えてよい。
特に Properties, Lemma \ref{properties-lemma-characterize-integral} により、
$Y$ は整である。$X_\eta$ の既約成分への分解を
$X_\eta = X_{1, \eta} \cup \ldots \cup X_{n, \eta}$ とする。
$X_i \subset X$ を、一般ファイバーが $X_{i, \eta}$ となる被約閉部分スキームとする。
$Z_{i, j} = X_i \cap X_j$ は $X_i$ の閉部分集合で、その一般ファイバー
$Z_{i, j, \eta}$ は $X_{i, \eta}$ において疎であることに注意する。
したがって $Y$ を縮小した後、すべての $y \in Y$ に対し $Z_{i, j, y}$ は
$X_{i, y}$ において疎であると仮定してよい。Lemma
\ref{more-morphisms-lemma-nowhere-dense-generic-fibre} を参照せよ。$Y$ をさらに縮小することにより、
$y \in Y$ に対し $X_y = \bigcup X_{i, y}$ であると仮定してよい。Lemma
\ref{more-morphisms-lemma-cover-generic-fibre-neighbourhood} を参照せよ。さらに $Y$ を縮小すれば、
各 $X_i \to Y$ は平坦かつ有限表示であると仮定してよい。Morphisms, Proposition
\ref{morphisms-proposition-generic-flatness} を参照せよ。射 $X_i \to Y$ は開である。
Morphisms, Lemma \ref{morphisms-lemma-fppf-open} を参照せよ。したがって、各 $i$ に
対して $f(X_i)$ に含まれる $\eta$ の開近傍 $V$ が存在する。各 $y \in V$ に対して
スキーム $X_{i, y}$ は $X_y$ の非空閉部分集合であり、
$X_y = \bigcup X_{i, y}$ で、交わり
$Z_{i, j, y} = X_{i, y} \cap X_{j, y}$ は $X_{i, y}$ において稠密でない。
これは明らかに $X_y$ が少なくとも $n$ 個の既約成分をもつことを意味する。
\end{proof}

\begin{lemma}
\label{more-morphisms-lemma-base-change-fibres-geometrically-irreducible}
$f : X \to Y$ をスキームの射とする。$g : Y' \to Y$ を任意の射とし、
$f' : X' \to Y'$ を $f$ の基底変換と記す。このとき
\begin{align*}
\{y' \in Y' \mid X'_{y'}\text{ が幾何学的に既約}\} \\
= g^{-1}(\{y \in Y \mid X_y\text{ が幾何学的に既約}\}).
\end{align*}
\end{lemma}

\begin{proof}
これは、像を $y \in Y$ とする $y' \in Y'$ に対し、ファイバー
$X'_{y'} = X_y \times_y y'$ が $\kappa(y')$ 上幾何学的に既約であることと、
$X_y$ が $\kappa(y)$ 上幾何学的に既約であることが同値だという主張に帰着する。
これは Varieties, Lemma
\ref{varieties-lemma-geometrically-irreducible-check-after-extension} から従う。
\end{proof}

\begin{lemma}
\label{more-morphisms-lemma-base-change-fibres-nr-geometrically-irreducible-components}
$f : X \to Y$ をスキームの射とする。
$$
n_{X/Y} : Y \to \{0, 1, 2, 3, \ldots, \infty\}
$$
を、$y \in Y$ に対して $(X_y)_K$ の既約成分の個数を対応させる関数とする。
ここで $K$ は $\kappa(y)$ の分離閉拡大である。これは良定義であり、
$g : Y' \to Y$ が射ならば
$$
n_{X'/Y'} = n_{X/Y} \circ g
$$
である。ここで $X' \to Y'$ は $f$ の基底変換である。
\end{lemma}

\begin{proof}
$y' \in Y'$ の像が $y \in Y$ であるとする。
$K \supset \kappa(y)$ と $K' \supset \kappa(y')$ を分離閉拡大とする。このとき体の可換図式
$$
\xymatrix{
K \ar[r] & K'' & K' \ar[l] \\
\kappa(y) \ar[u] \ar[rr] & & \kappa(y') \ar[u]
}
$$
を選べる。スキームの射
$$
\xymatrix{
(X'_{y'})_{K'} &
(X'_{y'})_{K''} = (X_y)_{K''} \ar[l] \ar[r] &
(X_y)_K
}
$$
が既約成分間の全単射を誘導することから結果が従う。Varieties, Lemma
\ref{varieties-lemma-separably-closed-field-irreducible-components} を参照せよ。
\end{proof}

\begin{lemma}
\label{more-morphisms-lemma-irreducible-polynomial-over-domain}
$A$ を分数体 $K$ をもつ整域とし、$P \in A[x_1, \ldots, x_n]$ とする。
% P05-EN-CH37-0123: the frozen English naming construction omits "by"; see the sidecar.
$\overline{K}$ を $K$ の代数閉包と記す。$P$ は
$\overline{K}[x_1, \ldots, x_n]$ において既約であると仮定する。このとき、
体へのすべての準同型 $\varphi : A_f \to \kappa$ に対し
$P^\varphi \in \kappa[x_1, \ldots, x_n]$ が既約となる $f \in A$ が存在する。
\end{lemma}

\begin{proof}
$A$ 上の $A[x_1, \ldots, x_n]$ の自己同型 $\Psi$ で、$a \not = 0$ かつ
$\Psi(P) = ax_n^d +$ $x_n$ に関する低次の項、となるものが存在する。
Algebra, Lemma \ref{algebra-lemma-helper-polynomial} を参照せよ。$P$ を $\Psi(P)$ で、
$A$ を $A_a$ で置き換えてよい。したがって $P$ は $x_n$ に関して次数
$d > 0$ のモニック多項式であると仮定してよい。$i = 1, \ldots, n - 1$ に対し、
$d_i$ を $P$ の $x_i$ に関する次数とする。このとき、$\kappa$ が体である任意の準同型
$\varphi : A \to \kappa$ に対し、$P^\varphi$ は $x_n$ に関して次数 $d$ の
モニック多項式であり、$x_i$ に関する次数は $\leq d_i$ であることに注意する。
したがって $P^\varphi$ が可約ならば
$$
P^\varphi = Q_1 Q_2
$$
% P05-EN-CH37-0124: allowing e_1 or e_2 = 0 creates the trivial factorization later excluded; see the sidecar.
と書ける。ここで $Q_1, Q_2$ は $x_n$ に関してそれぞれ次数
$e_1, e_2 \geq 0$ のモニック多項式で、$e_1 + e_2 = d$ を満たし、
$i = 1, \ldots, n - 1$ に対して $x_i$ に関する次数は $\leq d_i$ である。
言い換えれば
\begin{equation}
\label{more-morphisms-equation-factors}
Q_j = x_n^{e_j} + \sum\nolimits_{0 \leq l < e_j}
\left( \sum\nolimits_{L \in \mathcal{L}} a_{j, l, L} x^L \right) x_n^l
\end{equation}
と書ける。ここで和は、$0 \leq l_i \leq d_i$ を満たす形
$L = (l_1, \ldots, l_{n - 1})$ の多重指数 $L$ の集合 $\mathcal{L}$ 上で取る。
$e_1 + e_2 = d$ を満たす任意の $e_1, e_2 \geq 0$ に対し、$A$-代数
$$
B_{e_1, e_2} =
A[\{a_{1, l, L}\}_{0 \leq l < e_1, L \in \mathcal{L}},
\{a_{2, l, L}\}_{0 \leq l < e_2, L \in \mathcal{L}}]/(\text{関係式})
$$
を考える。ここで $(\text{関係式})$ は多項式
$$
P - Q_1Q_2 \in
A[\{a_{1, l, L}\}_{0 \leq l < e_1, L \in \mathcal{L}},
\{a_{2, l, L}\}_{0 \leq l < e_2, L \in \mathcal{L}}][x_1, \ldots, x_n]
$$
の係数で生成されるイデアルであり、$Q_1$ と $Q_2$ は
(\ref{more-morphisms-equation-factors}) のように定義する。さて、$P$ が $\overline{K}$ 上既約である
という仮定から、$A$-代数準同型 $B_{e_1, e_2} \to \overline{K}$ は存在しない。
Hilbert の零点定理（Algebra, Theorem \ref{algebra-theorem-nullstellensatz}）により、
これは $B_{e_1, e_2} \otimes_A K = 0$ を意味する。$B_{e_1, e_2}$ は有限生成
$A$-代数なので、$(B_{e_1, e_2})_{f_{e_1, e_2}} = 0$ を満たす
$f_{e_1, e_2} \in A$ が存在する。構成により、これは体への準同型
$\varphi : A_{f_{e_1, e_2}} \to \kappa$ に対し、$P^\varphi$ は
$Q_1$ の $x_n$ に関する次数が $e_1$、$Q_2$ の $x_n$ に関する次数が $e_2$ となる
因数分解 $P^\varphi = Q_1 Q_2$ をもたないことを意味する。したがって
$f = \prod_{e_1, e_2 \geq 0, e_1 + e_2 = d} f_{e_1, e_2}$ と取れば証明が完了する。
\end{proof}

\begin{lemma}
\label{more-morphisms-lemma-geom-irreducible-generic-fibre}
$f : X \to Y$ をスキームの射とする。次を仮定する。
\begin{enumerate}
\item $Y$ は既約で一般点を $\eta$ とし、
\item $X_\eta$ は幾何学的に既約であり、かつ
\item $f$ は有限型である。
\end{enumerate}
このとき、ある非空開部分スキーム $V \subset Y$ が存在して、
$X_V \to V$ は幾何学的に既約なファイバーをもつ。
\end{lemma}

\begin{proof}[Lemma \ref{more-morphisms-lemma-geom-irreducible-generic-fibre} の第1証明]
この補題に二つの証明を与える。両者は本質的に同値であり、第2証明の方が
自己完結的であるが少し長い。Lemma
\ref{more-morphisms-lemma-make-generic-fibre-geometrically-reduced} にあるような図式
$$
\xymatrix{
X' \ar[d]_{f'} \ar[r]_{g'} & X_V \ar[r] \ar[d] & X \ar[d]^f \\
Y' \ar[r]^g & V \ar[r] & Y
}
$$
を選ぶ。$f'$ の一般ファイバーは $f$ の一般ファイバーの被約化であり（Lemma
\ref{more-morphisms-lemma-reduction-generic-fibre} を参照せよ）、したがって幾何学的に既約であることに
注意する。射 $f'$ について補題が成り立つと仮定する。このとき $V$ を縮小すれば、
$f'$ のすべてのファイバーは幾何学的に既約となる。
$X' = (Y' \times_V X_V)_{red}$ なので、これは $Y' \times_V X_V$ のすべての
ファイバーが幾何学的に既約であることを意味する。したがって Lemma
\ref{more-morphisms-lemma-base-change-fibres-geometrically-irreducible} により、$X_V \to V$ の
すべてのファイバーは幾何学的に既約であり、証明が完了する。このようにして、
一般ファイバーは幾何学的に既約であるだけでなく幾何学的に被約でもあり、
$A$ が整域で $Y = \Spec(A)$ であると仮定してよい。

\medskip\noindent
$x \in X_\eta$ を一般点とする。$X_\eta$ は幾何学的に既約かつ被約なので、
$L = \kappa(x)$ は $K = \kappa(\eta)$ の有限生成拡大で、幾何学的に被約かつ
幾何学的に既約である。Varieties, Lemmas
\ref{varieties-lemma-geometrically-reduced-at-point} および
\ref{varieties-lemma-geometrically-irreducible-function-field} を参照せよ。特に体拡大
$L/K$ は分離的である。Algebra, Lemma
\ref{algebra-lemma-characterize-separable-field-extensions} を参照せよ。したがって、
$L$ を $K$ 上生成し、$x_1, \ldots, x_r$ が $L$ の $K$ 上の超越基底となる元
$x_1, \ldots, x_{r + 1} \in L$ を取れる。Algebra, Lemma
\ref{algebra-lemma-generating-finitely-generated-separable-field-extensions} を参照せよ。
$P \in K(x_1, \ldots, x_r)[T]$ を $x_{r + 1}$ の最小多項式とする。分母を払うことにより、
$P$ の係数は $A[x_1, \ldots, x_r]$ に属すると仮定してよい。$L$ は $K$ 上
幾何学的に被約かつ幾何学的に既約なので、多項式 $P$ は
$\overline{K}[x_1, \ldots, x_r, T]$ において既約であることに注意する。ここで
$\overline{K}$ は $K$ の代数閉包である。
% P05-EN-CH37-0125: the frozen source uses a bare "Denote" before this display and repeats the construction later; see the sidecar.
$$
B' = A[x_1, \ldots, x_{r + 1}]/(P(x_{r + 1}))
$$
と記し、$X' = \Spec(B')$ とおく。構成により $B'$ の分数体は、$K$-拡大として
$L = \kappa(x)$ と同型である。したがって開集合 $U \subset X$、開集合
$U' \subset X'$、および $Y$-同型 $U \to U'$ が存在する。Morphisms, Lemma
\ref{morphisms-lemma-common-open} を参照せよ。図式で表すと
$$
\xymatrix{
X \ar[rd] &
U \ar[l] \ar@{=}[r] \ar[d] &
U' \ar[r] \ar[d] &
X' \ar[ld] \ar@{=}[r] & \Spec(B') \\
& Y \ar@{=}[r] & Y &
}
$$
$U_\eta \subset X_\eta$ と $U'_\eta \subset X'_\eta$ は開稠密であることに注意する。
したがって Lemma \ref{more-morphisms-lemma-nowhere-dense-generic-fibre} を適用して $Y$ を縮小すれば、
すべての $y \in Y$ に対して $U_y$ は $X_y$ において稠密で、$U'_y$ は $X'_y$ において
稠密となる。ゆえに $X' \to Y$ について補題を証明すれば十分であり、これは Lemma
\ref{more-morphisms-lemma-irreducible-polynomial-over-domain} の内容である。
\end{proof}

\begin{proof}[Lemma \ref{more-morphisms-lemma-geom-irreducible-generic-fibre} の第2証明]
$Y' \subset Y$ を $Y$ の被約化とし、$X' \to X$ を $X$ の被約化とする。
$X' \to X  \to Y$ は $Y'$ を経由することに注意する。Schemes, Lemma
\ref{schemes-lemma-map-into-reduction} を参照せよ。Morphisms, Lemma
\ref{morphisms-lemma-reduction-universal-homeomorphism} により $Y' \to Y$ と
$X' \to X$ は普遍同相なので、$X' \to Y'$ について補題を証明すれば十分である。
したがって $X$ と $Y$ は被約であると仮定してよい。特に Properties, Lemma
\ref{properties-lemma-characterize-integral} により $Y$ は整である。
Morphisms, Proposition \ref{morphisms-proposition-generic-flatness} により、
$X_V \to V$ が平坦かつ有限表示となる非空アフィン開集合 $V \subset Y$ が存在する。
$Y$ を $V$ で置き換えることにより、(1), (2), (3) に加えて、$Y$ は整アフィン、
$X$ は被約、$f$ は平坦かつ有限表示であると仮定してよい。特に $f$ は普遍的に開である。
Morphisms, Lemma \ref{morphisms-lemma-fppf-open} を参照せよ。

\medskip\noindent
非空アフィン開集合 $U \subset X$ を選ぶ。このとき $U \to Y$ は平坦かつ有限表示で、
一般ファイバーは幾何学的に既約である。補集合 $X_\eta \setminus U_\eta$ は疎である。
したがって $Y$ を縮小した後、すべての $y \in Y$ に対し $U_y \subset X_y$ は
開稠密であると仮定してよい。Lemma \ref{more-morphisms-lemma-nowhere-dense-generic-fibre} を参照せよ。
ゆえに $X$ を $U$ で置き換えてよく、$Y$ は整アフィン、$X$ は被約アフィンで
$Y$ 上平坦かつ有限表示、一般ファイバー $X_\eta$ は幾何学的に既約である場合に帰着する。

\medskip\noindent
$X = \Spec(B)$, $Y = \Spec(A)$ と書く。このとき $A$ は整域、$B$ は被約、
$A \to B$ は平坦かつ有限表示で、$B_K$ は $A$ の分数体 $K$ 上幾何学的に既約である。
特に $B_K$ は整域である。$L$ を $B_K$ の分数体とする。$B$ は有限表示
$A$-代数なので、$L$ は $K$ の有限生成体拡大であることに注意する。
$(L \otimes_K K')_{red}$ が分離生成体拡大となるような有限純非分離拡大
$K'/K$ を取る。Algebra, Lemma \ref{algebra-lemma-make-separable} を参照せよ。
体拡大 $K'$ を $K$ 上生成し、ある素数冪 $q_i$ に対して $x_i^{q_i} \in A$ を
満たす元 $x_1, \ldots, x_n \in K'$ を選ぶ。
% P05-EN-CH37-0126: the frozen source explicitly omits the existence proof for these scaled generators; see the sidecar.
（$x_i$ の存在証明は省略する。）$A'$ を $x_1, \ldots, x_n$ で生成される
$K'$ の $A$-部分代数とする。このとき $A'$ は分数体が $K'$ である有限
$A$-部分代数 $A' \subset K'$ である。$\Spec(A') \to \Spec(A)$ は普遍同相であることに
注意する。Algebra, Lemma \ref{algebra-lemma-p-ring-map} を参照せよ。したがって
$\Spec(A')$ への基底変換後に結果を証明すれば十分である。$A$ を $A'$ で、$B$ を
$(B \otimes_A A')_{red}$ で置き換え、$L$ が $K$ の分離生成体拡大である状況に
移る。もちろん $(B \otimes_A A')_{red}$ は $A'$ 上平坦または有限表示でなくなる
可能性があるが、適切な $f \in A'$ に対して $A'$ を $A'_f$ で置き換えれば
修正できる。Algebra, Lemma \ref{algebra-lemma-generic-flatness} を参照せよ。

\medskip\noindent
この時点で、$A$ は整域、$B$ は被約、$A \to B$ は平坦かつ有限表示であり、
$B_K$ は整域で、その分数体 $L$ は $A$ の分数体 $K$ の分離生成体拡大である。
Algebra, Lemma
\ref{algebra-lemma-generating-finitely-generated-separable-field-extensions} により、
$L = K(x_1, \ldots, x_{r + 1})$ と書ける。ここで $x_1, \ldots, x_r$ は
$K$ 上代数的独立で、$x_{r + 1}$ は $K(x_1, \ldots, x_r)$ 上分離的である。
分母を払うことにより、$x_{r + 1}$ の $K(x_1, \ldots, x_r)$ 上の最小多項式
$P \in K(x_1, \ldots, x_r)[T]$ は $A[x_1, \ldots, x_r]$ に係数をもつと
仮定してよい。$L/K$ は分離的で、$L$ は $K$ 上幾何学的に既約なので、多項式
$P$ は $K$ の代数閉包 $\overline{K}$ 上既約であることに注意する。
% P05-EN-CH37-0125: repeated bare "Denote" before the display; see the sidecar.
$$
B' = A[x_1, \ldots, x_{r + 1}]/(P(x_{r + 1})).
$$
と記す。構成により $B$ と $B'$ の分数体は $K$-拡大として同型である。
したがって適切な $h \in B$ と $h' \in B'$ に対して $A$-代数の同型
$B_h \cong B'_{h'}$ が存在する。Morphisms, Lemma
\ref{morphisms-lemma-common-open} を参照せよ。言い換えれば、$X$ と
$X' = \Spec(B')$ は共通のアフィン開集合 $U$ をもつ。図式で表すと
$$
\xymatrix{
X = \Spec(B) \ar[rd] &
U \ar[l] \ar[r] \ar[d] &
\Spec(B') = X' \ar[ld] \\
& Y = \Spec(A) &
}
$$
$X$ における $Z = X \setminus U$ と $X'$ における $Z' = X' \setminus U$ に
Lemma \ref{more-morphisms-lemma-nowhere-dense-generic-fibre} を適用して $Y$ をもう一度縮小すれば、
すべての $y \in Y$ に対して $U_y$ は $X_y$ において稠密であり、$U_y$ は
$X'_y$ において稠密である。したがって $X' \to Y$ について補題を証明すれば十分で、
これは Lemma \ref{more-morphisms-lemma-irreducible-polynomial-over-domain} の内容である。
\end{proof}

\begin{lemma}
\label{more-morphisms-lemma-nr-geom-irreducible-components-good}
$f : X \to Y$ をスキームの射とする。$n_{X/Y}$ を、Lemma
\ref{more-morphisms-lemma-base-change-fibres-nr-geometrically-irreducible-components} で導入した、
$f$ のファイバーの幾何学的既約成分の個数を数える $Y$ 上の関数とする。
% P05-EN-CH37-0128: the frozen English omits "is" in two "Assume f of ..." clauses; see the sidecar.
$f$ は有限型であると仮定する。$y \in Y$ を点とする。このとき、
$n_{X/Y}|_V$ が定数となる非空開集合 $V \subset \overline{\{y\}}$ が存在する。
\end{lemma}

\begin{proof}
$Z$ を $\overline{\{y\}}$ に誘導される被約スキーム構造とする。
$f_Z : X_Z \to Z$ を $f$ の基底変換とする。$f_Z$ と $Z$ の一般点について補題を
証明すれば明らかに十分である。したがって Properties, Lemma
\ref{properties-lemma-characterize-integral} により、$Y$ は整スキームであると
仮定してよい。この場合の目標は、$n_{X/Y}|_V$ が定数となる非空開集合
$V \subset Y$ を構成することである。

\medskip\noindent
Lemma \ref{more-morphisms-lemma-make-components-generic-fibre-geometrically-irreducible} を
$f : X \to Y$ に適用し、$g : Y' \to V \subset Y$ を得る。
$g : Y' \to V$ は全射な有限 \'etale 射、特に開写像なので（Morphisms, Lemma
\ref{morphisms-lemma-etale-open} を参照せよ）、$n_{X'/Y'}|_{V'}$ が定数となる開集合
% P05-EN-CH37-0127: the frozen source omits the nonempty condition required for this reduction; see the sidecar.
$V' \subset Y'$ が存在することを証明すれば十分である。Lemma
\ref{more-morphisms-lemma-base-change-fibres-nr-geometrically-irreducible-components} を参照せよ。
したがって一般ファイバー $X_\eta$ のすべての既約成分は $\kappa(\eta)$ 上
幾何学的に既約であると仮定してよい。

\medskip\noindent
ここで
$X_\eta = X_{1, \eta} \bigcup \ldots \bigcup X_{n, \eta}$
が一般ファイバーの（幾何学的に）既約な成分への分解であるとする。
特に $n_{X/Y}(\eta) = n$ である。$X_i$ を $X$ における $X_{i, \eta}$ の閉包とする。
$Y$ を縮小した後、$X = \bigcup X_i$ であると仮定してよい。Lemma
\ref{more-morphisms-lemma-cover-generic-fibre-neighbourhood} を参照せよ。$Y$ をさらに縮小すれば、
$f$ の各ファイバーは少なくとも $n$ 個の既約成分をもつ。Lemma
\ref{more-morphisms-lemma-irreducible-components-in-neighbourhood} を参照せよ。したがって、
すべての $y \in Y$ に対して $n_{X/Y}(y) \geq n$ である。$Y$ をさらに縮小すれば、
各 $i$ とすべての $y \in Y$ に対して $X_{i, y}$ は幾何学的に既約となる。Lemma
\ref{more-morphisms-lemma-geom-irreducible-generic-fibre} を参照せよ。
$X_y = \bigcup X_{i, y}$ なので $n_{X/Y}(y) \leq n$ となり、証明が完了する。
\end{proof}

\begin{lemma}
\label{more-morphisms-lemma-nr-geom-irreducible-components-constructible}
$f : X \to Y$ をスキームの射とする。$n_{X/Y}$ を、Lemma
\ref{more-morphisms-lemma-base-change-fibres-nr-geometrically-irreducible-components} で導入した、
$f$ のファイバーの幾何学的既約成分の個数を数える $Y$ 上の関数とする。
% P05-EN-CH37-0128: repeated missing "is" in the finite-presentation assumption; see the sidecar.
$f$ は有限表示であると仮定する。このとき $n_{X/Y}$ のレベル集合
$$
E_n = \{y \in Y \mid n_{X/Y}(y) = n\}
$$
は $Y$ において局所構成可能である。
\end{lemma}

\begin{proof}
$n$ を固定し、$y \in Y$ とする。$Y$ における $y$ のある開近傍 $V$ で、
$E_n \cap V$ が $V$ において構成可能となるものが存在することを示さなければならない。
したがって $Y$ はアフィンであると仮定してよい。$Y = \Spec(A)$ と書き、
この環を有限型 $\mathbf{Z}$-代数の有向余極限 $A = \colim A_i$ として書く。
Limits, Lemma \ref{limits-lemma-descend-finite-presentation} により、ある $i$ と
有限表示射 $f_i : X_i \to \Spec(A_i)$ で、その $Y$ への基底変換が $f$ を
復元するものが存在する。Lemma
\ref{more-morphisms-lemma-base-change-fibres-nr-geometrically-irreducible-components} により、
$f_i$ について補題を証明すれば十分である。したがって $Y$ が Noetherian 環の
スペクトルである場合に帰着する。

\medskip\noindent
$Y$ が Noetherian スキームの場合に $E_n$ が構成可能であることを示すため、
Topology, Lemma \ref{topology-lemma-characterize-constructible-Noetherian} の
判定法を用いる。実際、$Z \subset Y$ を既約閉部分スキームとする。
$E_n \cap Z$ が非空開部分集合を含むか、または $Z$ において稠密でないことを
示さなければならない。$\xi \in Z$ を一般点とする。このとき Lemma
\ref{more-morphisms-lemma-nr-geom-irreducible-components-good} により、$n_{X/Y}$ は $Z$ における
$\xi$ の近傍で定数となる。これは明らかに所望の結論を与える。
\end{proof}













\section{ファイバーの連結成分}
\label{more-morphisms-section-connected}

\begin{lemma}
\label{more-morphisms-lemma-connected-components-in-neighbourhood}
$f : X \to Y$ をスキームの射とする。$Y$ は既約で一般点を $\eta$ とし、
$f$ は有限型であると仮定する。$X_\eta$ が $n$ 個の連結成分をもつならば、
ある非空開集合 $V \subset Y$ が存在して、すべての $y \in V$ に対しファイバー
$X_y$ は少なくとも $n$ 個の連結成分をもつ。
\end{lemma}

\begin{proof}
問題は純粋に位相的なので、$X$ と $Y$ をそれぞれの被約化で置き換えてよい。
特に Properties, Lemma \ref{properties-lemma-characterize-integral} により、
$Y$ は整である。$X_\eta$ の連結成分への分解を
$X_\eta = X_{1, \eta} \cup \ldots \cup X_{n, \eta}$ とする。
$X_i \subset X$ を、一般ファイバーが $X_{i, \eta}$ となる被約閉部分スキームとする。
$Z_{i, j} = X_i \cap X_j$ は $X$ の閉部分集合で、その一般ファイバー
$Z_{i, j, \eta}$ は空であることに注意する。したがって $Y$ を縮小した後、
$Z_{i, j} = \emptyset$ であると仮定してよい。Lemma
\ref{more-morphisms-lemma-empty-generic-fibre} を参照せよ。$Y$ をさらに縮小することにより、
$y \in Y$ に対し $X_y = \bigcup X_{i, y}$ であると仮定してよい。Lemma
\ref{more-morphisms-lemma-cover-generic-fibre-neighbourhood} を参照せよ。さらに $Y$ を縮小すれば、
各 $X_i \to Y$ は平坦かつ有限表示であると仮定してよい。Morphisms, Proposition
\ref{morphisms-proposition-generic-flatness} を参照せよ。射 $X_i \to Y$ は開である。
Morphisms, Lemma \ref{morphisms-lemma-fppf-open} を参照せよ。したがって、各 $i$ に
対して $f(X_i)$ に含まれる $\eta$ の開近傍 $V$ が存在する。各 $y \in V$ に対して
スキーム $X_{i, y}$ は $X_y$ の非空閉部分集合であり、
$X_y = \bigcup X_{i, y}$ で、交わり
$Z_{i, j, y} = X_{i, y} \cap X_{j, y}$ は空である。これは明らかに $X_y$ が
少なくとも $n$ 個の連結成分をもつことを意味する。
\end{proof}

\begin{lemma}
\label{more-morphisms-lemma-base-change-fibres-geometrically-connected}
$f : X \to Y$ をスキームの射とする。$g : Y' \to Y$ を任意の射とし、
$f' : X' \to Y'$ を $f$ の基底変換と記す。このとき
\begin{align*}
\{y' \in Y' \mid X'_{y'}\text{ が幾何学的に連結}\} \\
= g^{-1}(\{y \in Y \mid X_y\text{ が幾何学的に連結}\}).
\end{align*}
\end{lemma}

\begin{proof}
これは、像を $y \in Y$ とする $y' \in Y'$ に対し、ファイバー
$X'_{y'} = X_y \times_y y'$ が $\kappa(y')$ 上幾何学的に連結であることと、
$X_y$ が $\kappa(y)$ 上幾何学的に連結であることが同値だという主張に帰着する。
これは Varieties, Lemma
\ref{varieties-lemma-geometrically-connected-check-after-extension} から従う。
\end{proof}

\begin{lemma}
\label{more-morphisms-lemma-base-change-fibres-nr-geometrically-connected-components}
$f : X \to Y$ をスキームの射とする。
$$
n_{X/Y} : Y \to \{0, 1, 2, 3, \ldots, \infty\}
$$
を、$y \in Y$ に対して $(X_y)_K$ の連結成分の個数を対応させる関数とする。
ここで $K$ は $\kappa(y)$ の分離閉拡大である。これは良定義であり、
$g : Y' \to Y$ が射ならば
$$
n_{X'/Y'} = n_{X/Y} \circ g
$$
である。ここで $X' \to Y'$ は $f$ の基底変換である。
\end{lemma}

\begin{proof}
$y' \in Y'$ の像が $y \in Y$ であるとする。
$K \supset \kappa(y)$ と $K' \supset \kappa(y')$ を分離閉拡大とする。このとき体の可換図式
$$
\xymatrix{
K \ar[r] & K'' & K' \ar[l] \\
\kappa(y) \ar[u] \ar[rr] & & \kappa(y') \ar[u]
}
$$
を選べる。スキームの射
$$
\xymatrix{
(X'_{y'})_{K'} &
(X'_{y'})_{K''} = (X_y)_{K''} \ar[l] \ar[r] &
(X_y)_K
}
$$
が連結成分間の全単射を誘導することから結果が従う。Varieties, Lemma
\ref{varieties-lemma-separably-closed-field-connected-components} を参照せよ。
\end{proof}

\begin{lemma}
\label{more-morphisms-lemma-geometrically-connected-generic-fibre}
$f : X \to Y$ をスキームの射とする。次を仮定する。
\begin{enumerate}
\item $Y$ は既約で一般点を $\eta$ とし、
\item $X_\eta$ は幾何学的に連結であり、かつ
\item $f$ は有限型である。
\end{enumerate}
このとき、ある非空開部分スキーム $V \subset Y$ が存在して、
$X_V \to V$ は幾何学的に連結なファイバーをもつ。
\end{lemma}

\begin{proof}
Lemma \ref{more-morphisms-lemma-make-components-generic-fibre-geometrically-irreducible} にあるような図式
$$
\xymatrix{
X' \ar[d]_{f'} \ar[r]_{g'} & X_V \ar[r] \ar[d] & X \ar[d]^f \\
Y' \ar[r]^g & V \ar[r] & Y
}
$$
を選ぶ。$f'$ の一般ファイバーは幾何学的に連結であることに注意する（例えば Lemma
\ref{more-morphisms-lemma-base-change-fibres-nr-geometrically-connected-components} による）。
射 $f'$ について補題が成り立つと仮定する。すなわち、$W$ 上の
$X' \to Y'$ の各ファイバーが幾何学的に連結となる非空開集合 $W \subset Y'$ が
存在するとする。Morphisms, Lemma \ref{morphisms-lemma-etale-open} により $g$ は
開写像なので、非空開集合 $V = g(W)$ の点における $f$ のすべてのファイバーは
幾何学的に連結である。Lemma
\ref{more-morphisms-lemma-base-change-fibres-nr-geometrically-connected-components} を参照せよ。
このようにして、一般ファイバー $X_\eta$ の既約成分は幾何学的に既約であると
仮定してよい。

\medskip\noindent
$Y'$ を $Y$ の被約化とし、$X' = Y' \times_Y X$ とおく。このとき射
$X' \to Y'$ について補題を証明すれば十分である（例えばもう一度 Lemma
\ref{more-morphisms-lemma-base-change-fibres-nr-geometrically-connected-components} による）。
$X' \to Y'$ の一般ファイバーは $X \to Y$ の一般ファイバーと同じなので、
$Y$ は既約かつ被約、すなわち整であり（Properties, Lemma
\ref{properties-lemma-characterize-integral} を参照せよ）、一般ファイバー
$X_\eta$ の既約成分は幾何学的に既約であると仮定してよい。

\medskip\noindent
ここで
$X_\eta = X_{1, \eta} \bigcup \ldots \bigcup X_{n, \eta}$
が一般ファイバーの（幾何学的に）既約な成分への分解であるとする。
$X_i$ を $X$ における $X_{i, \eta}$ の閉包とする。$Y$ を縮小した後、
$X = \bigcup X_i$ であると仮定してよい。Lemma
\ref{more-morphisms-lemma-cover-generic-fibre-neighbourhood} を参照せよ。
$Z_{i, j} = X_i \cap X_j$ とおく。また
$$
\{1, \ldots, n\} \times \{1, \ldots, n\} = I \amalg J
$$
とする。ここで $Z_{i, j, \eta} = \emptyset$ ならば $(i, j) \in I$、
$Z_{i, j, \eta} \not = \emptyset$ ならば $(i, j) \in J$ とする。
$Y$ を縮小した後、すべての $(i, j) \in I$ に対し $Z_{i, j} = \emptyset$ であると
仮定してよい。Lemma \ref{more-morphisms-lemma-empty-generic-fibre} を参照せよ。$Y$ を縮小すれば、
各 $i$ とすべての $y \in Y$ に対して $X_{i, y}$ は幾何学的に既約となる。
Lemma \ref{more-morphisms-lemma-geom-irreducible-generic-fibre} を参照せよ。$Y$ をさらに縮小すれば、
すべての $(i, j) \in J$ に対して各 $Z_{i, j} \to Y$ は平坦かつ有限表示となる。
Morphisms, Proposition \ref{morphisms-proposition-generic-flatness} を参照せよ。
これは $f(Z_{i, j}) \subset Y$ が開であることを意味する。Morphisms, Lemma
\ref{morphisms-lemma-fppf-open} を参照せよ。ここで
$$
V  = \bigcap\nolimits_{(i, j) \in J} f(Z_{i, j})
$$
が所望のものである、すなわち各 $y \in V$ に対して $X_y$ は幾何学的に連結であると
主張する。実際、$X_\eta$ が連結であることから、$J$ の組で生成される同値関係は
ただ一つの同値類をもつ。いま $y \in V$ とし、$K \supset \kappa(y)$ を分離閉拡大と
すると、$(X_y)_K$ の既約成分はファイバー $(X_{i, y})_K$ である。さらに構成と
$y \in V$ により、$(X_{i, y})_K$ と $(X_{j, y})_K$ が交わることと
$(i, j) \in J$ は同値である。したがって同値類に関する上の考察から $(X_y)_K$ は
連結であり、証明が完了する。
\end{proof}

\begin{lemma}
\label{more-morphisms-lemma-nr-geom-connected-components-good}
$f : X \to Y$ をスキームの射とする。$n_{X/Y}$ を、Lemma
\ref{more-morphisms-lemma-base-change-fibres-nr-geometrically-connected-components} で導入した、
$f$ のファイバーの幾何学的連結成分の個数を数える $Y$ 上の関数とする。
% P05-EN-CH37-0133: the frozen English omits "is" in two "Assume f of ..." clauses; see the sidecar.
$f$ は有限型であると仮定する。$y \in Y$ を点とする。このとき、
$n_{X/Y}|_V$ が定数となる非空開集合 $V \subset \overline{\{y\}}$ が存在する。
\end{lemma}

\begin{proof}
$Z$ を $\overline{\{y\}}$ に誘導される被約スキーム構造とする。
$f_Z : X_Z \to Z$ を $f$ の基底変換とする。$f_Z$ と $Z$ の一般点について補題を
証明すれば明らかに十分である。したがって Properties, Lemma
\ref{properties-lemma-characterize-integral} により、$Y$ は整スキームであると
仮定してよい。この場合の目標は、$n_{X/Y}|_V$ が定数となる非空開集合
$V \subset Y$ を構成することである。

\medskip\noindent
Lemma \ref{more-morphisms-lemma-make-components-generic-fibre-geometrically-irreducible} を
$f : X \to Y$ に適用し、$g : Y' \to V \subset Y$ を得る。
$g : Y' \to V$ は全射な有限 \'etale 射、特に開写像なので（Morphisms, Lemma
\ref{morphisms-lemma-etale-open} を参照せよ）、$n_{X'/Y'}|_{V'}$ が定数となる開集合
% P05-EN-CH37-0131: the frozen source omits the nonempty condition required for this reduction; see the sidecar.
$V' \subset Y'$ が存在することを証明すれば十分である。
% P05-EN-CH37-0132: the frozen citation points to irreducible-, not connected-component base change; see the sidecar.
Lemma \ref{more-morphisms-lemma-base-change-fibres-nr-geometrically-irreducible-components} を参照せよ。
したがって一般ファイバー $X_\eta$ のすべての既約成分は $\kappa(\eta)$ 上
幾何学的に既約であると仮定してよい。Varieties, Lemma
\ref{varieties-lemma-irreducible-components-geometrically-irreducible} により、
これは $X_\eta$ の連結成分も幾何学的に連結であることを意味する。

\medskip\noindent
ここで
$X_\eta = X_{1, \eta} \bigcup \ldots \bigcup X_{n, \eta}$
が一般ファイバーの（幾何学的に）連結な成分への分解であるとする。
特に $n_{X/Y}(\eta) = n$ である。$X_i$ を $X$ における $X_{i, \eta}$ の閉包とする。
$Y$ を縮小した後、$X = \bigcup X_i$ であると仮定してよい。Lemma
\ref{more-morphisms-lemma-cover-generic-fibre-neighbourhood} を参照せよ。$Y$ をさらに縮小すれば、
$f$ の各ファイバーは少なくとも $n$ 個の連結成分をもつ。Lemma
\ref{more-morphisms-lemma-connected-components-in-neighbourhood} を参照せよ。したがって、
すべての $y \in Y$ に対して $n_{X/Y}(y) \geq n$ である。$Y$ をさらに縮小すれば、
各 $i$ とすべての $y \in Y$ に対して $X_{i, y}$ は幾何学的に連結となる。Lemma
\ref{more-morphisms-lemma-geometrically-connected-generic-fibre} を参照せよ。
$X_y = \bigcup X_{i, y}$ なので $n_{X/Y}(y) \leq n$ となり、証明が完了する。
\end{proof}

\begin{lemma}
\label{more-morphisms-lemma-nr-geom-connected-components-constructible}
$f : X \to Y$ をスキームの射とする。$n_{X/Y}$ を、Lemma
\ref{more-morphisms-lemma-base-change-fibres-nr-geometrically-connected-components} で導入した、
% P05-EN-CH37-0134: frozen prose says "geometric connected" rather than "geometrically connected"; see the sidecar.
$f$ のファイバーの幾何学的連結成分の個数を数える $Y$ 上の関数とする。
% P05-EN-CH37-0133: repeated missing "is" in the finite-presentation assumption; see the sidecar.
$f$ は有限表示であると仮定する。このとき $n_{X/Y}$ のレベル集合
$$
E_n = \{y \in Y \mid n_{X/Y}(y) = n\}
$$
は $Y$ において局所構成可能である。
\end{lemma}

\begin{proof}
$n$ を固定し、$y \in Y$ とする。$Y$ における $y$ のある開近傍 $V$ で、
$E_n \cap V$ が $V$ において構成可能となるものが存在することを示さなければならない。
したがって $Y$ はアフィンであると仮定してよい。$Y = \Spec(A)$ と書き、
この環を有限型 $\mathbf{Z}$-代数の有向余極限 $A = \colim A_i$ として書く。
Limits, Lemma \ref{limits-lemma-descend-finite-presentation} により、ある $i$ と
有限表示射 $f_i : X_i \to \Spec(A_i)$ で、その $Y$ への基底変換が $f$ を
復元するものが存在する。Lemma
\ref{more-morphisms-lemma-base-change-fibres-nr-geometrically-connected-components} により、
$f_i$ について補題を証明すれば十分である。したがって $Y$ が Noetherian 環の
スペクトルである場合に帰着する。

\medskip\noindent
$Y$ が Noetherian スキームの場合に $E_n$ が構成可能であることを示すため、
Topology, Lemma \ref{topology-lemma-characterize-constructible-Noetherian} の
判定法を用いる。実際、$Z \subset Y$ を既約閉部分スキームとする。
$E_n \cap Z$ が非空開部分集合を含むか、または $Z$ において稠密でないことを
示さなければならない。$\xi \in Z$ を一般点とする。このとき Lemma
\ref{more-morphisms-lemma-nr-geom-connected-components-good} により、$n_{X/Y}$ は $Z$ における
$\xi$ の近傍で定数となる。これは明らかに所望の結論を与える。
\end{proof}

\begin{lemma}
\label{more-morphisms-lemma-connected-flat-over-dvr}
\begin{slogan}
非連結スキームの平坦退化は、非連結であるか非被約である。
\end{slogan}
$f : X \to S$ をスキームの射とする。次を仮定する。
\begin{enumerate}
\item $S$ は離散付値環のスペクトルであり、
\item $f$ は平坦であり、
\item $X$ は連結であり、
\item 閉ファイバー $X_s$ は被約である。
\end{enumerate}
このとき一般ファイバー $X_\eta$ は連結である。
\end{lemma}

\begin{proof}
$S = \Spec(R)$ と書き、$\pi \in R$ を一意化元とする。矛盾を導くため、
$X_\eta$ は非連結であると仮定する。これは非自明な冪等元
$e \in \Gamma(X_\eta, \mathcal{O}_{X_\eta})$ が存在することを意味する。
$U = \Spec(A)$ を $X$ の任意のアフィン開集合とする。$A$ は $R$ 上平坦なので、
$\pi$ は $A$ 上の非零因子であることに注意する。例えば More on Algebra, Lemma
\ref{more-algebra-lemma-flat-torsion-free} を参照せよ。このとき $e|_{U_\eta}$ は
ある元 $e \in A[1/\pi]$ に対応する。$n \geq 0$ が最小で
$e = z/\pi^n$ となる元 $z \in A$ を取る。$z^2 = \pi^nz$ であることに注意する。
これは $n > 0$ ならば $z \bmod \pi A$ が冪零であることを意味する。仮定より
$A/\pi A$ は被約なので、$n$ の最小性から $n = 0$ が従う。したがって
$e \in A$ である。言い換えれば $e \in \Gamma(X, \mathcal{O}_X)$ である。
$X$ は連結なので $e$ は自明な冪等元であり、これは矛盾である。
\end{proof}







\section{切断と交わる連結成分}
\label{more-morphisms-section-connected-components}

\noindent
本節の結果は、特に群スキーム $G \to S$ とその単位切断
$e : S \to G$ に適用できる。

\begin{situation}
\label{more-morphisms-situation-connected-along-section}
% P05-EN-CH37-0135: the frozen English construction begins "Here ... be" without "let"; see the sidecar.
ここで $f : X \to Y$ をスキームの射とし、
$s : Y \to X$ を $f$ の切断とする。
% P05-EN-CH37-0136: the frozen English omits "by" in the naming construction; see the sidecar.
各 $y \in Y$ に対し、$s(y)$ を含む $X_y$ の連結成分を $X^0_y$ と表す。
最後に $X^0 = \bigcup_{y \in Y} X^0_y$ とおく。
\end{situation}

\begin{lemma}
\label{more-morphisms-lemma-base-change-connected-along-section}
$f : X \to Y$、$s : Y \to X$ を
Situation \ref{more-morphisms-situation-connected-along-section} のとおりとする。
$g : Y' \to Y$ を任意の射とし、基底変換図式
$$
\xymatrix{
X' \ar[r]_{g'} \ar[d]^{f'} & X \ar[d]_f \\
Y' \ar@/^1pc/[u]^{s'} \ar[r]^g & Y \ar@/_1pc/[u]_s
}
$$
を考え、それによって $(X')^0 \subset X'$ を得る。このとき
$(X')^0 = (g')^{-1}(X^0)$ である。
\end{lemma}

\begin{proof}
像を $y \in Y$ とする $y' \in Y'$ をとる。例えば Morphisms,
Definition \ref{morphisms-definition-scheme-structure-connected-component}
により、$X^0_y$ を $X_y$ の閉部分スキームとみなせる。
$s(y) \in X^0_y$ なので、Varieties, Lemma
\ref{varieties-lemma-geometrically-connected-if-connected-and-point}
から、$X_y^0$ は $\kappa(y)$ 上幾何学的に連結なスキームである。
したがって $X_y^0 \times_y y' \to X'_{y'}$ は $s'(y')$ を含む連結な
閉部分スキームである。ゆえに $X_y^0 \times_y y' \subset (X'_{y'})^0$ である。
逆の包含 $X_y^0 \times_y y' \supset (X'_{y'})^0$ は明らかである。実際、
$(X'_{y'})^0$ の $X_y$ における像は $s(y)$ を含む $X_y$ の連結部分集合である。
\end{proof}

\begin{lemma}
\label{more-morphisms-lemma-connected-along-section-good}
$f : X \to Y$、$s : Y \to X$ を
Situation \ref{more-morphisms-situation-connected-along-section} のとおりとする。
% P05-EN-CH37-0137: the frozen finite-type assumption omits "is"; see the sidecar.
$f$ は有限型であると仮定する。$y \in Y$ を一点とする。このとき、
ある非空開集合 $V \subset \overline{\{y\}}$ が存在して、$X^0$ の基底変換
$X_V$ における逆像は $X_V$ において開かつ閉である。
\end{lemma}

\begin{proof}
$Z \subset Y$ を、$\overline{\{y\}}$ に誘導される被約閉部分スキーム構造とする。
$f_Z : X_Z \to Z$ と $s_Z : Z \to X_Z$ を、それぞれ $f$ と $s$ の基底変換とする。
Lemma \ref{more-morphisms-lemma-base-change-connected-along-section} により
$(X_Z)^0 = (X^0)_Z$ である。
% P05-EN-CH37-0138: the frozen reduction replaces a base point by an X_Z-point over the generic point; see the sidecar.
したがって、射 $X_Z \to Z$ と、$Z$ の一般点に写る点 $x \in X_Z$ について
補題を証明すれば十分である。言い換えれば、問題は $Y$ が一般点 $\eta$ をもつ
整スキームである場合に帰着された（Properties, Lemma
\ref{properties-lemma-characterize-integral} を参照せよ）。目標は、$Y$ を縮小した後、
部分集合 $X^0$ が $X$ の開かつ閉な部分集合となることを示すことである。

\medskip\noindent
スキーム $X_\eta$ は体上有限型であり、したがって Noetherian であることに注意する。
ゆえにその連結成分は開かつ閉である。したがって、$X_\eta$ のある開かつ閉な部分集合
$T_\eta$ によって $X_\eta = X_\eta^0 \amalg T_\eta$ と書ける。
次に、$T \subset X$ を $T_\eta$ の閉包とし、$X^{00} \subset X$ を
$X_\eta^0$ の閉包とする。$T_\eta$ と $X^0_\eta$ は、それぞれ $T$ と $X^{00}$ の
一般ファイバーであることに注意する。Lemma
\ref{more-morphisms-lemma-cover-generic-fibre-neighbourhood} に先立つ議論を参照せよ。
さらに同補題により、$Y$ を縮小した後、集合論的に
$X = X^{00} \cup T$ であると仮定してよい。
$(T \cap X^{00})_\eta = T_\eta \cap X^0_\eta = \emptyset$ に注意する。
したがって $Y$ を縮小した後、$T \cap X^{00} = \emptyset$ であると仮定してよい。
Lemma \ref{more-morphisms-lemma-empty-generic-fibre} を参照せよ。特に $X^{00}$ は $X$ において開である。
$X^0_\eta$ は連結で、有理点 $s(\eta)$ をもつため幾何学的に連結であることに注意する。
Varieties, Lemma
\ref{varieties-lemma-geometrically-connected-if-connected-and-point} を参照せよ。
したがって $Y$ を縮小した後、$X^{00} \to Y$ のすべてのファイバーが幾何学的に
連結であると仮定してよい。Lemma
\ref{more-morphisms-lemma-geometrically-connected-generic-fibre} を参照せよ。
この時点で、ファイバー $X^{00}_y$ は $s(y)$ を含む $X_y$ の開かつ閉で連結な
部分集合であることが従う。ゆえに $X^0 = X^{00}$ であり、証明が完了する。
\end{proof}

\begin{lemma}
\label{more-morphisms-lemma-connected-along-section-locally-constructible}
$f : X \to Y$、$s : Y \to X$ を
Situation \ref{more-morphisms-situation-connected-along-section} のとおりとする。
$f$ が有限表示ならば、$X^0$ は $X$ において局所構成可能である。
\end{lemma}

\begin{proof}
$x \in X$ をとる。$X^0 \cap U$ が $U$ において構成可能となるような $x$ の
開近傍 $U$ が存在することを示さなければならない。これにより $Y$ がアフィンである
場合に帰着される。$Y = \Spec(A)$ と書き、$A = \colim A_i$ を有限型
$\mathbf{Z}$-代数の有向余極限として書く。Limits, Lemma
\ref{limits-lemma-descend-finite-presentation} により、ある $i$ と有限表示射
$f_i : X_i \to \Spec(A_i)$ で切断 $s_i : \Spec(A_i) \to X_i$ を備え、
$Y$ への基底変換が $f$ と切断 $s$ を復元するものが存在する。Lemma
\ref{more-morphisms-lemma-base-change-connected-along-section} により、$f_i, s_i$ について
補題を証明すれば十分である。したがって $Y$ が Noetherian 環のスペクトルである
場合に帰着される。

\medskip\noindent
$Y$ は Noetherian アフィンスキームであると仮定する。$f$ は有限表示、すなわち
有限型なので、$X$ も Noetherian スキームである。Morphisms, Lemma
\ref{morphisms-lemma-finite-type-noetherian} を参照せよ。この場合に補題を証明するには、
任意の既約閉部分集合 $Z \subset X$ に対し、交わり $Z \cap X^0$ が $Z$ の非空開集合を
含むか、または $Z$ において稠密でないことを示せば十分である。Topology, Lemma
\ref{topology-lemma-characterize-constructible-Noetherian} を参照せよ。
$x \in Z$ を一般点とし、$y = f(x)$ とおく。Lemma
\ref{more-morphisms-lemma-connected-along-section-good} により、$X^0 \cap X_V$ が $X_V$ において
開かつ閉となる非空開部分集合 $V \subset \overline{\{y\}}$ が存在する。
$f(Z) \subset \overline{\{y\}}$ かつ $f(x) = y \in V$ なので、
$W = f^{-1}(V) \cap Z$ は $Z$ の非空開部分集合である。したがって
$X^0 \cap W$ は $W$ において開かつ閉である。$W$ は既約なので、
$X^0 \cap W$ は空であるか $W$ に等しい。これで補題が証明された。
\end{proof}

\begin{lemma}
\label{more-morphisms-lemma-connected-along-section-open-neighbourhood}
$f : X \to Y$、$s : Y \to X$ を
Situation \ref{more-morphisms-situation-connected-along-section} のとおりとする。
$y \in Y$ を一点とする。次を仮定する。
\begin{enumerate}
\item $f$ は有限表示かつ平坦であり、
\item ファイバー $X_y$ は幾何学的に被約である。
\end{enumerate}
このとき $X^0$ は $X$ における $X^0_y$ の近傍である。
\end{lemma}

\begin{proof}
$Y$ を $y$ のアフィン開近傍で置き換えてよい。$Y = \Spec(A)$ と書き、
$A = \colim A_i$ を有限型 $\mathbf{Z}$-代数の有向余極限として書く。
Limits, Lemma \ref{limits-lemma-descend-finite-presentation} により、ある $i$ と
有限表示射 $f_i : X_i \to \Spec(A_i)$ で切断 $s_i : \Spec(A_i) \to X_i$ を備え、
$Y$ への基底変換が $f$ と切断 $s$ を復元するものが存在する。必要なら $i$ を増大させ、
$f_i$ も平坦であると仮定してよい。Limits, Lemma
\ref{limits-lemma-descend-flat-finite-presentation} を参照せよ。
$y_i$ を $Y_i$ における $y$ の像とする。$X_y = (X_{i, y_i}) \times_{y_i} y$ に注意する。
したがって $X_{i, y_i}$ は幾何学的に被約である。Varieties, Lemma
\ref{varieties-lemma-geometrically-reduced-upstairs} を参照せよ。Lemma
\ref{more-morphisms-lemma-base-change-connected-along-section} により、組 $f_i, s_i, y_i \in Y_i$ に
ついて補題を証明すれば十分である。したがって $Y$ が Noetherian 環のスペクトルである
場合に帰着される。

\medskip\noindent
$Y$ は Noetherian 環のスペクトルであると仮定する。$f$ は有限表示、すなわち有限型なので、
$X$ も Noetherian スキームである。Morphisms, Lemma
\ref{morphisms-lemma-finite-type-noetherian} を参照せよ。$x \in X^0$ を $y$ 上の点とする。
Topology, Lemma \ref{topology-lemma-constructible-neighbourhood-Noetherian} により、
$x$ を通る任意の既約閉集合 $Z \subset X$ に対して交わり $X^0 \cap Z$ が $Z$ において
稠密であることを証明すれば十分である。特に、一般点 $x' \in Z$ が $X^0$ に属することを
証明すれば十分である。Properties, Lemma
\ref{properties-lemma-locally-Noetherian-specialization-dvr} により、離散付値環 $R$ と、
閉点を $x$ に、一般点を $x'$ に写す射 $\Spec(R) \to X$ が存在する。合成
$\Spec(R) \to X \to Y$ によって、$\Spec(R)$ を $Y$ 上のスキームとみなす。
Lemma \ref{more-morphisms-lemma-base-change-connected-along-section} により、$(X_R)^0$ は $X^0$ の
逆像である。構成により、構造射 $X_R \to \Spec(R)$ は、$s$ の基底変換 $s_R$ のほかに
第二の切断 $t : \Spec(R) \to X_R$ をもつ。ここで $t(\eta_R)$ は $x'$ に写る $X_R$ の点、
$t(0_R)$ は $x$ に写る $X_R$ の点である。$t(0_R)$ は $(X_R)^0$ に属し、
$t(\eta_R) \leadsto t(0_R)$ であることに注意する。したがって、これから
$t(\eta_R) \in (X_R)^0$ が従うことを証明すれば十分である。ゆえに $Y$ が離散付値環の
スペクトルで、$y$ がその閉点である場合に補題を証明すれば十分である。

\medskip\noindent
$Y$ は離散付値環のスペクトルで、$y$ はその閉点であると仮定する。目標は $X^0$ が
$X^0_y$ の近傍であることを証明することである。$X_y$ は有限個の既約成分をもつので、
$X^0_y$ は $X_y$ において開かつ閉であることに注意する。したがって補集合
$C = X_y \setminus X_y^0$ は $X$ において閉である。ゆえに $U = X \setminus C$ は
$X^0_y$ の開近傍であり、$U^0 = X^0$ である。したがって射 $U \to Y$ について結果を
証明すれば十分である。言い換えれば、$X_y$ は連結であると仮定してよい。
$X$ が非連結であるとし、$X = X_1 \amalg \ldots \amalg X_n$ を連結成分への分解とする。
このとき $s(Y)$ はいずれか一つの $X_i$ に完全に含まれる。$s(Y) \subset X_1$ とする。
すると $X^0 \subset X_1$ である。ゆえに $X$ を $X_1$ で置き換え、$X$ は連結であると
仮定してよい。この時点で Lemma \ref{more-morphisms-lemma-connected-flat-over-dvr} により $X_\eta$ は
連結である。すなわち $X^0 = X$ であり、証明が完了する。
\end{proof}

\begin{lemma}
\label{more-morphisms-lemma-connected-along-section-open}
$f : X \to Y$、$s : Y \to X$ を
Situation \ref{more-morphisms-situation-connected-along-section} のとおりとする。
次を仮定する。
\begin{enumerate}
\item $f$ は有限表示かつ平坦であり、
\item $f$ のすべてのファイバーは幾何学的に被約である。
\end{enumerate}
このとき $X^0$ は $X$ において開である。
\end{lemma}

\begin{proof}
これは Lemma \ref{more-morphisms-lemma-connected-along-section-open-neighbourhood} の
直接の帰結である。
\end{proof}






\section{ファイバーの次元}
\label{more-morphisms-section-dimension}

\begin{lemma}
\label{more-morphisms-lemma-dimension-in-neighbourhood}
$f : X \to Y$ をスキームの射とする。
% P05-EN-CH37-0139: the frozen paired hypotheses both omit "is"; see the sidecar.
$Y$ は一般点 $\eta$ をもつ既約スキームで、$f$ は有限型であると仮定する。
$X_\eta$ の次元が $n$ ならば、ある非空開集合 $V \subset Y$ が存在して、
すべての $y \in V$ に対しファイバー $X_y$ の次元は $n$ である。
\end{lemma}

\begin{proof}
$Z = \{x \in X \mid \dim_x(X_{f(x)}) > n \}$ とおく。Morphisms, Lemma
\ref{morphisms-lemma-openness-bounded-dimension-fibres} により、これは $X$ の
閉部分集合である。仮定より $Z_\eta = \emptyset$ である。したがって Lemma
\ref{more-morphisms-lemma-empty-generic-fibre} により、$Y$ を縮小して $Z = \emptyset$ と仮定してよい。
次に
$Z' = \{x \in X \mid \dim_x(X_{f(x)}) > n - 1 \} =
\{x \in X \mid \dim_x(X_{f(x)}) = n\}$ とおく。前と同様に、これは $X$ の
閉部分集合である。仮定より $Z'_\eta \not = \emptyset$ である。したがって $Y$ を
縮小した後、$Z' \to Y$ は全射であると仮定してよい。Lemma
\ref{more-morphisms-lemma-nonempty-generic-fibre} を参照せよ。これで証明が完了する。
\end{proof}

\begin{lemma}
\label{more-morphisms-lemma-base-change-dimension-fibres}
$f : X \to Y$ を有限型射とする。
% P05-EN-CH37-0141: the frozen codomain omits -infinity, the dimension of an empty fibre; see the sidecar.
$$
n_{X/Y} : Y \to \{0, 1, 2, 3, \ldots, \infty\}
$$
を、$y \in Y$ に $X_y$ の次元を対応させる関数とする。
$g : Y' \to Y$ が射ならば
$$
n_{X'/Y'} = n_{X/Y} \circ g
$$
である。ここで $X' \to Y'$ は $f$ の基底変換である。
\end{lemma}

\begin{proof}
これは Morphisms, Lemma
\ref{morphisms-lemma-dimension-fibre-after-base-change} から従う。
\end{proof}

\begin{lemma}
\label{more-morphisms-lemma-dimension-fibres-constructible}
$f : X \to Y$ をスキームの射とする。$n_{X/Y}$ を、Lemma
\ref{more-morphisms-lemma-base-change-dimension-fibres} で導入した、$f$ のファイバーの次元を
与える $Y$ 上の関数とする。
% P05-EN-CH37-0140: the frozen finite-presentation hypothesis omits "is"; see the sidecar.
$f$ は有限表示であると仮定する。このとき $n_{X/Y}$ の等位集合
$$
E_n = \{y \in Y \mid n_{X/Y}(y) = n\}
$$
は $Y$ において局所構成可能である。
\end{lemma}

\begin{proof}
$n$ を固定する。$y \in Y$ をとる。$E_n \cap V$ が $V$ において構成可能となるような
$Y$ における $y$ の開近傍 $V$ が存在することを示さなければならない。したがって
$Y$ はアフィンであると仮定してよい。$Y = \Spec(A)$ と書き、$A = \colim A_i$ を
有限型 $\mathbf{Z}$-代数の有向余極限として書く。Limits, Lemma
\ref{limits-lemma-descend-finite-presentation} により、ある $i$ と有限表示射
$f_i : X_i \to \Spec(A_i)$ で、その $Y$ への基底変換が $f$ を復元するものが存在する。
Lemma \ref{more-morphisms-lemma-base-change-dimension-fibres} により、$f_i$ について補題を証明すれば
十分である。したがって $Y$ が Noetherian 環のスペクトルである場合に帰着される。

\medskip\noindent
$Y$ が Noetherian スキームの場合に $E_n$ が構成可能であることを証明するため、
Topology, Lemma \ref{topology-lemma-characterize-constructible-Noetherian} の
判定法を用いる。このため $Z \subset Y$ を既約閉部分スキームとする。
$E_n \cap Z$ が非空開部分集合を含むか、または $Z$ において稠密でないことを
示さなければならない。$\xi \in Z$ を一般点とする。このとき Lemma
\ref{more-morphisms-lemma-dimension-in-neighbourhood} により、$n_{X/Y}$ は $Z$ における $\xi$ の
近傍上で定数である。これは求める主張を与える。
\end{proof}

\begin{lemma}
\label{more-morphisms-lemma-dimension-fibres-flat}
$f : X \to Y$ を有限表示な平坦射とする。$n_{X/Y}$ を、Lemma
\ref{more-morphisms-lemma-base-change-dimension-fibres} で導入した、$f$ のファイバーの次元を与える
$Y$ 上の関数とする。このとき $n_{X/Y}$ は下半連続である。
\end{lemma}

\begin{proof}
Lemmas \ref{more-morphisms-lemma-flat-finite-presentation-CM-open} と
\ref{more-morphisms-lemma-flat-finite-presentation-CM-pieces} で構成された開集合を
$W \subset X$, $W = \coprod_{d \geq 0} U_d$ とする。$y \in Y$ を一点とする。
$n_{X/Y}(y) = \dim(X_y) = n$ ならば、$y$ は $U_n \to Y$ の像に属する。
Morphisms, Lemma \ref{morphisms-lemma-fppf-open} により、$f(U_n)$ は $Y$ において
開である。したがって、$n_{X/Y}$ が $\geq n$ となる $y$ の開近傍が存在する。
\end{proof}

\begin{lemma}
\label{more-morphisms-lemma-dimension-fibres-proper}
$f : X \to Y$ を固有射とする。$n_{X/Y}$ を、Lemma
\ref{more-morphisms-lemma-base-change-dimension-fibres} で導入した、$f$ のファイバーの次元を与える
$Y$ 上の関数とする。このとき $n_{X/Y}$ は上半連続である。
\end{lemma}

\begin{proof}
$Z_d = \{x \in X \mid \dim_x(X_{f(x)}) > d\}$ とおく。Morphisms, Lemma
\ref{morphisms-lemma-openness-bounded-dimension-fibres} により、$Z_d$ は $X$ の
閉部分集合である。$f$ は固有なので $f(Z_d)$ は閉である。
$y \in f(Z_d) \Leftrightarrow n_{X/Y}(y) > d$ であるから、補題が従う。
\end{proof}

\begin{lemma}
\label{more-morphisms-lemma-dimension-fibres-proper-flat}
$f : X \to Y$ を有限表示な固有平坦射とする。$n_{X/Y}$ を、Lemma
\ref{more-morphisms-lemma-base-change-dimension-fibres} で導入した、$f$ のファイバーの次元を与える
$Y$ 上の関数とする。このとき $n_{X/Y}$ は局所定数である。
\end{lemma}

\begin{proof}
これは Lemmas \ref{more-morphisms-lemma-dimension-fibres-flat} と
\ref{more-morphisms-lemma-dimension-fibres-proper} の直接の帰結である。
\end{proof}

















\section{弱相対 Noether 正規化}
\label{more-morphisms-section-relative-noether-normalization}

\noindent
本節の目標は Lemma \ref{more-morphisms-lemma-weak-relative-noether-normalization} を
証明することである。

\begin{lemma}
\label{more-morphisms-lemma-silly}
$R$ を環とする。$\mathfrak p_1, \ldots, \mathfrak p_r$ を $R$ の素イデアルで、
$i \not = j$ ならば $\mathfrak p_i \not \subset \mathfrak p_j$ であるものとする。
$k_i \subset \kappa(\mathfrak p_i)$ を部分体で、拡大
$\kappa(\mathfrak p_i)/k_i$ が代数的でないものとする。$J \subset R$ を、どの
$\mathfrak p_i$ にも含まれないイデアルとする。このとき元 $x \in J$ が存在して、
各 $i = 1, \ldots, r$ に対し $\kappa(\mathfrak p_i)$ における $x$ の像は
$k_i$ 上超越的である。
\end{lemma}

\begin{proof}
イデアル $J_i = J \mathfrak p_1 \ldots \hat{\mathfrak p}_i \ldots \mathfrak p_r$ は
$\mathfrak p_i$ に含まれない。Algebra, Lemma
\ref{algebra-lemma-product-ideals-in-prime} を参照せよ。したがって
% P05-EN-CH37-0142: the frozen fraction-field identity uses an undefined ring B instead of R; see the sidecar.
$\kappa(\mathfrak p_i) = \text{Frac}(B/\mathfrak p_i)$
の任意の元 $\xi$ は、$a, b \in J_i$ かつ $b \not \in \mathfrak p_i$ として
$\xi = a/b$ の形に書ける。$k_i$ 上超越的な $\xi$ を選ぶと、$a$ または $b$ の
どちらか一方は $\kappa(\mathfrak p_i)$ において $k_i$ 上超越的な元に写る。
したがって各 $i = 1, \ldots, r$ に対し、
$x_i \in J_i = J \mathfrak p_1 \ldots \hat{\mathfrak p}_i \ldots \mathfrak p_r$
であって $\kappa(\mathfrak p_i)$ における像が $k_i$ 上超越的なものが存在する。
このとき $x = x_1 + \ldots + x_r$ が求める元である。
\end{proof}

\begin{lemma}
\label{more-morphisms-lemma-sum-is-ok}
$R \to S$ を有限型環準同型とする。$d \geq 0$ とし、$a, b \in S$ とする。
$x$ を $a$ に写す $R$-代数準同型 $R[x] \to S$ が与える射
$$
f_a : \Spec(S) \longrightarrow \mathbf{A}^1_R
$$
のファイバーの次元が $\leq d$ であると仮定する。このとき、ある $n_0$ が存在して、
$n \geq n_0$ に対し、$x$ を $a^n + b$ に写す $R$-代数準同型 $R[x] \to S$ が与える射
$$
f_{a^n + b} : \Spec(S) \longrightarrow \mathbf{A}^1_R
$$
のファイバーの次元は $\leq d$ である。
\end{lemma}

\begin{proof}
この段落では、$R \to S$ が有限表示である場合に帰着する。実際、$S$ において $A$ と $B$ が
$a$ と $b$ に写るものとして、
% P05-EN-CH37-0143: the frozen presentation places J in a smaller polynomial ring that omits A and B; see the sidecar.
$S = R[A, B, x_1, \ldots, x_n]/J$、$J \subset R[x_1, \ldots, x_n]$ と書く。
$J$ はその有限生成イデアル $J_\lambda \subset J$ の合併である。
$S_\lambda = R[A, B, x_1, \ldots, x_n]/J_\lambda$ とおき、
% P05-EN-CH37-0144: the frozen naming construction omits "by"; see the sidecar.
$A$ と $B$ の像を $a_\lambda, b_\lambda \in S_\lambda$ と表す。すると、ある $\lambda$ に対し
$$
f_{a_\lambda} : \Spec(S_\lambda) \longrightarrow \mathbf{A}^1_R
$$
のファイバーの次元は $\leq d$ である。Limits, Lemma
\ref{limits-lemma-limit-dimension} を参照せよ。そのような $\lambda$ を固定する。
$R \to S_\lambda$, $a_\lambda$, $b_\lambda$ に対して条件を満たす $n_0$ が見つかれば、
同じ $n_0$ が $R \to S$ に対しても条件を満たす。実際、
$f_{a_\lambda^n + b_\lambda} : \Spec(S_\lambda) \to \mathbf{A}^1_R$ のファイバーは
$f_{a^n + b} : \Spec(S) \to \mathbf{A}^1_R$ のファイバーを含む。
これにより、次の段落で扱う場合に帰着される。

\medskip\noindent
$R \to S$ は有限表示であると仮定する。この段落では、$R$ が $\mathbf{Z}$ 上有限型である
場合に帰着する。Algebra, Lemma
\ref{algebra-lemma-limit-module-finite-presentation} により、有向集合 $\Lambda$ と、
余極限が $R \to S$ である $\Lambda$ 上の環準同型の系 $R_\lambda \to S_\lambda$ で、
$\mu \geq \lambda$ に対して $S_\mu = S_\lambda \otimes_{R_\lambda} R_\mu$ となり、
各 $R_\lambda$ と $S_\lambda$ が $\mathbf{Z}$ 上有限型であるものが存在する。
$\lambda_0 \in \Lambda$ と、$a, b \in S$ に写る元
$a_{\lambda_0}, b_{\lambda_0} \in S_{\lambda_0}$ を選ぶ。
$\lambda \geq \lambda_0$ に対し、$a_{\lambda_0}, b_{\lambda_0}$ の像を
$a_\lambda, b_\lambda \in S_\lambda$ と表す。十分大きい $\lambda \geq \lambda_0$ に対して
$$
f_{a_\lambda} : 
\Spec(S_\lambda) \longrightarrow \mathbf{A}^1_{R_\lambda}
$$
のファイバーの次元は $\leq d$ である。Limits, Lemma
\ref{limits-lemma-descend-dimension-d} を参照せよ。そのような $\lambda$ を固定する。
$R_\lambda \to S_\lambda$, $a_\lambda$, $b_\lambda$ に対して条件を満たす $n_0$ が
見つかれば、同じ $n_0$ が $R \to S$ に対しても条件を満たす。実際、Morphisms, Lemma
\ref{morphisms-lemma-dimension-fibre-after-base-change} により、
$f_{a^n + b} : \Spec(S) \to \mathbf{A}^1_R$ の任意のファイバーは
$f_{a_\lambda^n + b_\lambda} : \Spec(S_\lambda) \to \mathbf{A}^1_{R_\lambda}$ の
あるファイバーと同じ次元をもつ。
% P05-EN-CH37-0145: the frozen transition twice says "reduces us the the case"; see the sidecar.
これにより、次の段落で扱う場合に帰着される。

\medskip\noindent
$R$ と $S$ は $\mathbf{Z}$ 上有限型であると仮定する。特に $R$ の次元は有限であり、
$\dim(R)$ に関する帰納法を用いてよい。したがって、$R'$ と $S'$ が $\mathbf{Z}$ 上
有限型で $\dim(R') < \dim(R)$ である、補題と同様のすべての組
$R' \to S'$, $a$, $b$ に対して結果が成り立つと仮定してよい。

\medskip\noindent
主張は $S$ のスペクトルの位相に関するものなので、$S$ は被約であると仮定してよい。
$S^\nu$ を $S$ の正規化とする。$S$ は優秀なので、$S \subset S^\nu$ は有限拡大である。
Algebra, Proposition \ref{algebra-proposition-ubiquity-nagata} と
Morphisms, Lemma \ref{morphisms-lemma-nagata-normalization} を参照せよ。
したがって $\Spec(S^\nu) \to \Spec(S)$ は全射かつ有限である（Algebra, Lemma
\ref{algebra-lemma-integral-overring-surjective}）。ゆえに $R \to S^\nu$ と
$S^\nu$ における $a$, $b$ の像について結果が成り立てば、$R \to S$, $a$, $b$ に
ついても結果が成り立つ。
% P05-EN-CH37-0146: the frozen proof explicitly omits three finite-surjective/product/leading-term details; see the sidecar.
（細部は省略する。）これにより、次の段落で扱う場合に帰着される。

\medskip\noindent
$R$ と $S$ は $\mathbf{Z}$ 上有限型で、$S$ は正規であると仮定する。このとき、ある正規整域
$S_i$ によって $S = S_1 \times \ldots \times S_r$ と書ける。各 $R \to S_i$ と
$S_i$ における $a$, $b$ の像について結果が成り立てば、$R \to S$, $a$, $b$ についても
結果が成り立つ。（細部は省略する。）これにより、次の段落で扱う場合に帰着される。

\medskip\noindent
% P05-EN-CH37-0147: the frozen text twice omits "is" before "a normal domain"; see the sidecar.
$R$ と $S$ は $\mathbf{Z}$ 上有限型で、$S$ は正規整域であると仮定する。
$R$ を $S$ における $R$ の像で置き換えてよい（これは $R$ の次元を増加させない）。
これにより、次の段落で扱う場合に帰着される。

\medskip\noindent
$R \subset S$ は $\mathbf{Z}$ 上有限型で、$S$ は正規整域であると仮定する。射
$$
f_a : \Spec(S) \to \mathbf{A}^1_R
$$
を考える。仮定によれば、$f_a$ のファイバーの次元は $\leq d$ である。
したがって $f : \Spec(S) \to \Spec(R)$ のファイバーの次元は $\leq d + 1$ である
（Morphisms, Lemma \ref{morphisms-lemma-dimension-fibre-at-a-point-additive}）。
次に整スキームの射
$$
\phi : \Spec(S) \to \mathbf{A}^2_R = \Spec(R[x, y])
$$
$x$ を $a$ に、$y$ を $b$ に写す $R$-代数準同型 $R[x, y] \to S$ に対応するものを考える。
考えるべき場合は次の二つである。
\begin{enumerate}
\item $\phi$ は支配的である。
\item $\phi$ は支配的でない。
\end{enumerate}
どちらの場合にも、ある整数 $n_0$ と非空開集合 $V \subset \Spec(R)$ が存在して、
$n \geq n_0$ ならば点 $q \in \mathbf{A}^1_V$ における $f_{a^n + b}$ のファイバーの
次元が $\leq d$ となることを主張する。

\medskip\noindent
場合 (1) における主張を証明する。$f_{a^n + b} = \pi_n \circ \phi$ である。ここで
$$
\pi_n :  \mathbf{A}^2_R \to \mathbf{A}^1_R
$$
は、$x$ を $x^n + y$ に写す $R$-代数準同型 $R[x] \to R[x, y]$ に対応する平坦射である。
$\phi$ は支配的なので、$\phi|_U : U \to \mathbf{A}^2_R$ が平坦となる稠密開集合
$U \subset \Spec(S)$ が存在する（例えば一般平坦性から従う。Morphisms, Proposition
\ref{morphisms-proposition-generic-flatness} を参照せよ）。このとき合成
$$
f_{a^n + b}|_U :
U \xrightarrow{\phi|_U}
\mathbf{A}^2_R
\xrightarrow{\pi_n}
\mathbf{A}^1_R
$$
も平坦である。したがって Morphisms, Lemma
\ref{morphisms-lemma-dimension-fibre-at-a-point-additive} により、この射のファイバーは
$f|_U : U \to \Spec(R)$ のファイバーにおいて少なくとも余次元 $1$ をもつ。
言い換えれば、$f_{a^n + b}|_U$ のファイバーの次元は $\leq d$ である。
一方、$U$ は $\Spec(S)$ において稠密なので、すべての $p \in V$ に対して
$U \cap f^{-1}(p) \subset f^{-1}(p)$ が稠密となる非空開集合
$V \subset \Spec(R)$ が存在する（例えば Lemma
\ref{more-morphisms-lemma-nowhere-dense-generic-fibre} を参照せよ）。したがって
$\dim(f^{-1}(p) \setminus U \cap f^{-1}(p)) \leq d$ であり、主張が従う。
% P05-EN-CH37-0148: the frozen parenthetical has singular/plural disagreement; see the sidecar.
実際、$f_{a^n + b} : \Spec(S) \to \mathbf{A}^1_R$ の任意のファイバーは
$f$ の一つのファイバーに含まれる。

\medskip\noindent
場合 (2) を考える。この場合、$\Im(\phi) \subset V(g)$ となる非零元
$g = \sum c_{ij} x^i y^j$ が $R[x, y]$ に存在する。実際、$g$ は
$\text{Frac}(R)$ 上既約であると仮定してよい。$g \in R[x]$ で、その最高次係数を $c$ と
すると、$V = D(c) \subset \Spec(R)$ 上で $f$ のファイバーの次元はすでに $\leq d$ である。
実際、$f_a$ の像は $V(g) \subset \mathbf{A}^1_R$ に含まれ、後者は $V$ 上有限な
ファイバーをもつ。したがって $g$ は $R[x]$ に含まれないと仮定してよい。
$s \geq 1$ を $y$ の多項式としての $g$ の次数とし、$t$ を $x$ の多項式としての
$\sum c_{is} x^i$ の次数とする。このとき $c_{ts}$ は非零であり、
$$
g(x, -x^n) = (-1)^s c_{ts} x^{t + sn} + l.o.t.
$$
である。ただし $n$ は $x$ の多項式としての $g$ の次数より大きいとする
（細部は省略する）。そのような $n$ に対して $g(x, -x^n)$ は $x$ の非零多項式であり、
% P05-EN-CH37-0149: the frozen nonvanishing condition reverses the coefficient indices c_ts to c_st; see the sidecar.
$c_{st} \not \in \mathfrak p$ であるすべての $\mathfrak p \subset R$ に対し、
$\kappa(\mathfrak p)[x]$ の非零多項式に写る。
% P05-EN-CH37-0150: the frozen proof checks only g(x,-x^n) and omits the uniform arbitrary-fibre substitution; see the sidecar.
したがって $V$ を $\Spec(R)$ の主開集合 $D(c_{ts})$ とすれば主張が成り立つ。

\medskip\noindent
以上により、$\dim(R)$ に関する帰納法を用いることができる。実際、
$V(I) = \Spec(R) \setminus V$ となるイデアル $I \subset R$ をとる。$R$ は整域なので
$\dim(R/I) < \dim(R)$ であることに注意する。$R/I \to S/IS$ と $S/IS$ における
$a$, $b$ の像に対して $\dim(R)$ に関する帰納法から得られる整数を $n'_0$ とする。
このとき $\max(n_0, n'_0)$ が条件を満たす。
\end{proof}

\begin{lemma}
\label{more-morphisms-lemma-weak-relative-noether-normalization}
$R \to S$ を有限型環準同型とする。
% P05-EN-CH37-0151: for the zero ring S all fibres are empty, so the stated maximum d is undefined/nonintegral; see the sidecar.
$d$ を $\Spec(S) \to \Spec(R)$ のファイバーの次元の最大値とする。
このとき準有限環準同型
$R[t_1, \ldots, t_d] \to S$
が存在する。
\end{lemma}

\begin{proof}
この段落では、$R \to S$ が有限表示である場合に帰着する。実際、あるイデアル
$J \subset R[x_1, \ldots, x_n]$ によって $S = R[x_1, \ldots, x_n]/J$ と書く。
$J$ はその有限生成イデアル $J_\lambda \subset J$ の合併である。
$S_\lambda = R[x_1, \ldots, x_n]/J_\lambda$ とおく。すると、ある $\lambda$ に対し
$\Spec(S_\lambda) \to \Spec(R)$ のファイバーの次元は $\leq d$ である。
Limits, Lemma \ref{limits-lemma-limit-dimension} を参照せよ。そのような $\lambda$ を固定する。
準有限な $R[t_1, \ldots, t_d] \to S_\lambda$ が見つかれば、合成
$R[t_1, \ldots, t_d] \to S$ ももちろん準有限である。
これにより、次の段落で扱う場合に帰着される。

\medskip\noindent
$R \to S$ は有限表示であると仮定する。この段落では、$R$ が $\mathbf{Z}$ 上有限型である
場合に帰着する。Algebra, Lemma
\ref{algebra-lemma-limit-module-finite-presentation} により、有向集合 $\Lambda$ と、
余極限が $R \to S$ である $\Lambda$ 上の環準同型の系 $R_\lambda \to S_\lambda$ で、
$\mu \geq \lambda$ に対して $S_\mu = S_\lambda \otimes_{R_\lambda} R_\mu$ となり、
各 $R_\lambda$ と $S_\lambda$ が $\mathbf{Z}$ 上有限型であるものが存在する。
十分大きい $\lambda$ に対して、$\Spec(S_\lambda) \to \Spec(R_\lambda)$ のファイバーの
次元は $\leq d$ である。Limits, Lemma
\ref{limits-lemma-descend-dimension-d} を参照せよ。そのような $\lambda$ を固定する。
準有限環準同型 $R_\lambda[t_1, \ldots, t_d] \to S_\lambda$ が見つかれば、その基底変換
$R[t_1, \ldots, t_d] \to S$ も準有限である（Algebra, Lemma
\ref{algebra-lemma-quasi-finite-base-change}）。これにより、次の段落で扱う場合に帰着される。

\medskip\noindent
$R$ と $S$ は $\mathbf{Z}$ 上有限型であると仮定する。$d = 0$ ならば環準同型は準有限であり、
証明は完了する。$d > 0$ と仮定する。$R$-代数準同型
$R[x] \to S$, $x \mapsto a$ のファイバーの次元が $< d$ となる元 $a \in S$ を
見つける。そうすれば帰納法によって証明が完了する。

\medskip\noindent
$a$ の存在を $\dim(R)$ に関する帰納法で証明する。

\medskip\noindent
$\mathfrak q_1, \ldots, \mathfrak q_r \subset S$ を、$S$ の極小素イデアルのうち
$\dim_{\mathfrak q_i}(S/R) = d$ となるものとする。記法については Algebra, Definition
\ref{algebra-definition-relative-dimension} を参照せよ。$\mathfrak q_i$ は素イデアル
$\mathfrak p_i \subset R$ の上にあるとする。$\mathfrak q_i$ はそのファイバーの一般点なので、
$\text{trdeg}_{\kappa(\mathfrak p_i)}(\kappa(\mathfrak q_i)) = d$ である。例えば
$S \otimes_R \kappa(\mathfrak p_i)$ に Algebra, Lemma
\ref{algebra-lemma-dimension-at-a-point-finite-type-field} を適用せよ。
したがって Lemma \ref{more-morphisms-lemma-silly} により元 $a \in S$ が存在して、
各 $i = 1, \ldots, r$ に対し $\kappa(\mathfrak q_i)$ における $a$ の像は
$\kappa(\mathfrak p_i)$ 上超越的である。射
$$
f_a : \Spec(S) \longrightarrow \mathbf{A}^1_R
$$
を考える。
% P05-EN-CH37-0152: the frozen sentence malformedly says "corresponding ... to mapping"; see the sidecar.
これは $x$ を $a$ に写す $R$-代数準同型 $R[x] \to S$ に対応する。
$U \subset \Spec(S)$ を、ファイバーの次元が $\leq d - 1$ となる開部分集合とする。
Morphisms, Lemma \ref{morphisms-lemma-openness-bounded-dimension-fibres} を参照せよ。
構成により $U$ は $\Spec(S)$ のすべての一般点を含む。特に、
$\Spec(S) \to \Spec(R)$ のすべての一般ファイバーのすべての一般点を $U$ が含むことが
分かる。そのような点は必ず $\Spec(S)$ の一般点だからである。
$T = \Spec(S) \setminus U$ を $\Spec(S)$ の被約閉部分スキームとみなす。
直前の記述と仮定 $\dim(S/R) \leq d$ から、$T \to \Spec(R)$ の一般ファイバーの
次元は $\leq d - 1$ である。したがって Lemma
\ref{more-morphisms-lemma-dimension-in-neighbourhood} を $\Spec(R)$ の一般点の開近傍を作るために
数回適用することにより、$T_V \to V$ のファイバーの次元が $\leq d - 1$ となる
稠密開集合 $V \subset \Spec(R)$ が存在する。ゆえに $q \in \mathbf{A}^1_V$ に対し、
$q$ 上の $f_a$ のファイバーの次元は $\leq d - 1$ である。実際、$f_a|_U$ の
ファイバーと $f_a|_T$ のファイバーの次元をともに評価した。

\medskip\noindent
素イデアル回避により、ある $f \in R$ に対して $V = D(f)$ であると仮定してよい。
すると環準同型 $R_f[x] \to S_f$, $x \mapsto a$ のファイバーの次元は
$\leq d - 1$ である。$a$ を $fa$ で置き換え、$a \in (f)$ と仮定してよい。
$\dim(R)$ に関する帰納法により、
% P05-EN-CH37-0153: the frozen quotient polynomial ring is parenthesized as R/fR[x] rather than (R/fR)[x]; see the sidecar.
$\Spec(S/fS) \to \Spec(R/fR[x])$, $x \mapsto \overline{b}$ のファイバーの次元が
$\leq d - 1$ となる元 $\overline{b} \in S/fS$ が存在する。
$b \in S$ を $\overline{b}$ の持ち上げとする。Lemma \ref{more-morphisms-lemma-sum-is-ok} により、
$a^n + b$ が $R_f \to S_f$ に対してなお条件を満たすような $n > 0$ が存在する。
一方、$S/fS$ における $a^n + b$ の像は $\overline{b}$ であり、証明が完了する。
\end{proof}













\section{Bertini の定理}
\label{more-morphisms-section-bertini}

\noindent
Varieties, Section \ref{varieties-section-bertini} で始めた議論を続ける。
本節では \cite{Jou} の卓越した議論に従い、幾何学的に既約な多様体の
一般超平面切断が幾何学的に既約であることを証明する。

\begin{lemma}
\label{more-morphisms-lemma-amazing-bertini-lemma}
\begin{reference}
\cite{Jou} の 71、72 ページを参照せよ。
\end{reference}
$K/k$ を幾何学的に既約な有限生成体拡大とする。$n \geq 1$ とする。
$g_1, \ldots, g_n \in K$ を、ある $c_1, \ldots, c_n \in k$ が存在して、
$$
x_1, \ldots, x_n, \sum g_ix_i, \sum c_ig_i \in K(x_1, \ldots, x_n)
$$
が $k$ 上代数的独立となるものとする。このとき $K(x_1, \ldots, x_n)$ は
$k(x_1, \ldots, x_n, \sum g_ix_i)$ 上幾何学的に既約である。
\end{lemma}

\begin{proof}
$c_1, \ldots, c_n \in k$ を補題の主張のとおりとする。
$\xi = \sum g_ix_i$、$\delta = \sum c_ig_i$ と書く。$a \in k$ に対し、$K$ 上恒等で
$$
\sigma_a(x_i) = x_i + a c_i
$$
によって与えられる $K(x_1, \ldots, x_n)$ の自己同型 $\sigma_a$ を考える。
$\sigma_a(\xi) = \xi + a \delta$ かつ $\sigma_a(\delta) = \delta$ に注意する。
体の塔
$$
K_0 = k(x_1, \ldots, x_n) \subset
K_1 = K_0(\xi) \subset
K_2 = K_0(\xi, \delta) \subset K(x_1, \ldots, x_n) = \Omega
$$
を考える。$\sigma_a(K_0) = K_0$ かつ $\sigma_a(K_2) = K_2$ に注意する。
$\theta \in \Omega$ は $K_1$ 上分離代数的であるとする。Algebra, Lemma
\ref{algebra-lemma-geometrically-irreducible-separable-elements} により、
$\theta \in K_1$ を示さなければならない。

\medskip\noindent
% P05-EN-CH37-0154: the frozen naming construction omits "by"; see the sidecar.
$\Omega$ における $K_2$ の分離代数閉包を $K'_2$ と表す。
$K'_2/K_2$ は有限（Algebra, Lemma
\ref{algebra-lemma-make-geometrically-irreducible}）かつ分離なので、$K'_2$ と $K_2$ の
間にある体は有限個しかない（Fields, Lemma \ref{fields-lemma-primitive-element}）。
$k$ が無限ならば\footnote{有限体の場合は証明の最後の段落で扱う。}、
$$
K_2(\sigma_{a_1}(\theta)) = K_2(\sigma_{a_2}(\theta))
$$
が $\Omega$ の部分体として成り立つ相異なる元 $a_1, a_2$ が $k$ に存在する。
$\theta_i = \sigma_{a_i}(\theta)$、
$\xi_i = \sigma_{a_i}(\xi) = \xi + a_i \delta$ と書く。
$$
K_2 = K_0(\xi_1, \xi_2)
$$
に注意する。実際、$\xi_i = \xi + a_i \delta$、
$\xi = (a_2 \xi_1 - a_1 \xi_2)/(a_2 - a_1)$、
$\delta = (\xi_1 - \xi_2)/(a_1 - a_2)$ である。
$K_2/K_0$ は次数 $2$ の純超越拡大なので、$\xi_1$ と $\xi_2$ は $K_0$ 上代数的独立である。
$\theta_1$ は $K_0(\xi_1)$ 上代数的なので、$\xi_2$ は
$K_0(\xi_1, \theta_1)$ 上超越的である。

\medskip\noindent
仮定より $K/k$ は幾何学的に既約である。したがって
$K(x_1, \ldots, x_n)/K_0$ は幾何学的に既約である（Algebra, Lemma
\ref{algebra-lemma-geometrically-irreducible-base-change-transcendental}）。
さらに部分拡大 $K_0(\xi_1, \theta_1)/K_0$ も幾何学的に既約である（Algebra, Lemma
\ref{algebra-lemma-subalgebra-geometrically-irreducible}）。$\xi_2$ は
$K_0(\xi_1, \theta_1)$ 上超越的なので、
$K_0(\xi_1, \xi_2, \theta_1)/K_0(\xi_2)$ は幾何学的に既約である（Algebra, Lemma
\ref{algebra-lemma-geometrically-irreducible-add-transcendental}）。
上での $a_1, a_2$ の選び方により
$$
K_0(\xi_1, \xi_2, \theta_1) =
K_2(\sigma_{a_1}(\theta)) =
K_2(\sigma_{a_2}(\theta)) =
K_0(\xi_1, \xi_2, \theta_2)
$$
である。$\theta_2$ は $K_0(\xi_2)$ 上分離代数的なので、再び Algebra, Lemma
\ref{algebra-lemma-geometrically-irreducible-separable-elements} により
$\theta_2 \in K_0(\xi_2)$ である。
% P05-EN-CH37-0155: the frozen conclusion says "givens"; see the sidecar.
この関係に $\sigma_{a_2}^{-1}$ を施すと、求めるとおり
$\theta \in K_0(\xi) = K_1$ を得る。

\medskip\noindent
これで $k$ が無限の場合の証明が完了した。$k$ が有限ならば、変数 $t$ を選び、Algebra,
Lemma \ref{algebra-lemma-geometrically-irreducible-base-change-transcendental}
により幾何学的に既約な拡大 $K(t)/k(t)$ を考える。
% P05-EN-CH37-0156: the frozen finite-field transition says "it is still be true"; see the sidecar.
$x_1, \ldots, x_n, \sum g_ix_i, \sum c_ig_i$
は $K(t, x_1, \ldots, x_n)$ においてなお $k(t)$ 上代数的独立であるから、すでに述べた
議論により $K(t, x_1, \ldots, x_n)$ は
$k(t, x_1, \ldots, x_n, \sum g_ix_i)$ 上幾何学的に既約である。
そこで Algebra, Lemma
\ref{algebra-lemma-geometrically-irreducible-base-change-transcendental} をもう一度用いれば
証明が完了する。
\end{proof}

\begin{lemma}
\label{more-morphisms-lemma-algebraically-independent}
$A$ を体 $k$ 上有限型な整域とする。$n \geq 2$ とする。
$g_1, \ldots, g_n \in A$ を、$V(g_1, g_2)$ が次元 $\dim(A) - 2$ の既約成分を
もつような元とする。このとき、ある $c_1, \ldots, c_n \in k$ が存在して、
$$
x_1, \ldots, x_n, \sum g_ix_i, \sum c_ig_i \in
\text{Frac}(A)(x_1, \ldots, x_n)
$$
は $k$ 上代数的独立である。
\end{lemma}

\begin{proof}
$k$ 上の代数的独立性は、$y = \sum g_ix_i$ および $z = \sum c_ig_i$ によって与えられる射
$$
T = \Spec(A[x_1, \ldots, x_n])
\longrightarrow
\Spec(k[x_1, \ldots, x_n, y, z]) = S
$$
が支配的であることを意味する。$d = \dim(A)$ とおく。$T \to S$ が支配的でなければ、
像の次元は $< n + 2$ であり、したがって任意のファイバーの任意の既約成分の次元は
$> d + n - (n + 2) = d - 2$ である。Varieties, Lemma
\ref{varieties-lemma-dimension-fibres-locally-algebraic} を参照せよ。
$V(g_1, g_2)$ の次元 $d - 2$ の既約成分に含まれ、ほかのどの成分にも含まれない閉点
$u \in V(g_1, g_2)$ を選ぶ。$u$ 上にある $T$ の閉点
% P05-EN-CH37-0157: the frozen coordinate list omits a comma after \ldots; see the sidecar.
$t = (u, 1, 0, \ldots 0)$ を考える。$(c_1, \ldots, c_n) = (0, 1, 0, \ldots, 0)$ とおく。
すると $t$ は $S$ の点 $s = (1, 0, \ldots, 0)$ に写る。$s$ 上の $T \to S$ の
ファイバーは
$$
x_1 - 1, x_2, \ldots, x_n, \sum x_ig_i, g_2
$$
によって切り出され、したがって同値に
$$
x_1 - 1, x_2, \ldots, x_n, g_1, g_2
$$
によって切り出される。$g_1, g_2$ に関する仮定より、この部分スキームは次元
$d - 2$ の既約成分をもつ。
\end{proof}

\begin{lemma}
\label{more-morphisms-lemma-bertini-irreducible}
\begin{reference}
\cite[Theorem 6.3 part 4)]{Jou}
\end{reference}
Varieties, Situation \ref{varieties-situation-family-divisors} において、次を仮定する。
\begin{enumerate}
\item $X$ は $k$ 上有限型である。
\item $X$ は $k$ 上幾何学的に既約である。
\item $v_1, v_2, v_3 \in V$ と $H_{v_2} \cap H_{v_3}$ の既約成分 $Z$ で、
$Z \not \subset H_{v_1}$ かつ $\text{codim}(Z, X) = 2$ となるものが存在する。
\item $\bigcap_{v \in V} H_v$ の任意の既約成分 $Y$ は
$\text{codim}(Y, X) \geq 2$ を満たす。
\end{enumerate}
このとき、一般の $v \in V \otimes_k k'$ に対してスキーム $H_v$ は $k'$ 上
幾何学的に既約である。
\end{lemma}

\begin{proof}
% P05-EN-CH37-0162: the frozen proof explicitly omits the linear-independence argument; see the sidecar.
仮定 (3) が成り立つためには、元 $v_1, v_2, v_3$ は $V$ において $k$-線型独立で
なければならない（細部は省略する）。したがって、これらを最初の $3$ 元として含む
$V$ の基底 $v_1, \ldots, v_r$ を選べる。
$H_{univ} \subset \mathbf{A}^r_k \times_k X$ は「普遍因子」であった。
スキーム論的ファイバーが因子 $H_v$ である射影
$q : H_{univ} \to \mathbf{A}^r_k$ を考える。Lemma
\ref{more-morphisms-lemma-geom-irreducible-generic-fibre} により、$q$ の一般ファイバーが
幾何学的に既約であることを示せば十分である。この証明では $X$ をその被約化で
置き換えてよく、したがって $X$ は $k$ 上有限型な整スキームであると仮定してよい。

\medskip\noindent
$\mathcal{L}|_U \cong \mathcal{O}_U$ となる非空アフィン開集合 $U \subset X$ をとり、
$U = \Spec(A)$ と書く。
% P05-EN-CH37-0158: the frozen naming construction omits "by"; see the sidecar.
同型 $\mathcal{L}|_U \cong \mathcal{O}_U$ により切断 $\psi(v_i)|_U$ に対応する元を
$f_i \in A$ と表す。このとき $H_{univ} \cap (\mathbf{A}^r_k \times_k U)$ は
$$
H_U = \Spec(A[x_1, \ldots, x_r]/(x_1f_1 + \ldots + x_rf_r))
$$
で与えられる。基底の選び方により $f_1$ は零ではありえない。実際、
% P05-EN-CH37-0159: f_1=0 does not imply v_1=0 because psi need not be injective; see the sidecar.
これは $v_1 = 0$、したがって $H_{v_1} = X$ を意味し、仮定 (3) に矛盾する。
ゆえに $\sum x_if_i$ は $A[x_1, \ldots, x_r]$ の非零因子である。
したがって $d = \dim(X) = \dim(A)$ とすれば、$H_U$ の任意の既約成分の次元は
$d + r - 1$ である。$U' = U \cap D(f_1)$ とすれば
$$
H_{U'} =
\Spec(A_{f_1}[x_1, \ldots, x_r]/(x_1f_1 + \ldots + x_rf_r)) \cong
% P05-EN-CH37-0160: the frozen polynomial-variable list omits a comma after \ldots; see the sidecar.
\Spec(A_{f_1}[x_2, \ldots x_r]) =
\mathbf{A}^{r - 1}_k \times_k U'
$$
は既約である。一方、
$$
H_U \setminus H_{U'} =
\Spec(A/(f_1)[x_1, \ldots, x_r]/(x_2f_2 + \ldots + x_rf_r))
$$
の次元は高々 $d + r - 2$ である。実際、$i \not = 1$ に対し、スキーム
$(H_U \setminus H_{U'}) \times_U D(f_i)$ は空（$f_i = 0$ の場合）であるか、上と同じ
議論により $\Spec(A/(f_1))$ の開集合上の相対次元 $r - 1$ のアフィン空間に同型であり、
したがってその次元は高々 $d + r - 2$ である。一方、
$(H_U \setminus H_{U'}) \times_U V(f_2, \ldots, f_r)$ は
$\Spec(A/(f_1, \ldots, f_r))$ 上の相対次元 $r$ のアフィン空間であり、仮定 (4) により
その次元は高々 $d + r - 2$ である。以上から、上のような任意の $U$ に対し $H_U$ は
既約である。したがって $H_{univ}$ は既約である。

\medskip\noindent
したがって Varieties, Lemma
\ref{varieties-lemma-geometrically-irreducible-function-field} により、$H_{univ}$ の
一般点が $\mathbf{A}^r_k$ の一般点上幾何学的に既約であることを示せば十分である。
$X \setminus H_{v_1}$ に含まれ、仮定 (3) で存在が主張された
$H_{v_2} \cap H_{v_3}$ の既約成分 $Z$ と交わる $X$ の非空アフィン開集合
$U = \Spec(A)$ を選ぶ。上の記法により、体拡大
$$
\text{Frac}(A[x_1, \ldots, x_r]/(x_1f_1 + \ldots + x_rf_r))/
k(x_1, \ldots , x_r)
$$
が幾何学的に既約であることを証明しなければならない。$U$ の選び方により $f_1$ は
$A$ において可逆であることに注意する。
% P05-EN-CH37-0161: the frozen naming sentence redundantly says "Set K=... equal to"; see the sidecar.
$K = \text{Frac}(A)$ を $A$ の分数体とする。上と同様に変数 $x_1$ を消去すると、体拡大
$$
K(x_2, \ldots, x_r)/
k(x_2, \ldots, x_r, -\sum\nolimits_{i = 2, \ldots, r} f_1^{-1}f_i x_i)
$$
が幾何学的に既約であることを示せばよい。Lemma
\ref{more-morphisms-lemma-amazing-bertini-lemma} により、ある $c_2, \ldots, c_r \in k$ に対して
$$
x_2, \ldots, x_r, \sum\nolimits_{i = 2, \ldots, r} f_1^{-1}f_i x_i,
\sum\nolimits_{i = 2, \ldots, r} c_if_1^{-1}f_i
$$
が $A[x_2, \ldots, x_r]$ の分数体において $k$ 上代数的独立であることを示せば十分である。
これは Lemma \ref{more-morphisms-lemma-algebraically-independent} と、$Z \cap U$ が
$V(f_1^{-1}f_2, f_1^{-1}f_3) \subset U$ の既約成分であることから従う。
\end{proof}

\begin{remark}
\label{more-morphisms-remark-geometric-proof-bertini-irreducible}
Lemma \ref{more-morphisms-lemma-bertini-irreducible} の特別な場合の「幾何学的」証明を概説する。
$k$ は代数閉体で、$X \subset \mathbf{P}^n_k$ は次元 $\geq 2$ の滑らかで既約な
部分スキームであるとする。このとき $X \cap H$ が滑らかかつ既約となる超平面
$H \subset \mathbf{P}^n_k$ が存在すると主張する。実際、Varieties, Lemma
\ref{varieties-lemma-bertini} により、一般の
$v \in V = kT_0 \oplus \ldots \oplus kT_n$ に対応する超平面切断 $X \cap H_v$ は滑らかである。
一方、Enriques--Severi--Zariski によりスキーム $X \cap H_v$ は連結である。Varieties,
Lemma \ref{varieties-lemma-connectedness-ample-divisor} を参照せよ。
したがって $X \cap H_v$ は滑らかかつ既約である。
\end{remark}













\section{立方体定理}
\label{more-morphisms-section-theorem-cube}

\noindent
次の補題は、その仮定のもとで Picard 関手の対角射が局所閉埋め込みによって
表現可能であることを示す。

\begin{lemma}
\label{more-morphisms-lemma-diagonal-picard-flat-proper}
$f : X \to S$ を有限表示な平坦固有射とする。
$\mathcal{E}$ を有限局所自由 $\mathcal{O}_X$-加群とする。射 $g : T \to S$ に対し、
基底変換図式
$$
\xymatrix{
X_T \ar[d]_p \ar[r]_q & X \ar[d]^f \\
T \ar[r]^g & S
}
$$
を考える。すべての $g : T \to S$ に対し
$\mathcal{O}_T \to p_*\mathcal{O}_{X_T}$ が同型であると仮定する。このとき有限表示な
埋め込み $j : Z \to S$ で、射 $g : T \to S$ が $Z$ を経由することと、
$p^*\mathcal{N} \cong q^*\mathcal{E}$ を満たす有限局所自由
$\mathcal{O}_T$-加群 $\mathcal{N}$ が存在することが同値となるものが存在する。
\end{lemma}

\begin{proof}
仮定 $H^0(X_s, \mathcal{O}_{X_s}) = \kappa(s)$ により、$f$ のファイバー $X_s$ は
連結であることに注意する。したがって $\mathcal{E}$ の階数は各ファイバー上で一定である。
$f$ は開（Morphisms, Lemma \ref{morphisms-lemma-fppf-open}）かつ閉なので、
$S$ には開かつ閉な部分スキームへの分解 $S = \coprod S_r$ が存在して、$S_r$ の逆像上で
$\mathcal{E}$ の階数は一定値 $r$ をとる。ゆえに $\mathcal{E}$ は一定階数 $r$ をもつと
仮定してよい。
% P05-EN-CH37-0163: the frozen naming construction omits "by"; see the sidecar.
階数 $r$ の双対加群を
$\mathcal{E}^\vee = \SheafHom(\mathcal{E}, \mathcal{O}_X)$ と表す。

\medskip\noindent
コホモロジーと基底変換、より正確には Derived Categories of Schemes, Lemma
\ref{perfect-lemma-flat-proper-perfect-direct-image-general} により、
$E = Rf_*\mathcal{E}$ は $S$ の導来圏の完全対象であり、その構成は任意の基底変換と
可換である。同様のことが $E' = Rf_*\mathcal{E}^\vee$ にも成り立つ。次数 $< 0$ に
コホモロジーは存在しないので、ある $b$ に対し $E$ と $E'$ は局所的に $[0, b]$ に
tor 振幅をもつ。任意の $g : T \to S$ に対して
$p_*(q^*\mathcal{E}) = H^0(Lg^*E)$ かつ
$p_*(q^*\mathcal{E}^\vee) = H^0(Lg^*E')$ であることに注意する。
$j : Z \to S$ と $j' : Z' \to S$ を、Derived Categories of Schemes, Lemma
\ref{perfect-lemma-locally-closed-where-H0-locally-free} により $E$ と $E'$ に対して
$a = 0$、$r = r$ として構成される有限表示な埋め込みとする。大まかには、これらは
$H^0(Lj^*E)$ と $H^0((j')^*E')$ が引き戻しと可換な有限局所自由加群であるという
性質によって特徴づけられる。

\medskip\noindent
$g : T \to S$ を射とする。補題のような $\mathcal{N}$ が存在するならば、射影公式
Cohomology, Lemma \ref{cohomology-lemma-projection-formula} を用いて、加群
$$
p_*(q^*\mathcal{E}) \cong
p_*(p^*\mathcal{N}) \cong
\mathcal{N} \otimes_{\mathcal{O}_T} p_*\mathcal{O}_{X_T} \cong
\mathcal{N}\quad\text{同様に}\quad
p_*(q^*\mathcal{E}^\vee) \cong \mathcal{N}^\vee
$$
は階数 $r$ の有限局所自由加群であり、さらに任意の基底変換 $T' \to T$ の後もそうである。
したがってこの場合 $T \to S$ は $j$ と $j'$ の双方を経由する。ゆえに $S$ を
$Z \times_S Z'$ で置き換え、$f_*\mathcal{E}$ と $f_*\mathcal{E}^\vee$ は、その構成が
任意の基底変換と可換な階数 $r$ の有限局所自由 $\mathcal{O}_S$-加群であると
% P05-EN-CH37-0170: the reduction to the scheme-theoretic fibre product explicitly omits its base-change/universal-property detail; see the sidecar.
仮定してよい（細部は省略する）。

\medskip\noindent
% P05-EN-CH37-0164: the frozen conditional says "if g ... be"; see the sidecar.
この状況で $g : T \to S$ が射で、補題のような $\mathcal{N}$ が存在するならば、写像
（次数 $0$ のカップ積）
$$
p_*(q^*\mathcal{E})
\otimes_{\mathcal{O}_T}
p_*(q^*\mathcal{E}^\vee)
\longrightarrow \mathcal{O}_T
$$
は完全ペアリングである。逆に、このカップ積写像が完全ペアリングならば、$T$ 上局所的に
$p_*(q^*\mathcal{E})$ の切断の基底 $\sigma_1, \ldots, \sigma_r$ と
$p_*(q^*\mathcal{E}^\vee)$ の切断の基底 $\tau_1, \ldots, \tau_r$ で、積が
$\sigma_i \tau_j = \delta_{ij}$ を満たすものを選べる。$\sigma_i$ を $X_T$ 上の
$q^*\mathcal{E}$ の切断、$\tau_j$ を $X_T$ 上の $q^*\mathcal{E}^\vee$ の切断とみなすと、
$$
\sigma_1, \ldots, \sigma_r :
\mathcal{O}_{X_T}^{\oplus r}
\longrightarrow
q^*\mathcal{E}
$$
は同型であり、その逆は
$$
\tau_1, \ldots, \tau_r :
q^*\mathcal{E}
\longrightarrow
\mathcal{O}_{X_T}^{\oplus r}
$$
で与えられる。言い換えれば $p^*p_*q^*\mathcal{E} \cong q^*\mathcal{E}$ である。
一方、カップ積が非退化であるという条件は準コンパクト開部分スキーム、すなわち適当な
行列式が非零となる軌跡を切り出す。これで証明が完了する。
\end{proof}

\noindent
上の補題は特に、固有空間の固有平坦族においてベクトル束が各ファイバー上自明ならば、
それはベクトル束の引き戻しであることを示す。これを少し詳しく述べよう。

\begin{lemma}
\label{more-morphisms-lemma-trivial-on-fibres}
$f : X \to S$ を有限表示な平坦固有射で、$f_*\mathcal{O}_X = \mathcal{O}_S$ が成り立ち、
任意の基底変換後もこの等式が成り立つものとする。$\mathcal{E}$ を有限局所自由
$\mathcal{O}_X$-加群とする。次を仮定する。
\begin{enumerate}
\item すべての $s \in S$ に対し
$\mathcal{E}|_{X_s}$ は $\mathcal{O}_{X_s}^{\oplus r_s}$ に同型である。
\item $S$ は被約である。
\end{enumerate}
このとき、ある有限局所自由 $\mathcal{O}_S$-加群 $\mathcal{N}$ に対し
$\mathcal{E} = f^*\mathcal{N}$ である。
\end{lemma}

\begin{proof}
実際、この場合 Lemma \ref{more-morphisms-lemma-diagonal-picard-flat-proper} の局所閉埋め込み
$j : Z \to S$ は全単射であり、したがって閉埋め込みである。しかし $S$ は被約なので、
$j$ は同型である。
\end{proof}

\begin{lemma}
\label{more-morphisms-lemma-triviality-generic-fibre-valuation-ring}
$f : X \to S$ を有限表示な固有平坦射とする。
$\mathcal{L}$ を可逆 $\mathcal{O}_X$-加群とする。次を仮定する。
\begin{enumerate}
\item $S$ は付値環のスペクトルである。
\item $\mathcal{L}$ は $f$ の一般ファイバー $X_\eta$ 上自明である。
\item $f$ の閉ファイバー $X_0$ は整である。
\item $H^0(X_\eta, \mathcal{O}_{X_\eta})$ は $S$ の関数体に等しい。
\end{enumerate}
このとき $\mathcal{L}$ は自明である。
\end{lemma}

\begin{proof}
$S = \Spec(A)$ と書く。まず $A$ が離散付値環の場合に補題を証明する
（これは実際に最もよく用いられる場合である）。$\pi \in A$ を一意化元とする。
自明化切断 $s \in \Gamma(X_\eta, \mathcal{L}_\eta)$ をとる。必要なら $s$ を
$\pi^n s$ で置き換えることにより、$s \in \Gamma(X, \mathcal{L})$ と仮定してよい。
$s|_{X_0} = 0$ ならば、$s$ は $\pi$ で割り切れる。実際、$X_0$ はスキーム論的
ファイバーで、$X$ は $A$ 上平坦である。したがって $s|_{X_0}$ は非零であると仮定してよい。
すると $s$ の零点軌跡 $Z(s)$ は $X_0$ に含まれるが、$X_0$ の一般点を含まない
（$X_0$ は整だからである）。これは $Z(s)$ が $X$ において余次元 $\geq 2$ をもつことを意味し、
求めるとおり $Z(s) = \emptyset$ でない限り Divisors, Lemma
\ref{divisors-lemma-effective-Cartier-codimension-1} に矛盾する。

\medskip\noindent
一般の場合を証明する。付値環 $A$ は連接的なので（Algebra, Example
\ref{algebra-example-valuation-ring-coherent}）、$H^0(X, \mathcal{L})$ は連接
$A$-加群である。Derived Categories of Schemes, Lemma
\ref{perfect-lemma-cohomology-over-coherent-ring} を参照せよ。同値に、
$H^0(X, \mathcal{L})$ は有限表示 $A$-加群である（Algebra, Lemma
\ref{algebra-lemma-coherent-ring}）。$X$ は $A$ 上平坦なので
$H^0(X, \mathcal{L})$ は捩れなしであり、More on Algebra, Lemma
\ref{more-algebra-lemma-generalized-valuation-ring-modules} により、ある $n$ に対し
$H^0(X, \mathcal{L}) = A^{\oplus n}$ である。平坦基底変換（Cohomology of Schemes,
Lemma \ref{coherent-lemma-flat-base-change-cohomology}）により
$$
K = H^0(X_\eta, \mathcal{O}_{X_\eta}) \cong
H^0(X_\eta, \mathcal{L}_\eta) =
H^0(X, \mathcal{L}) \otimes_A K
$$
である。ここで $K$ は $A$ の分数体である。したがって $n = 1$ である。
生成元 $s \in H^0(X, \mathcal{L})$ を選ぶ。$\mathfrak m \subset A$ を極大イデアルとする。
% P05-EN-CH37-0165: when A is a field, the stated colimit over nonzero pi in m is empty and cannot equal kappa; see the sidecar.
このとき $\kappa = A/\mathfrak m = \colim A/\pi$ である。ここで右辺は
非零元 $\pi \in \mathfrak m$ にわたるフィルター付き余極限である（ここで $A$ が
付値環であることを用いる）。したがって
$X_0 = \lim X \times_S \Spec(A/\pi)$ である。$s|_{X_0}$ が零ならば、ある $\pi$ に対し
$s$ の $X \times_S \Spec(A/\pi)$ への制限が零となる。Limits, Lemma
\ref{limits-lemma-descend-section} を参照せよ。しかしこの場合 $\pi^{-1} s$ は
$\mathcal{L}$ の大域切断であり、$s$ が $H^0(X, \mathcal{L})$ の生成元であることに
矛盾する。ゆえに $s|_{X_0}$ は非零である。$Z(s) \subset X$ を $s$ の零点スキームとする。
$s|_{X_0}$ は非零で $X_0$ は整なので、$Z(s)_0 \subset X_0$ は有効 Cartier 因子である。
$f$ は固有で $S$ は局所的なので、$Z(s)$ の各点は $Z(s)_0$ の一点に特殊化する。
したがって Divisors, Lemma \ref{divisors-lemma-fibre-Cartier} の (3) により、
$Z(s)$ は相対有効 Cartier 因子であり、特に $Z(s) \to S$ は平坦である。
% P05-EN-CH37-0166: the frozen conclusion misspells "nonempty"; see the sidecar.
ゆえに $Z(s)$ が非空ならば $Z(s)_\eta$ も非空となり、$s|_{X_\eta}$ が
$\mathcal{L}_\eta$ の自明化であることに矛盾する。したがって求めるとおり
$Z(s) = \emptyset$ である。
\end{proof}

\begin{lemma}
\label{more-morphisms-lemma-get-a-closed}
$f : X \to S$ と $\mathcal{E}$ を Lemma
\ref{more-morphisms-lemma-diagonal-picard-flat-proper} のとおりとし、さらに $\mathcal{E}$ は可逆
$\mathcal{O}_X$-加群であると仮定する。さらに $f$ の幾何ファイバーが整ならば、
$Z$ は $S$ において閉である。
\end{lemma}

\begin{proof}
$j : Z \to S$ は有限表示なので、次を示せば十分である。分数体を $K$ とする付値環
$A$ と、$g(\Spec(K)) \in j(Z)$ を満たす任意の射 $g : \Spec(A) \to S$ に対して
$g(\Spec(A)) \subset j(Z)$ が成り立つ。Morphisms, Lemma
\ref{morphisms-lemma-reach-points-scheme-theoretic-image} を参照せよ。これは Lemma
\ref{more-morphisms-lemma-triviality-generic-fibre-valuation-ring} と、Lemma
\ref{more-morphisms-lemma-diagonal-picard-flat-proper} における $j : Z \to S$ の特徴づけから従う。
\end{proof}

\begin{lemma}
\label{more-morphisms-lemma-H1-O-picard-flat-proper}
スキームの可換図式
$$
\xymatrix{
X' \ar[rr] \ar[dr]_{f'} & & X \ar[dl]^f \\
& S
}
$$
を考える。$f' : X' \to S$ と $f : X \to S$ は Lemma
\ref{more-morphisms-lemma-diagonal-picard-flat-proper} の仮定を満たすとする。
$\mathcal{L}$ を可逆 $\mathcal{O}_X$-加群とし、$\mathcal{L}'$ をその $X'$ への引き戻しとする。
$Z \subset S$ と $Z' \subset S$ を、それぞれ $(f, \mathcal{L})$ と
$(f', \mathcal{L}')$ に対して Lemma \ref{more-morphisms-lemma-diagonal-picard-flat-proper} で
構成される局所閉部分スキームとし、$Z \subset Z'$ とする。$s \in Z$ で
$$
H^1(X_s, \mathcal{O}) \longrightarrow H^1(X'_s, \mathcal{O})
$$
が単射ならば、$s$ のある開近傍 $U$ に対し $Z \cap U = Z' \cap U$ である。
\end{lemma}

\begin{proof}
$S$ を $Z'$ で置き換えてよい。$S$ を $s$ のアフィン開近傍へ縮小した後、
$\mathcal{L}' = \mathcal{O}_{X'}$ と仮定してよい。
$E = Rf_*\mathcal{L}$、$E' = Rf'_*\mathcal{L}' = Rf'_*\mathcal{O}_{X'}$ とおく。
これらは、その構成が任意の基底変換と可換な完全複体である（Derived Categories of Schemes,
Lemma \ref{perfect-lemma-flat-proper-perfect-direct-image-general}）。
特に
$$
E \otimes_{\mathcal{O}_S}^\mathbf{L} \kappa(s) =
R\Gamma(X_s, \mathcal{L}_s) = R\Gamma(X_s, \mathcal{O}_{X_s})
$$
である。第二の等式は $s \in Z$ だからである。
$h_i = \dim_{\kappa(s)} H^i(X_s, \mathcal{O}_{X_s})$ とおく。
$S$ を縮小した後、$E$ を複体
$$
\mathcal{O}_S \to \mathcal{O}_S^{\oplus h_1} \to
\mathcal{O}_S^{\oplus h_2} \to \ldots
$$
によって表せる。More on Algebra, Lemma
\ref{more-algebra-lemma-lift-perfect-from-residue-field} を参照せよ（厳密には Derived
Categories of Schemes, Lemmas \ref{perfect-lemma-affine-compare-bounded} と
\ref{perfect-lemma-perfect-affine} も用いる）。同様に、$E'$ は複体
$$
\mathcal{O}_S \to \mathcal{O}_S^{\oplus h'_1} \to
\mathcal{O}_S^{\oplus h'_2} \to \ldots
$$
によって表されると仮定してよい。ここで
$h'_i = \dim_{\kappa(s)} H^i(X'_s, \mathcal{O}_{X'_s})$ である。
コホモロジーの関手性により、基底変換と可換な $D(\mathcal{O}_S)$ の写像
$$
E \longrightarrow E'
$$
が存在する。$E$ を表す複体は有限自由加群からなる有限複体で、$S$ はアフィンなので、
与えられた写像 $E \to E'$ を表す複体の写像
$$
\xymatrix{
\mathcal{O}_S \ar[r]_d \ar[d]_a &
\mathcal{O}_S^{\oplus h_1} \ar[r] \ar[d]_b &
\mathcal{O}_S^{\oplus h_2} \ar[r] \ar[d]_c & \ldots \\
\mathcal{O}_S \ar[r]^{d'} &
\mathcal{O}_S^{\oplus h'_1} \ar[r] &
\mathcal{O}_S^{\oplus h'_2} \ar[r] & \ldots
}
$$
を選べる。$s \in Z$ なので、$\mathcal{L}_s$ の自明化切断は
$\mathcal{L}'_s = \mathcal{O}_{X'_s}$ の自明化切断に引き戻される。したがって
$a \otimes \kappa(s)$ は同型であり、$S$ を縮小した後 $a$ は同型である。
最後に $H^1(X_s, \mathcal{O}) \to H^1(X'_s, \mathcal{O})$ が単射であるという仮定を
用いると、$b$ を定める行列には $\kappa(s)$ の非零元に写る
$h_1 \times h_1$ 小行列式が存在する。したがって $S$ を縮小した後、$b$ は単射であると
仮定してよい。一方、$\mathcal{L}' = \mathcal{O}_{X'}$ なので $d' = 0$ である。
ゆえに $d = 0$ である。こうして $\mathcal{L}_s$ の自明化切断が $X$ 上の
$\mathcal{L}$ の切断へ持ち上がることが分かる。
% P05-EN-CH37-0167: this proof and its multi-map variant explicitly omit the final properness/topological argument; see the sidecar.
直接的な位相的議論（省略する）により、必要なら $S$ をさらに少し縮小した後、
$\mathcal{L}$ は自明である。
\end{proof}

\begin{lemma}
\label{more-morphisms-lemma-H1-O-multiple-picard-flat-proper}
スキームの $n$ 個の可換図式
$$
\xymatrix{
X_i \ar[rr] \ar[dr]_{f_i} & & X \ar[dl]^f \\
& S
}
$$
を考える。$f_i : X_i \to S$ と $f : X \to S$ は Lemma
\ref{more-morphisms-lemma-diagonal-picard-flat-proper} の仮定を満たすとする。
$\mathcal{L}$ を可逆 $\mathcal{O}_X$-加群とし、$\mathcal{L}_i$ をその $X_i$ への
引き戻しとする。$Z \subset S$ と $Z_i \subset S$ を、それぞれ $(f, \mathcal{L})$ と
$(f_i, \mathcal{L}_i)$ に対して Lemma \ref{more-morphisms-lemma-diagonal-picard-flat-proper} で
構成される局所閉部分スキームとし、$Z \subset \bigcap_{i = 1, \ldots, n} Z_i$ とする。
$s \in Z$ で
$$
H^1(X_s, \mathcal{O}) \longrightarrow
\bigoplus\nolimits_{i = 1, \ldots, n} H^1(X_{i, s}, \mathcal{O})
$$
が単射ならば、$s$ のある開近傍 $U$ に対し
$Z \cap U = (\bigcap_{i = 1, \ldots, n} Z_i) \cap U$ である。ここで右辺の交わりは
スキーム論的交わりである。
\end{lemma}

\begin{proof}
この補題は Lemma \ref{more-morphisms-lemma-H1-O-picard-flat-proper} の変形であり、まずその証明を
読むことを強く勧める。本証明は基本的に、わずかな変更を加えたその証明の複製である。
$Z$ と $Z_i$ のスキーム値点の記述から、スキーム論的交わりをとって
$Z \subset \bigcap_{i = 1, \ldots, n} Z_i$ である。したがって $S$ をスキーム論的交わり
$\bigcap_{i = 1, \ldots, n} Z_i$ で置き換えてよい。$S$ を $s$ のアフィン開近傍へ
縮小した後、$i = 1, \ldots, n$ に対し $\mathcal{L}_i = \mathcal{O}_{X_i}$ と
仮定してよい。$E = Rf_*\mathcal{L}$、
$E_i = Rf_{i, *}\mathcal{L}_i = Rf_{i, *}\mathcal{O}_{X_i}$ とおく。
これらは、その構成が任意の基底変換と可換な完全複体である（Derived Categories of Schemes,
Lemma \ref{perfect-lemma-flat-proper-perfect-direct-image-general}）。
特に
$$
E \otimes_{\mathcal{O}_S}^\mathbf{L} \kappa(s) =
R\Gamma(X_s, \mathcal{L}_s) = R\Gamma(X_s, \mathcal{O}_{X_s})
$$
である。第2の等号は $s \in Z$ だからである。
$h_j = \dim_{\kappa(s)} H^j(X_s, \mathcal{O}_{X_s})$ とおく。
$S$ を縮小した後、$E$ を複体
$$
\mathcal{O}_S \to \mathcal{O}_S^{\oplus h_1} \to
\mathcal{O}_S^{\oplus h_2} \to \ldots
$$
で表せる。More on Algebra, Lemma
\ref{more-algebra-lemma-lift-perfect-from-residue-field}
を参照されたい（厳密には Derived Categories of Schemes, Lemmas
\ref{perfect-lemma-affine-compare-bounded} と
\ref{perfect-lemma-perfect-affine} も用いる）。同様に、$E_i$ は複体
$$
\mathcal{O}_S \to \mathcal{O}_S^{\oplus h_{i, 1}} \to
\mathcal{O}_S^{\oplus h_{i, 2}} \to \ldots
$$
で表されると仮定してよい。ここで
$h_{i, j} = \dim_{\kappa(s)} H^j(X_{i, s}, \mathcal{O}_{X_{i, s}})$ である。
コホモロジーの関手性により、射
$$
E \longrightarrow E_i
$$
を $D(\mathcal{O}_S)$ において得、その構成は基底変換と可換である。
$E$ を表す複体は有限自由加群からなる有限複体であり、$S$ はアフィンなので、複体の射
$$
\xymatrix{
\mathcal{O}_S \ar[r]_d \ar[d]_{a_i} &
\mathcal{O}_S^{\oplus h_1} \ar[r] \ar[d]_{b_i} &
\mathcal{O}_S^{\oplus h_2} \ar[r] \ar[d]_{c_i} & \ldots \\
\mathcal{O}_S \ar[r]^{d_i} &
\mathcal{O}_S^{\oplus h_{i, 1}} \ar[r] &
\mathcal{O}_S^{\oplus h_{i, 2}} \ar[r] & \ldots
}
$$
で、与えられた射 $E \to E_i$ を表すものを選べる。$s \in Z$ なので、
$\mathcal{L}_s$ の自明化切断は
$\mathcal{L}_{i, s} = \mathcal{O}_{X_{i, s}}$ の自明化切断へ引き戻される。したがって
$a_i \otimes \kappa(s)$ は同型であり、$S$ を縮小した後、$a_i$ は同型であると
仮定してよい。最後に、
$H^1(X_s, \mathcal{O}) \to
\bigoplus_{i = 1, \ldots, n} H^1(X_{i, s}, \mathcal{O})$
が単射であるという仮定から、$\oplus b_i$ を定める行列には、その $\kappa(s)$ における
像が非零となる $h_1 \times h_1$ 小行列式が存在する。ゆえに $S$ を縮小した後、
% P05-EN-CH37-0168: frozen direct-sum powers omit \oplus, and the following index range omits a comma; see the sidecar.
$(b_1, \ldots, b_n) : \mathcal{O}_S^{h_1}
\to \bigoplus_{i = 1, \ldots, n} \mathcal{O}_S^{h_{i, 1}}$
は単射であると仮定してよい。しかし $\mathcal{L}_i = \mathcal{O}_{X_i}$ なので、
$i = 1, \ldots n$ に対して $d_i = 0$ である。したがって
$(b_1, \ldots, b_n) \circ d = (\oplus d_i) \circ (a_1, \ldots, a_n)$
より $d = 0$ である。こうして $\mathcal{L}_s$ の自明化切断が $X$ 上の
$\mathcal{L}$ の切断へ持ち上がることが分かる。
% P05-EN-CH37-0167: the same final properness/topological argument is explicitly omitted here; see the sidecar.
直接的な位相的議論（省略する）により、必要なら $S$ をさらに少し縮小した後、
$\mathcal{L}$ は自明である。
\end{proof}

\begin{lemma}
\label{more-morphisms-lemma-pic-of-product}
$f : X \to S$ と $g : Y \to S$ を、Lemma
\ref{more-morphisms-lemma-diagonal-picard-flat-proper} の仮定を満たすスキームの射とする。
$\sigma : S \to X$ と $\tau : S \to Y$ を、それぞれ $f$ と $g$ の切断とする。
$s \in S$ とする。$\mathcal{L}$ を $X \times_S Y$ 上の可逆層とする。
$X$ 上の $(1 \times \tau)^*\mathcal{L}$、$Y$ 上の
$(\sigma \times 1)^*\mathcal{L}$、および $\mathcal{L}|_{(X \times_S Y)_s}$ が
自明ならば、$s$ のある開近傍 $U$ が存在して、$\mathcal{L}$ は
$(X \times_S Y)_U$ 上で自明である。
\end{lemma}

\begin{proof}
K\"unneth の公式（Varieties, Lemma \ref{varieties-lemma-kunneth}）により、射
$$
H^1(X_s \times_{\Spec(\kappa(s))} Y_s, \mathcal{O}) \to
H^1(X_s, \mathcal{O}) \oplus H^1(Y_s, \mathcal{O})
$$
は単射である。したがって2つの射
$$
1 \times \tau : X \to X \times_S Y
\quad\text{および}\quad
\sigma \times 1 : Y \to X \times_S Y
$$
に Lemma \ref{more-morphisms-lemma-H1-O-multiple-picard-flat-proper} を適用すれば結論を得る。
\end{proof}

\begin{theorem}[立方体定理]
\label{more-morphisms-theorem-of-the-cube}
$S$ をスキームとし、$X$、$Y$、$Z$ を $S$ 上のスキームとする。
$x : S \to X$ と $y : S \to Y$ を構造射の切断とする。
$\mathcal{L}$ を $X \times_S Y \times_S Z$ 上の可逆加群とする。次を仮定する。
\begin{enumerate}
\item $X \to S$ と $Y \to S$ は、幾何学的に整なファイバーをもつ有限表示の
平坦かつ固有な射である。
\item $\mathcal{L}$ の
$x \times \text{id}_Y \times \text{id}_Z$ と
$\text{id}_X \times y \times \text{id}_Z$
による引き戻しは、それぞれ $Y \times_S Z$ と $X \times_S Z$ 上で自明である。
\item ある点 $z \in Z$ が存在し、$\mathcal{L}$ の
$X \times_S Y \times_S z$ への制限は自明である。
\item $Z$ は連結である。
\end{enumerate}
このとき $\mathcal{L}$ は自明である。
\end{theorem}

\noindent
しばしば用いられる特別な場合は次のとおりである。
$k$ を体とし、$X, Y, Z$ を $k$-有理点 $x, y, z$ をもつ多様体とする。
$\mathcal{L}$ を $X \times Y \times Z$ 上の可逆加群とする。次を仮定する。
\begin{enumerate}
\item $\mathcal{L}$ は
$x \times Y \times Z$、$X \times y \times Z$、$X \times Y \times z$ 上で自明である。
\item $X$ と $Y$ は $k$ 上幾何学的に整かつ固有である。
\end{enumerate}
このとき $\mathcal{L}$ は自明である。

\begin{proof}
% P05-EN-CH37-0169: the frozen sentence lacks the verb "has" before geometrically integral fibres; see the sidecar.
射 $X \times_S Y \to S$ は、幾何学的に整なファイバーをもつ有限表示の平坦かつ
固有な射である（ファイバーに関する主張については Varieties, Lemmas
\ref{varieties-lemma-geometrically-integral},
\ref{varieties-lemma-bijection-irreducible-components},
および \ref{varieties-lemma-geometrically-reduced-any-base-change} を参照されたい）。
Derived Categories of Schemes, Lemma
\ref{perfect-lemma-proper-flat-geom-red-connected}
により、$X \to S$、$Y \to S$、または $X \times_S Y \to S$ による構造層の
前進は $S$ の構造層であり、このことは任意の基底変換後にも成り立つ。したがって射
$$
p : X \times_S Y \times_S Z \longrightarrow Z
$$
と可逆加群 $\mathcal{L}$ に Lemma \ref{more-morphisms-lemma-diagonal-picard-flat-proper} を適用し、
``普遍的'' な局所閉部分スキーム $Z' \subset Z$ を得る。これは
$\mathcal{L}|_{X \times_S Y \times_S Z'}$ が $Z'$ 上のある可逆加群
$\mathcal{N}$ の引き戻しとなるものである。$z$ の存在から $Z'$ は空でない。
Lemma \ref{more-morphisms-lemma-get-a-closed} により、$Z' \subset Z$ は閉部分スキームである。
$z' \in Z'$ を点とする。$p$ は積射
$$
(X \times_S Z) \times_Z (Y \times_S Z) \longrightarrow Z
$$
と書ける。したがって射 $p$、点 $z'$、および切断
$\sigma : Z \to X \times_S Z$ と $\tau : Z \to Y \times_S Z$
（それぞれ $x$ と $y$ により与えられる）に Lemma
\ref{more-morphisms-lemma-pic-of-product} を適用できる。ゆえに $Z'$ は開である。
したがって $Z' = Z$ であり、$Z$ 上のある可逆加群 $\mathcal{N}$ に対して
$\mathcal{L} = p^*\mathcal{N}$ である。射
$x \times y \times \text{id}_Z : Z  \to X \times_S Y \times_S Z$
で引き戻すと、一方では $\mathcal{N}$ を得、他方では仮定 (2) により自明な可逆加群を得る。
ゆえに $\mathcal{N} = \mathcal{O}_Z$ であり、証明は完了する。
\end{proof}









\section{極限論法}
\label{more-morphisms-section-limits}

\noindent
スキームの極限とネーター近似に関するいくつかの補題を与える。
主としてアフィンの場合に限って論じる。これらの補題の一部は、極限に関する章の補題の
特別な場合である。

\begin{lemma}
\label{more-morphisms-lemma-Noetherian-approximation}
$f : X \to S$ を有限表示なアフィンスキームの射とする。このときデカルト図式
$$
\xymatrix{
X_0 \ar[d]_{f_0} & X \ar[l]^g \ar[d]^f \\
S_0 & S \ar[l]
}
$$
で、次を満たすものが存在する。
\begin{enumerate}
\item $X_0$ と $S_0$ はアフィンスキームである。
\item $S_0$ は $\mathbf{Z}$ 上有限型である。
\item $f_0$ は有限型である。
\end{enumerate}
\end{lemma}

\begin{proof}
$S = \Spec(A)$、$X = \Spec(B)$ と書く。$f$ は有限表示なので、$B$ は
有限表示 $A$-代数である。Morphisms, Lemma
\ref{morphisms-lemma-locally-finite-presentation-characterize} を参照されたい。
したがって補題は Algebra, Lemma
\ref{algebra-lemma-limit-module-finite-presentation} から従う。
\end{proof}

\begin{lemma}
\label{more-morphisms-lemma-Noetherian-approximation-module}
$f : X \to S$ を有限表示なアフィンスキームの射とする。$\mathcal{F}$ を有限表示な
準連接 $\mathcal{O}_X$-加群とする。このとき Lemma
\ref{more-morphisms-lemma-Noetherian-approximation} のような図式で、$g^*\mathcal{F}_0 = \mathcal{F}$ を
満たす連接 $\mathcal{O}_{X_0}$-加群 $\mathcal{F}_0$ が存在するものがある。
\end{lemma}

\begin{proof}
$S = \Spec(A)$、$X = \Spec(B)$、$\mathcal{F} = \widetilde{M}$ と書く。
$f$ は有限表示なので、$B$ は有限表示 $A$-代数である。Morphisms, Lemma
\ref{morphisms-lemma-locally-finite-presentation-characterize} を参照されたい。
$\mathcal{F}$ は $\mathcal{O}_X$ 上有限表示なので、$M$ は有限表示 $B$-加群である。
Properties, Lemma \ref{properties-lemma-finite-presentation-module} を参照されたい。
したがって補題は Algebra, Lemma
\ref{algebra-lemma-limit-module-finite-presentation} から従う。
\end{proof}

\begin{lemma}
\label{more-morphisms-lemma-Noetherian-approximation-flat-module}
$f : X \to S$ を有限表示なアフィンスキームの射とする。$\mathcal{F}$ を有限表示な
準連接 $\mathcal{O}_X$-加群で、$S$ 上平坦なものとする。このとき Lemma
\ref{more-morphisms-lemma-Noetherian-approximation-module} のような図式と層 $\mathcal{F}_0$ を、
さらに $\mathcal{F}_0$ が $S_0$ 上平坦となるように選べる。
\end{lemma}

\begin{proof}
$S = \Spec(A)$、$X = \Spec(B)$、$\mathcal{F} = \widetilde{M}$ と書く。
$f$ は有限表示なので、$B$ は有限表示 $A$-代数である。Morphisms, Lemma
\ref{morphisms-lemma-locally-finite-presentation-characterize} を参照されたい。
$\mathcal{F}$ は $\mathcal{O}_X$ 上有限表示なので、$M$ は有限表示 $B$-加群である。
Properties, Lemma \ref{properties-lemma-finite-presentation-module} を参照されたい。
$\mathcal{F}$ は $S$ 上平坦なので、$M$ は $A$ 上平坦である。Morphisms, Lemma
\ref{morphisms-lemma-flat-module-characterize} を参照されたい。
したがって補題は Algebra, Lemma
\ref{algebra-lemma-flat-finite-presentation-limit-flat} から従う。
\end{proof}

\begin{lemma}
\label{more-morphisms-lemma-Noetherian-approximation-flat}
$f : X \to S$ を有限表示かつ平坦なアフィンスキームの射とする。このとき Lemma
\ref{more-morphisms-lemma-Noetherian-approximation} のような図式で、さらに $f_0$ が平坦となるものが存在する。
\end{lemma}

\begin{proof}
これは Lemma \ref{more-morphisms-lemma-Noetherian-approximation-flat-module} の特別な場合である。
\end{proof}

\begin{lemma}
\label{more-morphisms-lemma-Noetherian-approximation-smooth}
$f : X \to S$ を滑らかなアフィンスキームの射とする。このとき Lemma
\ref{more-morphisms-lemma-Noetherian-approximation} のような図式で、さらに $f_0$ が滑らかとなるものが存在する。
\end{lemma}

\begin{proof}
$S = \Spec(A)$、$X = \Spec(B)$ と書く。$f$ は滑らかなので、$B$ は滑らかな
$A$-代数である。Morphisms, Lemma \ref{morphisms-lemma-smooth-characterize} を
参照されたい。したがって補題は Algebra, Lemma
\ref{algebra-lemma-finite-presentation-fs-Noetherian} から従う。
\end{proof}

\begin{lemma}
\label{more-morphisms-lemma-Noetherian-approximation-geometrically-reduced}
$f : X \to S$ を、幾何学的に被約なファイバーをもつ有限表示なアフィンスキームの
射とする。このとき Lemma \ref{more-morphisms-lemma-Noetherian-approximation} のような図式で、
さらに $f_0$ が幾何学的に被約なファイバーをもつものが存在する。
\end{lemma}

\begin{proof}
Lemma \ref{more-morphisms-lemma-Noetherian-approximation} を適用してデカルト図式
$$
\xymatrix{
X_0 \ar[d]_{f_0} & X \ar[l]^g \ar[d]^f \\
S_0 & S \ar[l]_h
}
$$
を得る。ここでこれはアフィンスキームの図式であり、$X_0 \to S_0$ は
$\mathbf{Z}$ 上有限型なスキームの有限型射である。Lemma
\ref{more-morphisms-lemma-geometrically-reduced-constructible} により、$f_0$ のファイバーが
幾何学的に被約となる点の集合 $E \subset S_0$ は構成可能部分集合である。Lemma
\ref{more-morphisms-lemma-base-change-fibres-geometrically-reduced} により $h(S) \subset E$ である。
$S_0 = \Spec(A_0)$、$S = \Spec(A)$ と書く。有限型 $A_0$-代数の
有向余極限として $A = \colim_i A_i$ と書く。Limits, Lemma
\ref{limits-lemma-limit-contained-in-constructible} により、ある $i$ に対して
$\Spec(A_i) \to S_0$ の像は $E$ に含まれる。$S_0$ を $\Spec(A_i)$ で、
$X_0$ を $X_0 \times_{S_0} \Spec(A_i)$ で置き換えると、$f_0$ のすべての
ファイバーは幾何学的に被約である。
\end{proof}

\begin{lemma}
\label{more-morphisms-lemma-Noetherian-approximation-geometrically-irreducible}
$f : X \to S$ を、幾何学的に既約なファイバーをもつ有限表示なアフィンスキームの
射とする。このとき Lemma \ref{more-morphisms-lemma-Noetherian-approximation} のような図式で、
さらに $f_0$ が幾何学的に既約なファイバーをもつものが存在する。
\end{lemma}

\begin{proof}
Lemma \ref{more-morphisms-lemma-Noetherian-approximation} を適用してデカルト図式
$$
\xymatrix{
X_0 \ar[d]_{f_0} & X \ar[l]^g \ar[d]^f \\
S_0 & S \ar[l]_h
}
$$
を得る。ここでこれはアフィンスキームの図式であり、$X_0 \to S_0$ は
$\mathbf{Z}$ 上有限型なスキームの有限型射である。Lemma
\ref{more-morphisms-lemma-nr-geom-irreducible-components-constructible} により、$f_0$ の
ファイバーが幾何学的に既約となる点の集合 $E \subset S_0$ は構成可能部分集合である。
Lemma \ref{more-morphisms-lemma-base-change-fibres-geometrically-irreducible} により
$h(S) \subset E$ である。$S_0 = \Spec(A_0)$、$S = \Spec(A)$ と書く。
有限型 $A_0$-代数の有向余極限として $A = \colim_i A_i$ と書く。Limits, Lemma
\ref{limits-lemma-limit-contained-in-constructible} により、ある $i$ に対して
$\Spec(A_i) \to S_0$ の像は $E$ に含まれる。$S_0$ を $\Spec(A_i)$ で、
$X_0$ を $X_0 \times_{S_0} \Spec(A_i)$ で置き換えると、$f_0$ のすべての
ファイバーは幾何学的に既約である。
\end{proof}

\begin{lemma}
\label{more-morphisms-lemma-Noetherian-approximation-geometrically-connected}
$f : X \to S$ を、幾何学的に連結なファイバーをもつ有限表示なアフィンスキームの
射とする。このとき Lemma \ref{more-morphisms-lemma-Noetherian-approximation} のような図式で、
さらに $f_0$ が幾何学的に連結なファイバーをもつものが存在する。
\end{lemma}

\begin{proof}
Lemma \ref{more-morphisms-lemma-Noetherian-approximation} を適用してデカルト図式
$$
\xymatrix{
X_0 \ar[d]_{f_0} & X \ar[l]^g \ar[d]^f \\
S_0 & S \ar[l]_h
}
$$
を得る。ここでこれはアフィンスキームの図式であり、$X_0 \to S_0$ は
$\mathbf{Z}$ 上有限型なスキームの有限型射である。Lemma
\ref{more-morphisms-lemma-nr-geom-connected-components-constructible} により、$f_0$ の
ファイバーが幾何学的に連結となる点の集合 $E \subset S_0$ は構成可能部分集合である。
Lemma \ref{more-morphisms-lemma-base-change-fibres-geometrically-connected} により
$h(S) \subset E$ である。$S_0 = \Spec(A_0)$、$S = \Spec(A)$ と書く。
有限型 $A_0$-代数の有向余極限として $A = \colim_i A_i$ と書く。Limits, Lemma
\ref{limits-lemma-limit-contained-in-constructible} により、ある $i$ に対して
$\Spec(A_i) \to S_0$ の像は $E$ に含まれる。$S_0$ を $\Spec(A_i)$ で、
$X_0$ を $X_0 \times_{S_0} \Spec(A_i)$ で置き換えると、$f_0$ のすべての
ファイバーは幾何学的に連結である。
\end{proof}

\begin{lemma}
\label{more-morphisms-lemma-Noetherian-approximation-dimension-d}
$d \geq 0$ を整数とする。$f : X \to S$ を、すべてのファイバーの次元が $d$ である
有限表示なアフィンスキームの射とする。このとき Lemma
\ref{more-morphisms-lemma-Noetherian-approximation} のような図式で、さらに $f_0$ のすべての
ファイバーの次元が $d$ となるものが存在する。
\end{lemma}

\begin{proof}
Lemma \ref{more-morphisms-lemma-Noetherian-approximation} を適用してデカルト図式
$$
\xymatrix{
X_0 \ar[d]_{f_0} & X \ar[l]^g \ar[d]^f \\
S_0 & S \ar[l]_h
}
$$
を得る。ここでこれはアフィンスキームの図式であり、$X_0 \to S_0$ は
$\mathbf{Z}$ 上有限型なスキームの有限型射である。Lemma
\ref{more-morphisms-lemma-dimension-fibres-constructible} により、$f_0$ のファイバーの次元が
$d$ となる点の集合 $E \subset S_0$ は構成可能部分集合である。Lemma
\ref{more-morphisms-lemma-base-change-dimension-fibres} により $h(S) \subset E$ である。
$S_0 = \Spec(A_0)$、$S = \Spec(A)$ と書く。有限型 $A_0$-代数の
有向余極限として $A = \colim_i A_i$ と書く。Limits, Lemma
\ref{limits-lemma-limit-contained-in-constructible} により、ある $i$ に対して
$\Spec(A_i) \to S_0$ の像は $E$ に含まれる。$S_0$ を $\Spec(A_i)$ で、
$X_0$ を $X_0 \times_{S_0} \Spec(A_i)$ で置き換えると、$f_0$ のすべての
ファイバーの次元は $d$ である。
\end{proof}

\begin{lemma}
\label{more-morphisms-lemma-Noetherian-approximation-standard-syntomic}
$f : X \to S$ を標準 syntomic なアフィンスキームの射とする（Morphisms,
Definition \ref{morphisms-definition-syntomic} を参照されたい）。このとき Lemma
\ref{more-morphisms-lemma-Noetherian-approximation} のような図式で、さらに $f_0$ が標準
syntomic となるものが存在する。
\end{lemma}

\begin{proof}
この補題は Algebra, Lemma
\ref{algebra-lemma-relative-global-complete-intersection-Noetherian} の複製である。
\end{proof}

\begin{lemma}
\label{more-morphisms-lemma-Noetherian-approximation-combine}
（ネーター近似と性質の併合。）
$P$、$Q$ を基底変換で安定なスキームの射の性質とする。$f : X \to S$ を
有限表示なアフィンスキームの射とする。デカルト図式
$$
\vcenter{
\xymatrix{
X_1 \ar[d]_{f_1} & X \ar[l] \ar[d]^f \\
S_1 & S \ar[l]
}
}
\quad\text{および}\quad
\vcenter{
\xymatrix{
X_2 \ar[d]_{f_2} & X \ar[l] \ar[d]^f \\
S_2 & S \ar[l]
}
}
$$
が存在すると仮定する。ここでこれらはアフィンスキームの図式で、$S_1$ と $S_2$ は
$\mathbf{Z}$ 上有限型、$f_1$ と $f_2$ は有限型であり、$f_1$ は性質 $P$ を、
$f_2$ は性質 $Q$ をもつ。このときデカルト図式
$$
\xymatrix{
X_0 \ar[d]_{f_0} & X \ar[l] \ar[d]^f \\
S_0 & S \ar[l]
}
$$
で、$S_0$ が $\mathbf{Z}$ 上有限型、$f_0$ が有限型であり、$f_0$ が性質 $P$ と
性質 $Q$ の双方をもつものが存在する。
\end{lemma}

\begin{proof}
与えられた一対の図式は、環の余デカルト図式
$$
\vcenter{
\xymatrix{
B_1 \ar[r] & B \\
A_1 \ar[u] \ar[r] & A \ar[u]
}
}
\quad\text{および}\quad
\vcenter{
\xymatrix{
B_2 \ar[r] & B \\
A_2 \ar[u] \ar[r] & A \ar[u]
}
}
$$
に対応する。$A_1 \to A$ と $A_2 \to A$ の双方の像を含む、$A$ の有限型
$\mathbf{Z}$-部分代数 $A_0 \subset A$ をとる。このような部分代数は、仮定により
$A_1$ と $A_2$ の双方が $\mathbf{Z}$ 上有限型なので存在する。環
$B_{0, 1} = B_1 \otimes_{A_1} A_0$ と $B_{0, 2} = B_2 \otimes_{A_2} A_0$ は
有限型 $A_0$-代数であり、
$B_{0, 1} \otimes_{A_0} A \cong B \cong B_{0, 2} \otimes_{A_0} A$
を満たす。ここで同型は $A$-代数としてのものである。$A$ はその有限型
$A_0$-部分代数の有向余極限なので、Limits, Lemma
\ref{limits-lemma-descend-finite-presentation} により、$A_0$ を拡大した後、
$A_0$-代数の同型 $B_{0, 1} \cong B_{0, 2}$ が存在すると仮定してよい。
性質 $P$ と $Q$ は基底変換で安定と仮定しているので、$S_0 = \Spec(A_0)$ とおき、
$$
X_0 = X_1 \times_{S_1} S_0 =
\Spec(B_{0, 1}) \cong \Spec(B_{0, 2}) = X_2 \times_{S_2} S_0
$$
とおけばよい。
\end{proof}












\section{エタール近傍}
\label{more-morphisms-section-etale-neighbourhoods}

\noindent
射の性質の中には、エタール基底変換を行った後の方が調べやすいものがある。
そこで次の用語を導入しておくと便利である。

\begin{definition}
\label{more-morphisms-definition-etale-neighbourhood}
$S$ をスキームとし、$s \in S$ を点とする。
\begin{enumerate}
\item {\it $(S, s)$ のエタール近傍}とは、対 $(U, u)$ と、
$\varphi(u) = s$ を満たすスキームのエタール射 $\varphi : U \to S$ の組である。
\item $(S, s)$ の{\it エタール近傍の射} $f : (V, v) \to (U, u)$ とは、単に
$f(v) = u$ を満たす $S$-スキームの射 $f : V \to U$ のことである。
\item {\it 初等エタール近傍}とは、$\kappa(s) = \kappa(u)$ を満たすエタール近傍
$\varphi : (U, u) \to (S, s)$ のことである。
\end{enumerate}
\end{definition}

\noindent
% P05-EN-CH37-0171: the frozen literature sentence leaves the first citation parenthesis unclosed; see the sidecar.
初等エタール近傍の概念は文献によりさまざまな名称で呼ばれる。例えば
``エタール近傍''（\cite[Page 36]{Milne}）や
``強エタール''（\cite[Page 108]{KPR}）と呼ばれることがある。
ここでは論文 \cite{GruRay} の慣例に従い、初等エタール近傍と呼ぶ。

\medskip\noindent
$f : (V, v) \to (U, u)$ がエタール近傍の射ならば、$f$ は自動的にエタールである。
Morphisms, Lemma \ref{morphisms-lemma-etale-permanence} を参照されたい。
したがって $(V, v)$ は $(U, u)$ のエタール近傍となる。もちろんエタール射の合成は
エタールなので（Morphisms, Lemma \ref{morphisms-lemma-composition-etale}）、
逆に $(U, u)$ の任意のエタール近傍 $(V, v)$ は $(S, s)$ のエタール近傍でもある。
また、$U \subset S$ が $s$ の開近傍ならば、$(U, s) \to (S, s)$ は
エタール近傍である。これは開埋め込みがエタールであることから従う
（Morphisms, Lemma \ref{morphisms-lemma-open-immersion-etale}）。
本節では、以後これらの注意を断りなく用いる。

\medskip\noindent
$(U, u) \to (S, s)$ がエタール近傍ならば、$\kappa(u)/\kappa(s)$ は有限分離拡大で
あることに注意する。Morphisms, Lemma \ref{morphisms-lemma-etale-at-point} を
参照されたい。

\begin{lemma}
\label{more-morphisms-lemma-realize-prescribed-residue-field-extension-etale}
$S$ をスキームとし、$s \in S$ とする。$k/\kappa(s)$ を有限分離体拡大とする。
このとき、体拡大 $\kappa(u)/\kappa(s)$ が $k/\kappa(s)$ に同型となるエタール近傍
$(U, u) \to (S, s)$ が存在する。
\end{lemma}

\begin{proof}
$S$ はアフィンと仮定してよい。この場合、補題は Algebra, Lemma
\ref{algebra-lemma-make-etale-map-prescribed-residue-field} から従う。
\end{proof}

\begin{lemma}
\label{more-morphisms-lemma-etale-neighbourhoods-not-quite-filtered}
$S$ をスキームとし、$s$ を $S$ の点とする。エタール近傍の圏は次の性質をもつ。
\begin{enumerate}
\item $(U_i, u_i)_{i=1, 2}$ を $S$ における $s$ の2つのエタール近傍とする。
このとき第3のエタール近傍 $(U, u)$ と射
$(U, u) \to (U_i, u_i)$、$i = 1, 2$ が存在する。
\item $h_1, h_2: (U, u) \to (U', u')$ を $s$ のエタール近傍の間の2つの射とする。
$h_1$ と $h_2$ が剰余体の同じ写像 $\kappa(u') \to \kappa(u)$ を誘導すると仮定する。
このときエタール近傍 $(U'', u'')$ と射 $h : (U'', u'') \to (U, u)$ で、
$h_1$ と $h_2$ を等化するもの、すなわち $h_1 \circ h = h_2 \circ h$ を満たすものが
存在する。
\end{enumerate}
\end{lemma}

\begin{proof}
(1) について、ファイバー積 $U = U_1 \times_S U_2$ を考える。エタール射は
基底変換で保たれるので、これは $U_1$ と $U_2$ の双方の上でエタールである。
Morphisms, Lemma \ref{morphisms-lemma-base-change-etale} を参照されたい。
Schemes, Lemma \ref{schemes-lemma-points-fibre-product} のファイバー積の点の記述により、
例えば $u_1$ と $u_2$ の双方へ写る $U$ の点が存在する。
(2) について、$U''$ をファイバー積
$$
\xymatrix{
U'' \ar[r] \ar[d] & U \ar[d]^{(h_1, h_2)} \\
U' \ar[r]^-\Delta & U' \times_S U'.
}
$$
として定める。
% P05-EN-CH37-0172: the frozen proof claims a point over u' with residue field kappa(u'); the natural equalizer point lies over u and has residue field kappa(u); see the sidecar.
$h_1$ と $h_2$ は剰余体の同じ写像 $\kappa(u') \to \kappa(u)$ を誘導するので、
$u'$ 上にあり $\kappa(u'') = \kappa(u')$ を満たす点 $u'' \in U''$ が存在する。
特に $U'' \not = \emptyset$ である。さらに $U'$ は $S$ 上エタールなので、
ファイバー積 $U'\times_S U'$ もそうである（Morphisms, Lemmas
\ref{morphisms-lemma-base-change-etale} と
\ref{morphisms-lemma-composition-etale}）。したがって Morphisms, Lemma
\ref{morphisms-lemma-etale-permanence} により、縦の射 $(h_1, h_2)$ はエタールである。
ゆえに基底変換によって $U''$ は $U'$ 上エタールであり、したがって $S$ 上でも
エタールである（エタール射の合成はエタールだからである）。
よって $(U'', u'')$ は求める解である。
\end{proof}

\begin{lemma}
\label{more-morphisms-lemma-elementary-etale-neighbourhoods}
$S$ をスキームとし、$s$ を $S$ の点とする。$(S, s)$ の初等エタール近傍の圏は
余フィルター付きである（Categories, Definition
\ref{categories-definition-codirected} を参照されたい）。
\end{lemma}

\begin{proof}
これは定義と Lemma \ref{more-morphisms-lemma-etale-neighbourhoods-not-quite-filtered} から直ちに従う。
\end{proof}

\begin{lemma}
\label{more-morphisms-lemma-describe-henselization}
$S$ をスキームとし、$s \in S$ とする。このとき
$$
\mathcal{O}_{S, s}^h =
\colim_{(U, u)} \mathcal{O}(U)
$$
である。ここで余極限は、$(S, s)$ の初等エタール近傍 $(U, u)$ の圏の反対圏である
フィルター付き圏にわたる。
\end{lemma}

\begin{proof}
$\Spec(A) \subset S$ を $s$ のアフィン近傍とし、$\mathfrak p \subset A$ を $s$ に
対応する素イデアルとする。この選択により、標準的な同型
$\mathcal{O}_{S, s} = A_{\mathfrak p}$ と $\kappa(s) = \kappa(\mathfrak p)$ がある。
初等エタール近傍の共終系は、$U$ がアフィンで $U \to S$ が $\Spec(A)$ を経由する
初等エタール近傍 $(U, u)$ によって与えられる。言い換えれば、右辺は
$\colim_{(B, \mathfrak q)} B$ に等しい。ここで余極限は、$\mathfrak p$ 上にあり
$\kappa(\mathfrak p) = \kappa(\mathfrak q)$ を満たす素イデアル $\mathfrak q$ を備えた
エタール $A$-代数 $B$ にわたる。したがって補題は Algebra, Lemma
\ref{algebra-lemma-henselization-different} から従う。
\end{proof}

\noindent
ファイバー上の点のエタール近傍は全空間へ持ち上げられる。

\begin{lemma}
\label{more-morphisms-lemma-lift-etale-neighbourhood-fibre}
\begin{slogan}
ファイバー上の点のエタール近傍を全空間へ持ち上げる。
\end{slogan}
$X \to S$ をスキームの射とし、$x \in X$ の像を $s \in S$ とする。
$(V, v) \to (X_s, x)$ をエタール近傍とする。このときエタール近傍
$(U, u) \to (X, x)$ で、開埋め込みである $(X_s, x)$ のエタール近傍の射
$(U_s, u) \to (V, v)$ が存在するものがある。
\end{lemma}

\begin{proof}
$X$、$V$、$S$ はアフィンと仮定してよい。
% P05-EN-CH37-0173: the frozen affine-coordinate sentence omits the separating commas after its ring-map clauses; see the sidecar.
射 $X \to S$ は $A \to B$、点 $x$ は素イデアル $\mathfrak q \subset B$、
点 $s$ は $\mathfrak p = A \cap \mathfrak q$、射 $V \to X_s$ は
$B \otimes_A \kappa(\mathfrak p) \to C$ によって与えられるとする。
$\kappa(\mathfrak p)$ は $A/\mathfrak p$ の局所化なので、ある $f \in A$、
$f \not \in \mathfrak p$ とエタール環準同型
$B \otimes_A (A/\mathfrak p)_f \to D$ が存在して、
$$
C = (B \otimes_A \kappa(\mathfrak p))
\otimes_{B \otimes_A (A/\mathfrak p)_f} D
$$
を満たす。Algebra, Lemma \ref{algebra-lemma-etale} の (9) を参照されたい。
$A$ を $A_f$ で、$B$ を $B_f$ で置き換えた後、$D$ は
$B \otimes_A A/\mathfrak p = B/\mathfrak p B$ 上エタールであると仮定してよい。
そこで Algebra, Lemma \ref{algebra-lemma-lift-etale} を適用できる。これで補題が従う。
\end{proof}














\section{エタール近傍と分枝}
\label{more-morphisms-section-etale-nbhds-branches}

\noindent
ある点におけるスキームの（幾何学的）分枝数は Properties, Section
\ref{properties-section-unibranch} で定義した。Varieties, Section
\ref{varieties-section-number-of-branches} では、これを正規化射のファイバーと
関係づけた。本節では、この数のエタール近傍による特徴づけを論じる。

\begin{lemma}
\label{more-morphisms-lemma-nr-minimal-primes}
$R = \colim R_i$ を、遷移写像が忠実平坦である環の有向系の余極限とする。
このとき $R$ の最小素イデアルの個数を
$\{0, 1, 2, \ldots, \infty\}$ の元として数えたものは、$R_i$ の
最小素イデアルの個数の上限に等しい。
\end{lemma}

\begin{proof}
$A \to B$ が平坦な環準同型ならば、下降定理により
$\Spec(B) \to \Spec(A)$ は最小素イデアルを最小素イデアルへ写す
（Algebra, Lemma \ref{algebra-lemma-flat-going-down}）。
$A \to B$ が忠実平坦ならば、すべての最小素イデアルは最小素イデアルの像である
（Algebra, Lemma \ref{algebra-lemma-ff-rings} および
\ref{algebra-lemma-minimal-prime-image-minimal-prime} による）。
% P05-EN-CH37-0174: the frozen source reverses the monotonicity inequality for minimal-prime counts; see the sidecar.
したがって $i \leq i'$ ならば、最小素イデアルの個数について
$R_i$ の個数 $\geq$ $R_{i'}$ の個数である。
Algebra, Lemma \ref{algebra-lemma-colimit-faithfully-flat} により、各写像
$R_i \to R$ は忠実平坦であり、さらに最小素イデアルの個数について
$R$ の個数 $\geq$ $R_i$ の個数であることも分かる。
最後に、$\mathfrak q_1, \ldots, \mathfrak q_n$ を $R$ の相異なる最小素イデアルとする。
このとき、$R_i \cap \mathfrak q_1, \ldots, R_i \cap \mathfrak q_n$ が
（集合として、したがって素イデアルとして）相異なるような $i$ を見つけることができる。
以上で補題が従う。
\end{proof}

\begin{lemma}
\label{more-morphisms-lemma-nr-branches}
$X$ をスキーム、$x \in X$ を点とする。このとき次が成り立つ。
\begin{enumerate}
\item $x$ における $X$ の分枝数は、初等エタール近傍
$(U, u) \to (X, x)$ 全体にわたって取った、$u$ を通る $U$ の
既約成分の個数の上限に等しい。
\item $x$ における $X$ の幾何学的分枝数は、エタール近傍
$(U, u) \to (X, x)$ 全体にわたって取った、$u$ を通る $U$ の
既約成分の個数の上限に等しい。
\item $X$ が $x$ において単枝であるための必要十分条件は、すべての
初等エタール近傍 $(U, u) \to (X, x)$ に対して、$u$ を通る $U$ の
既約成分がちょうど一つ存在することである。
\item $X$ が $x$ において幾何学的単枝であるための必要十分条件は、すべての
エタール近傍 $(U, u) \to (X, x)$ に対して、$u$ を通る $U$ の
既約成分がちょうど一つ存在することである。
\end{enumerate}
\end{lemma}

\begin{proof}
(3) と (4) は、Properties, Lemma
\ref{properties-lemma-number-of-branches-1} により (1) と (2) から従う。

\medskip\noindent
(1) の証明。$\Spec(A)$ を $x$ のアフィン開近傍とし、
$\mathfrak p \subset A$ を $x$ に対応する素イデアルとする。
$X$ を $\Spec(A)$ で置き換えてよい。また、アフィン初等エタール近傍
$(U, u)$ は共終部分系をなすので、上限ではそれらだけを考えれば十分である。
Hensel 化 $A_\mathfrak p^h$ は、$B$ がエタール $A$-代数であり、
$\mathfrak q$ が $\mathfrak p$ 上にあって
$\kappa(\mathfrak q) = \kappa(\mathfrak p)$ を満たす素イデアルであるような
対 $(B, \mathfrak q)$ の圏にわたる環 $B_\mathfrak q$ の余極限であることを思い出そう。
Algebra, Lemma \ref{algebra-lemma-henselization-different} を参照されたい。
環とアフィン・スキームの対応により、これらの対 $(B, \mathfrak q)$ は
アフィン初等エタール近傍 $(U, u)$ とちょうど対応する。
$\mathfrak q$ を通る $\Spec(B)$ の既約成分は、ちょうど
$B_\mathfrak q$ の最小素イデアルである。Properties, Definition
\ref{properties-definition-number-of-branches} により、$A_\mathfrak p^h$ の
最小素イデアルの個数は $x$ における $X$ の分枝数である。
この系の遷移写像 $B_\mathfrak q \to B'_{\mathfrak q'}$ はすべて平坦である。
平坦な局所環準同型は忠実平坦なので（Algebra, Lemma
\ref{algebra-lemma-local-flat-ff}）、補題は Lemma
\ref{more-morphisms-lemma-nr-minimal-primes} から従う。

\medskip\noindent
(2) の証明。Algebra, Lemma
\ref{algebra-lemma-strict-henselization-different} を用いることを除けば、証明は
(1) と同じである。わずかな違いがある。$\kappa(x)$ の分離代数閉包
$\kappa^{sep}$ を一つ与えると、$\kappa(u)/\kappa(x)$ は有限分離拡大なので
（Morphisms, Lemma \ref{morphisms-lemma-etale-at-point}）、各エタール近傍
$(U, u)$ に対して $\kappa(x)$-埋め込み
$\phi : \kappa(u) \to \kappa^{sep}$ を選ぶことができる。
したがって、$(U, u) \to (X, x)$ がアフィン・エタール近傍であり、
$\phi : \kappa(u) \to \kappa^{sep}$ が $\kappa(x)$-埋め込みであるような
すべての三つ組 $(U, u, \phi)$ にわたる上限を考えることができる。
これらの三つ組は Algebra, Lemma
\ref{algebra-lemma-strict-henselization-different} の三つ組とちょうど対応し、
証明の残りもまったく同じである。
\end{proof}

\noindent
上の補題の相対版が必要になる。

\begin{lemma}
\label{more-morphisms-lemma-nr-branches-fibre}
$X \to S$ をスキームの射とし、$x \in X$ を像 $s$ をもつ点とする。
このとき次が成り立つ。
\begin{enumerate}
\item $x$ におけるファイバー $X_s$ の分枝数は、初等エタール近傍
$(U, u) \to (X, x)$ 全体にわたって取った、$u$ を通るファイバー $U_s$ の
既約成分の個数の上限に等しい。
\item $x$ におけるファイバー $X_s$ の幾何学的分枝数は、エタール近傍
$(U, u) \to (X, x)$ 全体にわたって取った、$u$ を通るファイバー $U_s$ の
既約成分の個数の上限に等しい。
\item ファイバー $X_s$ が $x$ において単枝であるための必要十分条件は、すべての
初等エタール近傍 $(U, u) \to (X, x)$ に対して、$u$ を通るファイバー $U_s$ の
既約成分がちょうど一つ存在することである。
% P05-EN-CH37-0175: the frozen relative statement names X rather than the fibre X_s in item (4); see the sidecar.
\item $X$ が $x$ において幾何学的単枝であるための必要十分条件は、すべての
エタール近傍 $(U, u) \to (X, x)$ に対して、$u$ を通る $U_s$ の
既約成分がちょうど一つ存在することである。
\end{enumerate}
\end{lemma}

\begin{proof}
Lemmas \ref{more-morphisms-lemma-nr-branches} と
\ref{more-morphisms-lemma-lift-etale-neighbourhood-fibre} を組み合わせればよい。
\end{proof}

\begin{lemma}
\label{more-morphisms-lemma-number-of-branches-and-smooth}
$X \to S$ をスキームの滑らかな射とする。
$x \in X$ の像を $s \in S$ とする。このとき次が成り立つ。
\begin{enumerate}
\item $x$ における $X$ の幾何学的分枝数は、$s$ における $S$ の
幾何学的分枝数に等しい。
\item $\kappa(x)/\kappa(s)$ が純非分離\footnote{実際には、
$\kappa(x)$ が $\kappa(s)$ 上幾何学的に既約であれば十分である。
これが必要になった場合には詳しい証明を加える。}体拡大ならば、
$x$ における $X$ の分枝数は $s$ における $S$ の分枝数に等しい。
\end{enumerate}
\end{lemma}

\begin{proof}
More on Algebra, Lemma
\ref{more-algebra-lemma-invariance-number-branches-smooth} と定義から直ちに従う。
\end{proof}




\section{不分岐射とエタール射}
\label{more-morphisms-section-unramified-is-etale}

\noindent
不分岐射が自動的にエタールとなることがある。

\begin{lemma}
\label{more-morphisms-lemma-unramified-dominant-unibranch-is-etale}
$f : X \to Y$ をスキームの射とする。$x \in X$ の像を $y \in Y$ とする。
次を仮定する。
\begin{enumerate}
\item $Y$ は整であり、$y$ において幾何学的単枝である。
\item $f$ は局所有限型である。
\item $f(x')$ が $Y$ の一般点となるような特殊化 $x' \leadsto x$ が存在する。
\item $f$ は $x$ において不分岐である。
\end{enumerate}
このとき $f$ は $x$ においてエタールである。
\end{lemma}

\begin{proof}
$X$ と $Y$ をそれぞれ $x$ と $y$ の適切なアフィン開近傍で置き換えてよい。
すると $Y$ は整域 $A$ のスペクトルであり、$X$ は有限型 $A$-代数 $B$ の
スペクトルである。$\mathfrak q \subset B$ を $x$ に対応する素イデアル、
$\mathfrak p \subset A$ を $y$ に対応する素イデアルとする。
局所環 $A_\mathfrak p = \mathcal{O}_{Y, y}$ は幾何学的単枝である。
環準同型 $A \to B$ は $\mathfrak q$ において不分岐である。
また、(3) の点 $x'$ は、$A \cap \mathfrak q' = (0)$ を満たす素イデアル
$\mathfrak q' \subset \mathfrak q$ に対応する。
したがって $A_\mathfrak p \to B_\mathfrak q$ は単射である。
結論は More on Algebra, Lemma
\ref{more-algebra-lemma-local-unramified-extension-unibranch-domain-is-etale} から従う。
\end{proof}

\begin{lemma}
\label{more-morphisms-lemma-global-unramified-dominant-unibranch-is-etale}
\begin{reference}
\cite[Expose I, Corollary 9.11]{SGA1}
\end{reference}
$f : X \to Y$ をスキームの射とする。次を仮定する。
\begin{enumerate}
\item $Y$ は整かつ幾何学的単枝である。
\item $X$ の少なくとも一つの既約成分が $Y$ を支配する。
\item $f$ は不分岐である。
\item $X$ は連結である。
\end{enumerate}
このとき $f$ はエタールであり、$X$ は既約である。
\end{lemma}

\begin{proof}
$X' \subset X$ を $Y$ を支配する既約成分とする。これは、$X'$ の一般点が
$Y$ の一般点へ写ることを意味する（たとえば Morphisms, Lemma
\ref{morphisms-lemma-dominant-finite-number-irreducible-components} を参照されたい）。
Lemma \ref{more-morphisms-lemma-unramified-dominant-unibranch-is-etale} により、$f$ は
$X'$ の各点でエタールである。特に、$f$ がエタールである開部分スキーム
$U \subset X$ は $X'$ を含む。$U$ の任意の準コンパクト開部分は有限個の
既約成分をもつことに注意する。Descent, Lemma
\ref{descent-lemma-locally-finite-nr-irred-local-fppf} を参照されたい。
一方、$Y$ は幾何学的単枝であり、$U$ は $Y$ 上エタールなので、スキーム $U$ は
幾何学的単枝である。特に、$U$ の各点を通る既約成分は高々一つである。
ここで単純な位相的議論により、$X' \subset U$ は開かつ閉である。
したがって当然 $X'$ は $X$ でも開かつ閉であり、連結性から所望のとおり
$X = U = X'$ を得る。
\end{proof}

\begin{lemma}
\label{more-morphisms-lemma-weird-permanence-etale}
$f : X \to Y$ と $g : Y \to Z$ をスキームの射とする。
$x \in X$ の像を $y \in Y$ とする。次を仮定する。
\begin{enumerate}
\item $Y$ は整であり、$y$ において幾何学的単枝である。
\item $g \circ f$ は $x$ においてエタールである。
\item $f(x')$ が $Y$ の一般点となるような特殊化 $x' \leadsto x$ が存在する。
\end{enumerate}
このとき $f$ は $x$ においてエタールであり、$g$ は $y$ においてエタールである。
\end{lemma}

\begin{proof}
$X$ を $x$ の開近傍で置き換えることにより、$g \circ f$ はエタールであると
仮定してよい。このとき Morphisms, Lemmas
\ref{morphisms-lemma-unramified-permanence} および
\ref{morphisms-lemma-etale-smooth-unramified} により、$f$ は不分岐である。
したがって Lemma \ref{more-morphisms-lemma-unramified-dominant-unibranch-is-etale} により、
$f$ は $x$ においてエタールである。さらにエタール降下により、$g$ は $y$ において
エタールである。たとえば Descent, Lemma
\ref{descent-lemma-syntomic-smooth-etale-permanence} を参照されたい。
\end{proof}

\begin{lemma}
\label{more-morphisms-lemma-global-weird-permanence-etale}
$f : X \to Y$ と $g : Y \to Z$ をスキームの射とする。次を仮定する。
\begin{enumerate}
\item $Y$ は整かつ幾何学的単枝である。
\item $g \circ f$ はエタールである。
\item $X$ のすべての既約成分は $Y$ を支配する。
\end{enumerate}
このとき $f$ はエタールであり、$g$ は $f$ の像のすべての点において
エタールである。
\end{lemma}

\begin{proof}
点ごとの版である Lemma \ref{more-morphisms-lemma-weird-permanence-etale} から直ちに従う。
\end{proof}










\section{滑らかな射の切断}
\label{more-morphisms-section-etale-over-smooth}

\noindent
本節では、大まかに言えば、スキーム $S$ の滑らかな被覆はエタール被覆によって
細分できるという結果を説明する。その証明法は切断の議論に基づく。

\begin{lemma}
\label{more-morphisms-lemma-slice-smooth-given-element}
$f : X \to S$ をスキームの射とする。
$x \in X$ を像 $s \in S$ をもつ点とする。
$h \in \mathfrak m_x \subset \mathcal{O}_{X, x}$ とする。次を仮定する。
\begin{enumerate}
\item $f$ は $x$ において滑らかである。
\item $\text{d}h$ の
$$
\Omega_{X_s/s, x} \otimes_{\mathcal{O}_{X_s, x}} \kappa(x) =
\Omega_{X/S, x} \otimes_{\mathcal{O}_{X, x}} \kappa(x)
$$
における像 $\text{d}\overline{h}$ は零でない。
\end{enumerate}
このとき $x$ のアフィン開近傍 $U \subset X$ で、$h$ が
$h \in \Gamma(U, \mathcal{O}_U)$ から来ており、$D = V(h)$ が
$x \in D$ を満たす $U$ の有効 Cartier 因子で、$D \to S$ が滑らかであるものが存在する。
\end{lemma}

\begin{proof}
$f$ は $x$ において滑らかなので、$X$ を $x$ の開近傍で置き換えて、$f$ は
滑らかであると仮定してよい。特に $f$ は平坦かつ局所有限表示である。
Lemma \ref{more-morphisms-lemma-slice-given-element} により、$x$ の開近傍 $U \subset X$ で、
$h$ が $h \in \Gamma(U, \mathcal{O}_U)$ から来ており、$D = V(h)$ が
$x \in D$ を満たす $U$ の有効 Cartier 因子で、$D \to S$ が平坦かつ
有限表示であるものが存在することはすでに分かっている。Morphisms, Lemma
\ref{morphisms-lemma-differentials-relative-immersion} により短完全列
$$
\mathcal{C}_{D/U} \to i^*\Omega_{U/S} \to \Omega_{D/S} \to 0
$$
がある。ここで $i : D \to U$ は閉埋め込みであり、$\mathcal{C}_{D/U}$ は
$U$ における $D$ の余法層である。$D$ は
$h \in \Gamma(U, \mathcal{O}_U)$ で切り出される有効 Cartier 因子なので、
% P05-EN-CH37-0176: the frozen conormal-sheaf identity uses O_S although the conormal sheaf is an O_D-module; see the sidecar.
$\mathcal{C}_{D/U} = h \cdot \mathcal{O}_S$ である。$U \to S$ は滑らかなので、
層 $\Omega_{U/S}$ は有限局所自由であり、したがってその引き戻し
$i^*\Omega_{U/S}$ も有限局所自由である。この列の最初の矢は自由生成元 $h$ を
$i^*\Omega_{U/S}$ の切断 $\text{d}h|_D$ へ写し、仮定によりこれはファイバー
$\Omega_{U/S, x} \otimes \kappa(x)$ において零でない。
$\otimes \kappa(x)$ の右完全性により、
$$
\dim_{\kappa(x)} \left( \Omega_{D/S, x} \otimes \kappa(x) \right)
=
\dim_{\kappa(x)} \left( \Omega_{U/S, x} \otimes \kappa(x) \right) - 1.
$$
を得る。Morphisms, Lemma \ref{morphisms-lemma-smooth-at-point} により、
$\Omega_{U/S, x} \otimes \kappa(x)$ は高々 $\dim_x(U_s)$ 個の元で生成できる。
表示した式により、$\Omega_{D/S, x} \otimes \kappa(x)$ は高々
$\dim_x(U_s) - 1$ 個の元で生成できる。さらに、たとえば Algebra, Lemma
\ref{algebra-lemma-one-equation} により
$\dim(\mathcal{O}_{D_s, x}) = \dim(\mathcal{O}_{U_s, x}) - 1$ であり
（また $D_s \subset U_s$ は有効 Cartier である。Divisors, Lemma
\ref{divisors-lemma-relative-Cartier} を参照されたい）、続いて Morphisms, Lemma
\ref{morphisms-lemma-dimension-fibre-at-a-point} を用いると
$\dim_x(D_s) = \dim_x(U_s) - 1$ であることに注意する。
したがって $\Omega_{D/S, x} \otimes \kappa(x)$ は高々 $\dim_x(D_s)$ 個の元で
生成でき、再び Morphisms, Lemma \ref{morphisms-lemma-smooth-at-point} により、
$D \to S$ は $x$ において滑らかである。$U$ を縮小すれば $D \to S$ は滑らかとなり、
証明が完了する。
\end{proof}

\begin{lemma}
\label{more-morphisms-lemma-slice-smooth-once}
$f : X \to S$ をスキームの射とする。
$x \in X$ を像 $s \in S$ をもつ点とする。次を仮定する。
\begin{enumerate}
\item $f$ は $x$ において滑らかである。
\item 写像
$$
\Omega_{X_s/s, x} \otimes_{\mathcal{O}_{X_s, x}} \kappa(x)
\longrightarrow
\Omega_{\kappa(x)/\kappa(s)}
$$
は零でない核をもつ。
\end{enumerate}
このとき、$x$ のアフィン開近傍 $U \subset X$ と、$x$ を含む有効 Cartier 因子
$D \subset U$ で、$D \to S$ が滑らかであるものが存在する。
\end{lemma}

\begin{proof}
$k = \kappa(s)$、$R = \mathcal{O}_{X_s, x}$ と書く。
$\mathfrak m$ を $R$ の極大イデアルとし、$\kappa = R/\mathfrak m$ と置くと、
$\kappa = \kappa(x)$ である。微分加群の形成は局所化と可換なので（Algebra,
Lemma \ref{algebra-lemma-differentials-localize} を参照されたい）、
$\Omega_{X_s/s, x} = \Omega_{R/k}$ である。Algebra, Lemma
\ref{algebra-lemma-differential-seq} により完全列
$$
\mathfrak m/\mathfrak m^2 \xrightarrow{\text{d}}
\Omega_{R/k} \otimes_R \kappa \to
\Omega_{\kappa/k} \to 0.
$$
がある。したがって (2) が成り立てば、$\text{d}\overline{h}$ が零でないような元
$\overline{h} \in \mathfrak m$ が存在する。$\overline{h}$ の持ち上げ
$h \in \mathcal{O}_{X, x}$ を選び、Lemma
\ref{more-morphisms-lemma-slice-smooth-given-element} を適用する。
\end{proof}

\begin{remark}
\label{more-morphisms-remark-necessary-condition-slice-smooth}
Lemma \ref{more-morphisms-lemma-slice-smooth-once} の第二条件は、$x$ が正次元ファイバーの
閉点である場合でさえ必要である。次が一例である。$k$ を標数 $p > 0$ の
非完全体とする。$a \in k$ を $p$ 乗でない元とする。
$\mathfrak m = (x, y^p - a) \subset k[x, y]$ と置く。これは
$X = \mathbf{A}^2_k$ の閉点 $w$ に対応する。$S = \mathbf{A}^1_k$ と置き、
$f : X \to S$ を $k[x] \to k[x, y]$ に対応する射とする。このとき、
$$
\xymatrix{
S' \ar[rr]_h \ar[rd]_g & & X \ar[ld]^f \\
& S
}
$$
という可換図式で、$g$ がエタールかつ $w$ が $h$ の像に入るものは存在しない。
実際、剰余体拡大 $\kappa(w)/\kappa(f(w))$ は純非分離である一方、
$g(s') = f(w)$ を満たす任意の $s' \in S'$ に対して拡大
$\kappa(s')/\kappa(f(w))$ は分離的となるからである。
\end{remark}

\noindent
剰余体拡大が分離的であると仮定すれば、Remark
\ref{more-morphisms-remark-necessary-condition-slice-smooth} の現象は生じない。
正確な結果は次のとおりである。

\begin{lemma}
\label{more-morphisms-lemma-slice-smooth-once-separable-residue-field-extension}
$f : X \to S$ をスキームの射とする。
$x \in X$ を像 $s \in S$ をもつ点とする。次を仮定する。
\begin{enumerate}
\item $f$ は $x$ において滑らかである。
\item 剰余体拡大 $\kappa(x)/\kappa(s)$ は分離的である。
\item $x$ は $X_s$ の一般点ではない。
\end{enumerate}
このとき、$x$ のアフィン開近傍 $U \subset X$ と、$x$ を含む有効 Cartier 因子
$D \subset U$ で、$D \to S$ が滑らかであるものが存在する。
\end{lemma}

\begin{proof}
$k = \kappa(s)$、$R = \mathcal{O}_{X_s, x}$ と書く。
$\mathfrak m$ を $R$ の極大イデアルとし、$\kappa = R/\mathfrak m$ と置くと、
$\kappa = \kappa(x)$ である。微分加群の形成は局所化と可換なので（Algebra,
Lemma \ref{algebra-lemma-differentials-localize} を参照されたい）、
$\Omega_{X_s/s, x} = \Omega_{R/k}$ である。仮定 (2) と Algebra, Lemma
\ref{algebra-lemma-computation-differential} により、写像
$$
\text{d} :
\mathfrak m/\mathfrak m^2
\longrightarrow
\Omega_{R/k} \otimes_R \kappa(\mathfrak m)
$$
は単射である。仮定 (3) から $\mathfrak m/\mathfrak m^2 \not = 0$ が従う。
したがって $\text{d}\overline{h}$ が零でないような元
$\overline{h} \in \mathfrak m$ が存在する。$\overline{h}$ の持ち上げ
$h \in \mathcal{O}_{X, x}$ を選び、Lemma
\ref{more-morphisms-lemma-slice-smooth-given-element} を適用する。
\end{proof}

\noindent
次の補題で構成する部分スキーム $Z$ は、実際には $x$ のアフィン開近傍における
完全交叉である。この事実が必要になった場合には、それを述べる別の補題を明示的に
定式化することにする。

\begin{lemma}
\label{more-morphisms-lemma-slice-smooth}
$f : X \to S$ をスキームの射とする。
$x \in X$ を像 $s \in S$ をもつ点とする。次を仮定する。
\begin{enumerate}
\item $f$ は $x$ において滑らかである。
\item $x$ は $X_s$ の閉点であり、$\kappa(s) \subset \kappa(x)$ は分離的である。
\end{enumerate}
このとき、$x$ を含む埋め込み $Z \to X$ で、次を満たすものが存在する。
\begin{enumerate}
\item $Z \to S$ はエタールである。
\item 集合論的に $Z_s = \{x\}$ である。
\end{enumerate}
\end{lemma}

\begin{proof}
$S$ を $s$ のアフィン開近傍で置き換える。さらに $X$ を、$X \to S$ が滑らかである
$x$ のアフィン開近傍で置き換える。アフィン・スキーム間の滑らかな射について、
$d = \dim_x(X_s)$ に関する帰納法で補題を証明する。

\medskip\noindent
$d = 0$ の場合。この場合、$Z$ を $x$ の開近傍に取れることを示す。
実際、$d = 0$ ならば $X \to S$ は $x$ において準有限である。Morphisms, Lemma
\ref{morphisms-lemma-locally-quasi-finite-rel-dimension-0} を参照されたい。
したがって $U \to S$ が準有限となるようなアフィン開近傍 $U \subset X$ が存在する。
Morphisms, Lemma \ref{morphisms-lemma-quasi-finite-points-open} を参照されたい。
ゆえに $X$ を $U$ で置き換えると、$X$ は $S$ 上準有限かつ滑らかであり、したがって
$S$ 上相対次元 $0$ で滑らか、したがって $S$ 上エタールである。
さらにファイバー $X_s$ は有限離散集合である。
% P05-EN-CH37-0177: the frozen source says an affine open neighbourhood of X where x is required; see the sidecar.
そこで $X$ をさらに $X$ のアフィン開近傍で置き換えると、
$f^{-1}(\{s\}) = \{x\}$ となる（$X_s$ の位相は $X$ の位相から誘導されるからである。
Schemes, Lemma \ref{schemes-lemma-fibre-topological} を参照されたい）。
これでこの場合の補題が証明された。

\medskip\noindent
次に $d > 0$ と仮定する。$x$ はそのファイバーの閉点なので、拡大
$\kappa(x)/\kappa(s)$ は有限であることに注意する（Hilbert の零点定理による。
Morphisms, Lemma
\ref{morphisms-lemma-closed-point-fibre-locally-finite-type} を参照されたい）。
代数的分離体拡大について成り立つことから、
$\Omega_{\kappa(x)/\kappa(s)} = 0$ である。したがって Lemma
\ref{more-morphisms-lemma-slice-smooth-once} を適用して図式
$$
\xymatrix{
D \ar[r] \ar[rrd] & U \ar[r] \ar[rd] & X \ar[d] \\
& & S
}
$$
で $x \in D$ を満たすものを得る。たとえば Algebra, Lemma
\ref{algebra-lemma-one-equation} により
$\dim(\mathcal{O}_{D_s, x}) = \dim(\mathcal{O}_{X_s, x}) - 1$ であり
（また $D_s \subset X_s$ は有効 Cartier である。Divisors, Lemma
\ref{divisors-lemma-relative-Cartier} を参照されたい）、続いて Morphisms, Lemma
\ref{morphisms-lemma-dimension-fibre-at-a-point} を用いると
$\dim_x(D_s) = \dim_x(X_s) - 1$ であることに注意する。
したがって射 $D \to S$ は滑らかで、
$\dim_x(D_s) = \dim_x(X_s) - 1 = d - 1$ を満たす。帰納法の仮定により、
所望の埋め込み $Z \to D$ を見つけることができ、証明が完了する。
\end{proof}

\begin{lemma}
\label{more-morphisms-lemma-etale-nbhd-dominates-smooth}
\begin{slogan}
滑らかな射はエタール局所切断をもつ。
\end{slogan}
\begin{reference}
\cite[Corollaire 17.16.3 (ii)]{EGA4} を参照されたい。
\end{reference}
$f : X \to S$ をスキームの滑らかな射とする。
$s \in S$ を $f$ の像に属する点とする。このときエタール近傍
$(S', s') \to (S, s)$ と $S$-射 $S' \to X$ が存在する。
\end{lemma}

\begin{proof}[Lemma \ref{more-morphisms-lemma-etale-nbhd-dominates-smooth} の第一証明]
仮定により $X_s \not = \emptyset$ である。Varieties, Lemma
\ref{varieties-lemma-smooth-separable-closed-points-dense} により、
$\kappa(x)$ が $\kappa(s)$ の有限分離体拡大であるような閉点 $x \in X_s$ が存在する。
したがって Lemma \ref{more-morphisms-lemma-slice-smooth} により、$Z \to S$ がエタールで
$x \in Z$ を満たす埋め込み $Z \to X$ が存在する。
$(S' , s') = (Z, x)$ と取ればよい。
\end{proof}

\begin{proof}[Lemma \ref{more-morphisms-lemma-etale-nbhd-dominates-smooth} の第二証明]
$f(x) = s$ を満たす点 $x \in X$ を選ぶ。図式
$$
\xymatrix{
X \ar[d] & U \ar[l] \ar[d] \ar[r]_-\pi & \mathbf{A}^d_V \ar[ld] \\
S & V \ar[l]
}
$$
で、$\pi$ がエタール、$x \in U$、かつ $V = \Spec(R)$ がアフィンであるものを選ぶ。
Morphisms, Lemma \ref{morphisms-lemma-smooth-etale-over-affine-space} を
参照されたい。特に $s \in V$ である。射 $\pi : U \to \mathbf{A}^d_V$ は開である。
Morphisms, Lemma \ref{morphisms-lemma-etale-open} を参照されたい。
したがって $W = \pi(U) \cap \mathbf{A}^d_s$ は $\mathbf{A}^d_s$ の空でない
開部分集合である。$\kappa(s) \subset \kappa(w)$ が有限分離となる点 $w \in W$ を取る。
Varieties, Lemma \ref{varieties-lemma-affine-space-over-field} を参照されたい。
Algebra, Lemma \ref{algebra-lemma-dim-affine-space} により、
$\kappa(s)[x_1, \ldots, x_d]$ において $w$ に対応する極大イデアルを生成する $d$ 個の元
$\overline{f}_1, \ldots, \overline{f}_d \in \kappa(s)[x_1, \ldots, x_d]$ が存在する。
$R$ を単項局所化で置き換えると、
$\overline{f}_1, \ldots, \overline{f}_d \in \kappa(s)[x_1, \ldots, x_d]$ へ写る
$f_1, \ldots, f_d \in R[x_1, \ldots, x_d]$ が存在すると仮定してよい。
$R$-代数
$$
R' = R[x_1, \ldots, x_d]/(f_1, \ldots, f_d)
$$
を考え、$S' = \Spec(R')$ と置く。構成により $V$ 上の閉埋め込み
$j : S' \to \mathbf{A}^d_V$ がある。構成により、$s$ 上の $S' \to V$ のファイバーは
一点 $s'$ で、その剰余体は $\kappa(s)$ 上有限分離である。
$\mathfrak q' \subset R'$ を対応する素イデアルとする。Algebra, Lemma
\ref{algebra-lemma-localize-relative-complete-intersection} により、ある
% P05-EN-CH37-0178: the frozen localization condition uses the undefined q instead of the just-defined q'; see the sidecar.
$g \in R'$、$g \not \in \mathfrak q$ に対して $(R')_g$ は $R$ 上の相対大域的完全交叉である。
したがって $S' \to V$ は $s'$ の近傍で平坦かつ有限表示である。Algebra, Lemma
\ref{algebra-lemma-relative-global-complete-intersection} を参照されたい。
構成により、$s$ 上の $S' \to V$ のスキーム論的ファイバーは
$\Spec(\kappa(s'))$ である。ゆえに Morphisms, Lemma
\ref{morphisms-lemma-etale-at-point} から、$S' \to S$ は $s'$ において
エタールである。ここで
$$
S'' = U \times_{\pi, \mathbf{A}^d_V, j} S'.
$$
と置く。構成により、射影 $p : S'' \to S'$ によって $s'$ へ写る点
$s'' \in S''$ が存在する。$p$ はエタール射 $\pi$ の基底変換なのでエタールである。
Morphisms, Lemma \ref{morphisms-lemma-base-change-etale} を参照されたい。
$S' \to S$ がエタールである $s' \in S'$ の空でない開近傍へ写るような、$s''$ の
小さなアフィン近傍 $S''' \subset S''$ を選ぶ。このときエタール近傍
$(S''', s'') \to (S, s)$ が補題の問題の解である。
\end{proof}

\noindent
次の補題は、滑らかな位相に関する層とエタール位相に関する層が同じものであることを示す。

\begin{lemma}
\label{more-morphisms-lemma-etale-dominates-smooth}
$S$ をスキームとする。$\mathcal{U} = \{S_i \to S\}_{i \in I}$ を $S$ の
滑らかな被覆とする。Topologies, Definition
\ref{topologies-definition-smooth-covering} を参照されたい。このとき
$\mathcal{U}$ を細分する（Sites, Definition
\ref{sites-definition-morphism-coverings} を参照されたい）エタール被覆
$\mathcal{V} = \{T_j \to S\}_{j \in J}$ が存在する（Topologies, Definition
\ref{topologies-definition-etale-covering} を参照されたい）。
\end{lemma}

\begin{proof}
各 $s \in S$ に対して、$s$ が $S_i \to S$ の像に入るような $i \in I$ が存在する。
Lemma \ref{more-morphisms-lemma-etale-nbhd-dominates-smooth} により、エタール射
$g_s : T_s \to S$ で、$s \in g_s(T_s)$ を満たし、かつ $g_s$ が
$S_i \to S$ を経由するものを見つけることができる。
したがって $\{T_s \to S\}$ は $\mathcal{U}$ を細分する $S$ のエタール被覆である。
\end{proof}

\begin{lemma}
\label{more-morphisms-lemma-cover-smooth-by-special}
$f : X \to S$ をスキームの滑らかな射とする。このときエタール被覆
$\{U_i \to X\}_{i \in I}$ で、$U_i \to S$ が $U_i \to V_i \to S$ と分解し、
$V_i \to S$ はエタール、$U_i \to V_i$ は切断をもち幾何学的に連結な
ファイバーをもつアフィン・スキーム間の滑らかな射であるものが存在する。
\end{lemma}

\begin{proof}
$s \in S$ とする。Varieties, Lemma
\ref{varieties-lemma-smooth-separable-closed-points-dense} により、
$\kappa(x)/\kappa(s)$ が分離的であるような閉点 $x \in X_s$ の集合は $X_s$ で稠密である。
したがって、$x$ を像に含むエタール射 $U \to X$ で、$U \to S$ が補題に記した仕方で
分解するものを構成すれば十分である。そのために、$x$ を通り、$Z \to S$ がエタールである
埋め込み $Z \to X$ を選ぶ（Lemma \ref{more-morphisms-lemma-slice-smooth}）。
$S$ を $Z$ で、$X$ を $Z \times_S X$ で置き換えると、$X \to S$ は
$\sigma(s) = x$ を満たす切断 $\sigma : S \to X$ をもつと仮定してよい。
そこでまず $S$ を $s$ のアフィン開近傍で置き換え、次に $X$ を $x$ の
アフィン開近傍で置き換える。最後に Section
\ref{more-morphisms-section-connected-components} の部分集合 $X^0 \subset X$ を考える。
Lemmas \ref{more-morphisms-lemma-connected-along-section-open} および
\ref{more-morphisms-lemma-connected-along-section-locally-constructible} により、これは $\sigma$ を含む
逆コンパクトな開部分スキームであり、$X^0 \to S$ のファイバーは幾何学的に連結である。
$X^0$ がアフィンでなければ、$x$ を含むアフィン開集合 $U \subset X^0$ を選ぶ。
$X^0 \to S$ は滑らかなので、$U$ の像は開である。$\sigma^{-1}(U)$ と
$U \to S$ の像に含まれる $s$ のアフィン開近傍 $V \subset S$ を選ぶ。
最後に、$U \cap f^{-1}(V) \to V$ が所望の性質をすべてもつことが分かる。
たとえば $U \cap f^{-1}(V)$ は $U \times_S V$ に等しく、アフィン・スキームの
ファイバー積なのでアフィンである。また、$U \cap f^{-1}(V) \to V$ の幾何学的
ファイバーは $X^0 \to S$ の既約ファイバーの空でない開部分スキームであり、
したがって連結である。いくつかの詳細は省略する。
\end{proof}











\section{エタール近傍と Artin 近似}
\label{more-morphisms-section-etale-nbhds-artin}

\noindent
本節では、二つの点付きスキームの完備局所環が同型ならば、それらは同型なエタール近傍を
もつ、という形の結果を証明する。Popescu の定理を用いる。Smoothing Ring Maps,
Theorem \ref{smoothing-theorem-popescu} を参照されたい。

\begin{lemma}
\label{more-morphisms-lemma-map-approximation}
$S$ を局所 Noetherian スキームとする。$X$、$Y$ を $S$ 上局所有限型な
スキームとする。$x \in X$ と $y \in Y$ を同じ点 $s \in S$ 上にある点とする。
$\mathcal{O}_{S, s}$ は G-環であると仮定する。さらに局所
$\mathcal{O}_{S, s}$-代数準同型
$$
\varphi : \mathcal{O}_{Y, y} \longrightarrow \mathcal{O}_{X, x}^\wedge
$$
が与えられているとする。各 $N \geq 1$ に対して、初等エタール近傍
$(U, u) \to (X, x)$ と、$u$ を $y$ へ写す $S$-射 $f : U \to Y$ で、図式
$$
\xymatrix{
\mathcal{O}_{X, x}^\wedge \ar[r] &
\mathcal{O}_{U, u}^\wedge \\
\mathcal{O}_{Y, y} \ar[r]^{f^\sharp_u} \ar[u]^\varphi &
\mathcal{O}_{U, u} \ar[u]
}
$$
が $\mathfrak m_u^N$ を法として可換となるものが存在する。
\end{lemma}

\begin{proof}
問題は $X$ 上局所的なので、$X$、$Y$、$S$ はアフィンであると仮定してよい。
$S = \Spec(R)$、$X = \Spec(A)$、$Y = \Spec(B)$ とする。
$B = R[x_1, \ldots, x_n]/(f_1, \ldots, f_m)$ と書く。
$\mathfrak p \subset A$ を $x$ に対応する素イデアルとする。More on Algebra,
Proposition \ref{more-algebra-proposition-finite-type-over-G-ring} により、局所環
$\mathcal{O}_{X, x} = A_\mathfrak p$ は G-環である。準同型 $\varphi$ は
$$
B_\mathfrak q^\wedge \longrightarrow A_\mathfrak p^\wedge
$$
である。ここで $\mathfrak q \subset B$ は $y$ に対応する素イデアルである。
$a_1, \ldots, a_n \in A_\mathfrak p^\wedge$ を、
$R[x_1, \ldots, x_n] \to B \to B_\mathfrak q^\wedge \to A_\mathfrak p^\wedge$
による $x_1, \ldots, x_n$ の像とする。Smoothing Ring Maps, Lemma
\ref{smoothing-lemma-approximation-property-variant} を適用すると、エタール環準同型
$A \to A'$、素イデアル $\mathfrak p' \subset A'$、および
$b_1, \ldots, b_n \in A'$ で、
$\kappa(\mathfrak p) = \kappa(\mathfrak p')$、
$a_i - b_i \in (\mathfrak p')^N(A'_{\mathfrak p'})^\wedge$、かつ
% P05-EN-CH37-0179: the frozen relation index runs to n although the presentation has m defining relations; see the sidecar.
$f_j(b_1, \ldots, b_n) = 0$ for $j = 1, \ldots, n$ を満たすものを得る。
これにより、$x_i$ の類を $b_i \in A'$ へ写す $R$-代数準同型 $B \to A'$ が定まる。
射 $f$ を $U = \Spec(A') \to \Spec(B)$ とし、$u = \mathfrak p'$ と取れば証明が完了する。
\end{proof}

\begin{lemma}
\label{more-morphisms-lemma-isomorphism-approximation}
$S$ を局所 Noetherian スキームとする。$X$、$Y$ を $S$ 上局所有限型な
スキームとする。$x \in X$ と $y \in Y$ を同じ点 $s \in S$ 上にある点とする。
$\mathcal{O}_{S, s}$ は G-環であると仮定する。完備局所環の間に
$\mathcal{O}_{S, s}$-代数同型
$$
\varphi : \mathcal{O}_{Y, y}^\wedge \longrightarrow \mathcal{O}_{X, x}^\wedge
$$
があると仮定する。このとき各 $N \geq 1$ に対して、$S$ 上の点付きスキームの射
$$
(X, x) \leftarrow (U, u) \rightarrow (Y, y)
$$
で、両方の矢が初等エタール近傍を定め、図式
$$
\xymatrix{
& \mathcal{O}_{U, u}^\wedge \\
\mathcal{O}_{Y, y}^\wedge \ar[rr]^\varphi \ar[ru] & &
\mathcal{O}_{X, x}^\wedge \ar[lu]
}
$$
が $\mathfrak m_u^N$ を法として可換となるものが存在する。
\end{lemma}

\begin{proof}
$N \geq 2$ と仮定してよい。Lemma \ref{more-morphisms-lemma-map-approximation} を適用して、
$(U, u) \to (X, x)$ と $f : (U, u) \to (Y, y)$ を得る。
$f$ は $u$ においてエタールであると主張する。これで証明が完了する。
実際、誘導される写像
$\mathcal{O}_{Y, y}^\wedge \to \mathcal{O}_{U, u}^\wedge$ が同型であることを示す。
これが示されれば、Lemma \ref{more-morphisms-lemma-lifting-along-artinian-at-point} により $f$ は
$u$ において滑らかであり、もちろん $f$ は $u$ において不分岐でもあるから、
Morphisms, Lemma \ref{morphisms-lemma-etale-smooth-unramified} により $f$ は
$u$ においてエタールである。局所環 $(R, \mathfrak m)$ に対して
$\text{Gr}_\mathfrak m(R) =
\bigoplus_{n \geq 0} \mathfrak m^n/\mathfrak m^{n + 1}$.
と置く。主張を証明するため、誘導される次数付き環の図式
$$
\xymatrix{
& \text{Gr}_{\mathfrak m_u}(\mathcal{O}_{U, u}) \\
\text{Gr}_{\mathfrak m_y}(\mathcal{O}_{Y, y}) \ar[rr]^\varphi \ar[ru] & &
\text{Gr}_{\mathfrak m_x}(\mathcal{O}_{X, x}) \ar[lu]
}
$$
を考える。$N \geq 2$ なので、表示した次数付き代数は次数 $1$ で生成されることから、
この図式は実際に可換である。仮定により下の矢は同型である。
たとえば More on Algebra, Lemma \ref{more-algebra-lemma-flat-unramified} により、写像
$\mathcal{O}_{X, x}^\wedge \to \mathcal{O}_{U, u}^\wedge$ は同型であり、したがって
図式の北西向きの矢は同型である。ゆえに $f$ は同型
% P05-EN-CH37-0180: the frozen associated-graded conclusion uses X/m_x as source and m_y on O_U; see the sidecar.
$\text{Gr}_{\mathfrak m_x}(\mathcal{O}_{X, x}) \to
\text{Gr}_{\mathfrak m_y}(\mathcal{O}_{U, u})$
を誘導する。両方の局所環について、帰納法と短完全列
$$
0 \to \text{Gr}^n_{\mathfrak m}(R) \to R/\mathfrak m^{n + 1} \to
R/\mathfrak m^n \to 0
$$
を用いると、蛇の補題により $f$ はすべての $n$ に対して同型
$\mathcal{O}_{Y, y}/\mathfrak m_y^n \to \mathcal{O}_{U, u}/\mathfrak m_u^n$
を誘導する。これは示すべきことであった。
\end{proof}

\begin{lemma}
\label{more-morphisms-lemma-relative-map-approximation-pre}
$X \to S$、$Y \to T$、$x$、$s$、$y$、$t$、$\sigma$、$y_\sigma$、および
$\varphi$ を次のように与える。スキームの射
$$
\vcenter{
\xymatrix{
X \ar[d] & Y \ar[d] \\
S & T
}
}
\quad\text{with points}\quad
\vcenter{
\xymatrix{
x \ar[d] & y \ar[d] \\
s & t
}
}
$$
があり、点も表示のとおりとする。ここで $S$ は局所 Noetherian、$T$ は
$\mathbf{Z}$ 上有限型である。射 $X \to S$ と $Y \to T$ は局所有限型である。
局所環 $\mathcal{O}_{S, s}$ は G-環である。写像
$$
\sigma : \mathcal{O}_{T, t} \longrightarrow \mathcal{O}_{S, s}^\wedge
$$
は局所準同型である。
$Y_\sigma = Y \times_{T, \sigma} \Spec(\mathcal{O}_{S, s}^\wedge)$ と置く。
次に $y_\sigma$ は、$y$ と $\Spec(\mathcal{O}_{S, s}^\wedge)$ の閉点へ写る
$Y_\sigma$ の点とする。最後に
$$
\varphi :
\mathcal{O}_{X, x}^\wedge
\longrightarrow
\mathcal{O}_{Y_\sigma, y_\sigma}^\wedge
$$
は $\mathcal{O}_{S, s}^\wedge$-代数の同型とする。この状況では可換図式
$$
\xymatrix{
X \ar[d] &
W \ar[l] \ar[rd] \ar[rr] & &
Y \times_{T, \tau} V \ar[r] \ar[ld] & Y \ar[d] \\
S & &
V \ar[ll] \ar[rr]^\tau & &
T
}
$$
と点 $w \in W$、$v \in V$ で、次を満たすものが存在する。
\begin{enumerate}
\item $(V, v) \to (S, s)$ は初等エタール近傍である。
\item $(W, w) \to (X, x)$ は初等エタール近傍である。
\item $\tau(v) = t$ である。
\end{enumerate}
$y_\tau \in Y \times_T V$ を、同一視
$(Y_\sigma)_s = (Y \times_T V)_v$ によって $y_\sigma$ に対応する点とする。
このとき
\begin{enumerate}
\item[(4)] $(W, w) \to (Y \times_{T, \tau} V, y_\tau)$ は初等エタール近傍である。
\end{enumerate}
\end{lemma}

\begin{proof}
$X_\sigma = X \times_S \Spec(\mathcal{O}_{S, s}^\wedge)$ と置き、
$x_\sigma \in X_\sigma$ を $x$ 上にある唯一の点とする。More on Algebra,
Proposition \ref{more-algebra-proposition-Noetherian-complete-G-ring} により、
$\mathcal{O}_{S, s}^\wedge$ は G-環であることに注意する。
Lemma \ref{more-morphisms-lemma-isomorphism-approximation} により
$$
(X_\sigma, x_\sigma) \leftarrow (U, u) \rightarrow (Y_\sigma, y_\sigma)
$$
で、両方の矢が初等エタール近傍であるものを選べる。

\medskip\noindent
$S$ を $s$ の開近傍で置き換えると、$S = \Spec(R)$ はアフィンであると仮定してよい。
$\mathcal{O}_{S, s}$ は G-環なので、Smoothing Ring Maps, Theorem
\ref{smoothing-theorem-popescu} により、環 $\mathcal{O}_{S, s}^\wedge$ は
滑らかな $R$-代数のフィルター付き余極限である。したがって
$$
\Spec(\mathcal{O}_{S, s}^\wedge) = \lim S_i
$$
と、$S$ 上滑らかなアフィン・スキーム $S_i$ の有向極限として書ける。
$s_i \in S_i$ を $\Spec(\mathcal{O}_{S, s}^\wedge)$ の閉点の像とする。
$\kappa(s) = \kappa(s_i)$ であることに注意する。
$X_i = X \times_S S_i$ と置き、$x_i \in X_i$ を $x$ へ写る唯一の点とする。
$\kappa(x) = \kappa(x_i)$ であることに注意する。$T$ は $\mathbf{Z}$ 上有限型なので、
Limits, Proposition \ref{limits-proposition-characterize-locally-finite-presentation} により、
ある $i$ と点付きスキームの射 $\sigma_i : (S_i, s_i) \to (T, t)$ で、
$\Spec(\mathcal{O}_{S, s}^\wedge) \to S_i$ との合成が $\sigma$ に等しいものを選べる。
$Y_i = Y \times_T S_i$ と置き、$y_i$ を $y_\sigma$ の像とする。
$\kappa(y_i) = \kappa(y_\sigma)$ であることに注意する。Limits, Lemma
\ref{limits-lemma-descend-finite-presentation} により、ある $i$ と図式
$$
\xymatrix{
X_i \ar[rd] &
U_i \ar[l] \ar[d] \ar[r] &
Y_i \ar[ld] \\
& S_i
}
$$
で、$\Spec(\mathcal{O}_{S, s}^\wedge)$ への基底変換が
$X_\sigma \leftarrow U \rightarrow Y_\sigma$ を復元するものを選べる。
Limits, Lemma \ref{limits-lemma-descend-etale} により、$i$ を大きくすれば、射
$X_i \leftarrow U_i \rightarrow Y_i$ はエタールであると仮定してよい。
$u_i \in U_i$ を $u$ の像とする。このとき $u_i \mapsto x_i$ なので
$\kappa(x) = \kappa(x_\sigma) = \kappa(u) \supset \kappa(u_i) \supset
\kappa(x_i) = \kappa(x)$ であり、$\kappa(u_i) = \kappa(x_i)$ を得る。
したがって $(X_i, x_i) \leftarrow (U_i, u_i)$ は初等エタール近傍である。
さらに $\kappa(y_i) = \kappa(y_\sigma) = \kappa(u)$ なので、
$(U_i, u_i) \to (Y_i, y_i)$ も初等エタール近傍である。

\medskip\noindent
この時点で、補題の主張にある形の図式
$$
\xymatrix{
X \ar[d] &
X \times_S S_i \ar[l] \ar[rd] &
U_i \ar[l] \ar[r] \ar[d] &
Y \times_T S_i \ar[r] \ar[ld] &
Y \ar[d] \\
S & &
S_i \ar[ll] \ar[rr] & &
T
}
$$
を構成した。ただし $S_i \to S$ は滑らかである。Lemma
\ref{more-morphisms-lemma-slice-smooth} により、$S_i$ を縮小すれば、$s_i$ を通る閉部分スキーム
$V \subset S_i$ で $V \to S$ がエタールであるものが存在すると仮定してよい。
$W$ を $U_i$ における $V$ のスキーム論的逆像に取れば結論を得る。
\end{proof}

\noindent
本節の残りは読み飛ばすことを読者に強く勧める。

\begin{lemma}
\label{more-morphisms-lemma-relative-map-approximation}
図式
$$
\vcenter{
\xymatrix{
X \ar[d] & Y \ar[d] \\
S & T \ar[l]
}
}
\quad\text{with points}\quad
\vcenter{
\xymatrix{
x \ar[d] & y \ar[d] \\
s & t \ar[l]
}
}
$$
を考える。ここで $S$ は局所 Noetherian スキームであり、各射は局所有限型であるとする。
$\mathcal{O}_{S, s}$ は G-環であると仮定する。さらに局所
$\mathcal{O}_{S, s}$-代数準同型
$$
\sigma : \mathcal{O}_{T, t} \longrightarrow \mathcal{O}_{S, s}^\wedge
$$
と局所 $\mathcal{O}_{S, s}$-代数準同型
$$
\varphi :
\mathcal{O}_{X, x}
\longrightarrow
\mathcal{O}_{Y_\sigma, y_\sigma}^\wedge
$$
が与えられているとする。ここで
$Y_\sigma = Y \times_{T, \sigma} \Spec(\mathcal{O}_{S, s}^\wedge)$ であり、
$y_\sigma$ は $y$ 上にある $Y_\sigma$ の唯一の点である。
各 $N \geq 1$ に対して、$S$ 上のスキームの可換図式
$$
\xymatrix{
X \ar[d] & X \times_S V \ar[l] \ar[rd] &
W \ar[l]^-f \ar[r] \ar[d] &
Y \times_{T, \tau} V \ar[r] \ar[ld] & Y \ar[d] \\
S & & V \ar[ll] \ar[rr]^\tau & & T
}
$$
と点 $w \in W$、$v \in V$ で、次を満たすものが存在する。
\begin{enumerate}
\item $v \mapsto s$、$\tau(v) = t$、$f(w) = (x, v)$、かつ $w \mapsto (y, v)$ である。
\item $(V, v) \to (S, s)$ は初等エタール近傍である。
\item 図式
$$
\xymatrix{
\mathcal{O}_{S, s}^\wedge \ar[r] & \mathcal{O}_{V, v}^\wedge \\
\mathcal{O}_{T, t} \ar[r]^{\tau^\sharp_v} \ar[u]_\sigma &
\mathcal{O}_{V, v} \ar[u]
}
$$
% P05-EN-CH37-0181: the frozen sentence says "commutes module" rather than "commutes modulo"; see the sidecar.
は $\mathfrak m_v^N$ を法として可換である。
\item $(W, w) \to (Y \times_{T, \tau} V, (y, v))$ は初等エタール近傍である。
\item 図式
$$
\xymatrix{
\mathcal{O}_{X, x} \ar[r]_\varphi &
\mathcal{O}_{Y_\sigma, y_\sigma}^\wedge \ar[r] &
\mathcal{O}_{Y_\sigma, y_\sigma}/\mathfrak m_{y_\sigma}^N \ar@{=}[r] &
\mathcal{O}_{Y \times_{T, \tau} V, (y, v)}/\mathfrak m_{(y, v)}^N
\ar[d]_{\cong} \\
\mathcal{O}_{X, x} \ar[r] \ar@{=}[u] &
\mathcal{O}_{X \times_S V, (x, v)} \ar[r]^{f^\sharp_w} &
\mathcal{O}_{W, w} \ar[r] &
\mathcal{O}_{W, w}/\mathfrak m_w^N
}
$$
は可換である。等号は、(2) と (3) により
$\mathcal{O}_{V, v}/\mathfrak m_v^N = \mathcal{O}_{S, s}/\mathfrak m_s^N$ 上で
$Y_\sigma$ と $Y \times_{T, \tau} V$ が標準的に同型であることから来る。
\end{enumerate}
\end{lemma}

\begin{proof}
$X$、$S$、$T$、$Y$ をアフィン開部分スキームで置き換えると、補題の主張の図式は
図式
$$
\vcenter{
\xymatrix{
A & B \\
R \ar[u] \ar[r] & C \ar[u]
}
}
\quad\text{with primes}\quad
\vcenter{
\xymatrix{
\mathfrak p_A & \mathfrak p_B \\
\mathfrak p_R \ar@{-}[u] \ar@{-}[r] & \mathfrak p_C \ar@{-}[u]
}
}
$$
に $\Spec$ を適用して得られると仮定してよい。ここで環は Noetherian で、環準同型は
有限型である。本証明では、各環 $E$ は $\mathfrak p_R$ 上にある素イデアル
$\mathfrak p_E$ を備えた Noetherian $R$-代数であり、すべての環準同型は指定された
素イデアルと両立する $R$-代数準同型であるものとする。さらに $E^\wedge$ と書くとき、
$E$ の $\mathfrak p_E$ における局所化の完備化を意味する。また More on Algebra,
Lemma \ref{more-algebra-lemma-flat-unramified} により、
$\kappa(\mathfrak p_{E_1}) = \kappa(\mathfrak p_{E_2})$ を満たすエタール環準同型
$E_1 \to E_2$ は同型 $E_1^\wedge = E_2^\wedge$ を誘導することを、以下では断りなく用いる。

\medskip\noindent
この記法のもとで、$\sigma$ と $\varphi$ は環準同型
$$
\sigma : C \to R^\wedge
\quad\text{and}\quad
\varphi : A \longrightarrow (B \otimes_{C, \sigma} R^\wedge)^\wedge
$$
に対応する。図で表すと
$$
\xymatrix{
A \ar@/^1em/[rrr]^\varphi &
B \ar[r] &
B \otimes_{C, \sigma} R^\wedge \ar[r] &
(B \otimes_{C, \sigma} R^\wedge)^\wedge \\
R \ar[r] \ar[u] &
C \ar[r]^\sigma \ar[u] & R^\wedge \ar[u] \ar[ru]
}
$$
である。More on Algebra, Proposition
\ref{more-algebra-proposition-Noetherian-complete-G-ring} により、$R^\wedge$ は
G-環であることに注意する。したがって More on Algebra, Proposition
\ref{more-algebra-proposition-finite-type-over-G-ring} により、
$B \otimes_{C, \sigma} R^\wedge$ は G-環である。Lemma
\ref{more-morphisms-lemma-map-approximation} を代数に翻訳して適用すると、同型
$\kappa(\mathfrak p_{B \otimes_{C, \sigma} R^\wedge})
\to \kappa(\mathfrak p_{B'})$
を誘導するエタール環準同型 $B \otimes_{C, \sigma} R^\wedge \to B'$ と、合成
$$
A \to B' \to (B')^\wedge = (B \otimes_{C, \sigma} R^\wedge)^\wedge
$$
が $(\mathfrak p_{(B \otimes_{C, \sigma} R^\wedge)^\wedge})^N$ を法として
$\varphi$ と一致するような $R$-代数準同型 $A \to B'$ が存在する。
$\varphi$ が結論に現れるのは補題の主張の (5) における可換性の要請だけなので、
$\varphi$ をこの合成で置き換えてよい。図で表すと
$$
\xymatrix{
& & B' \ar[r] & (B')^\wedge \ar@{=}[d] \\
A \ar[rru] &
B \ar[r] &
B \otimes_{C, \sigma} R^\wedge \ar[r] \ar[u] &
(B \otimes_{C, \sigma} R^\wedge)^\wedge \\
R \ar[r] \ar[u] &
C \ar[r]^\sigma \ar[u] & R^\wedge \ar[u] \ar[ru]
}
$$
次に、仮定により $R_{\mathfrak p_R}$ は G-環なので、$R^\wedge$ は滑らかな
$R$-代数のフィルター付き余極限であることを用いる（Smoothing Ring Maps, Theorem
\ref{smoothing-theorem-popescu}）。$C$ は $R$ 上有限表示なので、分解
$$
C \to R' \to R^\wedge
$$
で $R \to R'$ が滑らかであるものを得る。Algebra, Lemma
\ref{algebra-lemma-characterize-finite-presentation} を参照されたい。
$R'$ を大きくすれば、$B \otimes_{C, \sigma} R^\wedge$ への基底変換が $B'$ となる
エタール $B \otimes_C R'$-代数 $B''$ が存在すると仮定してよい。Algebra, Lemma
\ref{algebra-lemma-etale} を参照されたい。このとき $B'$ はこれらの $B''$ の
フィルター付き余極限であるから、$R'$ をさらに大きくすれば、
$A \to B'' \to B'$ が先に構成した写像となる $R$-代数準同型 $A \to B''$ が存在すると
仮定してよい（上と同じ参照先による）。図で表すと
$$
\xymatrix{
& & B'' \ar[r] & B' \ar[r] & (B')^\wedge \ar@{=}[d] \\
A \ar[rru] &
B \ar[r] &
B \otimes_C R' \ar[r] \ar[u] &
B \otimes_{C, \sigma} R^\wedge \ar[r] \ar[u] &
(B \otimes_{C, \sigma} R^\wedge)^\wedge \\
R \ar[r] \ar[u] &
C \ar[r] \ar[u] &
R' \ar[r] \ar[u] &
R^\wedge \ar[u] \ar[ru]
}
$$
であり、さらに
$$
B' = B'' \otimes_{(B \otimes_C R')} (B \otimes_{C, \sigma} R^\wedge)
$$
である。したがって $C$ を $R'$ で、$\sigma : C \to R^\wedge$ を
$R' \to R^\wedge$ で、$B$ を $B''$ で置き換え、図式を
$$
\xymatrix{
A \ar[r] &
B \ar[r] &
B \otimes_{C, \sigma} R^\wedge \\
R \ar[r] \ar[u] &
C \ar[r]^\sigma \ar[u] & R^\wedge \ar[u]
}
$$
へ簡約できる。ここで $\varphi$ は横の矢の合成と、それに続く
$B \otimes_{C, \sigma} R^\wedge$ からその完備化への標準写像に等しい。
証明の最後の段階では、$\Spec(R)$ 上の $\Spec(C)$ と $\Spec(R)$、および写像
$C \to R^\wedge$ に Lemma \ref{more-morphisms-lemma-map-approximation}（またはその証明）を
もう一度適用する。この補題により、$R \to D$ がエタールで、
$\kappa(\mathfrak p_R) = \kappa(\mathfrak p_D)$ を満たし、かつ
$$
C \to D \to D^\wedge = R^\wedge
$$
が $(\mathfrak p_{R^\wedge})^N$ を法として $\sigma : C \to R^\wedge$ に等しいような
環準同型 $C \to D$ を得る。そこで
$$
V = \Spec(D)
\quad\text{and}\quad
W = \Spec(B \otimes_C D)
$$
を補題の問題の解に取れる。実際、図式
$$
\xymatrix{
A \ar[r] &
B \otimes_{C, \sigma} R^\wedge \ar[r] &
B \otimes_{C, \sigma} R^\wedge/(\mathfrak p_{R^\wedge})^N \ar@{=}[r] &
B \otimes_C D/(\mathfrak p_D)^N \\
A \ar@{=}[u] \ar[r] &
A \otimes_R D \ar[r] &
B \otimes_R D \ar[r] &
B \otimes_C D/(\mathfrak p_D)^N \ar@{=}[u]
}
$$
は、$C \to D \to D^\wedge = R^\wedge$ が
$(\mathfrak p_{R^\wedge})^N$ を法として $\sigma$ に等しいので可換である。
これで (5) が証明され、他の性質は構成から直ちに従う。
\end{proof}

\begin{lemma}
\label{more-morphisms-lemma-control-agreement}
Noetherian スキームの有限型射 $T \to S$ を考える。
$t \in T$ は $s \in S$ へ写るとし、
$\sigma : \mathcal{O}_{T, t} \to \mathcal{O}_{S, s}^\wedge$
を局所 $\mathcal{O}_{S, s}$-代数準同型とする。各 $N \geq 1$ に対して、有限型射
$(T', t') \to (T, t)$ で、$\sigma$ が
$\mathcal{O}_{T, t} \to \mathcal{O}_{T', t'}$ を経由し、かつ
$\mathcal{O}_{T, t} \to \mathcal{O}_{T', t'}$ を経由する任意の局所
$\mathcal{O}_{S, s}$-代数準同型
$\sigma' : \mathcal{O}_{T, t} \to \mathcal{O}_{S, s}^\wedge$
に対して、$\sigma$ と $\sigma'$ が $\mathfrak m_s^N$ を法として一致するものが存在する。
\end{lemma}

\begin{proof}
$S$ と $T$ はアフィンであると仮定してよい。$S = \Spec(R)$、
$T = \Spec(C)$ とする。$c_1, \ldots, c_n \in C$ を $R$-代数としての $C$ の
生成元とする。$\mathfrak p \subset R$ を $s$ に対応する素イデアルとし、
$\mathfrak p = (f_1, \ldots, f_m)$ とする。$R$ を単項局所化で置き換えて
（$R_\mathfrak p$ の分母を払って）、$r_1, \ldots, r_n \in R$ と
$a_{i, I} \in \mathcal{O}_{S, s}^\wedge$ が存在すると仮定してよい。ここで
$I = (i_1, \ldots, i_m)$、$\sum i_j = N$ であり、
$$
\sigma(c_i) = r_i + \sum\nolimits_I a_{i, I} f_1^{i_1} \ldots f_m^{i_m}
$$
が $\mathcal{O}_{S, s}^\wedge$ で成り立つ。そこで
$$
C' = C[t_{i, I}]/
\left(c_i - r_i - \sum\nolimits_I t_{i, I} f_1^{i_1} \ldots f_m^{i_m}\right)
$$
を考える。
% P05-EN-CH37-0182: the frozen pointed prime pC'+(t) need not be compatible with t_{i,I} mapping to arbitrary coefficients a_{i,I}; see the sidecar.
$\mathfrak p' = \mathfrak pC' + (t_{i, I})$ とし、$t_{i, I}$ を $a_{i, I}$ へ写すことにより、
$\sigma : C \to \mathcal{O}_{S, s}^\wedge$ の $C'$ を経由する分解を与える。
$T' = \Spec(C')$ と取ればよい。実際、補題の主張にある任意の $\sigma'$ は、
$c_i$ を極大イデアルの $N$ 乗を法として $r_i$ へ写すからである。
\end{proof}

\begin{lemma}
\label{more-morphisms-lemma-control-graded}
Noetherian スキームの有限型射 $Y \to T \to S$ を考える。
$t \in T$ は $s \in S$ へ写るとし、
$\sigma : \mathcal{O}_{T, t} \to \mathcal{O}_{S, s}^\wedge$
を局所 $\mathcal{O}_{S, s}$-代数準同型とする。有限型射
$(T', t') \to (T, t)$ で、$\sigma$ が
$\mathcal{O}_{T, t} \to \mathcal{O}_{T', t'}$ を経由し、かつ
$\mathcal{O}_{T, t} \to \mathcal{O}_{T', t'}$ を経由する任意の局所
$\mathcal{O}_{S, s}$-代数準同型
$\sigma' : \mathcal{O}_{T, t} \to \mathcal{O}_{S, s}^\wedge$
に対して、閉埋め込み
$$
Y \times_{T, \sigma} \Spec(\mathcal{O}_{S, s}^\wedge) = Y_\sigma
\longleftarrow Y_t \longrightarrow
Y_{\sigma'} =
Y \times_{T, \sigma'} \Spec(\mathcal{O}_{S, s}^\wedge)
$$
の余法代数が同型となるものが存在する。
\end{lemma}

\begin{proof}
$\sigma$ の存在から $\kappa(s) = \kappa(t)$ であることは有用な注意である。
$s \in \Spec(\mathcal{O}_{S, s}^\wedge)$ 上の $Y_\sigma$ と $Y_{\sigma'}$ の
ファイバーはいずれも $Y_t$ と標準的に同型なので、主張は意味をもつ。
「$\sigma'$ が $\mathcal{O}_{T, t} \to \mathcal{O}_{T', t'}$ を経由する」という性質を
$\sigma'$ に対する制約と考える。このような制約が
$(T'_i, t'_i) \to (T, t)$、$i = 1, \ldots, n$ によって複数与えられていれば、
$(T'_1 \times_T \ldots \times_T T'_n, (t'_1, \ldots, t'_n)) \to
(T, t)$ を考えることでそれらを組み合わせられる。以下ではこれを断りなく用いる。

\medskip\noindent
Lemma \ref{more-morphisms-lemma-control-agreement} により、補題の主張にある任意の $\sigma'$ は
$\mathfrak m_s^2$ を法として $\sigma$ と同じであると仮定してよい。
$Y_\sigma$ における $Y_t$ の余法代数は、準連接次数付き $\mathcal{O}_{Y_t}$-代数
$$
\bigoplus\nolimits_{n \geq 0}
\mathfrak m_s^n\mathcal{O}_{Y_\sigma}/
\mathfrak m_s^{n + 1}\mathcal{O}_{Y_\sigma}
$$
にほかならず、$Y_{\sigma'}$ についても同様である。$\sigma$ と $\sigma'$ は
$\mathfrak m_s^2$ を法として一致するので、これら二つの代数は次数 $0$ と $1$ で同じである。
一方、これらの余法代数は次数 $0$ 上、次数 $1$ で生成される。したがって、いま次数
$0$ と $1$ で構成した同型を延長する同型が存在すれば、それは一意である。

\medskip\noindent
$S$ と $T$ はアフィンであると仮定してよい。$Y = Y_1 \cup \ldots \cup Y_n$ を
アフィン開被覆とする。各 $Y_i \to T \to S$ と $\sigma$ に対して所望の同型
（前段落を参照されたい）が存在するような、補題にある
$(T_i', t'_i) \to (T, t)$ を構成できれば、それらは一意性により貼り合わさり、
$Y \to T$ に対する結果を与える。したがって $Y$ はアフィンであると仮定してよい。

\medskip\noindent
$S = \Spec(R)$、$T = \Spec(C)$、$Y = \Spec(B)$ と書く。
表示 $B = C[x_1, \ldots, x_n]/(f_1, \ldots, f_m)$ を選ぶ。
$R^\wedge = \mathcal{O}_{S, s}^\wedge$ と書く。多項式
$a_{kj} \in R^\wedge[x_1, \ldots, x_n]$ で、
$$
\sum\nolimits_{j = 1, \ldots, m} a_{kj}\sigma(f_j) = 0,\quad
\text{for }k = 1, \ldots, K
$$
が $\sigma(f_j) \in R^\wedge[x_1, \ldots, x_n]$ の間の関係加群の生成元系となるものを取る。
したがって完全列
\begin{equation}
\label{more-morphisms-equation-resolution}
R^\wedge[x_1, \ldots, x_n]^{\oplus K} \to
R^\wedge[x_1, \ldots, x_n]^{\oplus m} \to
R^\wedge[x_1, \ldots, x_n] \to B \otimes_{C, \sigma} R^\wedge \to 0
\end{equation}
がある。$c$ を、この列の第一および第二の写像と、More on Algebra, Section
\ref{more-algebra-section-artin-rees} で定義されたイデアル
$\mathfrak m_{R^\wedge}R^\wedge[x_1, \ldots, x_n]$ に対する Artin--Rees の補題で
使える整数とする。多重指数記法で
$$
a_{kj} = \sum\nolimits_{I \in \Omega} a_{kj, I} x^I
\quad\text{and}\quad
f_j = \sum\nolimits_{I \in \Omega} f_{j, I} x^I
$$
と書く。ここで $a_{kj, I} \in R^\wedge$、$f_{j, I} \in C$ であり、$\Omega$ は
多重指数の有限集合である。このとき $R^\wedge$ において
$$
\sum\nolimits_{j = 1, \ldots, m,\ I, I' \in \Omega,\ I + I' = I''}
a_{kj, I} \sigma(f_{j, I'}) = 0,\quad
I''\text{ a multiindex}
$$
である。そこで
% P05-EN-CH37-0183: the frozen coefficient construction interchanges the j,k indices in its variable family and specialization; see the sidecar.
$$
C' = C[t_{jk, I}]/
\left(
\sum\nolimits_{j = 1, \ldots, m,\ I, I' \in \Omega,\ I + I' = I''}
t_{kj, I} f_{j, I'},\ I''\text{ a multiindex}\right)
$$
と取る。このとき $\sigma$ は $t_{kj, I}$ を $a_{jk, I}$ へ写す写像
$\tilde\sigma : C' \to R^\wedge$ を経由する。したがって $T' = \Spec(C')$ は、
$\sigma$ が $\mathcal{O}_{T, t} \to \mathcal{O}_{T', t'}$ を経由するような点
$t' \in T'$ をもつ。$C'[x_1, \ldots, x_n]$ において
$t_{kj} = \sum t_{kj, I} x^I$ と置く。このとき複体
\begin{equation}
\label{more-morphisms-equation-resolution-new}
C'[x_1, \ldots, x_n]^{\oplus K} \to
C'[x_1, \ldots, x_n]^{\oplus m} \to
C'[x_1, \ldots, x_n] \to B \otimes_C C' \to 0
\end{equation}
があり、これは $C'[x_1, \ldots, x_n]$ において完全で、$\tilde\sigma$ による基底変換は
(\ref{more-morphisms-equation-resolution}) を与える。

\medskip\noindent
Lemma \ref{more-morphisms-lemma-control-agreement} により、さらに射
$(T'', t'') \to (T', t')$ で、$\tilde\sigma$ が
$\mathcal{O}_{T', t'} \to \mathcal{O}_{T'', t''}$ を経由し、かつ
$\sigma' : C \to R^\wedge$ が $\mathcal{O}_{T'', t''}$ を経由すれば、誘導される写像
$\tilde \sigma' : C' \to R^\wedge$ が $\mathfrak m_s^{c + 1}$ を法として
$\tilde \sigma$ と一致するものを見つけることができる。したがって $\sigma'$ が
そのような写像ならば、$R^\wedge[x_1, \ldots, x_n]$ 上の複体
$$
R^\wedge[x_1, \ldots, x_n]^{\oplus K} \to
R^\wedge[x_1, \ldots, x_n]^{\oplus m} \to
R^\wedge[x_1, \ldots, x_n] \to B \otimes_{C, \sigma'} R^\wedge \to 0
$$
を、多項式 $t_{kj}$ と $f_j$ に $\tilde\sigma'$ を適用することによって得る。
言い換えれば、これは複体 (\ref{more-morphisms-equation-resolution-new}) の $\tilde\sigma'$ による
基底変換である。$\tilde \sigma$ と $\tilde \sigma'$ は
$\mathfrak m_s^{c + 1}$ を法として合同なので、この複体を定める行列は、複体
(\ref{more-morphisms-equation-resolution}) を定める行列と $\mathfrak m_s^{c + 1}$ を法として合同である。
(\ref{more-morphisms-equation-resolution}) は完全なので、More on Algebra, Lemma
\ref{more-algebra-lemma-approximate-complex-graded} を適用して
$$
\text{Gr}_{\mathfrak m_s}(B \otimes_{C, \sigma'} R^\wedge)
\cong
\text{Gr}_{\mathfrak m_s}(B \otimes_{C, \sigma} R^\wedge)
$$
を得る。これは所望の結論である。
\end{proof}

\begin{lemma}
\label{more-morphisms-lemma-relative-isomorphism-approximation}
Lemma \ref{more-morphisms-lemma-relative-map-approximation} の記法と仮定のもとで、
$\varphi$ は完備化上の同型を誘導すると仮定する。このとき $f$ がエタールとなるように
図式を選ぶことができる。
\end{lemma}

\begin{proof}
$N \geq 2$ と仮定してよく、Lemma \ref{more-morphisms-lemma-control-graded} にあるように
$(T, t)$ を $(T', t')$ で置き換えてよい。$(V, v) \to (S, s)$ は初等エタール近傍なので、
$(X \times_S V, (x, v)) \to (X, x)$ もそうである。したがって More on Algebra,
Lemma \ref{more-algebra-lemma-flat-unramified} により、
$\mathcal{O}_{X, x} \to \mathcal{O}_{X \times_S V, (x, v)}$ は完備化上の同型を誘導する。
$\mathcal{O}_{X, x} \to \mathcal{O}_{W, w}$ も完備化上の同型を誘導すると主張する。
これが示されれば、Lemma \ref{more-morphisms-lemma-lifting-along-artinian-at-point} により $f$ は
$w$ において滑らかであり、もちろん
% P05-EN-CH37-0184: the frozen proof says f is unramified at the undefined point u rather than at w; see the sidecar.
$f$ は $u$ において不分岐でもあるから、Morphisms, Lemma
\ref{morphisms-lemma-etale-smooth-unramified} により $f$ は $w$ においてエタールである。

\medskip\noindent
まず Lemma \ref{more-morphisms-lemma-relative-map-approximation} の (5) の可換性を用いると、
$i \leq N$ に対して可換図式
$$
\xymatrix{
\text{Gr}^i_{\mathfrak m_x}(\mathcal{O}_{X, x})
\ar[r]_-\varphi &
\text{Gr}^i_{\mathfrak m_{y_\sigma}}(\mathcal{O}_{Y_\sigma, y_\sigma}^\wedge)
\ar@{=}[r] &
\text{Gr}^i_{\mathfrak m_{(y, v)}}(\mathcal{O}_{Y \times_{T, \tau} V, (y, v)})
\ar[d]_{\cong} \\
\text{Gr}^i_{\mathfrak m_x}(\mathcal{O}_{X, x})
\ar[r]^-{\cong} \ar@{=}[u] &
\text{Gr}^i_{\mathfrak m_{(x, v)}}(\mathcal{O}_{X \times_S V, (x, v)})
\ar[r]^{f^\sharp_w} &
\text{Gr}^i_{\mathfrak m_w}(\mathcal{O}_{W, w})
}
$$
がある。これは、$f^\sharp_w$ が剰余体上の同型 $\kappa(x) \to \kappa(w)$ と
余接空間上の同型
$\mathfrak m_x/\mathfrak m_x^2 \to \mathfrak m_w/\mathfrak m_w^2$ を定めることを意味する。
したがって $f^\sharp_w$ は完備局所環の全射
$\mathcal{O}_{X, x}^\wedge \to \mathcal{O}_{W, w}^\wedge$ を定める。

\medskip\noindent
Lemma \ref{more-morphisms-lemma-control-graded} により同型
% P05-EN-CH37-0185: the frozen local-ring subscripts have inconsistent and once unbalanced outer parentheses; see the sidecar.
$\text{Gr}_{\mathfrak m_s}(\mathcal{O}_{(Y \times_{T, \tau} V, (y, v)})$
と
$\text{Gr}_{\mathfrak m_s}(\mathcal{O}_{Y_\sigma, y_\sigma})$
がある。これは余法層の同型の点 $y$ における茎を取ることで従う。
考えている局所環は Noetherian なので、$\mathfrak m_s$ に関する次数付き随伴環を取ることは
完備化と可換である。実際、イデアルに関する完備化は Noetherian 環上の有限加群に対する
完全関手だからである。これにより、次数付き環の図式
$$
\xymatrix{
\text{Gr}_{\mathfrak m_s}(\mathcal{O}_{W, w}^\wedge) &
\text{Gr}_{\mathfrak m_s}
(\mathcal{O}_{(Y \times_{T, \tau} V, (y, v)}^\wedge) \ar[l] \\
\text{Gr}_{\mathfrak m_s}(\mathcal{O}_{X, x}^\wedge)
\ar[r]^\varphi \ar[u] &
\text{Gr}_{\mathfrak m_s}(\mathcal{O}_{Y_\sigma, y_\sigma}^\wedge)
\ar[u]_{\cong}
}
$$
の右側の縦の同型を得る。この図式が可換であるとは主張しない。前段落の結果により、
左の矢は全射である。他の三つの矢は同型である。したがって左の矢は、互いに同型な
Noetherian 環の間の全射である。ゆえに Algebra, Lemma
\ref{algebra-lemma-surjective-endo-noetherian-ring-is-iso} により同型である
（Hilbert 関数を用いて直接論じることもできる）。特に
$\mathcal{O}_{X, x}^\wedge \to \mathcal{O}_{W, w}^\wedge$ は全射であるだけでなく
単射でもあり、証明が完了する。
\end{proof}















\section{有限局所自由射によるエタール被覆の局所的支配}
\label{more-morphisms-section-finite-free-over-etale}

\noindent
本節では、大まかに言えば、スキーム $S$ のエタール被覆は、$S$ の有限局所自由被覆の
Zariski 被覆によって細分できるという結果を説明する。

\begin{lemma}
\label{more-morphisms-lemma-dominate-etale-neighbourhood-finite-flat}
$S$ をスキームとし、$s \in S$ とする。
$f : (U, u) \to (S, s)$ をエタール近傍とする。このとき、アフィン開近傍
$s \in V \subset S$ と全射な有限局所自由射 $\pi : T \to V$ で、各
$t \in \pi^{-1}(s)$ に対して開近傍 $t \in W_t \subset T$ と可換図式
$$
\xymatrix{
T \ar[d]^\pi & W_t \ar[l] \ar[rr]_{h_t} \ar[rd] & & U \ar[dl] \\
V \ar[rr] & & S
}
$$
で $h_t(t) = u$ を満たすものが存在する。
\end{lemma}

\begin{proof}
問題は $S$ 上局所的なので、$S$ を $s$ の任意の開近傍で置き換えてよい。
$U$ も $u$ の開近傍で置き換えてよい。したがって Morphisms, Lemma
\ref{morphisms-lemma-etale-locally-standard-etale} により、$U \to S$ は
アフィン・スキームの標準エタール射であると仮定してよい。この場合、補題は
（$V = S$ として）Algebra, Lemma
\ref{algebra-lemma-standard-etale-finite-flat-Zariski} から従う。
\end{proof}

\begin{lemma}
\label{more-morphisms-lemma-dominate-etale-affine-finite-flat}
$f : U \to S$ をアフィン・スキームの全射エタール射とする。このとき、全射な
有限局所自由射 $\pi : T \to S$ と有限開被覆 $T = T_1 \cup \ldots \cup T_n$ で、
各 $T_i \to S$ が $U \to S$ を経由するものが存在する。図で表すと
$$
\xymatrix{
& \coprod T_i  \ar[rd] \ar[ld] & \\
T \ar[rd]^\pi & & U \ar[ld]_f \\
& S &
}
$$
であり、南西向きの矢は Zariski 被覆である。
\end{lemma}

\begin{proof}
これは Algebra, Lemma \ref{algebra-lemma-etale-finite-flat-zariski} の言い換えである。
\end{proof}

\begin{remark}
\label{more-morphisms-remark-topologies}
位相の言葉では、Lemmas \ref{more-morphisms-lemma-dominate-etale-neighbourhood-finite-flat} と
\ref{more-morphisms-lemma-dominate-etale-affine-finite-flat} は次を意味する。$S$ を任意のスキームとし、
$\{f_i : U_i \to S\}$ を $S$ のエタール被覆とする。このとき Zariski 開被覆
$S = \bigcup V_j$、各 $j$ に対する全射な有限局所自由射 $W_j \to V_j$、および
各 $j$ に対する Zariski 開被覆 $\{W_{j, k} \to W_j\}$ で、族
$\{W_{j, k} \to S\}$ が与えられたエタール被覆
$\{f_i : U_i \to S\}$ を細分するものが存在する。これは実際には何を意味するだろうか。
たとえば、Zariski 被覆と、$\pi$ が全射かつ有限局所自由であるような形
$\{\pi : T \to S\}$ の fppf 被覆について解けることが分かっている降下問題を考える。
このとき、この降下問題はエタール被覆についても肯定的な答えをもつ。この技巧は、
アフィン・スキーム $X$ に対する $\text{Br}(X) = \text{Br}'(X)$ の Gabber の証明で
用いられた。\cite{Hoobler} を参照されたい。
\end{remark}




\section{準有限射のエタール局所化}
\label{more-morphisms-section-etale-localization}

\noindent
ここでは「準有限射はエタール局所化の後に有限になる」という主題をめぐる
一連の補題を扱う。一般的な考え方は次のとおりである。スキームの射
$f : X \to S$ と点 $s \in S$ が与えられたとする。
$\varphi : (U, u) \to (S, s)$ を $S$ における $s$ の
エタール近傍とする。ファイバー積 $X_U = U \times_S X$ と基本図式
\begin{equation}
\label{more-morphisms-equation-basic-diagram}
\vcenter{
\xymatrix{
V \ar[r] \ar[dr] & X_U \ar[d] \ar[r] & X \ar[d]^f \\
& U \ar[r]^\varphi & S
}
}
\end{equation}
を考える。ここで $V \subset X_U$ は開である。
射 $f_U : X_U \to U$、または適当な開部分 $V$ に対する射
$V \to U$ に標準的なモデルはあるだろうか。
もちろん一般には答えは否定的である。しかし準有限射については
一定のことがいえる。

\begin{lemma}
\label{more-morphisms-lemma-etale-makes-quasi-finite-finite-at-point}
$f : X \to S$ をスキームの射とする。
$x \in X$ とし、$s = f(x)$ とおく。次を仮定する。
\begin{enumerate}
\item $f$ は局所有限型であり、
\item $x \in X_s$ は孤立点である\footnote{(1) の下では、これは
$f$ が $x$ において準有限であることを意味する。Morphisms,
Lemma \ref{morphisms-lemma-quasi-finite-at-point-characterize} を参照。}。
\end{enumerate}
このとき、次が存在する。
\begin{enumerate}
\item[(a)] 初等エタール近傍 $(U, u) \to (S, s)$、
\item[(b)] 開部分スキーム $V \subset X_U$
（\ref{more-morphisms-equation-basic-diagram} を参照）
\end{enumerate}
これらは次を満たす。
\begin{enumerate}
\item[(\romannumeral1)] $V \to U$ は有限射である。
\item[(\romannumeral2)] $U$ の $u$ に写る $V$ の点 $v$ が一意に存在する。
\item[(\romannumeral3)] 点 $v$ は射 $X_U \to X$ によって $x$ に写り、
$\kappa(x) = \kappa(v)$ を誘導する。
\end{enumerate}
さらに、任意の初等エタール近傍 $(U', u') \to (U, u)$ に対し、
$V' = U' \times_U V \subset X_{U'}$ とおけば、三つ組 $(U', u', V')$
も性質 (\romannumeral1)、(\romannumeral2)、(\romannumeral3) を満たす。
\end{lemma}

\begin{proof}
$f(Y) \subset W$ かつ $x \in Y$ となるアフィン開部分
$Y \subset X$、$W \subset S$ をとる。点 $x$ は射
$f|_Y : Y \to W$ のファイバーにおいても孤立点である。
$f|_Y : Y \to W$ に対して定理を証明できれば、$f$ に対する定理も従う。
したがって $f$ がアフィンスキームの射である場合に帰着される。この場合は
Algebra, Lemma \ref{algebra-lemma-etale-makes-quasi-finite-finite-one-prime}
である。
\end{proof}

\noindent
前の補題および次の補題では、射 $f$ が分離的であることを仮定しない。
したがって、そこで構成される開部分 $V$、$V_i$ は
$X_U$ において閉であるとは限らない。さらに、近傍 $U$ をアフィンに
選べば各 $V_i$ はアフィンであるが、非分離の場合には交わり
$V_i \cap V_j$ がアフィンであるとは限らない。

\begin{lemma}
\label{more-morphisms-lemma-etale-makes-quasi-finite-finite-multiple-points}
$f : X \to S$ をスキームの射とする。
$x_1, \ldots, x_n \in X$ を、$S$ において同じ像 $s$ をもつ点とする。
次を仮定する。
\begin{enumerate}
\item $f$ は局所有限型であり、
\item $i = 1, \ldots, n$ に対して $x_i \in X_s$ は孤立点である。
\end{enumerate}
このとき、次が存在する。
\begin{enumerate}
\item[(a)] 初等エタール近傍 $(U, u) \to (S, s)$、
\item[(b)] 各 $i$ に対する開部分スキーム $V_i \subset X_U$。
\end{enumerate}
各 $i$ に対して次が成り立つ。
\begin{enumerate}
\item[(\romannumeral1)] $V_i \to U$ は有限射である。
\item[(\romannumeral2)] $U$ の $u$ に写る $V_i$ の点 $v_i$ が一意に存在する。
\item[(\romannumeral3)] 点 $v_i$ は $X$ の $x_i$ に写り、
$\kappa(x_i) = \kappa(v_i)$ である。
\end{enumerate}
\end{lemma}

\begin{proof}
$n$ に関する帰納法を用いる。すなわち、$(U, u) \to (S, s)$ および
$V_i \subset X_U$、$i = 1, \ldots, n - 1$ が
$x_1, \ldots, x_{n - 1}$ に対して機能すると仮定する。
$\kappa(s) = \kappa(u)$ であるから、ファイバー $(X_U)_u = X_s$ である。
% P05-EN-CH37-0186
したがって $x_n \in X_s$ に対応する点
$x'_n \in X_u \subset X_U$ が一意に存在する。また $x'_n$ は
$X_u$ において孤立している。ゆえに
Lemma \ref{more-morphisms-lemma-etale-makes-quasi-finite-finite-at-point} により、
初等エタール近傍 $(U', u') \to (U, u)$ と、
$x'_n$ に対して、したがって $x_n$ に対して機能する開部分
$V_n \subset X_{U'}$ が存在する。
Lemma \ref{more-morphisms-lemma-etale-makes-quasi-finite-finite-at-point} の最後の主張により、
$i = 1, \ldots, n - 1$ に対する開部分スキーム
$V'_i = U'\times_U V_i$ は、$x_1, \ldots, x_{n - 1}$ に関して
引き続き機能する。これで示された。
\end{proof}

\noindent
非自明な体拡大 $\kappa(u)/\kappa(s)$、すなわち一般のエタール近傍を
許せば、次のように点を分裂させることができる。

\begin{lemma}
\label{more-morphisms-lemma-etale-makes-quasi-finite-finite-multiple-points-var}
$f : X \to S$ をスキームの射とする。
$x_1, \ldots, x_n \in X$ を、$S$ において同じ像 $s$ をもつ点とする。
次を仮定する。
\begin{enumerate}
\item $f$ は局所有限型であり、
\item $i = 1, \ldots, n$ に対して $x_i \in X_s$ は孤立点である。
\end{enumerate}
このとき、次が存在する。
\begin{enumerate}
\item[(a)] エタール近傍 $(U, u) \to (S, s)$、
\item[(b)] 各 $i$ に対する整数 $m_i$ と、$j = 1, \ldots, m_i$ に対する
開部分スキーム $V_{i, j} \subset X_U$。
\end{enumerate}
これらは次を満たす。
\begin{enumerate}
\item[(\romannumeral1)] 各 $V_{i, j} \to U$ は有限射である。
\item[(\romannumeral2)] $U$ の $u$ に写る $V_{i, j}$ の点
$v_{i, j}$ が一意に存在し、$\kappa(u) \subset \kappa(v_{i, j})$
は有限純非分離である。
% P05-EN-CH37-0187
\item[(\romannumeral4)] $v_{i, j} = v_{i', j'}$ ならば $i = i'$ かつ
$j = j'$ である。
\item[(\romannumeral3)] 点 $v_{i, j}$ は $X$ の $x_i$ に写り、
$(X_U)_u$ のほかの点は $x_i$ に写らない。
\end{enumerate}
\end{lemma}

\begin{proof}
この証明は、Algebra,
Lemma \ref{algebra-lemma-etale-makes-quasi-finite-finite-variant}
の証明をスキームの言葉で述べた変形である。
Morphisms, Lemma \ref{morphisms-lemma-quasi-finite-at-point-characterize}
により、射 $f$ は各点 $x_i$ において準有限である。
したがって各 $i$ に対して $\kappa(s) \subset \kappa(x_i)$ は有限である
（Morphisms, Lemma \ref{morphisms-lemma-residue-field-quasi-finite}）。
各 $i$ に対し、$L_i/\kappa(s)$ が分離的であり、
$\kappa(x_i)/L_i$ が純非分離となるような部分体
$\kappa(s) \subset L_i \subset \kappa(x_i)$ をとる。
$i = 1, \ldots, n$ に対して $\kappa(s)$-埋め込み $L_i \to L$
が存在するような有限 Galois 拡大 $L/\kappa(s)$ を選ぶ。
$\kappa(s)$-拡大として $L \cong \kappa(u)$ となるエタール近傍
$(U, u) \to (S, s)$ を選ぶ
（Lemma \ref{more-morphisms-lemma-realize-prescribed-residue-field-extension-etale}）。

\medskip\noindent
$y_{i, j}$、$j = 1, \ldots, m_i$ を、$x_i \in X$ と $u \in U$ の
上にある $X_U$ の点とする。Schemes,
Lemma \ref{schemes-lemma-points-fibre-product} により、これらの点
$y_{i, j}$ は環 $\kappa(u) \otimes_{\kappa(s)} \kappa(x_i)$ の
素イデアルとちょうど対応する。これにより点が有限個であることも分かる。
実際には $m_i = [L_i : \kappa(s)]$ であるが、これは必要ない。
$L$ の選び方と初等的な体論により、各組 $i, j$ に対して
$\kappa(u) \subset \kappa(y_{i, j})$ は有限純非分離である。
また、例えば Morphisms,
Lemma \ref{morphisms-lemma-base-change-quasi-finite} により、射
$X_U \to U$ はすべての $i, j$ に対して点 $y_{i, j}$ において
準有限である。

\medskip\noindent
射 $X_U \to U$、点 $u \in U$、および点
$y_{i, j} \in (X_U)_u$ に
Lemma \ref{more-morphisms-lemma-etale-makes-quasi-finite-finite-multiple-points}
を適用する。これにより、$\kappa(u) = \kappa(u')$ を満たす
エタール近傍 $(U', u') \to (U, u)$ と、その補題の性質
(\romannumeral1)、(\romannumeral2)、(\romannumeral3) を満たす開部分
$V_{i, j} \subset X_{U'}$ が得られる。エタール近傍
$(U', u') \to (S, s)$ と開部分 $V_{i, j} \subset X_{U'}$ が
本補題の問題の解であると主張する。確認は省略する。
\end{proof}

\begin{lemma}
\label{more-morphisms-lemma-etale-splits-off-quasi-finite-part-technical}
$f : X \to S$ をスキームの射とする。
$s \in S$ とし、$x_1, \ldots, x_n \in X_s$ とする。次を仮定する。
\begin{enumerate}
\item $f$ は局所有限型である。
\item $f$ は分離的である。
\item $x_1, \ldots, x_n$ は $X_s$ の相異なる孤立点である。
\end{enumerate}
このとき、初等エタール近傍 $(U, u) \to (S, s)$ と分解
$$
U \times_S X = W \amalg V_1 \amalg \ldots \amalg V_n
$$
が存在する。各項は開かつ閉な部分スキームであり、射
$V_i \to U$ は有限、$u$ 上の $V_i \to U$ のファイバーは
一点集合 $\{v_i\}$、各 $v_i$ は $x_i$ に写って
$\kappa(x_i) = \kappa(v_i)$ となり、$u$ 上の $W \to U$ の
ファイバーは、どの $x_i$ に写る点も含まない。
\end{lemma}

\begin{proof}
Lemma \ref{more-morphisms-lemma-etale-makes-quasi-finite-finite-multiple-points}
のように $(U, u) \to (S, s)$ と $V_i \subset X_U$ を選ぶ。
$X_U \to U$ は分離的であり
（Schemes, Lemma \ref{schemes-lemma-separated-permanence}）、
$V_i \to U$ は有限、したがって固有である
（Morphisms, Lemma \ref{morphisms-lemma-finite-proper}）。
ゆえに Morphisms,
Lemma \ref{morphisms-lemma-image-proper-scheme-closed} により
$V_i \subset X_U$ は閉である。したがって $V_i \cap V_j$ は
$v_i$ を含まない $V_i$ の閉部分集合である。よって
$U$ における $V_i \cap V_j$ の像は、$V_i \to U$ が固有であるため
閉集合であり、$u$ を含まない。したがって $U$ を縮小すれば、
% P05-EN-CH37-0188
すべての $i, j$ に対して $V_i \cap V_j = \emptyset$ と仮定できる。
これが補題の分解を与える。
\end{proof}

\noindent
次は、純非分離体拡大に帰着する変形である。

\begin{lemma}
\label{more-morphisms-lemma-etale-splits-off-quasi-finite-part-technical-variant}
$f : X \to S$ をスキームの射とする。
$s \in S$ とし、$x_1, \ldots, x_n \in X_s$ とする。次を仮定する。
\begin{enumerate}
\item $f$ は局所有限型である。
\item $f$ は分離的である。
\item $x_1, \ldots, x_n$ は $X_s$ の相異なる孤立点である。
\end{enumerate}
このとき、エタール近傍 $(U, u) \to (S, s)$ と分解
$$
U \times_S X =
W \amalg
\ \coprod\nolimits_{i = 1, \ldots, n}
\ \coprod\nolimits_{j = 1, \ldots, m_i}
V_{i, j}
$$
が存在する。各項は開かつ閉な部分スキームであり、射
$V_{i, j} \to U$ は有限、$u$ 上の $V_{i, j} \to U$ のファイバーは
一点集合 $\{v_{i, j}\}$、各 $v_{i, j}$ は $x_i$ に写り、
$\kappa(u) \subset \kappa(v_{i, j})$ は純非分離であり、
$u$ 上の $W \to U$ のファイバーは、どの $x_i$ に写る点も含まない。
\end{lemma}

\begin{proof}
証明は Lemma \ref{more-morphisms-lemma-etale-splits-off-quasi-finite-part-technical}
の証明とまったく同じである。ただし
Lemma \ref{more-morphisms-lemma-etale-makes-quasi-finite-finite-multiple-points}
の代わりに
Lemma \ref{more-morphisms-lemma-etale-makes-quasi-finite-finite-multiple-points-var}
を用いる。
\end{proof}

\noindent
次の形の方がいくらか読み取りやすいかもしれない。

\begin{lemma}
\label{more-morphisms-lemma-etale-splits-off-quasi-finite-part}
$f : X \to S$ をスキームの射とする。
$s \in S$ とし、次を仮定する。
\begin{enumerate}
\item $f$ は局所有限型である。
\item $f$ は分離的である。
\item $X_s$ の孤立点は高々有限個である。
\end{enumerate}
このとき、初等エタール近傍 $(U, u) \to (S, s)$ と分解
$$
U \times_S X = W \amalg V
$$
が存在する。両項は開かつ閉な部分スキームであり、射
$V \to U$ は有限で、射 $W \to U$ のファイバー $W_u$ は
孤立点を含まない。特に、$f^{-1}(s)$ が有限集合ならば
$W_u = \emptyset$ である。
\end{lemma}

\begin{proof}
$x_1, \ldots, x_n$ を $X_s$ の孤立点すべての集合として選び、
$V = \bigcup V_i$ とおけば、
Lemma \ref{more-morphisms-lemma-etale-splits-off-quasi-finite-part-technical}
から明らかである。
\end{proof}






\section{整射のエタール局所化}
\label{more-morphisms-section-etale-localization-integral}

\noindent
Section \ref{more-morphisms-section-etale-localization} の結果について、
整射の場合のいくつかの変形を与える。

\begin{lemma}
\label{more-morphisms-lemma-etale-makes-integral-split}
$R \to S$ を整環準同型とする。$\mathfrak p \subset R$ を
素イデアルとする。次を仮定する。
\begin{enumerate}
\item $\mathfrak p$ の上にある素イデアル
$\mathfrak q_1, \ldots, \mathfrak q_n$ は有限個である。
\item 各 $i$ に対し、最大分離部分拡大
% P05-EN-CH37-0189
$\kappa(\mathfrak q)/\kappa(\mathfrak q_i)_{sep}/\kappa(\mathfrak p)$
（Fields, Lemma \ref{fields-lemma-separable-first}）は
$\kappa(\mathfrak p)$ 上有限である。
\end{enumerate}
このとき、エタール環準同型 $R \to R'$ と、
$\mathfrak p$ の上にある素イデアル $\mathfrak p'$ が存在し、
$$
S \otimes_R R' = A_1 \times \ldots \times A_m
$$
となる。ここで $R' \to A_j$ は整であり、$\mathfrak p'$ の上にある
素イデアル $\mathfrak r_j$ を一意にもって、
$\kappa(\mathfrak r_j)/\kappa(\mathfrak p')$ は純非分離である。
\end{lemma}

\begin{proof}[第一の証明]
この証明では
Algebra, Lemma \ref{algebra-lemma-etale-makes-quasi-finite-finite-variant}.
を用いる。すなわち、この体の $\kappa(\mathfrak p)$ 上の生成元
$\theta_i \in \kappa(\mathfrak q_i)_{sep}$ を選ぶ
（Fields, Lemma \ref{fields-lemma-primitive-element}）。
ファイバー環 $S \otimes_R \kappa(\mathfrak p)$ のスペクトルは
有限離散であり、その点は $\mathfrak q_1, \ldots, \mathfrak q_n$
に対応する。中国剰余定理
（Algebra, Lemma \ref{algebra-lemma-chinese-remainder}）により、
$S \otimes_R \kappa(\mathfrak p) \to \prod \kappa(\mathfrak q_i)$
は全射である。したがって、ある $g \in R$、
$g \not \in \mathfrak p$ に対して $R$ を $R_g$ で置き換えれば、
$(0, \ldots, 0, \theta_i, 0, \ldots, 0) \in \prod \kappa(\mathfrak q_i)$
がある $x_i \in S$ の像であると仮定できる。
$S' \subset S$ を、選んだ $x_i$ たちが生成する $R$-部分代数とする。
$\Spec(S) \to \Spec(S')$ は全射であるから
（Algebra, Lemma \ref{algebra-lemma-integral-overring-surjective}）、
$\mathfrak q_i' = S' \cap \mathfrak q_i$ は $\mathfrak p$ の上にある
$S'$ の素イデアルである。$x_i$ の選び方により、これらの素イデアルは
% P05-EN-CH37-0190
相異なり、かつ
$\kappa(\mathfrak q'_i)_{sep} = \kappa(\mathfrak q_i)_{sep}$.
である。特に体拡大
$\kappa(\mathfrak q_i)/\kappa(\mathfrak q'_i)$ は純非分離である。
$R \to S'$ は有限なので、
Algebra, Lemma \ref{algebra-lemma-etale-makes-quasi-finite-finite-variant}.
% P05-EN-CH37-0191
を適用でき、$R \to R'$、$\mathfrak p'$、および分解
$$
S' \otimes_R R' = A'_1 \times \ldots \times A'_m \times B'
$$
を得る。ここで $R' \to A'_j$ は整であり、$\mathfrak p'$ の上にある
素イデアル $\mathfrak r'_j$ を一意にもち、
$\kappa(\mathfrak r'_j)/\kappa(\mathfrak p')$ は純非分離であり、
$B'$ は $\mathfrak p'$ の上にある素イデアルをもたない。
$R \to S'$ が有限なので $R' \to B'$ も有限である。したがって、
ある $g' \in R'$、$g' \not \in \mathfrak p'$ によって
$R'$ を局所化した後、$B' = 0$ と仮定できる。
写像 $S' \otimes_R R' \to S \otimes_R R'$ を通じて、
$S$ に対する対応する分解を得る。
\end{proof}

\begin{proof}[第二の証明]
この証明では狭義 Hensel 化を用いる。まず、$R$ は極大イデアル
% P05-EN-CH37-0192
$\mathfrak p$ をもつ狭義 Hensel 化であると仮定する。
このとき $S/\mathfrak p S$ は
$\mathfrak q_1, \ldots, \mathfrak q_n$ に対応する有限個の
素イデアルをもち、それらはすべて極大で、剰余体は
$\kappa(\mathfrak p)$ 上純非分離である。したがって
% P05-EN-CH37-0193
$S/\mathfrak p S$ は $\prod (S/\mathfrak p S)_{\mathfrak p_i}$
に等しい。More on Algebra,
Lemma \ref{more-algebra-lemma-characterize-henselian-pair} により、
この直積分解を、主張にあるような $S$ の直積分解へ持ち上げられる。
% P05-EN-CH37-0194

\medskip\noindent
一般の場合には、$R^{sh}$ を $R_\mathfrak p$ の狭義 Hensel 化とする。
第一段落の結果を $R^{sh} \to S \otimes_R R^{sh}$ に適用できる。
直積分解に対応する $S \otimes_R R^{sh}$ の、互いに直交する
$m$ 個の冪等元を考える。Algebra,
Lemma \ref{algebra-lemma-strict-henselization-different} により、
$R^{sh}$ はエタール環準同型
$(R, \mathfrak p) \to (R', \mathfrak p')$ のフィルター付き余極限である。
したがって、これらの冪等元は所望のようにある $R'$ へ降下する。
\end{proof}







\section{Zariski の主定理}
\label{more-morphisms-section-application-etale-neighbourhoods}

\noindent
本節では、Grothendieck によって再定式化された Zariski の主定理を証明する。
% P05-EN-CH37-0195
この文脈で「Zariski の主定理」というとき、多くの場合には
% P05-EN-CH37-0196
Lemma \ref{more-morphisms-lemma-finite-type-separated}、
Lemma \ref{more-morphisms-lemma-quasi-finite-separated-quasi-affine}、または
Lemma \ref{more-morphisms-lemma-quasi-finite-separated-pass-through-finite}
のいずれかを意味する。大部分の文献では、このうち最後のものを
Zariski の主定理と呼ぶ。

\medskip\noindent
代数的な形はすでに
Algebra, Theorem \ref{algebra-theorem-main-theorem}
で証明し、さらにこの代数的な形をスキームの言葉で述べ直している。
Morphisms, Theorem \ref{morphisms-theorem-main-theorem} を参照。
本節の形はより微妙である。完全な結果を得るため、
Section \ref{more-morphisms-section-etale-localization} のエタール局所化の手法を用いて
代数的な場合に帰着する。

\begin{lemma}
\label{more-morphisms-lemma-finite-type-separated}
$f : X \to S$ をスキームの射とする。
$f$ は有限型かつ分離的であると仮定する。
$S'$ を $X$ における $S$ の正規化とする。Morphisms,
Definition \ref{morphisms-definition-normalization-X-in-Y} を参照。
図式は
$$
\xymatrix{
X \ar[rd]_f \ar[rr]_{f'} & & S' \ar[ld]^\nu \\
& S &
}
$$
である。このとき、次を満たす開部分スキーム $U' \subset S'$ が存在する。
\begin{enumerate}
\item $(f')^{-1}(U') \to U'$ は同型である。
\item $(f')^{-1}(U') \subset X$ は、$f$ が準有限となる点の集合である。
\end{enumerate}
\end{lemma}

\begin{proof}
Morphisms, Lemma \ref{morphisms-lemma-quasi-finite-points-open} により、
$f$ が準有限となる点の部分集合 $U \subset X$ は開である。
補題は次と同値である。
\begin{enumerate}
\item[(a)] $U' = f'(U) \subset S'$ は開である。
\item[(b)] $U = (f')^{-1}(U')$ である。
\item[(c)] $U \to U'$ は同型である。
\end{enumerate}
任意の $x \in U$ をとる。$(f')^{-1}V \to V$ が同型となる
開近傍 $f'(x) \in V \subset S'$ が存在すると主張する。
まず、この主張から補題が従うことを示す。実際、このとき
$(f')^{-1}V \cong V$ は $X$ の開部分スキームなので $S$ 上局所有限型であり、
また $v \in V$ に対して剰余体拡大
$\kappa(v)/\kappa(\nu(v))$ は代数的である。後者は
$V \subset S'$ かつ $S'$ が $S$ 上整であることによる。
したがって $V \to S$ のファイバーは離散的であり
（Morphisms, Lemma
\ref{morphisms-lemma-algebraic-residue-field-extension-closed-point-fibre}）、
$(f')^{-1}V \to S$ は局所準有限である
（Morphisms, Lemma \ref{morphisms-lemma-locally-quasi-finite-fibres}）。
これは $(f')^{-1}V \subset U$ および $V \subset U'$ を意味する。
$x$ は任意であったから、(a)、(b)、(c) が成り立つ。

\medskip\noindent
$s = f(x)$ とする。$(T, t) \to (S, s)$ を初等エタール近傍とする。
$T$ への基底変換を下付き添字 ${}_T$ で表す。
$y = (x, t) \in X_T$ を、$x$ の上にあるファイバー $X_t$ の
一意な点とする。$U_T \subset X_T$ は $f_T$ が準有限となる点の集合である。
Morphisms, Lemma \ref{morphisms-lemma-base-change-quasi-finite} を参照。
また、
$$
X_T \xrightarrow{f'_T} S'_T \xrightarrow{\nu_T} T
$$
は $X_T$ における $T$ の正規化である。
Lemma \ref{more-morphisms-lemma-normalization-smooth-localization} を参照。
$y \in U_T \subset X_T \to S'_T \to T$ に対して上の主張が成り立つとする。
すなわち、$(f'_T)^{-1}V' \to V'$ が同型となる開近傍
$f'_T(y) \in V' \subset S'_T$ が得られると仮定する。
射 $S'_T \to S'$ はエタールなので、$V'$ の像 $V \subset S'$ は開である。
$f'_T(y) \in V'$ なので $f'(x) \in V$ である。また、
$$
\xymatrix{
(f'_T)^{-1}V' \ar[r] \ar[d] & (f')^{-1}(V) \ar[d] \\
V' \ar[r] & V
}
$$
はファイバー平方である。これは $S'_T \times_{S'} X = X_T$ による。
左の垂直射は同型であり、$\{V' \to V\}$ はエタール被覆なので、
Descent, Lemma \ref{descent-lemma-descending-property-isomorphism} により
右の垂直射も同型である。言い換えれば、
$x \in U \subset X \to S' \to S$ に対して主張が成り立つ。

\medskip\noindent
前段落の結果により、主張を証明する際には $S$ を
$s = f(x)$ の初等エタール近傍で置き換えてよい。したがって、
分解
$$
X = V \amalg W
$$
が存在すると仮定できる。ここで両項は開かつ閉な部分スキームであり、
$V \to S$ は有限、$x \in V$ である。
Lemma \ref{more-morphisms-lemma-etale-splits-off-quasi-finite-part-technical} を参照。
$X$ は $S$ 上 $V$ と $W$ の非交和であり、$V \to S$ は有限なので、
$X$ における $S$ の正規化は射
$$
X = V \amalg W \longrightarrow V \amalg W' \longrightarrow S
$$
である。ここで $W'$ は $W$ における $S$ の正規化である。Morphisms,
Lemmas \ref{morphisms-lemma-normalization-in-disjoint-union},
\ref{morphisms-lemma-finite-integral}、および
\ref{morphisms-lemma-normalization-in-integral}.
を参照。主張が従い、証明は完了する。
\end{proof}

\begin{lemma}
\label{more-morphisms-lemma-quasi-finite-separated-quasi-affine}
\begin{slogan}
準有限かつ分離的な射は準アフィンである
\end{slogan}
$f : X \to S$ をスキームの射とする。
$f$ は準有限かつ分離的であると仮定する。
$S'$ を $X$ における $S$ の正規化とする。Morphisms,
Definition \ref{morphisms-definition-normalization-X-in-Y} を参照。
図式は
$$
\xymatrix{
X \ar[rd]_f \ar[rr]_{f'} & & S' \ar[ld]^\nu \\
& S &
}
$$
である。このとき $f'$ は準コンパクトな開埋め込みであり、
$\nu$ は整である。特に $f$ は準アフィンである。
\end{lemma}

\begin{proof}
Lemma \ref{more-morphisms-lemma-finite-type-separated} から従う。実際、この補題により、
$(f')^{-1}(U') = X$ かつ $X \to U'$ が同型となる開部分スキーム
$U' \subset S'$ が存在する。言い換えれば、$f'$ は開埋め込みである。
$f$ は準コンパクトであり、$\nu : S' \to S$ は分離的なので、
$f'$ は準コンパクトである
（Schemes, Lemma \ref{schemes-lemma-quasi-compact-permanence}）。
したがって Morphisms,
Lemma \ref{morphisms-lemma-characterize-quasi-affine} により、
$f$ は準アフィンである。
\end{proof}

\begin{lemma}[Zariski の主定理]
\label{more-morphisms-lemma-quasi-finite-separated-pass-through-finite}
\begin{reference}
\cite[IV Corollary 18.12.13]{EGA}
\end{reference}
$f : X \to S$ をスキームの射とする。
$f$ は準有限かつ分離的であり、$S$ は準コンパクトかつ
準分離的であると仮定する。このとき分解
$$
\xymatrix{
X \ar[rd]_f \ar[rr]_j & & T \ar[ld]^\pi \\
& S &
}
$$
が存在し、$j$ は準コンパクトな開埋め込み、$\pi$ は有限である。
\end{lemma}

\begin{proof}
$X \to S' \to S$ を
Lemma \ref{more-morphisms-lemma-quasi-finite-separated-quasi-affine} の結論のものとする。
Properties, Lemma
\ref{properties-lemma-integral-algebra-directed-colimit-finite} により、
$\nu_*\mathcal{O}_{S'} = \colim_{i \in I} \mathcal{A}_i$ と書ける。
これは有限準連接 $\mathcal{O}_S$-代数
$\mathcal{A}_i \subset \nu_*\mathcal{O}_{S'}$ の有向余極限である。
このとき、各 $i$ に対して
$\pi_i : T_i = \underline{\Spec}_S(\mathcal{A}_i) \to S$
は有限射である。遷移射 $T_{i'} \to T_i$ はアフィンであり、
$S' = \lim T_i$ である。

\medskip\noindent
Limits, Lemma \ref{limits-lemma-descend-opens} により、ある $i$ と、
$S'$ における逆像が $f'(X)$ に等しい準コンパクトな開部分
$U_i \subset T_i$ が存在する。$i' \geq i$ に対し、
$U_{i'}$ を $T_{i'}$ における $U_i$ の逆像とする。このとき
$X \cong f'(X) = \lim_{i' \geq i} U_{i'}$ である。Limits,
Lemma \ref{limits-lemma-directed-inverse-system-has-limit} を参照。
Limits, Lemma \ref{limits-lemma-finite-type-eventually-closed} により、
ある $i' \geq i$ に対して $X \to U_{i'}$ は閉埋め込みである。
実際、十分大きい $i'$ に対して $X \cong U_{i'}$ であるが、
これは必要ない。したがって $X \to T_{i'}$ は埋め込みである。
Morphisms, Lemma \ref{morphisms-lemma-factor-quasi-compact-immersion}
により、これを $X \to T \to T_{i'}$ と分解できる。第一の射は
開埋め込み、第二の射は閉埋め込みである。これで示された。
\end{proof}

\begin{lemma}
\label{more-morphisms-lemma-quasi-finite-separated-pass-through-finite-addendum}
% P05-EN-CH37-0197
Lemma \ref{more-morphisms-lemma-quasi-finite-separated-pass-through-finite}
と同じ記法および仮定を用いる。さらに $f$ が局所有限表示であると
仮定する。このとき、$T$ が $S$ 上有限かつ有限表示となるように
分解を選ぶことができる。
\end{lemma}

\begin{proof}
Limits, Lemma
\ref{limits-lemma-finite-in-finite-and-finite-presentation} により、
$T = \lim T_i$ と書ける。すべての $T_i$ は $S$ 上有限かつ有限表示であり、
遷移射 $T_{i'} \to T_i$ は閉埋め込みである。Limits,
Lemma \ref{limits-lemma-descend-opens} により、ある $i$ と、
$T$ における逆像が $X$ である開部分スキーム $U_i \subset T_i$
が存在する。Limits, Lemma
\ref{limits-lemma-finite-type-eventually-closed} により、
十分大きい $i$ に対して $X \cong U_i$ である。
$T$ を $T_i$ で置き換えれば証明が完了する。
\end{proof}








\section{Zariski の主定理の応用 I}
\label{more-morphisms-section-applications-zmt}

\noindent
最初の応用は、有限射を有限ファイバーをもつ固有射として特徴づけることである。

\begin{lemma}
\label{more-morphisms-lemma-characterize-finite}
$f : X \to S$ をスキームの射とする。次は同値である。
\begin{enumerate}
\item $f$ は有限である。
\item $f$ は有限ファイバーをもつ固有射である。
\item $f$ は固有かつ局所準有限である。
\item $f$ は普遍閉、分離的、局所有限型であり、有限ファイバーをもつ。
\end{enumerate}
\end{lemma}

\begin{proof}
(1) から (2) は Morphisms, Lemmas \ref{morphisms-lemma-finite-proper},
\ref{morphisms-lemma-quasi-finite},
および \ref{morphisms-lemma-finite-quasi-finite} による。
(2) から (3) は Morphisms, Lemma \ref{morphisms-lemma-finite-fibre}
による。(3) から (4) は固有射の定義と Morphisms, Lemmas
\ref{morphisms-lemma-quasi-finite-locally-quasi-compact} および
\ref{morphisms-lemma-quasi-finite} による。

\medskip\noindent
(4) を仮定する。$s \in S$ をとる。Morphisms,
Lemma \ref{morphisms-lemma-finite-fibre} により、$X_s$ の有限個の点は
すべて $X_s$ における孤立点である。
Lemma \ref{more-morphisms-lemma-etale-splits-off-quasi-finite-part} のように、
初等エタール近傍 $(U, u) \to (S, s)$ と分解
$X_U = V \amalg W$ を選ぶ。$X_s$ のすべての点が孤立しているため、
$W_u = \emptyset$ である。$f$ は普遍閉なので、$U$ における
$W$ の像は $u$ を含まない閉集合である。$U$ を縮小すれば
$W = \emptyset$ と仮定できる。言い換えれば、$X_U = V$ は
$U$ 上有限である。$s \in S$ は任意であったから、像が $S$ を被覆する
エタール射の族 $\{U_i \to S\}$ で、基底変換
$X_{U_i} \to U_i$ が有限となるものが存在する。
$\{U_i \to S\}$ はエタール被覆である。Topologies,
Definition \ref{topologies-definition-etale-covering} を参照。
したがってこれは fpqc 被覆である。Topologies,
Lemma \ref{topologies-lemma-zariski-etale-smooth-syntomic-fppf-fpqc}
を参照。ゆえに Descent,
Lemma \ref{descent-lemma-descending-property-finite} により、
$f$ は有限である。
\end{proof}

\noindent
帰結として、次の有用な結果を得る。

\begin{lemma}
\label{more-morphisms-lemma-proper-finite-fibre-finite-in-neighbourhood}
$f : X \to S$ をスキームの射とし、$s \in S$ とする。
$f$ は固有であり、$f^{-1}(\{s\})$ は有限集合であると仮定する。
このとき、$f|_{f^{-1}(V)} : f^{-1}(V) \to V$ が有限となる
$s$ の開近傍 $V \subset S$ が存在する。
\end{lemma}

\begin{proof}
Morphisms, Lemma \ref{morphisms-lemma-finite-fibre} により、射 $f$ は
$f^{-1}(\{s\})$ のすべての点で準有限である。Morphisms,
Lemma \ref{morphisms-lemma-quasi-finite-points-open} により、
$f$ が準有限となる点の集合は開部分 $U \subset X$ である。
$Z = X \setminus U$ とおく。このとき $s \not \in f(Z)$ である。
$f$ は固有なので、集合 $f(Z) \subset S$ は閉である。
$f(Z) \cap V = \emptyset$ を満たす任意の $s$ の開近傍
$V \subset S$ を選ぶ。このとき $f^{-1}(V) \to V$ は局所準有限かつ
固有である。したがって準有限であり
（Morphisms, Lemma
\ref{morphisms-lemma-quasi-finite-locally-quasi-compact}）、
有限ファイバーをもち
（Morphisms, Lemma \ref{morphisms-lemma-quasi-finite}）、
Lemma \ref{more-morphisms-lemma-characterize-finite} により有限である。
\end{proof}

\begin{lemma}
\label{more-morphisms-lemma-flat-proper-family-cannot-collapse-fibre}
スキームの可換図式
$$
\xymatrix{
X \ar[rr]_h \ar[rd]_f & & Y \ar[ld]^g \\
& S
}
$$
を考える。$s \in S$ とし、次を仮定する。
\begin{enumerate}
\item $X \to S$ は固有射である。
\item $Y \to S$ は分離的かつ局所有限型である。
\item $X_s \to Y_s$ の像は有限である。
\end{enumerate}
このとき、$s$ を含む開部分スキーム
% P05-EN-CH37-0198
$U \subset S$ が存在し、$X_U \to Y_U$ は $U$ 上有限な閉部分スキーム
$Z \subset Y_U$ を経由する。
\end{lemma}

\begin{proof}
$Z \subset Y$ を $h$ のスキーム論的像とする。Morphisms,
Section \ref{morphisms-section-scheme-theoretic-image} を参照。
Morphisms, Lemma \ref{morphisms-lemma-scheme-theoretic-image-is-proper}
により、射 $X \to Z$ は全射で、$Z \to S$ は固有である。
したがって $X_s \to Z_s$ は全射である。
% P05-EN-CH37-0199
(3) により $Z_s$ は有限である。ゆえに
Lemma \ref{more-morphisms-lemma-proper-finite-fibre-finite-in-neighbourhood} により、
$s$ の開近傍上で $Z \to S$ は有限である。
\end{proof}





\section{Zariski の主定理の応用 II}
\label{more-morphisms-section-applications-zmt-II}

\noindent
本節では、準有限射の構造に関する Zariski の主定理の帰結を
さらにいくつか与える。

\begin{lemma}
\label{more-morphisms-lemma-separated-locally-quasi-finite-over-affine}
$f : X \to Y$ を分離的かつ局所準有限な射とし、$Y$ はアフィンであるとする。
このとき $X$ の任意の有限点集合は、$X$ のあるアフィン開部分に含まれる。
\end{lemma}

\begin{proof}
$x_1, \ldots, x_n \in X$ とする。$x_i \in U$ を満たす準コンパクトな
開部分 $U \subset X$ を選ぶ。このとき
Lemma \ref{more-morphisms-lemma-quasi-finite-separated-quasi-affine} により
$U \to Y$ は準アフィンである。したがって Properties,
Lemma \ref{properties-lemma-ample-finite-set-in-affine} により、
$x_1, \ldots, x_n$ を含むアフィン開部分 $V \subset U$ が存在する。
\end{proof}

\begin{lemma}
\label{more-morphisms-lemma-quasi-finite-finite-over-dense-open}
$f : Y \to X$ を準有限射とする。このとき、
$f|_{f^{-1}(U)} : f^{-1}(U) \to U$ が有限となる
稠密開部分 $U \subset X$ が存在する。
\end{lemma}

\begin{proof}
$U_i \subset X$、$i \in I$ を、制限
$f|_{f^{-1}(U_i)} : f^{-1}(U_i) \to U_i$ が有限となる開部分の族とする。
$U = \bigcup U_i$ とおけば、制限
$f|_{f^{-1}(U)} : f^{-1}(U) \to U$ は有限である。Morphisms,
Lemma \ref{morphisms-lemma-finite-local} を参照。
したがって問題は $X$ 上局所的であり、$X$ はアフィンと仮定してよい。

\medskip\noindent
$X$ はアフィンと仮定する。各 $V_j$ がアフィンとなるように
$Y = \bigcup_{j = 1, \ldots, m} V_j$ と書く。これは $f$ が準有限、
特に準コンパクトであるため可能である。各 $V_j \to X$ は準有限かつ
分離的である。$\eta \in X$ を $X$ のある既約成分の一般点とする。
Morphisms, Lemmas \ref{morphisms-lemma-quasi-finite} および
\ref{morphisms-lemma-generically-finite} により、
$f^{-1}(U_\eta) \cap V_j \to U_\eta$ が有限となる開近傍
$\eta \in U_\eta$ が存在する。各 $j = 1, \ldots, m$ に対して
機能するように $U_\eta$ を選べる。$X$ の一般点の集合は
$X$ において稠密である。したがって、各
$f^{-1}(W) \cap V_j \to W$ が有限となる稠密開部分
$W = \bigcup_\eta U_\eta$ が存在する。
$f|_{f^{-1}(U)} : f^{-1}(U) \to U$ が有限となる稠密開部分
$U \subset W$ の存在を示せば十分である。そこで $X$ を
$W$ のアフィン開部分スキームで置き換え、各 $V_j \to X$ が
有限であると仮定してよい。

\medskip\noindent
$X$ はアフィンで、$Y = \bigcup_{j = 1, \ldots, m} V_j$、
各 $V_j$ はアフィンであり、制限 $f|_{V_j} : V_j \to X$ は有限であると
仮定する。次のようにおく。
$$
\Delta_{ij} =
\Big(\overline{V_i \cap V_j} \setminus V_i \cap V_j\Big) \cap V_j.
$$
これは $V_j$ における開部分 $V_i \cap V_j$ の境界なので、
$V_j$ の疎な閉部分集合である。Morphisms,
Lemma \ref{morphisms-lemma-image-nowhere-dense-finite} により、
像 $f(\Delta_{ij})$ は $X$ の疎な閉部分集合である。Topology,
Lemma \ref{topology-lemma-nowhere-dense} により、合併
$T = \bigcup f(\Delta_{ij})$ は $X$ の疎な閉部分集合である。
したがって $U = X \setminus T$ は $X$ の稠密開部分である。
$f|_{f^{-1}(U)} : f^{-1}(U) \to U$ は有限であると主張する。
これを見るため、$U' \subset U$ をアフィン開部分とする。
$Y' = f^{-1}(U') = U' \times_X Y$、
$V_j' = Y' \cap V_j = U' \times_X V_j$ とおく。$f$ の制限
$$
f' = f|_{Y'} : Y' \longrightarrow U'
$$
を考える。この射について、
$Y' = \bigcup_{j = 1, \ldots, m} V'_j$ はアフィン開被覆であり、
各 $V'_j \to U'$ は有限で、$V_i' \cap V_j'$ は
$V'_i$ と $V'_j$ の両方において開かつ閉である。
したがって $V_i' \cap V_j'$ はアフィンであり、写像
$$
\mathcal{O}(V'_i) \otimes_{\mathbf{Z}} \mathcal{O}(V'_j)
\longrightarrow
\mathcal{O}(V'_i \cap V'_j)
$$
は全射である。これは $Y'$ が分離的であることを意味する。
Schemes, Lemma \ref{schemes-lemma-characterize-separated} を参照。
最後に、可換図式
$$
\xymatrix{
\coprod_{j = 1, \ldots, m} V'_j \ar[rd] \ar[rr] & & Y' \ar[ld] \\
& U' &
}
$$
を考える。南東向きの射は有限、したがって固有であり、水平射は全射、
南西向きの射は分離的である。ゆえに Morphisms,
Lemma \ref{morphisms-lemma-image-proper-is-proper} により
$Y' \to U'$ は固有である。これは準有限でもあるから、
Lemma \ref{more-morphisms-lemma-characterize-finite} により有限である。
これで示された。
\end{proof}

\begin{lemma}
\label{more-morphisms-lemma-stratify-flat-fp-lqf-universally-bounded}
$f : X \to S$ を、平坦かつ局所有限表示かつ分離的で、普遍的に有界な
ファイバーをもつ局所準有限射とする。このとき閉集合
$$
\emptyset = Z_{-1} \subset Z_0 \subset Z_1 \subset Z_2 \subset
\ldots \subset Z_n = S
$$
であって、$S_r = Z_r \setminus Z_{r - 1}$ とおくとき、層別化
$S = \coprod_{r = 0, \ldots, n} S_r$ が次の普遍性によって特徴づけられるものが
存在する。$g : T \to S$ が与えられたとき、射影
$X \times_S T \to T$ が次数 $r$ の有限局所自由射であることと、
$g(T) \subset S_r$ が集合論的に成り立つこととは同値である。
\end{lemma}

\begin{proof}
$n$ を $X \to S$ のファイバーの次数を上から抑える整数とする。
Morphisms, Lemma \ref{morphisms-lemma-base-change-universally-bounded}
により、任意の基底変換についてもファイバーの次数は $n$ で上から抑えられる。
とくに、この補題の主張に現れるすべての整数 $r$ は $\leq n$ である。
$n$ に関する帰納法で補題を証明する。基底の場合 $n = 0$ は明らかである。

\medskip\noindent
$\deg_{\kappa(s)}(X_s) = n$ を満たす点 $s \in S$ の集合は開部分集合
$S_n \subset S$ であり、$X \times_S S_n \to S_n$ は次数 $n$ の
有限局所自由射であると主張する。実際、$s \in S$ をそのような点とする。
Lemma \ref{more-morphisms-lemma-etale-splits-off-quasi-finite-part} のように、
初等的なエタール射 $(U, u) \to (S, s)$ と分解
$U \times_S X = W \amalg V$ を選ぶ。$V \to U$ は有限かつ平坦かつ
局所有限表示なので、$V \to U$ は有限局所自由である。Morphisms,
Lemma \ref{morphisms-lemma-finite-flat} を参照。
$U$ を $u$ のより小さい近傍に縮小すれば、$V \to U$ はある次数 $d$ の
有限局所自由射であるとしてよい。Morphisms,
Lemma \ref{morphisms-lemma-finite-locally-free} を参照。
$u \mapsto s$ かつ $W_u = \emptyset$ であるから $d = n$ である。
$n$ はファイバーの次数の最大値なので、$W = \emptyset$ である！
したがって $U \times_S X \to U$ は次数 $n$ の有限局所自由射である。
Descent, Lemma \ref{descent-lemma-descending-property-finite-locally-free}
により、$X \to S$ は $s$ の開近傍 $\Im(U \to S)$ 上で次数 $n$ の
有限局所自由射である（Morphisms,
Lemma \ref{morphisms-lemma-etale-open}）。これで主張が示された。

\medskip\noindent
$S' = S \setminus S_n$ に誘導される被約スキーム構造を入れ、
$X' = X \times_S S'$ とおく。$X' \to S'$ のファイバーの次数は
$n - 1$ によって普遍的に上から抑えられることに注意する。
帰納法により、射 $X' \to S'$ に適合する層別化
$S' = S_0 \amalg \ldots \amalg S_{n - 1}$ を得る。
$S = \coprod_{r = 0, \ldots, n} S_r$ が射 $X \to S$ に対して所望のものであると
主張する。スキームの射 $g : T \to S$ を取り、$X \times_S T \to T$ が
次数 $r$ の有限局所自由射であると仮定する。上で述べたように、これは
$r \leq n$ を意味する。$r = n$ なら、$T \to S$ が $S_n$ を経由することは
明らかである。$r < n$ なら、
% P05-EN-CH37-0200
$g(T) \subset S' = S \setminus S_d$ が集合論的に成り立つので、
$T_{red} \to S$ は $S'$ を経由する。Schemes,
Lemma \ref{schemes-lemma-map-into-reduction} を参照。
また、基底変換である $X \times_S T_{red} \to T_{red}$ も次数 $r$ の
有限局所自由射である。層別化
$S' = \coprod_{r = 0, \ldots, n - 1} S_r$ の普遍性により、
$g(T) = g(T_{red})$ は $S_r$ に含まれる。
逆に、$g(T) \subset S_r$ が集合論的に成り立つような
$g : T \to S$ が与えられたとする。$r = n$ なら $g$ は $S_n$ を経由し、
$X \times_S T \to T$ は基底変換として次数 $n$ の有限局所自由射である。
$r < n$ なら、$X \times_S T \to T$ は分離的かつ平坦かつ局所有限表示であり、
$T_{red}$ への制限は次数 $r$ の有限局所自由射である。
$T_{red} \to T$ は普遍同相なので、
$X \times_S T_{red} \to X \times_S T$ も普遍同相である。
したがって $X \times_S T \to T$ は普遍閉である（有限射
$X \times_S T_{red} \to T_{red}$ についてそうであるからである）。
ゆえに、たとえば Lemma \ref{more-morphisms-lemma-characterize-finite} により、
$X \times_S T \to T$ は有限である。そこで Morphisms,
Lemma \ref{morphisms-lemma-finite-flat} を用いると、
$X \times_S T \to T$ は有限局所自由である。最後に、すべてのファイバーの
次数が $r$ なので、その次数は $r$ である。
\end{proof}

\begin{lemma}
\label{more-morphisms-lemma-stratify-flat-fp-qf}
$f : X \to S$ を、平坦かつ局所有限表示かつ分離的な準有限スキーム射とする。
このとき閉集合
$$
\emptyset = Z_{-1} \subset Z_0 \subset Z_1 \subset Z_2 \subset
\ldots \subset S
$$
であって、$S_r = Z_r \setminus Z_{r - 1}$ とおくとき、層別化
$S = \coprod S_r$ が次の普遍性によって特徴づけられるものが存在する。
射 $g : T \to S$ が与えられたとき、射影 $X \times_S T \to T$ が
次数 $r$ の有限局所自由射であることと、$g(T) \subset S_r$ が
集合論的に成り立つこととは同値である。さらに、包含射 $S_r \to S$ は
準コンパクトである。
\end{lemma}

\begin{proof}
問題は $S$ 上局所的なので、$S$ はアフィンであるとしてよい。
Morphisms, Lemma
\ref{morphisms-lemma-locally-quasi-finite-qc-source-universally-bounded}
により、この場合 $f$ のファイバーは普遍的に有界である。
したがって、層別化の存在は
Lemma \ref{more-morphisms-lemma-stratify-flat-fp-lqf-universally-bounded} から従う。

\medskip\noindent
各 $r \geq 0$ について $U_r = S \setminus Z_r \to S$ が準コンパクトであることを
示す。これにより、初等的な位相の議論から最後の主張が従う。
準コンパクト写像の合成は準コンパクトなので、
$U_r \to U_{r - 1}$ が準コンパクトであることを示せば十分である。
アフィン開集合 $W \subset U_{r - 1}$ を選び、$W = \Spec(A)$ と書く。
このとき、あるイデアル $I \subset A$ に対して $Z_r \cap W = V(I)$ であり、
$X \times_S \Spec(A/I) \to \Spec(A/I)$ は次数 $r$ の有限局所自由射である。
$I_i \subset I$ が有限生成イデアル全体を走るとき
$A/I = \colim A/I_i$ であることに注意する。Limits,
Lemma \ref{limits-lemma-descend-finite-locally-free} により、ある $i$ について
$X \times_S \Spec(A/I_i) \to \Spec(A/I_i)$ は次数 $r$ の有限局所自由射である。
（ここでは、$X \to S$ が局所有限表示かつ分離的かつ準コンパクトなので
有限表示であることを用いた。）したがって
$\Spec(A/I_i) \to \Spec(A) = W$ は集合論的に $Z_r \cap W$ を経由する。
ゆえに $Z_r \cap W = V(I_i)$ は $A$ の有限個の元からなる部分集合の零点集合である。
これは $W \setminus Z_r$ が有限個の標準開集合の合併、したがって
準コンパクトであることを意味する。これで示された。
\end{proof}

\begin{lemma}
\label{more-morphisms-lemma-stratify-flat-fp-lqf}
$f : X \to S$ を、平坦かつ局所有限表示かつ分離的な局所準有限スキーム射とする。
このとき開部分スキーム
$$
S = U_0 \supset U_1 \supset U_2 \supset \ldots
$$
であって、体 $k$ に対する射 $\Spec(k) \to S$ が $U_d$ を経由することと、
$X \times_S \Spec(k)$ が $k$ 上次数 $\geq d$ であることとが同値になるものが
存在する。
\end{lemma}

\begin{proof}
この主張は、ファイバーの次数が $\geq d$ となる点の集合が開であるというだけである。
したがって $S$ 上局所的に議論でき、$S$ はアフィンであるとしてよい。
この場合、任意の準コンパクト開集合 $W \subset X$ に対し、
$W \to S$ のファイバーの次数が $\geq d$ となる点の集合 $U_d(W)$ は
Lemma \ref{more-morphisms-lemma-stratify-flat-fp-qf} により開である。
$U_d = \bigcup_W U_d(W)$ なので結論が従う。
\end{proof}

\begin{lemma}
\label{more-morphisms-lemma-go-down-with-annihilators}
$f : X \to S$ を、平坦かつ局所有限表示な局所準有限スキーム射とする。
% P05-EN-CH37-0201
$g \in \Gamma(X, \mathcal{O}_X)$ を非零元とする。このとき、
$g|_V \not = 0$ を満たす開集合 $V \subset X$、可換図式
$$
\xymatrix{
V \ar[r] \ar[d]_\pi & X \ar[d]^f \\
U \ar[r] & S,
}
$$
に入る開集合 $U \subset S$、準連接部分層
$\mathcal{F} \subset \mathcal{O}_U$、整数 $r > 0$、および像が $g|_V$ を含む
単射な $\mathcal{O}_U$-加群写像
$\mathcal{F}^{\oplus r} \to \pi_*\mathcal{O}_V$ が存在する。
\end{lemma}

\begin{proof}
$X$ と $S$ はアフィンであるとしてよい。
Lemmas \ref{more-morphisms-lemma-stratify-flat-fp-lqf-universally-bounded} および
\ref{more-morphisms-lemma-stratify-flat-fp-qf} のようなフィルトレーション
$\emptyset = Z_{-1} \subset Z_0 \subset Z_1 \subset Z_2 \subset \ldots
\subset Z_n = S$ を得る。$T \subset X$ を、有限 $\mathcal{O}_X$-加群
$\Im(g : \mathcal{O}_X \to \mathcal{O}_X)$ のスキーム論的支持とする。
$T$ は $\mathcal{O}_X$ の切断としての $g$ の支持であり
（Modules, Definition \ref{modules-definition-support}）、任意の開集合
$V \subset X$ に対して $g|_V \not = 0$ であることと
$V \cap T \not = \emptyset$ であることとは同値である。
$f(T) \subset Z_r$ が集合論的に成り立つような最小の整数を $r$ とする。
$T$ の既約成分の一般点 $\xi \in T$ であって、
$f(\xi) \not \in Z_{r - 1}$（したがって $f(\xi) \in Z_r$）を満たすものを取る。
$S$ を、$S \setminus Z_{r - 1}$ に含まれる $f(\xi)$ のアフィン近傍で
置き換えてよい。$S = \Spec(A)$ と書き、$V(I) = Z_r$ が集合論的に成り立つ
有限生成イデアル $I = (a_1, \ldots, a_m) \subset A$ を取る
（Algebra, Lemma \ref{algebra-lemma-qc-open} を参照）。
$r$ の選び方により $g$ の支持は $f^{-1}V(I)$ に含まれるので、ある整数 $N$ が存在して
$j = 1, \ldots, m$ に対し $a_j^N g = 0$ となる。
% P05-EN-CH37-0202
$a_j$ を $a_j^r$ で置き換えることにより、$Ig = 0$ としてよい。
任意の $A$-加群 $M$ に対し、$M$ の $I$-捩子を $M[I]$ と書く。すなわち
$M[I] = \{m \in M \mid Im = 0\}$ とする。$X = \Spec(B)$ と書けば
$g \in B[I]$ である。$A \to B$ は平坦なので
$$
B[I] = A[I] \otimes_A B \cong A[I] \otimes_{A/I} B/IB
$$
である。$Z_r$ の選び方により、$A/I$-加群 $B/IB$ は階数 $r$ の
有限局所自由加群である。したがって $S$ を $f(\xi)$ のより小さいアフィン開近傍で
置き換えれば、$A/I$-加群として
$B/IB \cong (A/IA)^{\oplus r}$ であるとしてよい。
$I$ を法として前文の同型に帰着する写像 $\psi : A^{\oplus r} \to B$ を選ぶ。
すると誘導される写像
$$
A[I]^{\oplus r} \longrightarrow B[I]
$$
は同型である。$\mathcal{F}$ を $A$-加群 $A[I]$ に対応する準連接層とし、
$\mathcal{F}^{\oplus r} \to \pi_*\mathcal{O}_V$ を
$A[I]^{\oplus r} \subset A^{\oplus r} \to B$ に対応する写像とすれば、
補題が従う。
\end{proof}

\begin{lemma}
\label{more-morphisms-lemma-there-is-a-scheme-integral-over}
$U \to X$ を全射なエタールスキーム射とし、$X$ は準コンパクトかつ
準分離的であると仮定する。このとき全射な整射 $Y \to X$ であって、
任意の $y \in Y$ に対して開近傍 $V \subset Y$ が存在し、
$V \to X$ が $U$ を経由するものが存在する。実際、$Y \to X$ は
有限かつ有限表示であるとしてよい。
\end{lemma}

\begin{proof}
$X$ は準コンパクトなので、$U' = \coprod U_i \to X$ が全射となるような
有限個のアフィン開集合 $U_i \subset U$ が存在する。$U$ を $U'$ で置き換えると、
$U$ はアフィンであるとしてよい。とくに $U \to X$ は分離的である
（Schemes, Lemma \ref{schemes-lemma-affine-separated}）。
したがって、$U \to X$ の幾何学的ファイバーの次数を上から抑える整数 $d$ が存在する
（Morphisms, Lemma
\ref{morphisms-lemma-locally-quasi-finite-qc-source-universally-bounded}).
準コンパクトかつ分離的なスキーム $U$ から $X$ への全射エタール射すべてについて、
$d$ に関する帰納法で補題を証明する。$d = 1$ なら $U = X$ であり、
$Y = U$ とすれば結論が成り立つ。$d > 1$ と仮定する。

\medskip\noindent
Lemma \ref{more-morphisms-lemma-quasi-finite-separated-quasi-affine} を適用して分解
$$
\xymatrix{
U \ar[rr]_j \ar[rd] & & Y \ar[ld]^\pi \\
& X
}
$$
を得る。ここで $\pi$ は整射、$j$ は準コンパクトな開埋め込みである。
$j(U)$ は $Y$ にスキーム論的に稠密であると仮定してよく、実際そう仮定する。
次の分解に注意する。
$$
U \times_X Y = U \amalg W
$$
ここで第1の直和因子は $U \to U \times_X Y$ の像であり
（Schemes, Lemma \ref{schemes-lemma-semi-diagonal} により閉であり、
$Y$ 上エタールなスキーム間の射としてエタールなので開でもある）、
第2の直和因子はその開かつ閉な補集合である。$W$ の像 $V \subset Y$ は、
$Y \setminus U$ を含む開部分スキームである。

\medskip\noindent
エタール射 $W \to Y$ の幾何学的ファイバーの濃度は $< d$ である。
実際、$U \subset Y$ の幾何学的点については直接見れば明らかである。
$U \subset Y$ は稠密なので、たとえば Lemma
\ref{more-morphisms-lemma-stratify-flat-fp-lqf-universally-bounded}
により、これは $Y$ のすべての幾何学的点について成り立つ
（準コンパクトで分離的なエタール射のファイバーの次数は特殊化のもとで増加しない）。
したがって $W \to V$ に帰納法の仮定を適用し、$Z$ がスキームで Zariski 局所的に
$W$ を経由するような全射な整射 $Z \to V$ を得る。
$Z' \to Y$ が整射、$Z \to Z'$ が開埋め込みとなる分解
$Z \to Z' \to Y$ を選ぶ
（Lemma \ref{more-morphisms-lemma-quasi-finite-separated-quasi-affine}）。
$Z'$ を、$Z'$ の中での $Z$ のスキーム論的閉包で置き換えれば、$Z$ は $Z'$ に
スキーム論的に稠密であるとしてよい。この操作の後には
$Z' \times_Y V = Z$ となる。最後に、$Y \setminus V$ 上に誘導される
被約閉部分スキーム構造を $T \subset Y$ とする。射
$$
Z' \amalg T \longrightarrow X
$$
を考える。これは構成により全射な整射である。$T \subset U$ なので、
射 $T \to X$ が $U$ を経由することは明らかである。一方、点 $z \in Z'$ を取る。
$z \not \in Z$ なら、$z$ は $Y \setminus V \subset U$ の点に写り、
射が $U$ を経由するような $z$ の近傍を得る。$z \in Z$ なら、
$W \subset U \times_X Y$ を、したがって $U$ を経由する近傍
$\Omega \subset Z$ が存在する。これで存在が示された。

\medskip\noindent
$U$ を Zariski 局所的に経由する整かつ全射な $Y \to X$ が得られたと仮定する。
$V_j \to X$ が $U$ を経由するような有限アフィン開被覆
$Y = \bigcup V_j$ を選ぶ。$Y_i \to X$ が有限かつ有限表示となる形で
$Y = \lim Y_i$ と書ける。Limits, Lemma
\ref{limits-lemma-integral-limit-finite-and-finite-presentation}.
十分大きい $i$ に対して、$Y$ への逆像が $V_j$ となるアフィン開集合
$V_{i, j} \subset Y_i$ を取れる。Limits, Lemma
\ref{limits-lemma-descend-opens} を参照。さらに $i$ を大きくすれば、
$X$ 上の射 $V_j \to U$ は $X$ 上の射 $V_{i, j} \to U$ から得られる。
Limits, Proposition
\ref{limits-proposition-characterize-locally-finite-presentation} を参照。
これで証明が完了する。
\end{proof}







\section{連結ファイバーをもつ射への応用}
\label{more-morphisms-section-application-geometrically-connected}

\noindent
本節では、すべてのファイバーが幾何学的に連結または幾何学的に整となる射を
構成するための補題をいくつか証明する。これらは後に有限型射の局所構造を
調べる際に有用となる。

\begin{lemma}
\label{more-morphisms-lemma-descent-connected-fibres}
スキームの射の図式
$$
\xymatrix{
Z \ar[r]_{\sigma} \ar[rd] & X \ar[d] \\
& Y
}
$$
% P05-EN-CH37-0203
および点 $y \in Y$ を考える。次を仮定する。
\begin{enumerate}
\item $X \to Y$ は有限表示かつ平坦である。
\item $Z \to Y$ は有限局所自由である。
\item $Z_y \not = \emptyset$,
\item $X \to Y$ のすべてのファイバーは幾何学的に被約である。
\item $X_y$ は $\kappa(y)$ 上幾何学的に連結である。
\end{enumerate}
このとき、$X^0_y = X_y$ を満たし、$X^0 \to Y$ のすべての非空ファイバーが
幾何学的に連結となる準コンパクト開集合 $X^0 \subset X$ が存在する。
\end{lemma}

\begin{proof}
この証明では、平坦、有限表示、有限局所自由という性質が基底変換および合成で
保たれることを用いる。また、有限局所自由射が開かつ閉であることも用いる。
これらの事実は Morphisms, Lemmas
\ref{morphisms-lemma-base-change-flat},
\ref{morphisms-lemma-base-change-finite-presentation},
\ref{morphisms-lemma-base-change-finite-locally-free},
\ref{morphisms-lemma-composition-flat},
\ref{morphisms-lemma-composition-finite-presentation},
\ref{morphisms-lemma-composition-finite-locally-free},
\ref{morphisms-lemma-fppf-open}, and
\ref{morphisms-lemma-finite-proper} にある。

\medskip\noindent
% P05-EN-CH37-0204
$X_Z \to Z$ は有限表示の平坦射であり、$\sigma$ から来る切断 $s$ をもつことに
注意する。$X_Z^0$ を Situation \ref{more-morphisms-situation-connected-along-section} で定義された
$X_Z$ の部分集合とする。Lemma \ref{more-morphisms-lemma-connected-along-section-open} により、
これは $X_Z$ の開部分集合である。

\medskip\noindent
$X$ の $Z \times_Y Z$ への引き戻し $X_{Z \times_Y Z}$ には2つの切断
$s_0, s_1$ が備わる。すなわち、これらは
$\text{pr}_0, \text{pr}_1 : Z \times_Y Z \to Z$ による $s$ の基底変換である。
Situation \ref{more-morphisms-situation-connected-along-section} の構成は2つの部分集合
$(X_{Z \times_Y Z})_{s_0}^0$ と $(X_{Z \times_Y Z})_{s_1}^0$ を与える。
Lemma \ref{more-morphisms-lemma-base-change-connected-along-section} により、これらは射
$1_X \times \text{pr}_0, 1_X \times \text{pr}_1 : X_{Z \times_Y Z} \to X_Z$
による $X_Z^0$ の逆像である。とくに、これらの部分集合は開である。

\medskip\noindent
$(Z \times_Y Z)_y = \{z_1, \ldots, z_n\}$ とする。
$X_y$ は幾何学的に連結なので、各 $z_i$ 上の
$(X_{Z \times_Y Z})_{s_0}^0$ と $(X_{Z \times_Y Z})_{s_1}^0$ の
ファイバーは一致する（いずれもファイバー全体に等しい）。言い換えると
$$
s_0(z_i) \in (X_{Z \times_Y Z})_{s_1}^0
\quad\text{and}\quad
s_1(z_i) \in (X_{Z \times_Y Z})_{s_0}^0.
$$
である。集合 $(X_{Z \times_Y Z})_{s_0}^0$ と $(X_{Z \times_Y Z})_{s_1}^0$ は
$X_{Z \times_Y Z}$ で開なので、$(Z \times_Y Z)_y$ の開近傍
$W \subset Z \times_Y Z$ であって
$$
s_0(W) \subset (X_{Z \times_Y Z})_{s_1}^0
\quad\text{and}\quad
s_1(W) \subset (X_{Z \times_Y Z})_{s_0}^0.
$$
を満たすものが存在する。このとき
Situation \ref{more-morphisms-situation-connected-along-section} の構成から直接
$$
p^{-1}(W) \cap (X_{Z \times_Y Z})_{s_0}^0
=
p^{-1}(W) \cap (X_{Z \times_Y Z})_{s_1}^0
$$
% P05-EN-CH37-0205
を得る。ここで $p : X_{Z \times_Y Z} \to Z \times_W Z$ は射影である。
$Z \times_Y Z \to Y$ は有限局所自由、したがって開かつ閉なので、
$y$ のアフィン開近傍 $V \subset Y$ であって $q^{-1}(V) \subset W$ を満たすものが
存在する。ここで $q : Z \times_Y Z \to Y$ は構造射である。
補題を証明するために $Y$ を $V$ で置き換えてよい。この置換の後、
% P05-EN-CH37-0206
$X_Z^0 \subset Y_Z$ は次を満たす開部分である。
$$
(1_X \times \text{pr}_0)^{-1}(X_Z^0) =
(1_X \times \text{pr}_1)^{-1}(X_Z^0).
$$
これは、$X_Z^0$ の像 $X^0 \subset X$ が開であり、
$(X_Z \to X)^{-1}(X^0) = X_Z^0$ を満たすことを意味する。Descent,
Lemma \ref{descent-lemma-open-fpqc-covering} を参照。
最後に、Lemma \ref{more-morphisms-lemma-connected-along-section-locally-constructible} により
$X_Z^0$ は準コンパクトなので、$X^0$ も準コンパクトである
（この時点で $Y$ はアフィンであり、したがって $X$ は準コンパクトかつ
準分離的である。そのため局所構成可能であることは構成可能であることと同じで、
とくに準コンパクトである。詳細は省略する）。以上により、$X^0$ は所望の性質を
すべてもつ。
\end{proof}

\begin{lemma}
\label{more-morphisms-lemma-fibre-geometrically-connected-reduced}
$h : Y \to S$ をスキームの射、$s \in S$ を点、
$T \subset Y_s$ を開部分スキームとする。次を仮定する。
\begin{enumerate}
\item $h$ は平坦かつ有限表示である。
\item $h$ のすべてのファイバーは幾何学的に被約である。
\item $T$ は $\kappa(s)$ 上幾何学的に連結である。
\end{enumerate}
このとき、アフィンな初等エタール近傍 $(S', s') \to (S, s)$ と
準コンパクト開集合 $V \subset Y_{S'}$ であって、次を満たすものが存在する。
\begin{enumerate}
\item[(a)] $V \to S'$ のすべてのファイバーは幾何学的に連結である。
\item[(b)] $V_{s'} = T \times_s s'$.
\end{enumerate}
\end{lemma}

\begin{proof}
問題は明らかに $S$ 上局所的なので、$S$ を $s$ のアフィン開近傍で置き換えてよい。
$Y_s$ の位相は $Y$ の位相から誘導される。Schemes,
Lemma \ref{schemes-lemma-fibre-topological} を参照。したがって、
$V_s = T$ を満たす準コンパクト開集合 $V \subset Y$ を取れる。
$h$ の $V$ への制限は、準コンパクト（$S$ はアフィンで $V$ は準コンパクトだから）、
準分離的、局所有限表示かつ平坦であり、したがって有限表示かつ平坦である。
ゆえに $Y$ を $V$ で置き換えることで、(1) と (2) に加えて
$Y_s = T$ かつ $S$ はアフィンであると仮定してよい。

\medskip\noindent
$h$ が $y$ で Cohen--Macaulay となる閉点 $y \in Y_s$ を選ぶ。
Lemma \ref{more-morphisms-lemma-flat-finite-presentation-CM-open} を参照。
Lemma \ref{more-morphisms-lemma-slice-CM} により図式
$$
\xymatrix{
Z \ar[r] \ar[rd] & Y \ar[d] \\
& S
}
$$
であって、$Z \to S$ が平坦かつ局所有限表示かつ局所準有限であり、
$Z_s = \{y\}$ を満たすものが存在する。
Lemma \ref{more-morphisms-lemma-etale-makes-quasi-finite-finite-at-point} を適用すると、
初等近傍 $(S', s') \to (S, s)$ と開集合
$Z' \subset Z_{S'} = S' \times_S Z$ であって、$Z' \to S'$ が有限で、
$s$ 上にある点 $z' \in Z'$ が一意であるものを得る。
$Z' \to S'$ は（$Z \to S$ の基底変換の開部分として）局所有限表示かつ平坦でも
あるため、$Z' \to S'$ は有限局所自由である。Morphisms,
Lemma \ref{morphisms-lemma-finite-flat} を参照。
$Y_{S'} \to S'$ は $h$ の基底変換として平坦かつ有限表示で、
幾何学的に被約なファイバーをもつことに注意する。また
$Y_{s'} = Y_s$ は幾何学的に連結である。
Lemma \ref{more-morphisms-lemma-descent-connected-fibres} を $S'$ 上の
$Z' \to Y_{S'}$ に適用すると、
% P05-EN-CH37-0207
(2) を満たす準コンパクト開集合 $V \subset Y_{S'}$ であって、$S'$ 上の
ファイバーが空または幾何学的に連結となるものを得る。
$V \to S'$ は開なので（Morphisms, Lemma \ref{morphisms-lemma-fppf-open}）、
$S'$ を $s'$ のアフィン開近傍で置き換えれば、$V \to S'$ は全射であるとしてよい。
したがって (1) が成り立つ。
\end{proof}

\begin{lemma}
\label{more-morphisms-lemma-cover-by-geometrically-connected}
$f : X \to S$ を、局所有限表示かつ平坦で、幾何学的に被約なファイバーをもつ
スキームの射とする。このときエタール被覆 $\{X_i \to X\}_{i \in I}$ であって、
$X_i \to S$ が $X_i \to S_i \to S$ と分解し、$S_i \to S$ はエタール、
$X_i \to S_i$ は有限表示かつ平坦で、幾何学的に連結かつ幾何学的に被約な
ファイバーをもつものが存在する。
\end{lemma}

\begin{proof}
像が $s \in S$ となる点 $x \in X$ を選ぶ。図式
$$
\xymatrix{
X' \ar[r] \ar[rd] & S' \times_S X \ar[r] \ar[d] & X \ar[d] \\
& S' \ar[r] & S
}
$$
と点 $s' \in S'$, $x' \in X'$, $y \in S' \times_S X$ であって、
$x'$ は $x$ に写り、$(S', s') \to (S, s)$ はエタール近傍、
$(X', x') \to (S' \times_S X, y)$ はエタール近傍
\footnote{実際、この証明は開集合 $X' \subset S' \times_S X$ を与える。}であり、
$X' \to S'$ が幾何学的に連結なファイバーをもつものを構成する。
各 $x \in X$ に対してこれを行えれば、こうして得られるエタール射
$X' \to X$ の族を被覆の要素として補題が従う。
% P05-EN-CH37-0208
最初に $X$ と $S$ をそれぞれ $x$ と $s$ のアフィン開近傍で置き換える。
これにより、$X$ と $S$ がアフィン（したがって $f$ は有限表示）である場合に
帰着する。

\medskip\noindent
% P05-EN-CH37-0209
$\kappa(s)$ の分離代数拡大 $\overline{k}$ を選ぶ。$X_s$ の基底変換を
$X_{\overline{k}}$ と書く。$x \in X_s$ に写る点 $\overline{x}$ を
$X_{\overline{k}}$ に選び、$\overline{x}$ の連結かつ
準コンパクトな開近傍 $\overline{V} \subset X_{\overline{k}}$ を選ぶ。
（体上局所有限型な任意のスキームは、局所 Noether 位相空間として局所連結なので、
これは可能である。）Varieties,
Lemma \ref{varieties-lemma-Galois-action-quasi-compact-open} により、
有限分離拡大 $k'/\kappa(s)$ と準コンパクト開集合 $V' \subset X_{k'}$ であって、
その基底変換が $\overline{V}$ となるものが存在する。とくに $V'$ は $k'$ 上
幾何学的に連結である。Varieties,
Lemma \ref{varieties-lemma-characterize-geometrically-connected} を参照。
Lemma \ref{more-morphisms-lemma-realize-prescribed-residue-field-extension-etale} により、
$\kappa(s')$ が $\kappa(s)$ の拡大として $k'$ と同型となるエタール近傍
$(S', s') \to (S, s)$ を取れる。$\overline{x}$ の像を
$x' \in (S' \times_S X)_{s'} = X_{k'}$ と書く。したがって、
$(S, s)$ を $(S', s')$ で、$(X, x)$ を $(S' \times_S X, x')$ で置き換えると、
次の段落で扱う場合に帰着する。

\medskip\noindent
$x$ を含み、幾何学的に連結な準コンパクト開集合 $V \subset X_s$ が存在すると
仮定する。このとき Lemma \ref{more-morphisms-lemma-fibre-geometrically-connected-reduced} を
適用すると、アフィンなエタール近傍 $(S', s') \to (S, s)$ と
準コンパクト開集合 $X' \subset S' \times_S X$ であって、
$X' \to S'$ が幾何学的に連結なファイバーをもち、$X'$ が $x$ に写る点を含むものを
得る。これで証明が完了する。
\end{proof}

\begin{lemma}
\label{more-morphisms-lemma-normal-morphism-irreducible}
$h : Y \to S$ をスキームの射、$s \in S$ を点、
$T \subset Y_s$ を開部分スキームとする。次を仮定する。
\begin{enumerate}
\item $h$ は有限表示である。
\item $h$ は正規射である。
\item $T$ は $\kappa(s)$ 上幾何学的に既約である。
\end{enumerate}
このとき、アフィンな初等エタール近傍 $(S', s') \to (S, s)$ と
準コンパクト開集合 $V \subset Y_{S'}$ であって、次を満たすものが存在する。
\begin{enumerate}
\item[(a)] $V \to S'$ のすべてのファイバーは幾何学的に整である。
\item[(b)] $V_{s'} = T \times_s s'$.
\end{enumerate}
\end{lemma}

\begin{proof}
Lemma \ref{more-morphisms-lemma-fibre-geometrically-connected-reduced} を適用すると、
アフィンな初等エタール近傍 $(S', s') \to (S, s)$ と準コンパクト開集合
$V \subset Y_{S'}$ であって、$V \to S'$ のすべてのファイバーが幾何学的に連結で、
$V_{s'} = T \times_s s'$ を満たすものを得る。$V$ は $h$ の基底変換の開部分なので、
$V \to S'$ のすべてのファイバーは幾何学的に正規である。
Lemma \ref{more-morphisms-lemma-normal} を参照。とくに、それらは幾何学的に被約である。
証明を終えるには、それらが幾何学的に既約であることを示せばよい。
ところが $t \in S'$ なら $V_t$ は $\kappa(t)$ 上有限型であり、したがって
$V_t \times_{\kappa(t)} \overline{\kappa(t)}$ は $\overline{\kappa(t)}$ 上有限型、
ゆえに Noetherian である。$S' \to S$ の選び方により、スキーム
$V_t \times_{\kappa(t)} \overline{\kappa(t)}$ は連結である。
したがって Properties, Lemma \ref{properties-lemma-normal-Noetherian} により
$V_t \times_{\kappa(t)} \overline{\kappa(t)}$ は既約である。これで示された。
\end{proof}








\section{有限型射の構造への応用}
\label{more-morphisms-section-application-finite-type}

\noindent
本節の結果は \cite{GruRay} に見いだせる。大まかに述べれば、有限型射は
始域と終域上エタール局所的に、有限射と、幾何学的に連結なファイバーをもち、
相対次元がもとの射のファイバー次元に等しい滑らかな射との合成になる、
という結果である。

\begin{lemma}
\label{more-morphisms-lemma-local-structure-finite-type}
$f : X \to S$ を射とする。$x \in X$ を取り、$s = f(x)$ とおく。
$f$ は局所有限型で $n = \dim_x(X_s)$ であると仮定する。
このとき可換図式
$$
\xymatrix{
X \ar[dd] & X' \ar[l]^g \ar[d]^\pi & x \ar@{|->}[dd] &
x' \ar@{|->}[l]  \ar@{|->}[d] \\
& Y \ar[d]^h & & y \ar@{|->}[d] \\
S \ar@{=}[r] & S & s & s \ar@{=}[l]
}
$$
と、$g(x') = x$ を満たす点 $x' \in X'$ であって、$y = \pi(x')$ とおくと
次を満たすものが存在する。
\begin{enumerate}
\item $h : Y \to S$ は相対次元 $n$ の滑らかな射である。
\item $g : (X', x') \to (X, x)$ は初等エタール近傍である。
\item $\pi$ は有限で、$\pi^{-1}(\{y\}) = \{x'\}$ である。
\item $\kappa(y)$ は $\kappa(s)$ の純超越拡大である。
\end{enumerate}
さらに、$f$ が局所有限表示なら $\pi$ は有限表示である。
\end{lemma}

\begin{proof}
問題は $X$ と $S$ 上局所的なので、$X$ と $S$ はアフィンであるとしてよい。
Algebra, Lemma \ref{algebra-lemma-refined-quasi-finite-over-polynomial-algebra}
により、$X$ を $X$ 内の $x$ の標準開近傍で置き換えれば、分解
$$
\xymatrix{
X \ar[r]^\pi & \mathbf{A}^n_S \ar[r] & S
}
$$
であって、$\pi$ が準有限で、$\kappa(\pi(x))$ が $\kappa(s)$ 上純超越となるものが
存在するとしてよい。Lemma \ref{more-morphisms-lemma-etale-makes-quasi-finite-finite-at-point}
により、初等エタール近傍
$$
(Y, y) \to (\mathbf{A}^n_S, \pi(x))
$$
と開集合 $X' \subset X \times_{\mathbf{A}^n_S} Y$ であって、$y$ 上の点
$x'$ を一意に含み、$X' \to Y$ が有限となるものが存在する。
% P05-EN-CH37-0210
これにより (1) -- (4) が成り立つ。最後の主張には Morphisms,
Lemma \ref{morphisms-lemma-finite-presentation-permanence} を用いる。
\end{proof}

\begin{lemma}
\label{more-morphisms-lemma-local-local-structure-finite-type}
\begin{slogan}
有限型射は、エタール近傍において滑らかな射上有限である。
\end{slogan}
$f : X \to S$ を射とする。$x \in X$ を取り、$s = f(x)$ とおく。
$f$ は局所有限型で $n = \dim_x(X_s)$ であると仮定する。
このとき可換図式
$$
\xymatrix{
X \ar[dd] & X' \ar[l]^g \ar[d]^\pi & x \ar@{|->}[dd] &
x' \ar@{|->}[l] \ar@{|->}[d] \\
& Y' \ar[d]^h & & y' \ar@{|->}[d] \\
S & S' \ar[l]_e & s & s' \ar@{|->}[l]
}
$$
と、$g(x') = x$ を満たす点 $x' \in X'$ であって、$y' = \pi(x')$、
$s' = h(y')$ とおくと次を満たすものが存在する。
\begin{enumerate}
\item $h : Y' \to S'$ は相対次元 $n$ の滑らかな射である。
\item $Y' \to S'$ のすべてのファイバーは幾何学的に整である。
\item $g : (X', x') \to (X, x)$ は初等エタール近傍である。
\item $\pi$ は有限で、$\pi^{-1}(\{y'\}) = \{x'\}$ である。
\item $\kappa(y')$ は $\kappa(s')$ の純超越拡大である。
\item $e : (S', s') \to (S, s)$ は初等エタール近傍である。
\end{enumerate}
さらに、$f$ が局所有限表示なら $\pi$ は有限表示である。
\end{lemma}

\begin{proof}
問題は $S$ 上局所的なので、$S$ を $s$ のアフィン開近傍で置き換えてよい。
次に Lemma \ref{more-morphisms-lemma-local-structure-finite-type} を適用して可換図式
$$
\xymatrix{
X \ar[dd] & X' \ar[l]^g \ar[d]^\pi & x \ar@{|->}[dd] &
x' \ar@{|->}[l] \ar@{|->}[d] \\
& Y \ar[d]^h & & y \ar@{|->}[d] \\
S \ar@{=}[r] & S & s & s \ar@{=}[l]
}
$$
を得る。ここで $h$ は相対次元 $n$ の滑らかな射で、$\kappa(y)$ は
$\kappa(s)$ の純超越拡大である。問題は $X$ 上も局所的なので、$Y$ を $y$ の
アフィン近傍で置き換えてよい（また $X'$ を $\pi$ によるその逆像で置き換える）。
$S$ はアフィンなので、これにより $Y \to S$ は準コンパクトかつ分離的かつ滑らか、
とくに有限表示となる。$T$ を $y$ を含む $Y_s$ の連結成分とする。
$Y_s$ は Noetherian なので $T$ は開である。また Varieties,
Lemma \ref{varieties-lemma-geometrically-connected-if-connected-and-point} により、
$T$ は $\kappa(s)$ 上幾何学的に連結である。さらに $T$ は $\kappa(s)$ 上
滑らかなので幾何学的に正規である。Varieties,
Lemma \ref{varieties-lemma-smooth-geometrically-normal} を参照。
したがって $T$ は $\kappa(s)$ 上幾何学的に既約である
（連結な Noetherian 正規スキームは既約である。Properties,
Lemma \ref{properties-lemma-normal-Noetherian} を参照）。
最後に、Lemma \ref{more-morphisms-lemma-smooth-normal} により、滑らかな射 $h$ は正規射である。
% P05-EN-CH37-0211
この時点で、射 $h : Y \to S$ と開集合 $T \subset Y_s$ に対して
Lemma \ref{more-morphisms-lemma-normal-morphism-irreducible} のすべての仮定が成り立つことを
確かめた。Lemma \ref{more-morphisms-lemma-normal-morphism-irreducible} を適用した結果、可換図式
$$
\xymatrix{
X \ar[dd] & X' \ar[l]^g \ar[d]^\pi & X' \times_Y Y' \ar[l] \ar[d] &
x \ar@{|->}[dd] & x' \ar@{|->}[l] \ar@{|->}[d] &
(x', s') \ar@{|->}[l] \ar@{|->}[d] \\
& Y \ar[d]^h & Y' \ar[d] \ar[l] & & y \ar@{|->}[d] &
(y, s') \ar@{|->}[l] \ar@{|->}[d] \\
S \ar@{=}[r] & S & S' \ar[l]_e & s & s \ar@{=}[l] & s' \ar@{|->}[l]
}
$$
に現れる $e : S' \to S$、$s' \in S'$、$Y'$ を得る。ここで
$e : (S', s') \to (S, s)$ は初等エタール近傍で、
$Y' \subset Y_{S'}$ は $S'$ 上のすべてのファイバーが幾何学的に既約となる
開近傍であり、同一視 $Y_s = Y_{S', s'}$ のもとで $Y'_{s'} = T$ を満たす。
$y \in T$ に対応する点を $(y, s') \in Y'$ とする。これはまた、$y$ 上にあり、
剰余体が $\kappa(y)$ に等しく、$S'$ の $s'$ に写る $Y \times_S S'$ の一意な点である。
同様に、$(x', s') \in X' \times_Y Y' \subset X' \times_S S'$ を、$x'$ 上にあり、
剰余体が $\kappa(x')$ に等しく、$s'$ 上にある一意な点とする。
すると、この図式の外側の部分が補題で提示された問題の解である。
細部を若干省略した。
\end{proof}

\begin{lemma}
\label{more-morphisms-lemma-local-local-structure-finite-type-affine}
% P05-EN-CH37-0212
仮定と記法を Lemma \ref{more-morphisms-lemma-local-local-structure-finite-type} と同じとする。
性質 (1) -- (6) に加えて、次も満たすようにできる。
\begin{enumerate}
\item[(7)] $S'$, $Y'$, $X'$ はアフィンである。
\end{enumerate}
\end{lemma}

\begin{proof}
$Y'$ がアフィンなら、$\pi$ は有限なので $X'$ もアフィンであることに注意する。
$s'$ のアフィン開近傍 $U' \subset S'$ を選び、$y'$ のアフィン開近傍
$V' \subset h^{-1}(U')$ を選ぶ。$W' = h(V')$ とおく。これは $U'$ に含まれる
$S'$ 内の $s'$ の開近傍である。Morphisms,
Lemma \ref{morphisms-lemma-smooth-open} を参照。
$s'$ のアフィン開近傍 $U'' \subset W'$ を選ぶ。このとき
$h^{-1}(U'') \cap V'$ は $U'' \times_{U'} V'$ に等しいのでアフィンである。
構成により、$h^{-1}(U'') \cap V' \to U''$ は全射な滑らかな射であり、
そのファイバーは $Y' \to S'$ の幾何学的に整なファイバーの（非空な）
開部分スキームなので、やはり幾何学的に整である。したがって $S'$ を $U''$ で、
$Y'$ を $h^{-1}(U'') \cap V'$ で置き換えてよい。
\end{proof}

\noindent
性質 $\pi^{-1}(\{y'\}) = \{x'\}$ の意義は、次の補題によって部分的に説明される。

\begin{lemma}
\label{more-morphisms-lemma-finite-morphism-single-point-in-fibre}
$\pi : X \to Y$ を有限射とする。$y = \pi(x)$ かつ
$\pi^{-1}(\{y\}) = \{x\}$ を満たす $x \in X$ を取る。このとき次が成り立つ。
\begin{enumerate}
\item $X$ における $x$ の任意の近傍 $U \subset X$ に対し、$y$ の近傍
$V \subset Y$ であって $\pi^{-1}(V) \subset U$ を満たすものが存在する。
\item 環写像 $\mathcal{O}_{Y, y} \to \mathcal{O}_{X, x}$ は有限である。
\item $\pi$ が有限表示なら、$\mathcal{O}_{Y, y} \to \mathcal{O}_{X, x}$ は
有限表示である。
% P05-EN-CH37-0213
\item 任意の準連接 $\mathcal{O}_X$-加群 $\mathcal{F}$ に対し、
$\mathcal{O}_{Y, y}$-加群として $\mathcal{F}_x = \pi_*\mathcal{F}_y$ である。
\end{enumerate}
\end{lemma}

\begin{proof}
第1の主張は純粋に位相的である。$\pi$ が
$\pi^{-1}(\{y\}) = \{x\}$ を満たす連続閉写像であることを用いればよい。
第2、第3の主張を示すためには $X = \Spec(B)$ および $Y = \Spec(A)$ と
仮定してよい。このとき $A \to B$ は有限環写像で、$y$ は $A$ の素イデアル
$\mathfrak p$ に対応し、$\mathfrak p$ の上にある $B$ の素イデアル
$\mathfrak q$ は一意である。このとき
$B_{\mathfrak q} = B_{\mathfrak p}$ である。Algebra,
Lemma \ref{algebra-lemma-unique-prime-over-localize-below} を参照。
言い換えると、写像 $A_{\mathfrak p} \to B_{\mathfrak q}$ は、
$A \to B$ を $\mathfrak p$ で局所化して得られる写像
$A_{\mathfrak p} \to B_{\mathfrak p}$ に等しい。
したがって (2) と (3) は局所化の簡単な性質から従う（細部を若干省略する）。
最後の主張について、ある $B$-加群 $M$ に対して
$\mathcal{F} = \widetilde M$ であると仮定する。このとき
% P05-EN-CH37-0214
$\mathcal{F} = M_{\mathfrak q}$ かつ
$\pi_*\mathcal{F}_y = M_{\mathfrak p}$ である。上の議論により、これらの局所化は
一致する。別の方法として、(1) と茎の定義を用いれば
$\mathcal{F}_x = \pi_*\mathcal{F}_y$ を直接示せる。
\end{proof}






\section{fppf 位相への応用}
\label{more-morphisms-section-application-fppf}

\noindent
上のエタール局所化の手法を用いると、fppf 位相は Zariski 被覆と、$f$ が全射な
有限局所自由射である形の被覆 $\{f : T \to S\}$ とによって「生成される」位相に
等しい、という次の結果を証明できる。

\begin{lemma}
\label{more-morphisms-lemma-dominate-fppf-etale-locally}
$S$ をスキームとし、$\{S_i \to S\}_{i \in I}$ を fppf 被覆とする。
このとき次が存在する。
\begin{enumerate}
\item エタール被覆 $\{S'_a \to S\}$。
\item 全射な有限局所自由射 $V_a \to S'_a$。
\end{enumerate}
さらに、fppf 被覆 $\{V_a \to S\}$ は与えられた被覆
$\{S_i \to S\}$ の細分である。
\end{lemma}

\begin{proof}
Lemma \ref{more-morphisms-lemma-qf-fp-flat-dominates-fppf} により、各 $S_i \to S$ は
局所準有限であると仮定してよい。

\medskip\noindent
点 $s \in S$ を固定する。$i \in I$ と、$s$ に写る点 $s_i \in S_i$ を選ぶ。
% P05-EN-CH37-0215
初等エタール近傍 $(S', s) \to (S, s)$ であって、開集合
$$
S_i \times_S S'  \supset V
$$
が存在し、この開集合が $s \in S'$ に写る点 $v \in V$ を一意に含み、
$V \to S'$ が有限となるものを選ぶ。
Lemma \ref{more-morphisms-lemma-etale-makes-quasi-finite-finite-at-point} を参照。
$V \to S'$ は有限局所自由である。実際、これは有限であり、また
$S_i \times_S S' \to S'$ は射 $S_i \to S$ の基底変換として平坦かつ
局所有限表示である。Morphisms, Lemmas
\ref{morphisms-lemma-base-change-finite-presentation},
\ref{morphisms-lemma-base-change-flat}, and
\ref{morphisms-lemma-finite-flat} を参照。
したがって $V \to S'$ は開であり、$S'$ を縮小すれば、$V \to S'$ は
全射な有限局所自由射であるとしてよい。$S$ の各点に対してこれを行えるので、
$\{S_i \to S\}$ は、各 $V_a \to S$ が
$V_a \to S'_a \to S$ と分解し、$S'_a \to S$ がエタール、
$V_a \to S'_a$ が全射な有限局所自由射となる形の被覆
$\{V_a \to S\}_{a \in A}$ によって細分される。
\end{proof}

\begin{lemma}
\label{more-morphisms-lemma-dominate-fppf}
$S$ をスキームとし、$\{S_i \to S\}_{i \in I}$ を fppf 被覆とする。
このとき次が存在する。
\begin{enumerate}
\item Zariski 開被覆 $S = \bigcup U_j$。
\item 全射な有限局所自由射 $W_j \to U_j$。
\item Zariski 開被覆 $W_j = \bigcup_k W_{j, k}$。
\item 全射な有限局所自由射 $T_{j, k} \to W_{j, k}$。
\end{enumerate}
さらに、fppf 被覆 $\{T_{j, k} \to S\}$ は与えられた被覆
$\{S_i \to S\}$ の細分である。
\end{lemma}

\begin{proof}
Lemma \ref{more-morphisms-lemma-dominate-fppf-etale-locally} で得られた fppf 被覆を
$\{V_a \to S\}_{a \in A}$ とする。すなわち、この被覆は
$\{S_i \to S\}$ を細分し、各 $V_a \to S$ は
$V_a \to S'_a \to S$ と分解する。ここで $S'_a \to S$ はエタール、
$V_a \to S'_a$ は全射な有限局所自由射である。

\medskip\noindent
Remark \ref{more-morphisms-remark-topologies} により、Zariski 開被覆 $S = \bigcup U_j$、
各 $j$ に対する全射な有限局所自由射 $W_j \to U_j$、および各 $j$ に対する
Zariski 開被覆 $\{W_{j, k} \to W_j\}$ であって、族
$\{W_{j, k} \to S\}$ がエタール被覆 $\{S'_a \to S\}$ を細分するものが存在する。
すなわち、各組 $j, k$ に対して $a(j, k)$ が存在し、射
$W_{j, k} \to S$ は
% P05-EN-CH37-0216
$W_{j, k} \to S'_a \to S$ と分解する。
$T_{j, k} = W_{j, k} \times_{S'_a} V_a$ とおけば、すべて明らかである。
\end{proof}

\begin{lemma}
\label{more-morphisms-lemma-extend-integral-surjective-morphisms}
$S$ をスキームとする。$U \subset S$ が開で、$V \to U$ が全射な整射なら、
全射な整射 $\overline{V} \to S$ であって、$\overline{V} \times_S U$ が
$U$ 上のスキームとして $V$ と同型となるものが存在する。
\end{lemma}

\begin{proof}
% P05-EN-CH37-0217
$V' \to S$ を $U$ における $S$ の正規化とする。Morphisms,
Section \ref{morphisms-section-normalization-X-in-Y} を参照。
構成により $V' \to S$ は整射である。Morphisms, Lemmas
\ref{morphisms-lemma-normalization-localization} および
\ref{morphisms-lemma-normalization-in-integral} により、$V'$ における $U$ の逆像は
$V$ である。$S \setminus U$ 上に誘導される被約スキーム構造を $Z$ とする。
このとき $\overline{V} = V' \amalg Z$ が所望のものである。
\end{proof}

\begin{lemma}
\label{more-morphisms-lemma-extend-finite-surjective-morphisms}
$S$ を準コンパクトかつ準分離的なスキームとする。$U \subset S$ が準コンパクトな
開部分で、$V \to U$ が全射な有限射なら、全射な有限射
$\overline{V} \to S$ であって、$\overline{V} \times_S U$ が $U$ 上のスキームとして
$V$ と同型となるものが存在する。
\end{lemma}

\begin{proof}
Zariski の主定理
（Lemma \ref{more-morphisms-lemma-quasi-finite-separated-pass-through-finite}）により、
$V$ は $S$ 上有限なスキーム $V'$ の準コンパクトな開部分であるとしてよい。
$V'$ を $V$ のスキーム論的像で置き換えると、$V$ は $V'$ に稠密であるとしてよい。
$V$ は $U$ 上有限なので $V \to V' \times_S U$ は閉であり、したがって
$V' \times_S U = V$ である。$S \setminus U$ 上に誘導される被約スキーム構造を
$Z$ とする。このとき $\overline{V} = V' \amalg Z$ が所望のものである。
\end{proof}

\begin{lemma}
\label{more-morphisms-lemma-dominate-fppf-integral}
$S$ をスキームとし、$\{S_i \to S\}_{i \in I}$ を fppf 被覆とする。
このとき全射な整射 $S' \to S$ と開被覆 $S' = \bigcup U'_\alpha$ であって、
各 $\alpha$ に対して射 $U'_\alpha \to S$ がある $i$ に対する
$S_i \to S$ を経由するものが存在する。
\end{lemma}

\begin{proof}
Lemma \ref{more-morphisms-lemma-dominate-fppf} のように $S = \bigcup U_j$、$W_j \to U_j$、
$W_j = \bigcup W_{j, k}$、および $T_{j, k} \to W_{j, k}$ を選ぶ。
Lemma \ref{more-morphisms-lemma-extend-integral-surjective-morphisms} により、
$W_j \to U_j$ を全射な整射 $\overline{W}_j \to S$ に延長できる。
次に $T_{j, k} \to W_{j, k}$ を全射な整射
$\overline{T}_{j, k} \to \overline{W}_j$ に延長できる。
$\overline{T}_j$ を $\overline{W}_j$ 上のすべてのスキーム
$\overline{T}_{j, k}$ の積に等しいものとする
（Limits, Lemma \ref{limits-lemma-infinite-product}）。
次に $S'$ を $S$ 上のすべてのスキーム $\overline{T}_j$ の積に等しいものとする。
$x \in S'$ なら、$S$ における $x$ の像が $U_j$ に入るような $j$ が存在する。
したがって、射影 $S' \to \overline{W}_j$ による $x$ の像が
$W_{j, k}$ に入るような $k$ が存在する。ゆえに、射影
$S' \to \overline{T}_j \to \overline{T}_{j, k}$ によって、点 $x$ は
$T_{j, k}$ に入る。そして $T_{j, k} \to S$ は、ある $i$ に対する $S_i$ を経由する。
最後に Limits, Lemmas \ref{limits-lemma-infinite-product-integral} および
\ref{limits-lemma-infinite-product-surjective} により、射 $S' \to S$ は整かつ全射である。
\end{proof}

\begin{lemma}
\label{more-morphisms-lemma-dominate-fppf-finite}
$S$ を準コンパクトかつ準分離的なスキームとし、
$\{S_i \to S\}_{i \in I}$ を fppf 被覆とする。このとき有限表示の全射な有限射
$S' \to S$ と開被覆 $S' = \bigcup U'_\alpha$ であって、各 $\alpha$ に対して
射 $U'_\alpha \to S$ がある $i$ に対する $S_i \to S$ を経由するものが存在する。
\end{lemma}

\begin{proof}
% P05-EN-CH37-0218
Lemma \ref{more-morphisms-lemma-dominate-fppf-integral} で得られた整かつ全射な射を
$Y \to X$ とする。$V_j \to X$ が $S_{i(j)}$ を経由するような有限アフィン開被覆
$Y = \bigcup V_j$ を選ぶ。$Y_\lambda \to X$ が有限かつ有限表示となる形で
$Y = \lim Y_\lambda$ と書ける。Limits,
Lemma \ref{limits-lemma-integral-limit-finite-and-finite-presentation} を参照。
十分大きい $\lambda$ に対して、$Y$ への逆像が $V_j$ となるアフィン開集合
$V_{\lambda, j} \subset Y_\lambda$ を取れる。Limits,
Lemma \ref{limits-lemma-descend-opens} を参照。さらに $\lambda$ を大きくすれば、
$X$ 上の射 $V_j \to S_{i(j)}$ は $X$ 上の射
$V_{\lambda, j} \to S_{i(j)}$ から得られる。Limits, Proposition
\ref{limits-proposition-characterize-locally-finite-presentation} を参照。
この $\lambda$ に対して $S' = Y_\lambda$ とおけば証明が完了する。
\end{proof}

\begin{lemma}
\label{more-morphisms-lemma-fppf-ph}
スキームの fppf 被覆は ph 被覆である。
\end{lemma}

\begin{proof}
$\{T_i \to T\}$ をスキームの fppf 被覆とする。Topologies,
Definition \ref{topologies-definition-fppf-covering} を参照。
$T_i \to T$ は局所有限型であることに注意する。$U \subset T$ をアフィン開集合とする。
$\{T_i \times_T U \to U\}$ が標準 ph 被覆によって細分されることを示せば十分である。
Topologies, Definition \ref{topologies-definition-ph-covering} を参照。
これは Lemma \ref{more-morphisms-lemma-dominate-fppf-finite} と、有限射が固有であるという事実
（Morphisms, Lemma \ref{morphisms-lemma-finite-proper}）から直ちに従う。
\end{proof}

\begin{remark}
\label{more-morphisms-remark-change-topologies}
Lemma \ref{more-morphisms-lemma-fppf-ph} の帰結として比較射
$$
\epsilon : (\Sch/S)_{ph} \longrightarrow (\Sch/S)_{fppf}
$$
を得る。これは基礎圏上の恒等関手によって与えられるサイトの射である
（Topologies, Remark \ref{topologies-remark-choice-sites} のようにサイトを
適切に選ぶ）。関手 $\epsilon_*$ は基礎前層上の恒等関手であり、
% P05-EN-CH37-0219
関手 $\epsilon^{-1}$ は fppf 層にその ph 層化を対応させる。
さらに合成により、ph 位相を syntomic、滑らか、エタール、および Zariski 位相と
比較できる。
\end{remark}




\section{準射影スキーム}
\label{more-morphisms-section-quasi-projective}

\noindent
「準射影スキーム」という用語はまだ定義していない。
一つの可能な定義は、豊富な可逆層をもつスキームとすることであろう。
しかし、$X$ が基礎スキーム $S$ 上のスキームであるとき、射
$X \to S$ が準射影的であるなら、{\it $X$ は $S$ 上準射影的である}
という（Morphisms, Definition
\ref{morphisms-definition-quasi-projective}）。任意のスキームの恒等射は
準射影的なので、$S$ 上準射影的なスキームが豊富な可逆層をもつとは
限らないことがわかる。このため、「準射影スキーム」という用語は
未定義のままにしておく方がよいと思われる。

\begin{lemma}
\label{more-morphisms-lemma-quasi-projective}
$S$ を豊富な可逆層をもつスキームとする。
$f : X \to S$ をスキームの射とする。次の条件は同値である。
\begin{enumerate}
\item $X \to S$ は準射影的である。
\item $X \to S$ は H-準射影的である。
\item $S$ 上のスキームの準コンパクトな開埋め込み $X \to X'$ であって、
$X' \to S$ が射影的となるものが存在する。
\item $X \to S$ は有限型であり、$X$ は豊富な可逆層をもつ。
\item $X \to S$ は有限型であり、$f$-非常に豊富な可逆層が存在する。
\end{enumerate}
\end{lemma}

\begin{proof}
(2) $\Rightarrow$ (1) は Morphisms, Lemma
\ref{morphisms-lemma-H-quasi-projective-quasi-projective} である。
(1) $\Rightarrow$ (2) は Morphisms, Lemma
\ref{morphisms-lemma-projective-over-quasi-projective-is-H-projective} である。
(2) $\Rightarrow$ (3) は Morphisms, Lemma
\ref{morphisms-lemma-H-quasi-projective-open-H-projective} である。

\medskip\noindent
(3) のような $X \subset X'$ が与えられていると仮定する。特に
$X \to S$ は有限型である。Morphisms, Lemma
% P05-EN-CH37-0220
\ref{morphisms-lemma-H-quasi-projective-open-H-projective} により、射
$X \to S$ は H-射影的である。したがって準コンパクトな埋め込み
$i : X \to \mathbf{P}^n_S$ が存在する。ゆえに
$\mathcal{L} = i^*\mathcal{O}_{\mathbf{P}^n_S}(1)$ は
$f$-非常に豊富である。$X \to S$ は準コンパクトなので、Morphisms,
Lemma \ref{morphisms-lemma-ample-very-ample} から $\mathcal{L}$ は
$f$-豊富であると結論する。したがって定義により $X \to S$ は
準射影的である。

\medskip\noindent
(4) $\Rightarrow$ (2) は Morphisms, Lemma
\ref{morphisms-lemma-quasi-projective-finite-type-over-S} である。

\medskip\noindent
(5) $\Rightarrow$ (4) は Morphisms, Lemma
\ref{morphisms-lemma-pullback-ample-tensor-relatively-ample} から従う。

\medskip\noindent
(1) と (5) の同値性は Morphisms, Lemmas
\ref{morphisms-lemma-ample-very-ample} および
\ref{morphisms-lemma-finite-type-ample-very-ample} から従う。
\end{proof}

\begin{lemma}
\label{more-morphisms-lemma-category-quasi-projective}
$S$ を豊富な可逆層をもつスキームとする。$\text{QP}_S$ を、
Lemma \ref{more-morphisms-lemma-quasi-projective} の同値な条件を満たす $S$ 上の
スキームからなる充満部分圏とする。
\begin{enumerate}
\item $S' \to S$ がスキームの射で、$S'$ が豊富な可逆層をもつなら、
基底変換は関手 $\text{QP}_S \to \text{QP}_{S'}$ を定める。
\item $X \in \text{QP}_S$ かつ $Y \in \text{QP}_X$ なら、
$Y \in \text{QP}_S$ である。
\item 圏 $\text{QP}_S$ はファイバー積に関して閉じている。
\item 圏 $\text{QP}_S$ は有限非交和に関して閉じている。
\item $X \to S$ が射影的なら、$X \in \text{QP}_S$ である。
\item $X \to S$ が有限型かつ準アフィンなら、$X$ は
$\text{QP}_S$ に属する。
\item $X \to S$ が準有限かつ分離的なら、$X \in \text{QP}_S$ である。
\item $X \to S$ が準コンパクトな埋め込みなら、
$X \in \text{QP}_S$ である。
% P05-EN-CH37-0221
\item ここにさらに項目を追加する。
\end{enumerate}
\end{lemma}

\begin{proof}
(1) は Morphisms, Lemma
\ref{morphisms-lemma-base-change-quasi-projective} から従う。

\medskip\noindent
(2) は Lemma \ref{more-morphisms-lemma-quasi-projective} の第4の特徴づけから従う。

\medskip\noindent
$X \to S$ と $Y \to S$ が準射影的なら、Morphisms, Lemma
\ref{morphisms-lemma-base-change-quasi-projective} により
$X \times_S Y \to Y$ は準射影的である。したがって (3) は (2) から従う。

\medskip\noindent
$X = Y \amalg Z$ がスキームの非交和であり、$\mathcal{L}$ が可逆
$\mathcal{O}_X$-加群で、$\mathcal{L}|_Y$ と $\mathcal{L}|_Z$ が豊富なら、
$\mathcal{L}$ は豊富である（詳細は省略する）。したがって (4) は
Lemma \ref{more-morphisms-lemma-quasi-projective} の第4の特徴づけから従う。

\medskip\noindent
(5) は Morphisms, Lemma
\ref{morphisms-lemma-projective-quasi-projective} から従う。

\medskip\noindent
(6) は Morphisms, Lemma
\ref{morphisms-lemma-quasi-affine-finite-type-quasi-projective} から従う。

\medskip\noindent
(7) は (6) および Lemma
\ref{more-morphisms-lemma-quasi-finite-separated-quasi-affine} から従う。

\medskip\noindent
(8) は (7) および Morphisms, Lemma
\ref{morphisms-lemma-immersion-locally-quasi-finite} から従う。
\end{proof}

\noindent
次の補題は本来この節に属するものではないが、ほかに置くのに適した場所も
見当たらない。

\begin{lemma}
\label{more-morphisms-lemma-integral-over-quasi-affine}
$X$ を準アフィン・スキームとする。$f : U \to X$ を整射とする。
このとき $U$ は準アフィンであり、図式
$$
\xymatrix{
U \ar[r] \ar[d] & \Spec(\Gamma(U, \mathcal{O}_U)) \ar[d] \\
X \ar[r] & \Spec(\Gamma(X, \mathcal{O}_X))
}
$$
はデカルトである。
\end{lemma}

\begin{proof}
整射はアフィンであり、アフィン射は準アフィンであり、スキームが
準アフィンであることと $\Spec(\mathbf{Z})$ への構造射が準アフィンで
あることは同値であり、さらに準アフィン射の合成は準アフィンなので、
スキーム $U$ は準アフィンである。最初の二つの主張は定義から直ちに従い、
三つ目は Morphisms, Lemma
\ref{morphisms-lemma-composition-quasi-affine} である。$U' =
X \times_{\Spec(\Gamma(X, \mathcal{O}_X))} \Spec(\Gamma(U, \mathcal{O}_U))$
とおき、拡張された図式
$$
\xymatrix{
U \ar[r]_j \ar[rd] & U' \ar[d] \ar[r] &
\Spec(\Gamma(U, \mathcal{O}_U)) \ar[d] \\
& X \ar[r] & \Spec(\Gamma(X, \mathcal{O}_X))
}
$$
を考える。整射は普遍閉であること（Morphisms, Lemma
\ref{morphisms-lemma-integral-universally-closed}）および図式の鉛直射が
% P05-EN-CH37-0222
分離的であることと合わせて、Morphisms, Lemma
\ref{morphisms-lemma-image-proper-scheme-closed} により射 $j$ は閉である。
一方、Properties, Lemma \ref{properties-lemma-quasi-affine} により
補題の図式の水平射は開なので、$j$ は開である。したがって $j$ は
$U$ を $U'$ の開かつ閉な部分スキームと同一視する。もし
$U \not = U'$ なら、$U$ は $U'$ で稠密でなく、まして
$\Gamma(U, \mathcal{O}_U)$ のスペクトルでも稠密でない。しかし、
$\Spec(\Gamma(U, \mathcal{O}_U))$ における $U$ のスキーム論的像は
$\Spec(\Gamma(U, \mathcal{O}_U))$ である。実際、$U$ が経由する閉部分
スキームを切り出す $\Gamma(U, \mathcal{O}_U)$ の任意のイデアルは
零でなければならない。したがって例えば Morphisms, Lemma
\ref{morphisms-lemma-quasi-compact-scheme-theoretic-image} により、$U$ は
$\Spec(\Gamma(U, \mathcal{O}_U))$ で稠密である。ゆえに $U = U'$ となり、
証明が完了する。
\end{proof}









\section{射影スキーム}
\label{more-morphisms-section-projective}

\noindent
この節は、射影射についての Section \ref{more-morphisms-section-quasi-projective} の
類似である。

\begin{lemma}
\label{more-morphisms-lemma-projective}
$S$ を豊富な可逆層をもつスキームとする。
$f : X \to S$ をスキームの射とする。次の条件は同値である。
\begin{enumerate}
\item $X \to S$ は射影的である。
\item $X \to S$ は H-射影的である。
\item $X \to S$ は準射影的かつ固有である。
\item $X \to S$ は H-準射影的かつ固有である。
\item $X \to S$ は固有であり、$X$ は豊富な可逆層をもつ。
\item $X \to S$ は固有であり、$f$-豊富な可逆層が存在する。
\item $X \to S$ は固有であり、$f$-非常に豊富な可逆層が存在する。
\item $\mathcal{A}_1$ が有限型 $\mathcal{O}_S$-加群で、
$\mathcal{A}_0$ 上 $\mathcal{A}_1$ により生成される準連接次数付き
$\mathcal{O}_S$-代数 $\mathcal{A}$ であって、
$X = \underline{\text{Proj}}_S(\mathcal{A})$ となるものが存在する。
\end{enumerate}
\end{lemma}

\begin{proof}
まず、いずれの条件のもとでも射 $f$ は固有であることに注意する。
Morphisms, Lemmas \ref{morphisms-lemma-H-projective} および
\ref{morphisms-lemma-locally-projective-proper} を参照せよ。したがって、
$f$ が固有射である場合にこれらの概念の同値性を証明すれば十分である。
以下ではこのことを断りなく用いる。

\medskip\noindent
(1) $\Leftrightarrow$ (3) および (2) $\Leftrightarrow$ (4) は
Morphisms, Lemma
\ref{morphisms-lemma-projective-is-quasi-projective-proper} である。

\medskip\noindent
(2) $\Rightarrow$ (1) は Morphisms, Lemma
\ref{morphisms-lemma-H-projective} である。

\medskip\noindent
(1) $\Rightarrow$ (2) および (3) $\Rightarrow$ (4) は Morphisms, Lemma
\ref{morphisms-lemma-projective-over-quasi-projective-is-H-projective} である。

\medskip\noindent
(1) $\Rightarrow$ (7) は Morphisms, Definitions
\ref{morphisms-definition-projective} および
\ref{morphisms-definition-very-ample} から直ちに従う。

\medskip\noindent
(3) と (6) は Morphisms, Definition
\ref{morphisms-definition-quasi-projective} により同値である。

\medskip\noindent
したがって (1) -- (4), (6) は同値であり、(7) を含意する。Lemma
\ref{more-morphisms-lemma-quasi-projective} により (3), (5), (7) は同値である。
よって (1) -- (7) は同値であることがわかる。

\medskip\noindent
Divisors, Lemma \ref{divisors-lemma-relative-proj-projective} により、
(8) は (1) を含意する。逆に (2) が成り立つなら、閉埋め込み
$$
i :
X
\longrightarrow
\mathbf{P}^n_S = \underline{\text{Proj}}_S(\mathcal{O}_S[T_0, \ldots, T_n]).
$$
を選べる。この等式については Constructions, Lemma
\ref{constructions-lemma-projective-space-bundle} を参照せよ。Divisors,
Lemma \ref{divisors-lemma-closed-subscheme-proj} により、$X$ は
$\mathcal{O}_S[T_0, \ldots, T_n]$ の準連接次数付き商代数
$\mathcal{A}$ の相対 Proj であることがわかる。このとき
$\mathcal{A}$ は (8) の条件を満たす。
\end{proof}

\begin{lemma}
\label{more-morphisms-lemma-category-projective}
$S$ を豊富な可逆層をもつスキームとする。$\text{P}_S$ を、Lemma
\ref{more-morphisms-lemma-projective} の同値な条件を満たす $S$ 上のスキームからなる
充満部分圏とする。
\begin{enumerate}
\item $S' \to S$ がスキームの射で、$S'$ が豊富な可逆層をもつなら、
基底変換は関手 $\text{P}_S \to \text{P}_{S'}$ を定める。
\item $X \in \text{P}_S$ かつ $Y \in \text{P}_X$ なら、
$Y \in \text{P}_S$ である。
\item 圏 $\text{P}_S$ はファイバー積に関して閉じている。
\item 圏 $\text{P}_S$ は有限非交和に関して閉じている。
\item $X \to S$ が有限なら、$X$ は $\text{P}_S$ に属する。
% P05-EN-CH37-0223
\item ここにさらに項目を追加する。
\end{enumerate}
\end{lemma}

\begin{proof}
(1) は Morphisms, Lemma
\ref{morphisms-lemma-base-change-projective} から従う。

\medskip\noindent
(2) は Lemma \ref{more-morphisms-lemma-projective} の第5の特徴づけと、固有射の合成が
固有であること（Morphisms, Lemma
\ref{morphisms-lemma-composition-proper}）から従う。

\medskip\noindent
$X \to S$ と $Y \to S$ が射影的なら、Morphisms, Lemma
\ref{morphisms-lemma-base-change-projective} により
$X \times_S Y \to Y$ は射影的である。したがって (3) は (2) から従う。

\medskip\noindent
$X = Y \amalg Z$ がスキームの非交和であり、$\mathcal{L}$ が可逆
$\mathcal{O}_X$-加群で、$\mathcal{L}|_Y$ と $\mathcal{L}|_Z$ が豊富なら、
$\mathcal{L}$ は豊富である（詳細は省略する）。したがって (4) は
Lemma \ref{more-morphisms-lemma-projective} の第5の特徴づけから従う。

\medskip\noindent
(5) は Morphisms, Lemma
\ref{morphisms-lemma-finite-projective} から従う。
\end{proof}

\noindent
次に、少し異なる種類の結果を与える。

\begin{lemma}
\label{more-morphisms-lemma-ample-in-neighbourhood}
\begin{reference}
\cite[IV Corollary 9.6.4]{EGA}
\end{reference}
$f : X \to Y$ をスキームの固有射とする。$\mathcal{L}$ を可逆
$\mathcal{O}_X$-加群とする。$y \in Y$ を、$\mathcal{L}_y$ が $X_y$ 上
豊富となる点とする。このとき $y$ の開近傍 $V \subset Y$ であって、
$\mathcal{L}|_{f^{-1}(V)}$ が $f^{-1}(V)/V$ 上豊富となるものが存在する。
\end{lemma}

\begin{proof}
$Y$ はアフィンであると仮定してよい。このとき、有向集合 $I$ と
スキームの射 $X_i \to Y_i$ の逆系であって、$Y_i$ は
$\mathbf{Z}$ 上有限型、遷移射 $X_i \to X_{i'}$ と $Y_i \to Y_{i'}$ は
アフィン、$X_i \to Y_i$ は固有であり、かつ
$X \to Y = \lim (X_i \to Y_i)$ となるものを得る。Limits, Lemma
\ref{limits-lemma-proper-limit-of-proper-finite-presentation-noetherian}
を参照せよ。$I$ を縮小したのち、$\mathcal{L}$ に引き戻される可逆
$\mathcal{O}_{X_i}$-加群 $\mathcal{L}_i$ の両立する系があると仮定してよい。
Limits, Lemma \ref{limits-lemma-descend-invertible-modules} を参照せよ。
$y_i \in Y_i$ を $y$ の像とする。このとき
$\kappa(y) = \colim \kappa(y_i)$ である。したがって Limits, Lemma
\ref{limits-lemma-limit-ample} により、ある $i$ について
$\mathcal{L}_{i, y_i}$ は $X_{i, y_i}$ 上豊富である。Cohomology of
Schemes, Lemma \ref{coherent-lemma-ample-in-neighbourhood} により、
$y_i$ の開近傍 $V_i \subset Y_i$ であって、$\mathcal{L}_i$ の
$f_i^{-1}(V_i)$ への制限が $V_i$ に関して豊富となるものを得る。
$V \subset Y$ を $V_i$ の逆像とすれば証明が完了する（ヒント：
Morphisms, Lemma \ref{morphisms-lemma-ample-base-change}、
$X \to Y \times_{Y_i} X_i$ がアフィンであること、およびアフィン射による
豊富な可逆層の引き戻しは Morphisms, Lemma
\ref{morphisms-lemma-pullback-ample-tensor-relatively-ample} により
豊富であることを用いる）。
\end{proof}







\section{Proj と Spec}
\label{more-morphisms-section-proj-spec}

\noindent
この節では、次数付き環の Proj とスペクトルの関係を明らかにする。

\medskip\noindent
$R$ を環とする。$A$ を次数付き $R$-代数とする。Algebra, Section
\ref{algebra-section-graded} を参照せよ。$m \geq 0$ に対し
$A_{\geq m} = \bigoplus_{d \geq m} A_d$ と書く。次数付き環
$$
B = \bigoplus\nolimits_{d \geq 0} A_{\geq d}
$$
を考える。$d' \geq d$ および $a \in A_{d'}$ に対し、$a$ に対応する
$B_d$ の元を $a^{(d)} \in B$ と書く。二つの明らかな環準同型を
$\sigma : A \to B$ および $\psi : A \to B$ と書く。すなわち
$a \in A_d$ なら、$\sigma(a) = a^{(0)}$ および
$\psi(a) = a^{(d)}$ である。このとき $\psi$ は次数付き環準同型であり、
$\sigma$ は $B$ を $A$ 上の次数付き代数にする。また、全射な次数付き
環準同型 $\tau : B \to A$ があり、$d' \geq d$ および
$a \in A_{d'}$ に対して、$d' > d$ なら $a^{(d)}$ を $0$ に、
$d' = d$ なら $a$ に送る。

\medskip\noindent
アフィン・スキームとスペクトル。
$X = \Spec(A)$ とおく。無関係イデアル $A_+$ は閉部分スキーム
$Z = V(A_+) = \Spec(A/A_+) = \Spec(A_0)$ を切り出す。
$U = X \setminus Z$ とおく。
$$
U \longrightarrow X \longrightarrow Z
$$
射影スキームと Proj。$P = \text{Proj}(A)$ とおく。$P$ を
$\Spec(A_0) = Z$ 上のスキームとみなすことができ、実際そのように
みなす。$L = \text{Proj}(B)$ とおく。$L$ を
$\Spec(B_0) = \Spec(A) = X$ 上のスキームとみなすことができ、実際
そのようにみなす。ここで $B_0$ と $A$ の同一視は $\sigma$ により
与えられることに注意する。全射 $\tau$ は閉埋め込み
$0 : P \to L$ を定める。$A \xrightarrow{\sigma} B \to A$ は写像
$A \to A_0 \to A$ に等しいので、図式
$$
\xymatrix{
P \ar[d] \ar[r]_0 & L \ar[d] \\
Z \ar[r] & X
}
$$
は可換である。

\medskip\noindent
$\psi$ は射 $L \to P$ を定めると主張する。これを見るには、
Constructions, Lemma \ref{constructions-lemma-morphism-proj} により、
$B_+ \not \subset \mathfrak p$ を満たす任意の斉次素イデアル
$\mathfrak p \subset B$ に対して
$\psi(A_+) \not \subset \mathfrak p$ を確かめれば十分である。まず、
斉次元 $g \in B_+$ で $g \not \in \mathfrak p$ となるものを選ぶ。
このとき、ある $d_i \geq d$ に対する $a_i \in A_{d_i}$ を用いて、
$g$ を有限和 $g = \sum a_i^{(d)}$ と書ける。したがって
$d' \geq d$ と $a \in A_{d'}$ であって
$a^{(d)} \not \in \mathfrak p$ となるものが存在する。このとき
$$
(a^{(d)})^{d'} =
(a^{d'})^{(d'd)} =
a^{(d)} (a^{d' - 1})^{(d(d' - 1))} =
\psi(a) (a^{d' - 1})^{(d(d' - 1))}
$$
（この記法には改善の余地がある）は $\mathfrak p$ に属さない。
したがって $\psi(a) \not \in \mathfrak p$ であり、主張が示された。
よって上の図式を可換図式
$$
\xymatrix{
P \ar[d] \ar[r]_0 & L \ar[d] \ar[r]_\pi & P \ar[d] \\
Z \ar[r] & X \ar[r] & Z
}
$$
へ拡張できる。ここで $X \to Z$ は $A_0 \to A$ により与えられる。
$\tau \circ \psi = \text{id}_A$ なので
$\pi \circ 0 = \text{id}_P$ であることがわかる。

\medskip\noindent
$\pi$ はアフィン射であることに注意する。これは Constructions, Lemma
\ref{constructions-lemma-morphism-proj} の構成から明らかである。実際、
ある $d > 0$ に対して $f \in A_d$ なら、$g = \psi(f)$ とおくと
$\pi^{-1}(D_+(f)) = D_+(g)$ である。この場合、斉次成分の等式
$$
(B[1/g])_{m'} = \bigoplus\nolimits_{m \geq m'} (A[1/f])_m
$$
を得る。この同型はさらなる局所化と両立する。$m' = 0$ ととると、
$\pi_*\mathcal{O}_L$ は $m \geq 0$ にわたる $\mathcal{O}_P(m)$ の
直和であることがわかる\footnote{同様に
$\pi_*\mathcal{O}_L(i) = \bigoplus_{m \geq -i} \mathcal{O}_P(m)$
が従う。}。したがって $L$ は相対スペクトル
$$
L = \underline{\Spec}_P
\left(
\bigoplus\nolimits_{m \geq 0} \mathcal{O}_P(m)
\right)
$$
と同一視される。特に $L \to P$ は錐である\footnote{$L$ はしばしば
$P$ 上の直線束となる。以下を参照せよ。}。Constructions, Section
\ref{constructions-section-cone} を参照せよ。さらに、
$0 : P \to L$ が錐の頂点であることは明らかである。

\medskip\noindent
前段落と同様に、ある $d > 0$ に対して $f \in A_d$ とし、
$g = \psi(f) \in B_d$ とする。環準同型
$$
\xymatrix{
A_0 \ar[r] \ar[d] & A \ar[d]^\sigma \ar[r] & A_0 \ar[d] \\
(A[1/f])_0 \ar[r]^-\psi &
(B[1/g])_0 = \bigoplus\nolimits_{m \geq 0} (A[1/f])_m
\ar[r]^-\tau &
(A[1/f])_0
}
$$
% P05-EN-CH37-0224
の構造を見ると、次数付き環におけるいくつかの計算\footnote{(1) と
(2) は明らかである。(3) を見るには、$a \in A_d$ なら
$\sigma(a) = \sigma(f) \psi(a/f)$ であることに注意する。(4) については、
$b/g^m$ が $\tau$ の核に属することと、$b \in A_{\geq md}$ が
$A_{md}$ で零に写ることは同値である。したがって、$m' > md$ かつ
$a \in A_{m'}$ なら、$a^{(md)}/g^m$ のある冪が $\sigma(f)$ により
生成されるイデアルに属することを示せば十分である。
$em' - emd \geq d$ を満たす $e$ をとる。このとき
$$
(a^{(md)}/g^m)^e = (a^e)^{(emd)}/g^{em} =
(fa^e)^{(emd + d)}/g^{em + 1} = \sigma(f) \cdot (a^e)^{(emd + d)}/g^{em + 1}
$$
となり、望む結果を得る（ひどい記法をお詫びする）。(5) を見るには
前と同様に論じ、$d = 1$ のとき
$a^{(md)}/g^m = \sigma(f) \cdot a^{(md + 1)}/g^{m + 1}$ であることに
注意する。} により、次が示される。
\begin{enumerate}
\item $\sigma(A_+)(B[1/g])_0 \subset \Ker(\tau : (B[1/g])_0 \to (A[1/f])_0)$,
\item $\sigma(f) \in (B[1/g])_0$ は非零因子である。
\item $\sigma(f) (B[1/g])_0 = \sigma(A_d) (B[1/g])_0$ が
イデアルとして成り立つ。
\item $\sigma(f) (B[1/g])_0$ と
$\Ker(\tau : (B[1/g])_0 \to (A[1/f])_0)$ は同じ根基をもつ。
\item $d = 1$ なら、
$\sigma(f) (B[1/g])_0 = \Ker(\tau : (B[1/g])_0 \to (A[1/f])_0)$.
\end{enumerate}
特に、
$$
0(D_+(f)) = V(\sigma(f)) \subset D_+(g) = \Spec((B[1/g])_0)
$$
が集合論的に成り立つことがわかる。言い換えると、$\sigma(A_d)$ により
生成されるイデアルは $D_+(g)$ 上の有効 Cartier 因子を切り出し、これは
集合論的には閉埋め込み $0 : P \to L$ の像に等しい。

\medskip\noindent
$L \to X$ は $U$ 上で同型であると主張する。実際、ある $d > 0$ に
対して $f \in A_d$ なら、
$$
\Spec(A_f) \times_X L = \text{Proj}(A_f \otimes_A B) =
\text{Proj}(B_{\sigma(f)})
$$
である。各 $e$ に対し
% P05-EN-CH37-0225
$(B_{\sigma(f)})_e = A_f \otimes_B B_e = A_f \otimes_A A_{\geq e} = A_f$,
であり、最後の等式は単射 $A_{\geq e} \subset A$ により誘導される。
したがって、$T$ の次数を $1$ として
$B_{\sigma(f)} \cong A_f[T]$ である。
$\text{Proj}(A_f[T]) \to \Spec(A_f)$ は同型なので、これで主張が
示された。以後、$U$ を $L$ の対応する開集合と同一視する。

\medskip\noindent
前段落で行った同一視により、制限 $\pi|_U : U \to P$ を考えられる。
上で何度か行ったように、ある $d > 0$ に対して $f \in A_d$ を選び、
$g = \psi(f) \in B_d$ とおく。このとき
$$
U \cap \pi^{-1}(D_+(f)) = U \cap D_+(g)
$$
は、$D_+(g)$ と $(B[1/g])_0$ のスペクトルとの同一視のもとで、
$\sigma(f) \in (B[1/g])_0$ の零点集合の補集合である。これは上の主張
(4) である。したがって $U \cap D_+(g)$ はアフィンであり、
$$
\mathcal{O}_L(U \cap D_+(g)) = (B[1/g])_0[1/\sigma(f)] =
\bigoplus\nolimits_{m \in \mathbf{Z}} (A[1/f])_m
$$
である。ここで最後の等号は、上で得た同一視
$(B[1/g])_0 = \bigoplus_{m \geq 0} (A[1/f])_m$ の自然な拡張である。
先に $\pi : L \to P$ について行ったのとまったく同様に、
$\pi|_U : U \to P$ はアフィンであり、$P$ 上のスキームとして
$$
U = \underline{\Spec}_P
\left(
\bigoplus\nolimits_{m \in \mathbf{Z}} \mathcal{O}_P(m)
\right)
$$
であると結論する。

\medskip\noindent
以上をまとめると、この構成はスキームの可換図式
\begin{equation}
\label{more-morphisms-equation-proj-and-spec}
\vcenter{
\xymatrix{
\underline{\Spec}_P \left(
\bigoplus\nolimits_{m \in \mathbf{Z}} \mathcal{O}_P(m)
\right)
\ar[r] \ar@{=}[d] &
L =
\underline{\Spec}_P \left(
\bigoplus\nolimits_{m \geq 0} \mathcal{O}_P(m)
\right) \ar[d]^\sigma \ar[r]_-\pi & P \ar[d] \\
U \ar[r] & X \ar[r] & Z
}
}
\end{equation}
を与える。ここで $\pi$ は錐であり、その零切断 $0 : P \to L$ は
集合論的に $L$ における $Z$ の逆像の上へ写る。

\medskip\noindent
$W \subset P$ を、$\mathcal{O}_P(1)|_W$ が可逆であり、自然な写像が
すべての $m \in \mathbf{Z}$ に対して同型
$\mathcal{O}_P(m)|_W \cong \mathcal{O}_P(1)^{\otimes m}|_W$ を誘導する
最大の開集合、すなわち Constructions, Lemma
\ref{constructions-lemma-where-invertible} において $d = 1$ とした開集合
とする。このとき $L|_W = \pi^{-1}(W) \to W$ は階数 $1$ のベクトル束
（Constructions, Section \ref{constructions-section-vector-bundle}）、すなわち
$$
L|_W = \mathbf{V}(\mathcal{O}_P(1)|_W)
$$
である。ここでは Grothendieck 流の記法を用いている。これは、上で
$L|_W$ が $\mathcal{O}_P(1)|_W$ 上の $\mathcal{O}_W$-対称代数の
相対スペクトルに等しいことを示したことから直ちに従う。このとき明らかに
射 $0|_W : W \to L|_W$ はこのベクトル束の零切断である。特に
$0(W)$ は $L|_W$ 上の有効 Cartier 因子である。さらに、開集合
$U|_W = (\pi|_U)^{-1}(W)$ は零切断の補集合である。

\medskip\noindent
$A$ が $A_0$ 上 $f_1, \ldots, f_r \in A_1$ により生成されるなら、
すべての $m \geq 0$ に対して
$(f_1, \ldots, f_r)^m = A_{\geq m}$ である。したがって上の $B$ は
$A_+ = (f_1, \ldots, f_r)$ の Rees 代数である。よってこの場合
$L \to X$ は $Z$ の吹き上げであり、前段落の $W$ について $W = P$ である。

\medskip\noindent
$P$ が準コンパクトなら、十分可除な $d$ に対して、
$\sigma(A_d)\mathcal{O}_L$ が切り出す閉部分スキーム $D \subset L$ は
有効 Cartier 因子であり、$0 : P \to L$ は $D$ を経由し、集合論的に
$0(P) = D$ である。これは Constructions, Lemma
\ref{constructions-lemma-proj-quasi-compact} と、上で証明した
(1), (2), (3), (4) から従う。（補題で得られる元の次数の最小公倍数で
割り切れる任意の $d$ をとればよい。）

\medskip\noindent
引き続き $P$ は準コンパクトであると仮定する。$\mathcal{F}$ を準連接
$\mathcal{O}_P$-加群とする。$\mathcal{F}_U = \pi^*\mathcal{F}|_U$
とおく。このとき
\begin{equation}
\label{more-morphisms-equation-cohomology-torsor}
R\Gamma(U, \mathcal{F}_U) =
\bigoplus\nolimits_{m \in \mathbf{Z}}
R\Gamma(P, \mathcal{F} \otimes_{\mathcal{O}_P} \mathcal{O}_P(m))
\end{equation}
が成り立つ。さらに、この直和分解は $\mathcal{F}$ に関して関手的であり、
右辺に誘導される $A$-加群構造は、$U \subset X$ から得られる左辺の
$A$-加群構造と同じである。この公式を証明するため、$\pi|_U$ は
アフィンであり、
$(\pi|_U)_*\mathcal{O}_U = \bigoplus_{m \in \mathbf{Z}} \mathcal{O}_P(m)$
なので、
\begin{align*}
R(\pi|_U)_*\mathcal{F}_U
& =
(\pi|_U)_*\mathcal{F}_U \\
& =
(\pi|_U)_*(\pi|_U)^*\mathcal{F} \\
& =
\mathcal{F} \otimes_{\mathcal{O}_P} 
\bigoplus\nolimits_{m \in \mathbf{Z}} \mathcal{O}_P(m) \\
& =
\bigoplus\nolimits_{m \in \mathbf{Z}}
\mathcal{F} \otimes_{\mathcal{O}_P} \mathcal{O}_P(m)
\end{align*}
を得る。Leray により
$R\Gamma(U, \mathcal{F}_U) =
R\Gamma(P, R(\pi|_U)_*\mathcal{F}_U)$ である。Cohomology, Lemma
\ref{cohomology-lemma-apply-Leray} を参照せよ。この場合コホモロジーを
とる操作は直和と可換なので、証明が完了する。Derived Categories of
Schemes, Lemma
\ref{perfect-lemma-quasi-coherence-pushforward-direct-sums} を参照せよ。
ここで $P$ の準コンパクト性を用いる。なお、Constructions, Lemma
\ref{constructions-lemma-proj-separated} により $P$ は分離的である。

\begin{lemma}
\label{more-morphisms-lemma-apply-proj-spec}
$R$ を環とする。$P$ を $R$ 上の固有スキームとし、$\mathcal{L}$ を
豊富な可逆 $\mathcal{O}_P$-加群とする。
$A = \bigoplus_{m \geq 0} \Gamma(P, \mathcal{L}^{\otimes m})$ とおく。
このとき $P = \text{Proj}(A)$ であり、図式
(\ref{more-morphisms-equation-proj-and-spec}) は図式
$$
\xymatrix{
\underline{\Spec}_P \left(
\bigoplus\nolimits_{m \in \mathbf{Z}} \mathcal{L}^{\otimes m}
\right)
\ar[r] \ar@{=}[d] &
L =
\underline{\Spec}_P \left(
\bigoplus\nolimits_{m \geq 0} \mathcal{L}^{\otimes m}
\right) \ar[d]^\sigma \ar[r]_-\pi & P \ar[d] \\
U \ar[r] & X \ar[r] & Z
}
$$
となり、上で説明した性質をもつ。
\end{lemma}

\begin{proof}
Morphisms, Lemma \ref{morphisms-lemma-proper-ample-is-proj} により
$P = \text{Proj}(A)$ である。さらに Properties, Lemma
\ref{properties-lemma-ample-gcd-is-one} により、この同一視を通じて
すべての $m \in \mathbf{Z}$ に対し
$\mathcal{O}_P(m) = \mathcal{L}^{\otimes m}$ である。
\end{proof}









\section{ファイバーの閉点}
\label{more-morphisms-section-closed-points-fibres}

\noindent
この節の内容の一部はプレプリント \cite{Osserman-Payne} から採られている。

\begin{lemma}
\label{more-morphisms-lemma-locally-principal-vertical}
$f : X \to S$ をスキームの射とする。$Z \subset X$ を閉部分スキームとし、
$s \in S$ とする。次を仮定する。
\begin{enumerate}
\item $S$ は既約で、一般点は $\eta$ である。
\item $X$ は既約である。
\item $f$ は支配的である。
\item $f$ は局所有限型である。
\item $\dim(X_s) \leq \dim(X_\eta)$,
\item $Z$ は $X$ において局所的に主である。
\item $Z_\eta = \emptyset$.
\end{enumerate}
このときファイバー $Z_s$ は、集合論的に $X_s$ の既約成分の和である。
\end{lemma}

\begin{proof}
$X_{red}$ を $X$ の被約化とする。このとき $Z \cap X_{red}$ は
$X_{red}$ の局所的に主な閉部分スキームである。Divisors, Lemma
\ref{divisors-lemma-pullback-locally-principal} を参照せよ。したがって
$X$ は被約であると仮定してよい。言い換えると $X$ は整である。
Properties, Lemma \ref{properties-lemma-characterize-integral} を参照せよ。
この場合、射 $X \to S$ は $S_{red}$ を経由する。Schemes, Lemma
\ref{schemes-lemma-map-into-reduction} を参照せよ。よって $S$ を
$S_{red}$ で置き換え、$S$ も整であると仮定してよい。

\medskip\noindent
$f$ が支配的であるという主張は、$X$ の一般点が $\eta$ に写ることを
意味する。Morphisms, Lemma
\ref{morphisms-lemma-dominant-finite-number-irreducible-components} を
参照せよ。さらに、スキーム $X_\eta$ は体 $\kappa(\eta)$ 上局所有限型な
整スキームである。したがって $d = \dim(X_\eta) \geq 0$ は
$X_\eta$ の任意の点 $\xi$ に対して $\dim_\xi(X_\eta)$ に等しい。
Algebra, Lemmas \ref{algebra-lemma-dimension-spell-it-out} および
\ref{algebra-lemma-dimension-at-a-point-finite-type-over-field} を参照せよ。
Morphisms, Lemma \ref{morphisms-lemma-openness-bounded-dimension-fibres}
と条件 (5) により、すべての $x \in X_s$ に対して
$\dim_x(X_s) = d$ であると結論する。

\medskip\noindent
Noetherian の場合、主張は次のように証明できる。
% P05-EN-CH37-0226
補題が成り立たないなら、$Z_s$ のある既約成分の一般点ではあるが
$X_s$ のどの既約成分の一般点でもない $x \in Z_s$ が存在する。
$\dim_x(X_s) = d$ であり、$X_s$ における $x$ のある近傍では閉部分
スキーム $Z_s$ は $X_s$ のどの既約成分も含まないので、
$\dim_x(Z_s) \leq d - 1$ であることがわかる。したがって $X$ を
$x$ の開近傍で置き換えた後、すべての $z \in Z$ に対して
$\dim_z(Z_{f(z)}) \leq d - 1$ と仮定してよい。Morphisms, Lemma
\ref{morphisms-lemma-openness-bounded-dimension-fibres} を参照せよ。
$\xi' \in Z$ を $Z$ のある既約成分の一般点とし、
% P05-EN-CH37-0227
$s' = f(\xi)$ とおく。$Z \not = X$ は局所的に主なので、
% P05-EN-CH37-0228
$\dim(\mathcal{O}_{X, \xi}) = 1$ であることがわかる。Algebra, Lemma
\ref{algebra-lemma-minimal-over-1} を参照せよ（ここで $X$ が
Noetherian であることを用いる）。$\xi \in X$ を $X$ の一般点とし、
$\xi_1$ を $\xi'$ を含む $X_{s'}$ の任意の既約成分の一般点とする。
このとき特殊化
$$
\xi \leadsto \xi_1 \leadsto \xi'.
$$
がある。$\dim(\mathcal{O}_{X, \xi}) = 1$ なので、二つの特殊化の
一方は等式でなければならない。仮定により $s' \not = \eta$ だから、
最初の特殊化は等式ではない。したがって $\xi' = \xi_1$ は
$X_{s'}$ のある既約成分の一般点である。Morphisms, Lemma
\ref{morphisms-lemma-openness-bounded-dimension-fibres} をもう一度
適用すると、これは
$\dim_{\xi'}(Z_{s'}) = \dim_{\xi'}(X_{s'}) \geq \dim(X_\eta) = d$
を含意し、望む矛盾を与える。

\medskip\noindent
一般の場合は次のように Noetherian の場合へ帰着する。補題が偽なら、
$s$ の上にある点 $x \in X$ であって、$x$ は $Z_s$ のある既約成分の一般点では
あるが $X_s$ のどの既約成分の一般点でもないものが存在する。
$U \subset S$ を $s$ のアフィン近傍とし、$V \subset X$ を
$f(V) \subset U$ を満たす $x$ のアフィン近傍とする。
$U = \Spec(A)$ および $V = \Spec(B)$ と書き、$f|_V$ は環準同型
$A \to B$ により与えられるとする。$x$, resp.\ $s$ に対応する素イデアルを
$\mathfrak q \subset B$, resp.\ $\mathfrak p \subset A$ とする。必要なら
$V$ を縮小して、$Z \cap V$ はある元 $g \in B$ により切り出されると
仮定してよい。$A$ の分数体を $K$ と書く。この時点でわかっていることは
次の通りである。
\begin{enumerate}
\item $A \subset B$ は整域の有限生成拡大である。
\item 元 $g \otimes 1$ は $B \otimes_A K$ において可逆である。
\item $d = \dim(B \otimes_A K) = \dim(B \otimes_A \kappa(\mathfrak p))$,
\item $g \otimes 1$ は $B \otimes_A \kappa(\mathfrak p)$ の単元ではない。
\item $g \otimes 1$ は $B \otimes_A \kappa(\mathfrak p)$ のどの
極小素イデアルにも属さない。
\end{enumerate}
矛盾を導くことが目標である。

\medskip\noindent
$A$ 上 $B$ を生成する元 $x_1, \ldots, x_n \in B$ を選ぶ。有限生成
$\mathbf{Z}$-代数 $A_0 \subset A$ に対し、$B_0 \subset B$ を
$x_1, \ldots, x_n$ により生成される $A_0$-部分代数とし、$A_0$ の
分数体を $K_0$ と書き、$\mathfrak p_0 = A_0 \cap \mathfrak p$ とおく。
$A_0$ が十分大きければ、系 $(A_0 \subset B_0, g, \mathfrak p_0)$ に
対しても (1) -- (5) が成り立つと主張する。

\medskip\noindent
各条件を順に証明する。(1) は構成から成り立つ。(2) について、
% P05-EN-CH37-0229
ある $h \otimes 1/a \in B \otimes_A K$ に対して
$(g \otimes 1) h = 1$ と書く。ある係数 $a_I, a'_I \in A$ を用いて
$g = \sum a_I x^I$, $h = \sum a'_I x^I$ と書く（多重指数記法）。
$A_0$ が $a$ およびすべての $a_I, a'_I$ を含めば、(2) が成り立つ。
実際、$B_0 \to B$ という単射の局所化として
$B_0 \otimes_{A_0} K_0 \subset B \otimes_A K$ である。(3) を得るため、
完全列
$$
0 \to I \to A[X_1, \ldots, X_n] \to B \to 0
$$
を考える。これは $I$ を定義し、第2の写像は $X_i$ を $x_i$ に送る。
$\otimes$ は右完全なので、$I \otimes_A K$, resp.\
$I \otimes_A \kappa(\mathfrak p)$ は全射
$K[X_1, \ldots, X_n] \to B \otimes_A K$, resp.\
$\kappa(\mathfrak p)[X_1, \ldots, X_n] \to
B \otimes_A \kappa(\mathfrak p)$ の核である。体上の多項式環は
Noetherian なので、$I \otimes_A K$ と
$I \otimes_A \kappa(\mathfrak p)$ を生成する有限個の元
$h_j \in I$, $j = 1, \ldots, m$ が存在する。
$h_j = \sum a_{j, I}X^I$ と書く。$A_0$ がすべての $a_{j, I}$ を
含めば、
$$
B_0 \otimes_{A_0} K_0 \otimes_{K_0} K = B \otimes_A K
\quad\text{and}\quad
B_0 \otimes_{A_0} \kappa(\mathfrak p_0)
\otimes_{\kappa(\mathfrak p_0)} \kappa(\mathfrak p)
=
B \otimes_A \kappa(\mathfrak p).
$$
となる。Morphisms, Lemma
\ref{morphisms-lemma-dimension-fibre-after-base-change} または
Algebra, Lemma \ref{algebra-lemma-dimension-preserved-field-extension}
により、(3) の次元の等式が満たされることがわかる。(4) は直ちに従う。
$B_0 \otimes_{A_0} \kappa(\mathfrak p_0) \subset
B \otimes_A \kappa(\mathfrak p)$ なので、Algebra, Lemma
\ref{algebra-lemma-image-dense-generic-points} により
$B_0 \otimes_{A_0} \kappa(\mathfrak p_0)$ の各極小素イデアルは
$B \otimes_A \kappa(\mathfrak p)$ のある極小素イデアルの下にある。
これは (5) を含意する。このようにして、問題は上で扱った Noetherian の
場合へ帰着される。
\end{proof}

\noindent
上の補題の代数的応用を与える。$A \to B$ が平坦なら補題の第4の仮定は
成り立つ。Lemma \ref{more-morphisms-lemma-equality-dimensions} を参照せよ。

\begin{lemma}
\label{more-morphisms-lemma-horizontal}
$A \to B$ を局所環の局所準同型とし、$g \in \mathfrak m_B$ とする。
次を仮定する。
\begin{enumerate}
\item $A$ と $B$ は整域であり、$A \subset B$ である。
\item $B$ は $A$ 上本質的有限型である。
\item $g$ は $\mathfrak m_AB$ の上のどの極小素イデアルにも属さない。
\item $\dim(B/\mathfrak m_AB) +
\text{trdeg}_{\kappa(\mathfrak m_A)}(\kappa(\mathfrak m_B)) =
\text{trdeg}_A(B)$.
\end{enumerate}
このとき $A \subset B/gB$ である。すなわち、$\Spec(A)$ の一般点は
射 $\Spec(B/gB) \to \Spec(A)$ の像に属する。
\end{lemma}

\begin{proof}
Algebra, Lemma \ref{algebra-lemma-image-dense-generic-points} により、
二つの主張は同値であることに注意する。証明を始めるため、$C$ を有限型
$A$-代数とし、$\mathfrak q$ を $B = C_{\mathfrak q}$ となる $C$ の
素イデアルとする。もちろん $C$ は整域であり $g \in C$ であると
仮定してよい。$C$ を局所化で置き換えると、
$\dim(C/\mathfrak m_AC) = \dim(B/\mathfrak m_AB) +
\text{trdeg}_{\kappa(\mathfrak m_A)}(\kappa(\mathfrak m_B))$ となることが
わかる。Morphisms, Lemma
\ref{morphisms-lemma-dimension-fibre-at-a-point} を参照せよ。$K$ を
$A$ の分数体とすると、同じ参照により
$\dim(C \otimes_A K) = \text{trdeg}_A(B)$ であることがわかる。
したがって仮定 (4) は、射 $\Spec(C) \to \Spec(A)$ の一般ファイバーと
閉ファイバーが同じ次元をもつことを意味する。

\medskip\noindent
補題が偽であると仮定する。このとき
$(B/gB) \otimes_A K = 0$ である。これは $g \otimes 1$ が
$B \otimes_A K = C_{\mathfrak q} \otimes_A K$ において可逆であることを
意味する。
% P05-EN-CH37-0230
$C_{\mathfrak q}$ は主局所化の極限なので、ある
$h \in C$, $h \not \in \mathfrak q$ に対して $g \otimes 1$ は
$C_h \otimes_A K$ において可逆であると結論する。したがって $C$ を
$C_h$ で置き換えた後、$(C/gC) \otimes_A K = 0$ と仮定してよい。
$C/\mathfrak m_AC$ の極小素イデアルが $B/\mathfrak m_AB$ の極小素
イデアルと一対一に対応するように、$C$ をもう一度置き換える。この時点で
Lemma \ref{more-morphisms-lemma-locally-principal-vertical} を
$X = \Spec(C) \to \Spec(A) = S$ および局所閉部分スキーム
% P05-EN-CH37-0231
$Z = \Spec(C/gC)$ に適用する。$Z_K = \emptyset$ なので、
$Z \otimes \kappa(\mathfrak m_A)$ は
$X \otimes \kappa(\mathfrak m_A) = \Spec(C/\mathfrak m_AC)$ の
ある既約成分を含まなければならないことがわかる。しかしこれは、$g$ が
$\mathfrak m_AB$ の上で極小な任意の素イデアルに属さないという仮定に
矛盾する。補題が従う。
\end{proof}

\begin{lemma}
\label{more-morphisms-lemma-equality-dimensions}
$A \to B$ を局所環の局所準同型とする。次を仮定する。
\begin{enumerate}
\item $A$ と $B$ は整域であり、$A \subset B$ である。
\item $B$ は $A$ 上本質的有限型である。
\item $B$ は $A$ 上平坦である。
\end{enumerate}
このとき
$$
\dim(B/\mathfrak m_AB) +
\text{trdeg}_{\kappa(\mathfrak m_A)}(\kappa(\mathfrak m_B)) =
\text{trdeg}_A(B).
$$
が成り立つ。
\end{lemma}

\begin{proof}
$C$ を有限型 $A$-代数とし、$\mathfrak q$ を $B = C_{\mathfrak q}$ と
なる $C$ の素イデアルとする。$C$ は整域であると仮定してよい。
Morphisms, Lemma \ref{morphisms-lemma-dimension-fibre-at-a-point} により
$\dim_{\mathfrak q}(C/\mathfrak m_AC) = \dim(B/\mathfrak m_AB) +
\text{trdeg}_{\kappa(\mathfrak m_A)}(\kappa(\mathfrak m_B))$ である。
$K$ を $A$ の分数体とすると、同じ参照により
$\dim(C \otimes_A K) = \text{trdeg}_A(B)$ であることがわかる。
したがって、実際に証明しようとしているのは
$\dim_{\mathfrak q}(C/\mathfrak m_AC) = \dim(C \otimes_A K)$ である。
$A$ を支配する $K$ の付値環 $A'$ を選ぶ。Algebra, Lemma
\ref{algebra-lemma-dominate} を参照せよ。$C' = C \otimes_A A'$ とおく。
$\mathfrak q$ の上にある $C'$ の素イデアル $\mathfrak q'$ を選ぶ。
そのような素イデアルは存在する。実際、
$$
C'/\mathfrak m_{A'}C' =
C/\mathfrak m_AC \otimes_{\kappa(\mathfrak m_A)} \kappa(\mathfrak m_{A'})
$$
であり、これは $C/\mathfrak m_AC \to C'/\mathfrak m_{A'}C'$ が忠実平坦
であることを示す。また、
$\dim_{\mathfrak q}(C/\mathfrak m_AC) =
\dim_{\mathfrak q'}(C'/\mathfrak m_{A'}C')$ も示す。Algebra, Lemma
\ref{algebra-lemma-dimension-at-a-point-preserved-field-extension} を
参照せよ。$B' = C'_{\mathfrak q'}$ は $B \otimes_A A'$ の局所化である
ことに注意する。したがって $B'$ は $A'$ 上平坦である。一般ファイバー
$B' \otimes_{A'} K$ は $B \otimes_A K$ の局所化である。よって $B'$ は
整域である。$A' \subset B'$ に対して補題を証明すれば、等式
$\dim_{\mathfrak q'}(C'/\mathfrak m_{A'}C') = \dim(C' \otimes_{A'} K)$
を得る。これは上で述べたことにより、望む等式
$\dim_{\mathfrak q}(C/\mathfrak m_AC) = \dim(C \otimes_A K)$
を含意する。これにより、補題は $A$ が付値環である場合へ帰着される。

\medskip\noindent
$A$ を付値環とし、$A \subset B$ は補題の通りであるとする。前と同様に、
$A$ 上有限型なある整域 $C$ に対して $B = C_{\mathfrak q}$ と書く。
Algebra, Lemma
\ref{algebra-lemma-finite-type-domain-over-valuation-ring-dim-fibres}
により $\dim(C/\mathfrak m_AC) = \dim(C \otimes_A K)$ を得て、証明が
完了する。
\end{proof}

\begin{lemma}
\label{more-morphisms-lemma-closed-point-nearby-fibre}
$f : X \to S$ をスキームの射とする。$x \leadsto x'$ を $X$ の点の
特殊化とする。$s = f(x)$ および $s' = f(x')$ とおく。次を仮定する。
\begin{enumerate}
\item $x'$ は $X_{s'}$ の閉点である。
\item $f$ は局所有限型である。
\end{enumerate}
このとき集合
$$
\{x_1 \in X
\text{ であって }
f(x_1) = s
\text{ かつ }
x_1\text{ は }X_s\text{ で閉であり、}
x \leadsto x_1 \leadsto x'
\}
$$
は $X_s$ における $x$ の閉包で稠密である。
\end{lemma}

\begin{proof}
特殊化 $x \leadsto x'$ に Schemes, Lemma
\ref{schemes-lemma-points-specialize} を適用する。これにより射
$\varphi : \Spec(B) \to X$ が得られる。ここで $B$ は付値環であり、
$\varphi$ は一般点を $x$ に、閉点を $x'$ に写す。また、$\kappa(x)$ は
$B$ の分数体であると仮定してよい。$A = B \cap \kappa(s)$ とおく。
これは付値環であり（Algebra, Lemma
\ref{algebra-lemma-valuation-ring-cap-field} を参照せよ）、
$\mathcal{O}_{S, s'} \to \kappa(s)$ の像を支配することに注意する。
可換図式
$$
\xymatrix{
\Spec(B) \ar[rd] \ar[r] &
X_A \ar[d] \ar[r] & X \ar[d] \\
& \Spec(A) \ar[r] & S
}
$$
を考える。$B$ の一般点（resp.\ 閉点）は、$\Spec(A)$ の一般点
（resp.\ 閉点）の上にある $X_A$ の点 $x_A$（resp.\ $x'_A$）へ写る。
Morphisms, Lemma
\ref{morphisms-lemma-base-change-closed-point-fibre-locally-finite-type}
により、$x'_A$ は $X_A$ の特殊ファイバーの閉点であることに注意する。
$X_A \to \Spec(A)$ の一般ファイバーは $X_s$ と同型であることにも
注意する。したがって補題は、$S$ が付値環のスペクトルで、
$s = \eta \in S$ が一般点、$s' \in S$ が閉点である場合へ帰着された。

\medskip\noindent
$\dim_x(X_\eta)$ に関する帰納法で補題を証明する。
$\dim_x(X_\eta) = 0$ なら、$x$ に特殊化する $X_\eta$ のほかの点はなく、
$x$ はそのファイバーで閉である。Morphisms, Lemma
\ref{morphisms-lemma-quasi-finite-at-point-characterize} を参照せよ。
したがって結果は成り立つ。$\dim_x(X_\eta) > 0$ と仮定する。

\medskip\noindent
$X' \subset X$ を、$X$ の既約閉部分スキーム $\overline{\{x\}}$ 上の
被約誘導スキーム構造とする。Schemes, Definition
\ref{schemes-definition-reduced-induced-scheme} を参照せよ。これは
$\dim_x(X_\eta)$ を減少させるだけなので、補題を証明する際に $X$ を
$X'$ で置き換えてよい。したがって $X$ は整スキームで、$x$ はその一般点
であるとも仮定してよい。さらに、$X$ を $x'$ のアフィン近傍で置き換えて
よい。よって $X = \Spec(B)$ であり、$A \subset B$ は整域の有限型拡大
である。この場合
$\dim_x(X_\eta) = \dim(X_\eta) = \dim(X_{s'})$ であり、実際
$X_{s'}$ は等次元であることに注意する。Algebra, Lemma
\ref{algebra-lemma-finite-type-domain-over-valuation-ring-dim-fibres}
を参照せよ。

\medskip\noindent
$W \subset X_\eta$ を真の閉部分集合（これが「避け」たい部分集合である）
とする。$X_s$ は体上有限型なので、$W$ は有限個の既約成分
$W = W_1 \cup \ldots \cup W_n$ をもつことがわかる。対応する素イデアルを
% P05-EN-CH37-0232
$\mathfrak q_j \subset B$, $j = 1, \ldots, r$ とする。点 $x'$ に対応する
極大イデアルを $\mathfrak q \subset B$ とする。
$\mathfrak p_1, \ldots, \mathfrak p_s \subset B$ を
$\mathfrak m_AB$ の上にある極小素イデアルとする。これらは Noetherian
スキーム $X_{s'}$ の既約成分に対応するので有限個である。さらに、これらの
既約成分はそれぞれ次元 $> 0$ をもつ（上を参照）ので、すべての $i$ に
対して $\mathfrak p_i \not = \mathfrak q$ である。すべての $j$ に対して
$g \not \in \mathfrak q_j$ かつすべての $i$ に対して
$g \not \in \mathfrak p_i$ となる元 $g \in \mathfrak q$ を選ぶ。
Algebra, Lemma \ref{algebra-lemma-silly} を参照せよ。$g$ が定める局所的に
主な閉部分スキームを $Z \subset X$ と書く。
$Z_\eta = Z_{1, \eta} \cup \ldots \cup Z_{n, \eta}$, $n \geq 0$ を
$Z$ の一般ファイバーの既約成分への分解とする（一般ファイバーは
Noetherian なので有限個である）。$Z_i \subset X$ を $Z_{i, \eta}$ の
閉包とする。$X$ をより小さいアフィン近傍で置き換え、各
$i = 1, \ldots, n$ に対して $x' \in Z_i$ と仮定してよい。構成により
$Z \cap X_{s'}$ は $X_{s'}$ のどの既約成分も含まない。したがって
Lemma \ref{more-morphisms-lemma-locally-principal-vertical} により
$Z_\eta \not = \emptyset$ と結論する。すなわち $n \geq 1$ である。
$x_1 \in Z_1$ を一般点とすると、$x_1 \leadsto x'$ かつ
$f(x_1) = \eta$ であることがわかる。また構成により
$Z_{1, \eta} \cap W_j \subset W_j$ は真の閉部分集合である。したがって
$Z_{1, \eta} \cap W_j$ の各既約成分は $X_\eta$ において余次元
$\geq 2$ をもつ一方、Algebra, Lemma
\ref{algebra-lemma-minimal-over-1} により
$\text{codim}(Z_{1, \eta}, X_\eta) = 1$ である。よって
$W \cap Z_{1, \eta}$ は真の閉部分集合である。この時点で、帰納法の仮定を
$Z_1 \to S$ と特殊化 $x_1 \leadsto x'$ に適用できることがわかる。
これにより、$W$ に含まれず $x'$ に特殊化する $Z_{1, \eta}$ の閉点
$x_2$ が得られる。したがって望み通り $x \leadsto x_2 \leadsto x'$ を
得て、点 $x_2$ は $X_\eta$ で閉であり、$x_2 \not \in W$ である。
\end{proof}

\begin{remark}
\label{more-morphisms-remark-full-specialization-sequence}
Lemma \ref{more-morphisms-lemma-closed-point-nearby-fibre} の証明は、実際には特殊化の列
$$
x \leadsto x_1 \leadsto x_2 \leadsto \ldots \leadsto x_d \leadsto x'
$$
が存在することを示している。ここで、すべての $x_i$ はファイバー
$X_s$ に属し、各特殊化は直近の特殊化であり、$x_d$ は $X_s$ の閉点
である。整数
$d = \text{trdeg}_{\kappa(s)}(\kappa(x)) = \dim(\overline{\{x\}})$
であり、ここで閉包は $X_s$ においてとる。さらに、点 $x$ を含まない
$X_s$ の任意の閉部分集合を避けるように点 $x_i$ を選べる。
\end{remark}

\noindent
Examples, Section \ref{examples-section-topology-finite-type} は、$A$ を
Noetherian と仮定しなければ次の補題が偽であることを示す。

\begin{lemma}
\label{more-morphisms-lemma-quasi-finite-quasi-section-meeting-nearby-open}
$\varphi : A \to B$ を局所環の局所環準同型とする。
$V \subset \Spec(B)$ を、$\mathfrak m_A$ の上にない素イデアルを少なくとも
一つ含む開部分スキームとする。$A$ は Noetherian、$\varphi$ は本質的
有限型であり、$A/\mathfrak m_A \subset B/\mathfrak m_B$ は有限であると
仮定する。このとき $\mathfrak q \in V$ で
$\mathfrak m_A \not = \mathfrak q \cap A$ を満たし、
$A \to B/\mathfrak q$ が準有限環準同型の局所化となるものが存在する。
\end{lemma}

\begin{proof}
$A$ は Noetherian で $A \to B$ は本質的有限型なので、$B$ も
Noetherian である。Properties, Lemma
\ref{properties-lemma-complement-closed-point-Jacobson} により、位相空間
$\Spec(B) \setminus \{\mathfrak m_B\}$ は Jacobson である。したがって
非空な開集合
$$
V \setminus \{\mathfrak q \subset B \mid \mathfrak m_A = \mathfrak q \cap A\}.
$$
に属する閉点 $\mathfrak q$ を選べる。（仮定により非空であり、
$\{\mathfrak m_A\}$ は $\Spec(A)$ の閉部分集合なので開である。）
このとき $\Spec(B/\mathfrak q)$ は $\mathfrak m_B$ と $\mathfrak q$ という
二つの点をもち、$\mathfrak q$ は $\mathfrak m_A$ の上にない。ある有限型
$A$-代数 $C$ と素イデアル $\mathfrak m$ に対して
$B/\mathfrak q = C_{\mathfrak m}$ と書く。このとき Algebra, Lemma
\ref{algebra-lemma-isolated-point-fibre} (2) により $A \to C$ は
$\mathfrak m$ において準有限である。したがって Algebra, Lemma
\ref{algebra-lemma-quasi-finite-open} により、$C$ を主局所化で
置き換えた後、環準同型 $A \to C$ は準有限であることがわかる。
\end{proof}

\begin{lemma}
\label{more-morphisms-lemma-quasi-finite-quasi-section-meeting-nearby-open-X}
$f : X \to S$ をスキームの射とする。$x \in X$ の像を $s \in S$ とする。
$U \subset X$ を開部分スキームとする。$f$ は局所有限型、$S$ は局所
Noetherian、$x$ は $X_s$ の閉点であり、$x' \leadsto x$ かつ
$f(x') \not = s$ を満たす点 $x' \in U$ が存在すると仮定する。このとき
閉部分スキーム $Z \subset X$ であって、(a) $x \in Z$、
(b) $f|_Z : Z \to S$ は $x$ において準有限、(c) $z \leadsto x$ かつ
$f(z) \not = s$ を満たす $z \in Z$, $z \in U$ が存在する、という条件を
満たすものが存在する。
\end{lemma}

\begin{proof}
これは Lemma \ref{more-morphisms-lemma-quasi-finite-quasi-section-meeting-nearby-open}
の言い換えである。実際、$A = \mathcal{O}_{S, s}$ および
$B = \mathcal{O}_{X, x}$ とおく。$V \subset \Spec(B)$ を $U$ の逆像と
書く。環準同型 $f^\sharp : A \to B$ は本質的有限型である。仮定により、
$\Spec(A)$ の閉点に写らない $V$ の点が少なくとも一つ存在する。
したがって Lemma
\ref{more-morphisms-lemma-quasi-finite-quasi-section-meeting-nearby-open} のすべての仮定が
成り立ち、$\mathfrak m_A$ の上になく、$A \to B/\mathfrak q$ が準有限環
準同型の局所化となる素イデアル $\mathfrak q \subset B$ を得る。点
$\mathfrak q$ の標準射 $\Spec(B) \to X$ による像を $z \in X$ とする。
$Z = \overline{\{z\}}$ に被約誘導スキーム構造を入れる。
$z \leadsto x$ なので、$x \in Z$ かつ
$\mathcal{O}_{Z, x} = B/\mathfrak q$ であることがわかる。構成により
$Z \to S$ は $x$ において準有限である。
\end{proof}

\begin{remark}
\label{more-morphisms-remark-alternative-closed-point-nearby-fibre}
$S$ が局所 Noetherian である場合、Lemma
\ref{more-morphisms-lemma-quasi-finite-quasi-section-meeting-nearby-open} またはその変形
Lemma \ref{more-morphisms-lemma-quasi-finite-quasi-section-meeting-nearby-open-X} を用いて、
Lemma \ref{more-morphisms-lemma-closed-point-nearby-fibre} の別証明を与えられる。
概略を述べる。まず $S$ を $s'$ における局所環のスペクトルで置き換える。
このとき $\dim(S)$ に関する帰納法を用いられる。
$\dim(S) = 0$ の場合は $s' = s$ なので自明である。$X$ を
$\overline{\{x\}}$ 上の被約誘導スキーム構造で置き換える。
Lemma \ref{more-morphisms-lemma-quasi-finite-quasi-section-meeting-nearby-open-X} を
$X \to S$、$x' \mapsto s'$、および $x$ を含む任意の非空開集合
$U \subset X$ に適用する。これにより、
% P05-EN-CH37-0233
閉部分スキーム $x' \in Z \subset X$ と点 $z \in Z$ であって、
$Z \to S$ は $x'$ において準有限であり、$f(z) \not = s'$ となるものを
得る。このとき $z$ は $X_{f(z)}$ の閉点で、$z \leadsto x'$ である。
$f(z) \not = s'$ なので
$\dim(\mathcal{O}_{S, f(z)}) < \dim(S)$ であることがわかる。$x$ は
$X$ の一般点なので $x \leadsto z$、したがって
$s = f(x) \leadsto f(z)$ である。
% P05-EN-CH37-0234
帰納法の仮定を $s \leadsto f(z)$ および $z \mapsto f(z)$ に適用して
証明を終える。
\end{remark}

\begin{lemma}
\label{more-morphisms-lemma-change-hypotheses}
$f : X \to S$ は局所有限型、$S$ は局所 Noetherian、$x \in X$ はその
ファイバー $X_s$ の閉点であり、$U \subset X$ は
$U \cap X_s = \emptyset$ かつ $x \in \overline{U}$ を満たす開部分スキーム
であると仮定する。このとき Lemma
\ref{more-morphisms-lemma-quasi-finite-quasi-section-meeting-nearby-open-X} の結論が
成り立つ。
\end{lemma}

\begin{proof}
次のように引用した補題へ帰着できる。まず $X$ と $S$ をそれぞれ $x$ と
$s$ のアフィン近傍で置き換える。このとき $X$ は Noetherian であり、
特に $U$ は準コンパクトである（Morphisms, Lemma
\ref{morphisms-lemma-finite-type-noetherian} および Topology, Lemmas
\ref{topology-lemma-Noetherian},
\ref{topology-lemma-Noetherian-quasi-compact} を参照せよ）。したがって
$x' \in U$ を満たす特殊化 $x' \leadsto x$ が存在する（Morphisms, Lemma
\ref{morphisms-lemma-reach-points-scheme-theoretic-image} を参照せよ）。
$f(x') \not = s$ であることに注意する。よって補題のすべての仮定が
満たされ、証明が完了する。
\end{proof}



\section{Stein 分解}
\label{more-morphisms-section-stein-factorization}

\noindent
Stein 分解とは、$f_*\mathcal{O}_X = \mathcal{O}_S$ を満たす固有射
$f : X \to S$ のファイバーが連結であるという主張である。

\begin{lemma}
\label{more-morphisms-lemma-stein-universally-closed}
$S$ をスキームとする。$f : X \to S$ を普遍閉かつ準分離的な射とする。
このとき、次の分解
$$
\xymatrix{
X \ar[rr]_{f'} \ar[rd]_f & & S' \ar[dl]^\pi \\
& S &
}
$$
であって、次の性質をもつものが存在する。
\begin{enumerate}
\item 射 $f'$ は普遍閉、準コンパクト、準分離的、かつ全射である。
\item 射 $\pi : S' \to S$ は整である。
\item $f'_*\mathcal{O}_X = \mathcal{O}_{S'}$ が成り立つ。
\item $S' = \underline{\Spec}_S(f_*\mathcal{O}_X)$ が成り立つ。
\item $S'$ は $X$ における $S$ の正規化である。Morphisms, Definition
\ref{morphisms-definition-normalization-X-in-Y} を参照せよ。
\end{enumerate}
分解 $f = \pi \circ f'$ の形成は平坦な基底変換と可換である。
\end{lemma}

\begin{proof}
Morphisms, Lemma \ref{morphisms-lemma-universally-closed-quasi-compact}
により、射 $f$ は準コンパクトである。したがって、$X$ における $S$ の
正規化 $S'$ は定義され（Morphisms, Definition
\ref{morphisms-definition-normalization-X-in-Y}）
、分解 $X \to S' \to S$ を得る。
Morphisms, Lemma \ref{morphisms-lemma-normalization-in-universally-closed}
により (2)、(4)、(5) が成り立つ。Morphisms, Lemma
\ref{morphisms-lemma-image-proper-scheme-closed} により、射 $f'$ は普遍閉である。
Schemes, Lemma \ref{schemes-lemma-quasi-compact-permanence}
によりこれは準コンパクトであり、
Schemes, Lemma \ref{schemes-lemma-compose-after-separated}
により準分離的である。

\medskip\noindent
残りの主張を示すには、基底スキーム $S$ がアフィン、すなわち
$S = \Spec(R)$ であると仮定してよい。このとき
$S' = \Spec(A)$ であり、$A = \Gamma(X, \mathcal{O}_X)$ は整な
$R$-代数である。したがって、$f'_*\mathcal{O}_X$ が
$\mathcal{O}_{S'}$ であることは明らかである（実際、
Schemes, Lemma
\ref{schemes-lemma-push-forward-quasi-coherent}
により $f'_*\mathcal{O}_X$ は準連接であり、それゆえ
$\widetilde{A}$ に等しい）。これで (3) が示された。

\medskip\noindent
$f'$ が全射であることを示そう。$f'$ は普遍閉であるから（上記参照）、
$f'$ の像は閉部分集合 $V(I) \subset S' = \Spec(A)$ である。
$h \in I$ をとる。このとき
% P05-EN-CH37-0235
$h|_X = f^\sharp(h)$ は $X$ の構造層の
大域切断であり、すべての点で消える。$X$ は準コンパクトなので、
$h|_X$ は冪零な切断、すなわちある $n > 0$ に対して
% P05-EN-CH37-0236
$h^n|X = 0$ である。しかし $A = \Gamma(X, \mathcal{O}_X)$ であるから
$h^n = 0$ である。$I$ のすべての元が冪零なので、望ましい等式
$V(I) = S'$ を得る。

\medskip\noindent
Cohomology of Schemes, Lemma \ref{coherent-lemma-flat-base-change-cohomology}
により、$f_*\mathcal{O}_X$ の形成は平坦な基底変換と可換である。
Constructions, Lemma \ref{constructions-lemma-spec-properties}
により、相対スペクトルの形成は任意の基底変換と可換である。
したがって、この分解の形成は平坦な基底変換と可換である。
\end{proof}

\begin{lemma}
\label{more-morphisms-lemma-stein-universally-closed-residue-fields}
Lemma \ref{more-morphisms-lemma-stein-universally-closed} において、さらに
$f$ が局所有限型であると仮定する。このとき $s \in S$ に対し、ファイバー
$\pi^{-1}(\{s\}) = \{s_1, \ldots, s_n\}$ は有限であり、体拡大
$\kappa(s_i)/\kappa(s)$ は有限である。
\end{lemma}

\begin{proof}
$\pi^{-1}(\{s\})$ の点どうしには特殊化が存在しないことを思い出そう。
Algebra, Lemma \ref{algebra-lemma-integral-no-inclusion} を参照せよ。
$f'$ は全射なので、$|X_s| \to \pi^{-1}(\{s\})$ は全射である。
$X_s$ は体上有限型の準分離スキームであることに注意する
（準コンパクト性は参照した補題の証明で示した）。したがって $X_s$ は
Noetherian である（Morphisms, Lemma
\ref{morphisms-lemma-finite-type-noetherian}）。ここで位相的な議論
（省略）により、$\pi^{-1}(\{s\})$ が有限であることが分かる。
各 $i$ に対し、$s_i$ に写る有限型の点 $x_i \in X_s$ をとることができる
（Morphisms, Lemma
\ref{morphisms-lemma-enough-finite-type-points}）。
$\kappa(s_i)/\kappa(s)$ が有限であることを示そう。有限型の点の定義により、
$x_i$ は有限型射 $\Spec(k_i) \to X_s$ で表せる。したがって
$\Spec(k_i) \to s = \Spec(\kappa(s))$ は有限型射の合成として有限型であり、
ゆえに $k_i/\kappa(s)$ は有限である（Morphisms, Lemma
\ref{morphisms-lemma-point-finite-type}）。
\end{proof}

\begin{lemma}
\label{more-morphisms-lemma-characterize-geometrically-connected-fibres}
$f : X \to S$ をスキームの射とし、$s \in S$ とする。このとき
$X_s$ が幾何学的に連結であるための必要十分条件は、任意のエタール近傍
$(U, u) \to (S, s)$ に対し、基底変換 $X_U \to U$ のファイバー
$X_u$ が連結であることである。
\end{lemma}

\begin{proof}
$X_s$ が幾何学的に連結なら、その任意の基底変換は連結である。
逆に、$X_s$ が幾何学的に連結でないと仮定する。このとき
Varieties, Lemma
\ref{varieties-lemma-characterize-geometrically-disconnected}
により、ある有限分離体拡大 $k/\kappa(s)$ に対して
$X_s \times_{\Spec(\kappa(s))} \Spec(k)$ は非連結である。
Lemma \ref{more-morphisms-lemma-realize-prescribed-residue-field-extension-etale}
により、$\kappa(u)/\kappa(s)$ が $k/\kappa(s)$ と同一視されるような
アフィン・エタール近傍 $(U, u) \to (S, s)$ が存在する。この場合
$X_u$ は非連結である。
\end{proof}

\begin{theorem}[Stein 分解；Noetherian の場合]
\label{more-morphisms-theorem-stein-factorization-Noetherian}
$S$ を局所 Noetherian スキームとし、$f : X \to S$ を固有射とする。
このとき、次の分解
$$
\xymatrix{
X \ar[rr]_{f'} \ar[rd]_f & & S' \ar[dl]^\pi \\
& S &
}
$$
であって、次の性質をもつものが存在する。
\begin{enumerate}
\item 射 $f'$ は固有で、そのファイバーは幾何学的に連結である。
\item 射 $\pi : S' \to S$ は有限である。
\item $f'_*\mathcal{O}_X = \mathcal{O}_{S'}$ が成り立つ。
\item $S' = \underline{\Spec}_S(f_*\mathcal{O}_X)$ が成り立つ。
\item $S'$ は $X$ における $S$ の正規化である。Morphisms, Definition
\ref{morphisms-definition-normalization-X-in-Y} を参照せよ。
\end{enumerate}
\end{theorem}

\begin{proof}
Lemma \ref{more-morphisms-lemma-stein-universally-closed} の分解を
$f = \pi \circ f'$ とする。Lemma
\ref{more-morphisms-lemma-stein-universally-closed} の結論に加えて、$f'$ は分離的
（Schemes, Lemma \ref{schemes-lemma-compose-after-separated}）
かつ有限型（Morphisms, Lemma
\ref{morphisms-lemma-permanence-finite-type}）であることに注意する。
したがって $f'$ は固有である。
Cohomology of Schemes, Proposition
\ref{coherent-proposition-proper-pushforward-coherent}
により、$f_*\mathcal{O}_X$ は連接 $\mathcal{O}_S$-加群である。
したがって $\pi$ は有限であり、すなわち (2) が成り立つ。

\medskip\noindent
これで、最も興味深い主張、すなわち $f'$ のすべてのファイバーが
幾何学的に連結であることを除くすべてが示された。上の議論から
$S$ を $S'$ で置き換えてよいことは明らかである。したがって、
$S$ は Noetherian かつアフィン、$f : X \to S$ は固有、さらに
$f_*\mathcal{O}_X = \mathcal{O}_S$ であると仮定してよい。
$s \in S$ を $S$ の点とする。$X_s$ が幾何学的に連結であることを
示さなければならない。Lemma
\ref{more-morphisms-lemma-characterize-geometrically-connected-fibres}
により、任意のエタール近傍 $(U, u) \to (S, s)$ に対して
$X_u$ が連結であることを示せば十分である。$U$ はアフィンと仮定してよい。
すると $U$ は Noetherian であり（Morphisms, Lemma
\ref{morphisms-lemma-finite-type-noetherian}）、基底変換
$f_U : X_U \to U$ は固有であり（Morphisms, Lemma
\ref{morphisms-lemma-base-change-proper}）、さらに
$(f_U)_*\mathcal{O}_{X_U} = \mathcal{O}_U$ も成り立つ
（Cohomology of Schemes, Lemma
\ref{coherent-lemma-flat-base-change-cohomology}）。
したがって $(f : X \to S, s)$ を基底変換
$(f_U : X_U \to U, u)$ で置き換えれば、
$f_*\mathcal{O}_X = \mathcal{O}_S$ のときにファイバー $X_s$ が
連結であることを証明すれば十分である。これは Derived Categories of
Schemes, Lemma
\ref{perfect-lemma-proper-idempotent-on-fibre}
から（$X_s$ の構造層の冪等元を調べることにより）従うが、
以下では直接の議論も与える。

\medskip\noindent
すなわち、形式関数定理、より正確には Cohomology of Schemes, Lemma
\ref{coherent-lemma-formal-functions-stalk}
を適用する。これにより
$$
\mathcal{O}^\wedge_{S, s} = (f_*\mathcal{O}_X)_s^\wedge =
\lim_n H^0(X_n, \mathcal{O}_{X_n})
$$
を得る。ここで $X_n$ は $X_s$ の第 $n$ 無限小近傍である。
$X_n$ の台となる位相空間は $X_s$ のものと等しい。したがって、
$X_s = T_1 \amalg T_2$ が空でない開かつ閉な部分スキームの非交和ならば、
同様にすべての $n$ に対して
$X_n = T_{1, n} \amalg T_{2, n}$ となる。このことは
$H^0(X_n, \mathcal{O}_{X_n})$ が非自明な冪等元 $e_{1, n}$、
すなわち $T_{1, n}$ 上で恒等的に $1$、$T_{2, n}$ 上で恒等的に $0$
である関数を含むことを意味する。$e_{1, n + 1}$ の $X_n$ への制限が
$e_{1, n}$ であることは明らかである。したがって
$e_1 = \lim e_{1, n}$ はこの極限の非自明な冪等元である。
これは $\mathcal{O}^\wedge_{S, s}$ が局所環であることに矛盾する。
ゆえに仮定は誤りであり、$X_s$ は連結である。
\end{proof}

\begin{theorem}[Stein 分解；一般の場合]
\label{more-morphisms-theorem-stein-factorization-general}
$S$ をスキームとし、$f : X \to S$ を固有射とする。
このとき、次の分解
$$
\xymatrix{
X \ar[rr]_{f'} \ar[rd]_f & & S' \ar[dl]^\pi \\
& S &
}
$$
であって、次の性質をもつものが存在する。
\begin{enumerate}
\item 射 $f'$ は固有で、そのファイバーは幾何学的に連結である。
\item 射 $\pi : S' \to S$ は整である。
\item $f'_*\mathcal{O}_X = \mathcal{O}_{S'}$ が成り立つ。
\item $S' = \underline{\Spec}_S(f_*\mathcal{O}_X)$ が成り立つ。
\item $S'$ は $X$ における $S$ の正規化である。Morphisms, Definition
\ref{morphisms-definition-normalization-X-in-Y} を参照せよ。
\end{enumerate}
\end{theorem}

\begin{proof}
Lemma \ref{more-morphisms-lemma-stein-universally-closed} を適用して、射
$f' : X \to S'$ を得る。Lemma
\ref{more-morphisms-lemma-stein-universally-closed} の結論に加えて、$f'$ は分離的
（Schemes, Lemma \ref{schemes-lemma-compose-after-separated}）
かつ有限型（Morphisms, Lemma
\ref{morphisms-lemma-permanence-finite-type}）であることに注意する。
したがって $f'$ は固有である。この時点で、$f'$ のファイバーが
幾何学的に連結であるという主張を除くすべての主張が証明された。

\medskip\noindent
$S = \Spec(R)$ はアフィンであると仮定してよい。
$R' = \Gamma(X, \mathcal{O}_X)$ とおく。このとき $S' = \Spec(R')$ である。
したがって $S$ を $S'$ で置き換え、
% P05-EN-CH37-0237
$S = \Spec(R)$ はアフィンであり $R = \Gamma(X, \mathcal{O}_X)$ であると
仮定してよい。次に $s \in S$ を点とする。$U \to S$ をアフィン・スキームの
エタール射とし、$u \in U$ を $s$ に写る点とする。
$X_U \to U$ を $X$ の基底変換とする。
Lemma \ref{more-morphisms-lemma-characterize-geometrically-connected-fibres}
により、$X_U \to U$ の $u$ 上のファイバーが連結であることを
示せば十分である。
Cohomology of Schemes, Lemma \ref{coherent-lemma-flat-base-change-cohomology}
により
$\Gamma(X_U, \mathcal{O}_{X_U}) = \Gamma(U, \mathcal{O}_U)$
を得る。したがって、次を示さなければならない。
% P05-EN-CH37-0238
$S = \Spec(R)$ がアフィン、$X \to S$ が
$\Gamma(X, \mathcal{O}_X) = R$ を満たす固有射、$s \in S$ が点ならば、
ファイバー $X_s$ は連結である。

\medskip\noindent
そのためには、$e \in H^0(X_s, \mathcal{O}_{X_s})$ なる冪等元が
$0$ と $1$ だけであることを示せば十分である
（Lemma \ref{more-morphisms-lemma-stein-universally-closed} により、
$X_s$ が空でないことはすでに分かっている）。Derived Categories of
Schemes, Lemma
\ref{perfect-lemma-proper-idempotent-on-fibre}
により、$R$ を主局所化で置き換えたのち、$e$ は $R$ の元の像であると
仮定してよい。$R \to H^0(X_s, \mathcal{O}_{X_s})$ は
$\kappa(s)$ を経由するので、結論を得る。
\end{proof}

\noindent
応用を一つ与える。

\begin{lemma}
\label{more-morphisms-lemma-geometrically-connected-fibres-towards-normal}
$f : X \to S$ をスキームの射とする。次を仮定する。
\begin{enumerate}
\item $f$ は固有である。
\item $S$ は整で、一般点を $\xi$ とする。
\item $S$ は正規である。
\item $X$ は被約である。
\item $X$ の各既約成分のすべての一般点は $\xi$ に写る。
\item $H^0(X_\xi, \mathcal{O}) = \kappa(\xi)$ が成り立つ。
\end{enumerate}
このとき $f_*\mathcal{O}_X = \mathcal{O}_S$ であり、$f$ のファイバーは
幾何学的に連結である。
\end{lemma}

\begin{proof}
Theorem \ref{more-morphisms-theorem-stein-factorization-general} を適用して
分解 $X \to S' \to S$ を得る。$S' = S$ を示せば十分である。
これは Morphisms, Lemma
\ref{morphisms-lemma-finite-birational-over-normal}
から従う。実際、$X$ は被約なので $S'$ も被約である
（Morphisms, Lemma \ref{morphisms-lemma-normalization-in-reduced}）。
上で引用した定理により、射 $S' \to S$ は整である。
Morphisms, Lemma \ref{morphisms-lemma-normalization-generic}
と仮定 (5) により、$S'$ のすべての一般点は $\xi$ の上にある。
一方、$S'$ は $f_*\mathcal{O}_X$ の相対スペクトルなので、
スキーム論的ファイバー $S'_\xi$ は
$H^0(X_\xi, \mathcal{O})$ のスペクトルであり、仮定によりこれは
$\kappa(\xi)$ に等しい。したがって $S'$ は整スキームであり、
その関数体は $S$ の関数体と等しい。これで証明が完了する。
\end{proof}

\noindent
もう一つ応用を与える。

\begin{lemma}
\label{more-morphisms-lemma-proper-flat-nr-geom-conn-comps-lower-semicontinuous}
$X \to S$ を有限表示の平坦固有射とする。
% P05-EN-CH37-0239
$n_{X/S}$ を、$f$ のファイバーの幾何学的連結成分の個数を数える
$S$ 上の関数であって、
Lemma \ref{more-morphisms-lemma-base-change-fibres-nr-geometrically-connected-components}
で導入されたものとする。このとき $n_{X/S}$ は下半連続である。
\end{lemma}

\begin{proof}
$s \in S$ とし、$n = n_{X/S}(s)$ とおく。$s$ における
$X \to S$ の幾何学的ファイバーは体上の固有スキームであり、
したがって Noetherian であり、それゆえ連結成分を有限個しかもたないので、
$n < \infty$ であることに注意する。
$n_{X/S}|_V \geq n$ となる $s$ の開近傍 $V$ を見いださなければならない。
Theorem \ref{more-morphisms-theorem-stein-factorization-general} による Stein 分解を
$X \to S' \to S$ とする。
Lemma \ref{more-morphisms-lemma-stein-universally-closed-residue-fields} により、
$s$ の上にある点 $s'_1, \ldots, s'_m \in S'$ は有限個であり、
拡大 $\kappa(s'_i)/\kappa(s)$ は有限である。
Lemma \ref{more-morphisms-lemma-etale-makes-integral-split} により、$S$ を $s$ の
エタール近傍で置き換えたのち、スキームとして
$S' = V_1 \amalg \ldots \amalg V_m$ であり、
$s'_i \in V_i$ かつ $\kappa(s'_i)/\kappa(s)$ は純非分離であると
仮定してよい。このときスキーム $X_{s_i'}$ は $\kappa(s)$ 上
幾何学的に連結なので、$m = n$ である。スキーム
$X_i = (f')^{-1}(V_i)$、$i = 1, \ldots, n$ は $S$ 上平坦かつ
有限表示である。したがって $X_i \to S$ の像は開である
（Morphisms, Lemma \ref{morphisms-lemma-fppf-open}）。
ゆえに $s$ のある近傍上で $n_{X/S}$ は少なくとも $n$ である。
\end{proof}

\begin{lemma}
\label{more-morphisms-lemma-proper-flat-geom-red}
$f : X \to S$ をスキームの射とする。次を仮定する。
\begin{enumerate}
\item $f$ は固有、平坦、かつ有限表示である。
\item $f$ の幾何学的ファイバーは被約である。
\end{enumerate}
% P05-EN-CH37-0240
このとき、$f$ のファイバーの幾何学的連結成分の個数を数える関数
$n_{X/S} : S \to \mathbf{Z}$ は局所定数である。
\end{lemma}

\begin{proof}
% P05-EN-CH37-0241
Lemma \ref{more-morphisms-lemma-proper-flat-nr-geom-conn-comps-lower-semicontinuous}
により、関数 $n_{X/S}$ は下半連続である。
$s \in S$ に対し、$\kappa(s)$-代数
$$
A_s = H^0(X_s, \mathcal{O}_{X_s})
$$
を考える。
% P05-EN-CH37-0242
$s$ を $\dim_{\kappa(s)} A_s$ に写す関数
$\beta_0 : X \to \mathbf{Z}$ を定める。Varieties, Lemma
\ref{varieties-lemma-proper-geometrically-reduced-global-sections}
と $X_s$ が幾何学的に被約であることにより、$A_s$ は
$\kappa(s)$ の有限分離拡大の有限積である。したがって
$A_s \otimes_{\kappa(s)} \kappa(\overline{s})$ は
$\kappa(\overline{s})$ の $\beta_0(s)$ 個のコピーの積である。
ゆえに $X_{\overline{s}}$ は $\beta_0(s)$ 個の連結成分をもつ。
言い換えれば、$S$ 上の関数として $n_{X/S} = \beta_0$ である。
したがって Derived Categories of Schemes, Lemma
\ref{perfect-lemma-jump-loci-geometric} により $n_{X/S}$ は上半連続である。
これで証明が完了する。
\end{proof}

\noindent
最後の応用を与える。

\begin{lemma}
\label{more-morphisms-lemma-split-off-proper-part-henselian}
\begin{reference}
進 Noetherian 基底の場合の参考文献は
\cite[III, Proposition 5.5.1]{EGA}
\end{reference}
$(A, I)$ を Hensel 対とする。$X \to \Spec(A)$ を分離的かつ有限型とし、
$X_0 = X \times_{\Spec(A)} \Spec(A/I)$ とおく。
$Y \subset X_0$ を開かつ閉な部分スキームで、
$Y \to \Spec(A/I)$ が固有であるものとする。このとき、$A$ 上固有で
$W \times_{\Spec(A)} \Spec(A/I) = Y$ を満たす開かつ閉な部分スキーム
$W \subset X$ が存在する。
\end{lemma}

\begin{proof}
$\Spec(A/I) \to \Spec(A)$ による基底変換を $T \mapsto T_0$ と表す。
Chow の補題（
Limits, Lemma \ref{limits-lemma-chow-finite-type}）
の形）により、$X'$ が $\mathbf{P}^n_A$ への埋め込みを許すような
全射固有射 $\varphi : X' \to X$ が存在する。
$Y' = \varphi^{-1}(Y)$ とおく。これは $X'_0$ の開かつ閉な部分スキームである。
$(X', Y')$ に対して補題が成り立つと仮定する。
$W' \subset X'$ を、$Y' = W'_0$ を満たす $A$ 上固有な開かつ閉な
部分スキームとする。Morphisms, Lemma
\ref{morphisms-lemma-image-proper-scheme-closed} により
$W = \varphi(W') \subset X$ と
$Q = \varphi(X' \setminus W') \subset X$ は閉部分集合であり、
Morphisms, Lemma \ref{morphisms-lemma-image-proper-is-proper}
により $W$ は $A$ 上固有である。$W \cap Q$ の $\Spec(A)$ における像は
閉である。$(A, I)$ は Hensel 対なので、$W \cap Q$ が空でなければ、
$W \cap Q$ は $\Spec(A/I)$ の上にある点をもつ。しかし、これは
$W'_0 = Y' = \varphi^{-1}(Y)$ に反する。したがって $W$ は
$W_0 = Y$ を満たす、$A$ 上固有な $X$ の開かつ閉な部分スキームである。
これで次段落に述べる場合に帰着した。

\medskip\noindent
$A$ 上の埋め込み $j : X \to \mathbf{P}^n_A$ が存在すると仮定する。
$\overline{X}$ を $j$ のスキーム論的像とする。$j$ は準コンパクトな射なので
（Schemes, Lemma \ref{schemes-lemma-quasi-compact-permanence}）、
$j : X \to \overline{X}$ は開埋め込みである
（Morphisms, Lemma \ref{morphisms-lemma-quasi-compact-immersion}）。
したがって基底変換 $j_0 : X_0 \to \overline{X}_0$ も開埋め込みである。
ゆえに $j_0(Y) \subset \overline{X}_0$ は開である。また Morphisms, Lemma
\ref{morphisms-lemma-image-proper-scheme-closed} により閉でもある。
$(\overline{X}, j_0(Y))$ に対して補題が成り立つと仮定する。
$\overline{W} \subset \overline{X}$ を、
$j_0(Y) = \overline{W}_0$ を満たす対応する $A$ 上固有な開かつ閉な
部分スキームとする。このとき
$T = \overline{W} \setminus j(X)$ は $\overline{W}$ で閉であり、
$\overline{W}$ が $A$ 上固有であることにより、その $\Spec(A)$ における像は
閉である。$(A, I)$ は Hensel 対なので、$T$ が空でなければ、
$T$ には $\Spec(A/I)$ に写る点が存在する。しかし
$j_0(Y) = \overline{W}_0$ は $j(X)$ に含まれるので、これは不可能である。
したがって $\overline{W}$ は $j(X)$ に含まれ、$W \subset X$ を
$j$ により $\overline{W}$ へ同型に写る一意的な開かつ閉な部分スキームと
定めることができる。これで次段落に述べる場合に帰着した。

\medskip\noindent
$X \subset \mathbf{P}^n_A$ が閉部分スキームであると仮定する。
このとき $X \to \Spec(A)$ は固有射である。
$Z = X_0 \setminus Y$ とおく。これは $X_0$ の開かつ閉な部分スキームであり、
$X_0 = Y \amalg Z$ である。Theorem
\ref{more-morphisms-theorem-stein-factorization-general} による Stein 分解を
$X \to X' \to \Spec(A)$ とする。$Y' \subset X'_0$ と
$Z' \subset X'_0$ をそれぞれ $Y$ と $Z$ の像とする。
% P05-EN-CH37-0243
$X \to Z$ のファイバーは幾何学的に連結なので、
$Y' \cap Z' = \emptyset$ である。$X \to X'$ は全射なので、
$X'_0 = Y' \amalg Z'$ となる。$X' \to \Spec(A)$ は整であるから、
$X'$ は $A$ 上整な $A$-代数のスペクトルである。
スペクトルの開かつ閉な部分集合は、対応する環の冪等元と
$1$ 対 $1$ に対応することを思い出そう。Algebra, Lemma
\ref{algebra-lemma-disjoint-decomposition} を参照せよ。したがって
More on Algebra, Lemma \ref{more-algebra-lemma-characterize-henselian-pair}
により、$W'$ と $V'$ が開かつ閉で、
$Y' = W'_0$ および $Z' = V'_0$ となるように
$X' = W' \amalg V'$ と書ける。
% P05-EN-CH37-0244
$W$ を $X$ における逆像とすれば、
証明が完了する。
\end{proof}









\section{一般平坦性による層別化}
\label{more-morphisms-section-generic-flatness-stratification}

\noindent
一般平坦性を用いると、与えられた加群が各層上で平坦になるような
基底の層別化を構成できる。

\begin{lemma}[一般平坦性による層別化]
\label{more-morphisms-lemma-generic-flatness-stratification}
$f : X \to S$ を、準コンパクトかつ準分離的なスキーム間の有限表示射とする。
$\mathcal{F}$ を有限表示の $\mathcal{O}_X$-加群とする。このとき、
$t \geq 0$ と閉部分スキーム
$$
S \supset S_0 \supset S_1 \supset \ldots \supset S_t = \emptyset
$$
であって、$S_i \to S$ は有限型のイデアル層で定義され、
$S_0 \subset S$ は肥厚であり、$\mathcal{F}$ の
$X \times_S (S_i \setminus S_{i + 1})$ への引き戻しが
$S_i \setminus S_{i + 1}$ 上平坦となるものが存在する。
\end{lemma}

\begin{proof}
デカルト図式
$$
\xymatrix{
X \ar[d] \ar[r] & X_0 \ar[d] \\
S \ar[r] & S_0
}
$$
と、$\mathcal{F}$ に引き戻される有限表示
$\mathcal{O}_{X_0}$-加群 $\mathcal{F}_0$ であって、
$X_0$ と $S_0$ が $\mathbf{Z}$ 上有限型となるものをとることができる。
Limits, Proposition \ref{limits-proposition-approximate} および
Lemmas \ref{limits-lemma-descend-finite-presentation}、
\ref{limits-lemma-descend-modules-finite-presentation} を参照せよ。
したがって、$X$ と $S$ は $\mathbf{Z}$ 上有限型であり、
$\mathcal{F}$ は連接 $\mathcal{O}_X$-加群であると仮定してよい。

\medskip\noindent
$X$ と $S$ は $\mathbf{Z}$ 上有限型であり、
$\mathcal{F}$ は連接 $\mathcal{O}_X$-加群であると仮定する。
この場合、すべての準連接イデアルは有限型なので、$S_i$ が有限型イデアルで
切り出されるという条件を確認する必要はない。
$S_0 = S_{red}$ を $S$ の被約化とする。
Morphisms, Proposition
\ref{morphisms-proposition-generic-flatness-reduced}
に述べられた一般平坦性により、$\mathcal{F}$ の
$X \times_S U_0$ への引き戻しが $U_0$ 上平坦となる稠密開集合
$U_0 \subset S_0$ が存在する。$S_1 \subset S_0$ を、その台となる
閉部分集合が $S \setminus U_0$ である被約閉部分スキームとする。
$S_1 \not = \emptyset$ である限り、この方法を続けて
$S_0 \supset S_1 \supset \ldots$ を得る。$S$ は Noetherian なので、
閉部分集合の任意の降鎖は安定する。したがって、ある $t \geq 0$ に対して
$S_t = \emptyset$ となる。
\end{proof}

\begin{lemma}
\label{more-morphisms-lemma-generic-flatness-stratification-scheme}
$f : X \to S$ を、準コンパクトかつ準分離的なスキーム間の有限表示射とする。
このとき、$t \geq 0$ と閉部分スキーム
$$
S \supset S_0 \supset S_1 \supset \ldots \supset S_t = \emptyset
$$
であって、$S_i \to S$ は有限型のイデアル層で定義され、
$S_0 \subset S$ は肥厚であり、
$X \times_S (S_i \setminus S_{i + 1})$ は
$S_i \setminus S_{i + 1}$ 上平坦となるものが存在する。
\end{lemma}

\begin{proof}
Lemma \ref{more-morphisms-lemma-generic-flatness-stratification} を
$\mathcal{F} = \mathcal{O}_X$ として適用する。
\end{proof}

\begin{lemma}[Longke Tang]
\label{more-morphisms-lemma-flatness-criterion}
$f : X \to S$ を全射な有限表示射とする。
$M$ を $H^i(M) = 0$（$i > 0$）を満たす
$D_\QCoh(\mathcal{O}_S)$ の対象とする。次の条件は同値である。
\begin{enumerate}
\item $M$ は次数 $0$ に置かれた平坦 $\mathcal{O}_S$-加群と同型である。
\item $Lf^*M$ は次数 $0$ に置かれた平坦
$\mathcal{O}_X$-加群と同型である。
\end{enumerate}
\end{lemma}

\begin{proof}
(1) が (2) を含意することの証明は省略する。(2) を仮定する。
(1) の証明は局所的な問題なので、$S$ はアフィンであると仮定してよい。
Lemma \ref{more-morphisms-lemma-generic-flatness-stratification-scheme} のような層別化
$S \supset S_0 \supset S_1 \supset \ldots \supset S_t = \emptyset$
を選ぶ。
% P05-EN-CH37-0245
$T_i = S_i \setminus S_{i - 1}$ および $Y_i = X \times_S T_i$ とおく。
包含射を $g_i : T_i \to S$、$h_i : Y_i \to X$ と表す。
$f$ は全射なので、
% P05-EN-CH37-0246
射 $f_i : X_i \to T_i$ は全射かつ平坦である。
$Lh_i^* \circ Lf^* = Lf_i^* \circ Lg_i^*$ なので、(2) から
$Lf_i^* Lg_i^*M$ は次数 $0$ に置かれた平坦
$\mathcal{O}_{Y_i}$-加群と同型である。$f_i$ は平坦かつ全射なので、
$Lg_i^*M$ は次数 $0$ に置かれた平坦
$\mathcal{O}_{T_i}$-加群と同型である。

\medskip\noindent
これが結論を含意することを $t$ に関する帰納法で示す。
次の記法を用いる。$S = \Spec(A)$ とし、$M$ は $D(A)$ の $N$ に対応する
$D_\QCoh(\mathcal{O}_S)$ の対象とする。Derived Categories of Schemes,
Lemma
\ref{perfect-lemma-affine-compare-bounded}
を参照せよ。

\medskip\noindent
基底段階として $t = 1$ とする。
$S_0 = \Spec(A/I)$ と書き、$I$ は有限生成イデアルとする。
% P05-EN-CH37-0247
$S_0 = S$ が集合論的に成り立つので、$I$ は冪零である。
仮定は $N \otimes_A^\mathbf{L} A/I$ の tor 振幅が $[0, 0]$ に
含まれるということである。More on Algebra, Lemma
\ref{more-algebra-lemma-check-tor-dimension-modulo-nilpotent}
により、$N$ についても同じことが成り立つ。

\medskip\noindent
帰納段階として $t > 1$ と仮定する。
$S_{t - 1} = \Spec(A/I)$ と書き、
$I = (f_1, \ldots, f_r)$ は有限生成イデアルとする。
$r$ に関する帰納法で議論する。$D(A_{f_r})$ における
$N_{f_r} = N \otimes_A^\mathbf{L} A_{f_r}$ は、帰納法の仮定により
tor 振幅が $[0, 0]$ に含まれることに注意する
（主開集合 $D(f_r)$ 上の層別化の長さは高々 $t - 1$ である）。
一方、$D(A/f_rA)$ における
$N' = N \otimes_A^\mathbf{L} A/f_rA$ を考える。
$S' = \Spec(A/f_rA)$ とおくと、層別化
$$
S' \supset S' \cap S_0 \supset S' \cap S_1 \supset \ldots \supset
S' \cap S_{t - 1} \supset
S' \cap S_t = \emptyset
$$
を得る。ここで $S' \cap S_{t - 1}$ は $S'$ において $r - 1$ 個の方程式で
切り出される。前と同様に、$M$ の各部分
$S' \cap S_i \setminus S' \cap S_{i + 1}$ への導来引き戻しの
tor 振幅は $[0, 0]$ に含まれる。したがって $r$ に関する帰納法により、
$N'$ の tor 振幅は $[0, 0]$ に含まれる。
最後に More on Algebra, Lemma \ref{more-algebra-lemma-f-flat}
により、$N$ の tor 振幅も $[0, 0]$ に含まれる。
\end{proof}

\begin{lemma}
\label{more-morphisms-lemma-cokernel-flat}
$R$ を Noetherian 整域とし、$R \to A \to B$ を有限型環準同型とする。
$M$ を有限 $A$-加群とし、
% P05-EN-CH37-0248
$N$ を有限 $B$-加群とする。$M \to N$ を $A$-線形写像とする。
% P05-EN-CH37-0249
$M_f \to N_f$ の余核が平坦 $R_f$-加群となるような非零元
$f \in R$ が存在する。
\end{lemma}

\begin{proof}
$M$ を $M \to N$ の像で置き換え、$M \subset N$ と仮定してよい。
ある素イデアル $\mathfrak q_i \subset B$ に対して
$N_i/N_{i - 1} = B/\mathfrak q_i$ となるようなフィルトレーション
$0 = N_0 \subset N_1 \subset \ldots \subset N_t = N$ を選ぶ。
Algebra, Lemma \ref{algebra-lemma-filter-Noetherian-module}
を参照せよ。$M_i = M \cap N_i$ とおく。このとき $Q = N/M$ は
部分加群 $Q_i = N_i/M_i$ によるフィルトレーションをもつ。
% P05-EN-CH37-0250
環 R のある非零元 $f$ で局所化した後に $Q_i/Q_{i - 1}$ が平坦になることを
示せば十分である（Algebra, Lemma \ref{algebra-lemma-flat-ses}
により平坦加群の拡大は平坦である）。
% P05-EN-CH37-0251
$Q_i/Q_{i - 1}$ は写像
$M_i/M_{i - 1} \to N_i/N_{i - 1}$ の余核と同型なので、
次段落で論じる場合に帰着する。

\medskip\noindent
$B$ は整域で、$M \subset N = B$ であると仮定する。
$A$ を $B$ における $A$ の像で置き換え、$A \subset B$ と仮定してよい。
一般平坦性により、$A$ と $B$ は $R$ 上平坦であると仮定してよい
（Algebra, Lemma \ref{algebra-lemma-generic-flatness-Noetherian}）。
ここで $R$ を主局所化で置き換えた後に $M \to B$ が
$R$-普遍的に単射となることを示せば十分である
（Algebra, Lemma \ref{algebra-lemma-ui-flat-domain}）。
一般自由性により、$B_g$ が自由 $A_g$-加群となるような非零元
$g \in A$ をとることができる
（Algebra, Lemma \ref{algebra-lemma-generic-flatness-Noetherian}）。
したがって、$A_g$-加群としての直和因子 $M' \subset B_g$ であって、
有限自由 $A_g$-加群であり、$M \to B \to B_g$ が $M'$ を経由するものを
選べる。明らかに、$R$ を主局所化で置き換えた後に
$M \to M'$ が $R$-普遍的に単射となることを示せば十分である。

\medskip\noindent
$M' = A_g^{\oplus n}$ とする。$M \subset M'$ は有限 $A$-加群なので、
ある $m \geq 0$ に対して $M$ は
$(1/g^m)A^{\oplus n}$ に含まれる。$M'$ の基底を取り替え、
$M \subset A^{\oplus n}$ と仮定してよい。このとき、$R$ を主局所化で
置き換えた後に $A^{\oplus n}/M$ と $A_g/A$ が $R$-平坦になることを
示せば十分である。実際、そのとき Algebra, Lemma
\ref{algebra-lemma-flat-tor-zero} により
$M \to A^{\oplus n}$ と $A^{\oplus n} \to A_g^{\oplus n}$ は
普遍的に単射であり、したがってその合成
$M \to M' = A_g^{\oplus n}$ も普遍的に単射である。

\medskip\noindent
一般平坦性（上の参照を見よ）により、加群 $A^{\oplus n}/M$ は
$R$-平坦であると仮定してよい。商 $A_g/A$ については、
$$
A_g/A = \colim (1/g^m)A/A \cong \colim A/g^mA
$$
であること、および加群 $A/g^mA$ が長さ $m$ のフィルトレーションをもち、
その逐次商が $A/gA$ と同型であることを用いる。
再び一般平坦性により $A/gA$ は $R$-平坦であると仮定してよい。
したがって各 $A/g^mA$ は $R$-平坦であり、それゆえ $A_g/A$ も
平坦である。
\end{proof}

\noindent
$f : X \to Y$ を基底スキーム $S$ 上のスキームの射とする。
$Z \subset Y$ を $f$ のスキーム論的像とする。
Morphisms, Section \ref{morphisms-section-scheme-theoretic-image}
を参照せよ。$g : S' \to S$ をスキームの射とし、
$f' : X \times_S S' \to Y \times_S S'$ を $g$ による $f$ の基底変換とする。
$Z \times_S S' \subset Y \times_S S'$ が $f'$ のスキーム論的像であるとは
限らない。上のようなすべての $g$ に対して $f'$ のスキーム論的像が
$Z \times_S S'$ に等しいとき、{\it $f/S$ のスキーム論的像の形成は
任意の基底変換と可換である}ということにする。

\begin{lemma}
\label{more-morphisms-lemma-base-change-scheme-theoretic-image}
$S$ を準コンパクトかつ準分離的なスキームとする。
$f : X \to Y$ を $S$ 上のスキームの射とし、$X$ と $Y$ はともに
$S$ 上有限表示であるとする。このとき、$t \geq 0$ と閉部分スキーム
$$
S \supset S_0 \supset S_1 \supset \ldots \supset S_t = \emptyset
$$
であって、次の性質をもつものが存在する。
\begin{enumerate}
\item $S_i \to S$ は有限型のイデアル層で定義される。
\item $S_0 \subset S$ は肥厚である。
\item $T_i = S_i \setminus S_{i + 1}$ とし、$f_i$ を $f$ の
$T_i$ への基底変換とすると、$f_i/T_i$ のスキーム論的像の形成は
任意の基底変換と可換である（補題の直前の議論を参照せよ）。
\end{enumerate}
\end{lemma}

\begin{proof}
可換図式
$$
\xymatrix{
X \ar[d] \ar[r] & Y \ar[d] \ar[r] & S \ar[d] \\
U \ar[r] & V \ar[r] & W
}
$$
であって、各正方形がデカルトであり、$U$、$V$、$W$ が
$\mathbf{Z}$ 上有限型となるものをとることができる。実際、まず
Limits, Proposition \ref{limits-proposition-approximate} を用いて
$S$ を、アフィンな遷移射をもつ $\mathbf{Z}$ 上有限型なスキームの
余フィルター付き極限として書き、次に Limits, Lemma
\ref{limits-lemma-descend-finite-presentation} を用いて射
$X \to Y$ を降下させる。これで次段落で論じる場合に帰着する。

\medskip\noindent
$S$ は Noetherian であると仮定する。この場合、すべての準連接イデアルは
有限型なので、$S_i$ が有限型イデアルで切り出されるという条件を
確認する必要はない。$S_0 = S_{red}$ を $S$ の被約化とする。
$\eta \in S_0$ を $S_0$ のある既約成分の一般点とする。
$S_0$ の台となる位相空間に関する Noetherian 帰納法により、
$\eta$ を含まない $S_0$ の任意の閉部分スキームに対して結論が
成り立つと仮定してよい。したがって、$f$ の $U_0$ への基底変換
$f_0$ が性質 (3) をもつような開近傍 $U_0 \subset S_0$ が存在することを
示せば十分である。

\medskip\noindent
$R$ を Noetherian 整域とする。$f : X \to Y$ を $R$ 上有限型な
スキームの射とする。前段落の議論により、ある非零元 $g \in R$ に対して
$R$ を $R_g$ で、$X$ と $Y$ をそれぞれ $R_g$ への基底変換で
置き換えた後に、$f/R$ のスキーム論的像の形成が任意の基底変換と
可換になることを示せば十分である。

\medskip\noindent
% P05-EN-CH37-0252
$Y = V_1 \cup \ldots V_n$ をアフィン開被覆とする。
$U_i = f^{-1}(V_i)$ とおく。$R$ 上の各射 $U_i \to V_i$ に対して
主張が成り立つなら、$f$ に対しても成り立つ。実際、Morphisms, Lemma
\ref{morphisms-lemma-quasi-compact-scheme-theoretic-image}
により、$U_i \to V_i$ のスキーム論的像は $V_i$ と
$f : X \to Y$ のスキーム論的像との共通部分である。
したがって $Y$ はアフィンであると仮定してよい。

\medskip\noindent
% P05-EN-CH37-0253
$X = U_1 \cup \ldots U_n$ をアフィン開被覆とする。このとき
$X \to Y$ のスキーム論的像は、
% P05-EN-CH37-0254
$\coprod U_i \to Y$ のスキーム論的像と同じである。
したがって $X$ はアフィンであると仮定してよい。

\medskip\noindent
% P05-EN-CH37-0255
$X = \Spec(A)$、$Y = \Spec(B)$ とし、$f$ は $R$-代数準同型
$\varphi : A \to B$ に対応するとする。このとき $f$ のスキーム論的像は
$\Spec(A/\Ker(\varphi))$ であり、基底変換後も同様である
（アフィン射による基底変換についてであるが、それらについて確認すれば
十分である）。したがって、すべての環準同型 $R \to R'$ に対して
$\Ker(\varphi \otimes_R R') = \Ker(\varphi) \otimes_R R'$ ならば、
スキーム論的像の形成は基底変換と可換である。

\medskip\noindent
$R$ の適当な非零元 $g$ に対し、$R$、$A$、$B$ をそれぞれ
$R_g$、$A_g$、$B_g$ で置き換えることにより、$A$ と $B$ は
$R$ 上平坦であると仮定してよい。Lemma \ref{more-morphisms-lemma-cokernel-flat}
により、$B/A$ も平坦 $R$-加群であると仮定してよい。
このとき
$0 \to \Ker(\varphi) \to A \to B \to B/A \to 0$ は
平坦 $R$-加群の完全列であり、これは望ましい基底変換の主張を含意する。
\end{proof}






\section{射の層別化}
\label{more-morphisms-section-stratifying-morphisms}

\noindent
$f : X \to S$ を準コンパクトかつ準分離的なスキームの有限表示射とする。
Section \ref{more-morphisms-section-generic-flatness-stratification} では、
$X$ が各層上平坦となるように $S$ を層別化できることを見た。
% P05-EN-CH37-0256
本節では、滑らかな層が得られるような $S$ と $X$ の層別化を探す。
これはそのままではうまくいかず、有限局所自由射による基底変換も必要になる。

\begin{lemma}
\label{more-morphisms-lemma-smoothness-stratification-at-generic-point}
$f : X \to S$ をスキームの有限表示射とする。
$\eta \in S$ を $S$ の既約成分の一般点とする。
$S$ は被約であると仮定する。このとき、次が存在する。
\begin{enumerate}
\item $\eta$ を含む開部分スキーム $U \subset S$。
\item 全射、普遍的に単射、かつ有限局所自由な射 $V \to U$。
\item $t \geq 0$ と閉部分スキーム
$$
X \times_S V \supset Z_0 \supset Z_1 \supset \ldots \supset Z_t = \emptyset
$$
であって、$Z_i \to X \times_S V$ は有限型のイデアル層で定義され、
$Z_0 \subset X \times_S V$ は肥厚であり、射
$Z_i \setminus Z_{i + 1} \to V$ は滑らかであるもの。
\end{enumerate}
\end{lemma}

\begin{proof}
$S$ を $\eta$ の開近傍で、$X$ をこの開集合への制限で置き換えてよいことは
明らかである。したがって $S = \Spec(A)$ であり、$A$ は被約環、
$\eta$ は極小素イデアル $\mathfrak p$ に対応すると仮定してよい。
この場合、局所環 $\mathcal{O}_{S, \eta} = A_\mathfrak p$ は
$\kappa(\mathfrak p)$ に等しいことを思い出そう。
Algebra, Lemma \ref{algebra-lemma-minimal-prime-reduced-ring}
を参照せよ。

\medskip\noindent
$k = \kappa(\eta)$ 上のスキーム $X_\eta$ に Varieties, Lemma
\ref{varieties-lemma-smooth-stratification} を適用する。
これによって得られる純非分離体拡大を $k'/k$ と表す。
次段落で $k' = k$ の場合に帰着する
（この段階は補題の主張における射 $V \to U$ を見いだすことに対応する。
特に $\kappa(\mathfrak p)$ の標数が零なら $V = U$ ととれる）。

\medskip\noindent
$k = \kappa(\mathfrak p)$ の標数が零なら $k' = k$ である。
$k = \kappa(\mathfrak p)$ の標数が $p > 0$ なら、
$p$ は $A_\mathfrak p = \kappa(\mathfrak p)$ で零に写る。
したがって $A$ を主局所化で置き換えることにより
（すなわち $S$ を縮小することにより）、$A$ において $p = 0$ と
仮定してよい。$k' \not = k$ なら、
% P05-EN-CH37-0257
$\beta \in k'$、$\beta \not \in k$、$\beta^p \in k$ を満たす
元が存在する。$A$ を主局所化で置き換え、
$\beta^p = a$ を満たす $a \in A$ が存在すると仮定してよい。
$A' = A[x]/(x^p - a)$ とおく。このとき
$S' = \Spec(A') \to \Spec(A) = S$ は有限局所自由、全射、かつ
普遍的に単射である。さらに、$\mathfrak p$ の上にある一意的な素イデアルを
$\mathfrak p' \subset A'$ と表すと、
$A'_{\mathfrak p'} = k(\beta)$ であり、$k'/k(\beta)$ の次数はより小さい。
したがって、$S$ を $S'$ で、$\eta$ を $\mathfrak p'$ に対応する点
$\eta'$ で置き換えると、$\eta$ の剰余体上の $k'$ の次数が減少する。
この操作を続け、帰納法により
$k' = \kappa(\mathfrak p) = \kappa(\eta)$ の場合に帰着する。

\medskip\noindent
したがって、$S$ はアフィンかつ被約であり、$t \geq 0$ と閉部分スキーム
$$
X_\eta \supset Z_{\eta, 0} \supset Z_{\eta, 1}
\supset \ldots \supset Z_{\eta, t} = \emptyset
$$
であって、$Z_{\eta, 0} = (X_\eta)_{red}$ であり、すべての $i$ に対して
$Z_{\eta, i} \setminus Z_{\eta, i + 1}$ が $\eta$ 上滑らかとなるものが
存在すると仮定してよい。
$\kappa(\eta) = \kappa(\mathfrak p) = A_\mathfrak p$ は、
$a \in A$、$a \not \in \mathfrak p$ に対する $A_a$ の
フィルター付き余極限であることを思い出そう。Algebra, Lemma
\ref{algebra-lemma-localization-colimit} を参照せよ。
したがって、ある $a \in A$、$a \not \in \mathfrak p$ に対し、
上の図式を $\Spec(A_a)$ 上の対応する図式へ降下させることができる。
より正確には、$S$ を $\Spec(A_a)$ で置き換えた後に、
$t \geq 0$ と閉部分スキーム
$$
X \supset Z_0 \supset Z_1 \supset \ldots \supset Z_t = \emptyset
$$
であって、$Z_i \to X$ は有限表示の閉埋め込み、
$Z_0 \to X$ は肥厚、$Z_i \setminus Z_{i + 1}$ は $S$ 上滑らかとなるものが
存在すると仮定してよい。言い換えれば、補題が成り立つ。
さらに詳しく述べると、まず
Limits, Lemma \ref{limits-lemma-descend-finite-presentation}
を用いて、$S$ 上有限表示な射
$$
Z_t \to Z_{t - 1} \to \ldots \to Z_0 \to X
$$
であって、$\eta$ への基底変換が上の与えられた閉部分スキーム間の包含を
与えるものを得る。$S$ をさらに縮小した後に、各射は閉埋め込みであると
仮定してよい。Limits, Lemma
\ref{limits-lemma-descend-closed-immersion-finite-presentation}
を参照せよ。$S$ を縮小した後に $Z_0 \to X$ は全射であり、
したがって肥厚であると仮定してよい。Limits, Lemma
\ref{limits-lemma-descend-surjective} を参照せよ。
$S$ をもう一度縮小した後に
$Z_i \setminus Z_{i + 1} \to S$ は滑らかであると仮定してよい。
Limits, Lemma \ref{limits-lemma-descend-smooth} を参照せよ。
これで証明が完了する。
\end{proof}

\begin{lemma}
\label{more-morphisms-lemma-smoothness-stratification}
$f : X \to S$ を準コンパクトかつ準分離的なスキーム間の有限表示射とする。
このとき、$t \geq 0$ と閉部分スキーム
$$
S \supset S_0 \supset S_1 \supset \ldots \supset S_t = \emptyset
$$
であって、次を満たすものが存在する。
\begin{enumerate}
\item $S_i \to S$ は有限型のイデアル層で定義される。
\item $S_0 \subset S$ は肥厚である。
\item 各 $i$ に対し、全射な有限局所自由射
$T_i \to S_i \setminus S_{i + 1}$ が存在する。
\item 各 $i$ に対し、$t_i \geq 0$ と閉部分スキーム
$$
X_i = X \times_S T_i \supset Z_{i, 0}
\supset Z_{i, 1} \supset \ldots \supset Z_{i, t_i} = \emptyset
$$
であって、$Z_{i, j} \to X_i$ は有限型のイデアル層で定義され、
$Z_{i, 0} \subset X_i$ は肥厚であり、射
$Z_{i, j} \setminus Z_{i, j + 1} \to T_i$ は滑らかであるものが存在する。
\end{enumerate}
\end{lemma}

\begin{proof}
デカルト図式
$$
\xymatrix{
X \ar[d] \ar[r] & X_0 \ar[d] \\
S \ar[r] & S_0
}
$$
であって、$X_0$ と $S_0$ が $\mathbf{Z}$ 上有限型となるものを
とることができる。Limits, Proposition
\ref{limits-proposition-approximate} および Lemma
\ref{limits-lemma-descend-finite-presentation} を参照せよ。
したがって $X$ と $S$ は $\mathbf{Z}$ 上有限型であると仮定してよい。
実際、$X_0 \to S_0$ に対する補題の問題の解を基底変換すれば
$S$ 上の解になる。詳細は省略する。

\medskip\noindent
$X$ と $S$ は $\mathbf{Z}$ 上有限型であると仮定する。
この場合、すべての準連接イデアルは有限型なので、$S_i$ が有限型イデアルで
切り出されるという条件を確認する必要はない。
$S_0 = S_{red}$ を $S$ の被約化とする。
$\eta \in S_0$ を既約成分の一般点とする。
Lemma \ref{more-morphisms-lemma-smoothness-stratification-at-generic-point} により、
開部分スキーム $U \subset S_0$、全射、普遍的に単射、かつ有限局所自由な射
$V \to U$、$t_0 \geq 0$、および閉部分スキーム
$$
X \times_S V \supset Z_{0, 0} \supset Z_{0, 1} \supset \ldots \supset
Z_{0, t_0} = \emptyset
$$
であって、$Z_{0, i} \to X \times_S V$ は有限型のイデアル層で定義され、
$Z_{0, 0} \subset X \times_S V$ は肥厚であり、
$Z_{0, i} \setminus Z_{0, i + 1} \to V$ は滑らかとなるものを
とることができる。次に $S_1 \subset S_0$ を、
$S_0 \setminus U$ 上の被約誘導スキーム構造とする。
$S$ の台となる位相空間に関する Noetherian 帰納法により、
$X \times_S S_1 \to S_1$ に対して補題が成り立つと仮定してよい。
これにより $t \geq 1$ と
$$
S_1 = S_1 \supset S_2 \supset \ldots \supset S_t = \emptyset
$$
および補題の主張にあるような $t_i$ と $Z_{i, j}$ が得られる。
これで補題が示された。
\end{proof}


















\section{相対次元一の射の改良}
\label{more-morphisms-section-make-good-curves}

\noindent
任意の曲線は、基礎体を拡大し、コンパクト化し、正規化することにより、
滑らかかつ射影的にできる。これは、一般ファイバーの次元が $1$ である
有限型射についての結果も含意する。

\begin{lemma}
\label{more-morphisms-lemma-make-good-curves}
$f : X \to S$ をスキームの射とする。$\eta \in S$ を $S$ の既約成分の
一般点とする。$f$ は分離的、有限表示で、
$\dim(X_\eta) \leq 1$ を満たすと仮定する。このとき可換図式
$$
\xymatrix{
\overline{Y}_1 \amalg \ldots \amalg \overline{Y}_n \ar[rd] &
Y_1 \amalg \ldots \amalg Y_n \ar[r]_-\nu \ar[d] \ar[l]^j &
X_V \ar[r] \ar[d] &
X_U \ar[r] \ar[d] &
X \ar[d]^f \\
& T_1 \amalg \ldots \amalg T_n \ar[r] &
V \ar[r] &
U \ar[r] &
S
}
$$
であって、次の性質をもつスキームの図式が存在する。
\begin{enumerate}
\item $U \subset S$ は $\eta$ の開近傍である。
\item $V \to U$ は有限、全射、かつ普遍的に単射な射である。
\item $X_U = U \times_S X$ と $X_V = V \times_S X$ は基底変換である。
\item $\nu$ は有限かつ全射であり、次を満たす開集合
$W \subset X_V$ が存在する。
\begin{enumerate}
\item $W$ は $X_V \to V$ のすべてのファイバーで稠密である。
\item $\nu^{-1}(W) \cap Y_i$ は $Y_i \to T_i$ のすべてのファイバーで
稠密である。
\item $\nu^{-1}(W) \to W$ は肥厚である。
\end{enumerate}
\item $j$ は開埋め込みである。
\item $T_i \to V$ は有限エタールである。
\item $Y_i \to T_i$ は全射かつ滑らかである。
\item $\overline{Y}_i \to T_i$ は滑らかかつ固有であり、
次元 $\leq 1$ の幾何学的に連結なファイバーをもつ。
\end{enumerate}
\end{lemma}

\begin{proof}
$S$ を $\eta$ の開近傍で、$X$ をこの開集合への制限で置き換えてよいことは
明らかである。さらに、$S$ をその被約化で、$X$ をこの被約化への
基底変換で置き換えてよい。したがって $S = \Spec(A)$ であり、
$A$ は被約環、$\eta$ は極小素イデアル $\mathfrak p$ に対応すると
仮定してよい。この場合、局所環
$\mathcal{O}_{S, \eta} = A_\mathfrak p$ は
$\kappa(\mathfrak p)$ に等しいことを思い出そう。Algebra, Lemma
\ref{algebra-lemma-minimal-prime-reduced-ring} を参照せよ。

\medskip\noindent
Varieties, Lemma
\ref{varieties-lemma-dim-1-projective-completion-after-insep}
を $k = \kappa(\eta)$ 上のスキーム $X_\eta$ に適用する。
これによって得られる純非分離体拡大を $k'/k$ と表す。
次段落で $k' = k$ の場合に帰着する
（この段階は補題の主張における射 $V \to U$ を見いだすことに対応する。
特に $\kappa(\mathfrak p)$ の標数が零なら $V = U$ ととれる）。

\medskip\noindent
$k = \kappa(\mathfrak p)$ の標数が零なら $k' = k$ である。
$k = \kappa(\mathfrak p)$ の標数が $p > 0$ なら、
$p$ は $A_\mathfrak p = \kappa(\mathfrak p)$ で零に写る。
したがって $A$ を主局所化で置き換えることにより
（すなわち $S$ を縮小することにより）、$A$ において $p = 0$ と
仮定してよい。$k' \not = k$ なら、
% P05-EN-CH37-0258
$\beta \in k'$、$\beta \not \in k$、$\beta^p \in k$ を満たす
元が存在する。$A$ を主局所化で置き換え、
$\beta^p = a$ を満たす $a \in A$ が存在すると仮定してよい。
$A' = A[x]/(x^p - a)$ とおく。このとき
$S' = \Spec(A') \to \Spec(A) = S$ は有限、全射、かつ普遍的に単射である。
さらに、$\mathfrak p$ の上にある一意的な素イデアルを
$\mathfrak p' \subset A'$ と表すと、
$A'_{\mathfrak p'} = k(\beta)$ であり、$k'/k(\beta)$ の次数はより小さい。
したがって、$S$ を $S'$ で、$\eta$ を $\mathfrak p'$ に対応する点
$\eta'$ で置き換えると、$\eta$ の剰余体上の $k'$ の次数が減少する。
この操作を続け、帰納法により
$k' = \kappa(\mathfrak p) = \kappa(\eta)$ の場合に帰着する。

\medskip\noindent
したがって、$S$ はアフィンかつ被約であり、図式
$$
\xymatrix{
\overline{Y}_{1, \eta} \amalg \ldots \amalg \overline{Y}_{n, \eta} \ar[rd] &
Y_{1, \eta} \amalg \ldots \amalg Y_{n, \eta} \ar[r]_-\nu \ar[d] \ar[l]^j &
X_\eta \ar[d] \\
& \Spec(k_1) \amalg \ldots \amalg \Spec(k_n) \ar[r] &
\eta
}
$$
であって、次の性質をもつスキームの図式が存在すると仮定してよい。
\begin{enumerate}
\item $\nu$ は $X_\eta$ の正規化である。
\item $j$ は稠密な像をもつ開埋め込みである。
\item $i = 1, \ldots, n$ に対し、
$k_i/\kappa(\eta)$ は有限分離拡大である。
\item $\overline{Y}_{i, \eta}$ は $k_i$ 上滑らか、射影的、かつ
幾何学的に既約であり、次元 $\leq 1$ である。
\end{enumerate}
$\kappa(\eta) = \kappa(\mathfrak p) = A_\mathfrak p$ は、
$a \in A$、$a \not \in \mathfrak p$ に対する $A_a$ の
フィルター付き余極限であることを思い出そう。Algebra, Lemma
\ref{algebra-lemma-localization-colimit} を参照せよ。
したがって、ある $a \in A$、$a \not \in \mathfrak p$ に対し、
上の図式を $\Spec(A_a)$ 上の対応する図式へ降下させることができる。
より正確には、$S$ を $\Spec(A_a)$ で置き換えた後に、可換図式
$$
\xymatrix{
\overline{Y}_1 \amalg \ldots \amalg \overline{Y}_n \ar[rd] &
Y_1 \amalg \ldots \amalg Y_n \ar[r]_-\nu \ar[d] \ar[l]^j &
X \ar[d] \\
& T_1 \amalg \ldots \amalg T_n \ar[r] & S
}
$$
であって、$\eta$ への基底変換が上の図式であり、
次の性質をもつスキームの図式が存在すると仮定してよい。
\begin{enumerate}
\item $\nu$ は有限な全射である。
\item $j$ は開埋め込みである。
\item $i = 1, \ldots, n$ に対し、$T_i \to S$ は有限エタールである。
\item $Y_i \to T_i$ は滑らかかつ全射である。
\item $\overline{Y}_i \to T_i$ は滑らかかつ固有であり、
次元 $\leq 1$ の幾何学的に連結なファイバーをもつ。
\end{enumerate}
このため、まず
Limits, Lemma \ref{limits-lemma-descend-finite-presentation}
を用いて、基底変換が前の図式となる図式を得る。次に Limits, Lemmas
\ref{limits-lemma-descend-etale}、
\ref{limits-lemma-descend-smooth},
\ref{limits-lemma-descend-finite-finite-presentation},
\ref{limits-lemma-limit-affine},
\ref{limits-lemma-descend-open-immersion},
\ref{limits-lemma-eventually-proper}、および
\ref{limits-lemma-descend-surjective}
を用いる。これにより、$\nu$ は有限かつ全射、$j$ は開埋め込み、
$T_i \to S$ は有限エタール、
% P05-EN-CH37-0259
$Y_i \to T$ は滑らか、$\overline{Y}_i \to T_i$ は固有かつ滑らかとなる。
$Y_i$ は空ではありえず、滑らかな射は開であり、$T_i \to S$ は
有限エタールなので、$S$ を縮小した後に
$Y_i \to T_i$ は全射であると仮定してよい。最後に、
$T_i$ の $\eta$ 上にある一意的な点 $\eta_i = \Spec(k_i)$ 上の
$\overline{Y}_i \to T_i$ のファイバーは幾何学的に連結である。
したがって、もう一度縮小することにより、すべてのファイバーについて
同じことが成り立つと仮定してよい。
Lemma \ref{more-morphisms-lemma-proper-flat-geom-red} を参照せよ。

\medskip\noindent
残るのは、(a)、(b)、(c) を満たす開集合 $W \subset X$ の存在を
証明することである。$\nu_\eta : \coprod Y_{i, \eta} \to X_\eta$ は
正規化射なので、
Varieties, Lemma \ref{varieties-lemma-normalization-locally-algebraic}
により、$\nu^{-1}(W_\eta) \to W_\eta$ が $W_\eta$ の被約化の
$W_\eta$ への包含に等しくなるような稠密開集合
$W_\eta \subset X_\eta$ が存在する。
$W \subset X$ を、その $\eta$ 上のファイバーが今見いだした開集合
$W_\eta$ である準コンパクトな開集合とする。
$A = \Gamma(S, \mathcal{O}_S)$ をさらに局所化で置き換え、
$\nu^{-1}(W) \to W$ は閉埋め込みであると仮定してよい。
Limits, Lemma
\ref{limits-lemma-descend-closed-immersion-finite-presentation} を参照せよ。
$\nu$ は全射でもあるので、$\nu^{-1}(W) \to W$ は肥厚である。
$W_i = \nu^{-1}(W) \cap Y_i$ とおく。
$S$ をもう一度縮小し、すべての $i$ に対して $W_i \to T_i$ は全射であると
仮定してよい（上と同じ議論による）。このとき $Y_i \to T_i$ は
幾何学的に既約なファイバーをもつので、
$W_i \subset Y_i$ は $Y_i \to T_i$ のすべてのファイバーで稠密である。
$\nu$ は有限かつ全射なので、
$W = \nu(\nu^{-1}(W))$ も $X \to S$ のすべてのファイバーで稠密である。
\end{proof}











\section{分離的かつ局所準有限な射の降下}
\label{more-morphisms-section-separated-locally-quasi-finite}

\noindent
本節では「分離的かつ局所準有限な射は fppf 被覆に関して降下を
満たす」ことを示す。用語については Descent, Definition
\ref{descent-definition-descending-types-morphisms} を参照されたい。
この結果は、Raynaud と Gruson による（多くの理由から）見事な論文の
\cite[Lemma 5.7.1]{GruRay} の証明中に含まれている。
また、射が局所有限表示であるという追加仮定のもとでは、
\cite{Murre-representation} および
\cite[Expos\'e X, Lemma 5.4]{SGA3} にも見いだせる。
形式的な主張は次のとおりである。

\begin{lemma}
\label{more-morphisms-lemma-separated-locally-quasi-finite-morphisms-fppf-descend}
$S$ をスキームとする。
$\{X_i \to S\}_{i\in I}$ を fppf 被覆とする。Topologies, Definition
\ref{topologies-definition-fppf-covering} を参照されたい。
$(V_i/X_i, \varphi_{ij})$ を $\{X_i \to S\}$ に関する降下データとする。
各射 $V_i \to X_i$ が分離的かつ局所準有限ならば、この降下データは
有効である。
\end{lemma}

\begin{proof}
分離性と局所準有限性は、任意の基底変換で保たれるスキームの射の
性質である。Schemes, Lemma \ref{schemes-lemma-separated-permanence} および
Morphisms, Lemma \ref{morphisms-lemma-base-change-quasi-finite} を参照されたい。
したがって Descent, Lemma \ref{descent-lemma-descending-types-morphisms} が
適用でき、fppf 被覆がアフィン・スキーム間の単一の平坦かつ全射な
% P05-EN-CH37-0260
有限表示射 $\{X \to S\}$ で与えられる場合に補題の主張を示せば十分である。
$X = \Spec(A)$、$S = \Spec(R)$ と書く。このとき $R \to A$ は忠実平坦な
環準同型である。$(V, \varphi)$ を $S$ 上の $X$ に関する降下データとし、
$\pi : V \to X$ は分離的かつ局所準有限であると仮定する。

\medskip\noindent
$W^1 \subset V$ を任意のアフィン開集合とする。
$W = \text{pr}_1(\varphi(W^1 \times_S X)) \subset V$ を考える。
対応する図式は次のとおりである。
$$
\xymatrix{
W^1 \times_S X \ar[rrrrr] \ar[ddd] \ar[rd]
& & & & &
\varphi(W^1 \times_S X) \ar[ddd] \ar[ld] \\
& V \times_S X \ar[rrr]^\varphi \ar[rd] \ar[dd]
& & &
X \times_S V \ar[ld] \ar[dd] & \\
& &
X \times_S X \ar[r]^1 \ar[d]_{\text{pr}_0}
&
X \times_S X \ar[d]^{\text{pr}_1}
& & \\
W^1 \ar[r] &
V \ar[r] &
X &
X &
V \ar[l] &
W \ar[l]
}
$$
さて、$X \to S$ は平坦かつ有限表示なので普遍的に開である
（Morphisms, Lemma \ref{morphisms-lemma-fppf-open}）。したがって $W$ は
開である。さらに、$W$ が準コンパクトであり $W^1 \subset W$ であることも
明らかである。また、$\varphi$ の余サイクル条件により
$\varphi(W \times_S X) = X \times_S W$ である。そこで $\varphi$ を
$W \times_S X$ に制限することにより、新しい降下データ
$(W, \varphi')$ を得る。射 $W \to X$ は準コンパクト、分離的、かつ
局所準有限である。したがって定義により分離的かつ準有限であり、
Lemma \ref{more-morphisms-lemma-quasi-finite-separated-quasi-affine} により準アフィンである。
ゆえに Descent, Lemma \ref{descent-lemma-quasi-affine} から、降下データ
$(W, \varphi')$ は有効である。

\medskip\noindent
言い換えれば、降下射 $\varphi$ に関して安定な準コンパクト開集合 $W_i$ による
開被覆 $V = \bigcup W_i$ が存在する。さらに、このような各準コンパクト
開集合 $W \subset V$ に対し、対応する降下データ
% P05-EN-CH37-0261
$(W, \varphi')$ は有効である。したがって、開集合 $W_i$ を降下させて得られる
スキームを貼り合わせることにより、もとの降下データも有効である。
Descent, Lemma \ref{descent-lemma-effective-for-fpqc-is-local-upstairs} を
参照されたい。
\end{proof}






\section{相対的有限表示}
\label{more-morphisms-section-finite-type-finite-presentation}

\noindent
$R \to A$ を有限型の環準同型とし、$M$ を $A$-加群とする。
More on Algebra,
Section \ref{more-algebra-section-relative-finite-presentation} では、
$M$ が $R$ に相対的に有限表示されることの意味を定義した。
また、この概念は局所化に関して良好な性質をもち、貼り合わせられることも
証明した。したがって、対応する大域的概念を次のように定義できる。

\begin{definition}
\label{more-morphisms-definition-relatively-finitely-presented-sheaf}
$f : X \to S$ を局所有限型なスキームの射とし、$\mathcal{F}$ を
準連接 $\mathcal{O}_X$-加群とする。アフィン開被覆
$S = \bigcup V_i$ が存在し、各 $i$ に対してアフィン開被覆
$f^{-1}(V_i) = \bigcup_j U_{ij}$ が存在して、$\mathcal{F}(U_{ij})$ が
$\mathcal{O}_S(V_i)$ に相対的に有限表示される
$\mathcal{O}_X(U_{ij})$-加群となるとき、$\mathcal{F}$ は
{\it $S$ に相対的に有限表示される}、または
{\it $S$ に相対的有限表示をもつ}という。
\end{definition}

\noindent
この条件から、$\mathcal{F}$ は有限型 $\mathcal{O}_X$-加群となることに
注意する。$X \to S$ が単に局所有限型である場合、$X \to S$ が
局所有限表示でなくても、$\mathcal{F}$ は $S$ に相対的に有限表示されうる。
$X \to S$ が局所有限表示であることと、$\mathcal{O}_X$ が $S$ に相対的に
有限表示されることが同値であると後で分かる。

\begin{lemma}
\label{more-morphisms-lemma-relative-finite-presentation-characterize}
$f : X \to S$ を局所有限型なスキームの射とし、$\mathcal{F}$ を
準連接 $\mathcal{O}_X$-加群とする。次は同値である。
\begin{enumerate}
\item $\mathcal{F}$ は $S$ に相対的に有限表示される。
\item % P05-EN-CH37-0262
$f(U) \subset V$ を満たすすべてのアフィン開集合 $U \subset X$、
$V \subset S$ に対し、$\mathcal{O}_X(U)$-加群 $\mathcal{F}(U)$ は
$\mathcal{O}_S(V)$ に相対的に有限表示される。
\end{enumerate}
さらに、これらが成り立つならば、
% P05-EN-CH37-0263
$f(U) \subset V$ を満たすすべての開部分スキーム $U \subset X$ および
$V \subset S$ に対し、制限 $\mathcal{F}|_U$ は $V$ に相対的に
有限表示される。
\end{lemma}

\begin{proof}
最後の主張は (1) と (2) の同値性から明らかである。(2) が (1) を
含意することも明らかなので、(1) が成り立つと仮定する。
Definition \ref{more-morphisms-definition-relatively-finitely-presented-sheaf} にある
アフィン開被覆を $S = \bigcup V_i$ および
$f^{-1}(V_i) = \bigcup U_{ij}$ とする。(2) における
$U \subset X$ と $V \subset S$ を取る。More on Algebra, Lemma
\ref{more-algebra-lemma-glue-relative-finite-presentation} により、
$\mathcal{F}(U_k)$ が $\mathcal{O}_S(V)$ に相対的に有限表示されるような
$U$ の標準開被覆 $U = \bigcup U_k$ を見つければ十分である。言い換えると、
各 $u \in U$ に対し、$\mathcal{F}(U')$ が $\mathcal{O}_S(V)$ に相対的に
有限表示されるような標準アフィン開集合 $u \in U' \subset U$ を
見つければ十分である。$f(u) \in V_i$ となる $i$ を選び、次に
$u \in U_{ij}$ となる $j$ を選ぶ。Schemes, Lemma
\ref{schemes-lemma-standard-open-two-affines} により、
% P05-EN-CH37-0264
$V'$ と $V_i$ の双方における標準アフィン開集合となる
$v \in V' \subset V \cap V_i$ を取れる。このとき
$f^{-1}V'  \cap U$ および $f^{-1}V' \cap U_{ij}$ は、それぞれ $U$ および
$U_{ij}$ の標準アフィン開集合である。補題を再び適用すると、
$f^{-1}V'  \cap U$ と $f^{-1}V' \cap U_{ij}$ の双方において標準アフィン
開集合となる $u \in U' \subset f^{-1}V' \cap U \cap U_{ij}$ を取れる。
したがって $U'$ は $U$ と $U_{ij}$ の双方の標準アフィン開集合でもある。
More on Algebra, Lemma
\ref{more-algebra-lemma-localize-relative-finite-presentation} により、
$\mathcal{F}(U_{ij})$ が $\mathcal{O}_S(V_i)$ に相対的に有限表示される
という仮定から、$\mathcal{F}(U')$ も $\mathcal{O}_S(V_i)$ に相対的に
有限表示される。ここで
$\mathcal{O}_X(U') =
\mathcal{O}_X(U') \otimes_{\mathcal{O}_S(V_i)} \mathcal{O}_S(V')$
なので、More on Algebra, Lemma
\ref{more-algebra-lemma-base-change-relative-finite-presentation} から
$\mathcal{F}(U')$ は $\mathcal{O}_S(V')$ に相対的に有限表示される。
More on Algebra, Lemma
\ref{more-algebra-lemma-localize-relative-finite-presentation} をもう一度
適用すると、$\mathcal{F}(U')$ は $\mathcal{O}_S(V)$ に相対的に
有限表示される。これで証明が完了する。
\end{proof}

\begin{lemma}
\label{more-morphisms-lemma-relative-finite-presentation}
$f : X \to S$ を局所有限型なスキームの射とし、$\mathcal{F}$ を
準連接 $\mathcal{O}_X$-加群とする。
\begin{enumerate}
\item $f$ が局所有限表示ならば、$\mathcal{F}$ が $S$ に相対的に
有限表示されることと、$\mathcal{F}$ が有限表示されることは同値である。
\item 射 $f$ が局所有限表示であることと、$\mathcal{O}_X$ が $S$ に
相対的に有限表示されることは同値である。
\end{enumerate}
\end{lemma}

\begin{proof}
定義から直ちに従う。More on Algebra, Definition
\ref{more-algebra-definition-relatively-finitely-presented} の後の議論を
参照されたい。
\end{proof}

\begin{lemma}
\label{more-morphisms-lemma-finite-morphism-relative-finite-presentation}
$\pi : X \to Y$ を、基底スキーム $S$ 上局所有限型なスキームの有限射とする。
$\mathcal{F}$ を準連接 $\mathcal{O}_X$-加群とする。このとき、
$\mathcal{F}$ が $S$ に相対的に有限表示されることと、
$\pi_*\mathcal{F}$ が $S$ に相対的に有限表示されることは同値である。
\end{lemma}

\begin{proof}
More on Algebra, Lemma \ref{more-algebra-lemma-finite-extension} の結果を
スキームの言葉に移したものである。
\end{proof}

\begin{lemma}
\label{more-morphisms-lemma-base-change-relative-finite-presentation}
$f : X \to S$ を局所有限型なスキームの射とし、$\mathcal{F}$ を
準連接 $\mathcal{O}_X$-加群とする。$S' \to S$ をスキームの射とし、
$X' = X \times_S S'$ とおき、
% P05-EN-CH37-0265
$\mathcal{F}'$ で $\mathcal{F}$ の $X'$ への引き戻しを表す。
$\mathcal{F}$ が $S$ に相対的に有限表示されるならば、
$\mathcal{F}'$ は $S'$ に相対的に有限表示される。
\end{lemma}

\begin{proof}
More on Algebra, Lemma
\ref{more-algebra-lemma-base-change-relative-finite-presentation} の結果を
スキームの言葉に移したものである。
\end{proof}

\begin{lemma}
\label{more-morphisms-lemma-pull-relative-finite-presentation}
$X \to Y \to S$ を局所有限型なスキームの射とし、$\mathcal{G}$ を
準連接 $\mathcal{O}_Y$-加群とする。$f : X \to Y$ が局所有限表示で、
% P05-EN-CH37-0266
$\mathcal{G}$ が $S$ に相対的に有限表示されるならば、
$f^*\mathcal{G}$ は $S$ に相対的に有限表示される。
\end{lemma}

\begin{proof}
More on Algebra, Lemma
\ref{more-algebra-lemma-pull-relative-finite-presentation} の結果を
スキームの言葉に移したものである。
\end{proof}

\begin{lemma}
\label{more-morphisms-lemma-composition-relative-finite-presentation}
$X \to Y \to S$ を局所有限型なスキームの射とし、$\mathcal{F}$ を
準連接 $\mathcal{O}_X$-加群とする。$Y \to S$ が局所有限表示で、
$\mathcal{F}$ が $Y$ に相対的に有限表示されるならば、$\mathcal{F}$ は
$S$ に相対的に有限表示される。
\end{lemma}

\begin{proof}
More on Algebra, Lemma
\ref{more-algebra-lemma-composition-relative-finite-presentation} の結果を
スキームの言葉に移したものである。
\end{proof}

\begin{lemma}
\label{more-morphisms-lemma-ses-relatively-finite-presentation}
$X \to S$ を局所有限型なスキームの射とする。
$0 \to \mathcal{F}' \to \mathcal{F} \to \mathcal{F}'' \to 0$ を
準連接 $\mathcal{O}_X$-加群の短完全列とする。
\begin{enumerate}
\item $\mathcal{F}', \mathcal{F}''$ が $S$ に相対的に有限表示されるならば、
$\mathcal{F}$ もそうである。
\item $\mathcal{F}'$ が有限型 $\mathcal{O}_X$-加群で、$\mathcal{F}$ が
$S$ に相対的に有限表示されるならば、$\mathcal{F}''$ も $S$ に相対的に
有限表示される。
\end{enumerate}
\end{lemma}

\begin{proof}
More on Algebra, Lemma
\ref{more-algebra-lemma-ses-relatively-finite-presentation} の結果を
スキームの言葉に移したものである。
\end{proof}

\begin{lemma}
\label{more-morphisms-lemma-sum-relatively-finite-presentation}
\begin{slogan}
加群の直和因子は、基底に相対的に有限表示されるという性質を受け継ぐ。
\end{slogan}
$X \to S$ を局所有限型なスキームの射とし、
$\mathcal{F}, \mathcal{F}'$ を準連接 $\mathcal{O}_X$-加群とする。
$\mathcal{F} \oplus \mathcal{F}'$ が $S$ に相対的に有限表示されるならば、
$\mathcal{F}$ と $\mathcal{F}'$ もそうである。
\end{lemma}

\begin{proof}
More on Algebra, Lemma
\ref{more-algebra-lemma-sum-relatively-finite-presentation} の結果を
スキームの言葉に移したものである。
\end{proof}










\section{相対的擬連接性}
\label{more-morphisms-section-relative-pseudo-coherence}

\noindent
本節は More on Algebra, Section
\ref{more-algebra-section-relative-pseudo-coherent} のスキーム版である。
まずそちらの節に目を通すことを強く勧める。本節の内容と、Cohomology,
Sections \ref{cohomology-section-strictly-perfect},
\ref{cohomology-section-pseudo-coherent},
\ref{cohomology-section-tor}, および
\ref{cohomology-section-perfect} における擬連接複体の内容は、任意の
$\mathcal{O}_X$-加群の複体に対して展開してある。しかし $X$ が
スキームならば、$D_\QCoh(\mathcal{O}_X)$ の対象だけを扱うことで、
% P05-EN-CH37-0267
多くの補題と議論が大幅に簡単になり、当面の問題がしばしば直ちに
代数的な対応物へ帰着される。さらに、最初に示すことの一つは、
% P05-EN-CH37-0268
相対的擬連接性からコホモロジー層の準連接性が従うことである。
Lemma \ref{more-morphisms-lemma-relative-pseudo-coherence} を参照されたい。したがって、
初読時には $D_\QCoh(\mathcal{O}_X)$ の対象だけを扱うことを勧める。

\begin{lemma}
\label{more-morphisms-lemma-relatively-pseudo-coherent}
$X \to S$ をアフィン・スキームの有限型射とし、$E$ を
$D(\mathcal{O}_X)$ の対象とする。$m \in \mathbf{Z}$ とする。
次は同値である。
\begin{enumerate}
\item ある閉埋め込み $i : X \to \mathbf{A}^n_S$ に対し、
$D(\mathcal{O}_{\mathbf{A}^n_S})$ の対象 $Ri_*E$ は $m$-擬連接である。
\item すべての閉埋め込み $i : X \to \mathbf{A}^n_S$ に対し、
$D(\mathcal{O}_{\mathbf{A}^n_S})$ の対象 $Ri_*E$ は $m$-擬連接である。
\end{enumerate}
\end{lemma}

\begin{proof}
$S = \Spec(R)$、$X = \Spec(A)$ と書く。$i$ は全射
$\alpha : R[x_1, \ldots, x_n] \to A$ に対応するとし、
% P05-EN-CH37-0269
$X \to \mathbf{A}^m_S$ は $\beta : R[y_1, \ldots, y_m] \to A$ に
対応するとする。$\alpha(f_j) = \beta(y_j)$ を満たす
$f_j \in R[x_1, \ldots, x_n]$ と、$\beta(g_i) = \alpha(x_i)$ を満たす
$g_i \in R[y_1, \ldots, y_m]$ を選ぶ。すると可換図式
$$
\xymatrix{
R[x_1, \ldots, x_n, y_1, \ldots, y_m]
\ar[d]^{x_i \mapsto g_i} \ar[rr]_-{y_j \mapsto f_j} & &
R[x_1, \ldots, x_n] \ar[d] \\
R[y_1, \ldots, y_m] \ar[rr] & & A
}
$$
を得る。これは閉埋め込みの可換図式
$$
\xymatrix{
\mathbf{A}^{n + m}_S & \mathbf{A}^n_S \ar[l] \\
\mathbf{A}^m_S \ar[u] & X \ar[u] \ar[l]
}
$$
に対応する。したがって、閉埋め込み
$$
f : \mathbf{A}^m_S \to \mathbf{A}^{n + m}_S
$$
のもとで、$D(\mathcal{O}_{\mathbf{A}^m_S})$ の対象 $E$ が
$m$-擬連接であることと $Rf_*E$ が $m$-擬連接であることが同値だと
示せば十分である。これは Derived Categories of Schemes, Lemma
\ref{perfect-lemma-closed-push-pseudo-coherent} と、
$f_*\mathcal{O}_{\mathbf{A}^m_S}$ が擬連接
$\mathcal{O}_{\mathbf{A}^{n + m}_S}$-加群であることから従う。
$f_*\mathcal{O}_{\mathbf{A}^m_S}$ の擬連接性は直接証明することも容易だが、
Derived Categories of Schemes, Lemma
\ref{perfect-lemma-pseudo-coherent-affine} および More on Algebra, Lemma
\ref{more-algebra-lemma-relatively-pseudo-coherent} からも従う。
\end{proof}

\noindent
% P05-EN-CH37-0270
$f : X \to S$ が局所有限型なスキームの射ならば、
$f(U) \subset V$ を満たすアフィン開集合 $U \subset X$ と
$V \subset S$ の各対に対し、環準同型
$\mathcal{O}_S(V) \to \mathcal{O}_X(U)$ は有限型であることを思い出そう
（Morphisms, Lemma
\ref{morphisms-lemma-locally-finite-type-characterize}）。したがって
閉埋め込み $U \to \mathbf{A}^n_V$ は常に存在し、次の定義は意味をもつ。

\begin{definition}
\label{more-morphisms-definition-relative-pseudo-coherence}
$f : X \to S$ を局所有限型なスキームの射とし、$E$ を
$D(\mathcal{O}_X)$ の対象、$\mathcal{F}$ を $\mathcal{O}_X$-加群とする。
$m \in \mathbf{Z}$ を固定する。
\begin{enumerate}
\item アフィン開被覆 $S = \bigcup V_i$ が存在し、各 $i$ に対して
アフィン開被覆 $f^{-1}(V_i) = \bigcup U_{ij}$ が存在して、各対
$(U_{ij} \to V_i, E|_{U_{ij}})$ が Lemma
\ref{more-morphisms-lemma-relatively-pseudo-coherent} の同値な条件を満たすとき、
$E$ は{\it $S$ に相対的に $m$-擬連接}であるという。
\item すべての $m \in \mathbf{Z}$ に対して $E$ が $S$ に相対的に
$m$-擬連接であるとき、$E$ は{\it $S$ に相対的に擬連接}であるという。
\item $\mathcal{F}$ を $D(\mathcal{O}_X)$ の対象とみなしたものが
$S$ に相対的に $m$-擬連接であるとき、$\mathcal{F}$ は
{\it $S$ に相対的に $m$-擬連接}であるという。
\item $\mathcal{F}$ を $D(\mathcal{O}_X)$ の対象とみなしたものが
$S$ に相対的に擬連接であるとき、$\mathcal{F}$ は
{\it $S$ に相対的に擬連接}であるという。
\end{enumerate}
\end{definition}

\noindent
$X$ が準コンパクトで、ある $m$ に対して $E$ が $S$ に相対的に
$m$-擬連接ならば、$E$ は上に有界である。$E$ が $S$ に相対的に
擬連接ならば、$E$ のコホモロジー層は準連接である。

\begin{lemma}
\label{more-morphisms-lemma-relative-pseudo-coherence}
$f : X \to S$ を局所有限型なスキームの射とする。
$D(\mathcal{O}_X)$ の $E$ が $S$ に相対的に $m$-擬連接ならば、
$i > m$ に対して $H^i(E)$ は準連接 $\mathcal{O}_X$-加群である。
$E$ が $S$ に相対的に擬連接ならば、$E$ は
$D_\QCoh(\mathcal{O}_X)$ の対象である。
\end{lemma}

\begin{proof}
各対 $(U_{ij} \to V_i, E|_{U_{ij}})$ が Lemma
\ref{more-morphisms-lemma-relatively-pseudo-coherent} の同値な条件を満たすような
アフィン開被覆 $S = \bigcup V_i$ と、各 $i$ に対するアフィン開被覆
$f^{-1}(V_i) = \bigcup U_{ij}$ を選ぶ。準連接性は $X$ 上局所的なので、
% P05-EN-CH37-0271
$Ri_*E$ が $\mathbf{A}^n_S$ 上 $m$-擬連接となる閉埋め込み
$i : X \to \mathbf{A}^n_S$ が存在すると仮定してよい。
Derived Categories of Schemes, Lemma
\ref{perfect-lemma-pseudo-coherent} により、これは $q > m$ に対して
$H^q(Ri_*E)$ が準連接であることを意味する。$i_*$ は完全関手なので、
$i_*H^q(E) = H^q(Ri_*E)$ は $\mathbf{A}^n_S$ 上準連接である。
Morphisms, Lemma \ref{morphisms-lemma-i-star-equivalence} により、
望みどおり $H^q(E)$ は準連接である（厳密には、ある準連接
$\mathcal{O}_X$-加群 $\mathcal{F}$ で
$i_*\mathcal{F} = i_*H^q(E)$ を満たすものが存在することが従い、
Modules, Lemma \ref{modules-lemma-i-star-equivalence} から
$\mathcal{F} \cong H^q(E)$ が従うので、結論を得る）。
\end{proof}

\noindent
次に、相対的擬連接性の条件が局所化に関して良好に振る舞うことを示す。

\begin{lemma}
\label{more-morphisms-lemma-localize-relative-pseudo-coherent}
$S$ をアフィン・スキーム、$V \subset S$ を標準開集合とする。
$X \to V$ をアフィン・スキームの有限型射、$U \subset X$ を
アフィン開集合、$E$ を $D(\mathcal{O}_X)$ の対象とする。
対 $(X \to V, E)$ が Lemma
\ref{more-morphisms-lemma-relatively-pseudo-coherent} の同値な条件を満たすならば、
対 $(U \to S, E|_U)$ も Lemma
\ref{more-morphisms-lemma-relatively-pseudo-coherent} の同値な条件を満たす。
\end{lemma}

\begin{proof}
$S = \Spec(R)$、$V = D(f)$、$X = \Spec(A)$、$U = D(g)$ と書く。
対 $(X \to V, E)$ が Lemma
\ref{more-morphisms-lemma-relatively-pseudo-coherent} の同値な条件を満たすと仮定する。

\medskip\noindent
全射 $R_f[x_1, \ldots, x_n] \to A$ を選ぶ。
$R_f = R[x_0]/(fx_0 - 1)$ と書く。このとき
$R[x_0, x_1, \ldots, x_n] \to A$ は全射であり、
$R_f[x_1, \ldots, x_n]$ は $R[x_0, \ldots, x_n]$-加群として
擬連接である。したがって
$$
X \to \mathbf{A}^n_V \to \mathbf{A}^{n + 1}_S
$$
があり、Derived Categories of Schemes, Lemma
\ref{perfect-lemma-closed-push-pseudo-coherent} を適用すると、
$E$ の $\mathbf{A}^{n + 1}_S$ への直像 $E'$ は $m$-擬連接である。

\medskip\noindent
$g \in A$ へ写る元 $g' \in R[x_0, x_1, \ldots, x_n]$ を選び、全射
$R[x_0, \ldots, x_{n + 1}] \to R[x_0, \ldots, x_n, 1/g']$ を考える。
図式
$$
\xymatrix{
X \ar[d] & U \ar[d] \ar[l] \ar[dr] \\
\mathbf{A}^{n + 1}_S & D(g')\ar[l] \ar[r] & \mathbf{A}^{n + 2}_S
}
$$
を得る。ここで左下の矢印は開埋め込み、右下の矢印は閉埋め込みである。
前と同様に、$E'|_{D(g')}$ の $\mathbf{A}^{n + 2}_S$ への直像は
$m$-擬連接である。これは $E|_U$ の $\mathbf{A}^{n + 2}_S$ への
直像でもあるので、補題が成り立つ。
\end{proof}

\begin{lemma}
\label{more-morphisms-lemma-glue-relative-pseudo-coherent}
$X \to S$ をアフィン・スキームの有限型射とし、$E$ を
$D(\mathcal{O}_X)$ の対象、$m \in \mathbf{Z}$ とする。
$X = \bigcup U_i$ を標準アフィン開被覆とする。次は同値である。
\begin{enumerate}
\item 各対 $(U_i \to S, E|_{U_i})$ が Lemma
\ref{more-morphisms-lemma-relatively-pseudo-coherent} の同値な条件を満たす。
\item 対 $(X \to S, E)$ が Lemma
\ref{more-morphisms-lemma-relatively-pseudo-coherent} の同値な条件を満たす。
\end{enumerate}
\end{lemma}

\begin{proof}
(2) $\Rightarrow$ (1) は Lemma
\ref{more-morphisms-lemma-localize-relative-pseudo-coherent} である。(1) を仮定する。
$S = \Spec(R)$、$X = \Spec(A)$、$U_i = D(f_i)$ と書き、
$A$ において $1 = \sum f_ig_i$ と書く。全射
$$
R[x_i, y_i, z_i] \to R[x_i, y_i, z_i]/(\sum y_iz_i - 1) \to A.
$$
で $y_i$ を $f_i$ へ、$z_i$ を $g_i$ へ写すものを考える。
$R[x_i, y_i, z_i]/(\sum y_iz_i - 1)$ は
$R[x_i, y_i, z_i]$-加群として擬連接である。したがって Derived Categories
of Schemes, Lemma \ref{perfect-lemma-closed-push-pseudo-coherent} により、
$E$ の $T = \Spec(R[x_i, y_i, z_i]/(\sum y_iz_i - 1))$ への直像が
$m$-擬連接であることを示せば十分である。各 $i_0$ に対し、この直像の
$W_{i_0} = \Spec(R[x_i, y_i, z_i, 1/y_{i_0}]/(\sum y_iz_i - 1))$
への制限が $m$-擬連接であることを示せば十分である。可換図式
$$
\xymatrix{
X \ar[d] & U_{i_0} \ar[l] \ar[d] \\
T & W_{i_0} \ar[l]
}
$$
があり、これにより $E$ の $T$ への直像を $W_{i_0}$ に制限したものは、
$E|_{U_{i_0}}$ の $W_{i_0}$ への直像である。
$R[x_i, y_i, z_i, 1/y_{i_0}]/(\sum y_iz_i - 1)$ は $R$ 上の
多項式環と同型なので、望む結論が従う。
\end{proof}

\begin{lemma}
\label{more-morphisms-lemma-relative-pseudo-coherence-characterize}
$f : X \to S$ を局所有限型なスキームの射とし、$E$ を
$D(\mathcal{O}_X)$ の対象とする。$m \in \mathbf{Z}$ を固定する。
次は同値である。
\begin{enumerate}
\item $E$ は $S$ に相対的に $m$-擬連接である。
% P05-EN-CH37-0272
\item $f(U) \subset V$ を満たすすべてのアフィン開集合
$U \subset X$ と $V \subset S$ に対し、対 $(U \to V, E|_U)$ が
Lemma \ref{more-morphisms-lemma-relatively-pseudo-coherent} の同値な条件を満たす。
\end{enumerate}
% P05-EN-CH37-0273
さらに、これらが成り立つならば、$f(U) \subset V$ を満たすすべての
開部分スキーム $U \subset X$ と $V \subset S$ に対し、制限 $E|_U$ は
$V$ に相対的に $m$-擬連接である。
\end{lemma}

\begin{proof}
最後の主張は (1) と (2) の同値性から明らかである。(2) が (1) を
含意することも明らかである。(1) を仮定する。Definition
\ref{more-morphisms-definition-relative-pseudo-coherence} にあるようなアフィン開被覆
$S = \bigcup V_i$ と $f^{-1}(V_i) = \bigcup U_{ij}$ をとる。
$U \subset X$ と $V \subset S$ を (2) のとおりとする。Lemma
\ref{more-morphisms-lemma-glue-relative-pseudo-coherent} により、各対
$(U_k \to V, E|_{U_k})$ が Lemma
\ref{more-morphisms-lemma-relatively-pseudo-coherent} の同値な条件を満たすような
$U$ の標準開被覆 $U = \bigcup U_k$ を見つければ十分である。
言い換えると、各 $u \in U$ に対し、対 $(U' \to V, E|_{U'})$ が
Lemma \ref{more-morphisms-lemma-relatively-pseudo-coherent} の同値な条件を満たすような
標準アフィン開集合 $u \in U' \subset U$ を見つければ十分である。
$f(u) \in V_i$ となる $i$ を選び、次いで $u \in U_{ij}$ となる $j$ を
選ぶ。Schemes, Lemma \ref{schemes-lemma-standard-open-two-affines} により、
% P05-EN-CH37-0274
$V'$ と $V_i$ の双方で標準アフィン開となる
$v \in V' \subset V \cap V_i$ を見つけられる。すると
$f^{-1}V'  \cap U$、resp.\ $f^{-1}V' \cap U_{ij}$ はそれぞれ $U$、
resp.\ $U_{ij}$ の標準アフィン開集合である。もう一度この補題を適用すると、
$f^{-1}V'  \cap U$ と $f^{-1}V' \cap U_{ij}$ の双方で標準アフィン開となる
$u \in U' \subset f^{-1}V' \cap U \cap U_{ij}$ を見つけられる。
したがって $U'$ は $U$ と $U_{ij}$ の双方の標準アフィン開集合でもある。
対 $(U_{ij} \to V_i, E|_{U_{ij}})$ が Lemma
\ref{more-morphisms-lemma-relatively-pseudo-coherent} の同値な条件を満たすという仮定から、
Lemma \ref{more-morphisms-lemma-localize-relative-pseudo-coherent} により、対
$(U' \to V, E|_{U'})$ も Lemma
\ref{more-morphisms-lemma-relatively-pseudo-coherent} の同値な条件を満たす。
\end{proof}

\noindent
コホモロジー層が準連接である導来圏の対象については、相対的な
$m$-擬連接性を More on Algebra, Definition
\ref{more-algebra-definition-relatively-pseudo-coherent} で定義された概念と
関係づけることができる。アフィン・スキーム $U = \Spec(A)$ に対して関手
$R\Gamma(U, -)$ が $D_\QCoh(\mathcal{O}_U)$ と $D(A)$ の間の同値を誘導する
という事実を用いる。Derived Categories of Schemes, Lemma
\ref{perfect-lemma-affine-compare-bounded} を参照されたい。この関手は
引き戻しと両立する。すなわち、$E$ を $D_\QCoh(\mathcal{O}_U)$ の対象とし、
$A \to B$ をアフィン・スキームの射
$g : V = \Spec(B) \to \Spec(A) = U$ に対応する環準同型とすれば、
$R\Gamma(V, Lg^*E) = R\Gamma(U, E) \otimes_A^\mathbf{L} B$ である。
Derived Categories of Schemes, Lemma
\ref{perfect-lemma-quasi-coherence-pullback} を参照されたい。

\begin{lemma}
\label{more-morphisms-lemma-qcoh-relative-pseudo-coherence-characterize}
$f : X \to S$ を局所有限型なスキームの射とし、$E$ を
$D_\QCoh(\mathcal{O}_X)$ の対象とする。$m \in \mathbf{Z}$ を固定する。
次は同値である。
\begin{enumerate}
\item $E$ は $S$ に相対的に $m$-擬連接である。
\item アフィン開被覆 $S = \bigcup V_i$ と、各 $i$ に対するアフィン開被覆
$f^{-1}(V_i) = \bigcup U_{ij}$ が存在し、
$\mathcal{O}_X(U_{ij})$-加群の複体 $R\Gamma(U_{ij}, E)$ は
$\mathcal{O}_S(V_i)$ に相対的に $m$-擬連接である。
% P05-EN-CH37-0275
\item $f(U) \subset V$ を満たすすべてのアフィン開集合
$U \subset X$ と $V \subset S$ に対し、$\mathcal{O}_X(U)$-加群の複体
$R\Gamma(U, E)$ は $\mathcal{O}_S(V)$ に相対的に $m$-擬連接である。
\end{enumerate}
\end{lemma}

\begin{proof}
(2) のとおりの $U$ と $V$ をとり、閉埋め込み
$i : U \to \mathbf{A}^n_V$ を選ぶ。Lemma
\ref{more-morphisms-lemma-relative-pseudo-coherence-characterize} を用いる形式的な議論により、
$Ri_*(E|_U)$ が $m$-擬連接であることと $R\Gamma(U, E)$ が
$\mathcal{O}_S(V)$ に相対的に $m$-擬連接であることが同値だと示せば
十分である。$U = \Spec(A)$、$V = \Spec(R)$、そして
% P05-EN-CH37-0276
$\mathbf{A}^n_V = \Spec(R[x_1, \ldots, x_n]$ とする。
補題に先立つ注意により、$E|_U$ は $D(A)$ の対象
$R\Gamma(U, E)$ に付随する $U$ 上の準連接層の複体と準同型である。また、
$i$ は閉埋め込み（したがって $i_*$ は完全）なので、
$R\Gamma(U, E) = R\Gamma(\mathbf{A}^n_V, Ri_*(E|_U))$ である。したがって
% P05-EN-CH37-0277
$Ri_*E$ は、$D(R[x_1, \ldots, x_n])$ の対象とみなした
$R\Gamma(U, E)$ に付随する。Derived Categories of Schemes, Lemma
\ref{perfect-lemma-pseudo-coherent-affine} により、$Ri_*(E|_U)$ の
$m$-擬連接性は $D(R[x_1, \ldots, x_n])$ における $R\Gamma(U, E)$ の
$m$-擬連接性と同値であり、これは定義により $R\Gamma(U, E)$ が
$R = \mathcal{O}_S(V)$ に相対的に $m$-擬連接であることと同値なので、
結論を得る。
\end{proof}

\begin{lemma}
\label{more-morphisms-lemma-closed-morphism-relative-pseudo-coherence}
% P05-EN-CH37-0278
$i : X \to Y$ を基底スキーム $S$ 上局所有限型なスキームの射とする。
$i$ が $X$ から $Y$ の閉部分集合への同相を誘導すると仮定する。
$E$ を $D(\mathcal{O}_X)$ の対象とする。このとき、$E$ が $S$ に相対的に
$m$-擬連接であることと $Ri_*E$ が $S$ に相対的に $m$-擬連接であることは
同値である。
\end{lemma}

\begin{proof}
Morphisms, Lemma \ref{morphisms-lemma-homeomorphism-affine} により、射 $i$ は
アフィンである。したがって $S$、$Y$、$X$ はアフィンだと仮定してよい。
$S = \Spec(R)$、$Y = \Spec(A)$、$X = \Spec(B)$ とする。
% P05-EN-CH37-0279
この条件は $A/\text{rad}(A) \to B/\text{rad}(B)$ が全射であることを
意味する。ここで $\text{rad}(A)$ と $\text{rad}(B)$ は $A$ と $B$ の
Jacobson 根基を表す。$B$ は $A$ 上有限型なので、$B$ を $A$-代数として
生成する $b_1, \ldots, b_m \in \text{rad}(B)$ を選べる。
$b_j^N = 0$ がすべての $j$ に対して成り立つとする。環の図式
$$
\xymatrix{
B & R[x_i, y_j]/(y_j^N) \ar[l] & R[x_i, y_j] \ar[l] \\
A \ar[u] & R[x_i] \ar[l] \ar[u] \ar[ru]
}
$$
を考える。これはアフィン・スキームの図式
$$
\xymatrix{
X \ar[d] \ar[r] & T \ar[d] \ar[r] & \mathbf{A}^{n + m}_S \ar[ld] \\
Y \ar[r] & \mathbf{A}^n_S
}
$$
に移る。Lemma \ref{more-morphisms-lemma-relative-pseudo-coherence-characterize} により、
$E$ が $S$ に相対的に $m$-擬連接であることと、その
$\mathbf{A}^{n + m}_S$ への直像が $m$-擬連接であることは同値である。
Derived Categories of Schemes, Lemma
\ref{perfect-lemma-closed-push-pseudo-coherent} により、これはその $T$ への
直像が $m$-擬連接であることと同値である。同じ補題により、これはさらに
$\mathbf{A}^n_S$ への直像が $m$-擬連接であることと同値である。
再び Lemma \ref{more-morphisms-lemma-relative-pseudo-coherence-characterize} により、これは
$Ri_*E$ が $S$ に相対的に $m$-擬連接であることと同値である。
\end{proof}

\begin{lemma}
\label{more-morphisms-lemma-finite-morphism-relative-pseudo-coherence}
$\pi : X \to Y$ を基底スキーム $S$ 上局所有限型なスキームの有限射とし、
$E$ を $D_\QCoh(\mathcal{O}_X)$ の対象とする。このとき、$E$ が $S$ に
相対的に $m$-擬連接であることと、$R\pi_*E$ が $S$ に相対的に
$m$-擬連接であることは同値である。
\end{lemma}

\begin{proof}
More on Algebra, Lemma
\ref{more-algebra-lemma-finite-extension-pseudo-coherent} の結果を
スキームの言葉に移したものである。Derived Categories of Schemes, Lemma
\ref{perfect-lemma-quasi-coherence-direct-image} により、$R\pi_*$ は実際に
$D_\QCoh(\mathcal{O}_X)$ を $D_\QCoh(\mathcal{O}_Y)$ へ写すことに注意する。
移し替えには Lemma \ref{more-morphisms-lemma-relative-pseudo-coherence-characterize} を用いる。
\end{proof}

\begin{lemma}
\label{more-morphisms-lemma-cone-relatively-pseudo-coherent}
$f : X \to S$ を局所有限型なスキームの射とし、$(E, E', E'')$ を
$D(\mathcal{O}_X)$ の完全三角形とする。$m \in \mathbf{Z}$ とする。
\begin{enumerate}
\item $E$ が $S$ に相対的に $(m + 1)$-擬連接であり、$E'$ が $S$ に
相対的に $m$-擬連接ならば、$E''$ は $S$ に相対的に $m$-擬連接である。
\item $E, E''$ が $S$ に相対的に $m$-擬連接ならば、$E'$ は $S$ に
相対的に $m$-擬連接である。
\item $E'$ が $S$ に相対的に $(m + 1)$-擬連接であり、$E''$ が $S$ に
相対的に $m$-擬連接ならば、$E$ は $S$ に相対的に
$(m + 1)$-擬連接である。
\end{enumerate}
さらに、$E, E', E''$ のうち二つが $S$ に相対的に擬連接ならば、
% P05-EN-CH37-0280
残る一つもそうである。
\end{lemma}

\begin{proof}
Lemma \ref{more-morphisms-lemma-relative-pseudo-coherence-characterize} と Cohomology, Lemma
\ref{cohomology-lemma-cone-pseudo-coherent} から直ちに従う。
\end{proof}

\begin{lemma}
\label{more-morphisms-lemma-rel-n-pseudo-module}
$X \to S$ を局所有限型なスキームの射とし、$\mathcal{F}$ を
$\mathcal{O}_X$-加群とする。このとき
\begin{enumerate}
\item すべての $m > 0$ に対して $\mathcal{F}$ は $S$ に相対的に
$m$-擬連接である。
\item $\mathcal{F}$ が $S$ に相対的に $0$-擬連接であることと、
$\mathcal{F}$ が有限型 $\mathcal{O}_X$-加群であることは同値である。
\item $\mathcal{F}$ が $S$ に相対的に $(-1)$-擬連接であることと、
$\mathcal{F}$ が準連接であり $S$ に相対的に有限表示であることは同値である。
\end{enumerate}
\end{lemma}

\begin{proof}
(1) は定義から直ちに従う。(3) を示すため、$X$ 上局所的に議論してよい
（両方の性質は局所的である）。したがって $X$ と $S$ はアフィンだと
仮定してよい。閉埋め込み $i : X \to \mathbf{A}^n_S$ を選ぶ。
すると $\mathcal{F}$ が $S$ に相対的に $(-1)$-擬連接であることと、
$i_*\mathcal{F}$ が $(-1)$-擬連接であることは同値であり、さらにこれは
$i_*\mathcal{F}$ が有限表示 $\mathcal{O}_{\mathbf{A}^n_S}$-加群であることと
同値である。Cohomology, Lemma
\ref{cohomology-lemma-finite-cohomology} を参照されたい。有限表示加群は
準連接である。Modules, Lemma
\ref{modules-lemma-finite-presentation-quasi-coherent} を参照されたい。
Morphisms, Lemma \ref{morphisms-lemma-i-star-equivalence} により、
$\mathcal{F}$ が準連接であることと $i_*\mathcal{F}$ が準連接であることは
同値である。以上から (3) が従う。(2) の証明も同様だが、より簡単である。
\end{proof}

\begin{lemma}
\label{more-morphisms-lemma-summands-relative-pseudo-coherent}
$X \to S$ を局所有限型なスキームの射とし、$m \in \mathbf{Z}$ とする。
$E, K$ を $D(\mathcal{O}_X)$ の対象とする。$E \oplus K$ が $S$ に
相対的に $m$-擬連接ならば、$E$ と $K$ もそうである。
\end{lemma}

\begin{proof}
Cohomology, Lemma \ref{cohomology-lemma-summands-pseudo-coherent} と
定義から従う。
\end{proof}

\begin{lemma}
\label{more-morphisms-lemma-complex-relative-pseudo-coherent-modules}
$X \to S$ を局所有限型なスキームの射とし、$m \in \mathbf{Z}$ とする。
$\mathcal{F}^\bullet$ を（局所的に）上に有界な $\mathcal{O}_X$-加群の複体で、
すべての $i$ に対して $\mathcal{F}^i$ が $S$ に相対的に
$(m - i)$-擬連接であるものとする。このとき $\mathcal{F}^\bullet$ は
$S$ に相対的に $m$-擬連接である。
\end{lemma}

\begin{proof}
Cohomology, Lemma \ref{cohomology-lemma-complex-pseudo-coherent-modules} と
定義から従う。
\end{proof}

\begin{lemma}
\label{more-morphisms-lemma-cohomology-relative-pseudo-coherent}
$X \to S$ を局所有限型なスキームの射、$m \in \mathbf{Z}$ とし、$E$ を
$D(\mathcal{O}_X)$ の対象とする。$E$ が（局所的に）上に有界であり、
すべての $i$ に対して $H^i(E)$ が $S$ に相対的に
$(m - i)$-擬連接ならば、$E$ は $S$ に相対的に $m$-擬連接である。
\end{lemma}

\begin{proof}
Cohomology, Lemma \ref{cohomology-lemma-cohomology-pseudo-coherent} と
定義から従う。
\end{proof}

\begin{lemma}
\label{more-morphisms-lemma-base-change-relative-pseudo-coherent}
$X \to S$ を局所有限型なスキームの射、$m \in \mathbf{Z}$ とし、$E$ を
$S$ に相対的に $m$-擬連接な $D(\mathcal{O}_X)$ の対象とする。
$S' \to S$ をスキームの射とする。$X' = X \times_S S'$ とおき、
% P05-EN-CH37-0281
$E'$ で $E$ の $X'$ への導来引き戻しを表す。$S'$ と $X$ が $S$ 上
Tor 独立ならば、$E'$ は $S'$ に相対的に $m$-擬連接である。
\end{lemma}

\begin{proof}
問題は $X$ と $X'$ 上局所的なので、$X$、$S$、$S'$、$X'$ は
アフィンだと仮定してよい。閉埋め込み $i : X \to \mathbf{A}^n_S$ を選び、
% P05-EN-CH37-0282
$i' : X' \to \mathbf{A}^n_{S'}$ でその $S'$ への基底変換を表す。
% P05-EN-CH37-0283
$g : X' \to X$ と $g' : \mathbf{A}^n_{S'} \to \mathbf{A}^n_S$ で射影を
表すと、$E' = Lg^*E$ である。$X$ と $S'$ は $S$ 上 Tor 独立なので、
基底変換写像（Cohomology, Remark
\ref{cohomology-remark-base-change}）は同型
$$
Ri'_*(Lg^*E) = L(g')^*Ri_*E
$$
を誘導する。実際、$x \in X$ 上にある点 $x' \in X'$ に対し、$x'$ の
茎における基底変換写像は、閉埋め込み $i$ と $i'$ から生じる写像
$$
E_x \otimes_{\mathcal{O}_{\mathbf{A}^n_S, x}}^\mathbf{L}
\mathcal{O}_{\mathbf{A}^n_{S'}, x'}
\longrightarrow
E_x \otimes_{\mathcal{O}_{X, x}}^\mathbf{L}
\mathcal{O}_{X', x'}
$$
である。始域は
$E_x \otimes_{\mathcal{O}_{S, s}}^\mathbf{L} \mathcal{O}_{S', s'}$ の
局所化と準同型である。一方、仮定により $\mathcal{O}_{X', x'}$ は
$\mathcal{O}_{X, x} \otimes_{\mathcal{O}_{S, s}}^\mathbf{L}
\mathcal{O}_{S', s'}$ の（同じ）局所化と同型なので、これは終域と同型である。
Cohomology, Lemma \ref{cohomology-lemma-pseudo-coherent-pullback} を適用して
補題を得る。
\end{proof}

\begin{lemma}
\label{more-morphisms-lemma-pull-relative-pseudo-coherent}
$f : X \to Y$ を基底 $S$ 上局所有限型なスキームの射とする。
$m \in \mathbf{Z}$ とし、$E$ を $D(\mathcal{O}_Y)$ の対象とする。次を仮定する。
\begin{enumerate}
\item $\mathcal{O}_X$ は $Y$ に相対的に擬連接である\footnote{これは
$f$ が擬連接であることを意味する。Definition
\ref{more-morphisms-definition-pseudo-coherent} を参照されたい。}。
\item $E$ は $S$ に相対的に $m$-擬連接である。
\end{enumerate}
このとき $Lf^*E$ は $S$ に相対的に $m$-擬連接である。
\end{lemma}

\begin{proof}
問題は $X$ 上局所的である。したがって $X$、$Y$、$S$ はアフィンだと
仮定してよい。More on Algebra, Lemma
\ref{more-algebra-lemma-pull-relative-pseudo-coherent} の証明と同様に議論すると、
可換図式
$$
\xymatrix{
X \ar[r]_i \ar[d]_f &
\mathbf{A}^d_Y \ar[r]_j \ar[ld]^p &
\mathbf{A}^{n + d}_S \ar[ld] \\
Y \ar[r] &
\mathbf{A}^n_S
}
$$
を得る。次に注意する。
$$
Ri_* Lf^*E = Ri_* Li^* Lp^*E =
% P05-EN-CH37-0284
Lp^*E \otimes_{\mathcal{O}_{\mathbf{A}_Y^n}}^\mathbf{L} Ri_*\mathcal{O}_X
$$
Cohomology, Lemma
\ref{cohomology-lemma-projection-formula-closed-immersion} による。
仮定と $Y$ がアフィンであることにより、
% P05-EN-CH37-0285
$Ri_*\mathcal{O}_X = i_*\mathcal{O}_X$ を、$q > 0$ に対して
$\mathcal{F}^q = 0$ を満たす有限自由
$\mathcal{O}_{\mathbf{A}_Y^n}$-加群の複体 $\mathcal{F}^\bullet$ で表せる
（詳細は省略する。Derived Categories of Schemes, Lemma
\ref{perfect-lemma-pseudo-coherent-affine}
および More on Algebra, Lemma
\ref{more-algebra-lemma-rel-n-pseudo-module} を用いる）。仮定により $E$ は
上に有界であり、$q > a$ に対して $H^q(E) = 0$ であるとする。
$q > a$ に対して $\mathcal{E}^q = 0$ を満たす
$\mathcal{O}_Y$-加群の複体 $\mathcal{E}^\bullet$ で $E$ を表す。
すると上の導来テンソル積は
$\text{Tot}(p^*\mathcal{E}^\bullet
\otimes_{\mathcal{O}_{\mathbf{A}_Y^n}} \mathcal{F}^\bullet)$ で表される。

\medskip\noindent
$j$ は閉埋め込みなので、関手 $j_*$ は完全であり、$Rj_*$ は任意の代表層複体に
$j_*$ を適用して計算される。したがって
% P05-EN-CH37-0284
$j_*\text{Tot}(p^*\mathcal{E}^\bullet \otimes_{\mathcal{O}_{\mathbf{A}_Y^n}}
\mathcal{F}^\bullet)$ が $\mathcal{O}_{\mathbf{A}^{n + m}_S}$-加群の複体として
$m$-擬連接であることを示さなければならない。次に、
% P05-EN-CH37-0284
$\text{Tot}(p^*\mathcal{E}^\bullet \otimes_{\mathcal{O}_{\mathbf{A}_Y^n}}
\mathcal{F}^\bullet)$ は部分複体によるフィルトレーションをもち、その逐次商は
複体
% P05-EN-CH37-0284
$p^*\mathcal{E}^\bullet
\otimes_{\mathcal{O}_{\mathbf{A}_Y^n}} \mathcal{F}^q[-q]$.
である。$q \ll 0$ に対して複体
% P05-EN-CH37-0284
$p^*\mathcal{E}^\bullet \otimes_{\mathcal{O}_{\mathbf{A}_Y^n}}
\mathcal{F}^q[-q]$
は次数 $\leq m$ でコホモロジーが零であり、したがって $m$-擬連接である。
ゆえに Lemma \ref{more-morphisms-lemma-cone-relatively-pseudo-coherent} と帰納法を適用すると、
すべての $q$ に対して
% P05-EN-CH37-0284
$p^*\mathcal{E}^\bullet \otimes_{\mathcal{O}_{\mathbf{A}_Y^n}}
\mathcal{F}^q[-q]$
が $S$ に相対的に擬連接であることを示せば十分である。
% P05-EN-CH37-0286
$q > 0$ に対して $\mathcal{F}^q = 0$ であることに注意する。
さらに $\mathcal{F}^q$ は有限自由なので、これは
$p^*\mathcal{E}^\bullet$ が $S$ に相対的に $m$-擬連接であることの証明に
帰着する。これは例えば Lemma
\ref{more-morphisms-lemma-base-change-relative-pseudo-coherent} から従う。
\end{proof}

\begin{lemma}
\label{more-morphisms-lemma-composition-relative-pseudo-coherent}
$f : X \to Y$ を基底 $S$ 上局所有限型なスキームの射とする。
$m \in \mathbf{Z}$ とし、$E$ を $D(\mathcal{O}_X)$ の対象とする。
$\mathcal{O}_Y$ が $S$ に相対的に擬連接であると仮定する\footnote{これは
$Y \to S$ が擬連接であることを意味する。Definition
\ref{more-morphisms-definition-pseudo-coherent} を参照されたい。}。次は同値である。
\begin{enumerate}
\item $E$ は $Y$ に相対的に $m$-擬連接である。
\item $E$ は $S$ に相対的に $m$-擬連接である。
\end{enumerate}
\end{lemma}

\begin{proof}
問題は $X$ 上局所的なので、$X$、$Y$、$S$ はアフィンだと仮定してよい。
More on Algebra, Lemma
\ref{more-algebra-lemma-pull-relative-pseudo-coherent} の証明と同様に議論すると、
% P05-EN-CH37-0287
可換図式
$$
\xymatrix{
X \ar[r]_i \ar[d]_f &
\mathbf{A}^m_Y \ar[r]_j \ar[ld]^p &
\mathbf{A}^{n + m}_S \ar[ld] \\
Y \ar[r] &
\mathbf{A}^n_S
}
$$
を得る。$\mathcal{O}_Y$ が $S$ に相対的に擬連接であるという仮定は、
$\mathcal{O}_{\mathbf{A}^m_Y}$ が $\mathbf{A}^m_S$ に相対的に擬連接であることを
含意する（平坦基底変換による。これは例えば Lemma
\ref{more-morphisms-lemma-base-change-relative-pseudo-coherent} を用いて分かる）。
これはさらに、
% P05-EN-CH37-0288
$j_*\mathcal{O}_{\mathbf{A}^n_Y}$ が
$\mathcal{O}_{\mathbf{A}^{n + m}_S}$-加群として擬連接であることを含意する。
すると補題の同値性は Derived Categories of Schemes, Lemma
\ref{perfect-lemma-closed-push-pseudo-coherent} から従う。
\end{proof}

\begin{lemma}
\label{more-morphisms-lemma-check-relative-pseudo-coherence-on-charts}
次の可換図式を考える。
$$
\xymatrix{
X \ar[rd] \ar[rr]_i & & P \ar[ld] \\
& S
}
$$
% P05-EN-CH37-0289
$i$ は閉埋め込みであり、$P \to S$ は平坦かつ局所有限表示であると仮定する。
$E$ を $D(\mathcal{O}_X)$ の対象とする。このとき次は同値である。
\begin{enumerate}
\item $E$ は $S$ に相対的に $m$-擬連接である。
\item $Ri_*E$ は $S$ に相対的に $m$-擬連接である。
\item $Ri_*E$ は $P$ 上 $m$-擬連接である。
\end{enumerate}
\end{lemma}

\begin{proof}
(1) と (2) の同値性は Lemma
\ref{more-morphisms-lemma-finite-morphism-relative-pseudo-coherence} である。
$\mathcal{O}_P$ が $S$ に相対的に擬連接であることを示せれば、
(2) と (3) の同値性は $\text{id} : P \to P$ に Lemma
\ref{more-morphisms-lemma-composition-relative-pseudo-coherent} を適用して従う。
必要な擬連接性は More on Algebra, Lemma
\ref{more-algebra-lemma-flat-finite-presentation-perfect} と定義から従う。
\end{proof}









\section{擬連接射}
\label{more-morphisms-section-pseudo-coherent}

\noindent
% P05-EN-CH37-0290
本節は何としても読まないようにされたい。この内容の一部が必要ならば、
まず対応する代数の各節に目を通されたい。すなわち
More on Algebra, Sections \ref{more-algebra-section-pseudo-coherent},
\ref{more-algebra-section-relative-pseudo-coherent}, および
\ref{more-algebra-section-pseudo-coherent-perfect-ring-map} である。
当面知っておくべきことは、環準同型 $A \to B$ が擬連接であることと、
$B = A[x_1, \ldots, x_n]/I$ であり、$A[x_1, \ldots, x_n]$-加群としての
$B$ が有限自由 $A[x_1, \ldots, x_n]$-加群による分解をもつことが
同値だということだけである。

\begin{lemma}
\label{more-morphisms-lemma-pseudo-coherent}
$f : X \to S$ をスキームの射とする。次は同値である。
\begin{enumerate}
\item アフィン開被覆 $S = \bigcup V_j$ と、各 $j$ に対するアフィン開被覆
$f^{-1}(V_j) = \bigcup U_{ji}$ が存在して、
% P05-EN-CH37-0291
$\mathcal{O}_S(V_j) \to \mathcal{O}_X(U_{ij})$ は擬連接環準同型である。
\item $f(U) \subset V$ を満たすすべてのアフィン開集合の対
$U \subset X$、$V \subset S$ に対し、環準同型
$\mathcal{O}_S(V) \to \mathcal{O}_X(U)$ は擬連接である。
\item $f$ は局所有限型であり、$\mathcal{O}_X$ は $S$ に相対的に
擬連接である。
\end{enumerate}
\end{lemma}

\begin{proof}
(1) と (2) の同値性を示すには、擬連接環準同型であるという性質について
Morphisms, Definition \ref{morphisms-definition-property-local} の条件
(1)(a), (b), (c) を確認すれば十分である。これらの性質は、局所化が
平坦であることを用いて More on Algebra, Lemmas
\ref{more-algebra-lemma-base-change-relative-pseudo-coherent},
\ref{more-algebra-lemma-localize-relative-pseudo-coherent}, および
\ref{more-algebra-lemma-glue-relative-pseudo-coherent} から従う。

\medskip\noindent
(1) が成り立つならば、擬連接環準同型は定義により有限型なので、$f$ は
局所有限型である。さらに Lemma
\ref{more-morphisms-lemma-qcoh-relative-pseudo-coherence-characterize} と定義により、
(1) から $\mathcal{O}_X$ が $S$ に相対的に擬連接であることが従う。
逆に (3) が成り立つならば、(2) のとおりのすべての $U$ と $V$ に対し、
環 $\mathcal{O}_X(U)$ は $\mathcal{O}_S(V)$ 上有限型であり、
$\mathcal{O}_X(U)$ は加群として $\mathcal{O}_S(V)$ に相対的に擬連接である。Lemmas
\ref{more-morphisms-lemma-relative-pseudo-coherence-characterize} および
\ref{more-morphisms-lemma-qcoh-relative-pseudo-coherence-characterize} を参照されたい。
これは擬連接環準同型の定義であるから、(2) と (1) が成り立つ。
\end{proof}

\begin{definition}
\label{more-morphisms-definition-pseudo-coherent}
スキームの射 $f : X \to S$ が Lemma \ref{more-morphisms-lemma-pseudo-coherent} の
同値な条件を満たすとき、この射は{\it 擬連接}であるという。この場合、
$X$ は $S$ 上擬連接であるともいう。
\end{definition}

\noindent
擬連接射の基底変換は一般には擬連接でないことに注意されたい。

\begin{lemma}
\label{more-morphisms-lemma-flat-base-change-pseudo-coherent}
擬連接射の平坦基底変換は擬連接である。
\end{lemma}

\begin{proof}
これは次の代数的結果に移る。$A \to B$ を擬連接環準同型とし、
$A \to A'$ を平坦とする。このとき $A' \to B \otimes_A A'$ は
擬連接である。これは、より一般的な More on Algebra, Lemma
\ref{more-algebra-lemma-base-change-relative-pseudo-coherent} から従う。
\end{proof}

\begin{lemma}
\label{more-morphisms-lemma-composition-pseudo-coherent}
スキームの擬連接射の合成は擬連接である。
\end{lemma}

\begin{proof}
これは次の代数的結果に移る。$A \to B \to C$ が合成可能な擬連接環準同型ならば、
$A \to C$ は擬連接である。これは More on Algebra, Lemma
\ref{more-algebra-lemma-pull-relative-pseudo-coherent} または More on Algebra,
Lemma \ref{more-algebra-lemma-composition-relative-pseudo-coherent} から従う。
\end{proof}

\begin{lemma}
\label{more-morphisms-lemma-pseudo-coherent-finite-presentation}
擬連接射は局所有限表示である。
\end{lemma}

\begin{proof}
定義から直ちに従う。
\end{proof}

\begin{lemma}
\label{more-morphisms-lemma-flat-finite-presentation-pseudo-coherent}
平坦かつ局所有限表示な射は擬連接である。
\end{lemma}

\begin{proof}
有限表示な平坦環準同型は擬連接（実際には完全）であるという事実から従う。
More on Algebra, Lemma
\ref{more-algebra-lemma-flat-finite-presentation-perfect} を参照されたい。
\end{proof}

\begin{lemma}
\label{more-morphisms-lemma-permanence-pseudo-coherent}
$f : X \to Y$ を、基底スキーム $S$ 上擬連接なスキームの射とする。
このとき $f$ は擬連接である。
\end{lemma}

\begin{proof}
これは次の代数的結果に移る。$R \to A \to B$ が合成可能な環準同型で、
$R \to A$ と $R \to B$ が擬連接ならば、
% P05-EN-CH37-0292
$R \to B$ は擬連接である。これは More on Algebra, Lemma
\ref{more-algebra-lemma-composition-relative-pseudo-coherent} から従う。
\end{proof}

\begin{lemma}
\label{more-morphisms-lemma-finite-pseudo-coherent}
$f : X \to S$ をスキームの有限射とする。このとき $f$ が擬連接であることと、
$f_*\mathcal{O}_X$ が $\mathcal{O}_S$-加群として擬連接であることは同値である。
\end{lemma}

\begin{proof}
代数に移すと、この補題は次を述べる。$R \to A$ が有限環準同型ならば、
$R \to A$ が環準同型として擬連接（定義により、これは $A$-加群としての
$A$ が $R$ に相対的に擬連接であることを意味する）であることと、$A$ が
$R$-加群として擬連接であることは同値である。これは、より一般的な
More on Algebra, Lemma
\ref{more-algebra-lemma-finite-extension-pseudo-coherent} から従う。
\end{proof}

\begin{lemma}
\label{more-morphisms-lemma-Noetherian-pseudo-coherent}
$f : X \to S$ をスキームの射とする。$S$ が局所 Noetherian ならば、
$f$ が擬連接であることと $f$ が局所有限型であることは同値である。
\end{lemma}

\begin{proof}
これは次の代数的結果に移る。
% P05-EN-CH37-0293
$R$ を Noetherian 環とし、$R \to A$ を有限型環準同型とすると、
$R \to A$ が擬連接であることと $R \to A$ が有限型であることは同値である。
これを見るため、擬連接環準同型は定義により有限型であることに注意する。
逆に $R \to A$ が有限型ならば $A = R[x_1, \ldots, x_n]/I$ と書け、
More on Algebra, Lemma
\ref{more-algebra-lemma-Noetherian-pseudo-coherent} により $A$ は
$R[x_1, \ldots, x_n]$-加群として擬連接である。すなわち $R \to A$ は
擬連接環準同型である。
\end{proof}

\begin{lemma}
\label{more-morphisms-lemma-descending-property-pseudo-coherent}
性質 $\mathcal{P}(f) =$``$f$ は擬連接である'' は基底上 fpqc 局所的である。
\end{lemma}

\begin{proof}
これを示すため Descent, Lemma
\ref{descent-lemma-descending-properties-morphisms} の判定条件を用いる。
Definition \ref{more-morphisms-definition-pseudo-coherent} により、擬連接性は基底上
Zariski 局所的である。Lemma \ref{more-morphisms-lemma-flat-base-change-pseudo-coherent} により、
擬連接性は平坦基底変換で保たれる。Descent, Lemma
\ref{descent-lemma-descending-properties-morphisms} の最後の仮定 (3) は、
次の代数的主張に移る。$A \to B$ を忠実平坦環準同型とし、
$C = A[x_1, \ldots, x_n]/I$ を $A$-代数とする。
% P05-EN-CH37-0294
$C \otimes_A B$ が $B[x_1, \ldots, x_n]$-加群として擬連接ならば、
% P05-EN-CH37-0295
$C$ は $A[x_1, \ldots, x_n]$-加群として擬連接である。これは
More on Algebra, Lemma \ref{more-algebra-lemma-flat-descent-pseudo-coherent}
である。
\end{proof}

\begin{lemma}
\label{more-morphisms-lemma-quotient-of-flat-finitely-presented}
$A \to B$ を有限表示な平坦環準同型とし、$I \subset B$ をイデアルとする。
このとき $A \to B/I$ が擬連接であることと、$I$ が $B$-加群として
擬連接であることは同値である。
\end{lemma}

\begin{proof}
表示 $B = A[x_1, \ldots, x_n]/J$ を選ぶ。Lemma
\ref{more-morphisms-lemma-flat-finite-presentation-pseudo-coherent} により $A \to B$ は
擬連接環準同型なので、$B$ は $A[x_1, \ldots, x_n]$-加群として擬連接である。
$A \to B/I$ が擬連接であることと $B/I$ が
$A[x_1, \ldots, x_n]$-加群として擬連接であることは同値である。
More on Algebra, Lemma \ref{more-algebra-lemma-finite-push-pseudo-coherent}
により、これは $B/I$ が $B$-加群として擬連接であることと同値である。
% P05-EN-CH37-0296
短完全列 $0 \to I \to B \to B/I \to 0$ により、$I$ が擬連接であることと
$B/I$ が擬連接であることは同値なので、補題が従う（More on Algebra,
Lemma \ref{more-algebra-lemma-two-out-of-three-pseudo-coherent} を参照されたい）。
\end{proof}

\noindent
次の補題は、より強い Lemma
\ref{more-morphisms-lemma-pseudo-coherent-fppf-local-source} によって不要になる。

\begin{lemma}
\label{more-morphisms-lemma-pseudo-coherent-syntomic-local-source}
性質 $\mathcal{P}(f) =$``$f$ は擬連接である'' は始域上 syntomic 局所的である。
\end{lemma}

\begin{proof}
これを示すため Descent, Lemma
\ref{descent-lemma-properties-morphisms-local-source} の判定条件を用いる。
Lemmas \ref{more-morphisms-lemma-flat-finite-presentation-pseudo-coherent} および
\ref{more-morphisms-lemma-composition-pseudo-coherent} により、擬連接性は局所有限表示な平坦射を
前合成しても保たれ、特に syntomic 射を前合成しても保たれる（Morphisms,
Lemmas \ref{morphisms-lemma-syntomic-flat} および
\ref{morphisms-lemma-syntomic-locally-finite-presentation} を参照されたい）。
Definition \ref{more-morphisms-definition-pseudo-coherent} から、擬連接性が始域と終域上
Zariski 局所的であることは明らかである。したがって、前述の Descent,
Lemma \ref{descent-lemma-properties-morphisms-local-source} により、次を示せば
十分である。$X' \to X \to Y$ をアフィン・スキームの射とし、
% P05-EN-CH37-0297
$X' \to X$ は syntomic、$X' \to Y$ は擬連接であると仮定する。このとき
$X \to Y$ は擬連接である。実際、Descent, Lemma
\ref{descent-lemma-flat-finitely-presented-permanence-algebra} により、いずれにせよ
$X \to Y$ は有限表示である。閉埋め込み $X \to \mathbf{A}^n_Y$ を選ぶ。
Algebra, Lemma \ref{algebra-lemma-lift-syntomic} により、
% P05-EN-CH37-0298
アフィン開被覆 $X' = \bigcup_{i = 1, \ldots, n} X'_i$ と、射
$X'_i \to X$ を持ち上げる syntomic 射 $W_i \to \mathbf{A}^n_Y$ を選べる。
すなわちファイバー積図式
$$
\xymatrix{
X'_i \ar[d] \ar[r] & W_i \ar[d] \\
X \ar[r] & \mathbf{A}^n_Y
}
$$
がある。$X'$ を $\coprod X'_i$ で置き換え、$W = \coprod W_i$ とおくと、
ファイバー積図式
$$
\xymatrix{
X' \ar[d] \ar[r] & W \ar[d]^h \\
X \ar[r] & \mathbf{A}^n_Y
}
$$
を得る。ここで $W \to \mathbf{A}^n_Y$ は平坦かつ有限表示であり、
$X' \to Y$ は依然として擬連接である。$W \to \mathbf{A}^n_Y$ は開写像であり
（Morphisms, Lemma \ref{morphisms-lemma-fppf-open}）、$X' \to X$ は全射なので、
$X \subset D(f) \subset \Im(h)$ を満たす
$f \in \Gamma(\mathbf{A}^n_Y, \mathcal{O})$ を選べる。
$Y = \Spec(R)$、$X = \Spec(A)$、$X' = \Spec(A')$、$W = \Spec(B)$、
$A = R[x_1, \ldots, x_n]/I$、$A' = B/IB$ と書く。このとき
$R \to A'$ は擬連接である。図式は
$$
\xymatrix{
A' = B/IB & B \ar[l] \\
A = R[x_1, \ldots, x_n]/I \ar[u] & R[x_1, \ldots, x_n] \ar[l] \ar[u]
}
$$
である。Lemma \ref{more-morphisms-lemma-quotient-of-flat-finitely-presented} により、$IB$ は
$B$-加群として擬連接である。上で選んだ $f$ により、環準同型
$R[x_1, \ldots, x_n]_f \to B_f$ は忠実平坦である。したがって
$I_f \subset R[x_1, \ldots, x_n]_f$ は擬連接である。More on Algebra,
Lemma \ref{more-algebra-lemma-flat-descent-pseudo-coherent} を参照されたい。
もう一度 Lemma \ref{more-morphisms-lemma-quotient-of-flat-finitely-presented} を適用すると、
$R \to A$ が擬連接であることが分かる。
\end{proof}

\begin{lemma}
\label{more-morphisms-lemma-pseudo-coherent-fppf-local-source}
性質 $\mathcal{P}(f) =$``$f$ は擬連接である'' は始域上 fppf 局所的である。
\end{lemma}

\begin{proof}
$f : X \to S$ をスキームの射とする。$\{g_i : X_i \to X\}$ を、各合成
$f \circ g_i$ が擬連接となる fppf 被覆とする。Lemma
\ref{more-morphisms-lemma-dominate-fppf} により、次のものが存在する。
\begin{enumerate}
\item Zariski 開被覆 $X = \bigcup U_j$。
\item 全射な有限局所自由射 $W_j \to U_j$。
\item Zariski 開被覆 $W_j = \bigcup_k W_{j, k}$。
\item 全射な有限局所自由射 $T_{j, k} \to W_{j, k}$。
\end{enumerate}
これらは、fppf 被覆 $\{h_{j, k} : T_{j, k} \to X\}$ が与えられた被覆
$\{X_i \to X\}$ を細分するように選べる。
% P05-EN-CH37-0299
$\psi_{j, k} : T_{j, k} \to X_{\alpha(j, k)}$ で、
$\{T_{j, k} \to X\}$ が与えられた被覆 $\{X_i \to X\}$ を細分することを
示す射を表す。
$T_{j, k} \to X$ は平坦かつ局所有限表示なので、Lemma
\ref{more-morphisms-lemma-flat-finite-presentation-pseudo-coherent} により $X_i$ と
$T_{j, k}$ はともに $X$ 上擬連接である。したがって
$\psi_{j, k} : T_{j, k} \to X_i$ は擬連接である。Lemma
\ref{more-morphisms-lemma-permanence-pseudo-coherent} を参照されたい。ゆえに
% P05-EN-CH37-0300
$T_{j, k} \to S$ は $\psi_{j, k}$ と $f \circ g_{\alpha(j, k)}$ の合成として
擬連接である。Lemma \ref{more-morphisms-lemma-composition-pseudo-coherent} を参照されたい。
以上により、補題は Zariski 開被覆の場合（これはよい）と、単一の全射な
有限局所自由射による被覆の場合とに帰着された。後者を次の段落で扱う。

\medskip\noindent
$X' \to X \to S$ をスキームの射の列とし、$X' \to X$ は全射な有限局所自由射、
% P05-EN-CH37-0301
$X' \to Y$ は擬連接であると仮定する。目標は $X \to S$ が擬連接であることを
示すことである。Descent, Lemma
\ref{descent-lemma-flat-finitely-presented-permanence} により、射 $X \to S$ は
局所有限表示である。問題が、$X'$、$X$、$S$ がアフィンであり、
$X' \to X$ がある階数 $r > 0$ の自由射である場合に帰着することは明らかである。
対応する代数的問題は次のとおりである。$R \to A \to A'$ を環準同型とし、
$R \to A'$ は擬連接、$R \to A$ は有限表示、さらに $A$-加群として
$A' \cong A^{\oplus r}$ であると仮定する。目標は $R \to A$ が擬連接である
ことを示すことである。$R \to A'$ が擬連接であるという仮定は、$A'$-加群
としての $A'$ が $R$ に相対的に擬連接であることを意味する。More on Algebra,
Lemma \ref{more-algebra-lemma-finite-extension-pseudo-coherent} により、これは
$A$-加群としての $A'$ が $R$ に相対的に擬連接であることを含意する。
$A$-加群として $A' \cong A^{\oplus r}$ なので、$A$-加群としての $A$ は
$R$ に相対的に擬連接である。More on Algebra, Lemma
\ref{more-algebra-lemma-summands-relative-pseudo-coherent} を参照されたい。
これは定義により $R \to A$ が擬連接であることを意味し、結論を得る。
\end{proof}








\section{完全射}
\label{more-morphisms-section-perfect}

\noindent
本節の内容を理解するには、擬連接射に関する節の内容を少しだけ理解しておく
必要がある。当面知っておくべきことは、環準同型 $A \to B$ が完全であることと、
それが擬連接であり、かつ $B$ が $A$-加群として有限 Tor 次元をもつことが
同値だということだけである。

\begin{lemma}
\label{more-morphisms-lemma-perfect}
$f : X \to S$ を局所有限型なスキームの射とする。次は同値である。
\begin{enumerate}
\item アフィン開被覆 $S = \bigcup V_j$ と、各 $j$ に対するアフィン開被覆
$f^{-1}(V_j) = \bigcup U_{ji}$ が存在して、
% P05-EN-CH37-0302
$\mathcal{O}_S(V_j) \to \mathcal{O}_X(U_{ij})$ は完全環準同型である。
\item $f(U) \subset V$ を満たすすべてのアフィン開集合の対
$U \subset X$、$V \subset S$ に対し、環準同型
$\mathcal{O}_S(V) \to \mathcal{O}_X(U)$ は完全である。
\end{enumerate}
\end{lemma}

\begin{proof}
(1) を仮定し、$U, V$ を (2) のとおりとする。Lemma
\ref{more-morphisms-lemma-pseudo-coherent} により、
$\mathcal{O}_S(V) \to \mathcal{O}_X(U)$ は擬連接である。したがって、
環準同型が「有限 Tor 次元をもつ」という性質が Morphisms, Definition
\ref{morphisms-definition-property-local} の条件 (1)(a), (b), (c) を満たすことを
示せば十分である。これらの性質は More on Algebra, Lemmas
\ref{more-algebra-lemma-flat-push-tor-amplitude},
\ref{more-algebra-lemma-flat-base-change-finite-tor-dimension}, および
\ref{more-algebra-lemma-glue-tor-amplitude} から従う。一部の詳細は省略する。
\end{proof}

\begin{definition}
\label{more-morphisms-definition-perfect}
スキームの射 $f : X \to S$ が Lemma \ref{more-morphisms-lemma-perfect} の同値な条件を
満たすとき、この射は{\it 完全}であるという。この場合、$X$ は $S$ 上完全である
ともいう。
\end{definition}

\noindent
完全射は特に擬連接であり、したがって局所有限表示であることに注意する。
完全射の基底変換は一般には完全でないことにも注意されたい。

\begin{lemma}
\label{more-morphisms-lemma-flat-base-change-perfect}
完全射の平坦基底変換は完全である。
\end{lemma}

\begin{proof}
これは次の代数的結果に移る。$A \to B$ を完全環準同型とし、$A \to A'$ を
平坦とする。このとき $A' \to B \otimes_A A'$ は完全である。擬連接環準同型に
ついてのこの結果は Lemma \ref{more-morphisms-lemma-flat-base-change-pseudo-coherent} で見た。
有限 Tor 次元についての対応する事実は More on Algebra, Lemma
\ref{more-algebra-lemma-flat-base-change-finite-tor-dimension} から従う。
\end{proof}

\begin{lemma}
\label{more-morphisms-lemma-composition-perfect}
スキームの完全射の合成は完全である。
\end{lemma}

\begin{proof}
これは次の代数的結果に移る。$A \to B \to C$ が合成可能な完全環準同型ならば、
$A \to C$ は完全である。擬連接性については Lemma
\ref{more-morphisms-lemma-composition-pseudo-coherent} とその証明で見た。仮定により、
$B$ の $A$ 上の Tor 次元が $\leq n$、$C$ の $B$ 上の Tor 次元が
$\leq m$ となる整数 $n$, $m$ が存在する。すると任意の $A$-加群 $M$ に対し
$$
M \otimes_A^{\mathbf{L}} C =
(M \otimes_A^{\mathbf{L}} B) \otimes_B^{\mathbf{L}} C
$$
であり、More on Algebra, Example \ref{more-algebra-example-tor} のスペクトル系列は、
望みどおり $p > n + m$ に対して $\text{Tor}^A_p(M, C) = 0$ であることを示す。
\end{proof}

\begin{lemma}
\label{more-morphisms-lemma-flat-finite-presentation-perfect}
$f : X \to S$ をスキームの射とする。次は同値である。
\begin{enumerate}
\item $f$ は平坦かつ完全である。
\item $f$ は平坦かつ局所有限表示である。
\end{enumerate}
\end{lemma}

\begin{proof}
(2) $\Rightarrow$ (1) は More on Algebra, Lemma
\ref{more-algebra-lemma-flat-finite-presentation-perfect} である。逆向きは、
擬連接射が局所有限表示であるという事実から従う。Lemma
\ref{more-morphisms-lemma-pseudo-coherent-finite-presentation} を参照されたい。
\end{proof}

\begin{lemma}
\label{more-morphisms-lemma-regular-target-perfect}
$f : X \to S$ をスキームの射とする。$S$ が正則であり、$f$ が局所有限型であると
仮定する。このとき $f$ は完全である。
\end{lemma}

\begin{proof}
More on Algebra, Lemma
\ref{more-algebra-lemma-regular-perfect-ring-map} を参照されたい。
\end{proof}

\begin{lemma}
\label{more-morphisms-lemma-regular-immersion-perfect}
スキームの正則埋め込みは完全である。
スキームの Koszul 正則埋め込みは完全である。
\end{lemma}

\begin{proof}
正則埋め込みは Koszul 正則埋め込みなので（Divisors, Lemma
\ref{divisors-lemma-regular-quasi-regular-immersion}）、第二の主張を示せば十分である。
これは次の代数的主張に移る。$I \subset A$ を、$A$ の Koszul 正則列
$f_1, \ldots, f_r$ により生成されるイデアルとする。このとき
$A \to A/I$ は完全環準同型である。$A \to A/I$ は全射なので、これは
$A$ 上の多項式代数による $A/I$ の表示である。したがって、$A/I$ が
$A$-加群として擬連接であり有限 Tor 次元をもつことを見れば十分である。
Koszul 列の定義により、Koszul 複体 $K(A, f_1, \ldots, f_r)$ は $A/I$ の
有限自由分解である。したがって $A/I$ は $A$-加群の完全複体であり、結論を得る。
\end{proof}

\begin{lemma}
\label{more-morphisms-lemma-perfect-permanence}
次のスキームの射の可換図式を考える。
$$
\xymatrix{
X \ar[rr]_f \ar[rd] & & Y \ar[ld] \\
& S
}
$$
$Y \to S$ は滑らかであり、$X \to S$ は完全であると仮定する。このとき
$f : X \to Y$ は完全である。
\end{lemma}

\begin{proof}
$f$ を合成
$$
X \longrightarrow X \times_S Y \longrightarrow Y
$$
に分解できる。ここで第一の射は写像 $i = (1, f)$、第二の射は射影である。
$Y \to S$ は平坦なので（Morphisms, Lemma
\ref{morphisms-lemma-smooth-flat}）、Lemma
\ref{more-morphisms-lemma-flat-base-change-perfect} により $X \times_S Y \to Y$ は完全である。
$Y \to S$ は滑らかなので、$X \times_S Y \to X$ も滑らかである。Morphisms,
Lemma \ref{morphisms-lemma-base-change-smooth} を参照されたい。したがって $i$ は
滑らかな射の切断であり、ゆえに $i$ は正則埋め込みである。Divisors, Lemma
\ref{divisors-lemma-section-smooth-regular-immersion} を参照されたい。
これにより $i$ は完全である。Lemma
\ref{more-morphisms-lemma-regular-immersion-perfect} を参照されたい。完全射の合成は完全なので、
$f$ は完全である。Lemma \ref{more-morphisms-lemma-composition-perfect} を参照されたい。
\end{proof}

\begin{remark}
\label{more-morphisms-remark-perfect-permanence}
% P05-EN-CH37-0303
基底 $S$ 上完全なスキーム $X, Y$ の間の射が完全であるとは限らない。
例えば $S = \Spec(k)$、$X = \Spec(k)$、
% P05-EN-CH37-0304
$Y = \Spec(k[x]/(x^2)$ とし、$X \to Y$ を唯一の $S$-射とする。
\end{remark}

\begin{lemma}
\label{more-morphisms-lemma-descending-property-perfect}
性質 $\mathcal{P}(f) =$``$f$ は完全である'' は基底上 fpqc 局所的である。
\end{lemma}

\begin{proof}
これを示すため Descent, Lemma
\ref{descent-lemma-descending-properties-morphisms} の判定条件を用いる。
Definition \ref{more-morphisms-definition-perfect} により、完全性は基底上 Zariski 局所的である。
Lemma \ref{more-morphisms-lemma-flat-base-change-perfect} により、完全性は平坦基底変換で
保たれる。Descent, Lemma
\ref{descent-lemma-descending-properties-morphisms} の最後の仮定 (3) は、
次の代数的主張に移る。$A \to B$ を忠実平坦環準同型とし、
$C = A[x_1, \ldots, x_n]/I$ を $A$-代数とする。
% P05-EN-CH37-0305
$C \otimes_A B$ が $B[x_1, \ldots, x_n]$-加群として完全ならば、
% P05-EN-CH37-0306
$C$ は $A[x_1, \ldots, x_n]$-加群として完全である。これは
More on Algebra, Lemma \ref{more-algebra-lemma-flat-descent-perfect} である。
\end{proof}

\begin{lemma}
\label{more-morphisms-lemma-check-perfect-stalks}
$f : X \to S$ をスキームの擬連接射とする。次は同値である。
\begin{enumerate}
\item $f$ は完全である。
\item $\mathcal{O}_X$ は局所的に $f^{-1}\mathcal{O}_S$-加群の層として
有限 Tor 次元をもつ。
\item すべての $x \in X$ に対し、環 $\mathcal{O}_{X, x}$ は
$\mathcal{O}_{S, f(x)}$-加群として有限 Tor 次元をもつ。
\end{enumerate}
\end{lemma}

\begin{proof}
問題は $X$ と $S$ 上局所的である。したがって $X = \Spec(B)$、
$S = \Spec(A)$ であり、$f$ は擬連接環準同型 $A \to B$ に対応すると
仮定してよい。

\medskip\noindent
(1) が成り立つならば、$B$ は $A$-加群として有限 Tor 次元 $d$ をもつ。
すると、$\mathfrak p = A \cap \mathfrak q$ を満たすすべての素イデアル
$\mathfrak q \subset B$ に対し、$B_\mathfrak q$ は
$A_\mathfrak p$-加群として Tor 次元 $d$ をもつ。More on Algebra, Lemma
\ref{more-algebra-lemma-tor-amplitude-localization} を参照されたい。すると
Cohomology, Lemma \ref{cohomology-lemma-tor-amplitude-stalk} により、
$\mathcal{O}_X$ は $f^{-1}\mathcal{O}_S$-加群の層として Tor 次元 $d$ をもつ。
したがって (1) は (2) を含意する。

\medskip\noindent
Cohomology, Lemma \ref{cohomology-lemma-tor-amplitude-stalk} により、
(2) は (3) を含意する。

\medskip\noindent
(3) を仮定する。$B_\mathfrak q$ の $A_\mathfrak p$ 上の Tor 次元が
$\mathfrak q$ に依存しない上界をもつとは仮定されていないので、結論を得るために
More on Algebra, Lemma \ref{more-algebra-lemma-tor-amplitude-localization} を
用いることはできない。表示 $A[x_1, \ldots, x_n] \to B$ を選ぶ。すると
% P05-EN-CH37-0307
$B$ は $A[x_1, \ldots, x_n]$-加群として擬連接である。
$\mathfrak q \subset A[x_1, \ldots, x_n]$ を
$\mathfrak p \subset A$ 上にある素イデアルとする。すると
$B_\mathfrak q$ は零であるか、仮定により $A_\mathfrak p$-加群として
有限 Tor 次元をもつ。$A \to A[x_1, \ldots, x_n]$ のファイバーは有限大域次元を
もつので、More on Algebra, Lemma
\ref{more-algebra-lemma-perfect-over-polynomial-ring} を
$A_\mathfrak p \to A[x_1, \ldots, x_n]_\mathfrak q$ に適用すると、
$B_\mathfrak q$ は完全 $A[x_1, \ldots, x_n]_\mathfrak q$-加群であることが分かる。
したがって More on Algebra, Lemma
\ref{more-algebra-lemma-check-perfect-stalks} により、$B$ は完全
$A[x_1, \ldots, x_n]$-加群である。ゆえに定義により $A \to B$ は
完全環準同型である。
\end{proof}

\begin{lemma}
\label{more-morphisms-lemma-perfect-closed-immersion-perfect-direct-image}
$i : Z \to X$ をスキームの完全閉埋め込みとする。このとき
$i_*\mathcal{O}_Z$ は完全 $\mathcal{O}_X$-加群、すなわち
$D(\mathcal{O}_X)$ の完全対象である。
\end{lemma}

\begin{proof}
これはほとんど定義から直ちに従う。実際、$U = \Spec(A)$ を $X$ のアフィン開集合と
する。すると、あるイデアル $I \subset A$ に対して
$i^{-1}(U) = \Spec(A/I)$ であり、More on Algebra, Lemma
\ref{more-algebra-lemma-perfect-ring-map} により $A/I$ は有限射影
$A$-加群による有限分解をもつ。したがって $i_*\mathcal{O}_Z|_U$ は
有限局所自由 $\mathcal{O}_U$-加群の有限長複体で表せる。これが示すべきことである。
Cohomology, Section \ref{cohomology-section-perfect} を参照されたい。
\end{proof}

\begin{lemma}
\label{more-morphisms-lemma-perfect-proper-perfect-direct-image}
$S$ を Noetherian スキームとし、$f : X \to S$ をスキームの完全かつ固有な射と
する。$E \in D(\mathcal{O}_X)$ を完全とする。このとき $Rf_*E$ は
$D(\mathcal{O}_S)$ の完全対象である。
\end{lemma}

\begin{proof}
Derived Categories of Schemes, Lemma
\ref{perfect-lemma-perfect-direct-image} を適用できると主張する。条件 (1) と
(2) は直ちに成り立つ。条件 (3) は $X$ 上局所的である。したがって $X$ と $S$ は
アフィンであり、$E$ は $\mathcal{O}_X$-加群の狭義完全複体で表されると仮定してよい。
ゆえに $\mathcal{O}_X$ が $f^{-1}\mathcal{O}_S$-加群の層として有限 Tor 次元を
もつことを示せば十分である。Lemma \ref{more-morphisms-lemma-check-perfect-stalks} により、
これは完全性と同値である。
\end{proof}

\begin{lemma}
\label{more-morphisms-lemma-perfect-fppf-local-source}
性質 $\mathcal{P}(f) =$``$f$ は完全である'' は始域上 fppf 局所的である。
\end{lemma}

\begin{proof}
$\{g_i : X_i \to X\}_{i \in I}$ をスキームの fppf 被覆とし、
$f : X \to S$ を、各 $f \circ g_i$ が完全となる射とする。Lemma
\ref{more-morphisms-lemma-pseudo-coherent-fppf-local-source} により、$f$ は擬連接である。
したがって Lemma \ref{more-morphisms-lemma-check-perfect-stalks} により、すべての
$x \in X$ に対して $\mathcal{O}_{X, x}$ が有限 Tor 次元の
$\mathcal{O}_{S, f(x)}$-加群であることを確認すれば十分である。
$x$ へ写る $i \in I$ と $x_i \in X_i$ を選ぶ。すると
$\mathcal{O}_{X_i, x_i}$ は $\mathcal{O}_{S, f(x)}$ 上有限 Tor 次元をもち、
$\mathcal{O}_{X, x} \to \mathcal{O}_{X_i, x_i}$ は忠実平坦である。
望む結論は More on Algebra, Lemma
\ref{more-algebra-lemma-flat-descent-tor-amplitude} から従う。
\end{proof}

\begin{lemma}
\label{more-morphisms-lemma-factor-regular-immersion}
$i : Z \to Y$ と $j : Y \to X$ をスキームの埋め込みとする。次を仮定する。
\begin{enumerate}
\item $X$ は局所 Noetherian である。
\item $j \circ i$ は正則埋め込みである。
\item $i$ は完全である。
\end{enumerate}
このとき $i$ と $j$ は正則埋め込みである。
\end{lemma}

\begin{proof}
$X$（したがって $Y$）は局所 Noetherian なので、4 種類の正則埋め込みは
すべて一致する。さらに、射が正則埋め込みであるかどうかは局所環の水準で
確認できる。Divisors, Lemma
\ref{divisors-lemma-Noetherian-scheme-regular-ideal} を参照されたい。
したがって結果は Divided Power Algebra, Lemma
\ref{dpa-lemma-regular-sequence} から従う。
\end{proof}









\section{局所完全交叉射}
\label{more-morphisms-section-lci}

\noindent
Divisors, Section \ref{divisors-section-regular-immersions} では、正則、
Koszul 正則、$H_1$-正則、準正則という 4 種類の正則埋め込みを定義した。
本節では、$X$ 上局所的に
$$
\xymatrix{
X \ar[rr]_i \ar[rd] & & \mathbf{A}^n_S \ar[ld] \\
& S
}
$$
と分解し、$* \in \{\emptyset, Koszul, H_1, quasi\}$ に対して $i$ が
$*$-正則埋め込みとなる射 $f : X \to S$ を考える。しかし
$* = \emptyset$、すなわち $i$ に正則埋め込みであることを要求する場合、この条件が
分解に依存しないことをどう証明するかは分からない。他方、局所完全交叉射には
完全であることを望むが、これは $* = Koszul$ または $* = \emptyset$ の場合にしか
成り立たない。そこで、$X$ 上局所的に上のような分解をもち、$i$ が
Koszul 正則埋め込みであるスキームの射 $f : X \to S$ を
{\it 局所完全交叉射}または{\it Koszul 射}と定義する。これが機能することを
見るため、まずこの条件が選んだ分解に依存しないことを示す。

\begin{lemma}
\label{more-morphisms-lemma-koszul-independence-factorization}
$S$ をスキームとし、$U$, $P$, $P'$ を $S$ 上のスキームとする。
$u \in U$ とし、$i : U \to P$、$i' : U \to P'$ を $S$ 上の埋め込みとする。
$P$ と $P'$ は $S$ 上滑らかであると仮定する。このとき次は同値である。
\begin{enumerate}
\item $i$ は $u$ のある近傍で Koszul 正則埋め込みである。
\item $i'$ は $u$ のある近傍で Koszul 正則埋め込みである。
\end{enumerate}
\end{lemma}

\begin{proof}
$i$ は $u$ のある近傍で Koszul 正則埋め込みであると仮定する。射
$j = (i, i') : U \to P \times_S P' = P''$ を考える。
$P'' = P \times_S P' \to P$ は滑らかなので、Divisors, Lemma
\ref{divisors-lemma-lift-regular-immersion-to-smooth} により $j$ は
Koszul 正則埋め込みである。すると Divisors, Lemma
\ref{divisors-lemma-push-regular-immersion-thru-smooth} により $i'$ は
Koszul 正則埋め込みである。
\end{proof}

\noindent
定義を述べる前に、次の簡単な注意をしておこう。$f : X \to S$ を局所有限型な
スキームの射とし、$x \in X$ とする。このとき開近傍 $U \subset X$ と、
$f|_U$ を埋め込み $i : U \to \mathbf{A}^n_S$ と滑らかな射影
$\mathbf{A}^n_S \to S$ の合成として表す分解が存在する。図式は
$$
\xymatrix{
X \ar[rd] & U \ar[l] \ar[d] \ar[r]_-i & \mathbf{A}^n_S = P \ar[ld]^\pi \\
& S
}
$$
である。実際、$X$ における $x$ の任意のアフィン開近傍 $U$ に対してこれを
行える。Morphisms, Lemma
\ref{morphisms-lemma-quasi-affine-finite-type-over-S} を参照されたい。

\begin{definition}
\label{more-morphisms-definition-lci}
$f : X \to S$ をスキームの射とする。
\begin{enumerate}
\item $x \in X$ とする。$f$ が $x$ で有限型であり、
% P05-EN-CH37-0308
ある開近傍 U と $f|_U$ の分解 $\pi \circ i$ が存在して、
$i : U \to P$ が Koszul 正則埋め込み、$\pi : P \to S$ が滑らかであるとき、
$f$ は{\it $x$ で Koszul} であるという。
\item $f$ がすべての点で Koszul であるとき、$f$ は{\it Koszul 射}、または
$f$ は{\it 局所完全交叉射}であるという。
\end{enumerate}
\end{definition}

\noindent
上で、分解 $f|_U = \pi \circ i$ の選択は無関係であることを見た。すなわち、
$f|_U$ を埋め込み $i$ とそれに続く滑らかな射 $\pi$ の合成として表す分解が
与えられたとき、$i$ が $x$ の近傍で Koszul 正則であるかどうかは、$x$ における
$f$ の内在的性質である。後で参照できるよう、これを補題として明記しておく。
% P05-EN-CH37-0309

\begin{lemma}
\label{more-morphisms-lemma-lci}
$f : X \to S$ を局所完全交叉射とし、$P$ を $S$ 上滑らかなスキームとする。
$U \subset X$ を開部分スキームとし、$i : U \to P$ を $S$ 上のスキームの
埋め込みとする。このとき $i$ は Koszul 正則埋め込みである。
\end{lemma}

\begin{proof}
これは局所完全交叉射の定義的性質である。上の議論を参照されたい。
\end{proof}

\noindent
すべての Koszul 射に共通する性質をいくつかここにまとめておく。

\begin{lemma}
\label{more-morphisms-lemma-lci-properties}
$f : X \to S$ を局所完全交叉射とする。このとき
\begin{enumerate}
\item $f$ は局所有限表示である。
\item $f$ は擬連接である。
\item $f$ は完全である。
\end{enumerate}
\end{lemma}

\begin{proof}
完全射は擬連接（完全環準同型は擬連接だから）であり、擬連接射は局所有限表示
（擬連接環準同型は有限表示だから）なので、最後の主張を示せば十分である。
完全性は局所的性質なので、$f$ は $\pi$ が滑らか、$i$ が Koszul 正則埋め込みで
あるように $\pi \circ i$ と分解すると仮定してよい。Koszul 正則埋め込みは完全で
ある。Lemma \ref{more-morphisms-lemma-regular-immersion-perfect} を参照されたい。滑らかな射は
平坦かつ局所有限表示なので完全である。Lemma
\ref{more-morphisms-lemma-flat-finite-presentation-perfect} を参照されたい。最後に完全射の合成は
完全である。Lemma \ref{more-morphisms-lemma-composition-perfect} を参照されたい。
\end{proof}

\begin{lemma}
\label{more-morphisms-lemma-affine-lci}
$f : X = \Spec(B) \to S = \Spec(A)$ をアフィン・スキームの射とする。
このとき $f$ が局所完全交叉射であることと、$A \to B$ が局所完全交叉準同型で
あることは同値である。More on Algebra, Definition
\ref{more-algebra-definition-local-complete-intersection} を参照されたい。
\end{lemma}

\begin{proof}
定義から直ちに従う。
\end{proof}

\noindent
Koszul 射の基底変換は一般には Koszul でないことに注意されたい。

\begin{lemma}
\label{more-morphisms-lemma-flat-base-change-lci}
局所完全交叉射の平坦基底変換は局所完全交叉射である。
\end{lemma}

\begin{proof}
省略する。ヒント：滑らかな射の基底変換は滑らかであり、Koszul 正則埋め込みの
平坦基底変換は Koszul 正則埋め込みであるため、これは成り立つ。Divisors,
Lemma \ref{divisors-lemma-regular-immersion-noetherian} を参照されたい。
\end{proof}

\begin{lemma}
\label{more-morphisms-lemma-composition-lci}
局所完全交叉射の合成は局所完全交叉射である。
\end{lemma}

\begin{proof}
$g : Y \to S$ と $f : X \to Y$ を局所完全交叉射とする。$x \in X$ とし、
$y = f(x)$ とおく。$y$ の開近傍 $V \subset Y$ と、ある Koszul 正則埋め込み
$i : V \to P$ および滑らかな射 $\pi : P \to S$ による分解
$g|_V = \pi \circ i$ を選ぶ。次に、$x \in X$ の開近傍 $U$ と、ある
Koszul 正則埋め込み $i' : U \to P'$ および滑らかな射 $\pi' : P' \to Y$ に
よる分解 $f|_U = \pi' \circ i'$ を選ぶ。実際、Definition
\ref{more-morphisms-definition-lci} の前後の議論により、$P' = \mathbf{A}^n_V$ と仮定してよい。
図式は次のとおりである。
$$
\xymatrix{
X \ar[d] & U \ar[l] \ar[r]_-{i'} & P' = \mathbf{A}^n_V \ar[d] \\
Y \ar[d] &  & V \ar[ll] \ar[r]_i & P \ar[d] \\
S & & & S \ar[lll]
}
$$
$P'' = \mathbf{A}^n_P$ とおく。すると $U \to P' \to P''$ は、Koszul 正則埋め込み
$i'$ と、$P'' \to P$ による $i$ の平坦基底変換との合成として Koszul 正則埋め込み
である。Divisors, Lemma
\ref{divisors-lemma-regular-immersion-noetherian} and Divisors, Lemma
\ref{divisors-lemma-composition-regular-immersion} を参照されたい。また
$P'' \to P \to S$ は滑らかな射の合成として滑らかである。Morphisms, Lemma
\ref{morphisms-lemma-composition-smooth} を参照されたい。したがって望みどおり
$X \to S$ は $x$ で Koszul である。
\end{proof}

\begin{lemma}
\label{more-morphisms-lemma-flat-lci}
\begin{slogan}
射が平坦かつ lci であることと syntomic であることは同値である。
\end{slogan}
$f : X \to S$ をスキームの射とする。次は同値である。
\begin{enumerate}
\item $f$ は平坦な局所完全交叉射である。
\item $f$ は syntomic である。
\end{enumerate}
\end{lemma}

\begin{proof}
アフィン局所的に議論すれば、これは More on Algebra, Lemma
\ref{more-algebra-lemma-syntomic-lci} である。より幾何学的な証明も与える。

\medskip\noindent
(2) を仮定する。Morphisms, Lemma
\ref{morphisms-lemma-syntomic-locally-standard-syntomic} により、$X$ の各点 $x$ に
対し、$x$ のアフィン開近傍 $U$ と $f(x)$ のアフィン開近傍 $V$ が存在し、
$f|_U : U \to V$ は標準 syntomic である。これは
$U = \Spec(R[x_1, \ldots, x_n]/(f_1, \ldots, f_c))
\to V = \Spec(R)$ であり、$R[x_1, \ldots, x_n]/(f_1, \ldots, f_c)$ が
$R$ 上の相対大域完全交叉であることを意味する。Algebra, Lemma
\ref{algebra-lemma-relative-global-complete-intersection-conormal} により、
すべての素イデアル $\mathfrak q \supset (f_1, \ldots, f_c)$ に対し、列
$f_1, \ldots, f_c$ は各局所環 $R[x_1, \ldots, x_n]_{\mathfrak q}$ で正則列で
ある。ホモロジー群 $H_i = H_i(K_\bullet)$ をもつ Koszul 複体
$K_\bullet = K_\bullet(R[x_1, \ldots, x_n], f_1, \ldots, f_c)$ を考える。
More on Algebra, Lemma \ref{more-algebra-lemma-regular-koszul-regular} により、
上の各 $\mathfrak q$ に対して $(H_i)_{\mathfrak q} = 0$, $i > 0$ である。
他方、More on Algebra, Lemma \ref{more-algebra-lemma-homotopy-koszul} により、
$H_i$ は $(f_1, \ldots, f_c)$ で零化される。したがって
$H_i = 0$, $i > 0$ であり、$f_1, \ldots, f_c$ は Koszul 正則列である。
これにより、望みどおり $U \to V$ は Koszul 正則埋め込み
$U \to \mathbf{A}^n_V$ とそれに続く滑らかな射の合成に分解する。

\medskip\noindent
(1) を仮定する。このとき $f$ は平坦かつ局所有限表示である
% P05-EN-CH37-0310
（Lemma \ref{more-morphisms-lemma-lci-properties}）。したがって Morphisms, Lemma
\ref{morphisms-lemma-syntomic-locally-standard-syntomic} により、局所環
$\mathcal{O}_{X_s, x}$ が局所完全交叉環であることを示せば十分である。
$X$ 上局所的に、ある Koszul 正則埋め込み $i : X \to P$ と滑らかな射
$\pi : P \to S$ による分解 $f = \pi \circ i$ を選ぶ。$X \to P$ は $S$ 上の
相対準正則埋め込みであることに注意する。Divisors, Definition
\ref{divisors-definition-relative-H1-regular-immersion} を参照されたい。
したがって Divisors, Lemma
\ref{divisors-lemma-relative-regular-immersion-flat-in-neighbourhood} により、
$X \to P$ は正則埋め込みであり、任意の基底変換後もそうである。
ゆえに各ファイバーは正則埋め込みであり、したがって $X$ のすべてのファイバーの
すべての局所環は局所完全交叉である。
\end{proof}

\begin{lemma}
\label{more-morphisms-lemma-regular-immersion-lci}
スキームの正則埋め込みは局所完全交叉射である。
スキームの Koszul 正則埋め込みは局所完全交叉射である。
\end{lemma}

\begin{proof}
正則埋め込みは Koszul 正則埋め込みなので（Divisors, Lemma
\ref{divisors-lemma-regular-quasi-regular-immersion}）、第二の主張を示せば
十分である。第二の主張は定義から直ちに従う。
\end{proof}

\begin{lemma}
\label{more-morphisms-lemma-lci-permanence}
次のスキームの射の可換図式を考える。
$$
\xymatrix{
X \ar[rr]_f \ar[rd] & & Y \ar[ld] \\
& S
}
$$
% P05-EN-CH37-0311
$Y \to S$ は滑らかであり、$X \to S$ は局所完全交叉射であると仮定する。
このとき $f : X \to Y$ は局所完全交叉射である。
\end{lemma}

\begin{proof}
定義から直ちに従う。
\end{proof}

\begin{lemma}
\label{more-morphisms-lemma-morphism-regular-schemes-is-lci}
$f : X \to Y$ をスキームの射とする。$f$ が局所有限型であり、$X$ と $Y$ が
正則ならば、$f$ は局所完全交叉射である。
\end{lemma}

\begin{proof}
第一の矢印が埋め込みである分解 $X \to \mathbf{A}^n_Y \to Y$ が存在すると
仮定してよい。$Y$ は正則なので、Algebra, Lemma
\ref{algebra-lemma-regular-goes-up} により $\mathbf{A}^n_Y$ も正則である。
したがって Divisors, Lemma
\ref{divisors-lemma-immersion-regular-regular-immersion} により
$X \to \mathbf{A}^n_Y$ は正則埋め込みである。
\end{proof}

\noindent
次の補題は性質が異なる。

\begin{lemma}
\label{more-morphisms-lemma-lci-avramov}
次のスキームの射の可換図式を考える。
$$
\xymatrix{
X \ar[rr]_f \ar[rd] & & Y \ar[ld] \\
& S
}
$$
次を仮定する。
\begin{enumerate}
\item $S$ は局所 Noetherian である。
\item $Y \to S$ は局所有限型である。
\item $f : X \to Y$ は完全である。
\item $X \to S$ は局所完全交叉射である。
\end{enumerate}
このとき $X \to Y$ は局所完全交叉射であり、すべての $x \in X$ に対して
$Y \to S$ は $f(x)$ で Koszul である。
\end{lemma}

\begin{proof}
この証明では、すべてのスキームは局所 Noetherian、すべての環は Noetherian と
する。この状況では正則列と Koszul 正則列が一致することを、以後断りなく用いる。
More on Algebra, Lemma
\ref{more-algebra-lemma-noetherian-finite-all-equivalent} を参照されたい。さらに、
イデアル（resp.\ イデアル層）の正則性は局所環（resp.\ 茎）で確認できる。
Algebra, Lemma \ref{algebra-lemma-regular-sequence-in-neighbourhood}
（resp.\ Divisors, Lemma
\ref{divisors-lemma-Noetherian-scheme-regular-ideal}）を参照されたい。
% P05-EN-CH37-0312

\medskip\noindent
問題は局所的である。したがって $S$, $X$, $Y$ はアフィンだと仮定してよい。
この状況では可換図式
$$
\xymatrix{
\mathbf{A}^{n + m}_S \ar[d] & X \ar[l] \ar[d] \\
\mathbf{A}^n_S \ar[d] & Y \ar[l] \ar[ld] \\
S
}
$$
で水平矢印が閉埋め込みであるものを選べる。点 $x \in X$ をとり、対応する局所環の
可換図式
$$
\xymatrix{
J \ar[r] &
\mathcal{O}_{\mathbf{A}^{n + m}_S, x} \ar[r] &
\mathcal{O}_{X, x} \\
I \ar[r] \ar[u] &
\mathcal{O}_{\mathbf{A}^n_S, f(x)} \ar[r] \ar[u] &
\mathcal{O}_{Y, f(x)} \ar[u]
}
$$
を考える。ここで $J$ と $I$ は水平矢印の核である。$X \to S$ は局所完全交叉射
なので、イデアル $J$ は正則列で生成される。$X \to Y$ は完全なので、環
$\mathcal{O}_{X, x}$ は $\mathcal{O}_{Y, f(x)}$ 上有限 Tor 次元をもつ。
したがって Divided Power Algebra, Lemma
\ref{dpa-lemma-perfect-map-ci} を適用すると、$I$ と $J/I$ は正則列で生成される
ことが分かる。最初の注意により、これで証明が完了する。
\end{proof}

\begin{lemma}
\label{more-morphisms-lemma-lci-to-regular}
次のスキームの射の可換図式を考える。
$$
\xymatrix{
X \ar[rr]_f \ar[rd] & & Y \ar[ld] \\
& S
}
$$
$S$ は局所 Noetherian、$Y \to S$ は局所有限型、$Y$ は正則、さらに
$X \to S$ は局所完全交叉射であると仮定する。このとき $f : X \to Y$ は
局所完全交叉射であり、すべての $x \in X$ に対して $Y \to S$ は
$f(x)$ で Koszul である。
\end{lemma}

\begin{proof}
Lemma \ref{more-morphisms-lemma-regular-target-perfect}（および Morphisms, Lemma
\ref{morphisms-lemma-permanence-finite-type}）により、これは Lemma
\ref{more-morphisms-lemma-lci-avramov} の特別な場合である。
\end{proof}

\begin{lemma}
\label{more-morphisms-lemma-perfect-conormal-free-lci}
$i : X \to Y$ を埋め込みとする。次を仮定する。
\begin{enumerate}
\item $i$ は完全である。
\item $Y$ は局所 Noetherian である。
\item 余法層 $\mathcal{C}_{X/Y}$ は有限局所自由である。
\end{enumerate}
このとき $i$ は正則埋め込みである。
\end{lemma}

\begin{proof}
代数に移すと、これは Divided Power Algebra, Proposition
\ref{dpa-proposition-regular-ideal} である。
\end{proof}

\begin{lemma}
\label{more-morphisms-lemma-lci-NL}
% P05-EN-CH37-0313
$f : X \to Y$ を局所完全交叉準同型とする。このとき素朴余接複体
$\NL_{X/Y}$ は Tor 振幅が $[-1, 0]$ にある $D(\mathcal{O}_X)$ の完全対象である。
\end{lemma}

\begin{proof}
代数に移すと、これは More on Algebra, Lemma
\ref{more-algebra-lemma-lci-NL} である。移し替えには Lemmas
\ref{more-morphisms-lemma-affine-lci} and \ref{more-morphisms-lemma-NL-affine} と、Derived Categories of
Schemes, Lemmas
\ref{perfect-lemma-affine-compare-bounded},
\ref{perfect-lemma-tor-dimension-affine} and
\ref{perfect-lemma-perfect-affine}.
を用いる。
\end{proof}

\begin{lemma}
\label{more-morphisms-lemma-perfect-NL-lci}
$f : X \to Y$ を局所 Noetherian スキームの完全射とする。次は同値である。
\begin{enumerate}
\item $f$ は局所完全交叉射である。
\item $\NL_{X/Y}$ は Tor 振幅が $[-1, 0]$ にある。
\item $\NL_{X/Y}$ は完全であり、Tor 振幅が $[-1, 0]$ にある。
\end{enumerate}
\end{lemma}

\begin{proof}
代数に移すと、これは Divided Power Algebra, Lemma
\ref{dpa-lemma-perfect-NL-lci} である。移し替えには Lemmas
\ref{more-morphisms-lemma-affine-lci} and \ref{more-morphisms-lemma-NL-affine} と、Derived Categories of
Schemes, Lemmas
\ref{perfect-lemma-affine-compare-bounded},
\ref{perfect-lemma-tor-dimension-affine} and
\ref{perfect-lemma-perfect-affine}.
を用いる。
\end{proof}

\begin{lemma}
\label{more-morphisms-lemma-flat-fp-NL-lci}
$f : X \to Y$ を有限表示な平坦射とする。次は同値である。
\begin{enumerate}
\item $f$ は局所完全交叉射である。
\item $f$ は syntomic である。
\item $\NL_{X/Y}$ は Tor 振幅が $[-1, 0]$ にある。
\item $\NL_{X/Y}$ は完全であり、Tor 振幅が $[-1, 0]$ にある。
\end{enumerate}
\end{lemma}

\begin{proof}
代数に移すと、これは Divided Power Algebra, Lemma
\ref{dpa-lemma-flat-fp-NL-lci} である。移し替えには Lemmas
\ref{more-morphisms-lemma-affine-lci} and \ref{more-morphisms-lemma-NL-affine} と、Derived Categories of
Schemes, Lemmas
\ref{perfect-lemma-affine-compare-bounded},
\ref{perfect-lemma-tor-dimension-affine} and
\ref{perfect-lemma-perfect-affine}.
を用いる。
\end{proof}

\noindent
次の補題は、滑らかな射を、対角射が完全な平坦射として特徴づける。

\begin{lemma}
\label{more-morphisms-lemma-smooth-diagonal-perfect}
$f : X \to Y$ を局所 Noetherian スキームの有限型射とする。
% P05-EN-CH37-0314
$\Delta : X \to X \times_Y X$ で対角射を表す。次は同値である。
\begin{enumerate}
\item $f$ は滑らかである。
\item $f$ は平坦であり、$\Delta : X \to X \times_Y X$ は正則埋め込みである。
\item $f$ は平坦であり、$\Delta : X \to X \times_Y X$ は局所完全交叉射である。
\item $f$ は平坦であり、$\Delta : X \to X \times_Y X$ は完全である。
\end{enumerate}
\end{lemma}

\begin{proof}
(1) を仮定する。Morphisms, Lemma \ref{morphisms-lemma-smooth-flat} により
$f$ は平坦である。Morphisms, Lemma
\ref{morphisms-lemma-base-change-smooth} により、射影
$X \times_Y X \to X$ は滑らかである。したがって対角射は滑らかな射の切断であり、
ゆえに正則埋め込みである。Divisors, Lemma
\ref{divisors-lemma-section-smooth-regular-immersion} を参照されたい。
したがって (1) $\Rightarrow$ (2) である。(2) $\Rightarrow$ (3) は Lemma
\ref{more-morphisms-lemma-regular-immersion-lci} である。(3) $\Rightarrow$ (4) は Lemma
\ref{more-morphisms-lemma-lci-properties} である。興味深い含意 (4) $\Rightarrow$ (1) は
Divided Power Algebra, Lemma \ref{dpa-lemma-perfect-diagonal} から直ちに従う。
\end{proof}

\begin{lemma}
\label{more-morphisms-lemma-descending-property-lci}
性質 $\mathcal{P}(f) =$``$f$ は局所完全交叉射である'' は基底上 fpqc 局所的である。
\end{lemma}

\begin{proof}
$f : X \to S$ をスキームの射とし、$\{S_i \to S\}$ を $S$ の fpqc 被覆とする。
$f$ の各基底変換 $f_i : X_i \to S_i$ が局所完全交叉射であると仮定する。
このことから特に $f$ は局所有限型であることに注意する。Lemma
\ref{more-morphisms-lemma-lci-properties} および Descent, Lemma
\ref{descent-lemma-descending-property-locally-finite-type} を参照されたい。
$x \in X$ とし、$x$ の開近傍 $U$ と $S$ 上の埋め込み
$j : U \to \mathbf{A}^n_S$ を選ぶ（Definition \ref{more-morphisms-definition-lci} に先立つ議論を
参照）。$j$ が Koszul 正則埋め込みであることを示さなければならない。
$f_i$ は局所完全交叉射なので、基底変換
$j_i : U \times_S S_i \to \mathbf{A}^n_{S_i}$ は Koszul 正則埋め込みである。
Lemma \ref{more-morphisms-lemma-lci} を参照されたい。$\{\mathbf{A}^n_{S_i} \to
\mathbf{A}^n_S\}$ は fpqc 被覆なので、Descent, Lemma
\ref{descent-lemma-descending-property-regular-immersion} により $j$ は望みどおり
Koszul 正則埋め込みである。
\end{proof}

\begin{lemma}
\label{more-morphisms-lemma-lci-syntomic-local-source}
性質 $\mathcal{P}(f) =$``$f$ は局所完全交叉射である'' は始域上 syntomic 局所的
である。
\end{lemma}

\begin{proof}
これを示すため Descent, Lemma
\ref{descent-lemma-properties-morphisms-local-source} の判定条件を用いる。
Lemmas \ref{more-morphisms-lemma-flat-lci} および \ref{more-morphisms-lemma-composition-lci} により、局所完全
交叉射であることは syntomic 射を前合成しても保たれる。Definition
\ref{more-morphisms-definition-lci} から、局所完全交叉射であることは始域と終域上 Zariski 局所的
である。したがって前述の Descent, Lemma
\ref{descent-lemma-properties-morphisms-local-source} により、次を示せば十分である。
$X' \to X \to Y$ をアフィン・スキームの射とし、$X' \to X$ は syntomic、
$X' \to Y$ は局所完全交叉射であると仮定する。このとき $X \to Y$ は局所完全交叉射
である。いずれにせよ $X \to Y$ は有限表示である。Descent, Lemma
\ref{descent-lemma-flat-finitely-presented-permanence-algebra} を参照されたい。
閉埋め込み $X \to \mathbf{A}^n_Y$ を選ぶ。Algebra, Lemma
\ref{algebra-lemma-lift-syntomic} により、アフィン開被覆
$X' = \bigcup_{i = 1, \ldots, n} X'_i$ と、射 $X'_i \to X$ を持ち上げる
syntomic 射 $W_i \to \mathbf{A}^n_Y$ を選べる。すなわちファイバー積図式
$$
\xymatrix{
X'_i \ar[d] \ar[r] & W_i \ar[d] \\
X \ar[r] & \mathbf{A}^n_Y
}
$$
がある。$X'$ を $\coprod X'_i$ で置き換え、$W = \coprod W_i$ とおくと、
アフィン・スキームのファイバー積図式
$$
\xymatrix{
X' \ar[d] \ar[r] & W \ar[d]^h \\
X \ar[r] & \mathbf{A}^n_Y
}
$$
を得る。ここで $h : W \to \mathbf{A}^n_Y$ は syntomic であり、$X' \to Y$ は
依然として局所完全交叉射である。$W \to \mathbf{A}^n_Y$ は開写像であり
（Morphisms, Lemma \ref{morphisms-lemma-fppf-open}）、$X' \to X$ は全射なので、
$X$ は $W \to \mathbf{A}^n_Y$ の像に含まれる。$\mathbf{A}^n_Y$ 上の閉埋め込み
$W \to \mathbf{A}^{n + m}_Y$ を選ぶ。すると図式は
$$
\xymatrix{
X' \ar[d] \ar[r] & W \ar[d]^h \ar[r] & \mathbf{A}^{n + m}_Y \ar[ld] \\
X \ar[r] & \mathbf{A}^n_Y
}
$$
となる。$h$ は syntomic であり、したがって局所完全交叉射なので（上を参照）、
$W \to \mathbf{A}^{n + m}_Y$ は Koszul 正則埋め込みである。$X' \to Y$ は局所完全
交叉射なので、$X' \to \mathbf{A}^{n + m}_Y$ も Koszul 正則埋め込みである。
Divisors, Lemma \ref{divisors-lemma-permanence-regular-immersion} により
$X' \to W$ は Koszul 正則埋め込みである。したがって Koszul 正則埋め込みは終域上
fpqc 局所的であること（Descent, Lemma
\ref{descent-lemma-descending-property-regular-immersion}）から、
$X \to \mathbf{A}^n_Y$ は Koszul 正則埋め込みである。これが示すべきことであった。
\end{proof}

\begin{lemma}
\label{more-morphisms-lemma-base-change-lci-fibres}
$S$ をスキームとする。$f : X \to Y$ を $S$ 上のスキームの射とする。
$X$ と $Y$ はともに $S$ 上平坦かつ局所有限表示であると仮定する。このとき集合
$$
\{x \in X \mid f\text{ Koszul at }x\}.
$$
は $X$ の開集合であり、その形成は任意の基底変換 $S' \to S$ と可換である。
\end{lemma}

\begin{proof}
集合が開であることは定義から従う（Definition \ref{more-morphisms-definition-lci} を参照）。
$S' \to S$ をスキームの射とする。$X' = S' \times_S X$、
$Y' = S' \times_S Y$ とおき、$f' : X' \to Y'$ で $f$ の基底変換を表す。
$f'$ が $x' \in X'$ で Koszul であるとする。$x'$ の像を
$s' \in S'$、$x \in X$、$y' \in Y'$、$y \in Y$、$s \in S$ と表す。
これは $x'$ の像である。
$f$ は局所有限表示であることに注意する。Morphisms, Lemma
\ref{morphisms-lemma-finite-presentation-permanence} を参照されたい。したがって
$x$ のアフィン近傍 $U \subset X$ と埋め込み $i : U \to \mathbf{A}^n_Y$ を
選べる。$U' = S' \times_S U$ とおき、$i' : U' \to \mathbf{A}^n_{Y'}$ で
$i$ の基底変換を表す。$f'$ が $x'$ で Koszul であるという仮定は、$i'$ が
$x'$ の近傍で Koszul 正則埋め込みであることを含意する。Lemma
\ref{more-morphisms-lemma-lci} を参照されたい。$X'$ は $X$ の基底変換なので $S'$ 上平坦かつ
局所有限表示である（Morphisms, Lemmas
\ref{morphisms-lemma-base-change-flat} and
\ref{morphisms-lemma-base-change-finite-presentation}）。したがって $i'$ は $x'$ の近傍で
$S'$ 上
相対 $H_1$-正則埋め込みである。Divisors, Definition
\ref{divisors-definition-relative-H1-regular-immersion} を参照されたい。
ゆえに基底変換 $i'_{s'} : U'_{s'} \to \mathbf{A}^n_{Y'_{s'}}$ は $x'$ の
開近傍で $H_1$-正則埋め込みである。Divisors, Lemma
\ref{divisors-lemma-relative-regular-immersion} と
\ref{divisors-definition-relative-H1-regular-immersion} に続く議論を
参照されたい。$s' = \Spec(\kappa(s')) \to \Spec(\kappa(s)) = s$ は全射平坦な
普遍開射なので（Morphisms, Lemma
\ref{morphisms-lemma-scheme-over-field-universally-open}）、基底変換
$i_s : U_s \to \mathbf{A}^n_{Y_s}$ は $x$ の近傍で $H_1$-正則埋め込みである。
Descent, Lemma \ref{descent-lemma-descending-property-regular-immersion} を参照。
最後に $\mathbf{A}^n_Y$ は $S$ 上平坦かつ局所有限表示なので、Divisors, Lemma
\ref{divisors-lemma-fibre-quasi-regular} により $i$ は $x$ の近傍で
(Koszul-)正則埋め込みである。
\end{proof}

\begin{lemma}
\label{more-morphisms-lemma-unramified-lci}
$f : X \to Y$ をスキームの局所完全交叉射とする。このとき $f$ が非分岐である
ことと、$f$ が形式的に非分岐であることは同値であり、この場合余法層
$\mathcal{C}_{X/Y}$ は $X$ 上有限局所自由である。
\end{lemma}

\begin{proof}
第一の主張は Lemma \ref{more-morphisms-lemma-unramified-formally-unramified} と、局所完全交叉射
が局所有限型であることから直ちに従う。$f$ の余法層を計算するため、$X$ 上局所的に
$f$ を $f = p \circ i$ と分解する。ここで $i : X \to V$ は Koszul 正則埋め込み、
$V \to Y$ は滑らかである。Lemma \ref{more-morphisms-lemma-two-unramified-morphisms-formally-smooth}
により、$\mathcal{C}_{X/Y}$ は $\mathcal{C}_{X/V}$ の局所直和因子である。
$i$ は Koszul 正則（したがって準正則）埋め込みなので、後者は有限局所自由である。
Divisors, Lemma \ref{divisors-lemma-quasi-regular-immersion} を参照されたい。
\end{proof}

\begin{lemma}
\label{more-morphisms-lemma-transitivity-conormal-lci}
$Z \to Y \to X$ をスキームの形式的に非分岐な射とする。
$Z \to Y$ は局所完全交叉射であると仮定する。Lemma
\ref{more-morphisms-lemma-transitivity-conormal} の完全列
$$
0 \to i^*\mathcal{C}_{Y/X} \to
\mathcal{C}_{Z/X} \to
\mathcal{C}_{Z/Y} \to 0
$$
は短完全列である。
\end{lemma}

\begin{proof}
問題は $Z$ 上局所的なので、射 $Z \to Y$ の分解
$Z \to \mathbf{A}^n_Y \to Y$ が存在すると仮定してよい。このとき可換図式
$$
\xymatrix{
Z \ar[r]_{i'} \ar@{=}[d] &
\mathbf{A}^n_Y \ar[r] \ar[d] &
\mathbf{A}^n_X \ar[d] \\
Z \ar[r]^i & Y \ar[r] & X
}
$$
を得る。$Z \to Y$ は局所完全交叉射なので、$Z \to \mathbf{A}^n_Y$ は
Koszul 正則埋め込みである。したがって Divisors, Lemma
\ref{divisors-lemma-transitivity-conormal-quasi-regular} により、列
$$
0 \to (i')^*\mathcal{C}_{\mathbf{A}^n_Y/\mathbf{A}^n_X} \to
\mathcal{C}_{Z/\mathbf{A}^n_X} \to
\mathcal{C}_{Z/\mathbf{A}^n_Y} \to 0
$$
は完全かつ局所的に分裂する。ここで
$i^*\mathcal{C}_{Y/X} = (i')^*\mathcal{C}_{\mathbf{A}^n_Y/\mathbf{A}^n_X}$
であることに注意する。これは Lemma
\ref{more-morphisms-lemma-universal-thickening-fibre-product-flat} による。さらに、図式
$$
\xymatrix{
(i')^*\mathcal{C}_{\mathbf{A}^n_Y/\mathbf{A}^n_X} \ar[r] &
\mathcal{C}_{Z/\mathbf{A}^n_X} \\
i^*\mathcal{C}_{Y/X} \ar[u]^{\cong} \ar[r] & \mathcal{C}_{Z/X} \ar[u]
}
$$
が可換であることに注意する。したがって下の水平射は局所的に分裂する単射である。
これで補題が証明された。
\end{proof}













\section{微分と余法層の完全列}
\label{more-morphisms-section-exact}

\noindent
この節では、余法層および微分の層の完全列に関するいくつかの結果をまとめる。
ある意味で、これらはすべて、合成可能なスキームの射の組に付随する余接複体の三角形の
実現である。

\medskip\noindent
$g : Z \to Y$ および $f : Y \to X$ をスキームの射とする。
\begin{enumerate}
\item 標準的な完全列
$$
g^*\Omega_{Y/X} \to \Omega_{Z/X} \to \Omega_{Z/Y} \to 0,
$$
が存在する。Morphisms, Lemma \ref{morphisms-lemma-triangle-differentials} を参照。
$g : Z \to Y$ が滑らか、より一般に形式的に滑らかならば、この列は短完全列である。
Morphisms, Lemma \ref{morphisms-lemma-triangle-differentials-smooth} または
Lemma \ref{more-morphisms-lemma-triangle-differentials-formally-smooth} を参照。
\item $g$ が埋め込み、より一般に形式的に非分岐ならば、標準的な完全列
$$
\mathcal{C}_{Z/Y} \to g^*\Omega_{Y/X} \to \Omega_{Z/X} \to 0,
$$
が存在する。Morphisms, Lemma \ref{morphisms-lemma-differentials-relative-immersion}
または Lemma \ref{more-morphisms-lemma-universally-unramified-differentials-sequence} を参照。
$f \circ g : Z \to X$ が滑らか、より一般に形式的に滑らかならば、この列は短完全列である。
Morphisms, Lemma \ref{morphisms-lemma-differentials-relative-immersion-smooth} または
Lemma \ref{more-morphisms-lemma-differentials-formally-unramified-formally-smooth} を参照。
\item $g$ および $f \circ g$ が埋め込み、より一般に形式的に非分岐ならば、標準的な完全列
$$
\mathcal{C}_{Z/X} \to \mathcal{C}_{Z/Y} \to g^*\Omega_{Y/X} \to 0,
$$
が存在する。Morphisms, Lemma \ref{morphisms-lemma-two-immersions} または
Lemma \ref{more-morphisms-lemma-two-unramified-morphisms} を参照。$f : Y \to X$ が滑らか、より一般に
形式的に滑らかならば、この列は短完全列である。Morphisms, Lemma
\ref{morphisms-lemma-two-immersions-smooth} または Lemma
\ref{more-morphisms-lemma-two-unramified-morphisms-formally-smooth} を参照。
\item $g$ および $f$ が埋め込み、より一般に形式的に非分岐ならば、標準的な完全列
$$
g^*\mathcal{C}_{Y/X} \to \mathcal{C}_{Z/X} \to \mathcal{C}_{Z/Y} \to 0.
$$
が存在する。Morphisms, Lemma \ref{morphisms-lemma-transitivity-conormal} または
Lemma \ref{more-morphisms-lemma-transitivity-conormal} を参照。$g : Z \to Y$ が正則埋め込み
\footnote{$g$ が $H_1$-正則埋め込みであれば十分である。局所完全交叉射である埋め込みは
Koszul 正則であることに注意。}
ならば、より一般に局所完全交叉射の場合にも、この列は短完全列である。Divisors,
Lemma \ref{divisors-lemma-transitivity-conormal-quasi-regular} または Lemma
\ref{more-morphisms-lemma-transitivity-conormal-lci} を参照。
\end{enumerate}










\section{Weakly \'etale morphisms}
\label{more-morphisms-section-weakly-etale}

\noindent
環準同型 $A \to B$ が弱エタールであるとは、$A \to B$ と
$B \otimes_A B \to B$ の両方が平坦であることをいう。More on Algebra,
定義 \ref{more-algebra-definition-weakly-etale} を参照せよ。
スキームの射に対する対応する概念は次のものである。

\begin{definition}
\label{more-morphisms-definition-weakly-etale}
スキームの射 $X \to Y$ が {\it 弱エタール} または
{\it 絶対平坦} であるとは、$X \to Y$ と対角射
$X \to X \times_Y X$ の両方が平坦であることをいう。
\end{definition}

\noindent
エタール射は弱エタールであり、逆に弱エタール射は実際に
ある意味でエタール射に似ていることが分かる。例えば
$X \to Y$ が弱エタールならば、これは
Cotangent, 補題 \ref{cotangent-lemma-when-zero} から従う
$L_{X/Y} = 0$ を満たす。
弱エタール射とエタール射を関係づける非常に精密な結果は後で証明する
(Pro-\'etale Cohomology, 節 \ref{proetale-section-weakly-etale} を参照)。
この節では基本事項にとどめる。

\begin{lemma}
\label{more-morphisms-lemma-check-weakly-etale-stalks}
スキームの射 $f : X \to Y$ に対して、次は同値である。
\begin{enumerate}
\item $X \to Y$ は弱エタールである。
\item 任意の $x \in X$ に対して、環写像
$\mathcal{O}_{Y, f(x)} \to \mathcal{O}_{X, x}$ は弱エタールである。
\end{enumerate}
\end{lemma}

\begin{proof}
(1) と (2) のいずれの仮定の下でも射 $f$ は平坦であることに注意する。
したがって $f$ が平坦であると仮定してよい。$Y$ における像が
$y = f(x)$ である $x \in X$ を取る。環の標準的な写像
$$
\mathcal{O}_{X, x} \otimes_{\mathcal{O}_{Y, y}} \mathcal{O}_{X, x}
\longrightarrow
\mathcal{O}_{X \times_Y X, \Delta_{X/Y}(x)}
\longrightarrow
\mathcal{O}_{X, x}
$$
ここで第一の写像は局所化（したがって平坦）であり、第二の写像は
全射である。(1) の条件は、すべての $x$ に対して第二の矢印が平坦である
ことを意味する。(2) は合成がすべての $x$ に対して平坦であるという条件である。
% P05-EN-CH37-0317
したがって同値性は Algebra, 補題 \ref{algebra-lemma-flat-localization}
第 (2) 部分による。
\end{proof}

\begin{lemma}
\label{more-morphisms-lemma-key}
スキームの射 $X \to Y$ が
$X \to X \times_Y X$ を平坦にするものとする。$\mathcal{F}$ を
$\mathcal{O}_X$-加群とする。$\mathcal{F}$ が $Y$ 上平坦ならば、
$\mathcal{F}$ は $X$ 上平坦である。
\end{lemma}

\begin{proof}
$Y$ における像が $y = f(x)$ である $x \in X$ を取る。
$X \to X \times_Y X$ は平坦なので、
$\mathcal{O}_{X, x} \otimes_{\mathcal{O}_{Y, y}} \mathcal{O}_{X, x} \to
\mathcal{O}_{X, x}$ は平坦である。したがって結論は
More on Algebra, 補題 \ref{more-algebra-lemma-key} と定義から従う。
\end{proof}

\begin{lemma}
\label{more-morphisms-lemma-weakly-etale-characterize}
スキームの射 $f : X \to S$ に対して、次は同値である。
\begin{enumerate}
\item 射 $f$ は弱エタールである。
% P05-EN-CH37-0315
\item $f(U) \subset V$ を満たす任意のアフィン開集合
$U \subset X$, $V \subset S$ に対して、環写像
$\mathcal{O}_S(V) \to \mathcal{O}_X(U)$ は弱エタールである。
\item 開被覆 $S = \bigcup_{j \in J} V_j$ と、開被覆
$f^{-1}(V_j) = \bigcup_{i \in I_j} U_i$ であって、
$j\in J, i\in I_j$ の各射 $U_i \to V_j$ が
弱エタールとなるものが存在する。
\item アフィン開被覆 $S = \bigcup_{j \in J} V_j$ と、アフィン開被覆
$f^{-1}(V_j) = \bigcup_{i \in I_j} U_i$ であって、
すべての $j\in J, i\in I_j$ に対し環写像
$\mathcal{O}_S(V_j) \to \mathcal{O}_X(U_i)$ が弱エタールとなるものが存在する。
\end{enumerate}
さらに、$f$ が弱エタールならば、$f(U) \subset V$ を満たす任意の
開部分スキーム $U \subset X$, $V \subset S$ に対して制限
$f|_U : U \to V$ は弱エタールである。
\end{lemma}

\begin{proof}
$f(U) \subset V$ を満たす開部分スキーム $U \subset X$, $V \subset S$ を
取る。このとき $U \times_V U \subset X \times_Y X$ は開であり
(Schemes, 補題 \ref{schemes-lemma-open-fibre-product})、
$f|_U : U \to V$ の対角 $\Delta_{U/V}$ は
$\Delta_{X/Y}|_U : U \to U \times_V U$ の制限である。
平坦性はスキームの射に関する局所的性質なので
(Morphisms, 補題 \ref{morphisms-lemma-flat-characterize})、
% P05-EN-CH37-0316
補題の最後の主張、ならびに (1) と (3) の同値性が従う。
$X$ と $Y$ がアフィンならば、$X \to Y$ が弱エタールであることと
$\mathcal{O}_Y(Y) \to \mathcal{O}_X(X)$ が弱エタールであることは同値である
(再び Morphisms, 補題 \ref{morphisms-lemma-flat-characterize} を用いる)。
したがって (1) と (3) は (2) と (4) にも同値である。
\end{proof}

\begin{lemma}
\label{more-morphisms-lemma-composition-weakly-etale}
スキームの射 $X \to Y \to Z$ を考える。
\begin{enumerate}
\item $X \to X \times_Y X$ と $Y \to Y \times_Z Y$ が平坦ならば、
$X \to X \times_Z X$ は平坦である。
\item $X \to Y$ と $Y \to Z$ が弱エタールならば、
$X \to Z$ は弱エタールである。
\end{enumerate}
\end{lemma}

\begin{proof}
(1) は対角射の次の分解
$$
X \to X \times_Y X \to X \times_Z X
$$
から従う。ここで $X$ の $Z$ 上の対角射について、
$$
X \times_Y X = (X \times_Z X) \times_{(Y \times_Z Y)} Y,
$$
平坦射の基底変換は平坦であること、ならびに平坦射の合成は平坦であること
(Morphisms, 補題 \ref{morphisms-lemma-base-change-flat} および
\ref{morphisms-lemma-composition-flat}) を用いた。
(2) は (1) と、今用いた平坦射の合成が平坦であるという事実から従う。
\end{proof}

\begin{lemma}
\label{more-morphisms-lemma-base-change-weakly-etale}
スキームの射 $X \to Y$ と $Y' \to Y$ を取り、
$X' = Y' \times_Y X$ を $X$ の基底変換とする。
\begin{enumerate}
\item $X \to X \times_Y X$ が平坦ならば、
$X' \to X' \times_{Y'} X'$ は平坦である。
\item $X \to Y$ が弱エタールならば、$X' \to Y'$ は弱エタールである。
\end{enumerate}
\end{lemma}

\begin{proof}
$X \to X \times_Y X$ が平坦であると仮定する。射
$X' \to X' \times_{Y'} X'$ は $Y' \to Y$ による
$X \to X \times_Y X$ の基底変換である。したがって
Morphisms, 補題 \ref{morphisms-lemma-base-change-flat} により平坦である。
これは (1) を示す。(2) は (1) と、今用いた平坦射の基底変換は平坦である
という事実から従う。
\end{proof}

\begin{lemma}
\label{more-morphisms-lemma-go-down}
スキームの射 $X \to Y \to Z$ を考える。$X \to Y$ が
平坦かつ全射であり、$X \to X \times_Z X$ が平坦であると仮定する。
このとき $Y \to Y \times_Z Y$ は平坦である。
\end{lemma}

\begin{proof}
可換図式
$$
\xymatrix{
X \ar[r] \ar[d] & X \times_Z X \ar[d] \\
Y \ar[r] & Y \times_Z Y
}
$$
において上の水平矢印は平坦であり、垂直矢印も平坦である。
したがって $X$ は $Y \times_Z Y$ 上平坦である。
Morphisms, 補題 \ref{morphisms-lemma-flat-permanence} により、
$Y$ は $Y \times_Z Y$ 上平坦であることが分かる。
\end{proof}

\begin{lemma}
\label{more-morphisms-lemma-weakly-etale-formally-unramified}
スキームの弱エタール射 $f : X \to Y$ を取る。
このとき $f$ は形式的不分岐、すなわち $\Omega_{X/Y} = 0$ である。
\end{lemma}

\begin{proof}
補題 \ref{more-morphisms-lemma-formally-unramified-differentials} により、
$f$ が形式的不分岐であることと $\Omega_{X/Y} = 0$ は同値である。
補題 \ref{more-morphisms-lemma-weakly-etale-characterize} と
Morphisms, 補題 \ref{morphisms-lemma-differentials-affine} を介して、
これは環の場合、すなわち More on Algebra, 補題
\ref{more-algebra-lemma-formally-unramified} から従う。
\end{proof}

\begin{lemma}
\label{more-morphisms-lemma-when-weakly-etale}
スキームの射 $f : X \to Y$ を取る。次のいずれの場合にも
$X \to Y$ は弱エタールである。
\begin{enumerate}
\item $X \to Y$ が平坦な単射である。
\item $X \to Y$ が開埋め込みである。
\item $X \to Y$ が平坦かつ不分岐である。
\item $X \to Y$ がエタールである。
\end{enumerate}
\end{lemma}

\begin{proof}
(1) の場合、$\Delta_{X/Y}$ は同型
(Schemes, 補題 \ref{schemes-lemma-monomorphism}) なので、
$f$ は明らかに弱エタールである。(2) は (1) の特別な場合である。
不分岐射の対角は開埋め込み
(Morphisms, 補題 \ref{morphisms-lemma-diagonal-unramified-morphism})
なので平坦である。したがって平坦な不分岐射は弱エタールである。
エタール射は平坦かつ不分岐
(Morphisms, 補題 \ref{morphisms-lemma-etale-smooth-unramified}) なので、
(4) は (3) から従う。
\end{proof}

\begin{lemma}
\label{more-morphisms-lemma-reduced-goes-up}
スキームの射 $f : X \to Y$ を取る。
$Y$ が被約で $f$ が弱エタールならば、$X$ は被約である。
\end{lemma}

\begin{proof}
補題 \ref{more-morphisms-lemma-weakly-etale-characterize} を介して、
これは環の場合、すなわち More on Algebra, 補題
\ref{more-algebra-lemma-absolutely-flat-over-absolutely-flat} から従う。
\end{proof}

\noindent
次の補題では弱エタールな環写像についての非自明な結果を用いる。

\begin{lemma}
\label{more-morphisms-lemma-weakly-etale-strictly-henselian-local-rings}
スキームの射 $f : X \to Y$ を取る。次は同値である。
\begin{enumerate}
\item $f$ は弱エタールである。
\item $x \in X$ に対する局所環写像
$\mathcal{O}_{Y, f(x)} \to \mathcal{O}_{X, x}$ が厳密ヘンゼル化上
同型を誘導する。
\end{enumerate}
\end{lemma}

\begin{proof}
$Y$ における像が $y = f(x)$ である点 $x \in X$ を取る。
$\kappa(x)$ の分離代数閉包 $\kappa^{sep}$ を選ぶ。
$\mathcal{O}_{X, x}^{sh}$ を $\kappa^{sep}$ に対応する厳密ヘンゼル化とし、
$\mathcal{O}_{Y, y}^{sh}$ を $\kappa^{sep}$ における $\kappa(y)$ の
分離代数閉包に関する厳密ヘンゼル化とする。可換図式
$$
\xymatrix{
\mathcal{O}_{X, x} \ar[r] & \mathcal{O}_{X, x}^{sh} \\
\mathcal{O}_{Y, y} \ar[u] \ar[r] & \mathcal{O}_{Y, y}^{sh} \ar[u]
}
$$
は局所環の局所準同型である。Algebra, 補題
\ref{algebra-lemma-strictly-henselian-functorial} を参照せよ。
厳密ヘンゼル化はエタール環写像のフィルター付き余極限なので、
More on Algebra, 補題 \ref{more-algebra-lemma-when-weakly-etale} により
水平写像は弱エタールである。さらに水平写像は More on Algebra, 補題
\ref{more-algebra-lemma-dumb-properties-henselization} により忠実平坦である。

\medskip\noindent
$f$ が弱エタールであると仮定する。補題
\ref{more-morphisms-lemma-check-weakly-etale-stalks} により左の垂直矢印は弱エタールである。
More on Algebra, 補題 \ref{more-algebra-lemma-composition-weakly-etale} と
\ref{more-algebra-lemma-weakly-etale-permanence} により右の垂直矢印も
弱エタールである。More on Algebra, 定理
\ref{more-algebra-theorem-olivier} により、右の垂直写像は同型であると結論する。

\medskip\noindent
$\mathcal{O}_{Y, y}^{sh} \to \mathcal{O}_{X, x}^{sh}$ が同型であると仮定する。
このとき $\mathcal{O}_{Y, y} \to \mathcal{O}_{X, x}^{sh}$ は弱エタールである。
$\mathcal{O}_{X, x} \to \mathcal{O}_{X, x}^{sh}$ は忠実平坦なので、
More on Algebra, 補題 \ref{more-algebra-lemma-go-down} により
$\mathcal{O}_{Y, y} \to \mathcal{O}_{X, x}$ は弱エタールであると結論する。
したがって補題 \ref{more-morphisms-lemma-check-weakly-etale-stalks} により (2) は (1) を含意する。
\end{proof}

\begin{lemma}
\label{more-morphisms-lemma-normal-goes-up}
スキームの射 $f : X \to Y$ を取る。$Y$ が正規スキームで
$f$ が弱エタールならば、$X$ は正規スキームである。
\end{lemma}

\begin{proof}
More on Algebra, 補題 \ref{more-algebra-lemma-henselization-normal} により、
スキーム $S$ が正規であることと、すべての $s \in S$ に対して
$\mathcal{O}_{S, s}$ の厳密ヘンゼル化が正規整域であることは同値である。
したがってこの補題は補題
\ref{more-morphisms-lemma-weakly-etale-strictly-henselian-local-rings} から従う。
\end{proof}

\begin{lemma}
\label{more-morphisms-lemma-weakly-etale-permanence}
スキーム $S$ を取る。$S$ 上のスキームの射
$f : X \to Y$ を取る。$X$, $Y$ が $S$ 上弱エタールならば、
$f$ は弱エタールである。
\end{lemma}

\begin{proof}
以下では Morphisms, 補題 \ref{morphisms-lemma-base-change-flat} と
\ref{morphisms-lemma-composition-flat} を断りなく用いる。
$X \to Y$ を $X \to X \times_S Y \to Y$ の合成として書く。
第二の射は平坦射 $X \to S$ の基底変換なので平坦である。
第一の射は、射 $X \times_S Y \to Y \times_S Y$ による
平坦射 $Y \to Y \times_S Y$ の基底変換なので平坦である。
したがって $X \to Y$ は平坦である。射
$X \times_Y X \to X \times_S X$ は埋め込みである。
よって補題 \ref{more-morphisms-lemma-key} により、$X$ は
% P05-EN-CH37-0318
$X \times_S X$ 上平坦なので、$X$ は $X \times_Y X$ 上も平坦である。
\end{proof}

\noindent
次は、分離的かつ純非分離的でもある体拡大は同型でなければならない
という観察の、スキーム論的な一般化である。

\begin{lemma}
\label{more-morphisms-lemma-weakly-etale-universal-homeomorphism}
スキームの射 $f : X \to Y$ を取る。$f$ が弱エタールかつ
普遍同相ならば、これは同型である。
\end{lemma}

\begin{proof}
$f$ は普遍同相なので、対角
$\Delta : X \to X \times_Y X$ は全射閉埋め込みである
(Morphisms, 補題 \ref{morphisms-lemma-homeomorphism-affine} および
\ref{morphisms-lemma-universally-injective})。$\Delta$ も平坦なので、
Morphisms, 補題 \ref{morphisms-lemma-characterize-flat-closed-immersions}
により $\Delta$ は同型でなければならない。言い換えれば $f$ は単射である
(Schemes, 補題 \ref{schemes-lemma-monomorphism})。
$f$ は普遍同相なので明らかに準コンパクトである。したがって Descent,
補題 \ref{descent-lemma-flat-surjective-quasi-compact-monomorphism-isomorphism}
により $f$ は同型であることが分かる。
\end{proof}

\noindent
次は \'Etale Morphisms, 補題
\ref{etale-lemma-relative-frobenius-etale} の弱エタール版である。

\begin{lemma}
\label{more-morphisms-lemma-relative-frobenius-weakly-etale}
標数 $p$ のスキーム $X$ に対するスキームの弱エタール射
$U \to X$ を取る。このとき相対フロベニウス
$F_{U/X} : U \to U \times_{X, F_X} X$ は同型である。
\end{lemma}

\begin{proof}
射 $F_{U/X}$ は Varieties, 補題
\ref{varieties-lemma-relative-frobenius} により普遍同相である。
$F_{U/X}$ は $X$ 上弱エタールなスキームの間の射なので、
補題 \ref{more-morphisms-lemma-weakly-etale-permanence} により弱エタールである。
したがって補題 \ref{more-morphisms-lemma-weakly-etale-universal-homeomorphism} により
$F_{U/X}$ は同型である。
\end{proof}






\section{被約ファイバー定理}
\label{more-morphisms-section-reduced-fibre-theorem}

\noindent
この節では、題名が示す種類の定理のうち最も単純なものを議論する。
結果の証明はある程度手間がかかるが、本質的には分岐の除去に関する
Epp の結果から直接的に従う。More on Algebra, 定理
\ref{more-algebra-theorem-epp} を参照せよ。

\medskip\noindent
$A$ を分数体 $K$ をもつ Dedekind 整域とする。
$X$ を $A$ 上平坦かつ有限型のスキームとする。
$L$ を $K$ の有限拡大とし、$B$ を $L$ における $A$ の整閉包とする。
このとき $B$ は Dedekind 整域である
(Algebra, 補題 \ref{algebra-lemma-integral-closure-Dedekind})。
$X_B = X \times_{\Spec(A)} \Spec(B)$ を基底変換とする。
すると $X_B \to \Spec(B)$ は有限型である
(Morphisms, 補題 \ref{morphisms-lemma-base-change-finite-type})。
したがって $X_B$ は Noetherian である
(Morphisms, 補題 \ref{morphisms-lemma-finite-type-noetherian})。
よって正規化 $\nu : Y \to X_B$ が存在する
(Morphisms, 定義 \ref{morphisms-definition-normalization} とその直後の議論を参照)。
図式は次のとおりである。
\begin{equation}
\label{more-morphisms-equation-normalized-base-change}
\xymatrix{
Y \ar[rd] \ar[r]_\nu & X_B \ar[r] \ar[d] & X \ar[d] \\
 & \Spec(B) \ar[r] & \Spec(A)
}
\end{equation}
$Y$ を $X$ の {\it 正規化基底変換} と呼ぶことがある。
一般には射 $\nu$ は有限とは限らない。しかし $A$ が Nagata 環
(実際上ほとんど常に満たされる条件) ならば、$\nu$ は有限であり
$Y$ は $B$ 上有限型である。Morphisms, 補題
\ref{morphisms-lemma-nagata-normalization} および
\ref{morphisms-lemma-finite-type-nagata} を参照せよ。

\medskip\noindent
正規化基底変換の形成は合成と可換である。
より正確には、$M/L/K$ が体の有限拡大で、整閉包が
$A \subset B \subset C$ であるとする。このとき
$M/L$ に関する $Y \to \Spec(B)$ の正規化基底変換 $Z$ は、
$M/K$ に関する $X \to \Spec(A)$ の正規化基底変換に等しい。

\begin{theorem}
\label{more-morphisms-theorem-normalized-base-change-with-reduced-fibre}
$A$ を分数体 $K$ をもつ Dedekind 環とする。
$X$ を $A$ 上平坦かつ有限型のスキームとする。
$A$ が Nagata 環であると仮定する。
すべてのファイバーのすべての一般点で正規化基底変換 $Y$ が
$\Spec(B)$ 上滑らかとなるような有限拡大 $L/K$ が存在する。
\end{theorem}

\begin{proof}
証明では、$X \to \Spec(A)$ のような（平坦かつ有限表示の）射が
滑らかとなる点の集合の形成は基底変換と可換することを繰り返し用いる。
Morphisms, 補題 \ref{morphisms-lemma-set-points-where-fibres-smooth} を参照せよ。

\medskip\noindent
まず $(X_L)_{red}$ が $L$ 上幾何学的に被約となるような有限拡大
$L/K$ を選ぶ。Varieties, 補題
\ref{varieties-lemma-finite-extension-geometrically-reduced} を参照せよ。
$Y \to (X_B)_{red}$ は双有理なので、Varieties, 補題
\ref{varieties-lemma-generic-points-geometrically-reduced} を適用すると
$Y_L$ も $L$ 上幾何学的に被約であることが分かる。
したがって Varieties, 補題
\ref{varieties-lemma-geometrically-reduced-dense-smooth-open} により、
$Y_L \to \Spec(L)$ は稠密開集合 $V \subset Y_L$ 上滑らかである。
よって射 $Y \to \Spec(B)$ の滑らかな軌跡 $U \subset Y$ は
(Morphisms, 定義 \ref{morphisms-definition-smooth} により) 開であり、
一般ファイバーに稠密である。$A$ を $B$ に、$X$ を $Y$ に置き換えると、
次の段落で扱う場合に帰着する。

\medskip\noindent
$X$ が正規で、$X \to \Spec(A)$ の滑らかな軌跡 $U \subset X$ が
一般ファイバーに稠密であると仮定する。これは $U$ が有限個を除く
すべてのファイバーに稠密であることを意味する。補題
\ref{more-morphisms-lemma-nowhere-dense-generic-fibre} を参照せよ。
$x_1, \ldots, x_r \in X \setminus U$ を、$X \setminus U$ の既約成分の
一般点であり、さらに $X \to \Spec(A)$ のファイバーの既約成分の
一般点でもある有限個の点とする。
$\mathcal{O}_i = \mathcal{O}_{X, x_i}$ と置く。$A_i$ を
$\Spec(A)$ における $x_i$ の像に対応する極大イデアルでの $A$ の局所化とする。
More on Algebra, 命題
\ref{more-algebra-proposition-epp-essentially-finite-type} により、
離散付値環の拡大 $A_i \to \mathcal{O}_i$ の解となる有限拡大
$K_i/K$ が存在する。すべての拡大 $K_i/K$ を支配する有限拡大
$L/K$ を取る。このとき More on Algebra, 補題
\ref{more-algebra-lemma-solution-goes-up} により、
$L/K$ も $A_i \to \mathcal{O}_i$ の解である。

\medskip\noindent
先ほど構成した拡大 $L/K$ に関する図式
(\ref{more-morphisms-equation-normalized-base-change}) を考える。$U_B \subset X_B$
は $B$ 上滑らかであり、したがって正規である（例えば Algebra, 補題
\ref{algebra-lemma-normal-goes-up} を用いる）。よって $Y \to X_B$ は
$U_B$ 上同型である。
$Y \to \Spec(B)$ のファイバーの既約成分の一般点で、極大イデアル
$\mathfrak m \subset B$ の上にある点 $y \in Y$ を取る。$y \not \in U_B$
と仮定する。このとき $y$ は点 $x_i$ のいずれかに写る。したがって
$\mathcal{O}_{Y, y}$ は $R(X) \otimes_K L$ における
$\mathcal{O}_i$ の整閉包の局所環である（詳細は省略する）。
$L/K$ は $A_i \to \mathcal{O}_i$ の解なので、
$B_\mathfrak m \to \mathcal{O}_{Y, y}$ は
$\mathfrak m_y$-進位相で形式的に滑らかであることが分かる
（これが「解」であることの定義である）。
言い換えると $\mathfrak m\mathcal{O}_{Y, y} = \mathfrak m_y$ であり、
剰余体拡大は分離的である。More on Algebra, 補題
\ref{more-algebra-lemma-extension-dvrs-formally-smooth} を参照せよ。
したがって $y$ におけるファイバーの局所環は $\kappa(y)$ である。
これは、例えば Algebra, 補題 \ref{algebra-lemma-separable-smooth} により、
$y$ におけるファイバーが $\kappa(\mathfrak m)$ 上滑らかであることを意味する。
これで証明は終わる。
\end{proof}

\begin{lemma}[曲線上の変種]
\label{more-morphisms-lemma-normalized-base-change-with-reduced-fibre-over-curve}
スキームの平坦有限型射 $f : X \to S$ を取る。
$S$ が関数体 $K$ をもつ Nagata 整スキームで、次元 $1$ の
正則スキームであると仮定する。このとき、図式
$$
\xymatrix{
Y \ar[rd]_g \ar[r]_-\nu & X \times_S T \ar[d] \ar[r] & X \ar[d]_f \\
& T \ar[r] & S
}
$$
において射 $g$ がすべてのファイバーの一般点で滑らかとなるような
有限拡大 $L/K$ が存在する。ここで $T$ は $\Spec(L)$ における
$S$ の正規化であり、$\nu : Y \to X \times_S T$ は正規化である。
\end{lemma}

\begin{proof}
有限アフィン開被覆 $S = \bigcup \Spec(A_i)$ を選ぶ。
すべての $i$ について $K$ は $A_i$ の分数体に等しい。
$X_i = X \times_S \Spec(A_i)$ と置く。
射 $X_i \to \Spec(A_i)$ に対して、定理
\ref{more-morphisms-theorem-normalized-base-change-with-reduced-fibre} のように
$L_i/K$ を選ぶ。$B_i \subset L_i$ を $A_i$ の整閉包とし、
$Y_i$ を $B_i$ への $X$ の正規化基底変換とする。
各 $L_i$ を支配する有限拡大 $L/K$ を取る。
$T_i \subset T$ を $\Spec(A_i)$ の逆像とする。
各 $i$ に対して可換図式
$$
\xymatrix{
g^{-1}(T_i) \ar[r] \ar[d] & Y_i \ar[r] \ar[d] & X \times_S \Spec(A_i) \ar[d] \\
T_i \ar[r] & \Spec(B_i) \ar[r] & \Spec(A_i)
}
$$
を得る。実際、左側の正方形はこの節の冒頭で議論した
正規化基底変換である。定理
\ref{more-morphisms-theorem-normalized-base-change-with-reduced-fibre} の証明で見たように、
$Y \to T$ の滑らかな軌跡は、$Y_i$ が $B_i$ 上滑らかとなる点の集合の
$g^{-1}(T_i)$ における逆像を含む。これで補題が従う。
\end{proof}

\begin{lemma}[分離拡大を伴う変種]
\label{more-morphisms-lemma-normalized-base-change-with-reduced-fibre-separable}
$A$ を分数体 $K$ をもつ Dedekind 環とする。
$X$ を $A$ 上平坦かつ有限型のスキームとする。
$A$ が Nagata 環で、$X$ の既約成分の各一般点
$\eta$ に対して体拡大 $\kappa(\eta)/K$ が分離的であると仮定する。
このとき、正規化基底変換 $Y$ がすべてのファイバーのすべての一般点で
$\Spec(B)$ 上滑らかとなるような有限分離拡大 $L/K$ が存在する。
\end{lemma}

\begin{proof}
これはいくつかの小さな変更を除けば、定理
\ref{more-morphisms-theorem-normalized-base-change-with-reduced-fibre} とまったく同じ方法で
証明される。最も重要な変更は More on Algebra, 命題
\ref{more-algebra-proposition-epp-essentially-finite-type} の代わりに
More on Algebra, 補題
\ref{more-algebra-lemma-epp-essentially-finite-type-separable} を用いることである。
証明では、$X \to \Spec(A)$ のような（平坦かつ有限表示の）射が
滑らかとなる点の集合の形成は基底変換と可換することを繰り返し用いる。
Morphisms, 補題 \ref{morphisms-lemma-set-points-where-fibres-smooth} を参照せよ。

\medskip\noindent
$X$ は $A$ 上平坦なので、$X$ の任意の一般点 $\eta$ は
$\Spec(A)$ の一般点に写る。
$X$ をその被約化で置き換えることにより、$X$ は被約であると仮定してよい。
この場合 Varieties, 補題
\ref{varieties-lemma-generic-points-geometrically-reduced} により
$X_K$ は $K$ 上幾何学的に被約である。
したがって Varieties, 補題
\ref{varieties-lemma-geometrically-reduced-dense-smooth-open} により、
$X_K \to \Spec(K)$ は稠密開集合上滑らかである。
よって射 $X \to \Spec(A)$ の滑らかな軌跡 $U \subset X$ は
(Morphisms, 定義 \ref{morphisms-definition-smooth} により) 開であり、
一般ファイバーに稠密である。これは次の段落の状況に帰着させる。

\medskip\noindent
$X$ が正規で、$X \to \Spec(A)$ の滑らかな軌跡 $U \subset X$ が
一般ファイバーに稠密であると仮定する。これは $U$ が有限個を除く
すべてのファイバーに稠密であることを意味する。補題
\ref{more-morphisms-lemma-nowhere-dense-generic-fibre} を参照せよ。
$x_1, \ldots, x_r \in X \setminus U$ を、$X \setminus U$ の既約成分の
一般点であり、さらに $X \to \Spec(A)$ のファイバーの既約成分の
一般点でもある有限個の点とする。
$\mathcal{O}_i = \mathcal{O}_{X, x_i}$ と置く。
$\mathcal{O}_i$ の分数体は $X$ の一般点の剰余体であることに注意する。
$A_i$ を $x_i$ の $\Spec(A)$ における像に対応する極大イデアルでの
$A$ の局所化とする。More on Algebra, 補題
\ref{more-algebra-lemma-epp-essentially-finite-type-separable} を適用すると、
$A_i \to \mathcal{O}_i$ の解となる有限分離拡大 $K_i/K$ が得られる。
すべての拡大 $K_i/K$ を支配する有限分離拡大 $L/K$ を取る。
More on Algebra, 補題 \ref{more-algebra-lemma-solution-goes-up} により、
$L/K$ はなお $A_i \to \mathcal{O}_i$ の解である。

\medskip\noindent
先ほど得た拡大 $L/K$ に関する図式
(\ref{more-morphisms-equation-normalized-base-change}) を考える。$U_B \subset X_B$
は $B$ 上滑らかであり、したがって正規である（例えば Algebra, 補題
\ref{algebra-lemma-normal-goes-up} を用いる）。よって $Y \to X_B$ は
$U_B$ 上同型である。
$Y \to \Spec(B)$ のファイバーの既約成分の一般点で、極大イデアル
$\mathfrak m \subset B$ の上にある点 $y \in Y$ を取る。$y \not \in U_B$
と仮定する。このとき $y$ は点 $x_i$ のいずれかに写る。したがって
$\mathcal{O}_{Y, y}$ は $R(X) \otimes_K L$ における
$\mathcal{O}_i$ の整閉包の局所環である（詳細は省略する）。
$L/K$ は $A_i \to \mathcal{O}_i$ の解なので、
$B_\mathfrak m \to \mathcal{O}_{Y, y}$ は形式的に滑らかである
（これが「解」であることの定義である）。
言い換えると $\mathfrak m\mathcal{O}_{Y, y} = \mathfrak m_y$ であり、
剰余体拡大は分離的である。したがって $y$ におけるファイバーの
局所環は $\kappa(y)$ である。これは、例えば Algebra, 補題
\ref{algebra-lemma-separable-smooth} により、$y$ におけるファイバーが
$\kappa(\mathfrak m)$ 上滑らかであることを意味する。
これで証明は終わる。
\end{proof}

\begin{lemma}[曲線上の分離拡大を伴う変種]
\label{more-morphisms-lemma-normalized-base-change-with-reduced-fibre-over-curve-separable}
スキームの平坦有限型射 $f : X \to S$ を取る。
$S$ が関数体 $K$ をもつ Nagata 整スキームで、次元 $1$ の
正則スキームであると仮定する。$X$ の既約成分の各一般点 $\eta$ に対して
体拡大 $\kappa(\eta)/K$ が分離的であると仮定する。このとき、図式
$$
\xymatrix{
Y \ar[rd]_g \ar[r]_-\nu & X \times_S T \ar[d] \ar[r] & X \ar[d]_f \\
& T \ar[r] & S
}
$$
において射 $g$ がすべてのファイバーの一般点で滑らかとなるような
有限分離拡大 $L/K$ が存在する。ここで $T$ は $\Spec(L)$ における
$S$ の正規化であり、$\nu : Y \to X \times_S T$ は正規化である。
\end{lemma}

\begin{proof}
これは補題
\ref{more-morphisms-lemma-normalized-base-change-with-reduced-fibre-separable} から、
補題 \ref{more-morphisms-lemma-normalized-base-change-with-reduced-fibre-over-curve} が
定理 \ref{more-morphisms-theorem-normalized-base-change-with-reduced-fibre} から従うのと
まったく同じ方法で従う。
\end{proof}








\section{ind-準アフィン射}
\label{more-morphisms-section-ind-quasi-affine}

\noindent
後で用いる理論を少し準備する。

\begin{definition}
\label{more-morphisms-definition-ind-quasi-affine}
スキーム $X$ が {\it ind-準アフィン} であるとは、$X$ の任意の
準コンパクト開部分スキームが準アフィンであることをいう。同様に、
スキームの射 $X \to Y$ が {\it ind-準アフィン} であるとは、
$Y$ の各アフィン開集合 $V$ に対して $f^{-1}(V)$ が
ind-準アフィンであることをいう。
\end{definition}

\noindent
アフィンスキームの開部分スキームは ind-準アフィンスキームの例である。
$X = \bigcup_{i \in I} U_i$ が準アフィン開部分スキームの合併で、
任意の二つの $U_i$ が第三のものに含まれるならば、$X$ は
ind-準アフィンである。ind-準アフィンスキーム $X$ は分離的である。
実際、任意の二つのアフィン開集合 $U, V$ は $X$ の分離的な開部分スキーム、
すなわち $U \cup V$ に含まれる。同様に ind-準アフィン射は分離的である。

\begin{lemma}
\label{more-morphisms-lemma-ind-quasi-affine-alternative-definition}
スキームの射 $f : X \to Y$ に対して、次は同値である。
\begin{enumerate}
\item $f$ は ind-準アフィンである。
\item 任意のアフィン開部分スキーム $V \subset Y$ と、任意の準コンパクト開部分スキーム
$U \subset f^{-1}(V)$ に対して、誘導射 $U \to V$ は準アフィンである。
\item $Y$ の準コンパクトかつ準分離な開部分スキーム
$V_j \subset Y$ による被覆 $\{ V_j \}_{j \in J}$ であって、
 すべての $j \in J$ と、任意の準コンパクト開部分スキーム
$U \subset f^{-1}(V_j)$ に対して、誘導射 $U \to V_j$ が準アフィンとなるものが存在する。
\item 任意の準コンパクトかつ準分離な開部分スキーム
 $V \subset Y$ と、任意の準コンパクト開部分スキーム
$U \subset f^{-1}(V)$ に対して、誘導射 $U \to V$ は準アフィンである。
\end{enumerate}
特に、ind-準アフィン射であるという性質は基底上 Zariski 局所的である。
\end{lemma}

\begin{proof}
(1) $\Leftrightarrow$ (2) の同値性は、定義と Morphisms, 補題
\ref{morphisms-lemma-characterize-quasi-affine} から従う。
(2) $\Rightarrow$ (4) について、(4) のような $U$, $V$ を取る。
Schemes, 補題 \ref{schemes-lemma-quasi-compact-permanence} により
誘導射 $U \to V$ は準コンパクトである。したがって任意のアフィン開集合
$V' \subset V$ に対してファイバー積 $V' \times_V U$ は準コンパクトであり、
(2) により誘導写像 $V' \times_V U \to V'$ は準アフィンである。
よって Morphisms, 補題 \ref{morphisms-lemma-characterize-quasi-affine} により
$U \to V$ も準アフィンである。
この議論は (3) $\Rightarrow$ (4) も与える。実際、同じ記号を用いると、
 ある $V_j$ に含まれる V のアフィン開集合 $V' \subset V$ が $V$ を被覆する。
したがって準コンパクト写像 $V' \times_V U \to V'$ が準アフィンであることを
示せばよい。しかし (3) により合成
$V' \times_V U \to V' \to V_j$ は準アフィンであり、さらに
Schemes, 補題 \ref{schemes-lemma-compose-after-separated} により
$V' \to V_j$ は準分離である。したがって Morphisms, 補題
\ref{morphisms-lemma-quasi-affine-permanence} により
$V' \times_V U \to V'$ は準アフィンである。
(4) $\Rightarrow$ (2) と (4) $\Rightarrow$ (3) は明らかである。
\end{proof}

\begin{lemma}
\label{more-morphisms-lemma-ind-quasi-affine-composition}
ind-準アフィン射であるという性質は合成で保たれる。
\end{lemma}

\begin{proof}
$f : X \to Y$ と $g : Y \to Z$ を ind-準アフィン射とする。
$V \subset Z$ と準コンパクト開集合
$U \subset f^{-1}(g^{-1}(V))$ を取り、$V$ はさらに準分離とする。
$f(U)$ は $g^{-1}(V)$ の準コンパクト部分集合なので、ある準コンパクト開集合
$W \subset g^{-1}(V)$（有限個のアフィンの合併）に含まれる。
分解 $U \to W \to V$ を得る。補題
\ref{more-morphisms-lemma-ind-quasi-affine-alternative-definition} により $W \to V$ は
準アフィンなので、特に $W$ は準分離である。再び補題
\ref{more-morphisms-lemma-ind-quasi-affine-alternative-definition} により
$U \to W$ も準アフィンである。したがって Morphisms, 補題
\ref{morphisms-lemma-composition-quasi-affine} により合成
$U \to V$ も準アフィンであり、最後にもう一度補題
\ref{more-morphisms-lemma-ind-quasi-affine-alternative-definition} を適用すればよい。
\end{proof}

\begin{lemma}
\label{more-morphisms-lemma-ind-quasi-affine-examples}
任意の準アフィン射は ind-準アフィンである。
任意の埋め込みは ind-準アフィンである。
\end{lemma}

\begin{proof}
第一の主張は定義から直ちに従う。特に閉埋め込みのようなアフィン射は
ind-準アフィンである。したがって補題
\ref{more-morphisms-lemma-ind-quasi-affine-composition} により、開埋め込みが
ind-準アフィンであることを示せばよい。しかしこれは定義から直ちに従う。
\end{proof}

\begin{lemma}
\label{more-morphisms-lemma-ind-quasi-affine-permanence}
スキームの射 $f : X \to Y$, $g : Y \to Z$ に対して、
$g \circ f$ が ind-準アフィンならば、$f$ は ind-準アフィンである。
\end{lemma}

\begin{proof}
補題 \ref{more-morphisms-lemma-ind-quasi-affine-alternative-definition} により、
$Z$ 上、次いで $Y$ 上 Zariski 局所的に考えてよい。したがって
$Z$、さらに $Y$ もアフィンであると仮定して一般性を失わない。
このとき $X$ の任意の準コンパクト開集合は準アフィンなので、
補題 \ref{more-morphisms-lemma-ind-quasi-affine-alternative-definition} が主張を与える。
\end{proof}

\begin{lemma}
\label{more-morphisms-lemma-base-change-ind-quasi-affine}
ind-準アフィンであるという性質は基底変換で保たれる。
\end{lemma}

\begin{proof}
ind-準アフィン射 $f : X \to Y$ を取る。
$f$ の任意の基底変換が ind-準アフィンであることを確認するには、
補題 \ref{more-morphisms-lemma-ind-quasi-affine-alternative-definition} により
$Y$ 上 Zariski 局所的に考えてよいので、$Y$ はアフィンと仮定する。
さらに基底変換射 $Z \to Y$ において $Z$ もアフィンと仮定してよい。
基底変換 $X \times_Y Z \to X$ はアフィン射なので、補題
\ref{more-morphisms-lemma-ind-quasi-affine-composition} と
\ref{more-morphisms-lemma-ind-quasi-affine-examples} により写像
$X \times_Y Z \to Y$ は ind-準アフィンである。次に補題
\ref{more-morphisms-lemma-ind-quasi-affine-permanence} により、基底変換
$X \times_Y Z \to Z$ は ind-準アフィンとなり、主張が従う。
\end{proof}

\begin{lemma}
\label{more-morphisms-lemma-descending-property-ind-quasi-affine}
ind-準アフィンであるという性質は基底上 fpqc 局所的である。
\end{lemma}

\begin{proof}
補題 \ref{more-morphisms-lemma-base-change-ind-quasi-affine} による基底変換での
ind-準アフィン性の安定性が一方向を与える。逆に、スキームの射
 $f : X \to Y$ と、基底変換 $f_i : X_i \to Y_i$ が
 すべての $i$ について ind-準アフィンとなる fpqc 被覆
 $\{g_i : Y_i \to Y\}$ を取る。$f$ が ind-準アフィンであることを示す。

\medskip\noindent
補題 \ref{more-morphisms-lemma-ind-quasi-affine-alternative-definition} により
$Y$ 上 Zariski 局所的に考えてよいので、$Y$ はアフィンと仮定する。
次に補題 \ref{more-morphisms-lemma-base-change-ind-quasi-affine} が保証する基底変換での
安定性を用いて被覆を細分し、単一のアフィン忠実平坦射
$g : Y' \to Y$ からなると仮定する。任意の準コンパクト開集合
$U \subset X$ に対して、その $Y'$-基底変換
$U \times_Y Y' \subset X \times_Y Y'$ も準コンパクトである。
残るのは、Descent, 補題
\ref{descent-lemma-descending-property-quasi-affine} により、
写像 $U \to Y$ が準アフィンであることと
$U \times_Y Y' \to Y'$ が準アフィンであることが同値だと確認することである。
\end{proof}

\begin{lemma}
\label{more-morphisms-lemma-etale-separated-ind-quasi-affine}
スキームの分離的な局所有限準有限射は ind-準アフィンである。
\end{lemma}

\begin{proof}
$f : X \to Y$ をスキームの分離的な局所有限準有限射とする。
$V \subset Y$ をアフィンとし、$U \subset f^{-1}(V)$ を準コンパクトな
開集合とする。$U$ が準アフィンであることを示せばよい。$U \to V$ は
スキームの分離的な準有限射なので、これは Zariski の主定理から従う。
補題 \ref{more-morphisms-lemma-quasi-finite-separated-quasi-affine} を参照。
\end{proof}




\section{スキームの圏における押し出し II}
\label{more-morphisms-section-pushouts-II}

\noindent
本節は第 \ref{more-morphisms-section-pushouts} 節の続きである。本節では、$Z \to X$ が
閉埋め込み、$Z \to Y$ が整射であり、さらに一つの条件が満たされる場合に、
$Y \leftarrow Z \rightarrow X$ の押し出しを構成する。詳しい議論については
\cite{Ferrand-Conducteur} を参照されたい。

\begin{situation}
\label{more-morphisms-situation-pushout-along-closed-immersion-and-integral}
ここで $S$ はスキームであり、$i : Z \to X$ と $j : Z \to Y$ は
$S$ 上のスキームの射である。次を仮定する。
\begin{enumerate}
\item $i$ は閉埋め込みである。
\item $j$ はスキームの整射である。
\item $y \in Y$ に対して、$j^{-1}(\{y\}) \subset i^{-1}(U)$ を満たす
アフィン開集合 $U \subset X$ が存在する。
\end{enumerate}
\end{situation}

\begin{lemma}
\label{more-morphisms-lemma-prepare-pushout-along-closed-immersion-and-integral}
状況 \ref{more-morphisms-situation-pushout-along-closed-immersion-and-integral} において、
$y \in Y$ に対し、$i^{-1}(U) = j^{-1}(V)$ かつ $y \in V$ を満たす
アフィン開集合 $U \subset X$ と $V \subset Y$ が存在する。
\end{lemma}

\begin{proof}
$y \in Y$ とする。$j^{-1}(\{y\}) \subset i^{-1}(U)$ を満たすアフィン開集合
$U \subset X$ を選ぶ（仮定により可能である）。$y$ のアフィン開近傍
$V \subset Y$ で $j^{-1}(V) \subset i^{-1}(U)$ を満たすものを選ぶ。
$j : Z \to Y$ は閉射であり（Morphisms, 補題
\ref{morphisms-lemma-integral-universally-closed}）、$i^{-1}(U)$ は $y$ 上の
ファイバーを含むので、これは可能である。$j$ は整射だから、スキーム論的
ファイバー $Z_y$ は体上整な代数のスペクトルである。Limits, 補題
\ref{limits-lemma-ample-profinite-set-in-principal-affine} により、
$Z_y \subset D(\overline{f}) \subset j^{-1}(V)$ を満たす
$\overline{f} \in \Gamma(i^{-1}(U), \mathcal{O}_{i^{-1}(U)})$ が取れる。
$i|_{i^{-1}(U)} : i^{-1}(U) \to U$ はアフィン間の閉埋め込みなので、
$i^{-1}(U)$ への制限が $\overline{f}$ となる
$f \in \Gamma(U, \mathcal{O}_U)$ を選べる。$U$ を主開集合 $D(f) \subset U$
で置き換えると、アフィン開集合 $y \in V \subset Y$ と $U \subset X$ で
$$
j^{-1}(\{y\}) \subset i^{-1}(U) \subset j^{-1}(V)
$$
を満たすものを得る。ここで（ある意味で）同じ議論を繰り返す。すなわち、
$y \in D(g)$ かつ $j^{-1}(D(g)) \subset i^{-1}(U)$ を満たす
$g \in \Gamma(V, \mathcal{O}_V)$ を選ぶ（上と同じ議論で可能である）。次に、
$i^{-1}(U) \to V$ による $g$ の引き戻しを $i^{-1}(U)$ への制限にもつ
$f \in \Gamma(U, \mathcal{O}_U)$ を選べる（これも上と同じ理由による）。
したがって最終的に、アフィン開集合 $y \in V' = D(g) \subset V \subset Y$ と
$U' = D(f) \subset U \subset X$ で次を満たすものを得る。
% P05-EN-CH37-0319: frozen source has the ill-typed right-hand inverse image i^{-1}(V').
$j^{-1}(V') = i^{-1}(V')$.
\end{proof}

\begin{proposition}
\label{more-morphisms-proposition-pushout-along-closed-immersion-and-integral}
\begin{reference}
\cite[Theorem 7.1 part iii]{Ferrand-Conducteur}
\end{reference}
状況 \ref{more-morphisms-situation-pushout-along-closed-immersion-and-integral} において、
押し出し $Y \amalg_Z X$ はスキームの圏で存在する。図式は
$$
\xymatrix{
Z \ar[r]_i \ar[d]_j & X \ar[d]^a \\
Y \ar[r]^-b & Y \amalg_Z X
}
$$
である。この図式はファイバー平方であり、射 $a$ は整射、射 $b$ は
閉埋め込みであり、
$$
\mathcal{O}_{Y \amalg_Z X} =
b_*\mathcal{O}_Y \times_{c_*\mathcal{O}_Z} a_*\mathcal{O}_X
$$
が環の層として成り立つ。ここで $c = a \circ i = b \circ j$ である。
\end{proposition}

\begin{proof}
位相空間としての $Y \amalg_Z X$ を、位相空間の圏におけるこの図式の
押し出しと定める（Topology, 第 \ref{topology-section-colimits} 節）。これは
台集合の押し出しに商位相を入れたものにほかならない（Topology, 補題
\ref{topology-lemma-colimits}）。$Y \amalg_Z X$ 上には環の層の射
$$
b_*\mathcal{O}_Y \longrightarrow c_*\mathcal{O}_Z
\longleftarrow a_*\mathcal{O}_X
$$
があり、
$$
\mathcal{O}_{Y \amalg_Z X} =
b_*\mathcal{O}_Y \times_{c_*\mathcal{O}_Z} a_*\mathcal{O}_X
$$
を環の層の圏におけるファイバー積として定義できる。これによってスキームが
得られることを証明するには、各点がアフィン開近傍をもつことを示せばよい。
$c$ の像は閉部分集合であり、その補集合は環付き空間として
$(Y \setminus j(Z)) \amalg (X \setminus i(Z))$ と同型なので、$c$ の像に
含まれない点については明らかである。

\medskip\noindent
$c$ の像の点は、$j$ の像にある一意な $y \in Y$ に対応する。補題
\ref{more-morphisms-lemma-prepare-pushout-along-closed-immersion-and-integral} により、
$y \in V$ かつ $i^{-1}(U) = j^{-1}(V)$ を満たすアフィン開集合
$U \subset X$ と $V \subset Y$ を得る。第一段落の構成は整合する開部分スキーム
への制限と明らかに両立するので、それがスキームを与えることを示す際には
$X$, $Y$, $Z$ がアフィンであると仮定してよい。

\medskip\noindent
$X = \Spec(A)$, $Y = \Spec(B)$, $Z = \Spec(C)$ がアフィンならば、
More on Algebra, 補題 \ref{more-algebra-lemma-points-of-fibre-product} により、
位相空間として $Y \amalg_Z X = \Spec(B \times_C A)$ である。
% P05-EN-CH37-0320: frozen source writes the ill-typed fibre product Y \times_Z X here.
$Y \times_Z X$ がスキームであるという証明を完結するには、
$\Spec(B \times_C A)$ 上の構造層が順像のファイバー積であることを示せばよい。
これは $\Spec(B \times_C A)$ の主アフィン開集合に More on Algebra, 補題
\ref{more-algebra-lemma-diagram-localize} を適用することから従う。

\medskip\noindent
上の議論から、スキーム $Y \amalg_Z X$ は、$U_i = a^{-1}(W_i)$,
$V_i = b^{-1}(W_i)$, $\Omega_i = c^{-1}(W_i)$ がそれぞれ $X$, $Y$, $Z$ の
アフィン開集合となるアフィン開被覆 $Y \amalg_Z X = \bigcup W_i$ をもつ。
したがって $a$ と $b$ はアフィンである。さらに $A_i$, $B_i$, $C_i$ を
$U_i$, $V_i$, $\Omega_i$ に対応する環とすると、$A_i \to C_i$ は全射であり、
$W_i$ は $B_i$ へ全射する $A_i \times_{C_i} B_i$ に対応する。ゆえに $b$ は
閉埋め込みである。環準同型 $A_i \times_{C_i} B_i \to A_i$ は More on Algebra,
補題 \ref{more-algebra-lemma-fibre-product-integral} により整なので、$a$ は
整射である。また、図式が Cartesian であるのは
$$
C_i \cong B_i \otimes_{B_i \times_{C_i} A_i} A_i
$$
が成り立つからである。実際、$B_i \times_{C_i} A_i \to B_i$ と
$A_i \to C_i$ は全射であり、その核は同じである。

\medskip\noindent
最後に、補題 \ref{more-morphisms-lemma-basic-example-pushout} と
\ref{more-morphisms-lemma-pushout-fpqc-local} を適用すれば、この構成がスキームの圏における
押し出しであると結論できる。
\end{proof}

\begin{lemma}
\label{more-morphisms-lemma-pushout-separated}
% P05-EN-CH37-0321: frozen source ends the opening situation phrase with a stray period.
状況 \ref{more-morphisms-situation-pushout-along-closed-immersion-and-integral} において、
$X$ と $Y$ が分離的ならば、押し出し $Y \amalg_Z X$（命題
\ref{more-morphisms-proposition-pushout-along-closed-immersion-and-integral}）も分離的である。
「$S$ 上分離的」、「準分離的」、「$S$ 上準分離的」についても同様である。
\end{lemma}

\begin{proof}
射 $Y \amalg X \to Y \amalg_Z X$ は全射かつ普遍閉である。したがって
Morphisms, 補題 \ref{morphisms-lemma-image-universally-closed-separated} を
適用できる。
\end{proof}

\begin{lemma}
\label{more-morphisms-lemma-pushout-finite-type}
状況 \ref{more-morphisms-situation-pushout-along-closed-immersion-and-integral} において、
$S$ は局所 Noether スキームであり、$X$, $Y$, $Z$ は $S$ 上局所有限型であると
仮定する。このとき押し出し $Y \amalg_Z X$（命題
\ref{more-morphisms-proposition-pushout-along-closed-immersion-and-integral}）は $S$ 上
局所有限型である。
\end{lemma}

\begin{proof}
アフィン開集合上で考えれば、More on Algebra, 補題
\ref{more-algebra-lemma-fibre-product-finite-type} の結果に帰着する。
\end{proof}

\begin{lemma}
\label{more-morphisms-lemma-pushout-functor}
状況 \ref{more-morphisms-situation-pushout-along-closed-immersion-and-integral} において、
可換図式
$$
\xymatrix{
Y' \ar[d]^g & Z' \ar[l]^{j'} \ar[r]_{i'} \ar[d]^h & X' \ar[d]^f \\
Y & Z \ar[l] \ar[r] & X
}
$$
が与えられ、二つの平方は Cartesian であり、$f, g, h$ は分離的かつ
局所準有限であるとする。このとき
\begin{enumerate}
\item 押し出し $Y \amalg_Z X$ と $Y' \amalg_{Z'} X'$ が存在する。
\item $Y' \amalg_{Z'} X' \to Y \amalg_Z X$ は分離的かつ局所準有限である。
\item 次の二つの平方
$$
\xymatrix{
Y' \ar[r] \ar[d] & Y' \amalg_{Z'} X' \ar[d] & X' \ar[l] \ar[d] \\
Y \ar[r] & Y \amalg_Z X & X \ar[l]
}
$$
は Cartesian である。
\end{enumerate}
\end{lemma}

\begin{proof}
押し出し $Y \amalg_Z X$ は命題
\ref{more-morphisms-proposition-pushout-along-closed-immersion-and-integral} により存在する。
押し出し $Y' \amalg_{Z'} X'$ の存在を示すため、状況
\ref{more-morphisms-situation-pushout-along-closed-immersion-and-integral} の条件 (3) が
$(X', Y', Z', i', j')$ に対して成り立つことを確認する。$y' \in Y'$ とし、
その像を $y \in Y$ と書く。$i(j^{-1}(y)) \subset U$ を満たすアフィン開集合
$U \subset X$ を選ぶ。$f^{-1}(U)$ に含まれ、準コンパクト部分集合
$i'((j')^{-1}(\{y'\}))$ を含む準コンパクト開集合 $U' \subset X'$ を選ぶ。
補題 \ref{more-morphisms-lemma-etale-separated-ind-quasi-affine} により $U'$ は準アフィンである。
$Z'_{y'}$ は体上整な代数のスペクトルなので、Limits, 補題
\ref{limits-lemma-ample-profinite-set-in-principal-affine} を適用でき、望みどおり
$i'((j')^{-1}(\{y'\}))$ を含む $U'$ のアフィン開部分スキームが存在する。

\medskip\noindent
存在を確認したので、残りの主張を調べる。アフィン局所的には、$B \to D$ と
$A' \to C'$ が局所準有限である More on Algebra, 補題
\ref{more-algebra-lemma-properties-algebras-over-fibre-product} の状況そのものである
\footnote{正確には、$X, Y, Z, Y \amalg_Z X, X', Y', Z', Y' \amalg_{Z'} X'$
はそれぞれ $A', B, A, B', C', D, C, D'$ に対応する。}。特に、射
$Y' \amalg_{Z'} X' \to Y \amalg_Z X$ は局所有限型である。
% P05-EN-CH37-0322: frozen source has the malformed phrase "squares in of the diagram".
図式の二つの平方は More on Algebra, 補題
\ref{more-algebra-lemma-module-over-fibre-product} により Cartesian である。
局所準有限性はファイバー上で確認できるので（Morphisms, 補題
\ref{morphisms-lemma-quasi-finite-at-point-characterize}）、
$Y' \amalg_{Z'} X' \to Y \amalg_Z X$ は局所準有限であると結論できる。

\medskip\noindent
なお $Y' \amalg_{Z'} X' \to Y \amalg_Z X$ が分離的であることを確認する必要が
ある。命題 \ref{more-morphisms-proposition-pushout-along-closed-immersion-and-integral} により、
$Y' \amalg X' \to Y' \amalg_{Z'} X'$ は普遍閉かつ全射である。また、射
$Y' \amalg X' \to Y \amalg_Z X$ も分離的である（これは
$Y' \amalg X' \to Y \amalg X \to Y \amalg_Z X$ と分解する）。したがって
Morphisms, 補題 \ref{morphisms-lemma-image-universally-closed-separated} により
結論が従う。
\end{proof}

\begin{lemma}
\label{more-morphisms-lemma-pushout-functor-equivalence-flat}
状況 \ref{more-morphisms-situation-pushout-along-closed-immersion-and-integral} において、
押し出し $Y \amalg_Z X$ 上平坦、分離的かつ局所準有限なスキームの圏は、
補題 \ref{more-morphisms-lemma-pushout-functor} のような $(X', Y', Z', i', j', f, g, h)$ で
$f, g, h$ が平坦であるものの圏と同値である。「平坦」を「\'etale」で
置き換えても同様である。
\end{lemma}

\begin{proof}
補題 \ref{more-morphisms-lemma-pushout-functor} のような $(X', Y', Z', i', j', f, g, h)$ から
始め、$f, g, h$ が平坦または \'etale ならば、More on Algebra, 補題
\ref{more-algebra-lemma-properties-algebras-over-fibre-product} により
$Y' \amalg_{Z'} X' \to Y \amalg_Z X$ は平坦または \'etale である。

\medskip\noindent
逆に、$W \to Y \amalg_Z X$ を分離的かつ局所準有限な射とする。
$X' = W \times_{Y \amalg_Z X} X$, $Y' = W \times_{Y \amalg_Z X} Y$,
$Z' = W \times_{Y \amalg_Z X} Z$ とおき、明らかな射 $i', j', f, g, h$ を
考える。押し出し $Y' \amalg_{Z'} X'$ を作る。押し出しの普遍性により、
$$
Y' \amalg_{Z'} X' \longrightarrow W
$$
% P05-EN-CH37-0323: frozen source names the malformed base Y \amalg_X Z.
$Y \amalg_X Z$ 上のスキームの射を得る。$W \to Y \amalg_Z X$ が平坦であると
仮定しなければ、一般にこの射は同型ではない（実際、More on Algebra, 補題
\ref{more-algebra-lemma-module-over-fibre-product-bis} により、表示された矢印は
閉埋め込みではあるが、一般には同型ではない）。しかし、
% P05-EN-CH37-0324: frozen source writes the ill-typed fibre product Y \times_Z X.
$W \to Y \times_Z X$ が平坦ならば、More on Algebra, 補題
\ref{more-algebra-lemma-properties-algebras-over-fibre-product} により同型である。
\end{proof}

\noindent
次に、二つの射がともに閉埋め込みである場合の存在を論じる。

\begin{lemma}
\label{more-morphisms-lemma-pushout-along-closed-immersions}
$i : Z \to X$ と $j : Z \to Y$ をスキームの閉埋め込みとする。このとき
押し出し $Y \amalg_Z X$ はスキームの圏で存在する。図式は
$$
\xymatrix{
Z \ar[r]_i \ar[d]_j & X \ar[d]^a \\
Y \ar[r]^-b & Y \amalg_Z X
}
$$
である。この図式はファイバー平方であり、射 $a$ と $b$ は閉埋め込みで、
短完全列
$$
0 \to \mathcal{O}_{Y \amalg_Z X} \to
a_*\mathcal{O}_X \oplus b_*\mathcal{O}_Y \to
c_*\mathcal{O}_Z \to 0
$$
がある。ここで $c = a \circ i = b \circ j$ である。
\end{lemma}

\begin{proof}
これは命題 \ref{more-morphisms-proposition-pushout-along-closed-immersion-and-integral} の
特別な場合である。$j$ のファイバーは一点集合なので、状況
\ref{more-morphisms-situation-pushout-along-closed-immersion-and-integral} の仮定 (3) は
直ちに従う。最後に二つの矢印の役割を逆にすれば、$a$ と $b$ がともに
閉埋め込みであると結論できる。
\end{proof}

\begin{lemma}
\label{more-morphisms-lemma-pushout-along-closed-immersions-properties-above}
$i : Z \to X$ と $j : Z \to Y$ をスキームの閉埋め込みとする。
$f : X' \to X$ と $g : Y' \to Y$ をスキームの射とし、
$\varphi : X' \times_{X, i} Z \to Y' \times_{Y, j} Z$ を $Z$ 上の
スキームの同型とする。射
$$
h :
X' \amalg_{X' \times_{X, i} Z, \varphi} Y'
\longrightarrow
X \amalg_Z Y
$$
を考える。このとき次が成り立つ。
\begin{enumerate}
\item $h$ が局所有限型であることと、$f$ と $g$ が局所有限型であることは
同値である。
\item $h$ が平坦であることと、$f$ と $g$ が平坦であることは同値である。
\item $h$ が平坦かつ局所有限表示であることと、$f$ と $g$ が平坦かつ
局所有限表示であることは同値である。
\item $h$ が滑らかであることと、$f$ と $g$ が滑らかであることは同値である。
\item $h$ が \'etale であることと、$f$ と $g$ が \'etale であることは同値である。
% P05-EN-CH37-0325: frozen source retains an editorial placeholder as an enumerated item.
\item 必要に応じてここに項目を追加する。
\end{enumerate}
\end{lemma}

\begin{proof}
補題 \ref{more-morphisms-lemma-pushout-along-closed-immersions} により押し出しが存在することは
分かっている。特に射 $h$ を得る。したがって、現れるすべてのスキームを
アフィンスキームで置き換えてよい。この場合、補題の各主張は More on Algebra,
補題 \ref{more-algebra-lemma-properties-algebras-over-fibre-product} の対応する
主張と同値である。
\end{proof}





\section{相対射}
\label{more-morphisms-section-relative-morphisms}

\noindent
本節では表現可能性の結果を証明する。これは Fundamental Groups, 第
\ref{pione-section-finite-etale} 節で、スキームの有限 \'etale 被覆の圏に
関する結果を証明するために用いる。本節の内容は Criteria for Representability,
第 \ref{criteria-section-relative-morphisms} 節で適切な一般性のもとに論じられる。

\medskip\noindent
$S$ をスキームとし、$Z$ と $X$ を $S$ 上のスキームとする。$S$ 上のスキーム
$T$ が与えられたとき、$S$ 上の射 $b : T \times_S Z \to T \times_S X$ を
考えることができる。図式は
\begin{equation}
\label{more-morphisms-equation-hom}
\vcenter{
\xymatrix{
T \times_S Z \ar[rd] \ar[rr]_b & &
T \times_S X \ar[ld] & Z \ar[rd] & & X \ar[ld] \\
& T \ar[rrr] & & & S
}
}
\end{equation}
である。もちろん、$b$ を、次の図式を可換にする射
$b : T \times_S Z \to X$ と考えることもできる。
$$
\xymatrix{
T \times_S Z \ar[r] \ar[d] \ar@/^1pc/[rrr]_-b &
Z \ar[rd] & & X \ar[ld] \\
T \ar[rr] & & S
}
$$
この状況で関手
\begin{equation}
\label{more-morphisms-equation-hom-functor}
\mathit{Mor}_S(Z, X) : (\Sch/S)^{opp} \longrightarrow \textit{Sets},
\quad
T \longmapsto \{b\text{ as above}\}
\end{equation}
を定義できる。まず基本的な表現可能性の結果を述べる。

\begin{lemma}
\label{more-morphisms-lemma-hom-from-finite-free-into-affine}
$Z \to S$ と $X \to S$ をアフィンスキームの射とする。
$\Gamma(Z, \mathcal{O}_Z)$ が有限自由 $\Gamma(S, \mathcal{O}_S)$-加群であると
仮定する。このとき $\mathit{Mor}_S(Z, X)$ は $S$ 上のアフィンスキームで
表現可能である。
\end{lemma}

\begin{proof}
$S = \Spec(R)$ と書く。$\Gamma(Z, \mathcal{O}_Z)$ の $R$ 上の基底
$\{e_1, \ldots, e_m\}$ を選ぶ。また表示
$$
\Gamma(X, \mathcal{O}_X) = R[\{x_i\}_{i \in I}]/(\{f_k\}_{k \in K}).
$$
を選ぶ。この商における $x_i$ の像を $\overline{x}_i$ と書く。また
$$
P = R[\{a_{ij}\}_{i \in I, 1 \leq j \leq m}].
$$
とおき、$R$-代数準同型
$$
\Psi :
R[\{x_i\}_{i \in I}]
\longrightarrow
P \otimes_R \Gamma(Z, \mathcal{O}_Z), \quad
x_i \longmapsto \sum\nolimits_j a_{ij} \otimes e_j.
$$
を考える。$c_{kj} \in P$ を用いて
$\Psi(f_k) = \sum c_{kj} \otimes e_j$ と書く。最後に、$J \subset P$ を元
$c_{kj}$（$k \in K$, $1 \leq j \leq m$）が生成するイデアルとする。
$W = \Spec(P/J)$ が関手 $\mathit{Mor}_S(Z, X)$ を表現すると主張する。

\medskip\noindent
まず構成から $P/J$ は $R$-代数であり、したがって射 $W \to S$ がある。
次に、構成から写像 $\Psi$ は $\Gamma(X, \mathcal{O}_X)$ を経由するので、
% P05-EN-CH37-0326: frozen source has the article error "an P/J-algebra homomorphism".
$P/J$-代数準同型
$$
P/J \otimes_R \Gamma(X, \mathcal{O}_X)
\longrightarrow
P/J \otimes_R \Gamma(Z, \mathcal{O}_Z)
$$
を得る。これは射 $b_{univ} : W \times_S Z \to W \times_S X$ を定める。
Yoneda の補題により、$b_{univ}$ は関手の変換
$W \to \mathit{Mor}_S(Z, X)$ を定め、これは同型であると主張する。同型である
ことを示すには、任意のアフィンスキーム $T$ に対して集合の全単射
$W(T) \to \mathit{Mor}_S(Z, X)(T)$ を誘導することを示せば十分である。

\medskip\noindent
$T = \Spec(R')$ を $S$ 上のアフィンスキームとし、
$b \in \mathit{Mor}_S(Z, X)(T)$ とする。構造射 $T \to S$ は $R'$ に
$R$-代数構造を定め、$b$ は $R'$-代数準同型
$$
b^\sharp :
R' \otimes_R \Gamma(X, \mathcal{O}_X)
\longrightarrow
R' \otimes_R \Gamma(Z, \mathcal{O}_Z).
$$
を定める。特に、ある $\alpha_{ij} \in R'$ を用いて
$b^\sharp(1 \otimes \overline{x}_i) = \sum \alpha_{ij} \otimes e_j$ と書ける。
これは規則 $a_{ij} \mapsto \alpha_{ij}$ で定まる $R$-代数準同型
$P \to R'$ に対応する。イデアル $J$ の構成により、この写像は商 $P/J$ を
経由して写像 $P/J \to R'$ を与える。これはさらに射 $T \to W$ に対応し、
$b$ は $b_{univ}$ の引き戻しとなる。細部は省略する。
\end{proof}

\begin{lemma}
\label{more-morphisms-lemma-hom-from-finite-locally-free-into-affine}
$Z \to S$ と $X \to S$ をスキームの射とする。$Z \to S$ が有限局所自由で、
$X \to S$ がアフィンならば、$\mathit{Mor}_S(Z, X)$ は $S$ 上アフィンな
スキームで表現可能である。
\end{lemma}

\begin{proof}
$\Gamma(Z \times_S U_i, \mathcal{O}_{Z \times_S U_i})$ が
$\mathcal{O}_S(U_i)$ 上有限自由となるアフィン開被覆 $S = \bigcup U_i$ を選ぶ。
$F_i \subset \mathit{Mor}_S(Z, X)$ を、$T \to S$ が $U_i$ を経由しないとき
$T/S$ に空集合を対応させ、それ以外では $\mathit{Mor}_S(Z, X)(T)$ を
対応させる部分関手とする。
% P05-EN-CH37-0327: frozen source has subject-verb disagreement in "collection ... satisfy".
これらの部分関手の族は Schemes, 補題 \ref{schemes-lemma-glue-functors} の条件
(2)(a), (2)(b), (2)(c) を満たし、補題が従う。条件 (2)(a) は補題
\ref{more-morphisms-lemma-hom-from-finite-free-into-affine} から従い、残り二つは直接的な議論から
従う。
\end{proof}

\noindent
次の補題における射 $f : X \to S$ に関する条件は、この種の主張を証明する際に
非常に有用である。この条件は、次のいずれかが成り立てば満たされる：$X$ が
準アフィンである、$f$ が準アフィンである、$f$ が準射影的である、$f$ が
局所射影的である、$X$ 上に豊富な可逆層が存在する、$X$ 上に $f$-豊富な
可逆層が存在する、または $X$ 上に $f$-非常に豊富な可逆層が存在する。

\begin{lemma}
\label{more-morphisms-lemma-hom-from-finite-locally-free-representable}
$Z \to S$ と $X \to S$ をスキームの射とし、次を仮定する。
\begin{enumerate}
\item $Z \to S$ は有限局所自由である。
\item $s \in S$ かつ $x_1, \ldots, x_d \in X_s$ であるすべての
$(s, x_1, \ldots, x_d)$ に対して、$x_1, \ldots, x_d \in U$ を満たす
アフィン開集合 $U \subset X$ が存在する。
\end{enumerate}
このとき $\mathit{Mor}_S(Z, X)$ はスキームで表現可能である。
\end{lemma}

\begin{proof}
$U \subset X$ と $V \subset S$ がアフィン開集合で、$U \to S$ が $V$ を
経由するような組 $(U, V)$ の集合を $I$ とする。$i \in I$ に対し、対応する組を
$(U_i, V_i)$ と書き、$F_i = \mathit{Mor}_{V_i}(Z_{V_i}, U_i)$ とおく。
$F_i$ が $\mathit{Mor}_S(Z, X)$ の部分関手であることは直ちに分かる。
% P05-EN-CH37-0328: frozen source omits the predicate after "we claim that conditions".
Schemes, 補題 \ref{schemes-lemma-glue-functors} の条件 (2)(a), (2)(b), (2)(c) が
満たされると主張する。これにより補題が従う。

\medskip\noindent
条件 (2)(a) は補題
\ref{more-morphisms-lemma-hom-from-finite-locally-free-into-affine} から従う。

\medskip\noindent
条件 (2)(b) を確認するため、$T/S$ と $b \in \mathit{Mor}_S(Z, X)(T)$ を考える。
$b$ を射 $T \times_S Z \to X$ とみなすと、開集合
$b^{-1}(U_i) \subset T \times_S Z$ を得る。明らかに $b \in F_i(T)$ であることと
$b^{-1}(U_i) = T \times_S Z$ であることは同値である。射影
$p : T \times_S Z \to T$ は有限、したがって閉なので、
$p^{-1}(\{t\}) \subset b^{-1}(U_i)$ を満たす点 $t \in T$ の集合
$U_{i, b} \subset T$ は開である。
% P05-EN-CH37-0329: frozen source writes the ill-typed composition b \circ f without the induced base-change map.
このとき $f : T' \to T$ が $U_{i, b}$ を経由することと
$b \circ f \in F_i(T')$ であることは同値であり、条件 (2)(b) の確認が終わる。

\medskip\noindent
最後に条件 (2)(c) を確認する。ここで $X \to S$ に関する仮定を用いる。
$T/S$ と $b \in \mathit{Mor}_S(Z, X)(T)$ を考える。各 $t \in T$ が前段落で
定めた開集合 $U_{i, b}$ のいずれかに含まれることを示せば十分である。
これは、射影を $p : T \times_S Z \to T$、与えられた射を
$b : T \times_S Z \to X$ としたとき、ある $i$ に対して
$b(p^{-1}(\{t\})) \subset U_i$ となることと同値である。$p$ は有限なので、
集合 $b(p^{-1}(\{t\})) \subset X$ は有限であり、$t$ の $S$ における像
$s$ 上の $X \to S$ のファイバーに含まれる。したがって $X \to S$ に関する
仮定は、適切な組が存在することをちょうど示す。
\end{proof}

\begin{lemma}
\label{more-morphisms-lemma-hom-from-finite-locally-free-separated-lqf}
$Z \to S$ と $X \to S$ をスキームの射とする。$Z \to S$ は有限局所自由で、
$X \to S$ は分離的かつ局所準有限であると仮定する。このとき
$\mathit{Mor}_S(Z, X)$ はスキームで表現可能である。
\end{lemma}

\begin{proof}
これは補題 \ref{more-morphisms-lemma-hom-from-finite-locally-free-representable} と
\ref{more-morphisms-lemma-separated-locally-quasi-finite-over-affine} から従う。
\end{proof}









\section{擬連接複体の特徴づけ III}
\label{more-morphisms-section-characterize-pseudo-coherent}

\noindent
本節では、コホモロジーによる擬連接複体の特徴づけを論じる。これは
Derived Categories of Schemes, 第 \ref{perfect-section-pseudo-coherent} 節の
続きである。基本的な道具は、Chow の補題の導来版を用いて射影空間の場合へ
帰着することである。補題 \ref{more-morphisms-lemma-derived-chow} を参照。

\begin{lemma}
\label{more-morphisms-lemma-case-of-tor-independence}
スキームの可換図式
$$
\xymatrix{
Z' \ar[d] \ar[r] & Y' \ar[d] \\
X' \ar[r] & S'
}
$$
を考える。$S \to S'$ を射とし、$X'$ と $Y'$ の $S$ への基底変換をそれぞれ
$X$ と $Y$ と書く。$Y' \to S'$ と $Z' \to X'$ は平坦であると仮定する。
このとき $X \times_S Y$ と $Z'$ は $X' \times_{S'} Y'$ 上 Tor 独立である。
\end{lemma}

\begin{proof}
問題は局所的なので、すべてのスキームがアフィンであると仮定してよい
（細部は省略する）。図式
$$
\xymatrix{
X \times_S Y \ar[r] \ar[d] & X' \times_{S'} Y' \ar[d] \\
X \ar[r] & X'
}
$$
は Cartesian で、縦の矢印は平坦である。
$X = \Spec(A)$, $X' = \Spec(A')$,
$X' \times_{S'} Y' = \Spec(B')$ と書く。このとき
$X \times_S Y = \Spec(A \otimes_{A'} B')$ である。また
$Z' = \Spec(C')$ と書く。示すべきことは
$$
\text{Tor}_p^{B'}(A \otimes_{A'} B', C') = 0,
\quad\text{for } p > 0
$$
$A' \to B'$ は平坦なので
$A \otimes_{A'} B' = A \otimes_{A'}^\mathbf{L} B'$ である。したがって
$$
(A \otimes_{A'} B') \otimes_{B'}^\mathbf{L} C' =
(A \otimes_{A'}^\mathbf{L} B') \otimes_{B'}^\mathbf{L} C' =
A \otimes_{A'}^\mathbf{L} C' =
A \otimes_{A'} C'
$$
% P05-EN-CH37-0330: frozen source has two consecutive sentence fragments, "The second equality by" and "The last equality because".
第二の等号は More on Algebra, 補題
\ref{more-algebra-lemma-double-base-change} により、最後の等号は
$A' \to C'$ が平坦であることにより従う。これで補題が証明された。
\end{proof}

\begin{lemma}[導来 Chow の補題]
\label{more-morphisms-lemma-derived-chow}
$A$ を環とし、$X$ を $A$ 上有限表示な分離スキームとする。$x \in X$ とする。
このとき、$x$ の開近傍 $U \subset X$、整数 $n \geq 0$、開集合
$V \subset \mathbf{P}^n_A$、閉部分スキーム
$Z \subset X \times_A \mathbf{P}^n_A$、点 $z \in Z$、および
$D(\mathcal{O}_{X \times_A \mathbf{P}^n_A})$ の対象 $E$ で、次を満たすものが
存在する。
\begin{enumerate}
\item $Z \to X \times_A \mathbf{P}^n_A$ は有限表示である。
\item $b : Z \to X$ は $U$ 上同型であり、$b(z) = x$ である。
\item $c : Z \to \mathbf{P}^n_A$ は $V$ 上閉埋め込みである。
\item $b^{-1}(U) = c^{-1}(V)$ であり、特に $c(z) \in V$ である。
\item $E|_{X \times_A V} \cong
(b, c)_*\mathcal{O}_Z|_{X \times_A V}$,
\item $E$ は擬連接で、$Z$ に台をもつ。
\end{enumerate}
\end{lemma}

\begin{proof}
有限型 $\mathbf{Z}$-部分代数 $A' \subset A$ と、$A'$ 上有限表示な分離スキーム
$X'$ で、その $A$ への基底変換が $X$ となるものが存在する。Limits, 補題
\ref{limits-lemma-descend-finite-presentation} および
\ref{limits-lemma-descend-separated-finite-presentation} を参照。$x$ の像を
$x' \in X'$ とする。$x' \in X'/A'$ に対して補題を証明できれば、
$x \in X/A$ に対する補題が従う。実際、$U', n', V', Z', z', E'$ が
$x' \in X'/A'$ に対する解を与えるならば、$U \subset X$ を $U'$ の逆像、
$n = n'$、$V \subset \mathbf{P}^n_A$ を $V'$ の逆像、
$Z \subset X \times \mathbf{P}^n$ を $Z'$ のスキーム論的逆像、$z \in Z$ を
$x$ に写る一意な点とし、$E$ を $E'$ の導来引き戻しとすればよい。
Cohomology, 補題 \ref{cohomology-lemma-pseudo-coherent-pullback} により $E$ は
擬連接である。残るのは (5) の確認だけである。そのため
$W = b^{-1}(U) = c^{-1}(V)$ とおき、
% P05-EN-CH37-0331: frozen source applies (b')^{-1} to U rather than the open U' of X'.
$W' = (b')^{-1}(U) = (c')^{-1}(V')$ とおいて、Cartesian 平方
$$
\xymatrix{
W \ar[d]_{(b, c)} \ar[r] & W' \ar[d]^{(b', c')} \\
X \times_A V \ar[r] & X' \times_{A'} V'
}
$$
を考える。補題 \ref{more-morphisms-lemma-case-of-tor-independence} により、スキーム
$X \times_A V$ と $W'$ は $X' \times_{A'} V'$ 上 Tor 独立である。したがって
Derived Categories of Schemes, 補題 \ref{perfect-lemma-compare-base-change} により、
$(b', c')_*\mathcal{O}_{W'}$ の $X \times_A V$ への導来引き戻しは
$(b, c)_*\mathcal{O}_W$ である。ここでは、$(b', c')$ が閉埋め込みなので
$R(b', c')_*\mathcal{O}_{Z'} = (b', c')_*\mathcal{O}_{Z'}$ であること、および
$(b, c)_*\mathcal{O}_Z$ についても同様であることを用いた。
$E'|_{U' \times_{A'} V'} = (b', c')_*\mathcal{O}_{W'}$ なので
$E|_{U \times_A V} = (b, c)_*\mathcal{O}_W$ を得て、(5) が成り立つ。
これにより次段落の状況へ帰着した。

\medskip\noindent
$A$ は $\mathbf{Z}$ 上有限型であると仮定する。$x$ のアフィン開近傍
$U \subset X$ を選ぶ。このとき $U$ は $A$ 上有限型である。閉埋め込み
$U \to \mathbf{A}^n_A$ を選び、これと開埋め込み
$\mathbf{A}^n_A \to \mathbf{P}^n_A$ の合成で得られる埋め込みを
$j : U \to \mathbf{P}^n_A$ と書く。$Z$ を
$$
(\text{id}_U, j) : U \longrightarrow X \times_A \mathbf{P}^n_A
$$
のスキーム論的閉包とする。射影 $X \times \mathbf{P}^n \to X$ は分離的なので、
Morphisms, 補題 \ref{morphisms-lemma-scheme-theoretic-image-of-partial-section} により
$b : Z \to X$ は $U$ 上同型である。$x$ 上にある一意な点を $z \in Z$ とする。

\medskip\noindent
$Y \subset \mathbf{P}^n_A$ を $j$ のスキーム論的閉包とする。このとき
$Z \subset X \times_A Y$ が $(\text{id}_U, j) : U \to X \times_A Y$ の
スキーム論的閉包であることは明らかである。$X$ は分離的なので、射
$X \times_A Y \to Y$ も分離的である。したがって上で用いたのと同じ補題により、
$Z \to Y$ は開部分スキーム $j(U) \subset Y$ 上同型である。
$V \cap Y = j(U)$ を満たす開集合 $V \subset \mathbf{P}^n_A$ を選ぶ。
これで (3) と (4) が成り立つ。

\medskip\noindent
$A$ は Noether 環なので、$X$ と $X \times_A \mathbf{P}^n_A$ は Noether
スキームである。したがってこの場合は $E = (b, c)_*\mathcal{O}_Z$ と取れる。
Derived Categories of Schemes, 補題
\ref{perfect-lemma-identify-pseudo-coherent-noetherian} を参照。これで証明が終わる。
\end{proof}

\begin{lemma}
\label{more-morphisms-lemma-compute-Fourier-Mukai-for-derived-chow}
$A$, $x \in X$ および $U, n, V, Z, z, E$ を補題 \ref{more-morphisms-lemma-derived-chow} の
とおりとする。任意の $K \in D_\QCoh(\mathcal{O}_X)$ に対して
$$
Rq_*(Lp^*K \otimes^\mathbf{L} E)|_V = R(U \to V)_*K|_U
$$
が成り立つ。ここで $p : X \times_A \mathbf{P}^n_A \to X$ と
$q : X \times_A \mathbf{P}^n_A \to \mathbf{P}^n_A$ は射影であり、射
$U \to V$ は有限表示な閉埋め込み $c \circ (b|_U)^{-1}$ である。
\end{lemma}

\begin{proof}
$b^{-1}(U) = c^{-1}(V)$ であり、$c$ は $V$ 上閉埋め込みなので、
$c \circ (b|_U)^{-1}$ は閉埋め込みである。$U$ と $V$ は $A$ 上有限表示だから
これは有限表示である。Morphisms, 補題
\ref{morphisms-lemma-finite-presentation-permanence} を参照。まず
$$
Rq_*(Lp^*K \otimes^\mathbf{L} E)|_V =
Rq'_*\left((Lp^*K \otimes^\mathbf{L} E)|_{X \times_A V}\right)
$$
が成り立つ。ここで $q' : X \times_A V \to V$ は射影である。これは全導来順像の
形成が局所化と可換だからである。$W = b^{-1}(U) = c^{-1}(V)$ とおき、閉埋め込み
$i = (b, c)|_W$ を $i : W \to X \times_A V$ と書く。このとき
$$
Rq'_*\left((Lp^*K \otimes^\mathbf{L} E)|_{X \times_A V}\right) =
Rq'_*(Lp^*K|_{X \times_A V} \otimes^\mathbf{L} i_*\mathcal{O}_W)
$$
を (5) により得る。$i$ は閉埋め込みなので
$i_*\mathcal{O}_W = Ri_*\mathcal{O}_W$ である。Derived Categories of Schemes,
補題 \ref{perfect-lemma-cohomology-base-change} を用いると、これは
$$
Rq'_* Ri_* Li^* Lp^*K|_{X \times_A V} =
R(q' \circ i)_* Lb^*K|_W =
R(U \to V)_* K|_U
$$
と書き直せる。これが求める等式である。
\end{proof}

\begin{lemma}
\label{more-morphisms-lemma-characterize-pseudo-coherent}
$A$ を環とし、$X$ を $A$ 上有限表示な分離スキームとする。
$K \in D_\QCoh(\mathcal{O}_X)$ とする。$D(\mathcal{O}_X)$ の任意の擬連接な
$E$ に対して $R\Gamma(X, E \otimes^\mathbf{L} K)$ が $D(A)$ で擬連接ならば、
$K$ は $A$ に関して擬連接である。
\end{lemma}

\begin{proof}
$K \in D_\QCoh(\mathcal{O}_X)$ であり、$D(\mathcal{O}_X)$ の任意の擬連接な
$E$ に対して $R\Gamma(X, E \otimes^\mathbf{L} K)$ が $D(A)$ で擬連接であると
仮定する。$x \in X$ とする。$x$ の近傍で $K$ が $A$ に関して擬連接であることを
示せば、補題が証明される。

\medskip\noindent
補題 \ref{more-morphisms-lemma-derived-chow} のように $U, n, V, Z, z, E$ を選ぶ。射影を
$p : X \times \mathbf{P}^n \to X$ および
$q : X \times \mathbf{P}^n \to \mathbf{P}^n_A$ と書く。このとき任意の
$i \in \mathbf{Z}$ に対して
\begin{align*}
& R\Gamma(\mathbf{P}^n_A,
Rq_*(Lp^*K \otimes^\mathbf{L} E)
\otimes^\mathbf{L}
\mathcal{O}_{\mathbf{P}^n_A}(i)) \\
& =
R\Gamma(X \times \mathbf{P}^n,
Lp^*K \otimes^\mathbf{L} E \otimes^\mathbf{L}
Lq^*\mathcal{O}_{\mathbf{P}^n_A}(i)) \\
& =
R\Gamma(X,
K \otimes^\mathbf{L} Rp_*(E \otimes^\mathbf{L}
Lq^*\mathcal{O}_{\mathbf{P}^n_A}(i)))
\end{align*}
が成り立つ。これは Derived Categories of Schemes, 補題
\ref{perfect-lemma-cohomology-base-change} による。Derived Categories of Schemes,
補題 \ref{perfect-lemma-flat-proper-pseudo-coherent-direct-image-general} により、
複体 $Rp_*(E \otimes^\mathbf{L} Lq^*\mathcal{O}_{\mathbf{P}^n_A}(i))$ は $X$ 上
擬連接である。したがって仮定により、表示式の対象は $D(A)$ の擬連接対象である。
Derived Categories of Schemes, 補題
\ref{perfect-lemma-pseudo-coherent-on-projective-space} により、
$Rq_*(Lp^*K \otimes^\mathbf{L} E)$ は $\mathbf{P}^n_A$ 上擬連接であると
結論できる。補題 \ref{more-morphisms-lemma-compute-Fourier-Mukai-for-derived-chow} により
$$
Rq_*(Lp^*K \otimes^\mathbf{L} E)|_{X \times_A V} =
R(U \to V)_*K|_U
$$
% P05-EN-CH37-0332: frozen source restricts Rq_* on P^n_A to X \times_A V rather than to V.
$U \to V$ は $\mathbf{P}^n_A$ の開部分スキームへの閉埋め込みなので、補題
\ref{more-morphisms-lemma-check-relative-pseudo-coherence-on-charts} により、これは $K|_U$ が
$A$ に関して擬連接であることを意味する。
\end{proof}

\begin{lemma}
\label{more-morphisms-lemma-characterize-pseudo-coh-improved}
$A$ を環とし、$X$ を $A$ 上有限表示な分離スキームとする。
$K \in D_\QCoh(\mathcal{O}_X).$ 任意の完全対象
$E \in D(\mathcal{O}_X)$ に対して $R \Gamma (X, E \otimes ^{\mathbf{L}} K)$ が
$D(A)$ で擬連接ならば、$K$ は $A$ に関して擬連接である。
\end{lemma}

\begin{proof}
補題 \ref{more-morphisms-lemma-characterize-pseudo-coherent} に照らせば、任意の擬連接対象
$E \in D(\mathcal{O}_X)$ に対して $R \Gamma (X, E \otimes ^{\mathbf{L}} K)$ が
$D(A)$ で擬連接であることを示せば十分である。Derived Categories of Schemes,
命題 \ref{perfect-proposition-detecting-bounded-above} により
$K \in D^-_\QCoh (\mathcal{O}_X)$ が従う。したがって結果は
Derived Categories of Schemes, 補題 \ref{perfect-lemma-perfect-enough} から従う。
\end{proof}

\begin{lemma}
\label{more-morphisms-lemma-characterize-relatively-perfect}
$A$ を環とし、$X$ を $A$ 上有限表示、平坦かつ分離的なスキームとする。
$K \in D_\QCoh(\mathcal{O}_X).$ 任意の完全対象
$E \in D(\mathcal{O}_X)$ に対して $R \Gamma (X, E \otimes^\mathbf{L} K)$ が
$D(A)$ で完全ならば、$K$ は $\Spec(A)$-完全である。
\end{lemma}

\begin{proof}
補題 \ref{more-morphisms-lemma-characterize-pseudo-coh-improved} により、$K$ は $A$ に関して
擬連接である。補題 \ref{more-morphisms-lemma-check-relative-pseudo-coherence-on-charts} により、
$K$ は $D( \mathcal{O}_X)$ で擬連接である。Derived Categories of Schemes,
命題 \ref{perfect-proposition-detecting-bounded-below} により、$K$ は
$D^-(\mathcal{O}_X)$ に属する。$\mathfrak{p}$ を $A$ の素イデアルとし、
$\mathfrak{p}$ 上のスキーム論的ファイバーの包含を $i : Y \to X$ と書く。
すなわち $Y$ は $\kappa(\mathfrak p)$ 上のスキームである。Derived Categories of
Schemes, 補題 \ref{perfect-lemma-bounded-on-fibres} により、$Li^*(K)$ が下に有界で
あることを示せば十分である。$D_\QCoh (\mathcal{O}_X)$ を生成する完全複体を
$G \in D_{perf} (\mathcal{O}_X)$ とする。Derived Categories of Schemes, 定理
\ref{perfect-theorem-bondal-van-den-Bergh} を参照。このとき
\begin{align*}
R\Hom _{\mathcal{O}_Y}(Li^*(G), Li^*(K))
& =
R\Gamma(Y, Li^*(G ^\vee \otimes ^\mathbf{L} K)) \\
& =
R\Gamma(X, G^\vee \otimes ^{\mathbf{L}} K)
\otimes^\mathbf{L}_A \kappa(\mathfrak{p})
\end{align*}
が成り立つ。第一の等号では、$Li^*$ が完全対象と双対を保つこと、および
Cohomology, 補題 \ref{cohomology-lemma-dual-perfect-complex} を用いる。細部は
省略する。第二の等号は、$X$ が $A$ 上平坦なので Derived Categories of Schemes,
補題 \ref{perfect-lemma-compare-base-change} から従う。仮定により、これは
$D(\kappa(\mathfrak{p}))$ の完全対象である。Derived Categories of Schemes, 備考
\ref{perfect-remark-pullback-generator} により、対象
$Li^*(G) \in D_{perf}(\mathcal{O}_Y)$ は $D_\QCoh(\mathcal{O}_Y)$ を生成する。
したがって Derived Categories of Schemes, 命題
\ref{perfect-proposition-detecting-bounded-below} により $Li^*(K)$ は下に有界であり、
結論が従う。
\end{proof}




\section{複体の有限性性質の降下}
\label{more-morphisms-section-descent-finiteness}

\noindent
本節は Derived Categories of Schemes, 第
\ref{perfect-section-descent-finiteness} 節の続きである。

\begin{lemma}
\label{more-morphisms-lemma-relative-pseudo-coherent-descends-fppf}
$X \to S$ を局所有限型とし、$\{f_i : X_i \to X\}$ をスキームの fppf 被覆と
する。$E \in D_\QCoh(\mathcal{O}_X)$, $m \in \mathbf{Z}$ とする。このとき
$E$ が $S$ に関して $m$-擬連接であることと、各 $Lf_i^*E$ が $S$ に関して
$m$-擬連接であることは同値である。
\end{lemma}

\begin{proof}
$E$ が $S$ に関して $m$-擬連接であると仮定する。補題
\ref{more-morphisms-lemma-flat-finite-presentation-pseudo-coherent} により射 $f_i$ は擬連接である。
したがって補題 \ref{more-morphisms-lemma-pull-relative-pseudo-coherent} により $Lf_i^*E$ は
$S$ に関して $m$-擬連接である。

\medskip\noindent
逆に、各 $i$ に対して $Lf_i^*E$ が $S$ に関して $m$-擬連接であると仮定する。
補題 \ref{more-morphisms-lemma-dominate-fppf} のように、$S = \bigcup U_j$,
$W_j \to U_j$, $W_j = \bigcup W_{j, k}$, $T_{j, k} \to W_{j, k}$ および
$S$ 上の射 $\alpha_{j, k} : T_{j, k} \to X_{i(j, k)}$ を選ぶ。
% P05-EN-CH37-0333: frozen source changes k to K and then uses the undefined T_{i,k}.
射 $T_{j, K} \to S$ は平坦かつ有限表示なので、補題
\ref{more-morphisms-lemma-permanence-pseudo-coherent} により $\alpha_{j, k}$ は擬連接である。
したがって
$$
L\alpha_{j, k}^*Lf_{i(j, k)}^*E = L(T_{i, k} \to S)^*E
$$
は補題 \ref{more-morphisms-lemma-pull-relative-pseudo-coherent} により $S$ に関して
$m$-擬連接である。ここからこの性質を被覆 $\{T_{j, k} \to W_{j, k}\}$,
$W_j = \bigcup W_{j, k}$, $\{W_j \to U_j\}$, $S = \bigcup U_j$ に沿って
降下させたい。Zariski 被覆については結果が成り立つ（$S$ に関する
$m$-擬連接性の定義による）ので、$L\pi^*E$ が $S$ に関して擬連接となる単一の
全射な有限局所自由射 $\pi : Y \to X$ がある場合を仮定してよい。
% P05-EN-CH37-0334: frozen source drops the m-qualifier in both descent assertions.
この場合、補題 \ref{more-morphisms-lemma-finite-morphism-relative-pseudo-coherence} により
$R\pi_*L\pi^*E$ は $S$ に関して擬連接である（ここで初めて $E$ のコホモロジー層が
準連接であることを用いる）。また
$R\pi_*L\pi^*E = E \otimes^\mathbf{L}_{\mathcal{O}_X} \pi_*\mathcal{O}_Y$
が成り立つ。例えば Derived Categories of Schemes, 補題
\ref{perfect-lemma-cohomology-base-change} による。さらに $X$ 上局所的に、写像
$\mathcal{O}_X \to \pi_*\mathcal{O}_Y$ は直和因子の包含である。したがって
補題 \ref{more-morphisms-lemma-summands-relative-pseudo-coherent} により結論が従う。
\end{proof}

\begin{lemma}
\label{more-morphisms-lemma-relative-pseudo-coherent-post-compose}
$X \to T \to S$ をスキームの射とする。$T \to S$ は平坦かつ局所有限表示で、
$X \to T$ は局所有限型であると仮定する。$E \in D(\mathcal{O}_X)$,
$m \in \mathbf{Z}$ とする。このとき $E$ が $S$ に関して $m$-擬連接であることと、
$E$ が $T$ に関して $m$-擬連接であることは同値である。
\end{lemma}

\begin{proof}
$X$ 上局所的に閉埋め込み $i : X \to \mathbf{A}^n_T$ を選べる。このとき
$\mathbf{A}^n_T \to S$ は平坦かつ局所有限表示である。したがって補題
\ref{more-morphisms-lemma-composition-relative-pseudo-coherent} を適用すれば同値性が従う。
\end{proof}

\begin{lemma}
\label{more-morphisms-lemma-relative-pseudo-coherent-descends-fppf-base}
$f : X \to S$ を局所有限型とし、$\{S_i \to S\}$ をスキームの fppf 被覆とする。
$f$ の基底変換を $f_i : X_i \to S_i$、射影を $g_i : X_i \to X$ と書く。
$E \in D_\QCoh(\mathcal{O}_X)$, $m \in \mathbf{Z}$ とする。このとき $E$ が
$S$ に関して $m$-擬連接であることと、各 $Lg_i^*E$ が $S_i$ に関して
$m$-擬連接であることは同値である。
\end{lemma}

\begin{proof}
これは補題 \ref{more-morphisms-lemma-relative-pseudo-coherent-descends-fppf} と
\ref{more-morphisms-lemma-relative-pseudo-coherent-post-compose} から形式的に従う。実際、$E$ が
$S$ に関して $m$-擬連接ならば、$Lg_i^*E$ は第一の補題により $S$ に関して
$m$-擬連接であり、したがって第二の補題により $Lg_i^*E$ は $S_i$ に関して
$m$-擬連接である。逆に、$Lg_i^*E$ が $S_i$ に関して $m$-擬連接ならば、
第二の補題により $Lg_i^*E$ は $S$ に関して $m$-擬連接であり、したがって
第一の補題により $E$ は $S$ に関して
$m$-擬連接である。
\end{proof}






\section{相対完全対象}
\label{more-morphisms-section-relatively-perfect}

\noindent
本節は Derived Categories of Schemes, 第
\ref{perfect-section-relatively-perfect} 節の議論の続きである。

\begin{lemma}
\label{more-morphisms-lemma-thickening-pseudo-coherent}
$i : X \to X'$ をスキームの有限次肥厚化とする。$K' \in D(\mathcal{O}_{X'})$ を、
$K = Li^*K'$ が擬連接となる対象とする。このとき $K'$ は擬連接である。
\end{lemma}

\begin{proof}
まず $K'$ が準連接コホモロジー層をもつことを証明する。そのため一次肥厚化の
場合へ帰着してよい。第 \ref{more-morphisms-section-thickenings} 節を参照。$X$ を切り出す
準連接イデアル層を $\mathcal{I} \subset \mathcal{O}_{X'}$ とする。短完全列
$$
0 \to \mathcal{I} \to \mathcal{O}_{X'} \to i_*\mathcal{O}_X \to 0
$$
と $K'$ のテンソル積を取ると、完全三角形
$$
K' \otimes_{\mathcal{O}_{X'}}^\mathbf{L} \mathcal{I}
\to K' \to
K' \otimes_{\mathcal{O}_{X'}}^\mathbf{L} i_*\mathcal{O}_X
\to
(K' \otimes_{\mathcal{O}_{X'}}^\mathbf{L} \mathcal{I})[1]
$$
を得る。$i_* = Ri_*$ であり、一次肥厚化なので $\mathcal{I}$ を準連接
$\mathcal{O}_X$-加群とみなせるから、これは
$$
i_*(K \otimes_{\mathcal{O}_X}^\mathbf{L} \mathcal{I})
\to K' \to
i_*K \to
i_*(K \otimes_{\mathcal{O}_X}^\mathbf{L} \mathcal{I})[1]
$$
と書き直せる。各項の同定には Cohomology, 補題
\ref{cohomology-lemma-projection-formula-closed-immersion} を用いる。
$K$ は $D_\QCoh(\mathcal{O}_X)$ に属するので、$K'$ は
$D_\QCoh(\mathcal{O}_{X'})$ に属すると結論できる。ここでは
Derived Categories of Schemes, 補題 \ref{perfect-lemma-pseudo-coherent},
\ref{perfect-lemma-quasi-coherence-tensor-product},
\ref{perfect-lemma-quasi-coherence-direct-image} を用いる。

\medskip\noindent
$K'$ は $D_\QCoh(\mathcal{O}_{X'})$ に属すると仮定する。問題は $X'$ 上局所的なので、
$X'$ はアフィンと仮定してよい。$X' = \Spec(A')$, $X = \Spec(A)$,
$A = A'/I$ とし、$I$ は冪零とする。このとき $K'$ はある対象
$M' \in D(A')$ から来る。Derived Categories of Schemes, 補題
\ref{perfect-lemma-affine-compare-bounded} を参照。したがって
Derived Categories of Schemes, 補題 \ref{perfect-lemma-pseudo-coherent-affine} と
$K$ に関する仮定により、$M = M' \otimes_{A'}^\mathbf{L} A$ は $D(A)$ の
擬連接対象である。よって More on Algebra, 補題
\ref{more-algebra-lemma-check-pseudo-coherent-modulo-nilpotent} により $M'$ は
擬連接である。これで証明が終わる。
\end{proof}

\begin{lemma}
\label{more-morphisms-lemma-thickening-relatively-perfect}
スキームの Cartesian 図式
$$
\xymatrix{
X \ar[r]_i \ar[d]_f & X' \ar[d]^{f'} \\
Y \ar[r]^j & Y'
}
$$
を考える。$X' \to Y'$ は平坦かつ局所有限表示で、$Y \to Y'$ は有限次肥厚化で
あると仮定する。$E' \in D(\mathcal{O}_{X'})$ とする。$E = Li^*(E')$ が
$Y$-完全ならば、$E'$ は $Y'$-完全である。
\end{lemma}

\begin{proof}
$E$ が $Y$-完全であるとは、$E$ が擬連接で、$f^{-1}\mathcal{O}_Y$-加群の
複体として局所有限 Tor 次元をもつことを意味する（Derived Categories of Schemes,
定義 \ref{perfect-definition-relatively-perfect}）。補題
\ref{more-morphisms-lemma-thickening-pseudo-coherent} により $E'$ は擬連接である。特に
$E'$ は $D_\QCoh(\mathcal{O}_{X'})$ に属する。Derived Categories of Schemes,
補題 \ref{perfect-lemma-pseudo-coherent} を参照。$E'$ が局所有限 Tor 次元を
もつことを証明するには $X'$ 上局所的に考えてよい。したがって
% P05-EN-CH37-0335: frozen source names S',S although the cartesian diagram uses Y',Y.
$X'$, $S'$, $X$, $S$ はアフィンで、それぞれ環 $A'$, $R'$, $A$, $R$ で
与えられると仮定してよい。Derived Categories of Schemes, 補題
\ref{perfect-lemma-affine-locally-rel-perfect} により可換環論版へ帰着する。
% P05-EN-CH37-0336: frozen source leaves the final commutative-algebra citation as a sentence fragment.
可換環論版は More on Algebra, 補題
\ref{more-algebra-lemma-thickening-relatively-perfect} である。
\end{proof}

\begin{lemma}
\label{more-morphisms-lemma-henselian-relatively-perfect}
$(R, I)$ を環とその Jacobson 根基に含まれるイデアル $I$ の組とする。
$S = \Spec(R)$, $S_0 = \Spec(R/I)$ とおく。$f : X \to S$ を固有、平坦かつ
有限表示とし、$X_0 = S_0 \times_S X$ と書く。$E \in D(\mathcal{O}_X)$ を
擬連接とする。$E$ の $X_0$ への導来制限 $E_0$ が $S_0$-完全ならば、
$E$ は $S$-完全である。
\end{lemma}

\begin{proof}
有限アフィン開被覆 $X = U_1 \cup \ldots \cup U_n$ を選ぶ。各 $i$ に対し閉埋め込み
$U_i \to \mathbf{A}^{d_i}_S$ を選べる。$U_{i, 0} = S_0 \times_S U_i$ とおく。
各 $i$ に対して、ある $a_i, b_i \in \mathbf{Z}$ が存在して複体
$E_0|_{U_{i, 0}}$ の Tor 振幅は $[a_i, b_i]$ に含まれる。点 $x \in X$ を取る。
ある $i$ に対し、$E_x$ の $R$ 上の Tor 振幅が $[a_i - d_i, b_i]$ に含まれることを
示す。Cohomology, 補題 \ref{cohomology-lemma-tor-amplitude-stalk} により Tor 振幅は
茎から読み取れるので、これで証明が完了する。

\medskip\noindent
$f$ は固有なので $f(\overline{\{x\}})$ は $S$ の閉部分集合である。$I$ は
Jacobson 根基に含まれるので、
% P05-EN-CH37-0337: frozen source uses the participle "meeting" where the finite verb "meets" is required.
$f(\overline{\{x\}})$ は閉部分集合 $S_0 \subset S$ と交わる。したがって
$x_0 \in X_0$ を満たす特殊化 $x \leadsto x_0$ がある。$x_0 \in U_i$ となる
$i$ を選ぶと、$x_0 \in U_{i, 0}$ である。以後この $i$ を固定する。
$U_i = \Spec(A)$ と書く。このとき $A$ は平坦かつ有限表示な $R$-代数で、
$d_i$ 変数多項式 $R$-代数の商である。Derived Categories of Schemes, 補題
\ref{perfect-lemma-affine-compare-bounded} と
\ref{perfect-lemma-pseudo-coherent-affine} により、制限 $E|_{U_i}$ は $D(A)$ の
擬連接対象 $K$ に対応する。$E_0$ は $K \otimes_A^\mathbf{L} A/IA$ に対応する。
$x \leadsto x_0$ に対応する素イデアルを
$\mathfrak q \subset \mathfrak q_0 \subset A$ とする。このとき
$E_x = K_{\mathfrak q}$ であり、$K_{\mathfrak q}$ は $K_{\mathfrak q_0}$ の
局所化である。したがって $K_{\mathfrak q_0}$ が $R$-加群の複体として
$[a_i - d_i, b_i]$ に含まれる Tor 振幅をもつことを示せば十分である。
$f(x_0)$ に対応する素イデアルを $I \subset \mathfrak p_0 \subset R$ とする。
このとき
\begin{align*}
K \otimes_R^\mathbf{L} \kappa(\mathfrak p_0)
& =
(K \otimes_R^\mathbf{L} R/I) \otimes_{R/I}^\mathbf{L}
\kappa(\mathfrak p_0) \\
& =
(K \otimes_A^\mathbf{L} A/IA) \otimes_{R/I}^\mathbf{L} \kappa(\mathfrak p_0)
\end{align*}
% P05-EN-CH37-0338: frozen source leaves "the second equality because" as a sentence fragment.
が成り立つ。第二の等号は $R \to A$ が平坦であることによる。$a_i, b_i$ の
選び方により、この複体のコホモロジーは区間 $[a_i, b_i]$ の次数でのみ非零である。
したがって最後に More on Algebra, 補題
\ref{more-algebra-lemma-lift-from-fibre-relatively-perfect} を $R \to A$,
$\mathfrak q_0$, $\mathfrak p_0$, $K$ に適用すれば結論が従う。
\end{proof}










\section{有理曲線の収縮}
\label{more-morphisms-section-contracting}

\noindent
本節では、ファイバーが次元 $\leq 1$ をもち、
$R^1f_*\mathcal{O}_X = 0$ を満たす固有射 $f : X \to Y$ を研究する。
本節の題名の意味については、Algebraic Curves の Sections
\ref{curves-section-contracting-rational-tails},
\ref{curves-section-contracting-rational-bridges}, and
\ref{curves-section-contracting-to-stable} を参照されたい。

\begin{lemma}
\label{more-morphisms-lemma-check-h1-fibre-zero}
スキームの固有射 $f : X \to Y$ と、
$\dim(X_y) \leq 1$ を満たす点 $y \in Y$ をとる。次のいずれかが成り立つとする：
\begin{enumerate}
\item $R^1f_*\mathcal{O}_X = 0$ である、または、より一般に、
\item $y$ が $g$ の像に含まれ、かつ
$R^1f'_*\mathcal{O}_{X'} = 0$ となる射 $g : Y' \to Y$ が存在する。
ここで $f' : X' \to Y'$ は $g$ による $f$ の基底変換である。
\end{enumerate}
このとき $H^1(X_y, \mathcal{O}_{X_y}) = 0$ である。
\end{lemma}

\begin{proof}
この補題を証明するには、$Y$ を $y$ の開近傍で置き換えてよい。
したがって、$Y$ はアフィンであり、$f$ のすべてのファイバーの次元が
$\leq 1$ であると仮定できる。Morphisms, Lemma
\ref{morphisms-lemma-openness-bounded-dimension-fibres} を参照せよ。
この場合、$R^1f_*\mathcal{O}_X$ は有限型準連接
$\mathcal{O}_Y$-加群であり、その構成は任意の基底変換と可換である。
Limits, Lemmas \ref{limits-lemma-proper-top-cohomology-finite-type} and
\ref{limits-lemma-higher-direct-images-zero-above-dimension-fibre} を参照せよ。
したがって主張は直ちに従う。
\end{proof}

\begin{lemma}
\label{more-morphisms-lemma-h1-fibre-zero}
スキームの固有射 $f : X \to Y$ と、
$\dim(X_y) \leq 1$ および $H^1(X_y, \mathcal{O}_{X_y}) = 0$ を満たす点
$y \in Y$ をとる。このとき、$y$ の開近傍 $V \subset Y$ であって
$R^1f_*\mathcal{O}_X|_V = 0$
となるものが存在し、任意の $Y' \to V$ による基底変換後にも同じことが成り立つ。
\end{lemma}

\begin{proof}
この補題を証明するには、$Y$ を $y$ の開近傍で置き換えてよい。
したがって、$Y$ はアフィンであり、$f$ のすべてのファイバーの次元が
$\leq 1$ であると仮定できる。Morphisms, Lemma
\ref{morphisms-lemma-openness-bounded-dimension-fibres} を参照せよ。
この場合、$R^1f_*\mathcal{O}_X$ は有限型準連接
$\mathcal{O}_Y$-加群であり、その構成は任意の基底変換と可換である。
Limits, Lemmas \ref{limits-lemma-proper-top-cohomology-finite-type} and
\ref{limits-lemma-higher-direct-images-zero-above-dimension-fibre} を参照せよ。
$Y = \Spec(A)$ とし、$y$ が素イデアル $\mathfrak p \subset A$ に対応し、
$R^1f_*\mathcal{O}_X$ が有限 $A$-加群 $M$ に対応するとする。
このとき、基底変換に関する上の主張により、
$H^1(X_y, \mathcal{O}_{X_y}) = 0$ は
$\mathfrak pM_\mathfrak p = M_\mathfrak p$ を意味する。
中山の補題から $M_\mathfrak p = 0$ を得る。$M$ は有限なので、
$M_f = 0$ を満たす $f \in A$, $f \not \in \mathfrak p$ が存在する。
したがって、$V$ を主開集合 $D(f)$ とすれば所望の結果を得る。
\end{proof}

\begin{lemma}
\label{more-morphisms-lemma-globally-generated-vanishing}
スキームの固有射 $f : X \to Y$ が、すべての $y \in Y$ に対して
$\dim(X_y) \leq 1$ および $H^1(X_y, \mathcal{O}_{X_y}) = 0$
を満たすとする。$X$ 上の準連接加群 $\mathcal{F}$ をとる。このとき
\begin{enumerate}
\item $p > 1$ に対して $R^pf_*\mathcal{F} = 0$ であり、
\item $Y$ 上の準連接加群 $\mathcal{G}$ からの全射
$f^*\mathcal{G} \to \mathcal{F}$ が存在すれば、
$R^1f_*\mathcal{F} = 0$ である。
\end{enumerate}
$Y$ がアフィンならば、さらに
\begin{enumerate}
\item[(3)] $p \not \in \{0, 1\}$ に対して $H^p(X, \mathcal{F}) = 0$ であり、
\item[(4)] $\mathcal{F}$ が大域生成ならば $H^1(X, \mathcal{F}) = 0$ である。
\end{enumerate}
\end{lemma}

\begin{proof}
(1) の消滅は Limits, Lemma
\ref{limits-lemma-higher-direct-images-zero-above-dimension-fibre} である。
(2) を証明するには $Y$ 上で局所的に議論してよいので、$Y$ はアフィンと仮定する。
すると $R^1f_*\mathcal{F}$ は、加群 $H^1(X, \mathcal{F})$ に付随する
$Y$ 上の準連接加群である。ここでは、$Y$ がアフィンであること、高次順像の
準連接性（Cohomology of Schemes, Lemma
\ref{coherent-lemma-quasi-coherence-higher-direct-images}), and
Cohomology of Schemes, Lemma
\ref{coherent-lemma-quasi-coherence-higher-direct-images-application}）
を用いている。
$Y$ はアフィンなので準連接加群 $\mathcal{G}$ は大域生成であり、したがって
$f^*\mathcal{G}$ および $\mathcal{F}$ も大域生成である。
これにより (4) から (2) が従う。
(3) は (1) と、先ほど述べた順像の準連接性から従う。したがって、
% P05-EN-CH37-0339
残るのは (4) の証明だけである。
$\mathcal{F}$ が大域生成ならば、全射
$\bigoplus_{i \in I} \mathcal{O}_X \to \mathcal{F}$ が存在する。
(1) とコホモロジーの長完全列により、これは $H^1$ 上の全射を誘導する。
Lemma \ref{more-morphisms-lemma-h1-fibre-zero} により
$R^1f_*\mathcal{O}_X = 0$、したがって $H^1(X, \mathcal{O}_X) = 0$ であり、
さらに $H^1(X, -)$ は直和と可換なので
(Cohomology, Lemma \ref{cohomology-lemma-quasi-separated-cohomology-colimit})
主張が従う。
\end{proof}

\begin{lemma}
\label{more-morphisms-lemma-h1-fibre-zero-check-h0-kappa}
スキームの固有射 $f : X \to Y$ が次を満たすとする：
\begin{enumerate}
\item すべての $y \in Y$ に対して $\dim(X_y) \leq 1$ かつ
$H^1(X_y, \mathcal{O}_{X_y}) = 0$ であり、
\item $\mathcal{O}_Y \to f_*\mathcal{O}_X$ は全射である。
\end{enumerate}
このとき、$f$ の任意の基底変換 $f' : X' \to Y'$ に対して
$\mathcal{O}_{Y'} \to f'_*\mathcal{O}_{X'}$ は全射である。
\end{lemma}

\begin{proof}
$Y$ と $Y'$ はアフィンと仮定してよい。このとき、$Y'' \to Y$ が
アフィンスキームの平坦射となるような閉埋め込み $Y' \to Y''$ を選べる。
平坦基底変換
(Cohomology of Schemes, Lemma \ref{coherent-lemma-flat-base-change-cohomology})
により、$X'' \to Y''$ に対して結果が成り立つ。
したがって、$Y'$ は $Y$ の閉部分スキームと仮定してよい。
$Y'$ を切り出すイデアルを $\mathcal{I} \subset \mathcal{O}_Y$ とする。
このとき短完全列
$$
0 \to \mathcal{I}\mathcal{O}_X \to
\mathcal{O}_X \to \mathcal{O}_{X'} \to 0
$$
があり、ここでは $\mathcal{O}_{X'}$ を $X$ 上の準連接加群とみなしている。
Lemma \ref{more-morphisms-lemma-globally-generated-vanishing} により
$H^1(X, \mathcal{I}\mathcal{O}_X) = 0$ である。したがって
$$
H^0(Y, \mathcal{O}_Y) \to
H^0(Y, f_*\mathcal{O}_X) = H^0(X, \mathcal{O}_X) \to
H^0(X, \mathcal{O}_{X'})
$$
は所望のとおり全射である。第1の矢印が全射なのは、$Y$ がアフィンであり、
$\mathcal{O}_Y \to f_*\mathcal{O}_X$ が全射と仮定したからである。
第2の矢印が全射なのは、上の短完全列に付随するコホモロジー長完全列と、
先ほど証明した消滅による。
\end{proof}

\begin{lemma}
\label{more-morphisms-lemma-h1-fibre-zero-h0-kappa}
スキームの射の可換図式
$$
\xymatrix{
X \ar[rr]_f \ar[rd] & & Y \ar[ld] \\
& S
}
$$
を考える。点 $s \in S$ をとり、次を仮定する：
\begin{enumerate}
\item $X \to S$ は有限表示で局所的であり、$X_s$ の各点で平坦である、
\item $f$ は固有である、
\item $f_s : X_s \to Y_s$ のファイバーの次元は $\leq 1$ であり、
$R^1f_{s, *}\mathcal{O}_{X_s} = 0$ である、
\item $\mathcal{O}_{Y_s} \to f_{s, *}\mathcal{O}_{X_s}$ は全射である。
\end{enumerate}
このとき、開集合 $Y_s \subset V \subset Y$ であって、
(a) $f^{-1}(V)$ は $S$ 上平坦、
(b) $y \in V$ に対して $\dim(X_y) \leq 1$、
(c) $R^1f_*\mathcal{O}_X|_V = 0$、
(d) $\mathcal{O}_V \to f_*\mathcal{O}_X|_V$ は全射、
となるものが存在し、(b), (c), (d) は任意の $Y' \to V$ による基底変換後にも成り立つ。
\end{lemma}

\begin{proof}
点 $y \in Y$ が $s$ の上にあるとする。
所望の性質をもつ $y$ の開近傍を見つければ十分である。
まず $Y$ を Lemma \ref{more-morphisms-lemma-h1-fibre-zero} で得られる開集合 $V$ で置き換え、
$R^1f_*\mathcal{O}_X$ が普遍的に零となるようにする
（この補題の仮定は Lemma \ref{more-morphisms-lemma-check-h1-fibre-zero} により成り立つ）。
さらに $Y$ を縮小し、$f$ のすべてのファイバーの次元が $\leq 1$ となるようにする
（Morphisms, Lemma
\ref{morphisms-lemma-openness-bounded-dimension-fibres}
と $f$ の固有性を用いる）。したがって、$V = Y$ として (b), (c) が成り立ち、
任意の基底変換 $Y' \to Y$ の後にも成り立つと仮定できる。
よって Lemma \ref{more-morphisms-lemma-h1-fibre-zero-check-h0-kappa} により、
$Y$ 上で (d) を示せば十分である。
さらに $Y$ を縮小してよく、実際 $Y$ と $S$ はアフィンと仮定する。

\medskip\noindent
Theorem \ref{more-morphisms-theorem-openness-flatness} により、$X \to S$ が平坦となる
開部分集合 $U \subset X$ があり、仮定からこれは $X_s$ を含む。
このとき $f(X \setminus U)$ は $y$ を含まない閉部分集合である。
したがって $Y$ を縮小すれば、$X$ は $S$ 上平坦と仮定できる。

\medskip\noindent
$S = \Spec(R)$ とする。集合 $E$ 上の多項式環 $R[x_e; e \in E]$ の
スペクトルを $Y'$ とし、閉埋め込み $Y \to Y'$ を選ぶ。
$f$ と $Y \to Y'$ の合成を $f' : X \to Y'$ と書く。
すると $f'$ と $s$ に対して、仮定 (1) -- (4) および (b), (c) が成り立つ。
% P05-EN-CH37-0340
$y$ の開近傍で $\mathcal{O}_{Y'} \to f'_*\mathcal{O}_X$ が全射であることを示せば、
$\mathcal{O}_Y \to f_*\mathcal{O}_X$ についても同じことが成り立つ。
したがって、$Y$ は $R[x_e; e \in E]$ のスペクトルと仮定してよい。

\medskip\noindent
ここで $X$ と $Y$ は $S$ 上平坦である。すると $Y_s$ と $X$ は
$Y$ 上 Tor 独立である。読者自身で証明を見つけることを勧めるが、これは
頂点が $X, Y, S, S$ である正方形と、その $s \to S$ による基底変換に
Lemma \ref{more-morphisms-lemma-case-of-tor-independence} を適用しても従う。したがって
$$
Rf_{s, *}\mathcal{O}_{X_s} = L(Y_s \to Y)^*Rf_*\mathcal{O}_X
$$
である。Derived Categories of Schemes, Lemma
\ref{perfect-lemma-compare-base-change} を参照せよ。
すでに示した消滅により、これは
$f_{s, *}\mathcal{O}_{X_s} = (Y_s \to Y)^*f_*\mathcal{O}_X$ を含意する。
したがって $\mathcal{O}_Y \to f_*\mathcal{O}_X$ は準連接
$\mathcal{O}_Y$-加群の射であり、その $Y_s$ への引き戻しは全射である。
$f_*\mathcal{O}_X$ は有限型 $\mathcal{O}_Y$-加群であると主張する。
これが正しければ、
$\mathcal{O}_Y \to f_*\mathcal{O}_X$
の余核 $\mathcal{F}$ は有限型準連接 $\mathcal{O}_Y$-加群であり、
$\mathcal{F}_y \otimes \kappa(y) = 0$ を満たす。
中山の補題 (Algebra, Lemma \ref{algebra-lemma-NAK}) により
$\mathcal{F}_y = 0$ である。よって $\mathcal{F}$ は $y$ のある開近傍で零であり
(Modules, Lemma \ref{modules-lemma-finite-type-stalk-zero})
証明が完了する。

\medskip\noindent
主張を証明する。
有限部分集合 $E' \subset E$ に対して $Y' = \Spec(R[x_e; e \in E'])$ とおく。
十分大きい $E'$ に対して射 $f' : X \to Y \to Y'$ は固有である。
Limits, Lemma \ref{limits-lemma-finite-type-eventually-proper} を参照せよ。
以下ではそのような $E'$ と $Y'$ を固定する。
有限型 $\mathbf{Z}$-部分代数の余極限として $R = \colim R_i$ と書く。
$S_i = \Spec(R_i)$ および $Y'_i = \Spec(R_i[x_e; e \in E'])$ とおく。
十分大きい $i$ に対して、図式
$$
\xymatrix{
X \ar[d] \ar[r]_{f'} & Y' \ar[d] \ar[r] & S \ar[d] \\
X_i \ar[r]^{f'_i} & Y'_i \ar[r] & S_i
}
$$
で、各正方形が Cartesian、$X_i$ が $S_i$ 上平坦、かつ
$X_i \to Y'_i$ が固有となるものを見つけられる。Limits, Lemmas
\ref{limits-lemma-descend-finite-presentation},
\ref{limits-lemma-descend-flat-finite-presentation}, and
\ref{limits-lemma-eventually-proper} を参照せよ。
上と同じ議論により $Y'$ と $X_i$ は $Y'_i$ 上 Tor 独立であり、したがって
$$
R\Gamma(X, \mathcal{O}_X) =
R\Gamma(X_i, \mathcal{O}_{X_i})
\otimes^\mathbf{L}_{R_i[x_e; e \in E']}
R[x_e; e \in E']
$$
である。ここでも上と同じ文献を用いた。Cohomology of Schemes, Lemma
\ref{coherent-lemma-proper-over-affine-cohomology-finite}
により、複体 $R\Gamma(X_i, \mathcal{O}_{X_i})$ は Noether 環
$R_i[x_e; e \in E']$ の導来圏で擬連接である
（More on Algebra, Lemma \ref{more-algebra-lemma-Noetherian-pseudo-coherent} を参照せよ）。
したがって $R\Gamma(X, \mathcal{O}_X)$ は $R[x_e; e \in E']$ の導来圏で
擬連接である。More on Algebra, Lemma
\ref{more-algebra-lemma-pull-pseudo-coherent} を参照せよ。
消滅しないコホモロジー加群は $H^0(X, \mathcal{O}_X)$ だけなので、これは有限
$R[x_e; e \in E']$-加群である。More on Algebra, Lemma
\ref{more-algebra-lemma-n-pseudo-module} を参照せよ。これで証明が完了する。
\end{proof}

\begin{lemma}
\label{more-morphisms-lemma-h1-fibre-zero-h0-flat}
スキームの射の可換図式
$$
\xymatrix{
X \ar[rr]_f \ar[rd] & & Y \ar[ld] \\
& S
}
$$
を考える。$X \to S$ は平坦、$f$ は固有、$y \in Y$ に対して
$\dim(X_y) \leq 1$、かつ $R^1f_*\mathcal{O}_X = 0$ と仮定する。
このとき $f_*\mathcal{O}_X$ は $S$ 上平坦であり、
$f_*\mathcal{O}_X$ の構成は任意の基底変換 $S' \to S$ と可換である。
\end{lemma}

\begin{proof}
$Y$ と $S$ はアフィンと仮定してよい。$S = \Spec(A)$ とする。
準連接 $\mathcal{O}_Y$-加群 $f_*\mathcal{O}_X$ が $S$ に関して平坦であることを
示すには、$H^0(X, \mathcal{O}_X)$ が $A$ 上平坦であることを示せば十分である
（いくつかの詳細は省略する）。
Lemma \ref{more-morphisms-lemma-globally-generated-vanishing} により、任意の $A$-加群 $M$ に対して
$H^1(X, \mathcal{O}_X \otimes_A M) = 0$ である。
さらに $\mathcal{O}_X$ は $A$ 上平坦なので、関手
$M \mapsto H^0(X, \mathcal{O}_X \otimes_A M)$ は完全である。
また、Cohomology, Lemma
\ref{cohomology-lemma-quasi-separated-cohomology-colimit} により、
この関手は直和と可換である。すると、
$H^0(X, \mathcal{O}_X \otimes_A M) = M \otimes_A H^0(X, \mathcal{O}_X)$
が $M$ に関して関手的に成り立つことは演習であり、これから所望の平坦性を得る。
最後に、$S' \to S$ が環準同型 $A \to A'$ で与えられるアフィンスキームの射ならば、
先ほど論じたアフィンの場合から
$$
H^0(X \times_S S', \mathcal{O}_{X \times_S S'}) =
H^0(X, \mathcal{O}_X \otimes_A A') = H^0(X, \mathcal{O}_X) \otimes_A A'
$$
を得る。これは $f_*\mathcal{O}_X$ の構成が任意の基底変換
$S' \to S$ と可換であることを示す。いくつかの詳細は省略する。
\end{proof}

\begin{lemma}
\label{more-morphisms-lemma-h1-fibre-zero-isom}
スキームの射の可換図式
$$
\xymatrix{
X \ar[rr]_f \ar[rd] & & Y \ar[ld] \\
& S
}
$$
を考える。点 $s \in S$ をとり、次を仮定する：
\begin{enumerate}
\item $X \to S$ は有限表示で局所的であり、$X_s$ の各点で平坦である、
\item $Y \to S$ は有限表示で局所的である、
\item $f$ は固有である、
\item $f_s : X_s \to Y_s$ のファイバーの次元は $\leq 1$ であり、
$R^1f_{s, *}\mathcal{O}_{X_s} = 0$ である、
\item $\mathcal{O}_{Y_s} \to f_{s, *}\mathcal{O}_{X_s}$ は同型である。
\end{enumerate}
このとき、開集合 $Y_s \subset V \subset Y$ であって、
(a) $V$ は $S$ 上平坦、
(b) $f^{-1}(V)$ は $S$ 上平坦、
(c) $y \in V$ に対して $\dim(X_y) \leq 1$、
(d) $R^1f_*\mathcal{O}_X|_V = 0$、
(e) $\mathcal{O}_V \to f_*\mathcal{O}_X|_V$ は同型、
となるものが存在し、(a) -- (e) は $f^{-1}(V) \to V$ を任意の
$S' \to S$ で基底変換した後にも成り立つ。
\end{lemma}

\begin{proof}
点 $y \in Y_s$ をとる。$Y$ はいつでも $y$ の開近傍で置き換えてよい。
したがって $Y$ と $S$ はアフィンと仮定できる。さらに、$X$ は $S$ 上平坦、
$y \in Y$ に対して $\dim(X_y) \leq 1$、
$R^1f_*\mathcal{O}_X = 0$ は普遍的に成り立ち、かつ
$\mathcal{O}_Y \to f_*\mathcal{O}_X$ は全射であると仮定できる。
Lemma \ref{more-morphisms-lemma-h1-fibre-zero-h0-kappa} を参照せよ
（ここで挙げたすべてを使うわけではない）。

\medskip\noindent
$S$ と $Y$ はアフィンと仮定する。
% P05-EN-CH37-0341
アフィン Noether スキーム $S_i$ の余フィルター極限として $S = \lim S_i$ と書く。
Limits, Lemma \ref{limits-lemma-descend-finite-presentation} により、
元 $0 \in I$ と図式
$$
\xymatrix{
X_0 \ar[rr]_{f_0} \ar[rd] & & Y_0 \ar[ld] \\
& S_0
}
$$
が存在する。これは有限型なスキームの射からなり、$S$ への基底変換が
補題の図式である。$0$ を大きく取り直すことで、$Y_0$ はアフィンであり、
% P05-EN-CH37-0342
$X_0 \to S_0$ は固有であると仮定できる。Limits, Lemmas
\ref{limits-lemma-eventually-proper} and
\ref{limits-lemma-limit-affine} を参照せよ。
$s$ の像を $s_0 \in S_0$ とする。
$Y_s$ はアフィンなので、$R^1f_{s, *}\mathcal{O}_{X_s} = 0$ は
$H^1(X_s, \mathcal{O}_{X_s}) = 0$ と同値である。
$X_s$ は忠実平坦写像 $\kappa(s_0) \to \kappa(s)$ による
$X_{0, s_0}$ の基底変換なので、
$H^1(X_{0, s_0}, \mathcal{O}_{X_{0, s_0}}) = 0$、したがって
$R^1f_{0, *}\mathcal{O}_{X_{0, s_0}} = 0$ である。
同様に、$\mathcal{O}_{Y_s} \to f_{s, *}\mathcal{O}_{X_s}$ が同型なので、
$\mathcal{O}_{Y_{0, s_0}} \to f_{0, *}\mathcal{O}_{X_{0, s_0}}$ も同型である。
$X_s \to Y_s$ のファイバーの次元は高々 $1$ なので、
射 $X_{0, s_0} \to Y_{0, s_0}$ についても同じことが成り立つ。
最後に、$X \to S$ は平坦なので、$0$ を大きく取り直すことで
$X_0$ は $S_0$ 上平坦と仮定できる。Limits, Lemma
\ref{limits-lemma-descend-flat-finite-presentation} を参照せよ。
したがって、$X_0 \to Y_0 \to S_0$ と点 $s_0$ に対して補題を証明すれば十分である。

\medskip\noindent
以上の簡約を組み合わせると、$S$ と $Y$ はアフィン、$S$ は Noether、
$f$ のファイバーの次元は $\leq 1$、かつ
$R^1f_*\mathcal{O}_X = 0$ は普遍的に成り立つ場合に帰着する。
点 $y \in Y_s$ をとる。次を主張する：
$$
\mathcal{O}_{Y, y} \longrightarrow (f_*\mathcal{O}_X)_y
$$
は同型である。この主張から補題が従う。実際、$f_*\mathcal{O}_X$ は連接なので
（Cohomology of Schemes, Proposition
\ref{coherent-proposition-proper-pushforward-coherent})
）、この主張により $Y$ を $y$ の開近傍で置き換えて、同型
$\mathcal{O}_Y \to f_*\mathcal{O}_X$ を得られる。
すると Lemma \ref{more-morphisms-lemma-h1-fibre-zero-h0-flat} により、$Y$ は $S$ 上平坦である。
最後に、Lemma \ref{more-morphisms-lemma-h1-fibre-zero-h0-flat} の最後の主張により、同型
$\mathcal{O}_Y \to f_*\mathcal{O}_X$ は任意の基底変換 $S' \to S$ の後にも同型のままである。

\medskip\noindent
主張を証明する。すでに、
$\mathcal{O}_{Y, y} \longrightarrow (f_*\mathcal{O}_X)_y$
は全射であり (Lemma \ref{more-morphisms-lemma-h1-fibre-zero-h0-kappa})、
$(f_*\mathcal{O}_X)_y$ は $\mathcal{O}_{S, s}$ 上平坦であり
(Lemma \ref{more-morphisms-lemma-h1-fibre-zero-h0-flat})、さらに誘導される写像
$$
\mathcal{O}_{Y_s, y} =  \mathcal{O}_{Y, y}/\mathfrak m_s\mathcal{O}_{Y, y}
\longrightarrow
(f_*\mathcal{O}_X)_y/\mathfrak m_s (f_*\mathcal{O}_X)_y
\to
(f_{s, *}\mathcal{O}_{X_s})_y
$$
は補題の仮定により単射であることを知っている。
よって Algebra, Lemma \ref{algebra-lemma-mod-injective} から、
$\mathcal{O}_{Y, y} \longrightarrow (f_*\mathcal{O}_X)_y$ は所望のとおり単射である。
\end{proof}

\begin{lemma}
\label{more-morphisms-lemma-bijection-on-Pic}
Noether スキームの固有射 $f : X \to Y$ が
$f_*\mathcal{O}_X = \mathcal{O}_Y$ を満たし、$f$ のファイバーの次元が
$\leq 1$ であり、すべての $y \in Y$ に対して
$H^1(X_y, \mathcal{O}_{X_y}) = 0$ であるとする。
このとき $f^* : \Pic(Y) \to \Pic(X)$ は、すべての $y \in Y$ に対して
$\mathcal{L}|_{X_y} \cong \mathcal{O}_{X_y}$ を満たす
$\mathcal{L} \in \Pic(X)$ のなす部分群への全単射である。
\end{lemma}

\begin{proof}
射影公式 (Cohomology, Lemma \ref{cohomology-lemma-projection-formula}) により、
$\mathcal{N} \in \Pic(Y)$ に対して
$f_*f^*\mathcal{N} \cong \mathcal{N}$ である。
すべての $y \in Y$ に対して
$\mathcal{L}|_{X_y} \cong \mathcal{O}_{X_y}$ を満たす
$\mathcal{L} \in \Pic(X)$ について、$\mathcal{N} = f_*\mathcal{L}$ は可逆であり、
$\mathcal{L} \cong f^*\mathcal{N}$ であると主張する。
これを示せば証明は完了する。

\medskip\noindent
$\mathcal{O}_Y$-加群 $\mathcal{N} = f_*\mathcal{L}$ は
Cohomology of Schemes, Proposition
\ref{coherent-proposition-proper-pushforward-coherent} により連接である。
したがって、これが可逆 $\mathcal{O}_Y$-加群であることを示すには、茎で確認すれば十分である
（Algebra, Lemma \ref{algebra-lemma-finite-projective}）。
Noether 局所環からその完備化への写像は忠実平坦なので、完備化
$(f_*\mathcal{L})_y^\wedge$ が自由であることを確認すれば十分である
（Algebra, Section \ref{algebra-section-completion-noetherian} and
Lemma \ref{algebra-lemma-finite-projective-descends} を参照せよ）。
このために、Cohomology of Schemes, Lemma
\ref{coherent-lemma-formal-functions-stalk} の形で述べられた形式関数定理を用いる。
$f_*\mathcal{O}_X = \mathcal{O}_Y$、したがって
$(f_*\mathcal{O}_X)_y^\wedge \cong \mathcal{O}_{Y, y}^\wedge$ なので、
% P05-EN-CH37-0343
各 $n$ に対して $\mathcal{L}|_{X_n} \cong \mathcal{O}_{X_n}$ が成り立ち、
$n$ を変えたときに互換であることを示せば十分である。
Lemma \ref{more-morphisms-lemma-picard-group-first-order-thickening} により、完全列
$$
H^1(X_y, \mathfrak m_y^n\mathcal{O}_X/\mathfrak m_y^{n + 1}\mathcal{O}_X)
\to \Pic(X_{n + 1}) \to \Pic(X_n)
$$
を得る。記号は形式関数定理のものとする。全射
$$
\mathcal{O}_{X_y}^{\oplus r_n} \cong
\mathfrak m_y^n/\mathfrak m_y^{n + 1} \otimes_{\kappa(y)} \mathcal{O}_{X_y}
\longrightarrow
\mathfrak m_y^n\mathcal{O}_X/\mathfrak m_y^{n + 1}\mathcal{O}_X
$$
がある整数 $r_n \geq 0$ に対して存在する。
$\dim(X_y) \leq 1$ なので、この全射は第1コホモロジー群上の全射を誘導する
（Cohomology, Proposition
\ref{cohomology-proposition-vanishing-Noetherian} による次数 $\geq 2$ の
コホモロジーの消滅を用いる）。したがって完全列中の $H^1$ は零であり、
遷移写像 $\Pic(X_{n + 1}) \to \Pic(X_n)$ は所望のとおり単射である。

\medskip\noindent
なお $f^*\mathcal{N} \cong \mathcal{L}$ を示す必要がある。
これは同じ方法で証明されるので、詳細を省略する。
\end{proof}









\section{アフィン層別}
\label{more-morphisms-section-affine-stratifications}

\noindent
この内容は \cite{RV} から採られている。本節を読む前に、
Topology, Section \ref{topology-section-stratifications} で
層別について少し読んでおいてほしい。

\medskip\noindent
$X$ がスキームならば、$X$ の層別とは通常、$X$ の台位相空間の層別をいう。
各層は局所閉部分集合である。Schemes, Remark
\ref{schemes-remark-reduced-induced-locally-closed} を用いて、
これらの層を $X$ の被約局所閉部分スキームとみなす。

\begin{definition}
\label{more-morphisms-definition-affine-stratification}
スキーム $X$ をとる。{\it アフィン層別}とは、局所有限な層別
$X = \coprod_{i \in I} X_i$ であって、各層 $X_i$ がアフィンであり、
包含射 $X_i \to X$ がアフィンとなるものをいう。
\end{definition}

\noindent
包含射 $X_i \to X$ が準コンパクトであるという条件のもとでは、層別
$X = \coprod X_i$ が局所有限であることは、各層が $X$ の局所構成可能部分集合で
あることと同値である。Properties, Lemma
\ref{properties-lemma-stratification-locally-finite-constructible} を参照せよ。

\medskip\noindent
% P05-EN-CH37-0344
$X_i \to X$ がアフィン射であるという条件は、局所閉部分集合 $X_i$ に
どのスキーム構造を入れるかに依存しない。Lemma
\ref{more-morphisms-lemma-thicken-property-morphisms} を参照せよ。さらに、$X$ が分離的
（または、より一般に対角射がアフィン）であり、$X = \coprod X_i$ が
アフィンな層をもつ局所有限層別ならば、射 $X_i \to X$ はアフィンである。
Morphisms, Lemma \ref{morphisms-lemma-affine-permanence} を参照せよ。
したがって、関心のある多くの場合には、包含射 $X_i \to X$ の
アフィン性という条件を考慮しなくてよい。

\medskip\noindent
層別の半順序添字集合 $I$ が有限である場合がしばしば重要となる。
半順序集合 $I$ の {\it 長さ}とは、$I$ の元からなる鎖
$i_0 < i_1 < \ldots < i_p$ の長さ $p$ の上限であることを思い出そう。

\begin{lemma}
\label{more-morphisms-lemma-improve-stratification}
スキーム $X$ と有限アフィン層別 $X = \coprod_{i \in I} X_i$ をとる。
$n$ を $I$ の長さとすると、添字集合が $\{0, \ldots, n\}$ である
アフィン層別が存在する。
\end{lemma}

\begin{proof}
$I$ には、すべての $i \in I$ に対して $X_i$ の閉包が
$\bigcup_{j \leq i} X_j$ に含まれるような半順序があることを思い出そう。
$I$ の極大添字全体を $I' \subset I$ とする。
$i \in I'$ ならば、他の層の閉包の和集合が $X_i$ の補集合なので、
$X_i$ は $X$ で開である。$U = \bigcup_{i \in I'} X_i$ を
$X$ の開部分スキームとみなすと、スキームとして
$U_{red} = \coprod_{i \in I'} X_i$ である。
Schemes, Lemma \ref{schemes-lemma-disjoint-union-affines} と
Lemma \ref{more-morphisms-lemma-thickening-affine-scheme} により、$U$ はアフィンスキームである。
さらに、各 $i \in I'$ に対して $X_i \to X$ がアフィンなので、
Lemma \ref{more-morphisms-lemma-thicken-property-morphisms} を用いる同じ議論により
射 $U \to X$ もアフィンである。補集合 $Z = X \setminus U$ に
誘導される被約スキーム構造を入れると、これはアフィン層別
$Z = \bigcup_{i \in I \setminus I'} X_i$ をもつ。
ここでは、スキームの射 $T \to Z$ がアフィンであることと、合成
$T \to X$ がアフィンであることが同値であることを用いた。これは
Morphisms, Lemmas \ref{morphisms-lemma-closed-immersion-affine},
\ref{morphisms-lemma-composition-affine}, and
\ref{morphisms-lemma-affine-permanence} から従う。
半順序集合 $I \setminus I'$ の長さは $I$ の長さよりちょうど1小さい。
したがって帰納法により、$Z$ はアフィン層別
$Z = Z_0 \amalg \ldots \amalg Z_{n - 1}$ をもち、その添字集合は
% P05-EN-CH37-0345
$\{1, \ldots, n\}$ である。$Z_n = U$ とおけば、$X$ の所望の層別を得る。
\end{proof}

\noindent
スキーム $X$ が有限アフィン層別をもてば、もちろん $X$ は準コンパクトである。
それほど明らかではないが、この条件は $X$ が準分離であることも含意する。

\begin{lemma}
\label{more-morphisms-lemma-qc-affine-stratification}
スキーム $X$ に対して、次は同値である：
\begin{enumerate}
\item $X$ は有限アフィン層別をもつ、
\item $X$ は準コンパクトかつ準分離である。
\end{enumerate}
\end{lemma}

\begin{proof}
$X = \bigcup X_i$ を有限アフィン層別とする。各 $X_i$ はアフィン、したがって
準コンパクトなので、$X$ は準コンパクトである。
アフィン開集合 $U, V \subset X$ をとる。$X_i \to X$ はアフィン射なので、
$U \cap X_i$ と $V \cap X_i$ は $X_i$ のアフィン開集合である。したがって
$U \cap V \cap X_i$ はアフィンスキーム $X_i$ のアフィン開集合である
（例えば Schemes, Lemma \ref{schemes-lemma-characterize-separated} を参照せよ）。
よって $U \cap V = \coprod U \cap V \cap X_i$ は有限個のアフィン層の和として
準コンパクトである。Schemes, Lemma
\ref{schemes-lemma-characterize-quasi-separated} により、$X$ は準分離である。

\medskip\noindent
$X$ は準コンパクトかつ準分離であると仮定する。$X$ が有限アフィン層別を
もつという主張を証明するため、Cohomology of Schemes, Lemma
\ref{coherent-lemma-induction-principle} の帰納原理を用いてよい。
$X$ が空ならば空のアフィン層別をもつ。$X$ が空でないアフィンスキームならば、
1つの層からなるアフィン層別をもつ。次に、$X = U \cup V$ で、$U$ は
準コンパクト開、$V$ はアフィン開であり、
% P05-EN-CH37-0346
有限アフィン層別 $U = \bigcup_{i \in I} U_i$ および
$U \cap V = \coprod_{j \in J} W_j$ が与えられていると仮定する。
$Z = X \setminus V$ および $Z' = X \setminus U$ と書く。
$Z$ は $U$ で閉であり、$Z'$ は $V$ で閉である。
$U_i \cap Z$ と $U_i \cap W_j = U_i \times_U W_j$ はアフィンスキームであり、
$U$ 上アフィンである。（ヒント：$U_i \times_U W_j \to W_j$ は
$U_i \to U$ の基底変換としてアフィンであり、したがって $U_i \cap W_j$ は
アフィン、ゆえに $U_i \cap W_j \to U_i$ はアフィン、したがって
$U_i \cap W_j \to U$ はアフィンである。）よって
$$
U = \coprod\nolimits_{i \in I} (U_i \cap Z) \amalg
\coprod\nolimits_{(i, j) \in I \times J} (U_i \cap W_j)
$$
は有限アフィン層別である。$I \amalg I \times J$ 上の半順序は、
$i' \leq (i, j) \Leftrightarrow i' \leq i$ および
$(i', j') \leq (i, j) \Leftrightarrow i' \leq i$ かつ $j' \leq j$
で与えられる。$(U_i \cap Z) \times_X V = \emptyset$ および
$(U_i \cap W_j) \times_X V = U_i \cap W_j$ はアフィンである。
したがって、射 $U_i \cap Z \to X$ と $U_i \cap W_j \to X$ はアフィンである。
実際、射のアフィン性は標的上局所的に確認でき
(Morphisms, Lemma \ref{morphisms-lemma-characterize-affine})
、しかも $U$ と $V$ の双方の上でアフィン性が成り立つ。
証明を終えるため、上の層別にもう1つの層 $Z'$ を加え、
その添字を半順序集合の極小元として加える。
\end{proof}

\begin{definition}
\label{more-morphisms-definition-affine-stratification-number}
空でない準コンパクト準分離スキーム $X$ をとる。
{\it アフィン層別数}とは、次の同値な条件を満たす最小の整数 $n \geq 0$ である：
\begin{enumerate}
\item $I$ の長さが $n$ である有限アフィン層別
$X = \coprod_{i \in I} X_i$ が存在する、
\item 添字集合 $\{0, \ldots, n\}$ をもつアフィン層別
$X = X_0 \amalg X_1 \amalg \ldots \amalg X_n$ が存在する。
\end{enumerate}
\end{definition}

\noindent
条件の同値性は Lemma \ref{more-morphisms-lemma-improve-stratification} による。
有限アフィン層別の存在は Lemma
\ref{more-morphisms-lemma-qc-affine-stratification} で証明されている。

\begin{lemma}
\label{more-morphisms-lemma-affine-stratification-number-bound}
$n + 1$ 個のアフィン開集合による開被覆をもつ分離スキーム $X$ をとる。
このとき $X$ のアフィン層別数は高々 $n$ である。
\end{lemma}

\begin{proof}
$X = U_0 \cup \ldots \cup U_n$ をアフィン開被覆とする。次のようにおく：
$$
X_i = (U_i \cup \ldots \cup U_n) \setminus
(U_{i + 1} \cup \ldots \cup U_n)
$$
すると $X_i$ は $U_i$ の閉部分スキームとしてアフィンである。
Morphisms, Lemma \ref{morphisms-lemma-affine-permanence} により、
射 $X_i \to X$ はアフィンである。最後に、
% P05-EN-CH37-0347
$\overline{X_i} \subset X_i \cup X_{i - 1} \cup \ldots X_0$ である。
\end{proof}

\begin{lemma}
\label{more-morphisms-lemma-affine-stratification-number-bound-Noetherian}
次元 $\infty > d \geq 0$ の Noether スキーム $X$ をとる。
このとき $X$ のアフィン層別数は高々 $d$ である。
\end{lemma}

\begin{proof}
$d$ に関する帰納法で証明する。$d = 0$ ならば $X$ はアフィンである。
Properties, Lemma \ref{properties-lemma-locally-Noetherian-dimension-0} を参照せよ。
$d > 0$ と仮定する。$X$ の既約成分の生成点を
$\eta_1, \ldots, \eta_n$ とする（Properties, Lemma
\ref{properties-lemma-Noetherian-irreducible-components}）。
$\eta_1, \ldots, \eta_n$ を含むアフィン開集合で $X$ を被覆できる。
Properties, Lemma \ref{properties-lemma-point-and-maximal-points-affine} を参照せよ。
$X$ は準コンパクトなので、すべての $j = 1, \ldots, m$ に対して
$\eta_1, \ldots, \eta_n \in U_j$ となる有限アフィン開被覆
$X = \bigcup_{j = 1, \ldots, m} U_j$ が存在する。
$\eta_1, \ldots, \eta_n$ を含むアフィン開集合
$U \subset U_1 \cap \ldots \cap U_m$ を選ぶ（先ほど引用した補題により可能である）。
すべての $j$ に対して $U \to U_j$ はアフィンなので、射 $U \to X$ はアフィンである。
Morphisms, Lemma \ref{morphisms-lemma-characterize-affine} を参照せよ。
$Z = X \setminus U$ とおく。構成により $\dim(Z) < \dim(X)$ である。
帰納法の仮定により、$n \leq \dim(Z)$ を満たす $Z$ のアフィン層別
% P05-EN-CH37-0348
$Z = \bigcup_{i \in \{0, \ldots, n\}} Z_i$ が存在する。
$U = X_{n + 1}$ および $i \leq n$ に対して $X_i = Z_i$ とおけば主張が従う。
\end{proof}

\begin{proposition}
\label{more-morphisms-proposition-vanishing-affine-stratification-number}
アフィン層別数 $n$ の、空でない準コンパクト準分離スキーム $X$ をとる。
このとき、任意の準連接 $\mathcal{O}_X$-加群 $\mathcal{F}$ と $p > n$ に対して
$H^p(X, \mathcal{F}) = 0$ である。
\end{proposition}

\begin{proof}
アフィン層別数 $n$ に関する帰納法で証明する。
$n = 0$ ならば $X$ はアフィンであり、結果は Cohomology of Schemes, Lemma
\ref{coherent-lemma-quasi-coherent-affine-cohomology-zero} である。
$n > 0$ と仮定する。Definition \ref{more-morphisms-definition-affine-stratification-number} により、
アフィンスキーム $U$ とアフィン開埋め込み $j : U \to X$ であって、
補集合 $Z$ のアフィン層別数が $n - 1$ となるものが存在する。
$U$ と $j$ はアフィンなので、$p > 0$ に対して
$H^p(X, j_*(\mathcal{F}|_U)) = 0$ である。Cohomology of Schemes, Lemmas
\ref{coherent-lemma-relative-affine-cohomology} and
\ref{coherent-lemma-relative-affine-vanishing} を参照せよ。
写像 $\mathcal{F} \to j_*(\mathcal{F}|_U)$ の核と余核をそれぞれ
$\mathcal{K}$ と $\mathcal{Q}$ と書く。すると、準連接
$\mathcal{O}_X$-加群の完全列
$$
0 \to \mathcal{K} \to \mathcal{F} \to j_*(\mathcal{F}|_U)
\to \mathcal{Q} \to 0
$$
を得る（Schemes, Section \ref{schemes-section-quasi-coherent} を参照せよ）。
この完全列を短完全列に分解してコホモロジー長完全列を用いる標準的な議論により、
$p \geq n$ に対して $H^p(X, \mathcal{K}) = 0$ および
$H^p(X, \mathcal{Q}) = 0$ を示せば十分である。
$\mathcal{F} \to j_*(\mathcal{F}|_U)$ の $U$ への制限は同型なので、
$\mathcal{K}$ と $\mathcal{Q}$ は $Z$ 上に台をもつ。
Properties, Lemma
\ref{properties-lemma-quasi-coherent-colimit-finite-type}
により、これらの加群をその有限型準連接部分加群のフィルター余極限として書ける。
$X$ 上の層のコホモロジーはフィルター余極限と可換なので
（Cohomology, Lemma \ref{cohomology-lemma-quasi-separated-cohomology-colimit}）、
台が $Z$ に含まれる有限型準連接加群 $\mathcal{G}$ に対して、
$p \geq n$ ならば $H^p(X, \mathcal{G}) = 0$ であることを示せば十分である。
$\mathcal{G} \oplus \mathcal{O}_Z$ のスキーム論的台を $Z' \subset X$ とする。
$\mathcal{G}$ を $Z'$ 上の準連接加群とみなしてよく、実際そうみなす。
Morphisms, Section \ref{morphisms-section-support} を参照せよ。
このとき $Z'$ と $Z$ の台位相空間は同じであり、したがってアフィン層別数も同じで、
$n - 1$ である。よって
$H^p(X, \mathcal{G}) = H^p(Z', \mathcal{G})$ であり
（この等式は Cohomology of Schemes, Lemma
\ref{coherent-lemma-relative-affine-cohomology})
、帰納法の仮定により $p \geq n$ に対して消滅する。
\end{proof}

\begin{example}
\label{more-morphisms-example-affine-stratification-number}
体 $k$ をとり、$X = \mathbf{P}^n_k$ を $k$ 上の $n$ 次元射影空間とする。
Constructions, Lemma
\ref{constructions-lemma-standard-covering-projective-space} により、
Lemma \ref{more-morphisms-lemma-affine-stratification-number-bound} を適用できる。
したがって $\mathbf{P}^n_k$ のアフィン層別数は高々 $n$ である。
一方、$\mathbf{P}^n_k$ 上のある準連接加群について次数 $n$ の
コホモロジーは非零である。Cohomology of Schemes, Lemma
\ref{coherent-lemma-cohomology-projective-space-over-ring} を参照せよ。
Proposition \ref{more-morphisms-proposition-vanishing-affine-stratification-number} により、
$\mathbf{P}^n_k$ のアフィン層別数は $n$ に等しい。
\end{example}








\section{普遍開射}
\label{more-morphisms-section-universally-open}

\noindent
普遍開射に関するいくつかの結果を述べる。

\begin{lemma}
\label{more-morphisms-lemma-test-universally-open}
スキームの射 $f : X \to S$ をとる。次は同値である：
\begin{enumerate}
\item $f$ は普遍開である、
\item 有限表示で局所的な任意の射 $S' \to S$ に対して、基底変換
$X_{S'} \to S'$ は開である、
\item 任意の $n$ に対して射
$\mathbf{A}^n \times X \to \mathbf{A}^n \times S$ は開である。
\end{enumerate}
\end{lemma}

\begin{proof}
(1) が (2) を、(2) が (3) を含意することは明らかである。
(3) が (1) を含意することを証明しよう。
あるスキームの射 $g : T \to S$ に対して、基底変換
$X_T \to T$ が開でないと仮定する。このとき、アフィン開集合
$V \subset S$, $U \subset X$, $W \subset T$ で、
$f(U) \subset V$ および $g(W) \subset V$ を満たし、
$U \times_V W \to W$ が開でないものを見つけられる。
これから $\mathbf{A}^n \times U \to \mathbf{A}^n \times V$ が開でないことを
示せれば、$\mathbf{A}^n \times X \to \mathbf{A}^n \times S$ も開でなく、
証明が完了する。したがって、次の段落で証明する結果に帰着する。

\medskip\noindent
環準同型 $A \to B$ をとり、ある $A$-代数 $A'$ に対して
$A' \to B' = A' \otimes_A B$ がスペクトルの開写像を誘導しないとする。
主開集合は $\Spec(B')$ の位相の基底をなすので、ある $g \in B'$ に対して
$D(g)$ の $\Spec(A')$ における像は開でない。次のように書く：
$g = \sum_{i = 1, \ldots, n} a'_i \otimes b_i$
ここで、ある $n$ に対して $a'_i \in A'$ かつ $b_i \in B$ である。
$B[x_1, \ldots, x_n]$ の元
$h = \sum_{i = 1, \ldots, n} x_i b_i$ を考える。矛盾を導くため、
射
$$
\Spec(B[x_1, \ldots, x_n]) \longrightarrow \Spec(A[x_1, \ldots, x_n])
$$
のもとで $D(h)$ が開部分集合に写ると仮定する。すると $D(h)$ は
準コンパクト開集合 $U \subset \Spec(A[x_1, \ldots, x_n])$ の上に全射的に写る。
$x_i$ を $a'_i$ に送る $A$-代数準同型
$A[x_1, \ldots, x_n] \to A'$ をとる。これはまた、$h$ を $g$ に送る
$B$-代数準同型 $B[x_1, \ldots, x_n] \to B'$ を誘導する。図式
$$
\xymatrix{
\Spec(B[x_1, \ldots, x_n]) \ar[d] &
\Spec(B') \ar[l] \ar[d] \\
\Spec(A[x_1, \ldots, x_n]) &
\Spec(A') \ar[l]
}
$$
は Cartesian なので、$D(g)$ の $\Spec(A')$ における像は
$U$ の $\Spec(A')$ における逆像に等しく、したがって開である。
これは所望の矛盾である。
\end{proof}

\begin{lemma}
\label{more-morphisms-lemma-quasi-finite-Noetherian-universally-open}
スキームの射 $f : X \to Y$ が次を満たすとする：
\begin{enumerate}
\item $f$ は局所準有限である、
\item $Y$ は幾何学的単枝かつ局所 Noether である、
\item $X$ の各既約成分は $Y$ のある既約成分を支配する。
\end{enumerate}
このとき $f$ は普遍開である。
\end{lemma}

\begin{proof}
Lemma \ref{more-morphisms-lemma-number-of-branches-and-smooth} と Properties, Lemma
\ref{properties-lemma-number-of-branches-1} により、任意の $n$ に対して
スキーム $\mathbf{A}^n \times Y$ は幾何学的単枝である。したがって、すべての
$n$ に対して射 $\mathbf{A}^n \times X \to \mathbf{A}^n \times Y$ は
補題の仮定を満たす。Lemma \ref{more-morphisms-lemma-test-universally-open} により、
$f$ が開であることを証明すれば十分である。Morphisms, Lemma
\ref{morphisms-lemma-locally-finite-presentation-universally-open} により、
一般化が $f$ に沿って持ち上がることを示せば十分である。
$Y$ の点の特殊化 $y' \leadsto y$ と、$y$ に写る点 $x \in X$ をとる。
Lemma \ref{more-morphisms-lemma-etale-makes-quasi-finite-finite-at-point} と同様に、図式
$$
\xymatrix{
u \ar[d] & U \ar[d] \ar[r] & X \ar[d] \\
v & V \ar[r] & Y
}
$$
を選ぶ。ここで $(V, v) \to (Y, y)$ は初等 \'etale 近傍、
$U \to V$ は有限、$u$ は $v$ に写る $U$ の唯一の点、
$U \subset V \times_Y X$ は開であり、$v \mapsto y$ かつ $u \mapsto x$ である。
$u$ を通る $U$ の既約成分 $E$ をとる（少なくとも1つ存在する）。
$U \to X$ は \'etale なので、$E$ は $X$ のある既約成分に写り、
仮定によりその成分は $Y$ のある既約成分を支配する。
$U \to V$ は有限、したがって閉なので、$E$ の像 $E' \subset V$ は
$v$ を通り、$Y$ のある既約成分を支配する既約閉部分集合である。
$V \to Y$ は \'etale なので、$E'$ は $v$ を通る $V$ の既約成分でなければならない。
$Y$ は幾何学的単枝なので、$E'$ は $v$ を通る $V$ の唯一の既約成分である
(Lemma \ref{more-morphisms-lemma-nr-branches})。$V$ は局所 Noether なので、$V$ を縮小して
$E' = V$（集合としての等号）と仮定できる。

\medskip\noindent
$V \to Y$ は \'etale なので、像が $y' \leadsto y$ である特殊化
$v' \leadsto v$ を見つけられる。上の結果により、$v'$ に写る
$u' \in U$ が存在する。$u$ は $v$ に写る $U$ の唯一の点であり、
$U \to V$ は閉なので、$u' \leadsto u$ である。最後に、$u'$ の像
$x' \in X$ は $x$ に特殊化し、$y'$ に写る点である。これで証明が完了する。
\end{proof}

\begin{lemma}
\label{more-morphisms-lemma-characterize-universally-open-finite}
環準同型 $A \to B$ をとる。$B$ は $A$-加群として
$b_1, \ldots, b_d \in B$ により生成されるとする。
$h = \sum x_ib_i \in B[x_1, \ldots, x_d]$ とおく。このとき、
$\Spec(B) \to \Spec(A)$ が普遍開であることと、$D(h)$ の
$\Spec(A[x_1, \ldots, x_d])$ における像が開であることは同値である。
\end{lemma}

\begin{proof}
$\Spec(B) \to \Spec(A)$ が普遍開ならば、もちろん $D(h)$ の像は開である。
逆に、$D(h)$ の像 $U$ が開であると仮定する。環準同型 $A \to A'$ をとる。
任意の主開集合 $D(g) \subset \Spec(A' \otimes_A B)$ の
$\Spec(A')$ における像が開であることを示せば十分である。
ある $a'_i \in A'$ によって
$g = \sum_{i = 1, \ldots, d} a'_i \otimes b_i$ と書ける。
% P05-EN-CH37-0349
$x_i$ を $a'_i$ に送る $A$-代数準同型
$A[x_1, \ldots, x_n] \to A'$ をとる。これはまた、$h$ を $g$ に送る
$B$-代数準同型 $B[x_1, \ldots, x_n] \to A' \otimes_A B$ を誘導する。図式
$$
\xymatrix{
\Spec(B[x_1, \ldots, x_n]) \ar[d] &
% P05-EN-CH37-0350
\Spec(B') \ar[l] \ar[d] \\
\Spec(A[x_1, \ldots, x_n]) &
\Spec(A') \ar[l]
}
$$
は Cartesian なので、$D(g)$ の $\Spec(A')$ における像は
$U$ の $\Spec(A')$ における逆像に等しく、したがって開である。
\end{proof}

\begin{lemma}
\label{more-morphisms-lemma-descend-quasi-finite-universally-open}
アフィン遷移射をもつスキームの有向系の極限を $S = \lim S_i$ とする。
$0 \in I$ をとり、$f_0 : X_0 \to Y_0$ を $S_0$ 上のスキームの射とする。
$S_0$, $X_0$, $Y_0$ は準コンパクトかつ準分離であると仮定する。
$f_0$ の $S_i$ への基底変換を $f_i : X_i \to Y_i$、
$f_0$ の $S$ への基底変換を $f : X \to Y$ とする。次を仮定する：
\begin{enumerate}
\item $f$ は局所準有限かつ普遍開である、
\item $f_0$ は有限表示で局所的である。
\end{enumerate}
このとき、$f_i$ が局所準有限かつ普遍開となる $i \geq 0$ が存在する。
\end{lemma}

\begin{proof}
Limits, Lemma \ref{limits-lemma-descend-quasi-finite} により、$0$ を大きく
取り直せば $f_0$ は局所準有限と仮定できる。$x \in X$ をとる。
準有限射の \'etale 局所化により、図式
$$
\xymatrix{
X \ar[d] & U \ar[l] \ar[d] \\
Y & V \ar[l]
}
$$
を見つけられる。ここで $V \to Y$ は \'etale、$U \subset X_V$ は開、
$U \to V$ は有限であり、$x$ は $U \to X$ の像に含まれる。
Lemma \ref{more-morphisms-lemma-etale-makes-quasi-finite-finite-at-point} を参照せよ。
$V$ を縮小して、$V$ と $U$ はアフィンと仮定できる。
$X$ は準コンパクトなので、このような $V$ と $U$ の有限非交和をとることで、
$U \to X$ が全射となる上の形の図式を作れる。Limits, Lemmas
\ref{limits-lemma-descend-finite-presentation},
\ref{limits-lemma-descend-opens},
\ref{limits-lemma-descend-surjective},
\ref{limits-lemma-descend-finite-finite-presentation},
\ref{limits-lemma-descend-etale}, and
\ref{limits-lemma-limit-affine} により、必要なら $0$ を大きく取り直して、図式
$$
\xymatrix{
X_0 \ar[d] & U_0 \ar[l] \ar[d] \\
Y_0 & V_0 \ar[l]
}
$$
があると仮定できる。ここで $V_0$ はアフィン、$V_0 \to Y_0$ は \'etale、
$U_0 \subset (X_0)_{V_0}$ は開、$U_0 \to V_0$ は有限、
$U_0 \to X_0$ は全射である。$V_i \to Y_i$ は \'etale、したがって普遍開なので、
% P05-EN-CH37-0351
十分大きい $i$ に対して $U_i \to V_i$ が普遍開であることを証明すれば十分である。
したがって、次の段落で論じる場合に帰着する。

\medskip\noindent
環のフィルター余極限を $A = \colim A_i$ とする。環準同型 $A_0 \to B_0$ をとり、
$B = A \otimes_{A_0} B_0$ および $B_i = A_i \otimes_{A_0} B_0$ とおく。
$A_0 \to B_0$ は有限かつ有限表示であり、$A \to B$ は普遍開と仮定する。
十分大きい $i$ に対して $A_i \to B_i$ が普遍開であることを示す必要がある。
$A_0$-加群として $B_0$ を生成する
$b_{0, 1}, \ldots, b_{0, d} \in B_0$ を選ぶ。
$B_0[x_1, \ldots, x_d]$ において
$h_0 = \sum_{j = 1, \ldots, d} x_jb_{0, j}$ とおく。
$h_0$ の $B[x_1, \ldots, x_d]$ における像を $h$、
$B_i[x_1, \ldots, x_d]$ における像を $h_i$ と書く。
$A \to B$ は普遍開なので、$D(h)$ の
$\Spec(A[x_1, \ldots, x_d])$ における像 $U$ は開である。
$U$ はアフィンスキームの像としてもちろん準コンパクトである。
十分大きい $i$ に対して、準コンパクト開集合
$U_i \subset \Spec(A_i[x_1, \ldots, x_d])$ で、その
$\Spec(A[x_1, \ldots, x_d])$ における逆像が $U$ となるものが存在する。
Limits, Lemma \ref{limits-lemma-descend-opens} を参照せよ。
$i$ を大きく取り直し、$D(h_i)$ が $U_i$ に写ると仮定できる。
これは、$D(h_i)$ における $U_i$ の引き戻しを考え、同じ補題を用いれば従う。
最後に、$i$ をさらに大きくすれば、すでに用いた Limits, Lemma
\ref{limits-lemma-descend-surjective} により、スキームの射
$D(h_i) \to U_i$ は全射となる。Lemma
\ref{more-morphisms-lemma-characterize-universally-open-finite} により、
$A_i \to B_i$ は普遍開である。
\end{proof}

\begin{lemma}
\label{more-morphisms-lemma-count-geometric-fibres}
局所準有限射 $f : X \to Y$ をとる。このとき
\begin{enumerate}
\item Lemmas \ref{more-morphisms-lemma-base-change-fibres-nr-geometrically-irreducible-components}
and \ref{more-morphisms-lemma-base-change-fibres-nr-geometrically-connected-components}
の関数 $n_{X/Y}$ は一致する、
\item $X$ が準コンパクトならば、$n_{X/Y}$ は最大値 $d < \infty$ をとる。
\end{enumerate}
\end{lemma}

\begin{proof}
局所準有限射の（幾何学的）ファイバーが離散的であることから、
2つの関数の一致は直ちに従う。Morphisms, Lemma
\ref{morphisms-lemma-locally-quasi-finite-fibres} を参照せよ。
有界性は Morphisms, Lemmas
\ref{morphisms-lemma-characterize-universally-bounded} and
\ref{morphisms-lemma-locally-quasi-finite-qc-source-universally-bounded} から従う。
\end{proof}

\begin{lemma}
\label{more-morphisms-lemma-count-geometric-fibres-universally-open}
スキームの分離的、局所準有限、普遍開な射 $f : X \to Y$ をとる。
$n_{X/Y}$ を Lemma \ref{more-morphisms-lemma-count-geometric-fibres} のものとする。
ある $y \in Y$ と $d \geq 0$ に対して $n_{X/Y}(y) \geq d$ ならば、
$y$ のある開近傍で $n_{X/Y} \geq d$ である。
\end{lemma}

\begin{proof}
問題は $Y$ 上局所的なので、$Y$ はアフィンと仮定してよい。
剰余体 $\kappa(y)$ の代数閉包を $K$ とする。仮定は、
$(X_y)_K$ が $\geq d$ 個の連結成分をもつということである。
すると、適当な準コンパクト開集合 $X' \subset X$ に対して、
スキーム $(X'_y)_K$ は $\geq d$ 個の連結成分をもつ。詳細は省略する。
$X$ を $X'$ で置き換えて、$X$ は準コンパクトと仮定してよい。
すると $f$ は準有限である。$y$ の上にある $X$ の点を
$x_1, \ldots, x_n$ とする。Lemma
\ref{more-morphisms-lemma-etale-splits-off-quasi-finite-part-technical-variant} を適用して、
\'etale 近傍 $(U, u) \to (Y, y)$ と分解
$$
U \times_Y X =
W \amalg
\ \coprod\nolimits_{i = 1, \ldots, n}
\ \coprod\nolimits_{j = 1, \ldots, m_i}
V_{i, j}
$$
を得る。記号は引用箇所のものとする。この状況では
$n_{X/Y}(y) = \sum_i m_i$ である。詳細は省略する。
$f$ は普遍開なので、すべての $i, j$ に対して $V_{i, j} \to U$ は開である。
したがって $U$ を縮小し、すべての $i, j$ に対して
$V_{i, j} \to U$ は全射と仮定できる。よって
$n_{U \times_Y X/U} \geq \sum_i m_i = n_{X/Y}(y) \geq d$.
$n_{X/Y}$ の構成は基底変換と可換なので、証明が完了する。
\end{proof}

\begin{lemma}
\label{more-morphisms-lemma-large-open}
スキームの分離的、局所準有限、普遍開な射 $f : X \to Y$ をとる。
$n_{X/Y}$ を Lemma \ref{more-morphisms-lemma-count-geometric-fibres} のものとする。
$n_{X/Y}$ が最大値 $d < \infty$ をとるならば、集合
$$
Y_d = \{y \in Y \mid n_{X/Y}(y) = d\}
$$
は $Y$ で開であり、射 $f^{-1}(Y_d) \to Y_d$ は有限である。
\end{lemma}

\begin{proof}
$Y_d$ の開性は Lemma
\ref{more-morphisms-lemma-count-geometric-fibres-universally-open} から直ちに従う。
$Y_d$ 上での有限性を証明するため、その補題の証明の議論を繰り返す。
すなわち、$y \in Y_d$ をとる。このとき $y$ の上にある $X$ の点は高々
$d$ 個である。$y$ の上にある $X$ の点を $x_1, \ldots, x_n$ とする。
Lemma \ref{more-morphisms-lemma-etale-splits-off-quasi-finite-part-technical-variant} を適用して、
\'etale 近傍 $(U, u) \to (Y, y)$ と分解
$$
U \times_Y X =
W \amalg
\ \coprod\nolimits_{i = 1, \ldots, n}
\ \coprod\nolimits_{j = 1, \ldots, m_i}
V_{i, j}
$$
を得る。記号は引用箇所のものとする。この状況では
$d = n_{X/Y}(y) = \sum_i m_i$ である。詳細は省略する。
$f$ は普遍開なので、すべての $i, j$ に対して $V_{i, j} \to U$ は開である。
したがって $U$ を縮小し、すべての $i, j$ に対して
$V_{i, j} \to U$ は全射であると仮定でき、さらに
% P05-EN-CH37-0352
$U$ は $W$ に写ると仮定できる。これにより
$n_{U \times_Y X/U} \geq \sum_i m_i = d$ である。
$n_{X/Y}$ の構成は基底変換と可換なので、
$n_{U \times_Y X/U} = d$ である。これは $W$ が空でなければならないことを意味し、
$U \times_Y X \to U$ は有限である。Descent, Lemma
\ref{descent-lemma-descending-property-finite} により、$X \to Y$ は
開射 $U \to Y$ の像の上で有限である。すなわち、所望のとおり
$f$ は $y$ のある開近傍の上で有限である。
\end{proof}






\section{重み付け}
\label{more-morphisms-section-weightings}

\noindent
本節の内容は \cite[Exposee XVII, 6.2.4]{SGA4} から採られている。

\medskip\noindent
有限ファイバーをもつスキームの局所準有限射 $\pi : U \to V$ をとる。
関数 $w : U \to \mathbf{Z}$ に対し、関数
$$
\textstyle{\int}_\pi w : V \longrightarrow \mathbf{Z},\quad
v \longmapsto
\sum\nolimits_{u \in U,\ \pi(u) = v} w(u) [\kappa(u) : \kappa(v)]_s
$$
を定義する。体拡大は有限であり
(Morphisms, Lemma \ref{morphisms-lemma-residue-field-quasi-finite}),
$[\kappa' : \kappa]_s$ は分離次数であり
(Fields, Definition \ref{fields-definition-insep-degree}),
、$\pi$ のファイバーは有限と仮定したので和も有限である。
点 $v \in V$ における $\int_\pi w$ の値は、次のようにも計算できる。
代数閉体 $k$ と、像が $v$ である射
$\overline{v} : \Spec(k) \to V$ を選ぶ。このとき
$$
(\textstyle{\int}_\pi w)(v) =
\sum\nolimits_{\overline{u} \in U_{\overline{v}}} w(\overline{u})
$$
である。ここで $w(\overline{u})$ はもちろん、射
$U_{\overline{v}} \to U$ のもとで点 $\overline{u}$ の像となる点 $u$ における
$w$ の値を表す。$\overline{u} \in U_{\overline{v}}$ を、
$\pi \circ \overline{u} = \overline{v}$ を満たす射
$\overline{u} : \Spec(k) \to U$ とみなしてよい。実際、$U \to V$ は
有限ファイバーをもつ局所準有限射なので、スキーム $U_{\overline{v}}$ は
% P05-EN-CH37-0353
$k$ 上有限次元代数のスペクトルであり、そのすべての素イデアルは
剰余体 $k$ をもつ極大イデアルである。上の等式を確認するには、
与えられた $u$ の上にある射 $\overline{u}$ の個数が
Fields, Lemma \ref{fields-lemma-separable-degree} により
$[\kappa(u) : \kappa(v)]_s$ に等しいことに注意すればよい。

\begin{lemma}
\label{more-morphisms-lemma-weighting-check-after-etale-base-change}
Cartesian 正方形
$$
\xymatrix{
U \ar[d]_\pi & U' \ar[l]^h \ar[d]^{\pi'} \\
V & V' \ar[l]_g
}
$$
が与えられ、$\pi$ は有限ファイバーをもつ局所準有限射であり、
$w : U \to \mathbf{Z}$ は関数であるとする。このとき
$(\int_\pi w) \circ g = \int_{\pi'} (w \circ h)$ である。
\end{lemma}

\begin{proof}
これは、上で与えた $\int_\pi w$ の第2の記述から直ちに従う。
定義から証明するには、次を用いる。$E/F$ を有限体拡大、$F'/F$ を別の体拡大とし、
$(E \otimes_F F')_{red} = \prod E'_i$ を $F'$ 上有限な体の積として書けば、
$$
[E : F]_s = \sum [E'_i : F']_s
$$
である。この等式を証明するため、代数閉体拡大 $\Omega/F'$ を選び、次に注意する：
\begin{align*}
[E : F]_s
& =
|\Mor_F(E, \Omega)| \\
& =
|\Mor_{F'}(E \otimes_F F', \Omega)| \\
& =
|\Mor_{F'}((E \otimes_F F')_{red}, \Omega)| \\
& =
\sum |\Mor_{F'}(E'_i, \Omega)| \\
& =
\sum [E'_i : F']_s
\end{align*}
ここでは Fields, Lemma \ref{fields-lemma-separable-degree} を用いた。
\end{proof}

\begin{definition}
\label{more-morphisms-definition-weighting}
局所準有限射 $f : X \to Y$ をとる。$f$ の {\it 重み付け}、または
{\it pond\'eration} とは、任意の図式
$$
\xymatrix{
X \ar[d]_f & U \ar[l]^h \ar[d]^\pi \\
Y & V \ar[l]_g
}
$$
で $V \to Y$ が \'etale、$U \subset X_V$ が開、$U \to V$ が有限であるものに対して、
関数 $\int_\pi (w \circ h)$ が局所定数となるような写像
$w : X \to \mathbf{Z}$ をいう。
\end{definition}

\noindent
もちろん $w = 0$ とすれば、任意の局所準有限射 $f$ の重み付けが得られるが、
あまり興味深いものではない。さまざまな目的には、{\it 正の}重み付け、すなわち
$w : X \to \mathbf{Z}_{> 0}$ が最も興味深いことが後で分かる。

\begin{lemma}
\label{more-morphisms-lemma-weighting-base-change}
局所準有限射 $f : X \to Y$ と重み付け $w : X \to \mathbf{Z}$ をとる。
$f' : X' \to Y'$ を、射 $Y' \to Y$ による $f$ の基底変換とする。
このとき、$w$ と射影 $X' \to X$ の合成
$w' : X' \to \mathbf{Z}$ は $f'$ の重み付けである。
\end{lemma}

\begin{proof}
射 $f'$ について Definition \ref{more-morphisms-definition-weighting} に現れる図式
$$
\xymatrix{
X' \ar[d]_{f'} & U' \ar[l]^{h'} \ar[d]^{\pi'} \\
Y' & V' \ar[l]_{g'}
}
$$
を考える。任意の $v' \in V'$ に対して、$\int_{\pi'} (w' \circ h')$ が
$v'$ のある開近傍で定数であることを示す必要がある。
Lemma \ref{more-morphisms-lemma-weighting-check-after-etale-base-change} と
\'etale 射が開であることにより、$V'$ を $v'$ の任意の \'etale 近傍で
置き換えてよい。$V'$ を $v'$ の \'etale 近傍で置き換えた後、
$U' = U'_1 \amalg \ldots \amalg U'_n$ であり、各 $U'_i$ は
$v'$ の上に唯一の点 $u'_i$ をもち、$\kappa(u'_i)/\kappa(v')$ が
純非分離であると仮定できる。Lemma
\ref{more-morphisms-lemma-etale-splits-off-quasi-finite-part-technical-variant} を参照せよ。
$\int_{U'_i \to V'} w'|_{U'_i}$ が $v'$ のある近傍で定数であることを
示せば十分であることは明らかである。したがって次の段落の場合に帰着する。

\medskip\noindent
$v' \in V'$ とし、$v'$ の上にある $U'$ の唯一の点 $u'$ があり、
$\kappa(u')/\kappa(v')$ は純非分離であるとする。$u'$ と $v'$ の像をそれぞれ
$x \in X$ と $y \in Y$ と書く。\'etale 近傍
$(V, v) \to (Y, y)$ と開集合 $U \subset X_V$ で、
$\pi : U \to V$ が有限であり、$v$ の上にある唯一の点 $u \in U$ が
射影 $h : U \to X$ により $x \in X$ に写り、さらに
$\kappa(u)/\kappa(v)$ が純非分離となるものを見つけられる。
これは上で用いた補題により可能である。射
$$
U'' = U \times_X U' \longrightarrow V \times_Y V' = V''
$$
を考える。$u$ と $u'$ はともに $x \in X$ に写るので、
$(u, u')$ に写る点 $u'' \in U''$ が存在する。$u''$ の像を
$v'' \in V''$ と書く。$V', v'$ を $V'', v''$ で置き換えると、合成
$V' \to Y' \to Y$ は \'etale 近傍の写像
$(V', v') \to (V, v)$ を経由し、誘導射
$X'_{V'} = X_{V'} \to X_V$ は $u'$ を $u$ に送ると仮定できる。
基底変換 $X'_{V'} = X_{V'}$ の内部には2つの開部分スキーム、すなわち
$U'$ と $U \subset X_V$ の逆像 $U_{V'}$ がある。構成により、両者は
$v'$ の上にある唯一の点を含み、その点はいずれも $u'$ である。
したがって $V'$ を縮小して、これらの開部分集合は同じであると仮定できる。
実際、$U' \setminus (U' \cap U_{V'})$ と
$U_{V'} \setminus (U' \cap U_{V'})$ の $V'$ における像は閉であり、
それらの像は $v'$ を含まない。よって $U' = U_{V'}$ であり、
Lemma \ref{more-morphisms-lemma-weighting-check-after-etale-base-change} の形の Cartesian 図式を得る。
仮定により $\int_\pi (w \circ h)$ は局所定数なので、主張が従う。
\end{proof}

\begin{lemma}
\label{more-morphisms-lemma-weighting-open}
局所準有限射 $f : X \to Y$ をとり、$w : X \to \mathbf{Z}$ を
$f$ の重み付けとする。$X' \subset X$ が開ならば、
$w|_{X'}$ は $f|_{X'} : X' \to Y$ の重み付けである。
\end{lemma}

\begin{proof}
定義から直ちに従う。
\end{proof}

\begin{lemma}
\label{more-morphisms-lemma-weighting-composition}
局所準有限射 $f : X \to Y$ と $g : Y \to Z$ をとる。
$w_f : X \to \mathbf{Z}$ を $f$ の重み付け、
$w_g : Y \to \mathbf{Z}$ を $g$ の重み付けとする。このとき関数
$$
X \longrightarrow \mathbf{Z},\quad
x \longmapsto w_f(x) w_g(f(x))
$$
は $g \circ f$ の重み付けである。
\end{lemma}

\begin{proof}
$x \in X$ に対して $w_{g \circ f}(x) = w_f(x) w_g(f(x))$ とおく。
図式
$$
\xymatrix{
X \ar[d]_{g \circ f} & U \ar[l] \ar[d]^\pi \\
Z & W \ar[l]
}
$$
で、$W \to Z$ は \'etale、$U \subset X_W$ は開、$U \to W$ は有限であるものを考える。
$\int_\pi w_{g \circ f}|_U$ が局所定数であることを示す必要がある。
点 $w \in W$ を選ぶ。Lemma
\ref{more-morphisms-lemma-weighting-check-after-etale-base-change} と \'etale 射が開であることにより、
$(W, w)$ を \'etale 近傍で置き換えた後に
$\int_\pi w_{g \circ f}|_U$ が定数であることを示せば十分である。
$(W, w)$ を \'etale 近傍で置き換えて、$U = U_1 \amalg \ldots \amalg U_n$ であり、
各 $U_i$ は $w$ の上に唯一の点 $u_i$ をもち、
$\kappa(u_i)/\kappa(w)$ が純非分離であると仮定できる。
Lemma \ref{more-morphisms-lemma-etale-splits-off-quasi-finite-part-technical-variant} を参照せよ。
$\int_{U_i \to W} w_{g \circ f}|_{U_i}$ が $w$ のある \'etale 近傍で
定数であることを示せば十分であることは明らかである。
したがって次の段落の場合に帰着する。

\medskip\noindent
$w \in W$ とし、$w$ の上にある唯一の点 $u \in U$ があり、
$\kappa(u)/\kappa(w)$ は純非分離であるとする。点 $v = f(u) \in Y$ を考える。
$(W, w)$ を初等 \'etale 近傍で置き換え、$v$ の開近傍
$V \subset Y_W$ で $V \to W$ が有限となるものが存在すると仮定できる。
Lemma \ref{more-morphisms-lemma-etale-makes-quasi-finite-finite-at-point} を参照せよ。
このとき $f_W^{-1}(V) \cap U$ は $u$ の開近傍である。ここで
$f_W : X_W \to Y_W$ は $f$ の $W$ への基底変換である。
したがって $W$ を Zariski 縮小して、$f_W(U) \subset V$ と仮定できる。
よって射
$$
U \xrightarrow{a} V \xrightarrow{b} W
$$
を得る。$V \to W$ は有限、したがって分離的なので、$U \to V$ は有限である。
$w_f$ と $w_g$ は $f$ と $g$ の重み付けなので、
$\int_a w_f|_U$ は $V$ 上局所定数であり、
$\int_b w_g|_V$ は $W$ 上局所定数である。したがって $W$ をもう一度縮小して、
これらの関数はそれぞれ値 $n$ と $m$ をもつ定数関数であると仮定できる。
直ちに
$\int_\pi w_{g \circ f}|_U = \int_{b \circ a} w_{g \circ f}|_U$
は値 $nm$ をもつ定数関数であり、所望の結果を得る。
\end{proof}

\begin{lemma}
\label{more-morphisms-lemma-weighting-universally-open}
局所準有限射 $f : X \to Y$ と重み付け $w : X \to \mathbf{Z}$ をとる。
すべての $x \in X$ に対して $w(x) > 0$ ならば、$f$ は普遍開である。
\end{lemma}

\begin{proof}
この性質は基底変換で保たれるので (Lemma
\ref{more-morphisms-lemma-weighting-base-change})、$f$ が開であることを証明すれば十分である。
さらに $X$ を $X$ の任意の開部分で置き換えてよいので、$f(X)$ が開であることを
証明すれば十分である。$y \in f(X)$ をとり、$f(x) = y$ となる
$x \in X$ を選ぶ。$f(X)$ が $y$ の開近傍を含むことを示せば十分であり、
そのためには $Y$ を $y$ の \'etale 近傍で置き換えてよい。
準有限射の \'etale 局所化 (Section \ref{more-morphisms-section-etale-localization}) により、
$x$ の開近傍 $U \subset X$ で、$\pi = f|_U : U \to Y$ が有限となるものが
存在すると仮定できる。このとき $\int_\pi w|_U$ は局所定数であり、
$y$ で正の値をもつ。したがって $\pi(U)$ は $y$ の開近傍を含み、証明が完了する。
\end{proof}

\begin{lemma}
\label{more-morphisms-lemma-weighting-flat-quasi-finite}
スキームの射 $f : X \to Y$ をとる。$f$ は局所準有限、有限表示で局所的、
かつ平坦であると仮定する。このとき $f$ には正の重み付け
$w : X \to \mathbf{Z}_{> 0}$ があり、$y \in Y$ の上にある $x \in X$ を
$$
w(x) =
\text{length}_{\mathcal{O}_{X, x}}
(\mathcal{O}_{X, x}/\mathfrak m_y \mathcal{O}_{X, x})
[\kappa(x) : \kappa(y)]_i
$$
に送る規則で与えられる。ここで $[\kappa' : \kappa]_i$ は非分離次数である
(Fields, Definition \ref{fields-definition-insep-degree}).
\end{lemma}

\begin{proof}
Definition \ref{more-morphisms-definition-weighting} に現れる図式を考える。
$u \in U$ の $X, Y, V$ における像を $x, y, v$ とする。次を主張する：
$$
\text{length}_{\mathcal{O}_{X, x}}
(\mathcal{O}_{X, x}/\mathfrak m_y \mathcal{O}_{X, x}) =
\text{length}_{\mathcal{O}_{U, u}}
(\mathcal{O}_{U, u}/\mathfrak m_v \mathcal{O}_{U, u})
$$
および
$$
[\kappa(x) : \kappa(y)]_i =
[\kappa(u) : \kappa(v)]_i
$$
第1の等式が成り立つのは、$\mathcal{O}_{X, x} \to \mathcal{O}_{U, u}$ が
平坦局所準同型であり、
$\mathfrak m_y \mathcal{O}_{U, u} = \mathfrak m_v \mathcal{O}_{U, u}$
かつ $\mathfrak m_x \mathcal{O}_{U, u} = \mathfrak m_u$ を満たすからである
（$\mathcal{O}_{Y, y} \to \mathcal{O}_{V, v}$ と
$\mathcal{O}_{X, x} \to \mathcal{O}_{U, u}$ は不分岐である）。したがって、
Algebra, Lemma \ref{algebra-lemma-pullback-module} により等式を得る。
第2の等式が成り立つのは、$\kappa(v)/\kappa(y)$ が有限分離拡大であり、
$\kappa(u)$ が $\kappa(x) \otimes_{\kappa(y)} \kappa(v)$ の因子なので、
非分離次数が変わらないからである。これにより、補題の関数の構成は
\'etale 基底変換と可換である。したがって問題は次の段落の議論に帰着する。

\medskip\noindent
$f$ は有限、平坦、かつ有限表示な射であると仮定する。
$\int_f w$ が $Y$ 上局所定数であることを示す必要がある。
実際、$f$ は有限局所自由であり
(Morphisms, Lemma \ref{morphisms-lemma-finite-flat})、
$\int_f w$ は $f$ の次数に等しいことを示す
（これは $Y$ 上の局所定数関数である）。すなわち $y \in Y$ に対して
\begin{align*}
(\textstyle{\int}_f w)(y)
& =
\sum\nolimits_{f(x) = y}
\text{length}_{\mathcal{O}_{X, x}}
(\mathcal{O}_{X, x}/\mathfrak m_y \mathcal{O}_{X, x})
[\kappa(x) : \kappa(y)]_i
[\kappa(x) : \kappa(y)]_s \\
& =
\sum\nolimits_{f(x) = y}
\text{length}_{\mathcal{O}_{X, x}}
(\mathcal{O}_{X, x}/\mathfrak m_y \mathcal{O}_{X, x})
[\kappa(x) : \kappa(y)] \\
& =
\text{length}_{\mathcal{O}_{Y, y}}((f_*\mathcal{O}_X)_y/
\mathfrak m_y (f_*\mathcal{O}_X)_y)
\end{align*}
% P05-EN-CH37-0356
最後の等式は Algebra, Lemma \ref{algebra-lemma-pushdown-module} による。
最後の数は $y$ における $f_*\mathcal{O}_X$ の階数であり、所望の結果を得る。
\end{proof}

\begin{lemma}
\label{more-morphisms-lemma-weighting-quasi-finite-Noetherian}
スキームの射 $f : X \to Y$ が次を満たすと仮定する：
\begin{enumerate}
\item $f$ は局所準有限である、
\item $Y$ は幾何学的単枝かつ局所 Noether である。
\end{enumerate}
このとき重み付け $w : X \to \mathbf{Z}_{\geq 0}$ が存在し、
$y \in Y$ の上にある $x \in X$ を、
$\mathcal{O}_{Y, y}^{sh}$ 上の $\mathcal{O}_{X, x}^{sh}$ の
「生成点での分離次数」に送る規則で与えられる。
\end{lemma}

\begin{proof}
Algebra, Lemma \ref{algebra-lemma-quasi-finite-strict-henselization} により、
$\mathcal{O}_{Y, y}^{sh} \to \mathcal{O}_{X, x}^{sh}$ は有限である。
$Y$ は幾何学的単枝なので、$\mathcal{O}_{Y, y}^{sh}$ には唯一の極小素イデアル
$\mathfrak p$ が存在する。More on Algebra, Lemma
\ref{more-algebra-lemma-geometrically-unibranch} を参照せよ。
$$
(\kappa(\mathfrak p) \otimes_{\mathcal{O}_{Y, y}^{sh}}
\mathcal{O}_{X, x}^{sh})_{red} =
\prod K_i
$$
を体の有限積として書く。$w(x) = \sum [K_i : \kappa(\mathfrak p)]_s$ とおく。

\medskip\noindent
この定義が \'etale 局所化に依存しないことは明らかなので、$w$ が重み付けであることを
示すには、$f$ が有限射ならば $\int_f w$ が局所定数であることを示す場合に帰着する。
$Y$ の生成点 $\eta$ における $\int_f w$ の値は、$\eta$ 上の
$X \to Y$ の幾何学的ファイバー $X_{\overline{\eta}}$ の点の個数にほかならない。
さらに $Y$ は単枝なので、$Y$ の点 $y$ は唯一の生成点 $\eta$ の特殊化である。
したがって $(\int_f w)(y)$ が $X_{\overline{\eta}}$ の点の個数に等しいことを
示せば十分である。$y$ のアフィン近傍に移ることで、$X \to Y$ は有限環準同型
$A \to B$ により与えられると仮定できる。
$\mathcal{O}_{Y, y}^{sh}$ は、代数閉体 $k$ への写像
$\kappa(y) \to k$ を用いて構成されているとする。このとき
$$
\mathcal{O}_{Y, y}^{sh} \otimes_A B =
\prod\nolimits_{f(x) = y}
\prod\nolimits_{\varphi \in \Mor_{\kappa(y)}(\kappa(x), k)}
\mathcal{O}_{X, x}^{sh}
$$
である。これは Algebra, Lemma \ref{algebra-lemma-finite-over-henselian} と
上で用いた補題による。$\mathcal{O}_{Y, y}^{sh}$ の極小素イデアル
$\mathfrak p$ は、$\eta$ に対応する $A$ の素イデアルに写る。
したがって、幾何学的ファイバーの点の個数は体拡大で変わらないことから、
所望の等式が成り立つ。
\end{proof}

\begin{lemma}
\label{more-morphisms-lemma-weighting-specialization}
スキームの局所準有限射 $f : X \to Y$ をとり、$w : X \to \mathbf{Z}$ を
$f$ の重み付けとする。$Y$ の点の特殊化 $y \leadsto y'$ をとる。
$f(x) = y$ を満たすすべての $x \in X$ に対して $w(x) = 0$ と仮定する。
このとき、$f(x') = y'$ を満たすすべての $x' \in X$ に対して
$w(x') = 0$ である。
\end{lemma}

\begin{proof}
$y'$ に写る点 $x' \in X$ をとる。Lemma
\ref{more-morphisms-lemma-etale-makes-quasi-finite-finite-at-point} により、図式
$$
\xymatrix{
U \ar[r] \ar[dr]_\pi & X_V \ar[d] \ar[r] & X \ar[d]^f \\
& V \ar[r]^\varphi & Y
}
$$
を選べる。ここで $V \to Y$ は \'etale、点 $v' \in V$ は $y'$ に写り、
% P05-EN-CH37-0354
$U \subset Y_V$ は開で、$U \to V$ は有限であり、$v'$ に写る唯一の点
% P05-EN-CH37-0355
$u' \in U$ が存在し、さらに $X$ で $x'$ に写る。
$V \to Y$ は開なので、$v'$ に特殊化し $y$ に写る点 $v \in V$ が存在する。
Definition \ref{more-morphisms-definition-weighting} により、関数
$\int_\pi (w \circ h)$ は局所定数である。ここで $h : U \to X$ は
上の水平矢印の合成である。$w$ に関する仮定により
$\int_\pi (w \circ h)$ は $v$ で零なので、$v'$ でも零である。
しかし $v'$ における値は $w(x')[\kappa(u'):\kappa(v')]_s$ であるから、主張が従う。
\end{proof}




\section{重み付けについてさらに}
\label{more-morphisms-section-more-weightings}

\noindent
重み付けの基本的性質をさらにいくつか証明する。重み付けは一見すると
非常に不規則になり得るように思われるが、実際には
Definition \ref{more-morphisms-definition-weighting} で課した条件はかなり強い。

\begin{lemma}
\label{more-morphisms-lemma-jumps-w}
$f : X \to Y$ を局所準有限射とし、
$w : X \to \mathbf{Z}$ を $f$ の重み付けとする。このとき
関数 $w$ のレベル集合は $X$ において局所構成可能である。
\end{lemma}

\begin{proof}
以下の証明では Lemma \ref{more-morphisms-lemma-weighting-open} と
\ref{more-morphisms-lemma-weighting-base-change} を断りなく用いる。
またスキームおよび位相空間の構成可能部分集合の初等的性質も用いる。
Topology, Section \ref{topology-section-constructible} および
Properties, Section \ref{properties-section-constructible} を参照せよ。
% P05-EN-CH37-0357
これらから問題が $X$ と $Y$ 上で局所的であることが分かるので、詳細は省略する。
したがって $X$ と $Y$ はアフィンであると仮定してよい。
有限表示な全射 $Y' \to Y$ で、$w$ のレベル集合が
$X' = Y' \times_Y X$ の局所構成可能部分集合へ引き戻されるものを見つければ、
Morphisms, Theorem \ref{morphisms-theorem-chevalley} により結論を得る。

\medskip\noindent
$X$ と $Y$ がアフィンであると仮定する。
$T \to Y$ が有限であるような埋め込み $X \to T$ を選べる。
Lemma \ref{more-morphisms-lemma-quasi-finite-separated-pass-through-finite} を参照せよ。
Morphisms, Lemma \ref{morphisms-lemma-massage-finite} により、
$Y$ 上全射かつ有限局所自由な $Y'$ で $Y$ を置き換え、
$X$ を $Y' \times_Y X$ で置き換え、さらに $T$ を
$Y' \times_Y T$ を閉部分スキームとして含む $Y'$ 上有限局所自由なスキームで
置き換えると、$T$ は $Y$ 上有限局所自由であり、$Y$ へ同型に写る
閉部分スキーム $T_i$ を含み、集合論的に
$T = \bigcup_{i = 1, \ldots, n} T_i$ であると仮定してよい。
$T_i \subset T$ は構成可能な閉部分集合である
（有限表示なスキーム射 $Y \to T$ の像だからである）。したがって
$I \subset \{1, \ldots, n\}$ に対し、交わり
$\bigcap_{i \in I} T_i$ は $T$ の構成可能な閉部分集合であり、
ゆえに $Y$ の構成可能な閉部分集合へ写る。

\medskip\noindent
空でない部分からなる非交和分解
$\{1, \ldots, n\} = I_1 \amalg \ldots \amalg I_r$ に対し、
$T_y = \{x_1, \ldots, x_r\}$ がちょうど $r$ 個の点からなり、
$x_j \in T_i \Leftrightarrow i \in I_j$ を満たすような点 $y \in Y$ 全体を
$Y_{I_1, \ldots, I_r} \subset Y$ とする。
上の考察により、これらは $Y$ の構成可能な分割をなす。
$Y' \to Y$ の像がちょうど $Y_{I_1, \ldots, I_r}$ となる
$Y$ 上有限表示なアフィンスキーム $Y'$ が存在する。
Algebra, Lemma \ref{algebra-lemma-constructible-is-image} を参照せよ。
したがって、ある非交和分解
$\{1, \ldots, n\} = I_1 \amalg \ldots \amalg I_r$ に対して
$Y = Y_{I_1, \ldots, I_r}$ であると仮定してよい。
この場合、$T(j) = \bigcap_{i \in I_j} T_i$ とおけば
$T = T(1) \amalg \ldots \amalg T(r)$ は $T$ の互いに素な閉部分集合
（したがって開部分集合）への分解である。
局所閉部分スキーム $X$ と交わらせると、開閉部分への同様の分解
$X = X(1) \amalg \ldots \amalg X(r)$ を得る。
% P05-EN-CH37-0358
射 $X(j) \to Y$ は埋め込みである。
$w$ は重み付けなので、$w|_{X(j)}$ は局所定数である
\footnote{実際、$f : X \to Y$ が埋め込みで $w$ が $f$ の重み付けならば、
$w$ が非零である軌跡への $f$ の制限は開写像である。}。これで結論を得る。
\end{proof}

\begin{lemma}
\label{more-morphisms-lemma-weighting-on-dense-set}
$f : X \to Y$ をスキームの局所準有限射とし、
$w : X \to \mathbf{Z}$ を $f$ の重み付けとする。
$E \subset X$ を稠密部分集合とする。
すべての $x \in E$ に対して $w(x) = 0$ と仮定する。このとき $w = 0$ である。
\end{lemma}

\begin{proof}
先に Lemma \ref{more-morphisms-lemma-weighting-specialization} を読むことを勧める。
$x \in X$ の像を $y \in Y$ とする。$w(x) = 0$ を示そう。
Lemma \ref{more-morphisms-lemma-etale-makes-quasi-finite-finite-at-point} により図式
$$
\xymatrix{
U \ar[r] \ar[dr]_\pi & X_V \ar[d] \ar[r] & X \ar[d]^f \\
& V \ar[r]^\varphi & Y
}
$$
を選べる。ここで $V \to Y$ は \'etale、点 $v \in V$ は $y$ へ写り、
% P05-EN-CH37-0359
$U \subset Y_V$ は開で $U \to V$ は有限であり、さらに
% P05-EN-CH37-0360
$v$ へ写り、かつ $X$ において $x$ へ写る一意な点 $u \in U$ が存在する。
$V$（したがって $U$）はアフィンであると仮定してよく、そう仮定する。
Definition \ref{more-morphisms-definition-weighting} により、関数
$\int_\pi (w \circ h)$ は局所定数である。ここで $h : U \to X$ は
上の水平矢印の合成である。
$w \circ h$ が $0$ となる点全体 $U_0$ は Lemma \ref{more-morphisms-lemma-jumps-w} により
$U$ の構成可能部分集合であり、$h^{-1}(E)$ を含むので $U$ で稠密である
（$h$ が開であることを用いよ）。
Topology, Lemma \ref{topology-lemma-dense-in-constructible} により、
$U_0$ は稠密開集合を含む。言い換えると、$w \circ h$ が非零である点全体は
ある疎な閉部分集合 $Z \subset U$ に含まれる。
すると Morphisms, Lemma \ref{morphisms-lemma-image-nowhere-dense-finite} により
$\pi(Z)$ は $V$ の疎な閉部分集合である。
したがって $\int_\pi (w \circ h)$ は稠密開集合上で零であり、ゆえに至る所で消える。
一方、$v$ における値は $w(x)[\kappa(u):\kappa(v)]_s$ であるから、結論を得る。
\end{proof}

\begin{lemma}
\label{more-morphisms-lemma-jumps-int-w}
$f : X \to Y$ を有限表示な局所準有限射とし、
$w : X \to \mathbf{Z}$ を $f$ の重み付けとする。このとき
関数 $\int_f w$ のレベル集合は $Y$ において局所構成可能である。
\end{lemma}

\begin{proof}
Lemma \ref{more-morphisms-lemma-weighting-check-after-etale-base-change} により
関数 $\int_f w$ の形成は任意の基底変換と可換であり、
% P05-EN-CH37-0361
Lemma \ref{more-morphisms-lemma-weighting-base-change} により基底変換後も重み付けを得る。
したがって有限表示な全射 $Y' \to Y$ を見つければ、
$Y'$ への基底変換後に結果を証明すれば十分である。
Morphisms, Theorem \ref{morphisms-theorem-chevalley} を参照せよ。

\medskip\noindent
問題は $Y$ 上局所的なので、$Y$ はアフィンであると仮定してよい。
このとき $f$ は有限表示だから、$X$ は準コンパクトかつ準分離である。
準コンパクト開集合からなる被覆 $X = U \cup V$ をとる。このとき
$$
\textstyle{\int}_f w =
\textstyle{\int}_{f|_U} w|_U +
\textstyle{\int}_{f|_V} w|_V -
\textstyle{\int}_{f|_{U \cap V}} w|_{U \cap V}
$$
したがって $w|_U$, $w|_V$, $w|_{U \cap V}$ について結果が分かれば、
$w$ についても分かる。帰納原理
（Cohomology of Schemes, Lemma \ref{coherent-lemma-induction-principle}）により、
$X$ がアフィンの場合に補題を証明すれば十分である。

\medskip\noindent
$X$ と $Y$ がアフィンであると仮定する。
$T \to Y$ が有限であるような開埋め込み $X \to T$ を選べる。
Lemma \ref{more-morphisms-lemma-quasi-finite-separated-pass-through-finite} を参照せよ。
適当な $Y' \to Y$ でなお基底変換できるので、
Morphisms, Lemma \ref{morphisms-lemma-massage-finite} を用いて、
射 $T \to Y$ が誘導するすべての剰余体拡大
% P05-EN-CH37-0362
（したがってなおさら $X \to Y$ が誘導するもの）が自明である場合に帰着できる。
この状況では $\int_f w$ は単に各ファイバーにおける $w$ の値の和をとる。
$w$ のレベル集合は $X$ で局所構成可能である
（Lemma \ref{more-morphisms-lemma-jumps-w}）。Properties, Lemma
\ref{properties-lemma-stratification-locally-finite-constructible} により、
関数 $w$ は有限個の値しかとらない。
したがって Morphisms, Theorem \ref{morphisms-theorem-chevalley} と
整数の和に関する初等的議論により結論を得る。
\end{proof}

\begin{lemma}
\label{more-morphisms-lemma-semicontinuous-w}
$f : X \to Y$ を局所準有限射とし、
$w : X \to \mathbf{Z}_{> 0}$ を $f$ の正の重み付けとする。
このとき $w$ は上半連続である。
\end{lemma}

\begin{proof}
$x \in X$ の像を $y \in Y$ とする。\'etale 近傍
$(V, v) \to (Y, y)$ と開集合 $U \subset X_V$ を、
$\pi : U \to V$ が有限で、$v$ へ写る一意な点 $u \in U$ が存在し、
$\kappa(u)/\kappa(v)$ が純非分離となるように選ぶ。
Lemma \ref{more-morphisms-lemma-etale-makes-quasi-finite-finite-multiple-points-var} を参照せよ。
このとき $(\int_\pi w|_U)(v) = w(u)$ である。
Definition \ref{more-morphisms-definition-weighting} により、$V$ を $v$ の近傍で置き換えた後、
% P05-EN-CH37-0363
すべての $u' \in U$ に対して $w|_U(u') \leq w|_U(u) = w(x)$ が成り立つ。
実際、$w|_U(u')$ は $(\int_\pi w|_U)(\pi(u'))$ の式の一つの項として現れる。
\'etale 射 $U \to X$ は開であるから、これで補題が証明された。
\end{proof}

\begin{lemma}
\label{more-morphisms-lemma-semicontinuous-int-w}
$f : X \to Y$ を有限ファイバーをもつ分離な局所準有限射とし、
$w : X \to \mathbf{Z}_{> 0}$ を $f$ の正の重み付けとする。
このとき $\int_f w$ は下半連続である。
\end{lemma}

\begin{proof}
$y \in Y$ とし、$x_1, \ldots, x_r \in X$ を $y$ の上にある点とする。
Lemma \ref{more-morphisms-lemma-etale-splits-off-quasi-finite-part-technical-variant} を適用して、
\'etale 近傍 $(U, u) \to (Y, y)$ と分解
$$
U \times_Y X =
W \amalg
\ \coprod\nolimits_{i = 1, \ldots, n}
\ \coprod\nolimits_{j = 1, \ldots, m_i}
V_{i, j}
$$
を前掲箇所のように得る。ここで $(\int_f w)(y) = \sum w(v_{i, j})$ であり、
$w(v_{i, j}) = w(x_i)$ であることに注意せよ。
定義により $\int_{V_{i, j} \to U} w|_{V_{i, j}}$ は局所定数なので、
$U$ を縮小した後、これらの関数は値 $w(v_{i, j})$ をもつ定数関数であると仮定してよい。
したがって
$$
\textstyle{\int}_{U \times_Y X \to U} w|_{U \times_Y X} =
\textstyle{\int}_{W \to U} w|_W +
\sum \textstyle{\int}_{V_{i, j} \to U} w|_{V_{i, j}} =
\textstyle{\int}_{W \to U} w|_W + (\int_f w)(y)
$$
これは $\geq (\int_f w)(y)$ である。$U \to Y$ は開であり、積分の形成は
基底変換と可換であるから結論を得る
（Lemma \ref{more-morphisms-lemma-weighting-check-after-etale-base-change}）。
\end{proof}

\begin{lemma}
\label{more-morphisms-lemma-max-int-w}
$f : X \to Y$ を $X$ が準コンパクトである局所準有限射とし、
$w : X \to \mathbf{Z}$ を $f$ の重み付けとする。
このとき $\int_f w$ は最大値をとる。
\end{lemma}

\begin{proof}
Lemma \ref{more-morphisms-lemma-jumps-w} と Properties, Lemma
\ref{properties-lemma-stratification-locally-finite-constructible} により、
$w$ は $X$ 上で有限個の値しかとらない。
Morphisms, Lemma
\ref{morphisms-lemma-locally-quasi-finite-qc-source-universally-bounded} により、
$X \to Y$ の幾何学的ファイバーの大きさは有界である。
したがって $\int_f w$ は有界である。
\end{proof}

\begin{lemma}
\label{more-morphisms-lemma-max-int-finite}
$f : X \to Y$ を分離な局所準有限射とし、
$w : X \to \mathbf{Z}_{> 0}$ を $f$ の正の重み付けとする。
% P05-EN-CH37-0364
$\int_w f$ が最大値 $d$ をとると仮定し、$(\int_f w)(y) = d$ を満たす点
$y$ 全体の開集合を $Y_d \subset Y$ とする。このとき射
$f^{-1}(Y_d) \to Y_d$ は有限である。
\end{lemma}

\begin{proof}
Lemma \ref{more-morphisms-lemma-semicontinuous-int-w} により $Y_d$ は開である。
$y \in Y_d$ とし、$x_1, \ldots, x_n$ を $y$ の上にある $X$ の点とする。
Lemma \ref{more-morphisms-lemma-etale-splits-off-quasi-finite-part-technical-variant} を適用して、
\'etale 近傍 $(U, u) \to (Y, y)$ と分解
$$
U \times_Y X =
W \amalg
\ \coprod\nolimits_{i = 1, \ldots, n}
\ \coprod\nolimits_{j = 1, \ldots, m_i}
V_{i, j}
$$
を前掲箇所のように得る。ここで $d = \sum w(v_{i, j})$ であり、
$w(v_{i, j}) = w(x_i)$ であることに注意せよ。
定義により $\int_{V_{i, j} \to U} w|_{V_{i, j}}$ は局所定数なので、
$U$ を縮小した後、これらの関数は値 $w(v_{i, j})$ をもつ定数関数であると仮定してよい。
したがって
$$
\textstyle{\int}_{U \times_Y X \to U} w|_{U \times_Y X} =
\textstyle{\int}_{W \to U} w|_W +
\sum \textstyle{\int}_{V_{i, j} \to U} w|_{V_{i, j}} =
\textstyle{\int}_{W \to U} w|_W + (\int_f w)(y)
$$
これは $\geq (\int_f w)(y) = d$ であるから、
% P05-EN-CH37-0365
$W$ は空でなければならない。したがって $U \times_Y X \to U$ は有限である。
Descent, Lemma \ref{descent-lemma-descending-property-finite} により、
$X \to Y$ は開射 $U \to Y$ の像の上で有限である。
言い換えると、望みどおり $f$ は $y$ のある開近傍上で有限である。
\end{proof}

\begin{lemma}
\label{more-morphisms-lemma-open-and-closed-in-finite}
$A \to B$ を有限かつ有限表示な環準同型とする。
有限表示な環準同型 $A \to A_{univ}$ と冪等元
$e_{univ} \in B \otimes_A A_{univ}$ で、任意の環準同型 $A \to A'$ と
冪等元 $e \in B \otimes_A A'$ に対し、$e_{univ}$ を $e$ へ写す環準同型
$A_{univ} \to A'$ が存在するようなものがある。
\end{lemma}

\begin{proof}
$B$ を $A$ 加群として生成する $b_1, \ldots, b_n \in B$ を選ぶ。
各 $i$ に対し、$B$ において $P_i(b_i) = 0$ となるモニック多項式
$P_i \in A[x]$ を選ぶ。Algebra, Lemma \ref{algebra-lemma-finite-is-integral} を参照せよ。
したがって $B$ は有限自由 $A$ 代数
$B' = A[x_1, \ldots, x_n]/(P_1(x_1), \ldots, P_n(x_n))$ の商である。
全射 $B' \to B$ の核を $J \subset B'$ とする。
% P05-EN-CH37-0366
$B$ は有限生成 $A$ 代数なので $J =(f_1, \ldots, f_m)$ は有限生成である。
Algebra, Lemma \ref{algebra-lemma-compose-finite-type} を参照せよ。
$B'$ の $A$ 基底 $b'_1, \ldots, b'_N$ を選び、代数
$$
A_{univ} = A[z_1, \ldots, z_N, y_1, \ldots, y_m]/I
$$
を考える。ここで $I$ は、式
$$
(\sum z_j b'_j)^2 - \sum z_j b'_j + \sum y_k f_k
$$
の $B'[z_1, \ldots, z_N, y_1, \ldots, y_m]$ における
基底要素 $b'_1, \ldots, b'_N$ の係数を、
% P05-EN-CH37-0367
$A[z_1, \ldots, z_n, y_1, \ldots, y_m]$ の元として生成されるイデアルである。
構成により、元 $\sum z_j b'_j$ は代数 $B \otimes_A A_{univ}$ の
冪等元 $e_{univ}$ へ写る。さらに $e \in B \otimes_A A'$ が冪等元ならば、
$e$ を $B' \otimes_A A'$ の形 $\sum b'_j \otimes a'_j$ の元へ持ち上げられ、
ある $a''_k \in A'$ を選んで
$$
(\sum b'_j \otimes a'_j)^2 - \sum b'_j \otimes a'_j + \sum f_k \otimes a''_k
$$
を $B' \otimes_A A'$ において零にできる。したがって $z_j$ を $a'_j$ へ、
$y_k$ を $a''_k$ へ写し、$e_{univ}$ を $e$ へ写す $A$ 代数準同型
% P05-EN-CH37-0368
$A_{univ} \to A$ を得る。これで証明は終わる。
\end{proof}

\begin{lemma}
\label{more-morphisms-lemma-open-and-closed-in-quasi-finite}
$X \to Y$ を準有限かつ有限表示なアフィンスキームの射とする。
有限表示な射 $Y_{univ} \to Y$ と開部分スキーム
$U_{univ} \subset Y_{univ} \times_Y X$ で、$U_{univ} \to Y_{univ}$ が有限であり、
次の性質をもつものが存在する：アフィンスキームの任意の射 $Y' \to Y$ と、
$U' \to Y'$ が有限であるような開部分スキーム
$U' \subset Y' \times_Y X$ が与えられると、$U_{univ}$ の逆像が $U'$ となる射
$Y' \to Y_{univ}$ が存在する。
\end{lemma}

\begin{proof}
有限型射が準有限であることと相対次元 $0$ であることは同値である。
Morphisms, Lemma \ref{morphisms-lemma-locally-quasi-finite-rel-dimension-0} を参照せよ。
Lemma \ref{more-morphisms-lemma-Noetherian-approximation-dimension-d} を $d = 0$ として適用し、
$X$ と $Y$ が Noetherian である場合に帰着する。
$X' \to Y$ が有限となる開埋め込み $X \to X'$ を選べる。
Algebra, Lemma \ref{algebra-lemma-quasi-finite-open-integral-closure} を参照せよ。
% P05-EN-CH37-0369
$Y' \to Y$ と（2）のような $U'$ があるとき、
$$
U' \to Y' \times_Y X \to Y' \times_Y X'
$$
は $Y'$ 上有限なスキーム間の開埋め込みであり、したがって閉でもある。
ゆえに $U'$ は
$$
\Gamma(Y', \mathcal{O}_{Y'})
\otimes_{\Gamma(Y, \mathcal{O}_Y)}
\Gamma(X', \mathcal{O}_{X'})
$$
の冪等元で、対応する開閉部分集合が開集合 $Y' \times_Y X$ に含まれるものに対応する。
$Y'_{univ} \to Y$ と冪等元
$$
e'_{univ} \in
% P05-EN-CH37-0370
\Gamma(Y_{univ}, \mathcal{O}_{Y_{univ}})
\otimes_{\Gamma(Y, \mathcal{O}_Y)}
\Gamma(X', \mathcal{O}_{X'})
$$
を、環準同型 $\Gamma(Y, \mathcal{O}_Y) \to \Gamma(X', \mathcal{O}_{X'})$ に対し
Lemma \ref{more-morphisms-lemma-open-and-closed-in-finite} で構成した組とする
（ここでは $Y$ が Noetherian であることから $X'$ が $Y$ 上有限表示であることを用いる）。
対応する開閉部分スキームを $U'_{univ} \subset Y'_{univ} \times_Y X'$ とする。このとき
$$
U'_{univ} \setminus Y'_{univ} \times_Y X
$$
は $U'_{univ}$ の閉部分集合であり、したがって閉な像
$T \subset Y'_{univ}$ をもつ。$Y_{univ} = Y'_{univ} \setminus T$ とおき、
$U_{univ}$ を $U'_{univ}$ の $Y_{univ} \times_Y X$ への制限とすれば、
補題が成り立つことが分かる。
\end{proof}

\begin{lemma}
\label{more-morphisms-lemma-descend-weighting}
$Y = \lim Y_i$ をアフィンスキームの有向極限とする。$0 \in I$ とし、
$f_0 : X_0 \to Y_0$ を準有限かつ有限表示なアフィンスキームの射とする。
$f : X  \to Y$ および $i \geq 0$ に対する $f_i : X_i \to Y_i$ を
$f_0$ の基底変換とする。$w : X \to \mathbf{Z}$ が $f$ の重み付けならば、
十分大きい $i$ に対し、$X$ への引き戻しが $w$ となる $f_i$ の重み付け
$w_i : X_i \to \mathbf{Z}$ が存在する。
\end{lemma}

\begin{proof}
Lemma \ref{more-morphisms-lemma-jumps-w} により、$w$ のレベル集合は $X$ の構成可能部分集合
$E_k$ である。Properties, Lemma
\ref{properties-lemma-stratification-locally-finite-constructible} により、
関数 $w$ は有限個の値しかとらない。したがって、すべての $k$ に対し
$E_k$ が $X_i$ の構成可能部分集合 $E_{i, k}$ に降下するような $i$ が存在する。
% P05-EN-CH37-0371
さらに $X_i = \coprod E_{i, k}$ と仮定してよい。
これは、$X$ の位相空間が、スペクトル遷移写像をもつ有向系に沿った
スペクトル空間 $X_i$ の、位相空間の圏における極限であることから従う。
Limits, Section \ref{limits-section-descent} および
Topology, Section \ref{topology-section-limits-spectral} を参照せよ。
レベル集合が構成可能集合 $E_{i, k}$ となるように
$w_i : X_i \to \mathbf{Z}$ を定める。

\medskip\noindent
Lemma \ref{more-morphisms-lemma-open-and-closed-in-quasi-finite} のように
$Y_{i, univ} \to Y_i$ と
$U_{i, univ} \subset Y_{i, univ} \times_{Y_i} X_i$ を選ぶ。
この構成の普遍性により、$w_i$ が重み付けであることを示すには、
$$
\tau_i = \textstyle{\int}_{U_{i, univ} \to Y_{i, univ}} w_i|_{U_{i, univ}}
$$
が $Y_{i, univ}$ 上局所定数であることを示せば十分である。
Lemma \ref{more-morphisms-lemma-jumps-int-w} により、この関数のレベル集合は構成可能だが、
（まだ）局所定数とは限らない。
$Y_{univ} = Y_{i, univ} \times_{Y_i} Y$ とおき、
$U_{univ} \subset Y_{univ} \times_Y X$ を $U_{i, univ}$ の逆像とする。
すると $w$ の $Y_{univ} \times_Y X$ への引き戻しは
$Y_{univ} \times_Y X \to Y_{univ}$ の重み付けなので
（Lemma \ref{more-morphisms-lemma-weighting-base-change}）、今度は
$$
\tau = \textstyle{\int}_{U_{univ} \to Y_{univ}} w_i|_{U_{univ}}
$$
% P05-EN-CH37-0372
は $Y_{univ}$ 上局所定数である。したがって $\tau$ のレベル集合は開閉である。
最後に
$Y_{univ} = \lim_{i' \geq i} Y_{i', univ}$
であり、$\tau$ のレベル集合は、同様に定義される $\tau_{i'}$ の
レベル集合の逆極限である。ゆえに上の参照結果から、十分大きい $i'$ に対して
$\tau_{i'}$ のレベル集合も開である。そのような添字 $i'$ に対し、
$w_{i'}$ は望みどおり $f_{i'}$ の重み付けである。
\end{proof}







\section{重み付けとアフィン層別数}
\label{more-morphisms-section-bounds-asn}

\noindent
本節では、アフィンスキームによるある種の被覆をもつスキームの
アフィン層別数に対する上界を与える。

\begin{lemma}
\label{more-morphisms-lemma-affineness-of-large-open}
$f : X \to Y$ を準有限かつ有限表示なアフィンスキームの射とする。
$w : X \to \mathbf{Z}_{> 0}$ を $f$ の正の重み付けとする。
$d < \infty$ を $\int_f w$ の最大値とする。$Y$ の開集合
$$
Y_d = \{y \in Y \mid (\textstyle{\int}_f w)(y) = d \}
$$
はアフィンである。
\end{lemma}

\begin{proof}
Lemma \ref{more-morphisms-lemma-max-int-w} により $\int_f w$ は最大値をとる。
また Lemma \ref{more-morphisms-lemma-semicontinuous-int-w} により $Y_d$ は開である。
したがって補題の主張は意味をもつ。

\medskip\noindent
Noetherian の場合への帰着。この段落は読み飛ばしてよい。
有限型射が準有限であることと相対次元 $0$ であることは同値である。
Morphisms, Lemma \ref{morphisms-lemma-locally-quasi-finite-rel-dimension-0} を参照せよ。
Lemma \ref{more-morphisms-lemma-Noetherian-approximation-dimension-d} を $d = 0$ として適用すると、
アフィン Noetherian スキームの準有限射 $f_0 : X_0 \to Y_0$ と射 $Y \to Y_0$ で、
$f$ が $f_0$ の基底変換となるものを見つけられる。
次に $Y = \lim Y_i$ を $Y_0$ 上有限型なアフィンスキームの有向極限として書ける。
Algebra, Lemma \ref{algebra-lemma-ring-colimit-fp} を参照せよ。
Lemma \ref{more-morphisms-lemma-descend-weighting} により、重み付け $w$ が $f_0$ の基底変換
$f_i : X_i \to Y_i$ の重み付け $w_i$ に降下するような $i$ を見つけられる。
$f_i,w_i$ に対して補題が成り立てば、$\int_f w$ の形成は基底変換と可換なので、
$f$ に対する補題も従う。Lemma \ref{more-morphisms-lemma-weighting-check-after-etale-base-change} を参照せよ。

\medskip\noindent
$X$ と $Y$ が Noetherian であると仮定する。射 $g : Y' \to Y$ による
$f$ の基底変換を $X' \to Y'$ とする。$\int_f w$ の形成、したがって開集合
$Y_d$ は基底変換と可換である。$g$ が有限かつ全射ならば、
$Y'_d \to Y_d$ も有限かつ全射である。この場合、$Y_d$ がアフィンであることと
$Y'_d$ がアフィンであることは同値である。Cohomology of Schemes, Lemma
\ref{coherent-lemma-image-affine-finite-morphism-affine-Noetherian} を参照せよ。

\medskip\noindent
$T$ が $Y$ 上有限であるような埋め込み $X \to T$ を選べる。
Lemma \ref{more-morphisms-lemma-quasi-finite-separated-pass-through-finite} を参照せよ。
有限射 $T \to Y$ に Morphisms, Lemma \ref{morphisms-lemma-massage-finite} を
適用する。この補題により、$Y' \times_Y T$ が、特別な形をもち $Y'$ 上有限な
スキーム $T'$ の閉部分スキームとなるような有限全射 $Y' \to Y$ が存在する。
最初の段落の議論により、$Y$ を $Y'$ で、$T$ を $T'$ で、$X$ を
$Y' \times_Y X$ で置き換えてよい。したがって、埋め込み $X \to T$
（開または閉とは限らない）と閉部分スキーム
$T_i \subset T$, $i = 1, \ldots, n$ が存在し、次を満たすと仮定してよい：
\begin{enumerate}
\item $T \to Y$ は有限（かつ局所自由）である。
\item $T_i \to Y$ は同型である。
\item 集合論的に $T = \bigcup_{i = 1, \ldots, n} T_i$ である。
\end{enumerate}
$Y$ の既約成分を $Y$ の整閉部分スキームとみなし、その非交和を
$Y' = \coprod Y_k$ とする。$Y$ が Noetherian であることを用いれば、
$Y' \to Y$ によりもう一度基底変換してよい。したがってさらに
$Y$ は整かつ Noetherian であると仮定してよい。

\medskip\noindent
重複を除くことにより、$i \not = j$ ならば $T_i \not = T_j$ であると
仮定してよく、そう仮定する。$Y$、したがってすべての $T_i$ は整なので、
$T_i$ と $T_j$ が交わるならば、その交わりは $Y$ の真の閉部分集合へ写る
閉部分集合である。

\medskip\noindent
$V_i = X \cap T_i$ は局所閉部分集合であり、さらに $X$ の閉部分スキームなので
アフィンである。$\eta \in Y$ と $\eta_i \in T_i$ をそれぞれ総称点とする。
$\eta \not \in Y_d$ ならば $Y_d = \emptyset$ であり、証明は終わる。
$\eta \in Y_d$ と仮定する。
% P05-EN-CH37-0373
$\eta_i \in V_i$ となる添字 $i$ 全体を $I \in \{1, \ldots, n\}$ と書く。
$i \in I$ に対し、局所閉部分集合 $V_i \subset T_i$ は既約空間 $T_i$ の
総称点を含むので開である。一方、$f$ は開であるから
（Lemma \ref{more-morphisms-lemma-weighting-universally-open}）、任意の $x \in X$ に対し、
ある $i \in I$ と特殊化 $\eta_i \leadsto x$ を見つけられる。
したがって $x \in T_i$、ゆえに $x \in V_i$ である。言い換えると、集合論的に
$X = \bigcup_{i \in I} V_i$ である。
$Y_d = \bigcap_{i \in I} \Im(V_i \to Y)$ と主張する。これが示されれば、
$Y$ のアフィン開集合 $\Im(V_i \to Y)$ の交わりはアフィンなので証明が終わる。

\medskip\noindent
$y \in Y$ に対し、$X$ において $f^{-1}(\{y\}) = \{x_1, \ldots, x_r\}$ とする。
各 $i \in I$ に対し、$\eta_i \leadsto x_{j(i)}$ となる
% P05-EN-CH37-0374
$j(i) \in \{1, \ldots, x_r\}$ は高々一つである。実際、$j(i)$ が存在して
$j$ に等しいことと $x_j \in V_i$ であることは同値である。
$i \in I$ に対して $j = j(i)$ が存在するならば、$V_i \to Y$ は
$x_j \mapsto y$ の近傍で同型である。したがって、始域と終域を
$x_j \mapsto y$ の近傍で置き換えた後、
$\bigcup_{i \in I,\ j(i) = j} V_i \to Y$ は有限である。
ゆえに重み付けの定義から
$w(x_j) = \sum_{i \in I,\ j(i) = j} w(\eta_i)$ を得る。したがって
$$
(\textstyle{\int}_f w)(\eta) =
\sum\nolimits_{i \in I} w(\eta_i) \geq
\sum\nolimits_{j(i)\text{ exists}} w(\eta_i) =
\sum\nolimits_j w(x_j) = (\textstyle{\int}_f w)(y)
$$
等号が成り立つことと $y$ が $\bigcap_{i \in I} \Im(V_i \to Y)$ に
含まれることは同値であり、これが示すべきことであった。
\end{proof}

\begin{proposition}
\label{more-morphisms-proposition-asn-weighting}
$f : X \to Y$ をスキームの全射な準有限射とし、
$w : X \to \mathbf{Z}_{> 0}$ を $f$ の正の重み付けとする。
$X$ はアフィン、$Y$ は分離\footnote{$Y$ の対角射がアフィンであれば十分である。}
かつ空でないと仮定する。このとき $Y$ のアフィン層別数は、
$\int_f w$ がとる相異なる値の個数から $1$ を引いた数以下である。
\end{proposition}

\begin{proof}
$Y$ は分離なので射 $X \to Y$ はアフィンである
（Morphisms, Lemma \ref{morphisms-lemma-affine-permanence}）。
Lemma \ref{more-morphisms-lemma-max-int-w} により関数 $\int_f w$ は最大値 $d$ をとる。
$d$ に関する帰納法を用いる。$Y$ の開部分スキーム
$Y_d = \{y \in Y \mid (\int_f w)(y) = d\}$ を考える。
Lemma \ref{more-morphisms-lemma-max-int-finite} により $f^{-1}(Y_d) \to Y_d$ は有限である。
Lemma \ref{more-morphisms-lemma-affineness-of-large-open} により、任意のアフィン開集合
$W \subset Y$ に対して $Y_d \cap W$ はアフィンである
（ここでは $W \times_Y X$ が $X$ 上アフィンなのでアフィンであることを用いる）。
したがって $Y_d \to Y$ はアフィン射である。ゆえに
$f^{-1}(Y_d) = Y_d \times_Y X$ は $X$ 上アフィンなのでアフィンスキームである。
さらに $f^{-1}(Y_d) \to Y_d$ は全射であるから、
Limits, Lemma \ref{limits-lemma-affine} により $Y_d$ はアフィンである。
$Y_d = Y$ ならば $Y$ のアフィン層別数は $0$ であり、
% P05-EN-CH37-0375
命題が成り立つ。そうでなければ、$Y' = Y \setminus Y_d$ を
$Y$ の被約閉部分スキームとみなし、$X' = Y' \times_Y X$ とおく。
$X'$ は $X$ で閉なのでアフィンであり、$Y'$ は $Y$ で閉なので分離である。
射 $f' : X' \to Y'$ は全射であり、$w$ は $f'$ の重み付け $w'$ を誘導する
（Lemma \ref{more-morphisms-lemma-weighting-base-change} を参照せよ）。また
$\int_{f'} w'$ は $\int_f w$ の $Y'$ への制限である。
帰納法により $Y'$ は、長さ $\leq$ $\int_{f'} w'$ の相異なる値の個数から
$1$ を引いた数、となるアフィン層別をもつ。これで証明は完了する。
\end{proof}









\section{完全分解射}
\label{more-morphisms-section-completely-decomposed}

\noindent
Nishnevich は、今日 Nishnevich 位相と呼ばれるものを定義するために、
\'etale 射の完全分解族という概念を研究した。例えば \cite{Nishnevich} を参照せよ。

\begin{definition}
\label{more-morphisms-definition-cd-morphism}
スキームの射 $f : X \to Y$ が {\it 完全分解}である
\footnote{これは標準的でない用語かもしれない。}とは、任意の点 $y \in Y$ に対し、
$f(x) = y$ かつ体拡大 $\kappa(x)/\kappa(y)$ が自明となる点 $x \in X$ が
存在することをいう。終域を固定したスキームの射の族
$\{f_i : X_i \to Y\}_{i \in I}$ が {\it 完全分解}であるとは、
% P05-EN-CH37-0376
$\coprod f_i : \coprod Y_i \to X$ が完全分解であることをいう。
\end{definition}

\noindent
いくつかの基本的な補題から始める。

\begin{lemma}
\label{more-morphisms-lemma-composition-cd}
スキームの二つの完全分解射の合成は完全分解である。
$\{f_i : X_i \to Y\}_{i \in I}$ が完全分解で、各 $i$ に対し完全分解族
$\{X_{ij} \to X_i\}_{j \in J_i}$ があるならば、族
$\{X_{ij} \to Y\}_{i \in I, j \in J_i}$ は完全分解である。
\end{lemma}

\begin{proof}
省略する。
\end{proof}

\begin{lemma}
\label{more-morphisms-lemma-base-change-cd}
スキームの完全分解射の基底変換は完全分解である。
$\{f_i : X_i \to Y\}_{i \in I}$ が完全分解で、$Y' \to Y$ が
スキームの射ならば、$\{X_i \times_Y Y' \to Y'\}_{i \in I}$ は完全分解である。
\end{lemma}

\begin{proof}
$f : X \to Y$ と $g : Y' \to Y$ をスキームの射とする。
$y' \in Y'$ を、$Y$ における像が $y = g(y')$ である点とする。
$x \in X$ が $f(x) = y$ かつ $\kappa(x) = \kappa(y)$ を満たすならば、
$Y'$ で $y'$ へ、$X$ で $x$ へ写り、さらに $\kappa(x') = \kappa(y')$ を満たす
一意な点 $x' \in X' = X \times_Y Y'$ が存在する。
Schemes, Lemma \ref{schemes-lemma-points-fibre-product} を参照せよ。
この事実から補題は容易に従うので、詳細は省略する。
\end{proof}

\begin{lemma}
\label{more-morphisms-lemma-decompose}
\begin{reference}
\cite[Lemma 2.1.2]{EHIK}
\end{reference}
$f : X \to Y$ をスキームの射とする。$f$ は完全分解であり、
$f$ は局所有限表示であり、
$Y$ は準コンパクトかつ準分離であると仮定する。
このとき $n \geq 0$ と射 $Z_i \to Y$, $i = 1, \ldots, n$ で、
次の性質をもつものが存在する：
\begin{enumerate}
\item $\coprod Z_i \to Y$ は全射である。
\item すべての $i$ に対し $Z_i \to Y$ は埋め込みである。
\item すべての $i$ に対し $Z_i \to Y$ は有限表示である。
\item すべての $i$ に対し、基底変換 $X \times_Y Z_i \to Z_i$ は切断をもつ。
\end{enumerate}
\end{lemma}

\begin{proof}
$y \in Y$ とする。仮定により $Y$ 上の射
$\sigma : \Spec(\kappa(y)) \to X$ が存在する。
$\Spec(\kappa(y))$ を、$Z \to Y$ が有限表示な埋め込みであるような
$Y$ 上のアフィンスキーム $Z$ の有向極限として書ける。
実際、アフィン開集合 $y \in \Spec(A) \subset Y$ を選び、$y$ が $A$ の
素イデアル $\mathfrak p$ に対応するとする。このとき $\kappa(\mathfrak p)$ は、
$I \subset \mathfrak p$ が有限生成イデアル、$f \in A$, $f \not \in \mathfrak p$
であるような環 $(A/I)_f$ のフィルター余極限である。
射 $Z = \Spec((A/I)_f) \to Y$ は有限表示な埋め込みである。
$Y$ は準分離なので $Z \to Y$ は準コンパクトである。
Schemes, Lemma \ref{schemes-lemma-quasi-compact-permanence} を参照せよ。
Limits, Proposition
\ref{limits-proposition-characterize-locally-finite-presentation} により、
そのようなある $Z$ に対し、$\kappa(\mathfrak p)$ のスペクトル上で $\sigma$ と
一致する $Y$ 上の射 $\sigma' : Z \to X$ が存在する。
$\sigma'$ は $Y$ 上の射なので、射影 $X \times_Y Z \to Z$ の切断を得る。
% P05-EN-CH37-0377

\medskip\noindent
したがって $Y$ は、$X \times_Y Z \to Z$ が切断をもつ有限表示な埋め込み
$Z \to Y$ の像の合併である。$Z \to Y$ の像は構成可能であり
（Morphisms, Lemma \ref{morphisms-lemma-chevalley}）、$Y$ は構成可能位相で
コンパクトであるから（Properties, Lemma
\ref{properties-lemma-quasi-compact-quasi-separated-spectral} および
Topology, Lemma \ref{topology-lemma-constructible-hausdorff-quasi-compact}）、
そのうち有限個で十分である。
\end{proof}

\begin{lemma}
\label{more-morphisms-lemma-descend-cd}
$S = \lim_{\lambda \in \Lambda} S_\lambda$ を、アフィン遷移射をもつ
スキームの有向系の極限とする。$0 \in \Lambda$ とし、
$f_0 : X_0 \to Y_0$ を $S_0$ 上のスキームの射とする。
$\lambda \geq 0$ に対し $f_\lambda : X_\lambda \to Y_\lambda$ を
$f_0$ の $S_\lambda$ への基底変換とし、$f : X \to Y$ を
$f_0$ の $S$ への基底変換とする。次を仮定する：
\begin{enumerate}
\item $f$ は完全分解である。
\item $Y_0$ は準コンパクトかつ準分離である。
\item $f_0$ は局所有限表示である。
\end{enumerate}
% P05-EN-CH37-0378
このとき $f_\lambda$ が完全分解となる $\lambda \geq 0$ が存在する。
\end{lemma}

\begin{proof}
$Y_0$ は準コンパクトかつ準分離であり、$Y$ は $Y_0$ 上アフィンなので、
これは準コンパクトかつ準分離である。Lemma \ref{more-morphisms-lemma-decompose} のように
$n \geq 0$ と $Z_i \to Y$, $i = 1, \ldots, n$ を選ぶ。
$Z_i$ の選択により存在する $Y$ 上の射を $\sigma_i : Z_i \to X$ と書く。
$0 \in \Lambda$ を大きくした後、$S$ への基底変換が射 $Z_i \to Y$ となる
有限表示な射 $Z_{i, 0} \to Y_0$ が存在すると仮定してよい。
Limits, Lemma \ref{limits-lemma-descend-finite-presentation} を参照せよ。
Limits, Lemma \ref{limits-lemma-descend-immersion} により、必要なら $0$ を
大きくした後、$Z_{i, 0} \to Y_0$ は埋め込みであると仮定してよい。
$\coprod Z_i \to Y$ は全射なので、必要なら $0$ を大きくした後、
$\coprod Z_{i, 0} \to Y_0$ は全射であると仮定してよい。
Limits, Lemma \ref{limits-lemma-descend-surjective} を参照せよ。ここで
$Z_i = \lim_{\lambda \geq 0} Z_{i, \lambda}$
であり、$Z_{i, \lambda} = Y_\lambda \times_{Y_0} Z_{i, 0}$ である。合成
$$
Z_i \xrightarrow{\sigma_i} X \to X_0
$$
を $Y_0$ 上の射とみなす。$f_0$ は局所有限表示なので、
Limits, Proposition
\ref{limits-proposition-characterize-locally-finite-presentation} により、
$Z_i \to Z_{i, \lambda}$ との前合成が表示された矢印となる $Y_0$ 上の射
$\sigma'_{i, \lambda} : Z_{i, \lambda} \to X_0$ が存在するような
$\lambda \geq 0$ を見つけられる。するともちろん $\sigma'_{i, \lambda}$ は
$Y_\lambda$ 上の射 $\sigma_{i, \lambda} : Z_{i, \lambda} \to
X_\lambda = X_0 \times_{Y_0} Y_\lambda$ を定める。
$\coprod Z_{i, \lambda} \to Y_\lambda$ は全射なので、
$X_\lambda \to Y_\lambda$ は完全分解である。
\end{proof}




\section{豊富な可逆加群の族}
\label{more-morphisms-section-families-ample-invertible-modules}

\noindent
Morphisms, Section \ref{morphisms-section-families-ample-invertible-modules} の
議論を続ける。

\begin{lemma}
\label{more-morphisms-lemma-ample-family-ample-relative}
$f : X \to Y$ をスキームの射とする。次を仮定する：
\begin{enumerate}
\item $Y$ は豊富な可逆加群の族をもつ。
\item $X$ 上に $f$-豊富な可逆加群が存在する。
\end{enumerate}
このとき $X$ は豊富な可逆加群の族をもつ。
\end{lemma}

\begin{proof}
$\mathcal{L}$ を $X$ 上の $f$-豊富な可逆加群とする。
特に $f$ は準コンパクトである。Morphisms, Definition
\ref{morphisms-definition-relatively-ample} を参照せよ。
Morphisms, Definition
\ref{morphisms-definition-family-ample-invertible-modules} により $Y$ は
準コンパクトなので、$X$ も準コンパクトである（したがって $X$ 自身は
Morphisms, Definition \ref{morphisms-definition-family-ample-invertible-modules} の
第1条件を満たす）。$x \in X$ の像を $y \in Y$ とする。
% P05-EN-CH37-0379
仮定 (1) により、可逆 $\mathcal{O}_Y$ 加群 $\mathcal{N}$ と切断
$t \in \Gamma(Y, \mathcal{N})$ で、$t$ が消えない軌跡 $Y_t$ が
アフィンとなるものを見つけられる。このとき $\mathcal{L}$ は
$f^{-1}(Y_t) = X_{f^*t}$ 上豊富なので、$(X_{f^*t})_s$ がアフィンかつ
$x$ を含むような切断 $s \in \Gamma(X_{f^*t}, \mathcal{L})$ を見つけられる。
Properties, Lemma \ref{properties-lemma-invert-s-sections} により、
ある $n \geq 0$ に対し積 $(f^*t)^n s$ は切断
$s' \in \Gamma(X, f^*\mathcal{N}^{\otimes n} \otimes \mathcal{L})$ に延長される。
最後に、$f^*\mathcal{N}^{\otimes n + 1} \otimes \mathcal{L}$ の切断
$s'' = f^* ts'$ は $X \setminus X_{f^*t}$ のすべての点で消える。
したがって望みどおり $X_{s''} = (X_{f^*t})_s$ はアフィンである。
\end{proof}

\begin{lemma}
\label{more-morphisms-lemma-resolution-property-goes-up-affine}
$f : X \to Y$ をスキームのアフィン射または準アフィン射とする。
$Y$ が豊富な可逆加群の族をもつならば、$X$ もそうである。
\end{lemma}

\begin{proof}
Morphisms, Lemma \ref{morphisms-lemma-quasi-affine-O-ample} により、これは
Lemma \ref{more-morphisms-lemma-ample-family-ample-relative} の特別な場合である。
\end{proof}






\section{吹き上げと豊富な可逆加群の族}
\label{more-morphisms-section-ample-family-by-blowing-up}

\noindent
\cite{Gross-thesis} の結果を証明する。

\begin{lemma}
\label{more-morphisms-lemma-get-ample-family}
$X$ をスキームとする。$X$ 上の有効 Cartier 因子 $D_1, \ldots, D_m$ と
可逆加群 $\mathcal{L}_1, \ldots, \mathcal{L}_m$ が与えられ、
$\bigcap D_i = \emptyset$ かつ $\mathcal{L}_i|_{X \setminus D_i}$ が
豊富であると仮定する。このとき $X$ は豊富な可逆加群の族をもつ。
\end{lemma}

\begin{proof}
$x \in X$ とする。$x \not \in D_i$ となる添字
$i \in \{1, \ldots, m\}$ を選び、$U_i = X \setminus D_i$ とおく。
% P05-EN-CH37-0380
$\mathcal{L}_i|_{U_i}$ は豊富なので、$s$ が消えない軌跡 $(U_i)_s$ が
アフィンとなるような $n \geq 1$ と切断
$s \in \Gamma(U_i, \mathcal{L}_i^{\otimes n})$ を見つけられる
（Properties, Definition \ref{properties-definition-ample}）。
$U_i$ は標準切断 $1 \in \mathcal{O}_X(D_i)$ が消えない軌跡なので、
Properties, Lemma \ref{properties-lemma-invert-s-sections} により、
$s$ が切断
$$
s' \in \Gamma(X, \mathcal{L}_i^{\otimes n} \otimes_{\mathcal{O}_X}
\mathcal{O}_X(N D_i))
$$
へ延長されるような $N \geq 0$ が存在する。$N$ を $N + 1$ で置き換えると、
$s'$ は $D_i$ のすべての点で消えるので、$X_{s'} = (U_i)_s$ はアフィンである。
これにより $X$ が豊富な可逆加群の族をもつことが証明された。
Morphisms, Definition
\ref{morphisms-definition-family-ample-invertible-modules} を参照せよ。
\end{proof}

\begin{lemma}
\label{more-morphisms-lemma-blow-up-ample-family}
\begin{reference}
\cite[Proposition 1.3.1]{Gross-thesis}
\end{reference}
$X$ を既約成分が有限個である準コンパクトかつ準分離なスキームとする。
準コンパクトな稠密開集合 $U \subset X$ と $U$-許容吹き上げ $X' \to X$ で、
スキーム $X'$ が豊富な可逆加群の族をもつようなものが存在する。
\end{lemma}

\begin{proof}
$\eta_1, \ldots, \eta_n \in X$ を $X$ の既約成分の総称点とする。
Properties, Lemma \ref{properties-lemma-point-and-maximal-points-affine} と
$X$ の準コンパクト性により、各 $U_i$ が $\eta_1, \ldots, \eta_n$ を含むような
有限アフィン開被覆 $X = U_1 \cup \ldots \cup U_m$ を見つけられる。
特に準コンパクト開部分集合 $U = U_1 \cap \ldots \cap U_m$ は $X$ で稠密である。
$Z_i = V(\mathcal{I}_i)$ として $U_i = X \setminus Z_i$ となる有限型準連接
イデアル層 $\mathcal{I}_i \subset \mathcal{O}_X$ をとる。
Properties, Lemma \ref{properties-lemma-quasi-coherent-finite-type-ideals} を参照せよ。
$$
f : X' \longrightarrow X
$$
をイデアル層 $\mathcal{I} = \mathcal{I}_1 \cdots \mathcal{I}_m$ における
$X$ の吹き上げとする。$V(\mathcal{I})$ は集合論的に $U$ と交わらない
$Z_i$ の合併なので、$f$ は $U$-許容吹き上げである。
また $f$ は射影的で、$\mathcal{O}_{X'}(1)$ は $f$-相対的に豊富である。
Divisors, Lemma \ref{divisors-lemma-blowing-up-projective} を参照せよ。
Divisors, Lemma \ref{divisors-lemma-blowing-up-two-ideals} により、各 $i$ に対し
% P05-EN-CH37-0381
射 $f'$ は $X' \to X'_i \to X$ と分解する。ここで $X'_i \to X$ は
$\mathcal{I}_i$ における吹き上げで、$X' \to X'_i$ は別の吹き上げ
（すなわち $\mathcal{I}_i$ を除いたイデアル $\mathcal{I}_j$ の積の
引き戻しにおける吹き上げ）である。これから
$D_i = f^{-1}(Z_i) \subset X'$ は有効 Cartier 因子である。
Divisors, Lemmas \ref{divisors-lemma-blow-up-pullback-effective-Cartier} および
\ref{divisors-lemma-blowing-up-gives-effective-Cartier-divisor} を参照せよ。
$X' \setminus D_i = f^{-1}(U_i)$ である。$\mathcal{O}_{X'}(1)$ は $f$-豊富なので、
$X' \setminus D_i$ への $\mathcal{O}_{X'}(1)$ の制限は豊富である。
Lemma \ref{more-morphisms-lemma-get-ample-family} により、$X'$ は豊富な可逆加群の族をもつ。
\end{proof}

\begin{proposition}
\label{more-morphisms-proposition-envelope-with-resolution-property}
$X$ を準コンパクトかつ準分離なスキームとする。有限表示、固有、かつ完全分解
（Definition \ref{more-morphisms-definition-cd-morphism}）である射 $f : Y \to X$ で、
スキーム $Y$ が豊富な可逆加群の族をもつようなものが存在する。
\end{proposition}

\begin{proof}
Limits, Proposition \ref{limits-proposition-approximate} により、$X_0$ が
$\mathbf{Z}$ 上有限型なスキームであるようなアフィン射 $X \to X_0$ が存在する。
以下で望ましい性質をすべてもつ射 $Y_0 \to X_0$ を構成する。
$Y = X \times_{X_0} Y_0$ とおき、$f$ を射影 $f : Y \to X$ とすれば結論を得る。
実際、$f$ は有限表示であり
（Morphisms, Lemma \ref{morphisms-lemma-base-change-finite-presentation}）、
$f$ は固有であり（Morphisms, Lemma \ref{morphisms-lemma-base-change-proper}）、
$f$ は完全分解である（Lemma \ref{more-morphisms-lemma-base-change-cd}）。最後に $Y \to Y_0$ は
$X \to X_0$ の基底変換としてアフィンなので、
Lemma \ref{more-morphisms-lemma-resolution-property-goes-up-affine} により $Y$ は
豊富な可逆加群の族をもつ。これで次の段落で扱う場合に帰着した。

\medskip\noindent
$X$ が $\mathbf{Z}$ 上有限型であると仮定する。特に $\dim(X) < \infty$ である。
$\dim(X)$ に関する帰納法を用いる。$\dim(X) = 0$ ならば、$X$ はアフィンであり、
% P05-EN-CH37-0382
分解性質をもつ。一般の場合、$X'$ が豊富な可逆加群の族をもつような
稠密開集合 $U \subset X$ と $U$-許容吹き上げ $X' \to X$ が存在する。
Lemma \ref{more-morphisms-lemma-blow-up-ample-family} を参照せよ。
$f : X' \to X$ は $U$ 上同型なので、$U$ の各点は同じ剰余体をもつ
$X'$ の点へ持ち上がる。$Z = X \setminus U$ に誘導される被約スキーム構造を
入れる。$U$ は $X$ で稠密なので（上を参照）、$\dim(Z) < \dim(X)$ である。
帰納法により、$W$ が豊富な可逆加群の族をもつような固有完全分解射
$W \to Z$ を得る。したがって $Y = X' \amalg W \to X$ が望む射である。
\end{proof}




\section{閉埋め込みの拡大判定法}
\label{more-morphisms-section-extensive-criterion-closed-immersions}

\noindent
本節では、スキームの射が閉埋め込みであるための判定法を与える。

\begin{lemma}
\label{more-morphisms-lemma-affine-extensive-criterion}
アフィンスキームの射 $f : X \to Y$ が閉埋め込みであるための必要十分条件は、
任意の単射な環準同型 $A \to B$ と可換正方形
$$
\xymatrix{
\Spec(B) \ar[d] \ar[r] & X \ar[d]^f \\
\Spec(A) \ar[r] \ar@{..>}[ur] & Y
}
$$
に対し、二つの三角形を可換にする持ち上げ $\Spec(A) \to X$ が存在することである。
\end{lemma}

\begin{proof}
射 $f$ が環準同型 $\phi : R \to S$ により与えられるとする。
$f$ が閉埋め込みであることと $\phi$ が全射であることは同値である。

\medskip\noindent
まず $\phi$ が全射であると仮定する。$\psi : A \to B$ を単射な環準同型とし、
可換図式
$$
\xymatrix{
R \ar[r]^\alpha \ar[d]^\phi & A \ar[d]^\psi \\
S \ar[r]^\beta \ar@{..>}[ur] & B
}
$$
が与えられたとする。$\phi(r) = s$ となる $r \in R$ を用いて
$s \mapsto \alpha(r)$ とおくことにより、持ち上げ $S \to A$ を定める。
$\psi$ は単射で正方形は可換なので、これは良定義である。
環のスペクトルをとることは可換環とアフィンスキームの間の反同値を定めるから、
$f$ は望む持ち上げ性質をもつ。

\medskip\noindent
次に、$\phi$ がすべての単射な環準同型 $\psi: A \to B$ に対して
持ち上げをもつと仮定する。$\phi(R)$ は $S$ の部分環なので、可換正方形
$$
\xymatrix{
R \ar[r] \ar[d]^\phi  & \phi(R) \ar[d] \\
S \ar@{=}[r] \ar@{..>}[ur] & S
}
$$
を得て、ここには持ち上げ $S \to \phi(R)$ が存在する。
したがって包含 $\phi(R) \to S$ は同型でなければならない。
ゆえに $\phi$ は全射であり、証明が終わる。
\end{proof}

\begin{lemma}
\label{more-morphisms-lemma-scheme-affine-if-aff-is-split-mono}
$X$ をスキームとする。Schemes, Lemma
\ref{schemes-lemma-morphism-into-affine} の標準射
$X \to \Spec(\Gamma(X, \mathcal{O}_X))$ がレトラクションをもつならば、
$X$ はアフィンスキームである。
\end{lemma}

\begin{proof}
$S = \Spec(\Gamma(X, \mathcal{O}_X))$ と書き、$f : X \to S$ を補題の射とする。
$s : S \to X$ をレトラクションとすれば、$\text{id}_X = sf$ である。
すると $f s f = \text{id}_S f$ である。$f$ は同型
$\Gamma(S, \mathcal{O}_S) \to \Gamma(X, \mathcal{O}_X)$ を誘導するので、
$fs$ と $\text{id}_S$ は $\Gamma(S, \mathcal{O}_S)$ 上同じ写像を誘導する。
$S$ はアフィンだから $f s = \text{id}_S$ である。
ゆえに $f$ は同型であり、示すべきとおり $X$ はアフィンスキームである。
\end{proof}

\begin{lemma}
\label{more-morphisms-lemma-aff-is-injective-on-sheaves}
$X$ をスキームとし、$f : X \to S = \Spec(\Gamma(X, \mathcal{O}_X))$ を
Schemes, Lemma \ref{schemes-lemma-morphism-into-affine} の標準射とする。
$f^\sharp : \mathcal{O}_S \to f_*\mathcal{O}_X$ の核に含まれる最大の
準連接 $\mathcal{O}_S$ 加群は零である。$X$ が準コンパクトならば、
$f^\sharp$ は単射である。特に $X$ が準コンパクトならば、$f$ は支配射である。
\end{lemma}

\begin{proof}
$f^\sharp$ の核に含まれる最大の準連接 $\mathcal{O}_S$ 加群に対応する
部分加群を $M \subset \Gamma(S, \mathcal{O}_S)$ とする。
任意の元 $a \in M$ は $f^\sharp$ により零へ写る。一方、$f^\sharp(a)$ は
$$
\Gamma(S, f_*\mathcal{O}_X) = \Gamma(X, \mathcal{O}_X) =
\Gamma(S, \mathcal{O}_S)
$$
の元で $a$ 自身に対応する。したがって $a = 0$ である。
ゆえに $M = 0$ であり、第1の主張が証明された。
これは射 $f : X \to S$ がスキーム論的全射であることと同値である。

\medskip\noindent
$X$ が準コンパクトならば、Morphisms, Lemma
\ref{morphisms-lemma-quasi-compact-scheme-theoretic-image} により
$\Ker(f^\sharp)$ は準連接である。したがって $\Ker(f^\sharp) = 0$ で、
$f^\sharp$ は単射である。この場合、Morphisms, Lemma
\ref{morphisms-lemma-quasi-compact-scheme-theoretic-image} の part (4) により
$f$ は支配射である。
\end{proof}

\begin{lemma}
\label{more-morphisms-lemma-extensive-criterion}
$f: X \to Y$ をスキームの準コンパクト射とする。このとき $f$ が
閉埋め込みであるための必要十分条件は、任意の単射な環準同型 $A \to B$ と
可換正方形
$$
\xymatrix{
\Spec(B) \ar[d] \ar[r] & X \ar[d]^f \\
\Spec(A) \ar[r] \ar@{..>}[ur] & Y
}
$$
に対し、図式を可換にする持ち上げ $\Spec A \to X$ が存在することである。
\end{lemma}

\begin{proof}
$f$ が閉埋め込みであると仮定する。$A \to B$ を単射な環準同型とし、
可換正方形
$$
\xymatrix{
\Spec(B) \ar[d] \ar[r]  & X \ar[d]^f \\
\Spec(A) \ar[r] \ar@{..>}[ur] & Y
}
$$
を考える。このとき $\Spec(A) \times_Y X \to \Spec(A)$ は閉埋め込みなので、
イデアル $I \subset A$ と可換図式
$$
\xymatrix{
\Spec(B) \ar[d] \ar[r] & \Spec(A/I) \ar[r] \ar[d] & X \ar[d]^f \\
\Spec(A) \ar[r] \ar@{..>}[ur] & \Spec(A) \ar[r] & Y
}
$$
を得る。Lemma \ref{more-morphisms-lemma-affine-extensive-criterion} により持ち上げを得る。

\medskip\noindent
$f$ が補題で述べた持ち上げ性質をもつと仮定する。
$f$ が閉埋め込みであることの証明は $Y$ 上局所的なので、
$Y$ はアフィンであると仮定してよく、そう仮定する。
特に $Y$ は準コンパクトであり、したがって $X$ も準コンパクトである。
ゆえに有限アフィン開被覆 $X = U_1 \cup \ldots \cup U_n$ が存在する。射
$$
\pi : U = \coprod U_i \longrightarrow X
$$
の始域はアフィンであり、誘導される環準同型
$\Gamma(X, \mathcal{O}_X) \to \Gamma(U, \mathcal{O}_U)$ は単射である。
仮定により図式
$$
\xymatrix{
U \ar[r]^\pi \ar[d] & X \ar[d]^f \\
\Spec(\Gamma(X, \mathcal{O}_X)) \ar[r]^-{f'} \ar@{..>}[ur]^h & Y
}
$$
に持ち上げが存在する。ここで $f'$ は環準同型
$\Gamma(Y, \mathcal{O}_Y) \to \Gamma(X, \mathcal{O}_X)$ に対応する
アフィンスキームの射である。$\pi$ がエピ射であることから、射 $h$ は
標準射 $X \to \Spec(\Gamma(X, \mathcal{O}_X))$ のレトラクションである。
詳細は省略する。したがって Lemma
\ref{more-morphisms-lemma-scheme-affine-if-aff-is-split-mono} により $X$ はアフィンである。
Lemma \ref{more-morphisms-lemma-affine-extensive-criterion} により $f$ は閉埋め込みである。
\end{proof}
